% 本文件由 LaTeX 排版 Agent 根据有效 translation.json 自动生成。
% author=Origen; work_id=0416; title=驳斥塞尔苏斯

\chapter{驳斥塞尔苏斯}

% structure: p0001 source=h1 target=omitted reason=page-title-already-rendered-by-parent

第一卷\newline{}第二卷\newline{}第三卷\newline{}第四卷\newline{}第五卷\newline{}第六卷\newline{}第七卷\newline{}第八卷\par

\SourceRecord{Translated by Frederick Crombie.；Ante-Nicene Fathers；Vol. 4.；Edited by Alexander Roberts, James Donaldson, and A. Cleveland Coxe.；Buffalo, NY: Christian Literature Publishing Co.,；1885.；<http://www.newadvent.org/fathers/0416.htm>.}

% structure: p0001 source=h1 target=section reason=page-title

\section{《驳斥凯尔苏斯》卷一}

% structure: p0002 source=h2 target=subsection reason=relative-heading-rank; base=h2

\subsection{前言}

1. 当假证人针对我们的主和救主耶稣基督作证时，祂保持了沉默；当不实的指控加于祂时，祂没有回答，深信祂在犹太人中的全部生活和行为，比任何对假证的答复或对指控的正式辩护，更能驳斥他们。虔诚的安博罗削，我不知道你为什么希望我写一篇答复，回应凯尔苏斯在其论著中对基督徒提出的虚假指控，以及他对教会信德的诘难；好像事实本身没有提供明显的驳斥，而教义本身比任何著作都更好的答复，因为它既消除了虚假陈述，又不让指控有任何可信性或有效性。关于我们的主在面对假证时保持沉默，目前只需引用玛窦的记载就够了，因为马尔谷的记载与此相同。玛窦的话如下：\BibleQuote{大司祭和公议会设法寻找假证，反对耶稣，为要处死祂，却没有找到，虽然有许多假证人前来。最后有两个人上前来说：\InnerQuote{这人曾说：我能拆毁天主的圣殿，三天内把它重建起来。}大司祭站起来对祂说：\InnerQuote{这些人作证反对祢的事，祢什么也不回答吗？}耶稣却默不作声。}关于祂在被冤枉时没有回答，记载如下：\BibleQuote{耶稣站在总督面前；总督问祂说：\InnerQuote{祢是犹太人的君王吗？}耶稣对他说：\InnerQuote{你说的是。}当祂被司祭长和长老控告时，祂什么也不回答。于是比拉多对祂说：\InnerQuote{祢没听见他们作证反对祢这么多事吗？}祂连一句话也不回答，以致总督大为惊奇。}\par

2. 确实，连普通智慧的人也感到惊奇：一个被假证指控和攻击的人，本来能够为自己辩护，表明自己对这些指控毫无罪责，并且可以列举自己一生中值得称赞的事迹和凭天主能力所行的奇迹，使法官有机会作出更光荣的判决，但祂却没有这样做，而是不屑于此，并以高贵的天性藐视控告者。如果祂提出辩护，法官会毫不犹豫地释放祂，这从关于比拉多的记载中可以清楚看出，他说：\BibleQuote{你们愿意我释放哪一个给你们？是巴辣巴，还是那称为基督的耶稣？}以及圣经补充的：\BibleQuote{因为他知道他们是出于嫉妒才把祂交来的。}然而，耶稣总是受到假证人的攻击，只要邪恶存于世上，祂就常遭指控。但即使现在，祂在这些事上仍然保持沉默，不作出声的回答，而是将祂的辩护置于祂真正门徒的生命中，这生命是卓越的见证，超越一切假证，并驳斥和推翻所有无根据的指控和控告。\par

3. 我敢说，你们要我写的这篇\Quote{辩护}，多少会削弱那基于事实的辩护，以及耶稣的能力——这能力对那些并非完全没有感知的人是显明的。尽管如此，为了不显得不情愿承担你所托付的任务，我们已尽力针对凯尔苏斯提出的每一点陈述，提出我们认为适于驳斥它们的答复，尽管他的论据无力动摇任何（真正）信者的信德。但愿不会有人，在分享了基督耶稣所显示的如此伟大的爱主之情后，竟被凯尔苏斯或任何类似之人的论据所动摇。因为保禄在列举那些通常使人远离基督之爱和基督耶稣内天主之爱的无数原因时（所有这些原因，他心中的爱都胜过它们），并没有把论据列为分离的理由。请注意他首先说：\BibleQuote{谁能使我们与基督的爱分离呢？是困苦吗？是窘迫吗？是迫害吗？是饥饿吗？是赤身露体吗？是危险吗？是刀剑吗？（正如经上所记：\InnerQuote{为了祢，我们终日被处死；我们被看作待宰的羊。}）然而，靠那爱我们的主，我们在这一切事情上大获全胜。}其次，当列出另一系列容易使那些信仰不坚定的人分离的原因时，他说：\BibleQuote{因为我深信，无论是死亡，是生命，是天使，是掌权者，是现存的，是未来的，是有能者，是崇高，是深远，或是任何受造之物，都不能使我们与天主的爱分离，这爱是在我们的主基督耶稣内的。}\par

如今，我们确实应当因患难或\Person{保禄}所列举的其他事物不能使我们与基督分离而欢欣鼓舞，但\Person{保禄}、其他宗徒以及任何类似他们的人却不应有此感觉，因为他们早已超越了这些，当他们说：\BibleQuote{借着那爱我们的主，我们在这一切事上\Emph{大}获全胜}，这比仅仅说\Quote{获胜}更为有力。但如果宗徒们因不与被爱我们的主耶稣基督内的天主的爱分离而感到欢欣鼓舞是合宜的，那么他们确实会有此感觉，因为无论是死亡、生命、天使、掌权者，还是其他任何事物，都不能使他们与主耶稣基督内的天主的爱分离。因此，我并不为那个因\Person{凯尔苏斯}——他早已不再与世人共处，早已离世——或任何貌似合理的论证而信心动摇的信徒感到庆幸。因为我不知该将那些需要借助书籍中的论证来回应\Person{凯尔苏斯}对基督徒的指控、以防信心动摇并加以坚固的人置于何等地位。然而，在众多被视为信徒的人中，或许会有一些人因\Person{凯尔苏斯}的著作而信心动摇甚至被推翻，但若能通过这样的答复来驳斥他的言论并彰显真理，他们便可得以保全。因此，我们认为理应遵从你的嘱托，对你寄来的论著作出答复——虽然我认为任何在哲学上稍有造诣的人都不会允许它被称为\WorkTitle{真道}，正如\Person{凯尔苏斯}所命名的。\par

\Person{保禄}观察到希腊哲学中有某些并非微不足道的东西，它们在众人眼中看似合理，却以虚假冒充真理，便就此说道：\BibleQuote{你们要谨慎，免得有人用哲学和虚空的诡诈，按人的传统，按世俗的基本原理，而不按基督，把你们掳去。}他看出世俗智慧的话语中有一种表面的伟大，便称哲学家的话是\BibleQuote{按世俗的基本原理}。然而，有理智的人不会说\Person{凯尔苏斯}的话是\BibleQuote{按世俗的基本原理}。那些含有某种欺骗成分的话语，宗徒称之为\BibleQuote{虚空的诡诈}，大概是区别于并非\Quote{虚空}的诡诈。\Person{耶肋米亚}先知注意到这一点，大胆对天主说：\BibleQuote{上主，祢欺骗了我，我被祢欺骗了；祢比我强，且胜过了我。}但在我看来，\Person{凯尔苏斯}的话中毫无欺骗成分，甚至没有那种\Quote{虚空的}诡诈——即那些创立哲学派别且在此追求上禀赋非凡之人的话语中所包含的诡诈。正如没有人会说几何证明中的任何普通错误是为了欺骗，或为了在类似问题上练习而如此描述；同样，那些被称为\BibleQuote{虚空的诡诈}、\BibleQuote{人的传统}和\BibleQuote{按世俗的基本原理}的观点，必须与那些创立哲学派别之人的见解有某种相似之处（如果这些称谓要恰当地适用于它们的话）。\par

在完成这部著作直至\Person{凯尔苏斯}引入犹太人与\Person{耶稣}辩论之处后，我决定将这篇前言置于卷首，以便阅读我们答复\Person{凯尔苏斯}的读者能首先看到它，并明白此书并非为那些信心坚固的信徒而作，而是为那些或对基督徒信仰全然不知、或如宗徒所称\BibleQuote{在信德上软弱的人}而作；关于后者，他说：\BibleQuote{你们要收纳在信德上软弱的人。}这篇前言必须作为我的道歉，因为我以某种方式开始答复\Person{凯尔苏斯}，却又以另一种方式继续。我最初的想法是指出他的主要反对意见，然后简要给出答复，随后将整个论述系统化。但后来，实际情况本身提示我要节省时间，满足于开头已陈述的内容，并在后续部分尽我所能与\Person{凯尔苏斯}的指控紧密交锋。因此，我恳请读者对前言之后靠近卷首的部分予以宽容。如果后续的有力论证不能打动你，那么同样请求对这些部分也予以宽容，并将你（若仍希望对\Person{凯尔苏斯}的反对意见进行论证性解决）引向那些比我更有智慧、能通过言语和论著推翻他对我们指控的人。但更好的是，即使遇到\Person{凯尔苏斯}的著作也完全不需要答复的人，他蔑视其全部内容，因为这些内容被每个在基督内的信徒通过内在的圣神所轻视，且理当如此。\par

% structure: p0009 source=h2 target=subsection reason=relative-heading-rank; base=h2

\subsection{第一章}

\Person{塞尔苏斯}为了诋毁基督教而提出的第一个论点是，基督徒彼此之间秘密结社，违反法律，他说：\Quote{有的团体是公开的，并且是符合法律的；另一些则是秘密的，并且是违背法律而维持的。}他希望诋毁基督徒所谓的\Quote{爱筵}，仿佛它们起源于共同的危险，并且比任何誓言更有约束力。既然他喋喋不休地谈论公法，声称基督徒的团体违反公法，我们当回答：若有人置身于\Person{斯基泰人}中，他们的法律是不敬神的，而他无法逃走，被迫生活在他们中间，那么此人为了真理的律法（\Person{斯基泰人}会视之为邪恶）之故，完全有理由与志同道合者一起，违反他们的法律而结社；同样，如果以真理为准，外邦人关于偶像以及无神的多神信仰的法律，就是\Quote{斯基泰的}法律，或者甚至比任何此类法律更不敬神。因此，为了真理的缘故而组织反对现有法律的团体，并不是不合理的。因为，正如那些为了杀死篡夺国家自由的暴君而组织秘密团体的人做得对一样，基督徒在被那称为魔鬼的以及谎言所暴虐统治时，也违反魔鬼的法律而结成联盟，以对抗他的势力，并为了拯救那些他们可能说服其摆脱这个犹如\Quote{斯基泰}且专制统治的人。\par

% structure: p0011 source=h2 target=subsection reason=relative-heading-rank; base=h2

\subsection{第二章}

\Person{塞尔苏斯}接着说，基督教所依据的教义体系，即犹太教，其起源是野蛮的。他表面上公允，并没有因为基督教起源于野蛮人而责备它，反而称赞野蛮人在发现这些教义方面的能力。然而，他接着说，希腊人在评判、确立和实行野蛮民族的发现方面比任何人都更擅长。以下是我们对他的指控的回答，以及我们对基督教所含真理的辩护：如果一个人从研究希腊的意见和习俗转向福音，他不仅会断定福音的教义是真的，而且会通过实践确立它们的真理，并从希腊的角度弥补它们在论证上似乎欠缺的地方，从而证实基督教的真理。此外，我们必须说，福音有其自身的论证，比任何由希腊辩证法建立的论证更神圣。宗徒称这更神圣的方法为\Quote{圣神与德能的表现}：关于\Quote{圣神}，是指预言，这些预言足以使任何读它们的人产生信德，尤其是在那些关于基督的事情上；关于\Quote{德能}，是指神迹与奇事，我们必须相信这些神迹与奇事已经发生过，这不仅基于许多其他理由，也基于这一点：那些按照福音规诫生活的人中，仍然保留着这些神迹与奇事的痕迹。\par

% structure: p0013 source=h2 target=subsection reason=relative-heading-rank; base=h2

\subsection{第三章}

此后，\Person{塞尔苏斯}继续谈论基督徒秘密教授并实践他们所喜爱的教义，并说他们这样做有一定的目的，因为他们逃脱了迫在眉睫的死刑惩罚，他将他们的危险与\Person{苏格拉底}等人为哲学所遭遇的危险相比；这里他本也可以提到\Person{毕达哥拉斯}和其他哲学家。但我们对此的回答是，在\Person{苏格拉底}的事例中，\Person{雅典人}随后立即后悔了；他们心里对他没有留下丝毫苦毒，\Person{毕达哥拉斯}的历史也是如此。后者的追随者确实在意大利称为\Place{大希腊}的那个地区建立学派，持续了相当长的时间；但在基督徒的事例中，\Place{罗马元老院}、当时的君主、军队、民众以及那些已成为信仰归信者的亲属，都向他们的教义开战，并且原本会阻止其进展，以如此强大的联盟来压倒它；然而，幸亏有天主的帮助，它逃脱了危险，凌驾于其上，最终击败了整个世界合谋反对它的阴谋。\par

% structure: p0015 source=h2 target=subsection reason=relative-heading-rank; base=h2

\subsection{第四章}

我们也要注意他是如何试图诋毁我们的道德体系，声称它只是我们与其他哲学家所共有的，并非什么可敬的或新的教导分支。对此我们必须说，除非所有人在本性上都自然而然地有健全的道德观念印在心上，否则那给自身招致天主公义审判者的罪人受惩罚的教义就会被排除。因此，同一位天主将祂借着先知和救主所教导的那些真理播种在所有人的心中，以便在审判之日，每个人都无可推诿，因他们心中刻有\BibleQuote{律法的要求刻在他心上}——圣经在\Person{希腊人}视为神话的故事中隐约暗示了这一真理，其中描述天主亲自用手指写下诫命，交给\Person{梅瑟}，而拜牛犊者的邪恶使\Person{梅瑟}将石版摔碎，仿佛邪恶的洪流将它们冲走了。但\Person{梅瑟}再次凿出石版，天主第二次写下诫命，交给了他；先知之言，在第一次过犯之后，预备灵魂接受天主第二次的书写。\par

% structure: p0017 source=h2 target=subsection reason=relative-heading-rank; base=h2

\subsection{第五章}

论及敬拜偶像之规条乃基督教所特有，\Person{克耳苏}肯定其正确，称基督徒不认为人手所造者为神，因依正理，由最卑劣邪恶之工匠所造、且多由恶人供奉之像，不得视为神。然其下文欲表明此乃共通之见，非基督教首创，遂引\Person{赫拉克利特}之言：\Quote{凡近无生命之像如敬神者，其行与入室与屋交谈者无异。}对此，我等当言：道德原则之观念已铭刻人心，不仅\Person{赫拉克利特}，任何希腊人或蛮族，皆可经思辨而得同一结论；彼又引\Person{希罗多德}为证，谓波斯人亦持此见。我等可增引\Place{克提翁}之\Person{芝诺}，其\WorkTitle{政制}有云：\Quote{无需建殿，因工匠与卑贱者之作，不得视为神圣、贵重或圣洁。}可见，此见（如其他见）亦由天主之指刻于人心，使人知所需职责。\par

% structure: p0019 source=h2 target=subsection reason=relative-heading-rank; base=h2

\subsection{第6章}

此后，不知何故，\Person{克耳苏}声言基督徒藉某些魔鬼之名与咒语而显奇能，意或指以咒语驱魔者之做法。此显然诋毁福音。盖基督徒胜魔非凭咒语，乃因\Person{耶稣}之名，并述其事迹；常诵此等事迹，尤其是以纯正真诚之心诵之者，屡能驱魔。\Person{耶稣}之名对邪魔之权能，乃至恶人呼之亦有效验，此\Person{耶稣}曾亲言：\Quote{当那日，必有许多人对我说：因祢的名，我们驱了魔鬼，行了许多异能。}\Person{克耳苏}之忽略此言，出于恶意抑或无知，吾不得知。其复指控救主本身，谓其行奇事乃藉巫术；且预见他人将得同样知识而行同样事，却矜言藉天主权能，故禁其入国。彼责曰：若彼等被逐为义，而救主自行同样之事，则为恶人；若救主行此非恶，则行同样事者亦非恶。纵不能显明\Person{耶稣}以何权能行此奇迹，然基督徒不用符咒，唯用\Person{耶稣}之名及圣经中某些信靠之言，此乃显然。\par

% structure: p0021 source=h2 target=subsection reason=relative-heading-rank; base=h2

\subsection{第7章}

再者，因彼屡称基督教教义为秘密体系，我等亦当驳之。盖世人知晓基督徒所传扬者，远胜于哲学名家所珍视之见解。谁不知\Person{耶稣}为童贞女所生、被钉十字架、其复活为多人所信之信条，以及将来有公审判，恶人按其所应受罚，义人得适当赏报？然复活之奥秘，因不被理解，竟成不信者之笑柄。在此情形下，称基督教教义为秘密体系，实属荒谬。然有某些教义，在外传教义讲授之后，并不公开于众人，此非基督教独有，亦见于哲学体系，其中某些真理为外传，某些为内传。\Person{毕达哥拉斯}之某些听众，满足于其\OriginalTerm{他本人之言}{ipse dixit}；而另一些则被秘密传授那些不宜传与世俗及未预备之耳的教义。此外，希腊及蛮族各处所有之奥秘仪式，虽行于秘密，均未受诋毁。故彼欲诋毁基督教之秘密教义，实属徒劳，因其未正确理解基督教之本质。\par

% structure: p0023 source=h2 target=subsection reason=relative-heading-rank; base=h2

\subsection{第8章}

\Person{塞而苏斯}确实以某种雄辩的口才，为那些因死亡而见证基督信仰真理的人辩护，说了以下话：\Quote{我并不主张，如果一个人采纳了良好教义的体系，并因此面临来自人类的危险，他就应当背教、假装背教或公开否认自己的观点。} 他又谴责那些持基督信仰观点却假装不持或否认它们的人，说：\Quote{持有某种观点的人不应假装放弃原意或公开否认它。} 在此，\Person{塞而苏斯}必须被指证为自相矛盾。因为从他的其他著作中可以确知他是\OriginalTerm{伊壁鸠鲁派}{Epicurean}；但在这里，因为他认为不公开承认\Person{伊壁鸠鲁}的观点能更有效地攻击基督信仰，他便假装在人身上有一个比其世俗本性更好的东西，这东西与天主相近，并说：\Quote{那些拥有这个要素（即灵魂）且它处于健康状态的人，总是在寻求他们同类的本性，即天主，总是渴望听到关于祂的事并常常记起祂。} 现在请注意他性格的虚伪！他稍早前说过：\Quote{采纳了良好教义体系的人，即使因此面临来自人类的危险，也不应否认它、假装否认它，更不应公开否认它。} 现在他却陷于各种矛盾之中。因为他知道，如果承认自己是伊壁鸠鲁派，他在指责那些多少引入了上智安排道理并将一位神置于世界之上的人时，将不会获得任何信任。而且我们听说有两位名叫\Person{塞而苏斯}的人，两人都是伊壁鸠鲁派；较早的一位生活在\Person{尼禄}时代，而这一位则生活在\Person{哈德良}时代及以后。\par

% structure: p0025 source=h2 target=subsection reason=relative-heading-rank; base=h2

\subsection{第九章}

他接着建议，在采纳意见时应遵循理性和理性的向导，因为不遵循此道而同意意见的人很容易受骗。他将轻率的信徒比作\OriginalTerm{墨特拉巨泰}{Metragyrtæ}、占卜者、\OriginalTerm{密特拉}{Mithræ}、\OriginalTerm{萨巴迪安}{Sabbadians}，以及任何其他可能遇到的东西，并比作\OriginalTerm{赫卡忒}{Hecate}或其他任何鬼神的幻影。因为正如这类人中常可发现恶人，他们利用那些易受欺骗之人的无知，将其引向任何他们想要的地方，他说，在基督徒中也是如此。而且他断言，某些既不愿为信仰提出理由也不愿接受理由的人，不断重复：\Quote{不要探究，只要相信！} 以及\Quote{你的信德将拯救你！} 他还声称，这些人还说：\Quote{今世的智慧是坏的，而愚蠢是好的！} 对此我们必须回答：如果所有人都能放下俗务、专心哲学，那么任何人都不应采用其他方法，而应只采用这一种。因为在基督信仰体系中也同样会发现——毫不傲慢地说——对信仰条目的探究、对先知书中隐晦言论的解释、对福音书中的比喻以及无数其他以象征意义叙述或施行之事的解释（正如在其他体系中一样）。但由于所提及的途径是不可能的，部分由于生活的必需，部分由于人的软弱（因为只有极少数人认真致力于学习），那么为了帮助大众，还有什么比耶稣传给外邦人的方法更好的呢？我们不妨询问，对于大批信徒——他们已洗去从前沉溺的邪恶泥沼——是毫无理由地相信，并因而通过相信人因罪受罚、因善行受赏而得到改造和行为改善更好呢，还是不应允许他们仅凭信德皈依，而要等到能够彻底审视（必要的）理由时才行动？因为显然，按照这样的计划，几乎所有的人（极少数例外）都不会获得他们通过单纯的信德所获得的行为改善，而是会继续留在邪恶的生活中。现在，无论可以提出其他什么证据来证明基督信仰的仁爱计划并非没有天主的干预而被引入人间，这一点也必须补充。因为一个虔诚的人不会相信，即使是身体的医生（他使病人恢复健康）能够没有天主的许可而定居在任何城市或国家，因为没有天主的帮助，人就不会有好事发生。如果那位治愈了许多人的\Emph{身体}、或使他们恢复健康的人，没有天主的帮助就无法完成他的治疗，那么那位治愈了许多人的\Emph{灵魂}、使他们转向美德、改善了他们的本性、使他们依附于万有之上的天主、教导他们使一切行为都指向祂的喜悦、并避开所有令祂不悦的事——甚至他们最小的言语、行为，或内心的念头——的那位，更是何等超越呢？\par

% structure: p0027 source=h2 target=subsection reason=relative-heading-rank; base=h2

\subsection{第十章}

其次，既然我们的对手不断重复那些关于信仰的言论，我们必须说，我们认为信仰对大众是有益的，我们承认我们教导那些无法放下一切其他事务、专心于研究论证的人在没有理由的情况下相信。而我们的对手，尽管他们不承认，实际上也这样做。因为有谁在投身哲学研究、加入某个学派时——无论是偶然为之，还是因为有一位该学派的老师指导——，会出于其他任何原因，而不是因为他\Emph{相信}自己的学派优于其他任何学派呢？因为他并不等待聆听所有其他哲学家和所有不同学派的论证，以及谴责一个体系、支持另一个体系的理由，就由此选择成为，例如斯多葛派、柏拉图派、逍遥派、伊壁鸠鲁派，或其他某个学派的追随者，这样他便被一种非理性的冲动（尽管他们不承认）带向实践，比如说，斯多葛主义，而忽视其他学派；或者轻视柏拉图主义，认为它比其他学派更具谦卑；或者轻视逍遥派，认为它更人性化，并且比其他体系更公平地承认人类生活中的祝福。还有一些人，起初因看到世上恶人与善人的遭遇而对天主眷顾的教义感到震惊，便鲁莽地断定根本没有天主眷顾，从而采纳了伊壁鸠鲁和凯尔苏斯的观点。\par

% structure: p0029 source=h2 target=subsection reason=relative-heading-rank; base=h2

\subsection{第十一章}

既然如此，正如理性教导的，我们必须信赖希腊人或蛮族人中那些曾开创学派的人中的某一位，那么我们为什么不信赖那位统驭万有的天主，以及那位教导唯独应崇拜天主、其他事物要么不存在、要么虽存在且值得尊敬但不值得崇拜和敬礼的那一位呢？关于这些事情，不仅相信而且以理性之眼沉思的人，会陈述他所想到的、经过仔细研究的证明。既然所有人类事务都依赖于信仰，为什么不信赖天主而不是信赖他们，岂不更为合理？因为谁在出海航行、缔结婚姻、生儿育女或向田地撒种时，不是相信这样做会带来更好的结果，尽管相反的情况可能且有时确实会发生？然而，相信更好的结果——甚至符合他们愿望的——会随之而来，使得所有人都冒险从事不确定的事业，这些事业的结果可能与他们预期的不同。如果对更好未来的希望和信仰支撑着每一件不确定的事业中的生活，那么为什么这种信仰不应被那位更有理由相信存在一位创造万有的天主、并相信那位以超凡的智慧和神圣的伟大胆量向世界各地的人宣布这一教义、甘冒极大危险和被视为耻辱的死亡、为人类而承受的那一位的人理性地接受呢？祂还教导那些起初被说服而接受祂教义的人，在面临各种危险和随时降临的死亡的情况下，前往世界各地确保人类的救恩。\par

% structure: p0031 source=h2 target=subsection reason=relative-heading-rank; base=h2

\subsection{第十二章}

其次，当塞耳苏斯明说：\Quote{如果他们回答我，不是像我问询信息那样——因为我对他们所有的意见都很熟悉——而是因为我同样关注它们全部，那将很好。如果他们不这样，而是一再重复他们常说的\InnerQuote{不要探究}等等，那么，}他继续说道：\Quote{至少向我解释他们所说的是些什么性质的东西，以及它们从何而来，}等等。现在，关于他声称\Quote{熟悉我们所有的教义}这一点，我们必须说，这是一个自夸而大胆的断言；因为如果他特别阅读了先知书——其中充满公认的难题和对大众隐晦的宣告——并且如果他通读了福音中的比喻、律法和犹太史的其他著作以及宗徒的言论，并且以诚恳的态度，渴望理解其含义，那么他就不会那样大胆地表示，说他\Quote{熟悉他们所有的教义}。甚至我们自己，虽然对这些著作花了大量研究，也不会说\Quote{熟悉一切}，因为我们尊重真理。我们中没有一个人会断言：\Quote{我熟知伊壁鸠鲁的全部教义，}或自信知道柏拉图的所有教义，因为知道这些体系的诠释者之间存在如此多的分歧。有谁敢说他了解斯多葛学派或逍遥学派的所有意见呢？除非，他确实从某些无知而愚昧的人那里听到这种自夸：\Quote{我都知道，}这些人并未察觉自己的无知，从而因有这样的人作导师，就以为知道了全部。在我看来，这样的人很像一个去过埃及的人——那里埃及的\OriginalTerm{学者}{savans}精通本国文献，极好地哲学思辨关于他们视为神圣的事物，但平民听到某些神话，不明其理，却因自以为有知识而十分得意——于是以为他通晓全部埃及知识，因为他只作了无知之徒的门徒，并未与任何司铎交往，也未从其他途径学习埃及的奥秘。我对埃及学者与无知者所说的话，也可适用于波斯人：他们之中有奥秘，由学者以理性原则进行，但被较浅薄的群众以象征意义理解。同样的评论也适用于叙利亚人、印度人以及所有拥有文献和神话的民族。\par

% structure: p0033 source=h2 target=subsection reason=relative-heading-rank; base=h2

\subsection{第十三章}

但既然塞耳苏斯宣称，许多基督徒有一句谚语：\Quote{今生的智慧是坏东西，愚拙是好的，}我们必须回答，他诽谤了福音，没有照保罗著作中实际出现的字句引用，那里是这样说的：\BibleQuote{若有人在这世界自以为有智慧，倒不如变为愚拙，好成为有智慧的。因为这世界的智慧，在天主面前是愚拙。}所以宗徒并没有简单地说\Quote{智慧在天主面前是愚拙}，而是\Quote{这\Emph{世界}的智慧}。同样，他没有说\Quote{你们中间若有人自以为有智慧，就让他普遍变为愚拙}，而是说\Quote{让他在这\Emph{世界}变为愚拙，好成为有智慧的。}因此，我们把\Quote{这世界的智慧}称为每一种虚假的哲学体系，按照圣经，这体系终归无有；我们称愚拙为好的，并非无条件，而是当一个人对这\Emph{世界}变为愚拙时。就好像我们说，相信灵魂不死和\OriginalTerm{灵魂转世}{metempsychosis}教义的柏拉图主义者，会被抛弃此意见的斯多葛学派、喋喋不休谈论柏拉图微妙之处的逍遥学派以及称引入眷顾和将天主置于万物之上为迷信的伊壁鸠鲁学派指责为愚拙。此外，更符合基督教精神的是，基于理性和智慧而不是仅凭信德来同意教义，只有在某些情况下基督教才要求后者，为不使人完全无助，这是耶稣的真门徒保罗所表明的，他说：\BibleQuote{因为世人凭自己的智慧，既不认识天主，天主就乐意用人所当作愚拙的道理，拯救那些信的人。}这些话清楚地表明，应当凭天主的智慧认识天主。但因这结果未发生，天主第二次乐意拯救那些相信的人，不是通过\Quote{愚拙}\Emph{普遍地}，而是通过那种取决于宣讲的愚拙。因为宣讲被钉十字架的耶稣基督，就是宣讲的\Quote{愚拙}，正如保罗也领悟到，他说：\BibleQuote{我们却是传被钉十字架的基督，这话在犹太人看来是绊脚石，在希腊人看来是愚拙；但蒙召的，无论是犹太人，还是希腊人，基督是天主的能力，天主的智慧。}\par

% structure: p0035 source=h2 target=subsection reason=relative-heading-rank; base=h2

\subsection{第十四章}

策尔苏斯认为在许多民族中可以找到一种普遍的教义关系，便列举了所有产生此类见解的民族；但出于某种我所不知的原因，他轻视犹太人，不将他们包括在其他民族之中，好像他们没有与其他民族一同努力并得出相同结论，或在许多问题上持有类似见解似的。因此，理应问他：为什么他相信异邦人和希腊人关于他们所谈论的那些民族古老性的记载，却唯独将这一民族的历史视为虚假？如果各自的作者本着真理精神记述了这些作品中所发现的事件，我们为何独独不信犹太人的先知？如果梅瑟和先知们出于偏爱自身体系的愿望而在其历史中记载了许多事情，我们为何不能对其他国家的史家也这样说？或者，当埃及人或其历史讲述犹太人的坏话时，他们在那一点上应当被相信；而犹太人说出同样关于埃及人的话，声称自己曾遭受他们极大的不义，并因此受到天主的惩罚，却要被指控为说谎？这不仅适用于埃及人，也适用于其他人；因为我们会发现亚述人与犹太人之间存在联系，这一点在亚述人的古老历史中有所记载。同样，犹太史家（我避免使用\Quote{先知}一词，以免显得预断案情）也记述了亚述人是犹太人的敌人。那么，请看此人任意的做法：他相信这些民族的历史，认为它们有学问，却谴责其他民族完全无知。因为请听策尔苏斯的陈述，他说：\Quote{从一开始就有一个权威的记述，所有最有学问的民族、城邦和人均对此一直保持一致。}然而，他却不愿像对待埃及人、亚述人、印度人、波斯人、奥德里西人、萨莫色雷斯人和厄琉息斯人那样，称犹太人为有学问的民族。\par

% structure: p0037 source=h2 target=subsection reason=relative-heading-rank; base=h2

\subsection{第15章}

毕达哥拉斯学派努美尼乌斯比策尔苏斯公正得多，他多次证明自己是一位极善言辞的人，仔细检验了许多见解，并从许多来源收集了具有真理外貌的东西；在他论\WorkTitle{论善}的第一卷书中，谈到那些采纳天主无形体之见解的民族时，他也将犹太人列入持这种观点的人中；他不仅毫不迟疑地在其论述中使用他们先知的语言，而且赋予其隐喻意义。此外，据说赫密波斯在其第一卷\WorkTitle{论立法者}中记载，毕达哥拉斯正是从犹太民族获得了他在希腊人中所引入的哲学。还有史家赫卡泰乌斯现存一部关于犹太人的著作，其中对该民族的学识给予了如此高的评价，以至于赫伦尼乌斯·斐洛在其论犹太人的著作中，首先怀疑这部作品是否真是那位史家的作品；其次，他说如果真是他的作品，那么很可能他是被犹太历史表面上的合理性所吸引，从而赞同了他们的体系。\par

% structure: p0039 source=h2 target=subsection reason=relative-heading-rank; base=h2

\subsection{第16章}

我必须表示惊讶，策尔苏斯竟将奥德里西人、萨莫色雷斯人、厄琉息斯人和极北人列入最古老、最有学问的民族之列，却认为犹太人不值得在此类民族中占有一席之地，无论是在学问还是古老性方面，尽管在埃及人、腓尼基人和希腊人中流传着许多著作证明他们是一个古老的民族，但我认为无需引用。因为任何人若愿意，都可以阅读弗拉维乌斯·约瑟夫斯在其两卷\WorkTitle{论犹太人的古老性}中所记载的内容，他在那里汇集了大量见证犹太民族古老性的作者；并且还存在年轻的塔提安的\WorkTitle{致希腊人演讲}，其中他以极大学问列举了那些论述犹太民族和梅瑟古老性的史家。因此，策尔苏斯作出这些陈述似乎并非出于对真理的热爱，而是出于仇恨，其目的是诋毁与犹太教相关的基督信仰的起源。不仅如此，他把荷马笔下的\OriginalTerm{加拉科法吉人}{Galactophagi}、高卢人的德鲁伊特人和格泰人称为最有学问和最古老的部落，因为他们的传统与犹太人的传统相似，尽管我不知道他们是否有任何历史幸存下来；但唯独希伯来人，他尽其所能地剥夺了他们古老性和学问的荣誉。再者，当列出那些凭借其行为对同时代人、凭借其著作对后代有所裨益的古代有学问之人时，他将梅瑟排除在外；而策尔苏斯在其列表中给予首位的利诺斯，既没有法律也没有讲论使任何部落变得更好；然而一个分散在整个世界上的整个民族都遵守梅瑟的法律。那么，请想一想，他岂不是出于公开的恶意，将梅瑟从其有学问之人的名单中驱逐出去，同时却断言利诺斯、穆赛乌斯、俄耳甫斯、斐瑞居德斯、波斯人琐罗亚斯德和毕达哥拉斯讨论了这些主题，并且他们的见解被记载在书中，从而一直保存至今？而且，他也是有意忽略以俄耳甫斯为主要润色的神话，其中描述诸神具有人性的弱点和情欲。\par

% structure: p0041 source=h2 target=subsection reason=relative-heading-rank; base=h2

\subsection{第17章}

随后，策尔苏斯攻击梅瑟的历史，指责那些以比喻和寓意解释它的人。在此，有人可以对这位伟人说，他在自己的作品上题写了\WorkTitle{真言}的标题：\Quote{善士啊，你为何夸耀那些记载，说诸神竟做出你们博学的诗人和哲学家所描述的冒险，犯下可憎的阴谋，与自己的父亲开战，割除自己的私处，竟敢犯下和遭受如此滔天恶行；而梅瑟，他关于天主没有这样的记载，甚至关于圣天使也没有，他关于人所叙述的罪行也轻得多（因为在他的著作中，无人敢于犯下\Person{克洛诺斯}对\Person{乌拉诺斯}，或\Person{宙斯}对父亲，或诸神与人类之父与自己女儿交合那样的罪行），却应当被视为欺骗了那些受他律法管辖的人，并引导他们走向错误？}此处，策尔苏斯在我看来有些像柏拉图派哲学家\Person{特拉西马库斯}那样，他不允许苏格拉底按自己的意愿回答关于正义的问题，而是说：\Quote{小心不要说效用就是正义，或义务，或诸如此类的东西。}因为同样地，策尔苏斯（如他所认为的）攻击梅瑟的历史，指责那些以寓意理解它们的人，同时却也称赞这样做的人，说他们更为公允（比不这样做的人）；于是，他仿佛以挑剔阻止那些能够揭示真相的人提出他们想要提出的辩护。\par

% structure: p0043 source=h2 target=subsection reason=relative-heading-rank; base=h2

\subsection{第十八章}

我挑战书籍与书籍的比较，会说：\Quote{来吧，善士，取下\Person{利努斯}、\Person{穆赛乌斯}、\Person{俄耳甫斯}的诗篇，以及\Person{斐瑞居德斯}的著作，仔细将它们与梅瑟的律法相比较——历史与历史相比，伦理论述与法律和诫命相比——看哪一方更适合当场改变听者的品性，哪一方更使其陷入邪恶；并注意，你们那一系列作家很少关心那些要立刻独立阅读的读者，而是将他们的哲学（如你们所称）写给那些能理解其比喻和寓意意义的人；而梅瑟，如同一位杰出的演说家深思某种修辞手法，并在每部分小心引入双重意义的语言，他在五部书中做到了这一点：在涉及道德的段落中，既不给他犹太臣民行恶的借口，也不给那些拥有更大智慧并能探究其意义的少数人一部缺乏思辨内容的论著。但是，你们那些博学的诗人的作品似乎不再被保存，尽管如果读者察觉到任何益处，它们本应被精心珍藏；而梅瑟的作品却感动了许多人，甚至包括那些对犹太习俗陌生的人，使他们相信，正如这些著作所证明，最初制定这些法律并交给梅瑟的，是创造世界的那位天主。因为宇宙的创造者在为治理宇宙制定法律后，应当赋予祂的言语一种力量，能征服大地上各处的所有人。我这样主张，尚未对耶稣进行任何探究，但已证明梅瑟，远远低于主，正如\WorkTitle{真言}将显示，远胜于你们那些智慧的诗人和哲学家。}\par

% structure: p0045 source=h2 target=subsection reason=relative-heading-rank; base=h2

\subsection{第十九章}

在这些陈述之后，策尔苏斯出于一种隐瞒的愿望，想要玷污梅瑟关于创造的记载，该记载教导说世界还不满一万年，而是远少于这个数字，尽管他掩盖了自己的愿望，却暗示他同意那些认为世界非受造的人。因为他主张，从永远以来就有许多大火和许多洪水，最近在\Person{丢卡利翁}时代发生的洪水相对现代，他清楚地告诉那些能理解他的人，他认为世界是非受造的。但是让这位基督教信仰的攻击者告诉我们，他是被什么论据迫使他接受有许多大火和许多大洪水的陈述，以及\Person{丢卡利翁}时代的洪水和\Person{法厄同}时代的大火比其他任何都更近期。如果他提出\Person{柏拉图}的对话作为这些问题的证据，我们会对他说，我们也可以相信，在梅瑟纯净而虔诚的灵魂中——他上升到所有受造物之上，与宇宙的创造者联合，并以比\Person{柏拉图}或那些希腊和罗马的智者更清晰的方式宣布了神圣事物——居住着一种神圣的精神。如果他要求我们给出这样信的理由，让他先给出他自己无根据主张的理由，然后我们将表明我们这一观点是正确的。\par

% structure: p0047 source=h2 target=subsection reason=relative-heading-rank; base=h2

\subsection{第二十章}

然而，\Person{策尔苏斯}违背自己的意愿，被迫作证说世界是相对年轻的，尚不满一万年，因为他提到希腊人认为那些事物是古老的，原因是由于大洪水与大火灾，他们未能目睹或接受任何更远古事件的记录。但让\Person{策尔苏斯}以他所认为的埃及人中最博学的人士作为关于大火灾与大洪水的神话的权威来源吧，这些人的智慧痕迹在对无理性动物的崇拜以及证明此类崇拜\OriginalTerm{天主}{God}符合理性并具有隐秘神秘特性的论据中可见。那么，埃及人当他们夸耀地讲述自己对动物神性的见解时，被认为是明智的；但如果某个犹太人，表明了自己对法律和立法者的忠诚，将一切归于宇宙的造物主和唯一的天主，那么在\Person{策尔苏斯}及其同类看来，他却被视为低于将神性不仅贬低到理性的、有死的动物层次，甚至贬低到无理性动物层次的人！——这种观点远远超越了神话式的灵魂转世论，该论认为灵魂从天穹之巅坠落，进入驯养和野生的野兽体内！如果埃及人讲述这类寓言，他们被认为通过其谜团和奥秘传达了哲学意义；但如果\Person{梅瑟}为整个民族撰写并留下历史和法律，则被视为空洞的寓言，其语言不容许任何寓意解释！\par

% structure: p0049 source=h2 target=subsection reason=relative-heading-rank; base=h2

\subsection{第二十一章}

以下是\Person{策尔苏斯}和伊壁鸠鲁派的观点：\Person{策尔苏斯}说：\Quote{梅瑟，他说，学到了存在于明智民族和雄辩之士中的教义，获得了神性的名声。}对此，我们必须回答说：可以允许他承认\Person{梅瑟}确实听到了一门相当古老的教义，并将其传给了希伯来人；如果他听到的教义是虚假的，既不虔诚也不可敬，而他仍然接受了并传给属下的人，那么他应受谴责；但如果，如你所断言，他坚持了明智真实的见解，并以此教育了他的人民，那么请问，他做了什么该受谴责的事呢？但愿不仅伊壁鸠鲁，而且亚里士多德（他对眷顾的看法不如前者那样不虔敬）以及斯多葛派（他们断言\OriginalTerm{天主}{God}是一物体）也听到了这样的教义！那么世界就不会充满那些要么否认或削弱眷顾作用，要么引入腐败的物质原则的见解，根据斯多葛派，他们的神是一物体，他们不怕说他会变化，可以在所有部分中被改变和转化，并且一般来说，他可能腐败，如果有人要他腐败，但他幸运地逃避了腐败，因为没有人要他腐败。然而，犹太人和基督徒的教义保存了神性的不变性和不能改变性，却被污蔑为不虔敬，不是因为它分享了那些对天主概念不虔敬者的亵渎，而是因为它向神性的祈祷中说：\BibleQuote{祢是同样的，}而且，据信仰条文，天主曾说过：\BibleQuote{我不改变。}\par

% structure: p0051 source=h2 target=subsection reason=relative-heading-rank; base=h2

\subsection{第二十二章}

此后，\Person{策尔苏斯}虽未谴责犹太人所行的割礼，却声称这一习俗源自埃及人；因此他相信埃及人，而不相信说\Person{亚巴郎}是首先实行此仪式的人的\Person{梅瑟}。不仅\Person{梅瑟}提及\Person{亚巴郎}的名字，赋予他与天主极大的亲密关系；而且许多从事召叫邪灵实践的人，在他们的咒语中也使用\Quote{亚巴郎的天主}这一表达，藉此名字指出那义人与天主的友谊。然而，他们虽然使用\Quote{亚巴郎的天主}这个短语，却不知道\Person{亚巴郎}是谁！\Person{依撒格}、\Person{雅各伯}和\Person{以色列}也是如此；这些名字尽管明确是希伯来文，却常被那些声称凭知识产生奇妙效果的埃及人使用。然而，始于\Person{亚巴郎}的割礼仪式，后来被\Person{耶稣}废止，祂不希望祂的门徒遵行，现在不待我们解释；因为当前时机不宜谈论这些事，而是要努力驳斥\Person{策尔苏斯}对犹太人教义的指控，他认为如果通过攻击基督宗教在犹太教中的起源，能够证明犹太教也是虚假的，那么他就更容易确立基督宗教的虚假性。\par

% structure: p0053 source=h2 target=subsection reason=relative-heading-rank; base=h2

\subsection{第二十三章}

在此之后，策尔苏斯接着断言：\Quote{那些追随梅瑟为领袖的牧人和牧羊人，被庸俗的欺骗迷惑了心智，以致以为只有一位天主。}那么，请他证明，在他看来，这些牧人和牧羊人从崇拜众神非理性地离开之后，他自己如何能够确立希腊人或其他被称为野蛮人的民族中所存在的众多神祇。因此，请他确立由宙斯生下的缪斯之母谟涅摩叙涅的存在，或时序女神之母特弥斯的存在；或者请证明永远赤身露体的美惠三女神能够有真实、实质的存在。但他无法从她们的任何行动中证明，这些希腊人的虚构形象，看似具有身体，实际上却是（真正的）神祇。为什么希腊人关于神祇的寓言比例如埃及人的寓言更真实？埃及人在他们的语言中完全不知道九位缪斯之母谟涅摩叙涅，也不知道时序女神之母特弥斯，也不知道美惠三女神之一的欧佛洛绪涅，也不知道任何其他这些名字。通过我们在世界令人赞叹的秩序中所见而产生确信，我们敬拜世界的造物主，视之为一个效果的唯一作者，这个效果因其自身完全和谐，因此不可能由许多作者造成，这岂不是更明显（且比所有这些发明好得多！）？并且我们相信整个天穹并非由许多灵魂的运动维系，因为一个就够了，它承载着整个不游动的球体从东到西，并在自身内包含世界所需的一切，而这些都不是自存的！因为一切都是世界的部分，而天主不是整体的部分。但天主不能是不完美的，如同部分是不完美的一样。也许更深沉的思考会表明，既然天主不是部分，那么祂也不恰当地是整体，因为整体由部分构成；理性不允许我们相信统摄万有的天主是由部分构成的，其中每个部分都不能做所有其他部分所能做的事。\par

% structure: p0055 source=h2 target=subsection reason=relative-heading-rank; base=h2

\subsection{第二十四章}

此后他继续说道：\Quote{这些牧人和牧羊人得出结论，只有一位天主，名为至高者，或阿多奈，或天上的，或撒巴特，或他们乐于给予这个世界的其他一些名字；此外他们一无所知。}在本书的后面部分，他说：\Quote{统摄万有的天主，无论被冠以希腊人中流行的宙斯之名，还是例如印度人或埃及人所使用的名字，都没有区别。}对此，我们须指出，这涉及一个深奥而神秘的课题——即关于名字的本质问题：究竟是如亚里士多德所想，名字是约定授予的，还是如斯多葛派所持，名字是自然产生的，最初的话语是事物的模仿，据此名字形成，并符合他们引入的某些词源学原则；抑或如伊壁鸠鲁所教导（在此与斯多葛派不同），名字是自然赋予的——最初的人根据他们所处环境说出了某些话语。如果我们能够针对前述陈述，确立强大名字的本质，其中一些名字被埃及的学者、波斯的贤士、被称为婆罗门的印度哲学家、萨马纳人以及不同国家的其他人所使用；并且能够表明所谓的魔法，并不像伊壁鸠鲁和亚里士多德的追随者所认为的那样是完全不确定的东西，而是如那些精通魔法的人所证明，是一个连贯的体系，拥有极少数人知道的话语；那么我们就说，\OriginalTerm{撒巴特}{Sabaoth}和\OriginalTerm{阿多奈}{Adonai}的名号，以及其他在希伯来人中受到如此尊崇的名字，并非适用于任何普通受造物，而是属于指向万有塑造者的奥秘神学。因此，这些名字，当伴随着与其本性相称的适当环境而被念出时，拥有巨大的力量；而其他在埃及语言中流行的名字，对某些只能做特定事情的恶魔有效；波斯语言中的其他名字对其他灵体有相应的力量；如此在每个不同的民族中，各有不同的用途。由此会发现，地上各种恶魔被分派到不同地方，每个都带有与地域和国家的各种方言相称的名字。因此，对这些事物有哪怕一点点更高贵观念的人，会小心避免将不同的名字应用于不同的事物；免得他像那些错误地将天主之名应用于无生命物质的人，或者将\Quote{善}的称号从第一因、或从美德与卓越拖下来，应用于盲目的普路托斯，以及健康的、比例得当的肉、血和骨头的混合物，或者被认为高贵出身的东西。\par

% structure: p0057 source=h2 target=subsection reason=relative-heading-rank; base=h2

\subsection{第二十五章}

也许有一种危险，如同将\Quote{天主}、或\Quote{善}之名用于不当对象那样大，即按照一套隐秘体系改换天主的名号，把属于较低存在者的名称用于较高者，\Emph{反之亦然}。我不再赘述：当提起\Person{宙斯}之名时，同时也听到\Person{克洛诺斯}和\Person{瑞亚}之子、\Person{赫拉}的丈夫、\Person{波塞冬}的兄弟、\Person{雅典娜}和\Person{阿尔忒弥斯}的父亲，且与自己的女儿\Person{珀耳塞福涅}乱伦；或者\Person{阿波罗}立即让人想到\Person{勒托}和\Person{宙斯}之子、\Person{阿尔忒弥斯}的兄弟、\Person{赫尔墨斯}的同父异母兄弟；如此类推，所有其他名字都是这些\Person{凯尔苏斯}的智者所发明，他们是这些意见的始祖，也是希腊的古代神学家。有什么理由决定一方面他应恰当地被称为\Person{宙斯}，另一方面又不该有\Person{克洛诺斯}为父、\Person{瑞亚}为母？同样的论证适用于所有被称为神者。但这项指控完全不适用于那些因某种奥秘理由而将\Quote{萨巴奥特}、\Quote{阿多奈}或任何其他名号归于（真正的）天主的人。当人能够就名字的奥秘进行哲学思考时，他将会就天主的使者的称号论说很多，其中一位称\Person{弥额尔}，另一位称\Person{加俾额尔}，另一位称\Person{辣法耳}，这些称号恰合他们在世间履行职务的方式，皆按万有天主的旨意。类似的名字哲学也适用于我们的\Person{耶稣}，祂的名字已经以一种不容置疑的方式，将无数的恶灵从（人的）灵魂和身体中驱逐出去，祂的名号对被逐者所施展的能力如此巨大。仍在名字这个话题上，我们必须提及精通咒语的人讲述，同一咒语以其原本语言诵念能完成咒文所宣称的功效；但译成另一种语言时，则变得无效而无力。因此，并非所指的事物，而是词语的性质和特性，对某种目的拥有某种能力。因此，基于这些理由，我们为基督徒的行为辩护，他们甚至斗争至死，避免以\Person{宙斯}之名称呼天主，或从任何其他语言中给祂取名。因为他们要么笼统地使用通常的名号——天主——要么加上诸如\Quote{万物的创造者}、\Quote{天地的造物主}之类的附加语——祂向人类派遣了那些善人，将天主之名加于他们名下，于是在人间行出某种大能的事工。此外，关于名字这一主题，还可说更多，以驳斥那些认为我们应不关心如何使用名字的人。如果\Person{柏拉图}在\WorkTitle{斐利布斯篇}中的话令我们惊讶，他说：\Quote{我的恐惧啊，\Person{普罗泰戈拉}，关于诸神的名字可不小，}因为\Person{斐利布斯}在与\Person{苏格拉底}讨论时称快乐为\Quote{神}，我们难道不更应赞许基督徒的虔诚，他们完全不将神话中使用的任何名号用于世界的造物主？关于此话题，目前就说这么多。\par

% structure: p0059 source=h2 target=subsection reason=relative-heading-rank; base=h2

\subsection{第二十六章}

但让我们看看这位自称无所不知的\Person{策尔苏斯}是如何诬蔑犹太人的：他声称\Quote{他们崇拜天使，沉迷于巫术，而\Person{梅瑟}是他们的导师}。现在，让他这个自称通晓基督信仰和犹太教一切的人告诉我们，他在\Person{梅瑟}的著作中哪里找到立法者规定崇拜天使；也让他说明，在那些接受了\Person{梅瑟}法律并读到\BibleQuote{不可求问巫师，以免被你们玷污}禁令的人中，巫术如何可能存在。此外，他还承诺稍后要展示\Quote{犹太人如何因无知而受骗，陷入错误}。如果他发现犹太人对基督的无知是由于他们没有听到关于祂的预言，那么他会真实地说明犹太人如何陷入错误。但他根本没有意愿让这一点显明，反而把根本不是错误的东西视为犹太人的错误。\Person{策尔苏斯}既已承诺在他作品的后半部分让我们了解犹太教的教义，首先开始谈论我们的救主，即祂是我们这一世代（就我们基督徒而言）的领袖，并说\Quote{几年前祂开始教导这教义，被基督徒视为天主子}。关于这一点——即祂几年前才存在——我们必须作如下评论。如果没有天主的助佑，\Person{耶稣}在这几年中渴望传播祂的话和教导，能如此成功，以至于世界各地，不论希腊人或野蛮人、学者或无知者，都有不少的人接受了祂的教义，甚至为此奋斗至死而不肯否认，这是任何其他体系都从未有过的事吗？我确实不想谄媚基督信仰，而是渴望彻底考察事实，我要说，即使那些从事医治众多病人的人，如果没有天主的帮助，也无法达到治愈身体的目的；如果有人成功地将灵魂从邪恶、放荡、不义以及对天主的轻视洪流中解救出来，并展示出有一百个人（假定数目这么大）的性情改善了，作为这样的证据，没有人会合理地说，他在那一百个人心中植入能够除去那么多邪恶的教义是没有天主助佑的。如果有人坦诚考虑这些事情，承认人类任何改善没有天主帮助都无法实现，那么当他把许多皈依祂教义的人以前的生活与后来的行为相比较，并想想他们过去沉溺于怎样的放荡、不义和贪婪，直到像\Person{策尔苏斯}和他同感的人所说\Quote{他们受骗了}，接受了据这些人说是有害于人类生命的教义；但从他们接受之日起，他们变得在一些方面更温和、更虔诚、更坚定，以至于其中一些人，出于极度贞洁的渴望和更纯洁地敬拜天主的意愿，甚至戒绝了（合法的）爱情的允许放纵。\par

% structure: p0061 source=h2 target=subsection reason=relative-heading-rank; base=h2

\subsection{第二十七章}

任何考察此事的人都会看到，\Person{耶稣}尝试并成功完成了超出人类能力的事工。因为尽管从起初，所有事物——当时的君王、他们的首领和将军，以及大体上拥有最小权力的人，此外还有各城的统治者、士兵和人民——都反对祂教义在世界的传播，但它却证明是得胜的，因为它是天主圣言，其本性是不受阻挠的；它比所有这样的敌人都更强大，征服了整个\Place{希腊}和相当一部分蛮族地区，并召集了无数灵魂归向祂的宗教。虽然在皈依基督信仰的众多人中，单纯无知者必然多于较有智慧者，如同前一类人总是多于后一类人一样，但\Person{策尔苏斯}不愿注意这一点，认为这种仁慈的教义——它触及太阳下每一个灵魂——是粗俗的，并且因其粗俗和缺乏理性能力，仅获得了无知之人的接受。然而他自己也承认，并非只有单纯的人被\Person{耶稣}的教义引导而接受祂的宗教；因为他承认其中有中等智力、性情温和、有理解力并能理解比喻的人。\par

% structure: p0063 source=h2 target=subsection reason=relative-heading-rank; base=h2

\subsection{第二十八章}

既然他模仿修辞学家训练门徒的方式，引入一个与耶稣进行个人辩论的犹太人，并以极其幼稚的语气说话，完全不相称于哲学家的白发，让我尽力按我的能力审查他的陈述，并表明他在整个讨论中没有保持与犹太人角色相称的一致性。因为他描写他与耶稣争辩，并自以为在许多点上驳倒了祂；首先，他指责祂\Quote{发明了他从童贞女所生}，并辱骂祂\Quote{生于某个犹太村庄，一个贫穷的乡下妇女，她以纺织维持生计，被她的丈夫——一个木匠——赶出家门，因为她被判定犯有通奸罪；被丈夫驱逐后，她流浪了一段时间，羞辱地生下了耶稣，一个私生子；他因贫穷在埃及做雇工，并在那里获得了一些神奇的权能——埃及人对此深感自豪——然后回到自己的国家，因这些权能而得意洋洋，并藉此宣称自己是天主。}现在，既然我不能允许不信者所说的任何话不加审查，而必须从头调查一切，我认为所有这些都恰当地与耶稣是天主之子的预言相和谐。\par

% structure: p0065 source=h2 target=subsection reason=relative-heading-rank; base=h2

\subsection{第二十九章}

因为出身有助于一个人成名、显赫和被人议论；即当一个人的父母恰好有地位和影响，拥有财富，并能够用于儿子的教育，而当出生地又伟大而光荣时；但当一个人拥有所有这些不利条件，却仍然能够克服这些障碍，使自己被认识，给听到他的人留下印象，并变得杰出，被全世界看见——全世界以前未曾如此谈论他——我们怎能不钦佩这样的本性，认为它本身高尚，致力于伟大事业，并拥有一种绝不可轻视的勇气呢？如果更全面地考察这样一个人的历史，为什么不探究：他如何在节俭和贫穷中长大，没有接受完整的教育，没有学习那些能使他自信地与群众交往、煽动民意、吸引许多听众的体系和观点，却仍然致力于教导新的见解，向人类引入一种不仅颠覆了犹太人的习俗（同时保持对他们先知的适当尊重），而且尤其推翻了希腊人关于天主的既定惯例的学说呢？这样一个受过如此教养、正如他的诽谤者承认从未从人那里学到任何伟大事物的人，如何能以绝非可轻视的方式教导关于神圣审判、惩治邪恶的惩罚、以及赏赐给德行的奖赏的教义，以至于不仅粗野无知的人被他的话所吸引，而且也不少以智慧著称的人——他们能够辨别那些被认为更普通的流传教义的隐藏意义，这隐藏意义可以说包裹着某种更深奥的含义？柏拉图笔下的塞里福斯人在地米斯托克利因军事才能成名后责备他，说他的名声不是由于自己的功劳，而是由于有幸出生在希腊最著名的城邦；这位好脾气的雅典人看到自己的祖国确实有助于他的声誉，回答道：\Quote{若我是塞里福斯人，我也不会像我这样杰出；你若是有幸是雅典人，你也不会成为地米斯托克利！}如今，我们的耶稣，被指责出生在一个村庄，且不是希腊的，也不属于任何广受尊敬的民族，被视为一个贫穷劳动妇女的儿子而受鄙视，因贫穷离开祖国在埃及做雇工——用刚才的例子来说，不仅是塞里福斯人（一个非常小且无名的岛上的居民），而且可以说是塞里福斯人中最卑微的——却能够震动整个有人居住的世界，不仅远远超过地米斯托克利这位雅典人，甚至超过毕达哥拉斯、柏拉图或世界上任何地方的任何其他智者，或任何王子或将军的成就。\par

% structure: p0067 source=h2 target=subsection reason=relative-heading-rank; base=h2

\subsection{第三十章}

如今，凡以平常心探究这些事实本质的人，岂不都对该人物之胜利深感惊异？他竟能以声名完全克服一切招致恶名之因，并超越世上所有其他杰出人物，这岂非奇迹？然而，卓著之士鲜能同时赢得数种声誉。有人因其智慧受钦仰，有人因其军事才能，有些外邦人因其不可思议的咒术，有人因某一品质，另有人因另一品质；但能同时因多种德行受钦仰并得声望者不多。唯独此人，除其他优点外，既因其智慧，又因其奇迹，又因其治理能力而令人惊叹。因为他劝服一些人脱离他们自己的律法，归附于他——并非如暴君所为，亦非如强盗以武装追随者攻击他人；亦非如富人施恩于来投者；亦非如那些公然应受指责之人；而是作为关于万有之天主的教义、属于天主的敬拜以及一切能取悦至高天主、使遵行之人蒙恩的道德训诲的导师。对于\Person{特米斯托克利}或任何其他显赫人物，没有什么事妨碍他们的声誉；但对于此人，除了我们已经列举的、足以使最高贵之人的灵魂蒙羞的那些之外，还有那表面上耻辱的十字架之死，这足以抹去他先前获得的荣耀，并使那些如否认其教义者所声称的、曾受他欺骗的人放弃他们的迷误，并谴责他们自己的欺骗者。\par

% structure: p0069 source=h2 target=subsection reason=relative-heading-rank; base=h2

\subsection{第三十一章}

此外，人们或许惊奇：\Person{耶稣}的门徒——如果如诽谤\Person{耶稣}者所说，他们并未在他死后看见复活，也未被说服相信他的神性——为何不惧怕与他们的师傅同受苦难、置身危险、离开家乡，按照\Person{耶稣}的意愿去教导他所传授给他们的教义？因为我认为，任何坦诚考察事实的人都不会说，这些人为了\Person{耶稣}的教义而献身于危险的生活，除非他们内心深处坚信其真理性——相信\Person{耶稣}不仅教导他们遵守他的诫命，也教导他人，并且当那些敢于在所有地方、所有听众面前引进这些新见解的人，以及那些不能保留任何仍持旧见旧习的人作为朋友时，显然有丧命之虞。\Person{耶稣}的门徒难道没有看到，当他们不仅用先知性圣经向犹太人证明这位被先知所称道之人，也向其他外邦民族证明，这位前一天或前两天被钉十字架者是为人类种族自愿受死的吗？——这类似于那些为国捐躯以消除瘟疫、荒歉或风暴的人。因为很可能在事物的本性中，出于某些难以被大众理解的奥秘原因，有这样一种德能：一个义人自愿为公益而死，就能驱除邪灵——那些邪灵是瘟疫、荒歉、风暴或类似灾祸的原因。因此，那些不愿相信\Person{耶稣}为人类钉死于十字架的人，就问一问他们是否也不接受希腊人和外邦人中流传的许多关于为公众利益牺牲生命以消除降临城邦和国家的灾祸的记载？或者他们愿说这些事实际发生过，但那个关于这个被称为人者为摧毁一个强大的恶灵——即恶灵统治者，他曾奴役地上所有灵魂——而死的故事不足采信？\Person{耶稣}的门徒看到了这一点和更多（很可能\Person{耶稣}私下教导过他们），并且充满了神圣的能力（因为赋予他们力量和勇气的并非诗中的贞女，而是天主的真智慧与真理解），他们竭尽全力\Quote{在所有人中崭露头角}，不仅在\Place{阿尔戈斯}人中，而且在所有希腊人和外邦人中，并且\Quote{为自己赢得光荣的声誉}。\par

% structure: p0071 source=h2 target=subsection reason=relative-heading-rank; base=h2

\subsection{第三十二章}

但现在让我们回到那位犹太人被引入的地方，他谈论耶稣的母亲，说：\Quote{当她怀孕时，她被曾许配给他的木匠赶出门外，因为她犯了奸淫，而且她生了一个孩子，给了某个名叫\Person{潘特拉}的士兵。}让我们看看，那些盲目编造关于童贞玛利亚与\Person{潘特拉}奸淫以及她被木匠抛弃的寓言的人，是否不是为了推翻祂由\OriginalTerm{圣神}{Holy Ghost}奇妙受孕而编造了这些故事？因为他们本可以以另一种方式伪造历史，因为这件事极其奇妙，而且他们本可以不像勉强承认的那样，承认耶稣不是由普通的婚姻所生。的确，那些不相信耶稣奇妙诞生的人，预料他们会编造一些谎言。但他们没有以可信的方式这样做，而是保留了童贞玛利亚怀孕耶稣不是来自若瑟这个事实，这使得谎言对那些能够理解和识破这种编造的人非常明显。那么，一个敢于为人类做如此多事的人，为的是尽可能使所有期待神圣审判的希腊人和蛮族离弃邪恶，并以取悦\OriginalTerm{宇宙的创造者}{Creator of the world}的方式规范他们全部行为，难道他的诞生不应该奇妙，而是最卑劣、最可耻的吗？我要以希腊人的身份问他们，特别是问\Person{塞尔苏斯}，他要么持有\Person{柏拉图}的观点，要么不持有，总之他引用了它们：那位将灵魂遣入人身的天主，难道使这个敢于做出如此伟大事迹、教导如此多人、从世上的大量邪恶中改革如此多人的灵魂，降临到一种比任何其他诞生都更可耻的诞生中，而不是通过合法的婚姻将他引入世界吗？或者，难道不是更符合理性：每个灵魂，为了某些神秘的原因（我现在根据\Person{毕达哥拉斯}、\Person{柏拉图}和\Person{恩培多克勒}的意见说话，\Person{塞尔苏斯}经常提到他们），都被引入一个身体，并且根据其功德和先前的行为被引入？因此，很可能这个灵魂，它在肉身中的居留所带来的益处比许多人（为避免偏见，我不说\Quote{所有人}）要多，它需要的身体不仅优于其他身体，而且具备所有卓越的品质。\par

% structure: p0073 source=h2 target=subsection reason=relative-heading-rank; base=h2

\subsection{第三十三章}

现在，如果某个灵魂，为了某些神秘的原因，不值得被置于一个完全无理性的身体中，也不值得被置于一个纯粹理性的身体中，而是被赋予一个怪异的身体，使得理性无法在这样的身体中正常运作，这个身体头部与其它部分不成比例且整体太短；而另一个灵魂则接受一个这样的身体，使得灵魂比前一个稍微理性一些；另一个则更理性一些，身体的性质或多或少地阻碍了理性原则的接受；那么，为什么不应也有一个灵魂接受了完全奇妙的身体，既拥有一些与其他人类共同的特质，以便能和他们一起度过一生，也拥有一些优越的品质，使得灵魂能够保持不受罪的玷污呢？如果面相学家——无论是\Person{佐皮洛斯}、\Person{洛克斯}、\Person{波莱蒙}，还是其他任何写过此类主题的人——的学说有任何真实性，他们以某种奇妙的方式宣称所有身体都适应灵魂的习性，那么那个将以奇妙能力居住在人中间并施行伟大事迹的灵魂，难道必须有一个像\Person{塞尔苏斯}所认为的、由\Person{潘特拉}与童贞玛利亚通奸所生的身体吗？！不，从这种不圣洁的结合中，必然会产生一个愚人，给人类带来伤害——一个教导放荡、邪恶和其他罪恶的教师，而不是教导节制、正义和其他德行！\par

% structure: p0075 source=h2 target=subsection reason=relative-heading-rank; base=h2

\subsection{第三十四章}

但正如先知们所预言的，将有一位从贞女而生，按所预许的记号，祂将使这事实以其名为标记，表明在祂诞生时天主与人同在。如今，在我看来，一位犹太人引用依撒意亚的预言，说厄玛奴耳将由贞女诞生，是适宜的。然而，自称无所不知的凯尔苏斯却没有这样做，或出于无知，或出于不愿（如果他读过并故意沉默不语）提供一个可能挫败他意图的论据。这预言是这样说的：\BibleQuote{主又对阿哈次说：\InnerQuote{你向上主你的天主要求一个记号，或求诸阴府深处，或求诸上天高处。}但阿哈次说：\InnerQuote{我不要求，我也不试探上主。}他说：\InnerQuote{如今，达味家族，你们听着！你们使人厌烦还算小事，还要使我的天主厌烦吗？因此，主要亲自给你们一个记号：看，有位贞女要怀孕生子，给他起名叫厄玛奴耳，意思是\InnerQuote{天主与我们同在。}}}而且，凯尔苏斯故意怀着恶意没有引用这预言，对我而言是清楚的，因为尽管他大量引用玛窦福音的经文，如基督诞生时出现的星以及其他奇迹，他却只字未提。现在，如果有一位犹太人咬文嚼字，说经文不是\Quote{看，一位贞女}，而是\Quote{看，一位少女}，我们回答说：\OriginalTerm{Olmah}{Olmah}一词——七十贤士译本译作\Quote{贞女}，其他译本译作\Quote{少女}——按他们的说法，出现在申命纪中，用于\Quote{贞女}，如下文：\BibleQuote{如果有一位贞女的少女已许配给一个丈夫，有人在城里遇见她，与她同寝，你们应将他们二人带到该城门口，用石头砸死他们；那少女因为她在城里没有喊叫，那男人因为他污辱了邻人的妻子。}又如：\BibleQuote{如果有人在田野遇见一个已许配的少女，强行与她同寝，那么只有与她同寝的男人该死；但少女，你不可对她做什么；在她身上没有该死的罪。}\par

% structure: p0077 source=h2 target=subsection reason=relative-heading-rank; base=h2

\subsection{第三十五章}

但为了不使我们因为一个希伯来词，而试图说服那些无法决定是否应当相信的人，认为先知预言了这个人的出生是由贞女所生，因为在他出生时说了\Quote{天主与我们同在}这句话，就让我们从这些话本身来论证我们的观点。记载说主对阿哈次说：\BibleQuote{为你自己向上主你的天主求一个记号，或求诸深渊，或求诸高处；}随后给了这记号：\BibleQuote{看，一位贞女要怀孕生子。}那么，这会是怎样的一个记号呢——一个不是贞女的少女生育孩子？两者中哪一个更适合作为厄玛奴耳（即\Quote{天主与我们同在}）的母亲呢？是一个与男子同寝、按妇女方式怀孕的女子，还是一个仍为纯洁神圣的贞女？显然，只有后者才适合产生这样的存在：在他出生时被称为\Quote{天主与我们同在}。如果他吹毛求疵地说，那命令是向阿哈次发出的：\Quote{为你自己向上主你的天主求一个记号，}我们要反问，在阿哈次的年代，有谁生了一个儿子，在他出生时使用了\Quote{厄玛奴耳}即\Quote{天主与我们同在}这个说法呢？如果找不到任何人，那么显然对阿哈次所说的话是对达味家族说的，因为经上记载救主按肉身生于达味家族；并且这记号被称为\Quote{在深处或在高处}，因为\BibleQuote{那下降的，就是上升到诸天之上、充满万有的那一位。}我这些论证是针对一个相信预言的犹太人而用的。现在让凯尔苏斯或任何与他同见的人告诉我，先知说出这些关于未来的陈述，或预言中包含的其他句子，是带着什么意思？是出于对未来的预见，还是不是？如果是出于对未来的预见，那么先知们是受天主默感的；如果不是出于对未来的预见，那么就请他解释一个如此大胆谈论未来、并因先知能力而受犹太人景仰的人的含义。\par

% structure: p0079 source=h2 target=subsection reason=relative-heading-rank; base=h2

\subsection{第三十六章}

现在，既然我们已经触及先知的主题，我们即将提出的论点不仅对相信他们是由天主灵感而发言的犹太人有益，而且对希腊人中的较开明者也有益。对这些我们这样说：我们必须承认犹太人有先知，如果他们要在所赐予他们的法律体系下被保守，并相信他们所学到的世界的创造主，并且就法律而言，没有任何借口背教到外邦人的多神教去。我们以下列方式建立这种必要性。因为\BibleQuote{外邦人}，正如犹太人的法律本身所记载的，\BibleQuote{他们必听从占星术士和占卜者}；但对那民族却说：\BibleQuote{至于你们，主你们的天主不许你们这样做。}此外还附加了应许：\BibleQuote{主你们的天主要从你们兄弟中给你们兴起一位先知。}既然外邦人使用占卜方式，或通过神谕，或通过征兆，或通过鸟占，或通过腹语者，或通过那些自称献祭术的人，或通过迦勒底谱系家——所有这些做法都被禁止给犹太人——那么这个民族，如果没有方法获得对未来的知识，被人类普遍的求知未来的热情所驱使，就会轻视自己的先知，认为他们没有任何神圣的成分；并且不会接受任何在\Person{梅瑟}之后的先知，也不会将他们的言语记录下来，而是会自发地转向外邦人的占卜习俗，或试图在自己中间建立类似的做法。因此，他们的先知甚至对无关紧要的事件发出预言，以安抚那些渴望这类事情的人，这并无荒谬之处，就如\Person{撒慕尔}预言关于丢失的三头母驴，或在列王纪第三卷中提到关于一位王子生病的事。为什么那些渴望从偶像那里得到预兆的人不应受到犹太人律法管理者的严厉责备呢？——正如我们发现\Person{厄里亚}责备\Person{阿哈齐雅}，说：\BibleQuote{因为在以色列没有天主吗，你竟去求问\Place{厄刻龙}的神巴耳则步布？}\par

% structure: p0081 source=h2 target=subsection reason=relative-heading-rank; base=h2

\subsection{第三十七章}

我认为，那么，我们已经相当充分地证明了不仅我们的救主应由童贞女诞生，而且犹太人中也有先知，他们不仅发出关于未来的普遍预言——例如，关于基督、世界的王国、以色列将要发生的事，以及那些将要相信救主的民族，还有许多其他关于祂的事——而且也有关于具体事件的预言，例如，克士的驴如何丢失了，将要被找回，以及关于降在以色列王子身上的疾病，以及其他类似记录的细节。但作为对希腊人的进一步回答，他们不相信耶稣由童贞女诞生，我们必须说：创造主通过几种动物的生育表明，祂在一种动物身上所做的，如果祂愿意，也能在其他动物身上做，甚至在人本身身上做。因为已经确定有一种雌性动物不与雄性交配（如动物学家所说，秃鹫就是这种情况），而这种动物没有性交就能延续种族。那么，假设如果天主想派遣一位神圣的导师到人类中间，祂以不同于通常的方式使他诞生，这有什么不可信的呢？不，根据希腊人自己的说法，并非所有的人都是由一男一女所生。因为如果世界是被创造的，正如许多希腊人也乐于承认的那样，那么最初的人必定不是通过性交产生的，而是从大地中产生的，其中含有种子元素；然而，我认为这比耶稣像其他人一样出生更不可信，就他诞生的一半而言。而且，用希腊历史来回答希腊人，以表明我们并不是唯一使用这类奇迹叙述的人，这并无荒谬之处。因为有些人认为，不仅是在古代和英雄的叙述中，而且在最近发生的事件中，可以叙述一件可能的事：\Person{柏拉图}是\Person{安菲克提翁}的儿子，\Person{阿里斯顿}被阻止与他的妻子行房，直到她生下她由\Person{阿波罗}所怀的孩子。然而，这些都是真实的传说，它们导致了关于一个人的这类故事的发明，这个人被他们认为拥有比大众更多的智慧和力量，并且他的身体物质的开始来自于比其他人更好、更神圣的元素，因为他们认为这对于那些太伟大以至于不能成为人类的人是适当的。既然\Person{凯尔苏斯}引入了犹太人与耶稣辩论，并且像他想象的那样，撕碎了耶稣由童贞女诞生的虚构，同时比较了关于\Person{达娜厄}、\Person{墨拉尼珀}、\Person{奥革}和\Person{安提俄珀}的希腊传说，我们的回答是：这样的语言适合一个小丑，而不是一个以严肃语气写作的人。\par

% structure: p0083 source=h2 target=subsection reason=relative-heading-rank; base=h2

\subsection{第三十八章}

此外，他接纳了按照\BibleBook{玛窦}{福音}所载的吾主下到\Place{埃及}的历史，却拒绝相信伴随此事的神迹奇事，即：或是天使传达了神圣的启示，或是吾主离开\Place{犹太}并居住在\Place{埃及}是一件有意义的事；他却凭空捏造了完全不同的事情，既承认\Person{耶稣}所行的奇事——借此他引导群众跟随他作为\Person{基督}——又想要贬低这些奇事，认为是借着魔法而非神能成就的；因为他断言\Quote{他（\Person{耶稣}）自幼被当作私生子养大，曾在\Place{埃及}作佣工，然后习得某些神奇能力，从那里回到本国，并借着这些能力自称是神。}我不明白，一个行邪术的人怎会竭力教导这样一种教义，它始终劝勉我们当\Person{天主}将按各人的行为审判来行事；又怎会以同样的信念训练他的门徒，就是他后来用作传播其教义的工具。因为门徒们究竟是教了他们如何行奇迹之后才打动听众，还是没有借助这些？因此，断言他们根本没有行奇迹，而是被一些毫无说服力的论证（如同希腊辩证法的智慧）所说服，然后投身于教导他们所寄居之人的新教义，这完全是荒谬的。因为他们教导教义、传播新思想时，所凭恃的是什么呢？但如果他们确实行了奇迹，那么怎会相信行邪术的人会冒如此大的风险去引入一个禁止行邪术的教义呢？\par

% structure: p0085 source=h2 target=subsection reason=relative-heading-rank; base=h2

\subsection{第三十九章}

我不认为有必要与一种并非严肃而是嘲笑的态度所提出的论点纠缠，例如：\Quote{如果\Person{耶稣}的母亲是美丽的，那么那本性不爱可朽身体的神，就因为她美丽而与她交合；}或\Quote{那神极不可能对她产生情欲，因为她既不富有也无王室地位，甚至她邻舍中也没有人认识她。}他也是以同样的嘲讽口气补充道：\Quote{当她为丈夫所憎恶并被赶出家门时，她并未被神能所救，她的故事也无人相信。这些事，}他说，\Quote{与\OriginalTerm{天国}{kingdom of heaven}毫无关系。}这样的言论与那些在街头辱骂他人、其言语不值得任何认真关注的人有何区别？\par

% structure: p0087 source=h2 target=subsection reason=relative-heading-rank; base=h2

\subsection{第四十章}

在这些断言之后，他从\BibleBook{玛窦}{福音}，或许也从其他福音中，取了关于鸽子在\Person{若翰}的洗礼中降临到我们的救主身上的记载，并企图贬低这一陈述，声称该叙述是虚构的。当他想象着自己已经完全驳斥了主由童贞女所生的故事之后，他并没有按顺序处理后续的记载；因为激情和仇恨不守秩序，愤怒和报复的人会随着情绪来临而诽谤他们所恨的人，他们的激情阻止他们以仔细有序的计划安排指控。因为如果他遵循了适当的次序，他会拿起福音书，为了攻击它，先反对第一个叙述，然后转到第二个，依此类推。但现在，在童贞女生子之后，这位自诩通晓我们所有历史的\Person{塞耳苏}，攻击了在洗礼中以鸽子形式显现的圣神的记载。然后，在那之后，他试图诋毁关于主将要来到世界的预言。接下来，他跳到了紧接\Person{耶稣}诞生叙述之后的内容——关于那颗星以及从\Place{东方}来朝拜婴儿的\OriginalTerm{贤士}{wise men}的记载。而你若肯费心，也可以发现\Person{塞耳苏}在他的整本书中许多混乱的陈述；以至于即使在这一点上，那些懂得观察并要求有条理安排的人，可以判定他极其鲁莽和狂妄，因为他竟在他的作品上题写了\WorkTitle{真言}的标题——这是博学的哲学家从未做过的事。因为\Person{柏拉图}说，对于一个有智慧的人来说，对有些不确定的事情做出强烈的断言并非明智的表现；就连著名的\Person{克吕西波}，也常常陈述他据以作出决定的理由，并将我们引向那些我们认为比自己更能言善辩的人。然而，这个人，比上述各位以及所有其他希腊人都更聪明，依他自诩通晓一切，竟在他的书上题写了\WorkTitle{真言}二字！\par

% structure: p0089 source=h2 target=subsection reason=relative-heading-rank; base=h2

\subsection{第四十一章}

但是，为了不使我们看起来是因无力反驳而故意忽视他的指控，我们决心尽己所能逐一回答每一项指控，不关注事物本身性质的联系和顺序，而是按照本书中出现的主题顺序。因此，让我们注意他关于圣神以鸽子形体向我们的救主显现的诋毁言论。有一个犹太人向我们所承认的主\Person{耶稣}说了以下的话：\Quote{你在洗礼时，那犹太人说，在\Person{若翰}旁边，你说有一只鸟从空中降到你身上。}然后这同一个犹太人继续质问道：\Quote{有什么可信的证人看到了这个显现？或者谁听见天上传来声音，宣告你是\OriginalTerm{天主子}{Son of God}？除了你自己的断言和另一个与你一同受罚的人的话之外，还有什么证据？}\par

% structure: p0091 source=h2 target=subsection reason=relative-heading-rank; base=h2

\subsection{第四十二章}

在我们开始答辩之前，必须先指出：试图证明几乎任何历史（不论多么真实）确实发生过，并对其形成合理的理解，是所能尝试的最困难的工作之一，而且在某些情况下是不可能做到的。例如，假设有人声称从未有过任何特洛伊战争，主要是由于其中交织着难以置信的叙述，说某位\Person{阿基里斯}是海神\Person{忒提斯}和凡人\Person{珀琉斯}之子，或\Person{萨耳珀冬}是\Person{宙斯}之子，或\Person{阿斯卡拉福斯}和\Person{伊阿耳墨诺斯}是\Person{阿瑞斯}之子，或\Person{埃涅阿斯}是\Person{阿佛洛狄忒}之子，我们该如何证明这些事是真实的呢？尤其是，不知何故，虚构附在普遍流行的看法上，即伊利昂确实发生过一场希腊人与特洛伊人之间的战争？又假设有人不相信\Person{俄狄浦斯}和\Person{伊俄卡斯忒}以及他们的两个儿子\Person{厄忒俄克勒斯}和\Person{波吕尼刻斯}的故事，因为故事中引入了一种半处女的怪物——斯芬克斯，我们又如何证明这类事情的真实性呢？同样，关于\Person{厄庇戈诺伊}的历史，尽管其中没有交织这样神奇的事件，或关于\Person{赫拉克勒斯}后裔的回归，以及无数其他历史事件，也是如此。但那些坦诚对待历史、并希望自己不被历史欺骗的人，会运用自己的判断力，决定同意哪些陈述，接受哪些陈述为比喻，努力发现这类虚构作者的含义，以及拒绝相信哪些陈述，因为这些陈述是为了取悦某些人而写的。我们预先说了这些，是关于\WorkTitle{福音}中有关耶稣的全部历史，并不是邀请敏锐的人简单地、不加理性地相信，而是希望表明，阅读的人需要坦诚，需要大量研究，并且，可以这样说，需要对作者含义的洞察力，以便发现每件事被记载下来的目的。\par

% structure: p0093 source=h2 target=subsection reason=relative-heading-rank; base=h2

\subsection{第四十三章}

因此，我们首先要说：如果那个不信圣神以鸽形显现的人被描述为伊壁鸠鲁派、德谟克利特的追随者或逍遥派，那么这种说法是与这类反对者的性格一致的。但如今连这位最聪明的\Person{策尔苏斯}也没有察觉，他竟然把这样的反对意见加给一个犹太人，而这犹太人相信先知书中所含的内容比鸽子显现的叙述更难以置信！因为有人会对那犹太人（当他表示不信这显现，并想抨击它为虚构时）说：\Quote{先生，您如何能证明主对\Person{亚当}、\Person{厄娃}、\Person{加音}、\Person{诺厄}、\Person{亚巴郎}、\Person{依撒格}和\Person{雅各伯}所说的话，就是记载中祂对他们所说的那些话呢？}而且，为了比较历史与历史，我会对那犹太人说：\Quote{就连你们自己的\Person{厄则克耳}也写道，说：\BibleQuote{天开了，我看见了一个天主的异象。}叙述之后，他又说：\BibleQuote{这是上主荣耀的样子的显现；祂对我说，}等等。}如今，如果关于耶稣的记载是假的（因为我们不能像你们所设想的，清楚证明它是真的，它只被祂自己看见或听见，并且如你们似乎观察到的，也被一个受罚者看见），我们为何不反而说\Person{厄则克耳}在说\BibleQuote{天开了}等等时也是在讲奇迹呢？甚至\Person{依撒意亚}也断言：\BibleQuote{我看见万军的上主坐在高高举起的宝座上；色辣芬们侍立在旁：一个有六个翅膀，另一个也有六个翅膀。}我们如何知道他是真的看见了还是没有呢？啊，犹太人，你相信这些异象是真实的，并且不仅是由一个更神圣的圣神显示给先知，而且也是由同一位圣神说出并记录的。那么，当宣称天在他面前开了，他听见一个声音，或看见万军的上主坐在高高举起的宝座上时，\Person{依撒意亚}和\Person{厄则克耳}与耶稣相比，谁更值得相信呢？在前者，确实没有发现什么工作能与后者相提并论；而耶稣的善行不仅局限于祂在肉身内居住的时期，直到现在，祂的能力仍然在那些通过祂而相信天主的人心中产生皈依和生活的改善。一个明显的证据表明这些事是由祂的能力所做的，那就是：虽然正如祂自己所说的，并且也被公认，没有足够的工人来收割灵魂的庄稼，但是实际上仍然有如此大量的庄稼被聚集并运送到处处存在的禾场和天主的教会中。\par

% structure: p0095 source=h2 target=subsection reason=relative-heading-rank; base=h2

\subsection{第四十四章}

我用这些论点回答那犹太人。我作为基督徒，并不否定\Person{厄则克耳}和\Person{依撒意亚}，而是非常希望基于我们共同的信仰基础，表明这位（耶稣）比他们更值得信赖，因为祂说祂看见了这样的景象，并且很可能将自己所见的异象和所听见的声音告诉了祂的门徒。但另一方可能会反对说，并非所有叙述鸽子显现和天上声音的人都从耶稣那里听说过这些事，而是那教导\Person{梅瑟}关于他以前的事件（从创造开始，一直下到他的祖先\Person{亚巴郎}）的同一圣神，也教导了\WorkTitle{福音}的作者们在耶稣受洗时所发生的奇迹。并且，那享有称为\Quote{智慧之言}的神恩的人，也会解释天开、鸽子显现的原因，以及为何圣神以鸽子的形态、而非任何其他生物的形态向耶稣显现。但我们目前的主题并不要求我们解释这一点，我们的目的是要表明，\Person{策尔苏斯}将犹太人描绘为基于这种理由不相信一个事实——该事实比他所坚定相信的许多事件具有更大的可能性——是没有表现出一丁点健全判断的。\par

% structure: p0097 source=h2 target=subsection reason=relative-heading-rank; base=h2

\subsection{第四十五章}

我记起有一次，在当着许多审判官的面与某些被称为博学的犹太人辩论时，我使用了以下论证：\Quote{告诉我，诸位，}我说，\Quote{既然有两位人物曾降临到人类，关于他们流传着超越人类能力的奇迹——即你们自己的立法者梅瑟，他写了关于自己的事，以及我们的导师耶稣，他没有留下关于自己的著作，但门徒们在福音中为他作证——那么，有什么理由断定梅瑟是可信的、说的是真理，尽管埃及人毁谤他是巫师，似乎他靠骗术行了大能的事，而耶稣却因为你们控告他而不被相信？然而，有各民族为两者作证：犹太人为梅瑟作证；基督徒并不否认梅瑟的先知使命，反而从那同一来源证明关于耶稣的陈述是真实的，并接受他的门徒所记载的关于他的奇迹。现在，如果你们要我们说明信仰耶稣的理由，请你们先给出信仰比耶稣更早的梅瑟的理由，然后我们再将接受后来的耶稣的理由给你们。但如果你们退缩、回避论证，那么，我们效法你们的榜样，此刻也拒绝提出任何论证。不过，请承认你们无法为梅瑟提出证明，然后再听我们从法律和先知为耶稣所作的辩护。现在请看这几乎难以置信的事！从法律和先知中关于耶稣的宣言可以看出，梅瑟和先知们确实都是天主的先知。}\par

% structure: p0099 source=h2 target=subsection reason=relative-heading-rank; base=h2

\subsection{第四十六章}

因为法律和先知中充满了与耶稣受洗时所记载的奇迹相似的奇迹，即关于鸽子与来自天上的声音。而且我认为耶稣所行的奇事是圣神那时以鸽子的形像显现的证明，尽管刻尔苏斯出于贬低它们的意图，声称耶稣只做了他在埃及学到的那些。我不仅会提到他的奇迹，而且也会适当地提到耶稣宗徒们的奇迹。因为他们若不借助奇迹奇事，就无法说服那些听到他们新教义和新教导的人，使他们抛弃自己的民族习惯，并冒着生命危险接受他们的训诲。在基督徒中，至今仍保留着那以鸽子形像显现的圣神的痕迹。他们驱逐恶灵，行许多治愈，并按照\OriginalTerm{圣言}{\GreekText{λόγος}}的旨意预见某些事件。尽管刻尔苏斯，或他引入的犹太人，可能会嘲笑我将要说的话，但我仍要说：许多人几乎是违背自己意愿地归向了基督教，某种灵突然改变了他们的心，使他们从憎恨这教义转变为愿意为它而死，并以醒时的异象或夜晚的梦显现给他们。我们已知许多这样的例子，如果我们把它们写下来，尽管是我们亲眼所见、亲耳所闻，我们也会给不信者提供极大的嘲笑机会，他们会认为我们与那些他们认为是编造这些事情的人一样，自己也做了同样的事。但天主见证我们良心的愿望，不是用虚假的陈述，而是用各种不同的见证来确立耶稣教义的神性。既然是一个犹太人对圣神以鸽子的形像降临到耶稣身上的叙述感到困惑，我们要对他说：\Quote{先生，在依撒意亚书中是谁说：\InnerQuote{现在主派遣了我和他的圣神。}}在这句话中，意思有歧义——即圣父与圣神派遣了耶稣，还是圣父派遣了基督与圣神——后一种解释是正确的。因为救主被派遣后，圣神也被派遣了，为使先知的话得以应验；由于这预言的应验必须为后代所知，耶稣的门徒因此将结果记载下来。\par

% structure: p0101 source=h2 target=subsection reason=relative-heading-rank; base=h2

\subsection{第四十七章}

我想对刻尔苏斯说，他描写那犹太人多少承认了洗者若翰是给耶稣施洗的，但洗者若翰的存在——他为赦罪而施洗——是由一位距若翰与耶稣时代不远的人记载的。因为在《\WorkTitle{犹太人古事记}》第十八卷中，若瑟夫证实若翰是一位施洗者，并应许那领受洗礼的人得洁净。这位作者虽然不相信耶稣是基督，但在寻求耶路撒冷沦陷和圣殿被毁的原因时——他本应说，对耶稣的谋害是这些灾难降临到百姓身上的原因，因为他们处死了基督，一位先知——却说（虽不情愿，但离真相不远）：这些灾祸临到犹太人，作为对义人雅各伯之死的惩罚，雅各伯是耶稣（称为基督）的兄弟，犹太人将他处死，尽管他是一个以正义著称的人。耶稣的真门徒保禄说，他视这位雅各伯为主内的兄弟，与其说是由于血缘关系或一同长大，不如说是由于他的德行与教导。若若瑟夫说耶路撒冷的荒凉是因雅各伯而临到犹太人，那岂不更合理地说，是因耶稣基督（的死）而发生的？基督神性的见证者是如此众多的教会，它们由从罪孽洪流中聚集而来的人组成，他们归向了造物主，并将自己的一切行为都归于他的善意。\par

% structure: p0103 source=h2 target=subsection reason=relative-heading-rank; base=h2

\subsection{第四十八章}

纵然那犹太人不能为厄则克耳和依撒意亚的例子中天开于耶稣及祂所听见的声音作任何辩护，当我们比较厄则克耳和依撒意亚或任何其他先知书中所记载的类似情形时，我们仍将尽力支持我们的立场，主张：既然许多人相信在\Emph{梦中}有印象被带到了他们心中，有的是关于神圣之事，有的是关于现世未来之事，且或清晰或晦涩——凡接受\OriginalTerm{天主眷顾}{providence}道理的人都明白这一事实——那么，说那能在梦中接受印象的心智，也能在醒着的异象中受到感动，以造福于受感动者或将要从他听闻此事的人，又有何荒谬呢？正如在梦中我们自以为听见，觉得听觉器官真的被触动，又自以为用眼睛看见——虽然身体上的视觉或听觉器官并未受影响，而只是心智产生这些感觉——同样，当记载说先知们目睹了相当奇妙的事件，如听见上主的话语或看见天开了时，相信类似之事发生在它们身上也并无荒谬。因为我不认为那可见的天真的被打开了，其物理结构被分裂，为的是让厄则克耳能记下这样的事。那么，凡聪明的福音听众岂不应同样相信救主所经历之事？虽然这样的事可能会成为简单之人的绊脚石，他们天真地想要移动整个世界，并撕裂那整个天的坚固而庞大的身体。但更深究此事的人会说：如圣经所称，有一种神圣的普遍知觉，只有有福之人知道如何发现，按撒罗满的话说：\BibleQuote{你必寻得认识天主的知识。}这种知觉力量有多种形式，例如一种视力本能地看见比身体更美好的事物，其中包括\OriginalTerm{革鲁宾}{cherubim}和\OriginalTerm{色辣芬}{seraphim}；一种听力能听见不在空中存在的声响；一种味觉能享用那从天上降下、赐生命于世界的活粮；同样，一种嗅觉能嗅到使保禄说他是在天主前属于基督的馨香之气；还有一种触觉，若望说他\BibleQuote{亲手摸过生命之言}。有福的先知们发现了这种神圣知觉，以这种神圣方式看见、听见、品尝、嗅闻（可谓如此），没有感官知觉器官，借着信德抓住\OriginalTerm{圣言}{Logos}，以致从它而来的医治之流临到他们身上；他们以此方式看见了他们所记下的，听见了他们所说听见的，并如同他们所描述的那样，当吃下交给他们的书卷时，也有了类似的感受。同样，依撒格闻到他儿子神圣衣服的香气，并在属灵祝福中加了这句话：\BibleQuote{看，我儿子的香气，如同上主祝福的丰盛田地的香气。}类似地，且更应理解为心智而非感官所感知的，耶稣触摸了癞病人以洁净他，依我看，是双重的意义——不仅如众人所听见的，藉可见的接触使他脱离可见的癞病，也藉祂真正神圣的触摸使他脱离另一种癞病。正是为此，若翰作证说：\BibleQuote{我看见圣神仿佛鸽子从天降下，停在祂身上。我不认识祂，但那派遣我来以水施洗的，对我说：你看见圣神降下，停在谁身上，谁就是用圣神施洗的。我看见了，且作证：这就是天主之子。}那时天向耶稣开了；据记载，除了若翰，没有人看见天开了。关于这天开，救主预告门徒此事将要发生，且他们要看见，说：\BibleQuote{我实实在在告诉你们：你们要看见天开了，天主的使者在人子身上上去下来。}同样，保禄曾被提到第三层天，之前他看见天开了，因为他是耶稣的门徒。但解释保禄为何说\BibleQuote{或在身内，我不知道；或在身外，我不知道：天主知道}不属于我们现在的目标。但我要在论证中加入刻尔苏斯所想象的那些点，即耶稣亲自讲述了天开以及圣神仿佛鸽子降在祂身上于约旦河的事，虽然圣经并未声称祂说自己看见了。因为这位伟人没有看出，这与那在山上显现时命令门徒\BibleQuote{不可将所见告诉任何人，除非人子从死人中复活}的那位不相称，竟然把自己所见的和在约旦河边若翰所听见的告诉门徒。因为可以观察耶稣品格的特色，祂在所有场合避免不必要的自述，为此祂说：\BibleQuote{我若为自己作证，我的见证就不真。}由于祂避免不必要的自述，宁愿以行动而非言语表明祂就是基督，犹太人才对祂说：\BibleQuote{祢若是基督，请明明地告诉我们。}既然在刻尔苏斯的著作中，一位犹太人对耶稣说了关于圣神以鸽子形状显现的话：\Quote{这是祢自己的见证，除了祢所举出的一位与祢同受刑罚的人外，别无支持。}我们有必要向他表明，这样的话不适合放在犹太人嘴里。因为犹太人并不把若翰与耶稣联系起来，也不把若翰的刑罚与基督的刑罚联系起来。以此例子，这位自夸博学的人被证明不知道在犹太人与耶稣辩论时，应该让那犹太人说什么话。\par

% structure: p0105 source=h2 target=subsection reason=relative-heading-rank; base=h2

\subsection{第四十九章}

此后，他不知为何，故意撇开了证实耶稣言论的最有力证据，即祂的来临已由犹太先知——梅瑟以及这位立法者前后相继的众先知——所预言；我认为，这是因为无法应对这一论点：无论犹太人或是任何异端教派，都不否认基督是预言的主题。但也许他对有关基督的预言并不熟悉。因为一个熟悉基督徒陈述（许多先知预言了救主的来临）的人，绝不会将更适合于撒玛黎雅人或撒杜塞人发表的言论归于一位犹太人；对话中的犹太人也不会以像如下这样的语言表达自己：\Quote{但我的先知曾在耶路撒冷宣告，天主子将作为义人的审判者和恶人的惩罚者来临。} 并非只有一位先知预言了基督的来临。虽然接受梅瑟五书的撒玛黎雅人和撒杜塞人会说，其中包含关于基督的预言，但预言绝不是在那时候尚未被提及的耶路撒冷发出的。真希望所有控告基督教的人，都像策尔苏斯一样——不仅不了解事实，甚至不了解圣经的字面意义——这样他们攻击基督教时，他们的论点就不会有丝毫动摇——我不是说信德——而是那些不稳定和暂时的信徒的微小信德。然而，一位犹太人绝不会承认任何先知使用了\Quote{\InnerQuote{天主子}将要求临}的说法；因为他们使用的术语是\Quote{\InnerQuote{天主的基督？}将要求临}。确实，他们常常直接质问我们关于\Quote{天主子}，说没有这样一个存在，或者它从未成为预言的主题。当然，我们不否认\Quote{天主子}是预言的主题；但我们断言，他极不恰当地将如下的言论归于那位犹太辩士（他不会承认祂是）：\Quote{我的先知曾在耶路撒冷宣告，\InnerQuote{天主子}将要求临。}\par

% structure: p0107 source=h2 target=subsection reason=relative-heading-rank; base=h2

\subsection{第五十章}

其次，仿佛唯一被预言的事件就是祂将成为\Quote{义人的审判者和恶人的惩罚者}，仿佛祂的出生地、祂在犹太人手中所受的苦难、祂的复活以及祂所行的奇事都未曾成为预言的主题，他接着说：\Quote{为什么单单是你，而不是其他无数成千上万在预言发表之后存在的人，这些预言才适用于他们？} 并且，我不知道他为何想暗示别人可能认为他们自己就是先知所指的人，他说：\Quote{有些人被狂热驱使，另一些人聚集了一大批跟随者，宣称天主子已从天降下。} 我们并没有确定这类事情发生在犹太人中间。那么，我们首先要指出：许多先知以各种方式发表关于基督的预言；有些通过隐晦的说法，有些通过比喻或以其他方式，还有些则以直白的话语。正如他在下文以犹太人的口吻对来自他本民族的皈依者说话，并强调且怀恶意地重复说：\Quote{有关祂生平事件的预言也可能同样适合其他事件。} 我们将从更多的预言论中列举少数几个；对于这些，任何人可以选择说出他认为能够确保反驳它们的话，并使有智识的信徒远离信德。\par

% structure: p0109 source=h2 target=subsection reason=relative-heading-rank; base=h2

\subsection{第五十一章}

圣经论及救主诞生之地，即那将出自\Place{白冷}的统治者，如此说：\Quote{你，白冷，厄弗辣大家，在犹大诸城中并非最小；因为将由你中为我出一位在以色列中作统治者的人；他的来历源于亘古，出于永远。} 这个预言不适用于任何一位像\Person{克耳苏斯}之犹太所说的那些人，即那些狂热分子和群众领袖，且宣称自己来自天上的人，除非明确显示他出生于\Place{白冷}，或如另一人所说，出自\Place{白冷}成为民众的领袖。关于耶稣在\Place{白冷}的诞生，如果有人希望，在\Person{弥该亚}的预言和耶稣的门徒们所记载的福音史之外，从其他来源获得更多证据，那么让他知道，与福音中关于他诞生的叙述相符，在\Place{白冷}展示着他出生的山洞，以及洞内他裹着襁褓的马槽。这一景象在周围地区广为谈论，甚至在信仰的敌人中间也是如此，据说在这个山洞里诞生了那位被基督徒敬拜和尊崇的耶稣。此外，我认为，在基督来临之前，民众的司祭长和经师们，因这预言清晰明了，曾教导说基督将在\Place{白冷}诞生。这种意见在犹太人中也很盛行；为此记载说，当\Person{黑落德}向民众的司祭长和经师询问时，听到他们说基督要在\Place{犹太地}的\Place{白冷}诞生，\Quote{因为达味曾在那里}。在\BibleBook{若望}{福音}中也记载，犹太人说基督要在\Place{白冷}诞生，\Quote{因为达味曾在那里}。但在我们主来临之后，那些忙于推翻他的诞生地从起初就是预言主题之信仰的人，向人民隐瞒了这教导；他们的行径类似那些人，即收买了那些守卫坟墓的卫兵、亲眼目睹他从死者中复活的人，并指使这些目击者如此报告：\Quote{就说：他的门徒趁我们睡觉的时候，来把他偷走了。如果这事被总督听见，我们会说服他，并确保你们无事。}\par

% structure: p0111 source=h2 target=subsection reason=relative-heading-rank; base=h2

\subsection{第五十二章}

争论与偏见是强有力的工具，使人甚至忽视那些极其明显的事；因此，那些已经习惯某种见解的人，这些见解深深印入他们的心灵，并给他们打上某种特征，就不会放弃这些见解。因为一个人较容易放弃在其他事情上的习惯（尽管摆脱习惯可能是困难的），而放弃他的意见则更难。不，即使是前者，那些习惯它的人也不容易放弃；因此，无论是房屋、城市、村庄，还是亲密的熟人，当我们对它们有好感时，我们都不愿意离开。因此，这就是为何当时许多犹太人忽视了预言、耶稣所行的奇迹以及他据称所受的苦难的明证。人类的天性如此受影响，对于那些观察到一旦对最卑劣和微不足道的祖先及同胞传统产生偏见的人，难以摆脱它们的人来说，这是显而易见的。例如，没有人能轻易说服一个埃及人轻视他从其父亲那里学到的，以致不再视这个或那个非理性的动物为神，或不再避免吃这种动物的肉，甚至冒着死亡的危险。现在，如果在相当长的篇幅中继续这一主题的考察，我们列举了关于\Place{白冷}及其预言的要点，我们认为我们有义务这样做，以作为对那些人的反驳，这些人宣称：如果流传在犹太人中间关于耶稣的预言如我们所说的那样清楚，为什么他们在他来临时不赞同他的道理，并投身于他所指出的更美好的生活？然而，没有人应该对信徒提出这样的责难，因为可以看到那些学会陈述理由的人为他们对耶稣的信仰提出了有力的理由。\par

% structure: p0113 source=h2 target=subsection reason=relative-heading-rank; base=h2

\subsection{第五十三章}

若要问第二个清楚指向耶稣的预言，我们可以引用梅瑟在基督降世前许多年所写的话，即当雅各伯临别尘世时，为他每个儿子发出预言，关于犹大，他说：\BibleQuote{权杖必不离犹大，统治者的杖必不离他脚间，直到那为他保留的来到。}任何读到这预言的人——这预言实际上远在梅瑟之前就已存在，以至于一个不信者可能怀疑不是梅瑟所写——都会惊讶梅瑟竟能预言犹太人的首领（在十二支派中）应出自犹大支派，并成为人民的统治者；正因如此，整个民族被称为犹太人，得名于这个统治支派。第二，一个坦诚思考这预言的人会惊讶，在宣布人民首领和统治者将出自犹大支派之后，又划定他们统治的界限，说：\BibleQuote{权杖必不离犹大，统治者必不离他脚间，直到那为他保留的来到，他是外邦人的期望。}因为那为他保留的已经来了，即天主的基督，天主预许的首领。显然，他是先前所有——我敢说也包括后来所有——人中唯一作为外邦人的期望的那一位；因为所有外邦民族的皈依者都通过他相信了天主，这正符合依撒意亚的预言：外邦人曾寄望于他的名：\BibleQuote{外邦人将寄望于你的名。}这个人也对比喻为囚徒的人——每个人都为自己罪的锁链所俘虏——说：\BibleQuote{出来吧！}对无知的民众说：\BibleQuote{进入光明！}这些事也早已预言：\BibleQuote{我立你作人民的盟约，复兴地业，使荒芜的产业归人承受；对囚徒说：出来吧！对在黑暗中的人说：显现出来吧！}我们可以在这个人显现时，通过世界各地单纯相信他的人，看到这一预言的应验：\BibleQuote{他们将在路上牧放，他们的草场必在一切大道上。}\par

% structure: p0115 source=h2 target=subsection reason=relative-heading-rank; base=h2

\subsection{第五十四章}

由于塞尔苏斯虽然自称通晓福音，却因救主的苦难而责难他，说他未从圣父得到任何帮助，或无力帮助自己；我们必须指出，他的苦难连同其原因都是预言的题目；因为他为人类益处而死，为他所受的定案而受鞭伤。此外，预言说有些外邦人要认识他（其中不包括先知）；并且预言他将以一种在人间被视为可耻的形象显现。预言的话是这样：\BibleQuote{看哪，我的仆人必有智慧，必被高举、受光荣，且被提升到极高。同样，许多人将因你而惊奇；所以你的容貌在人间必无尊荣，你的光荣在世人间也必如此。看哪，许多民族要因他而惊奇；君王要闭口无言：因为那些未曾得到关于他的消息的人将要看见他；那些未曾听过他的人将要认识他。}\BibleQuote{主啊，谁相信了我们的报道？上主的臂膀向谁启示了呢？我们在他面前好像婴孩，如同干地上的根。他没有容貌也没有光荣；我们看见他，他既无容貌也无美丽；但他的外表毫无光彩，比所有人都更逊色。他是一个受苦的人，懂得忍受病痛；因为他的面容被转开，他受到轻视，被视为无足轻重。这个人担当了我们的罪，为我们而忍受痛苦；我们以为他受罚、受苦、受虐待。但他为了我们的过犯被刺透，为了我们的罪孽被压伤。他受的惩罚使我们得平安，因他受的鞭伤我们得医治。我们众人都像羊一样迷了路。一个人迷了路，上主因我们的罪将他交付；他因受恶待，不开口。他像羊被牵到宰杀之地，又像羊羔在剪毛者面前无声，所以他也不开口。他受屈辱时，他的审判被夺去。谁能述说他的世代？因为他的生命从地上被夺去；因我人民的罪孽，他被引到死地。}\par

% structure: p0117 source=h2 target=subsection reason=relative-heading-rank; base=h2

\subsection{第五十五章}

现在我记起，有一次在与某些被视为有智慧的犹太人进行辩论时，我引用了这些预言；而我的犹太对手回答说，这些预言关乎整个民族，被视为一个个体，并且处于分散和受苦的状态，以便许多外邦人因犹太人在众多外邦民族中的分散而成为归信者。他用这种方式解释了这些话：\Quote{你的形象在人前必无尊荣；}然后，\Quote{未曾有人向他报信的人将看见；}以及\Quote{一个受苦的人}。在那次讨论中，我们用许多论证来证明这些关于某个特定人物的预言，并不像他们那样适用于整个民族。我问他们，\Quote{这个人承担我们的罪，并为我们受苦，}和\Quote{但祂为我们的罪过受伤，为我们的愆尤被压伤，}这些话适用于哪个角色？以及\Quote{因祂的创伤我们得了痊愈}这句话正当地属于谁？因为很明显，正是那些曾是罪人、并因救主的苦难而得医治的人（无论是犹太民族还是外邦归信者），在预知这些事件的先知著作中使用这样的语言，并且这位先知在圣神的感动下将这些话应用于一个人。但我们似乎用这句话最有力地驳倒了他们：\Quote{因我人民的罪过，祂被引至死亡。}因为如果按照他们的说法，人民是预言的主体，那么这个人怎么会因天主的子民的罪过而被引至死亡呢？除非祂与天主的子民不同？而这个人除了耶稣基督之外还能是谁？因祂的创伤，信祂的人得了痊愈，当\Quote{祂解除了管辖我们的执政者和掌权者，并公然在十字架上羞辱了他们}时。我们或许会在另一次机会详细解释预言的各个部分，不留下任何一个未经审查。但这些事已经较为详细地讨论了，我认为这是必要的，因为凯尔苏作品中引用的犹太人的言论。\par

% structure: p0119 source=h2 target=subsection reason=relative-heading-rank; base=h2

\subsection{第五十六章}

凯尔苏和他所引入的犹太人，以及所有不信耶稣的人，都没有注意到预言谈到基督的两次来临：第一次以人类的苦难和谦卑为特征，为的是基督与人同在时，能揭示通向天主的道路，并使今世中无人有借口说自己不知道未来的审判；第二次则唯独以荣耀和神性为显著，没有任何人性的软弱与祂的神性伟大交织其中。详细引用预言会显得冗长；我认为目前引用第四十五首圣咏的一部分就足够了，这首圣咏除了其他标题外，还有这样的题词：\Quote{为挚爱者的圣咏，}其中显然对天主说：\Quote{恩宠倾注于你的嘴唇：因此天主永远祝福你。大能者啊，将你的剑佩在腰间，带着你的荣美和威严。愿你向前驰骋，为真理、谦逊和公义而统治；你的右手将奇妙地引领你。大能者啊，你的箭尖锐；仇敌心中的人民必仆倒在你以下。}但要仔细注意下文，其中称祂为天主：\Quote{天主，你的宝座永远长存；你国度的权杖是公义的权杖。你喜爱正义，憎恶不义；因此天主，你的天主，用喜乐油膏你，超过你的同伴。}请注意，先知以亲密的语气对天主说话，\Quote{你的宝座永远长存，} \Quote{你国度的权杖是公义的权杖，}说这位天主被一位作为祂的天主的神所膏，被膏是因为祂比祂的同伴更喜爱正义、憎恶不义。我记得我正是用这段经文严厉地追问那位被视为有学问的犹太人；他对此感到困惑，便按照他犹太主义的观点回答说，\Quote{天主，你的宝座永远长存，你国度的权杖是公义的权杖，}这些话是论及万有的天主；而\Quote{你喜爱正义，憎恶不义，因此你的天主膏了你，}等等，是指默西亚。\par

% structure: p0121 source=h2 target=subsection reason=relative-heading-rank; base=h2

\subsection{第五十七章}

此外，犹太人向救主这样说：\Quote{如果你说，按天主圣意所生的每一个人都是天主子，那你与其他人在哪方面不同呢？}对此我们回答说：凡如\Person{保禄}所说的，不再处于恐惧之下，如同有师保，而只因善本身而择善者，就是\Quote{天主子}；但此人之远胜于凡因德性而被称为\OriginalTerm{天主子}{son of God}的人，因为祂，如同是这一切的来源和起始。\Person{保禄}的话如下：\BibleQuote{因为你们没有领受奴隶的精神，以致再恐惧；你们领受的却是义子的圣神，因此我们呼叫：阿爸，父啊！}但是，按照\Person{克尔苏}笔下的犹太人，\Quote{无数的人会指证耶稣为谎者，声称那些关于祂的预言其实是指着他们自己说的。}我们确实不知道\Person{克尔苏}是否知道有人来到这个世界，并渴望像\Person{耶稣}那样行事，就自称也是\Quote{天主子}或天主的\Quote{德能}。既然我们本着真理的精神查考每一段经文，我们将提及：在基督诞生前，犹太人中有个名叫\Person{特乌达}的，自命不凡，他死后他所迷惑的追随者就完全溃散了。在他之后，在户口登记的日子，那时\Person{耶稣}似乎已诞生，有个加里肋亚人\Person{犹大}聚集了许多犹太民众，自称是一个智慧的人和新奇教理的教师。当他也因叛乱受到惩罚后，他的教理被推翻，只影响了极少数人，且都是最卑微的人。在\Person{耶稣}的时代之后，撒玛黎雅人\Person{多思推}想使撒玛黎雅人相信他就是梅瑟所预言的基督；他似乎赢得了某些人赞同他的观点。但是，引述《\WorkTitle{宗徒大事录}》中\Person{加玛里耳}极其明智的观察，以显示上述那些人与此应许无关，他们既不是\Quote{天主子}也不是天主的\Quote{德能}，而基督\Person{耶稣}才是真正的\OriginalTerm{天主子}{Son of God}，这并非不合理的。\Person{加玛里耳}在所述段落中说：\BibleQuote{如果这计划或这事业是由人来的，必要消灭}（正如上述那些人死后他们的图谋也归于消灭），\BibleQuote{但若是从天主来的，你们就不能推翻这个教导，以免你们也被发现与天主对抗。}还有撒玛黎雅人\Person{西满}，一个行邪术的，他想用他的邪术引诱一些人。那时他成功了；但如今，我想，在全世界上不可能找到三十个他的追随者，或许我还多说了。在\Place{巴力斯坦}只有极少数；而在世界其他地方，他本想在那里传播他名字的光荣，却无处提说。即使有提及，也是从《\WorkTitle{宗徒大事录}》中引述的；因此，他之所以被人提及，全归功于基督徒，而不可辩驳的事实证明\Person{西满}丝毫没有神性。\par

% structure: p0123 source=h2 target=subsection reason=relative-heading-rank; base=h2

\subsection{第五十八章}

在此事之后，克尔苏笔下的这位犹太人，代替福音中所提的贤士，说：\Quote{耶稣说迦勒底人是被诱导在他出生时前来，在他还是婴儿时就拜他为神，并将此事告知分封侯黑落德；后者就派人杀掉了所有约在同一时间出生的婴孩，以为这样就能保证他在众人中死去；他这样做是出于害怕，如果耶稣活到足够年纪，就会得到王位。}请看此例，此人错误百出，既不能区分贤士与迦勒底人，也看不出他们所称述的不同，因而歪曲了福音叙述。此外，我不知道他为何对引导贤士前来的原因保持沉默，也没有按照圣经记载说是一颗他们在东方看见的星。现在我们来看对此应作何回答。我们相信东方所见的那颗星是一颗新星，不同于任何其他已知的行星体，无论是天上的星辰还是低处的星体，但具有那些有时出现的天体之性质，例如彗星，或那些类似木梁、胡须、酒缸的流星，或希腊人习惯用来描述其不同外貌的其他名称。我们以下列方式确立我们的立场。\par

% structure: p0125 source=h2 target=subsection reason=relative-heading-rank; base=h2

\subsection{第五十九章}

据观察，在重大事件发生和世间巨大变革时，常有这类星星出现，预示王朝的更迭或战争的爆发，或可能引起地上骚乱的情势。但我们在斯多葛派的\Person{凯雷蒙}所著的\WorkTitle{论彗星}中读到，有时在\Emph{好事}将要发生时，彗星也会出现；他列举了这样的实例。那么，如果在新的王朝开始或其他重要事件发生时，有所谓的彗星或任何类似的天体出现，那么为何在祂的诞生——这位要将新教理带给人类，并将其教导不只向犹太人，也向希腊人和许多野蛮民族宣告者——时，出现一颗星星就值得惊奇呢？现在我要说，关于彗星，并无流行的预言说某颗彗星要在某个王国或某个时期出现；但关于\Person{耶稣}诞生时出现一颗星星，却有\Person{巴郎}经由\Person{梅瑟}所记的如下预言：\BibleQuote{由雅各伯将出现一颗星，由以色列将兴起一位首领。}现在，如果有必要查考关于贤士和\Person{耶稣}诞生时星星出现的叙述，以下是我们必须说的，部分是对希腊人的回答，部分是对犹太人的回答。\par

% structure: p0127 source=h2 target=subsection reason=relative-heading-rank; base=h2

\subsection{第六十章}

对希腊人，我须说：\Person{贤士}们既与邪魔亲密，按他们的知识和愿望所及，为达到某些目的而呼求它们，所成就的结果，仅不超出邪魔及其呼求咒语的超自然力量和强度所能成就的范围；但若出现更大的天主性的显示，则邪魔的力量便被推翻，因为它们不能抗拒天主性的光明。因此，很可能，在耶稣诞生时，\Quote{一大队天军}——如\Person{路加}所记，而我相信——\Quote{赞美天主说：\InnerQuote{天主在天受光荣，主爱的人在世享平安。}}邪魔便因此变得软弱，失去力量，其巫术的虚假得以显明，其能力被打破；这推翻不仅因天使为耶稣诞生而降世，也因耶稣自身的能力及他与生俱来的天主性。\Person{贤士}们想要产生惯常的结果——即他们以前藉某些咒语和巫术所行的事——却发觉失效，便想探知失败的原因，猜测缘由重大；及见天上有神迹，便欲明了其意义。因此我认为，他们既拥有\Person{巴郎}的预言（\Person{梅瑟}亦记载了这些预言，因\Person{巴郎}以这类预测著称），并在其中发现关于星辰的预言以及这些话：\Quote{我将他指给他看，但不是现在；我认为他有福，虽然他不在近处。}他们便猜测，那与星辰一同被预言出现的人，实际上已经来到了世界；并已断定他的能力超越一切魔鬼、一切寻常显现和权能，便决定向他致敬。于是他们来到\Place{犹太}，确信某位君王已诞生，却不知他将统治什么国度，也不知他出生之地；他们带着礼物，献给那一位其本性——如果我可以这么说——兼具天主和凡人：黄金，如献于君王；没药，如献于必死者；乳香，如献于天主。他们是在得知祂出生地之后才献上这些礼物的。但祂既是天主、人类救主，远超所有服务于人的天使，便有一位天使酬报\Person{贤士}们对他的敬拜之虔诚，指示他们不要回到\Person{黑落德}那里，而从另一条路返回家乡。\par

% structure: p0129 source=h2 target=subsection reason=relative-heading-rank; base=h2

\subsection{第六十一章}

\Person{黑落德}谋害这孩子（虽然\Person{克尔苏}笔下的犹太人不相信这事真正发生），并不足为奇。因为罪恶在某种程度上是盲目的，它想击败命运，仿佛它比命运更强。\Person{黑落德}正是处于这种状态：他既相信一位犹太人的君王已经诞生，却又抱有与这种信念相矛盾的目的；他没有看到这孩子要么确实是一位君王，将登基为王，要么他不会成为君王，那么杀他也就毫无意义。因此他想杀死他，他的心灵因罪恶而受着对立情感的搅扰，并被那盲目而邪恶的魔鬼所煽动，这魔鬼从一开始就谋害救主，认为他是一位并将成为一位伟大的人物。然而，一位天使察知事态的进程，便指示\Person{若瑟}——尽管\Person{克尔苏}可能不相信——要他带着孩子和他的母亲退往\Place{埃及}，而\Person{黑落德}杀害了\Place{白冷}及其周围边界的所有婴孩，希望以此也杀死那位生来就是犹太人之王的人。因为他没有看见那不寐的守护力量环绕着那些配得保护与保存以造福人类的人，耶稣是其中第一位，在尊荣和卓越上超越所有其他人，他确实要成为一位君王，但并非\Person{黑落德}所认为的那种君王，而是天主赐予国度的那一种——即为了那些将受他统治的人的利益，他要赐予他的臣民并非普通而不重要的福惠（如果可以这么说的话），而是要训练他们，使他们服从真正来自天主的法律。耶稣深知这一点，并否认他是群众所期待的那种君王，反而声明他国度的卓越性，他说：\Quote{如果我的国属于这世界，我的臣仆必要为我作战，使我不被交给犹太人；但如今我的国不属于这世界。}现在，如果\Person{克尔苏}看到了这一点，他就不会说：\Quote{但如果这事发生是为了使你长大成人后不代替他作王，那么你既已长大，为何不成为君王，反而你，天主子，如此卑微地流浪，因恐惧而躲藏，过着颠沛流离的痛苦生活呢？}其实，避免使自己暴露于危险之中并非不光彩，而是谨慎防范，这样做不是出于怕死，而是出于渴望通过留在世上使他人受益，直到那担任人性者为人类利益而死的时候来到。这一点，凡思考耶稣为使人得益而死者——我们在前文中已尽力论述——都会明白。\par

% structure: p0131 source=h2 target=subsection reason=relative-heading-rank; base=h2

\subsection{第六十二章}

在发表这些言论之后，显示出他对宗徒人数的无知，他接着这样说：\Quote{耶稣聚集了十或十一个臭名昭著的人，是最恶的名声狼藉的税吏和船夫，与他们一同到处逃窜，并以可耻和纠缠的方式谋生。}让我们尽我们所能看一看这种说法中有什么真理。对于我们所有拥有福音叙述的人——\Person{刻尔苏}似乎甚至没有读过这些叙述——明显的是，\Person{耶稣}拣选了十二宗徒，其中只有\Person{玛窦}是税吏；当他不加区分地称他们为船夫时，他可能指的是\Person{雅各伯}和\Person{若望}，因为他们离开了船和父亲\Person{载伯德}，跟随了\Person{耶稣}；因为\Person{伯多禄}和他的兄弟\Person{安德肋}，他们用网来获取必要的生活所需，必须被归类为船夫，而是正如圣经所描述的，他们是渔夫。\Person{肋伯斯}也是\Person{耶稣}的跟随者，可能曾是税吏；但他不属于宗徒之列，除非按照\BibleBook{马尔谷}{福音}某个抄本中的说法。我们尚未确定其余门徒在成为\Person{耶稣}门徒之前的职业，他们靠什么谋生。因此，我针对上述这类说法断言，所有能够进行明智而公正考察\Person{耶稣}宗徒历史的人都很清楚，正是借着天主的能力，这些人教导了基督信仰，并成功地带领其他人接受了天主的圣言。因为使他们说服听众的有效原因，并不是任何言谈的能力，也不是按照希腊辩证术或修辞学来有条理地安排他们的信息。是的，我认为，如果\Person{耶稣}拣选了一些按群众看法是明智的人，这些人既适合思考也适合说话以取悦他们，并用这样的人作为祂教义的仆役，那么祂就最应被怀疑使用了技巧，就像那些某些学派的领袖的哲学家一样，因此关于祂道理神性的许诺就不会显现出来；因为如果教义和宣讲在于有说服力的言辞和安排，那么信德，就像世界上的哲学家对他们的看法一样，就会是出于人的智慧，而不是出于天主的能力。现在，有谁看到渔夫和税吏——他们没有获得哪怕最基础的知识（正如福音所描述的，\Person{刻尔苏}相信他们在此事上说真话，因为他们记录了自己的无知）——不仅大胆地在犹太人中讲论对\Person{耶稣}的信德，而且还成功地在外邦人中宣讲祂，难道不会追问他们从哪里获得这种说服力，因为他们肯定不是普通人惯用的方法？谁会不说，\Person{耶稣}在祂宗徒的历史中所应许的\BibleQuote{跟随我，我要使你们成为捕人的渔夫}，是借着一种神圣的能力应验了？\Person{保禄}也说，正如我们刚才提到的，称赞这一点说：\BibleQuote{我的言论和我的宣讲不是用人的智慧动听的言词，而是用圣神和能力的明证，为使你们的信德不凭人的智慧，而凭天主的能力。}因为按照先知们预言福音宣讲的预言，\BibleQuote{主以大能将话语赐给传布它的人，甚至所爱者的权能之王}，为使那预言得以应验：\BibleQuote{祂的话必急速奔馳。}我们看见\BibleQuote{耶稣宗徒的声音已传遍天下，他们的言语已达到地极。}因此，凡听见被大能地宣讲的圣言的人，都充满了能力，这能力既表现在他们的心性和生活上，也表现在为真理奋斗甚至死亡上；而有些人则完全是空洞的，尽管他们声称相信天主借着\Person{耶稣}，因为他们没有任何神圣的能力，只有表面上皈依天主的圣言。虽然我先前已经提过救主所讲的一句福音的话，但我仍要在此再引用它，因为它适合当前的情形：它证实了救主对祂福音宣讲的神性预见，以及祂圣言的能力，这能力不用教师的帮助就能战胜那些在神圣能力伴随的劝说之下屈服的人；所引用的\Person{耶稣}的话是：\BibleQuote{庄稼多，工人少；因此你们要祈求庄稼的主，派遣工人到祂的庄稼里去。}\par

% structure: p0133 source=h2 target=subsection reason=relative-heading-rank; base=h2

\subsection{第六十三章}

由于\Person{凯尔苏斯}称耶稣的宗徒为臭名昭著的人，说他们是劣迹斑斑的税吏和渔夫，关于这一指责，我们必须指出，他似乎是为了控诉基督教，才只在他乐意的时候相信福音记载，并表达他对这些记载的不信，以免被迫承认同一批书中所述的天主性的显示；然而，一个人若看到指导作者们的真理之神，就应该从他们对较次要事物的叙述中，也相信对神圣事物的叙述。\newline{}现在，在\WorkTitle{巴尔纳伯通函}中——凯尔苏斯或许正是从此处取得了宗徒是臭名昭著的恶人这一说法——记载着：\Quote{耶稣拣选了自己的宗徒，作为比所有其他作恶者更有罪的人。}\newline{}而在\BibleBook{路加}{福音}中，\Person{伯多禄}对耶稣说：\BibleQuote{主，离开我，因为我是个罪人。}\newline{}此外，\Person{保禄}自己后来也成了耶稣的宗徒，他在写给\Person{弟茂德}的\BibleBook{弟茂德}{前书}中写道：\BibleQuote{这话是确实的：耶稣基督到世界来，是为拯救罪人，而我是其中最大的。}\newline{}我不知道\Person{凯尔苏斯}怎么会忘记或没想到要谈论保禄——那位在耶稣之后成为基督内各教会创始人的宗徒。\newline{}他大概看出，任何关于那位宗徒的话都需要解释，以符合这一事实：保禄在成为天主的教会的迫害者、信徒的刻毒反对者，甚至将耶稣的门徒交付处死之后，竟发生了如此巨大的改变，以致他从\Place{耶路撒冷}直到\Place{依里黎谷}宣讲耶稣的福音，并切望在不必建造在他人根基上的地方，即在完全没有宣讲过天主在基督里的福音之处，传播这喜讯。\newline{}因此，若耶稣为向人类显示祂所拥有的治愈灵魂的权能，而拣选了臭名昭著、邪恶的人，并将他们提升到如此卓越的道德境界，以致他们成为所有经由他们而归化于基督福音之人的最纯粹德行的典范，这又何荒谬之有呢？\par

% structure: p0135 source=h2 target=subsection reason=relative-heading-rank; base=h2

\subsection{第六十四章}

但如果我们用他们归化前的生活来责备那些被呼召的人，那么我们也有理由责备\Person{斐多}，即使在他成为哲学家之后；因为据史书记载，他是被\Person{苏格拉底}从声名狼藉的场所引向哲学研究的。不，连\Person{色诺克拉底}的继承人\Person{波勒蒙}的放荡生活也会成为哲学的耻辱；而即使在这样的事例中，我们也应当视之为一种赞誉：因为推理能力凭借这些人的说服力量，能使那些先前被恶习缠身的人远离此等恶习。\newline{}在希腊人中，只有一位斐多——我不知道是否有第二位——和一位波勒蒙，在过着放荡而极其邪恶的生活后，投身于哲学；但在耶稣那里，不仅在我们所说的那个时候有十二位门徒，而且历代有更多人，他们组成一群节制之人，这样谈论他们的从前的生活：\BibleQuote{因为我们从前也是昏愚的，悖逆的，被迷惑的，服侍各种私欲和逸乐，生活在恶毒和嫉妒中，是可憎恶的，彼此仇恨。但当我们救主天主的良善和慈爱显现时，藉着重生之洗和圣神的更新，祂厚厚地倾注在我们身上的}我们成了现在这样。\newline{}因为\BibleQuote{天主发出祂的圣言，治好了他们，拯救他们出离坟墓}，正如先知在\BibleBook{}{圣咏集}中所教导的。\newline{}除了已经说的，我还要补充以下内容：\Person{克吕西波}在他的著作\WorkTitle{情欲的治疗}中，在努力抑制人类灵魂的情欲时，并不妄图确定哪些意见是正确的，而是说根据各宗派的原理，那些被情欲支配的人应得到治疗，他接着说：\Quote{倘若快乐是目的，那么就必须藉快乐来治愈情欲；倘若有三种主要的善，照样，根据这一学说，那些被情欲支配的人也必须以同样的方式摆脱它们；}然而，基督教的攻击者们没有看到，有多少人的情欲受到了抑制，恶习的洪流被阻止，野蛮的习性被福音所缓和。因此，那些总是夸耀自己热心公益的人，理应公开感谢那用新方法引导人们放弃许多恶习的教义，至少向它表示赞赏：即便它不是真理，至少对人类有益。\par

% structure: p0137 source=h2 target=subsection reason=relative-heading-rank; base=h2

\subsection{第六十五章}

因为耶稣教导门徒不可鲁莽，给他们训诫说：\Quote{如果他们在这城里迫害你们，就逃到另一城；如果在那城里迫害你们，就再逃到第三城。}祂不但这样教导，还以一致生活的榜样加以印证，行事不鲁莽、不合时宜或没有充分理由地使自己陷于危险。瑟尔苏借此机会提出恶意的诽谤指控——他引入的一个犹太人对耶稣说：\Quote{你和你的门徒一起到各处躲藏。}我们说，类似这样被用来诽谤耶稣及其门徒的行为，在亚里士多德的生平中也记录着。这位哲学家看到法庭即将因他不敬神的罪名传唤他，因他某些哲学观点被雅典人视为不敬神，便离开雅典，将学校设在卡尔基斯，向他的朋友们辩护自己的做法说：\Quote{我们离开雅典吧，免得给雅典人提供第二次犯罪的机会，就像从前对待苏格拉底那样，阻止他们对哲学犯下第二次不敬神的行为。}他还说：\Quote{耶稣与门徒四处走动，以可耻和强求的方式谋生。}让他指出他们谋生方式中可耻和强求之处何在。因为福音中记载，有些妇女曾因疾病被治愈，其中包括苏撒纳，她们用自己的财产供给门徒的生活所需。在哲学家中，有谁在为熟人服务时不习惯从他们那里接受生活必需品呢？难道只有他们这样做才是适当和得体的，而当耶稣的门徒也这样做时，瑟尔苏就指控他们以可耻的强求谋生吗？\par

% structure: p0139 source=h2 target=subsection reason=relative-heading-rank; base=h2

\subsection{第六十六章}

此外，瑟尔苏的这个犹太人随后对耶稣说：\Quote{再者，你还在婴孩时期，有什么必要被带到埃及去？是为了逃避被杀吗？但一位天主不应该害怕死亡；然而却有一位天使从天降下，命令你和你的朋友们逃跑，以免被逮住杀死！难道那已经为你派了两位天使的伟大天主，不能保全你——祂的独生子——在那里安全吗？}从这些话看来，瑟尔苏似乎认为耶稣的人性与灵魂中没有任何神性元素，甚至祂的身体也不像荷马神话中所描述的那样；并且他还嘲笑耶稣在十字架上流出的血，补充说那并不是\par

\Quote{灵气，正如流淌在蒙福诸神血管中的那种。}\par

我们现在，相信耶稣自己关于其神性所说的话：\Quote{我是道路、真理、生命}，以及祂关于自己穿着人身所说的话：\Quote{现在你们却想杀我，一个把真理告诉你们的人}，从而得出结论：祂是一种复合的存在。因此，祂作为一个人为在世上的寄居作准备，就不应不合时宜地暴露于死亡的危险中。同样，祂有必要由父母带走，遵从来自天上的一位天使的指示，天使向他们传达了天主的旨意，第一次说：\Quote{若瑟，达味之子，不要怕娶你的妻子玛利亚，因为她所怀孕的是出于圣神}；第二次说：\Quote{起来，带着孩子和祂的母亲，逃往埃及，住在那里，直到我告诉你；因为黑落德要寻找孩子，把祂杀死。}现在，这些话所记载的在我看来一点也不奇怪。因为在两段圣经中都说明天使是在梦中说这些话的；而且在梦中，某些人可能被指示要做某些事，这是许多人都常经历的事——脑海中产生的印象要么来自天使，要么来自其他事物。那么，相信那位已经降生成人的，也会受到人的引导来避开危险，这有何荒谬呢？并非因为不可能有其他情况，而是从道德适宜性来说，应该使用方法来确保耶稣的安全。当然，婴孩耶稣逃脱黑落德的陷阱，与父母住在埃及直到阴谋者死去，这比天主的眷顾阻止黑落德自由意志杀死孩子，或者使用神话中冥王哈得斯的隐形头盔，或对耶稣做类似的事情，或者来杀祂的人像索多玛人那样被击打失明，要好得多。因为以一种非常神奇且不必要公开的方式帮助祂，对祂毫无益处。祂希望表明，作为一个天主为之作证的人，在人们肉眼所见的形体之内，祂拥有更高层次的神性——那正是圣子——天主的圣言——天主的德能和天主的智慧——那位被称为基督的。但现在不是讨论降生成人的耶稣复合本性的适当时机；对信徒来说，这样的探讨可以说是某种私下的问题。\par

% structure: p0143 source=h2 target=subsection reason=relative-heading-rank; base=h2

\subsection{第六十七章}

在这之后，塞尔苏斯笔下的这个犹太人，如同一位爱好学识且熟谙希腊文学的希腊人，继续说道：\Quote{古老的寓言神话，将神性起源归于\Person{珀耳修斯}、\Person{安菲翁}、\Person{埃阿科斯}和\Person{米诺斯}，我们并不相信。然而，为了不让这些故事显得不足为信，他们将这些人物的事迹描绘得伟大而奇妙，真正超出人力所及；但你们在言行上，做了什么崇高或奇妙的事呢？你们没有向我们显明什么，尽管他们在圣殿中挑战你，要你显示一个确凿的记号，证明你是天主子。}对此，我们不得不说：让希腊人向我们指出，在所列举的那些人中，有谁的事迹具有惠及后世的效用与光辉，并如此伟大，以致让人相信那些将他们描写为神裔的寓言？但这些希腊人无法向我们提供任何关于他们所说之人的事物，甚至远远不及耶稣所行的；除非他们把我们带回到他们的寓言和历史中，希望我们毫无根据地相信它们，同时否认福音记载，尽管后者有最清晰的证据。因为我们断言，整个人类居住的世界都包含了耶稣之事的证据，即那些因祂而藉由从无数罪恶中悔改皈依之人所建立的天主教会。耶稣的名仍然能消除人心的困扰，驱逐魔鬼，除去疾病；并在那些不因生计或满足任何世俗需要而假装基督徒，而是真诚接受关于天主和基督以及将来审判之教义的人身上，产生奇妙的温顺精神、彻底的品性改变，以及仁爱、良善和温和。\par

% structure: p0145 source=h2 target=subsection reason=relative-heading-rank; base=h2

\subsection{第六十八章}

此后，塞尔苏斯疑心有人会提出耶稣所行的大事（我们已从众多事迹中略举几例），便假装承认关于祂的医治、复活、以少许饼使群众饱食并有余，以及其他塞尔苏斯认为门徒们以惊奇笔调记载的故事，可能为真。他接着说道：\Quote{好吧，我们相信这些确实为你所行。}但随即将其与变戏法者的把戏相比，这些人自称能行更奇妙的事，并与那些受过埃及人教导者所表演的技艺相比较，他们在市集中央为几个小钱就能传授最受尊崇的技艺：驱逐人身上的魔鬼、祛除疾病、召唤英雄的灵魂、展示并不存在的丰盛宴席、餐桌、盘盏和美味，以及使没有生命的动物活动起来，如同活物一般。他问道：\Quote{那么，既然这些人能行这样的奇事，我们是否必须断定他们是\InnerQuote{天主子}？还是必须承认这些是恶人受邪灵驱使的行径？}可见，通过这些话，他仿佛承认了魔法的存在。但我不知道他是否就是那位写了多卷反对魔法书籍的人。不过，为助成其论点，他拿耶稣的奇迹与魔法产生的结果相比。若耶稣像那些施行魔法的人一样，仅为炫耀而行其事，两者确有相似之处；但如今，没有一个变戏法者藉其表演邀请观众改过迁善，或训练那些因所见而惊异的人敬畏天主，也没有人试图说服他们如此生活，如同那些将要被天主成义的人。变戏法者不做这些事，因为他们既无能力也无意愿，更无心思去操劳人的改造，因为他们自己的生活充满了最严重、最臭名昭著的罪恶。然而，那藉其所行的奇迹，引导看见这些卓越成果的人着手改造自己品行的人，岂不应当向祂真正的门徒，也向其他人，显明自己是最有德行之生活的模范，好让祂的门徒致力于教导人天主的心意，而其他人则藉祂的言行（比藉祂的奇迹更彻底）学习如何指导自己的生活，使他们在一切行为中常以普世天主的喜悦为依归？如果耶稣的生活是如此，人怎能合理地将祂与骗子之流相比，而不相反地，按照应许，相信祂是天主，以人形显现来造福我们的种族呢？\par

% structure: p0147 source=h2 target=subsection reason=relative-heading-rank; base=h2

\subsection{第六十九章}

之后，\Person{策尔苏斯}将基督教教义与某个异端派别的意见混淆在一起，并把这些作为适用于所有相信天主圣言者的指控提出来，说：\Quote{像你这样的身体不可能属于天主。} 对此我们必须回答说：\Person{耶稣}进入世界时，作为由女人所生者，取了一个人的身体，一个能承受自然死亡的身体。为此，除其他理由外，我们也说祂是一位伟大的角力者；由于祂的人性身体，祂在各方面都像其他人一样受过试探，但不像人那样有罪作为后果，而是完全没有罪。因为我们清楚地知道：\BibleQuote{祂没有犯过罪，口中也找不到诡诈；作为不知罪的一位}，天主把祂当作纯洁的祭品交付出来，为所有犯了罪的人。然后\Person{策尔苏斯}说：\Quote{神的身体不会像你一样，哦\Person{耶稣}，如此生成。} 祂还看到，如果祂像经上所记那样是出生的，祂的身体在某种程度上或许比众人的更神圣，并在某种意义上是一个神的身体。但他不相信祂因圣神受孕的叙述，而相信祂是由一个名叫\Person{潘特拉}的人所生，此人玷污了童贞玛利亚，\Quote{因为神的身体不会像你一样如此生成。} 但我们在前文更详细地谈论了这些事。\par

% structure: p0149 source=h2 target=subsection reason=relative-heading-rank; base=h2

\subsection{第七十章}

此外，他断言：\Quote{神的身体不会用这样的食物（像\Person{耶稣}那样）来滋养，} 因为他能从福音叙述中证明祂既吃了食物，而且是特定种类的食物。好吧，就算如此。让他断言祂同门徒吃了逾越节晚餐，当时祂不仅说了：\BibleQuote{我渴望又渴望，同你们吃这一次逾越节晚餐，} 而且确实吃了同一筵席。也让他说，祂在\Place{雅各伯井}旁感到口渴，并喝了井里的水。这些事实在哪些方面与我们关于祂身体本质的论述相矛盾呢？此外，毫无疑问，祂复活后吃了一片鱼；因为按照我们的观点，祂取得了一个（真实的）身体，如同由女人所生那样。\Quote{但是，} \Person{策尔苏斯}反驳说，\Quote{神的身体不使用像\Person{耶稣}那样的声音，也不使用像祂那样的说服方法。} 这些确实是微不足道、完全可鄙的反对意见。因为我们对他的答复将是：在希腊人中被视为神的，即\Person{皮提亚和狄迪玛的阿波罗}，在\Place{德尔斐}为自己的皮提亚女祭司，以及在\Place{米利都}为女先知使用这样的声音；然而，无论是皮提亚还是狄迪玛的阿波罗，都未被希腊人指控不是神，也没有任何其他其崇拜固定在一处的希腊神祇受到指控。而且，神使用一种因带有权能而发出、在听众心中产生一种难以言喻的说服力的声音，这当然要好得多。\par

% structure: p0151 source=h2 target=subsection reason=relative-heading-rank; base=h2

\subsection{第七十一章}

继续对\Person{耶稣}大加辱骂，称祂因不敬和邪恶意见而可以说是被天主所恨，他断言：\Quote{祂的这些信条是一个邪恶且为天主所恨的巫师的信条。} 然而，如果对名称与事实加以恰当考察，就会发现人不可能被天主所恨，因为天主爱一切存在之物，且\BibleQuote{不恨祂所造的任何东西，} 因为祂不是在仇恨的精神中创造任何事物的。如果先知书中的某些表达给人这种印象，那么应当根据圣经运用这种语言对待天主的一般原则来解释，好像天主受人情感影响一样。但是，对于这样一个声称提出可信陈述、却认为自己必须对\Person{耶稣}使用诽谤和诋毁（好像祂是一个邪恶巫师）的人，还需要作何答复呢？这并非寻求证明自己案情之人的做法，而是一个处于无知和不哲学心态之人的做法，因为正确的做法是陈述案情，并坦诚地加以探究；尽自己所能，提出与之相关的想法。但是，既然\Person{策尔苏斯}笔下的犹太人已藉上述言论结束了对\Person{耶稣}的指控，那么我们也在这里结束我们第一本书中对他答复的内容。如果天主赐予那摧毁一切虚假的真理之恩，符合祈祷中的话：\BibleQuote{以你的真理消灭他们，} 我们将在接下来的内容中开始考虑犹太人的第二次出场，按照\Person{策尔苏斯}的描绘，他对那些已成为\Person{耶稣}信徒的人说话。\par

\SourceRecord{Translated by Frederick Crombie.；Ante-Nicene Fathers；Vol. 4.；Edited by Alexander Roberts, James Donaldson, and A. Cleveland Coxe.；Buffalo, NY: Christian Literature Publishing Co.,；1885.；<http://www.newadvent.org/fathers/04161.htm>.}

% structure: p0001 source=h1 target=section reason=page-title

\section{驳\Person{塞尔索}，第二卷}

% structure: p0002 source=h2 target=subsection reason=relative-heading-rank; base=h2

\subsection{第一章}

我们答复\Person{塞尔索}所著名为\WorkTitle{真确论述}的论著的第一卷，以那犹太人向\Person{耶稣}讲话的情景作结，现已足够长度；我们打算将当前部分作为对他所提出的指控的回应，这些指控是针对那些从犹太教改信基督信仰的人。我们首先提请注意这一特殊问题：即\Person{塞尔索}既然决定在他的书中引入一些人物，为何不让那犹太人对改信外教的人讲话，反而对改信犹太教的人讲话？但这位自称拥有普世知识的人，大概不懂得这类叙述中的适当性；所以让我们继续看他向改信犹太教的人所说的话。他断言：\Quote{他们离弃了祖传的法律，因为他们的心智被\Person{耶稣}掳获；他们被极其荒谬地欺骗了，以致成为另一个名字和另一种生活方式的叛徒。}在这里他没有注意到，犹太改信者并没有离弃祖传的法律，因为他们按照其规诫生活，从字面意思来说，他们的名字源自法律的贫穷；因为\OriginalTerm{厄彼雍}{Ebion}在犹太人中意为\Quote{贫穷}，那些接受\Person{耶稣}为基督的犹太人被称为\OriginalTerm{厄彼雍派}{Ebionites}。不仅如此，\Person{伯多禄}本人似乎也遵守了\Person{梅瑟}法律所规定的犹太规诫相当长一段时间，尚未从\Person{耶稣}那里学到从按字面解释的法律上升到按精神解释的法律——我们从\BibleBook{宗徒}{大事录}得知这一事实。因为当天主的使者向\Person{科尔乃略}显现后，提示他\BibleQuote{派人到\Place{约培}去，请那称为\Person{伯多禄}的\Person{西满}来}，\Person{伯多禄}\BibleQuote{约在第六时辰上屋顶祈祷。他很饿，想吃东西；但人们在预备时，他陷入了神魂超拔，看见天开了，有一个器具降下，好像一块大布，四角系着，缒到地上，里面有各种四足兽、地上的爬虫和天空的飞鸟。有声音对他说：\Person{伯多禄}，起来，宰了吃！\Person{伯多禄}却说：主啊，绝对不可！因为我从未吃过污秽不洁之物。第二次又有声音对他说：天主所洁净的，你不可称为污秽。}现在请注意，通过这个例子，\Person{伯多禄}被描绘成仍然遵守犹太人关于洁净和不洁净动物的习俗。从后续的叙述可以看出，他作为犹太人，按照他们的传统生活，并轻视犹太教以外的人，需要异象来引导他将信仰之言传给\Person{科尔乃略}（并非按肉身是以色列人）以及与他在一起的人。此外，在\BibleBook{}{迦拉达书}中，\Person{保禄}说，\Person{伯多禄}因害怕犹太人，在\Person{雅各伯}到来后就停止与外邦人一起吃饭，\BibleQuote{他因怕那些受割礼的人，就退避，疏远了他们}；其余的犹太人和\Person{巴尔纳伯}也随同他这样做。那些被派去服务受割礼的人，当然不应该停止遵守犹太习俗，当那些\BibleQuote{被视为柱石}的人与\Person{保禄}和\Person{巴尔纳伯}行右手相交之礼，以便他们自己致力于受割礼的人，而后者则可向外邦人宣讲。我为何要提那些向外邦人宣讲的人退避并离开外邦人呢？就连\Person{保禄}自己也\BibleQuote{对犹太人，我就成为犹太人，为要赢得犹太人}。因此，在\BibleBook{宗徒}{大事录}中也记载他甚至献祭在祭台上，以使犹太人相信他并非离弃他们的法律。现在，如果\Person{塞尔索}知道所有这些情况，他就不会让那犹太人对改信犹太教的人说出这样的话：\Quote{同胞们，是什么促使你们离弃祖传的法律，让自己的心智被我们刚才与之交谈的那人掳获，以致被最荒谬地欺骗，从而离开我们成为另一个名字和另一种生活方式的叛徒？}\par

% structure: p0004 source=h2 target=subsection reason=relative-heading-rank; base=h2

\subsection{第二章}

既然我们正在讨论伯多禄以及向受割损者传教的老师，我认为在此引用一段取自\BibleBook{若望}{福音}的耶稣言论并加以解释，并非不合时宜。因为那里记载耶稣说：\Quote{我还有许多事要告诉你们，但你们现在负担不了。诚然，当祂，那\OriginalTerm{真理之神}{Spirit of truth}，来到时，祂要引导你们进入一切真理：因为祂不凭自己说话；但无论祂听见什么，祂将说出来。}当我们探究耶稣在那段经文中要对门徒说，但他们当时未能承担的\Quote{许多事}是什么时，我必须指出：很可能因为宗徒是犹太人，自幼受梅瑟法律的字句训练，祂不能向他们说明什么是真正的法律，以及犹太敬拜如何只是天上事物的模型和影子，未来福分又如何通过关于食物、饮料、节期、新月和安息日的诫命预先被影射。这些都是祂必须向他们解释的许多课题；但祂看到，要从心中根除几乎与生俱来、并在成长过程中根深蒂固、且使接受者相信它们是神圣的并推翻它们即为不敬的那些见解，是一项极其困难的任务；要借着基督教教义的优越性，即真理，以说服听者的方式，证明这些见解不过是\Quote{损失和废物}，祂便将此事推迟到将来——即祂受难和复活之后的时期。因为不合时宜地帮助那些尚未能接受的人，可能会颠覆他们已经形成的关于耶稣是基督和永生天主之子的观念。请看看，对于\Quote{我还有许多事要告诉你们，但你们现在负担不了}这一说法，是否有些充分的理由？因为在法律中有许多点需要以属灵意义加以解释和澄清，而这些门徒由于生在犹太人中间并在他们中间长大，在某种程度上无法承担。此外，我认为，因为这些礼仪是预像，而真理是圣神要教导他们的，所以加了这些话：\Quote{当那真理之神来到时，祂要引导你们进入一切真理。}好像祂在说：进入关于那些事物的全部真理，这些事物对你们而言只是预像，而你们却相信构成了你们向天主所行的真正敬拜。于是，按照耶稣的应许，真理之神来到\Person{伯多禄}那里，就地上的四足走兽、爬虫和飞鸟对他说：\Quote{\Person{伯多禄}，起来，宰了吃！}当\Person{伯多禄}仍处于迷信无知时，圣神来对他说。他回答天主的命令说：\Quote{主，绝对不可！因为我从未吃过任何凡俗或不洁之物。}然而，圣神通过说\Quote{天主所洁净的，你不可称为凡俗}教导他关于食物的真正属灵意义。因此，在那异象之后，引导\Person{伯多禄}进入一切真理的真理之神，告诉了他许多在耶稣还带着肉身与他们同在时他无法承担的事。但我将另有机会解释那些与梅瑟法律字面理解相关的问题。\par

% structure: p0006 source=h2 target=subsection reason=relative-heading-rank; base=h2

\subsection{第三章}

然而，我们当前的目标是揭露刻尔苏斯的无知，他让他笔下的犹太人以下列方式对他的同族和入教的以色列人说话：\Quote{是什么促使你们离弃了你们祖先的法律？}等等。那么，那些惯于责备不听从法律命令的人，说：\Quote{告诉我，你这诵读法律的人，你不听见法律吗？因为经上记载，亚伯拉罕有两个儿子……等等，直到\InnerQuote{这些都是\OriginalTerm{寓意}{allegories}}，}他们怎么会离弃了祖先的法律呢？那些常常用这样的话述说祖先习俗的人，又怎么会离弃了祖先的法律呢？\Quote{法律不也是说这些事吗？因为梅瑟法律上记载：\InnerQuote{牛在场上踹谷时，不可笼住它的嘴。}天主岂是关心牛吗？还是全为我们说的？的确，是为我们写的……}等等。那么，犹太人在这些问题上的推理是多么混乱（尽管他本可以更有力地说话），当他说：\Quote{你们中有些人借口解释和寓意离弃了我们祖先的习俗；有些人尽管佯装以属灵方式解释它们，却仍然遵守我们祖先的习俗；还有一些人，没有这种解释，却愿意接受耶稣作为预言的应验，并按照祖先的习俗遵守梅瑟法律，好像在其中拥有圣神的全部心意。}那么，刻尔苏斯怎么会在此处如此清楚地看到这些事，而在其后来的作品中，他却提到某些与耶稣的教义完全无关且不虔敬的异端，甚至还有一些完全撇开造物主的异端，而且他似乎不知道有皈依基督教的以色列人，他们没有离弃他们祖先的法律？他不是本着真理的精神来审查一切，并接受任何他认为有用的事物；而是怀着敌人的精神来编写这些陈述，并且一听到就企图推翻一切。\par

% structure: p0008 source=h2 target=subsection reason=relative-heading-rank; base=h2

\subsection{第四章}

那位犹太人继续对他本族中的皈依者这样说道：\Quote{昨天和前天，当我们惩罚那迷惑你们的人时，你们背离了祖传的法律。}他这番话（正如我们刚才所证明的）丝毫不能表明对真理的精确认识。但他后面所说的话似乎有些力量，他说：\Quote{你们怎么从我们的敬拜中吸取你们教义的开端，而当你们有所进展时，却又对它表示不敬，尽管你们除了我们的法律之外，拿不出别的作为你们学说的基础？}当然，基督宗教的入门是通过梅瑟敬拜和先知著作；入门之后，正是在对这些著作的解释和阐明中取得进步，而那些入门者则深入探究\Quote{那照着启示的奥秘，自创世以来隐藏着，如今却在先知的书信中}并藉我们的主耶稣基督的显现而显明出来。但那些在基督信仰知识上长进的人，并不像你所断言的那样，对法律中所记载的东西表示不敬。相反，他们赋予它们更高的尊荣，表明这些著作中含有何等深邃的智慧和奥秘的理由，而这些是犹太人未能完全领悟的，他们只作肤浅的理解，甚至在一定程度上视之为神话。我们的体系——即福音——以法律为基础，这有什么荒谬呢？因为主耶稣本人曾对那些不愿信祂的人说：\BibleQuote{你们如果信梅瑟，也必信我，因为他曾论及我。如果你们不信他所写的，怎能信我的话呢？}而且，一位福音作者——马尔谷——说：\BibleQuote{耶稣基督福音的开始，正如先知依撒意亚书上所记：\InnerQuote{看，我派遣我的使者在你面前，他要在你前面预备你的道路。}}这表明福音的开端与犹太人的著作相连。那么，凯尔索的犹太人所提出的反对意见有什么力量呢？他说：\Quote{如果有人预言天主子要降临人间，那他就是我们的先知之一，是我们天主的先知。}或者说，若翰为耶稣施洗是犹太人，这怎能算是对基督宗教的指控？因为尽管他是犹太人，但并不因此每个信徒——无论是从外邦还是从犹太教皈依的——都必须对梅瑟法律字句照办。\par

% structure: p0010 source=h2 target=subsection reason=relative-heading-rank; base=h2

\subsection{第五章}

在这些事之后，虽然凯尔索在关于耶稣的论述中变得重复，第二次说\Quote{祂因自己的罪行被犹太人惩罚}，但我们不再进行辩护，因为我们对自己已经说过的话感到满意。但接下来，由于他笔下的这位犹太人贬低关于死人复活、天主审判、义人应得的赏报以及吞噬恶人的火的道理，认为这些都是老生常谈，并以为断言在教导这些事上并无新意就能推翻基督信仰，我们必须对他说：我们的主看到犹太人的行为完全不符合先知的教导，便用比喻训诲说，天主的国将从他们那里夺去，赐给外邦的皈依者。因此，现在我们也能真实地看到，现今犹太人的一切教义都只是琐事和寓言，因为他们没有从圣经知识而来的光明；而基督徒的教义却是真理，有能力提升和升华人的灵魂与理智，并说服人寻求一个天上而非地上属世犹太人的公民身份。这一结果在那些能够看清法律和先知所含概念之伟大，并能将其推荐给他人的人身上显明出来。\par

% structure: p0012 source=h2 target=subsection reason=relative-heading-rank; base=h2

\subsection{第六章}

但就算耶稣遵守了犹太人的一切习俗，包括他们的祭祀礼仪，这又怎能妨碍我们承认祂是天主子呢？耶稣，就是那位赐下法律和先知的天主子；而我们这些属于教会的人，并未违反法律，而是摆脱了犹太人的神话化，并通过法律和先知的奥秘默观而使我们的心灵受到陶冶和教化。因为先知们本身并非将这些话的意义局限于它们所叙述的普通历史，也不是按照字面和条文来理解法律的规定，他们在即将叙述历史时，曾这样表达：\BibleQuote{我要开口说比喻，要说古时的谜语。}在另一处，当祈祷祈求理解法律（这法律是隐晦的，需要天主的帮助才能明白）时，他们献上这样的祈祷：\BibleQuote{求你开启我的眼睛，使我看见你法律中的奇妙。}\par

% structure: p0014 source=h2 target=subsection reason=relative-heading-rank; base=h2

\subsection{第七章}

此外，让他们指出，何处能找到耶稣所发出的、出自傲慢的言语，哪怕只是表面现象？因为一个傲慢的人怎能这样表达自己：\Quote{你们跟我学习吧，因为我是良善心谦的，你们必会找到灵魂的安息？}又怎能称祂为傲慢者，祂吃完晚餐后，在门徒面前脱去外衣，拿一条毛巾束腰，把水倒在盆里，就开始洗每个门徒的脚，并责备那个不愿让祂洗脚的人说：\Quote{我若不洗你，你就与我无分？}又怎能称祂为傲慢者，祂曾说：\Quote{我在你们中间，不是像坐席的，而是像服事人的？}请任何人指出祂所说的谎言是什么，并指出哪些是大的谎言，哪些是小的谎言，以便证明耶稣犯了大的谎言。还有另一种驳斥他的方式：因为一个谎言并不比另一个更虚假或更不虚假，同样，一个真理也不比另一个更真实或更不真实。至于他对耶稣提出了哪些不敬天主的指控，让采尔苏斯的犹太人尤其提出来吧。难道不敬天主是：放弃肉身的割礼、字面的安息日、字面的庆节、字面的月朔、以及洁净与不洁净的食物，而将心思转向良善、真实和属天主的法律，同时，那位为基督作使者的人知道如何向犹太人成为犹太人，为要赢得犹太人，向在法律下的人成为在法律下的，为要赢得在法律下的人吗？\par

% structure: p0016 source=h2 target=subsection reason=relative-heading-rank; base=h2

\subsection{第八章}

他说，此外，\Quote{许多其他人物，对于甘愿受骗的人来说，会显得像耶稣一样。}那么让这位采尔苏斯的犹太人给我们指出，不是许多人，也不是少数，而是一个像耶稣那样的个体，他以显现在祂身上的能力，在人类中引入了一个有益于人生、使人从邪恶行为中改邪归正的教义和思想体系。他还说，基督徒皈依者向犹太人提出这项指责，说他们没有相信耶稣是天主。在这一点上，我们已经在前面几页中做了初步辩护，同时指出了我们在哪些方面认为祂是天主，在哪些方面认为祂是人。\Quote{我们这些，}他接着说，\Quote{向所有人宣告有一位从天主而来、将要惩罚恶人的人，怎能在他来时却轻视他呢？}对于这个极为愚蠢的论点，我认为不值得做任何回答。这就像有人说：\Quote{我们这些教导节制的人，怎能做出放荡的事？或者我们这些正义的使者，怎能犯罪呢？}因为这些矛盾在人中间存在，所以，说他们相信先知们关于基督未来降临的预言，而当基督按照先知的话来临时，他们却拒绝相信祂，这完全符合人性。既然我们还要加上另一个理由，我们就说，这个结果本身就由先知预言过。依撒意亚明确宣告：\BibleQuote{你们听了听，却不明白；看了看，却不领悟，因为这人民的心变迟钝了，}等等。让他们解释，为什么向犹太人预言了，他们尽管听了看了，却不明白所说的话，也不领悟所当看见的。因为显然，当他们看到耶稣时，他们并没有看见祂是谁；当他们听到祂时，他们从祂的话中没有理解祂里面的天主性，而这天主性将天主迄今为止对犹太人的眷顾转移到了祂从外邦人中所招聚的人。因此我们可以看到，在耶稣来临之后，犹太人被完全抛弃了，现在不再拥有任何他们古代的光荣，没有任何天主居留于他们中间的迹象。因为他们不再有先知或奇迹，而这些痕迹在很大程度上仍然可以在基督徒中找到，其中一些比犹太人中所存在的更显著；如果我们的见证可以被接受的话，这些我们自己也亲眼见过。但采尔苏斯的犹太人喊道：\Quote{为什么我们羞辱了事先宣告的那一位？是为了使我们比其他人受更多惩罚吗？}对此我们必须回答，由于他们的不信和对耶稣的其他侮辱，犹太人不仅在那被认为将要临到世界的审判中比其他人受更多的苦，而且已经遭受了这样的苦难。因为除了犹太人之外，还有哪个民族是从自己的都城，从祖先敬拜的圣地流亡的呢？他们遭受这些灾难，是因为他们是一个极其邪恶的民族，尽管犯了许多其他罪，却没有任何一项罪像他们犯在我们耶稣身上的那些罪那样受到如此严厉的惩罚。\par

% structure: p0018 source=h2 target=subsection reason=relative-heading-rank; base=h2

\subsection{第九章}

犹太人继续他的论辩说：\Quote{我们怎能认为他是天主呢？他不仅在其它方面——正如传闻所说的——没有实现他的任何承诺，而且在我们定他的罪、判他应受惩罚之后，他竟被发现企图隐藏自己，并以最可耻的方式试图逃跑，还被那些他称为门徒的人出卖了？然而，}他继续说，\Quote{一位天主既不能逃跑，也不能被带走成为囚徒；尤其是，他不能被那些曾是同伴、共享一切、以他为老师的人所抛弃和交出，而他被视为救主、至大天主的儿子、以及一位天使。}对此我们回答说：我们并不认为\Person{耶稣}的身体（那是当时可看见、可感知的）是天主。我为何说他的身体呢？不，甚至不认为他的灵魂——有记载说：\BibleQuote{我的心灵忧闷得要死}——是天主。但按照犹太人的说法：\BibleQuote{我是主，是一切有血肉者的天主}，以及\BibleQuote{在我以前没有天主被形成，在我以后也不会有}，天主被认为是使用先知之灵魂和身体作为工具的那一位；并且，按照希腊人的说法，那位说：\par

\Quote{我知道沙子的数目和大海的尺度，\newline{}我也理解哑巴，并听见不说话的人，}\par

他说话时，并借\Place{皮提亚}女祭司使自己被听见时，被视为一位天主；同样，按照我们的看法，是\OriginalTerm{圣言}{\GreekText{Λόγος}}天主，万有之天主圣子，在\Person{耶稣}内说了这些话：\BibleQuote{我是道路、真理、生命}，以及这些：\BibleQuote{我是门}，以及这些：\BibleQuote{我是从天上降下来的生活的食粮}，以及其他类似的言词。因此我们指责犹太人，不承认他为天主，而先知们在许多经文里曾为他作证，说他是一种大能，是仅次于万有之父天主的另一位天主。因为我们断言，在\Person{梅瑟}的创世记述中，圣父发出命令时，就是向祂说的：\BibleQuote{有光}，和\BibleQuote{有穹苍}，并对其他受造物发出命令；并且也对祂说了：\BibleQuote{让我们照我们的肖像，按我们的模样造人}；而圣言奉命令时，服从了圣父的一切旨意。我们这样说，不是出于自己的臆测，而是因为我们相信在犹太人流传的预言中，关于天主和创造工程，有明确的话说：\BibleQuote{祂说了，它们就造成；祂命令了，它们就创造出来。}如果天主下了命令，受造物就形成了，那么按照先知之神的看法，除了那可以说是活生生的圣言和真理的那一位，谁还能完成圣父的这些命令呢？福音书并不认为在\Person{耶稣}内说这些话的那一位\BibleQuote{我就是道路、真理、生命}，其本性如此受限，以至于除了\Person{耶稣}的灵魂和身体之外便无存在；这一点从许多考虑以及我们将要引用的几个例子中可以明显看出。\Person{若翰洗者}在预言天主子即将显现时，说祂不是在那身体和灵魂中，而是随处显现自己，关于祂说：\BibleQuote{在你们中间站着一位你们不认识的人，祂在我以后来。}因为如果他以为天主子只在那里，即\Person{耶稣}可见的身体那里，他怎能说\BibleQuote{在你们中间站着一位你们不认识的人}呢？而\Person{耶稣}自己，为提升门徒对天主子更崇高的思想，说：\BibleQuote{哪里有两个或三个人因我的名字聚在一起，我就在他们中间。}同样性质的是祂对门徒的许诺：\BibleQuote{看，我天天与你们在一起，直到今世的终结。}我们引用这些经文，并不在天主子和\Person{耶稣}之间作区分。因为\Person{耶稣}的灵魂和身体，在\OriginalTerm{救恩计划}{\GreekText{οἰκονομία}}之后，与天主的圣言成为一体。如果按照\Person{保禄}的教导，\BibleQuote{那与主结合的，便是一个心神}，那么凡理解与主结合为何物并实际与祂结合的人，就与主是一个心神；那么那个曾与天主的圣言结合的存在，难道不应在更大更神圣的程度上是一体吗？祂确实藉着所行的奇迹，在犹太人中间显明自己是天主的能力；\Person{凯尔苏斯}怀疑这些奇迹是藉巫术完成的，但当时的犹太人却不知为何将它们归于\Person{贝耳则步}，说：\BibleQuote{祂是靠魔王贝耳则步驱魔。}但我们的救主斥责这些人说了最大的谬论，因为邪恶的王国尚未终结。这对所有有智慧的福音叙述读者将是显而易见的，现在不是解释的时候。\par

% structure: p0022 source=h2 target=subsection reason=relative-heading-rank; base=h2

\subsection{第十章}

但是耶稣曾作过什么应许而没有实践呢？让\Person{策尔苏斯}举出任何这样的实例，并证实他的指控。但他将无法做到，尤其是因为他是从误解福音叙述或犹太人的故事中，认为可以找出他对耶稣或对我们提出的指控根据。此外，这位犹太人又说：\Quote{我们判定他有罪，并宣判他应受死刑。}让他们指出，那些设法捏造伪证反对他的人，是如何证明他有罪的。控告耶稣的人所提出的主要指控，不就是祂所说的\BibleQuote{我能拆毁天主的圣殿，并在三天内把它重建起来}吗？但祂这样说，是指祂身体的圣殿；而他们因为不能理解说话者的意思，以为祂指的是石造的圣殿，犹太人对其尊敬程度超过那本当受尊敬者——那真正的天主圣殿，就是\OriginalTerm{圣言}{Word}、\OriginalTerm{智慧}{Wisdom}和\OriginalTerm{真理}{Truth}。谁能说\Quote{耶稣曾试图可耻地躲藏起来以逃脱}呢？让任何人指出一件可称为可耻的行为。当他接着说\Quote{祂被俘了}时，我会说，如果被俘意味着违背自己意愿的行为，那么耶稣并没有被俘；因为在合适的时候，祂并没有阻止自己落入人手中，如同天主的羔羊，为除去世人的罪。因为祂知道将要临到自己的一切事，就上前去，对他们说：\BibleQuote{你们找谁？}他们回答：\BibleQuote{纳匝肋人耶稣。}祂对他们说：\BibleQuote{我就是。}出卖祂的犹大也同他们站在一起。所以，当祂对他们说：\BibleQuote{我就是}时，他们就退后，倒在地上。祂又问他们：\BibleQuote{你们找谁？}他们又说：\BibleQuote{纳匝肋人耶稣。}耶稣对他们说：\BibleQuote{我告诉你们，我就是；你们既然找我，就让这些人走吧。}不但如此，祂甚至对那想要帮助祂、击伤大司祭仆人并削掉他耳朵的人说：\BibleQuote{把你的剑放回原处；因为凡动剑的，必死于剑下。你想我不能求我父，祂就立即给我超过十二营的天使吗？但经上记载必须如此，又怎能应验呢？}若有人认为这些叙述是福音作者的编造，那么为什么不更应将那些出于仇恨和敌意反对耶稣和基督徒的叙述视为编造呢？而这些真理出自那些通过对耶稣的感情表现真诚的人，他们为祂的话忍受一切、无论什么考验。因为门徒们领受了这种毅力和决心，甚至至死，并有一种心态，不会为他们的导师编造不实之事，这对于所有公正的判断者来说都是一个非常明显的证据，证明他们完全确信自己所写的是真理，因为他们为所相信为天主子的那一位，忍受了如此众多而严峻的考验。\par

% structure: p0024 source=h2 target=subsection reason=relative-heading-rank; base=h2

\subsection{第十一章}

其次，祂被那些祂称为门徒的人出卖，这是\Person{责尔苏}笔下的犹太人从\WorkTitle{福音书}中得知的事；他称其中一个为犹达斯，然而又说\Quote{许多门徒}，以此似乎加重指控。他并不费心注意所有关于犹达斯的记载：这个犹达斯对于祂的师傅持有相反且矛盾的意见，既没有全心全意地反对祂，也没有全心全意地保持学生对教师应有的尊敬。因为出卖祂的人给来捉拿\Person{耶稣}的群众一个暗号，说：\BibleQuote{我亲吻谁，谁就是；抓住他}——仍然保留着对祂的师傅的某种尊敬；因为若不如此，他就会公开地出卖祂，没有任何感情上的伪装。因此，关于犹达斯的意图，这一情形将使所有人满意：即在他贪婪的本性和出卖祂的师傅的邪恶计划中，他心里仍有一种混合的情感，是由\Person{耶稣}的话在他心中产生的，可以说是某种美善的残余。因为记载说：\BibleQuote{出卖祂的犹达斯见祂已被定罪，就后悔了，把那三十块银钱还给大司祭和长老，说：我犯了罪，因为我出卖了无辜的血。但他们说：这与我们何干？你自己看着办吧。}——并且他把钱扔在圣殿里，就出去上吊自尽了。但如果这个贪婪的犹达斯，也曾偷窃放在钱袋中用于救济穷人的钱，悔改了，并把那三十块银钱还给大司祭和长老，那么显然\Person{耶稣}的教诲能够在他心中产生某种悔改的感觉，并没有被这个叛徒完全轻视和憎恶。不仅如此，那句声明：\BibleQuote{我犯了罪，因为我出卖了无辜的血。}是他对自己罪行的公开承认。也要注意，从那悔改中产生的对罪的悲伤是何等强烈，以至于不容他再活下去；他怎样把钱扔在圣殿里，就离开走掉，上吊自尽了：因为他给自己定了罪，显示出\Person{耶稣}的教导对这个罪人犹达斯——这个贼和叛徒——具有何等大的力量，使他不能始终轻蔑地对待他从\Person{耶稣}所学的一切。难道\Person{责尔苏}和他的朋友们现在要说，那些证明犹达斯的背弃并非完全背弃（甚至在他图谋反抗祂的师傅之后）的证据都是捏造的，只有这一点是真的，即祂的一个门徒出卖了祂；并且他们要在圣经记载上再加上：他也是全心全意地出卖了祂吗？在同一部著作中，以这种敌意来决定哪些该信、哪些不信，是荒谬的。如果我们必须就此发表一番使对手蒙羞的陈述，我们要说，在\WorkTitle{圣咏集}中，第108篇全篇是关于犹达斯的预言，开头是这样：\BibleQuote{天主，在我的赞美前，求你不要缄默；因为罪人的口和狡诈之人的口，已经张开攻击我。}在这篇圣咏中，也预言了犹达斯因自己的罪而离开宗徒的行列，另有一位被选来代替他的职位；这一点由以下话语表明：\BibleQuote{愿他的主教职被别人取代。}但现在假设祂被祂的一个门徒出卖了，这个门徒被比犹达斯更恶劣的邪灵所附，并完全倾泻出（可以这样说）他从\Person{耶稣}那里听到的一切话语，这对控告\Person{耶稣}或基督宗教有什么帮助呢？这又如何证明其教义是虚假的呢？我们在前一章中已经回复了紧随其后的那些陈述，表明\Person{耶稣}并不是在试图逃跑时被逮捕，而是为我们众人自愿交付了自己。由此可知，即使祂被捆绑，也是按照祂自己的意愿被捆绑；从而教导我们一课：为了宗教的缘故，我们应在毫不勉强的心情下承担类似的事。\par

% structure: p0026 source=h2 target=subsection reason=relative-heading-rank; base=h2

\subsection{第十二章}

以下在我看来是幼稚的断言，即：\Quote{从来没有善良的将军和伟大群众的领袖被出卖过；甚至没有一个邪恶的强盗头目和极恶之人的指挥官，他似乎对其同伙有些用处；但耶稣，被其下属出卖后，既没有像一位好将军那样统率，也没有在欺骗他的门徒之后，在他欺骗的受害者心中产生那种，可以这么说，会向一个强盗首领表现出来的善意。}现在人们可以找到许多关于将军被自己士兵出卖的记载，以及强盗头目因那些不守约定的人而被捕的记载。但即使没有将军或强盗头目曾被出卖，这与事实——即耶稣的一位门徒成了他的背叛者——作为对耶稣的指控有什么关联呢？既然塞耳苏斯炫耀哲学，我要问他：那么，如果亚里士多德在做了柏拉图二十年的学生后离去，攻击他关于灵魂不朽的学说，并把柏拉图的理念称为纯粹的琐碎，这是对柏拉图的指控吗？如果我仍有疑问，我会继续这样问：在亚里士多德离开后，柏拉图在辩证法上不再强大，不能再捍卫他的观点了吗？因此，柏拉图的意见就是假的吗？或者，尽管柏拉图是正确的——正如他的哲学的学生会主张的那样——亚里士多德是否对他的老师犯了邪恶和忘恩负义之罪？不仅如此，克里西波也在他的著作的许多地方似乎攻击克莱安特，引入与他观点对立的新奇意见，尽管后者在他年轻时是他的老师，并且他开始学习哲学。据说亚里士多德确实是柏拉图的学生二十年，而克里西波在克莱安特的学园中度过了不短的时间；而犹达斯与耶稣在一起的时间还不到三年。但从哲学家生平的故事中，我们可以找到许多与塞耳苏斯因犹达斯而对耶稣提出的指控相似的例子。甚至毕达哥拉斯学派也为那些在投身哲学后又堕回无知生活方式的人建造衣冠冢；然而，毕达哥拉斯及其追随者并不因此而在论证和证明上软弱。\par

% structure: p0028 source=h2 target=subsection reason=relative-heading-rank; base=h2

\subsection{第十三章}

这个刻尔苏斯笔下的犹太人继续以上内容，如此说道：\Quote{虽然他能够就耶稣生平的事迹陈述许多真实的情况，与门徒所记载的不同，却甘愿略过不提。}那么，那些与福音记载不同的真实陈述，刻尔苏斯笔下的犹太人略过不提的，到底是什么呢？还是说他只是在运用一种貌似修辞的手法，假装有话要说，而实际上他除了福音叙述之外，并没有任何能够使听者感受到其真实性，并为控告耶稣及其教义提供明确依据的内容？他还指控门徒编造了耶稣预知并预言了发生在他身上的一切这一说法；但我们将透过救主所说的许多其他预言——他在其中预告了后世基督徒将要遭遇的事——来确证这一说法的真实性，尽管刻尔苏斯或许不喜欢。有谁不为这一预言而感到惊奇呢：\BibleQuote{你们要因我的缘故，被带到总督和君王面前，对他们和外邦人作证}；以及他可能说过的其他关于门徒将来受迫害的预言？因为世人中曾有过什么思想体系，因其缘故而使他受到惩罚，以至于耶稣的控告者中有人可以说，他预见到自己思想的不敬神或虚假将成为控告他们的理由，他认为这会有助于他的声誉，说他早已事先预言了这一点？如今，若有人因自己的见解而该被带到总督和君王面前，除了\OriginalTerm{伊壁鸠鲁学派}{Epicureans}——他们完全否认\OriginalTerm{天主眷顾}{providence}的存在——之外，还有谁呢？还有\OriginalTerm{逍遥学派}{Peripatetics}，他们说祈祷毫无用处，向神明献祭也无用。但有人会说，撒玛黎雅人因他们的宗教而遭受迫害。对此我们将指出，\OriginalTerm{西卡里人}{Sicarians}因施行割损——即自残身体，违反既定法律以及仅允许犹太人的习俗——而被处死。你从未听说过法官会问，一个努力按此既定宗教生活的西卡里人，若背弃信仰是否会获得释放，若坚持不变则被处死；因为割损的证据足以确保受割损者必死无疑。唯独基督徒，按照他们救主的预言\BibleQuote{你们要因我的缘故被带到总督和君王面前}，被法官们催逼至最后一口气，要他们否认基督信仰，并按照公家习俗献祭；然后发誓弃绝，才可回家安全度日。且看这句话是否不是以极大的权柄说出的：\BibleQuote{所以，凡在人前承认我的，我在我天上的父前也必承认他；凡在人前否认我的}，等等。请与我一同回想耶稣说这些话的时候，看他那时的预言尚未应验。或许你会以一种不信的口气说，他在说胡话，白费唇舌，因为他的话不会应验；或者，在犹豫是否相信他的话时，你会说，如果这些预言应验了，耶稣的教义得以确立，以至于总督和君王都想着要消灭那些承认耶稣的人，那么我们就相信他说这些预言，是作为一个从\OriginalTerm{天主}{God}那里领受了极大能力，将这一教义植入人类种族之中的人，并且相信它必将得胜。有谁不会充满惊奇呢？当他回想当时教导并说\BibleQuote{这福音必传遍普天下，对他们和外邦人作证}的那一位，并看到按照他的话，耶稣基督的福音在天穹之下传遍了普天下，传给希腊人和野蛮人、聪明人和愚拙人？因为那用权能说出的话，已经征服了各种本性的人，并且不可能看到有任何种族的人未曾接受耶稣的教导。但让这个刻尔苏斯笔下的犹太人不相信他预知了一切发生在他身上的事吧，他当思考：当耶路撒冷尚存，全犹太人的敬拜在其中举行时，耶稣是如何预言它将遭罗马人的手毁灭的。因为他们不会主张，耶稣的熟人和门徒本人没有将记在福音中的教导写下来流传后世，从而使他的门徒没有得到他们作品中所包含的耶稣回忆录。在这些作品中记载着：\BibleQuote{当你们看见耶路撒冷被军队围困时，那时就知道她的荒凉近了。}但在那时，并没有军队围绕耶路撒冷，包围、围困它；因为围攻开始于尼禄在位时期，持续到维斯帕先执政，他的儿子提图斯毁灭了耶路撒冷，据\Person{若瑟夫}所说，是因了称为基督的耶稣的兄弟义者雅各的缘故，但实际上，正如真理所显明的，是因了天主之子耶稣基督的缘故。\par

% structure: p0030 source=h2 target=subsection reason=relative-heading-rank; base=h2

\subsection{第十四章}

然而，\Person{策尔苏斯}或可接受或承认耶稣预知了将临于祂的一切，但或会想轻描淡写地接受这一事实，正如他在论及奇迹时所为——彼时他声称奇迹是藉巫术而行；因为他或会说，许多人藉占卜，或鸟占、或卜兆、或牺牲、或星命术，得以预知将来之事。但他不会作出此等让步，因这让步过大；尽管他隐约承认耶稣行了奇迹，却想以巫术之指控削弱其力量。如今，\Person{弗勒贡}在其\WorkTitle{编年史}（我认为是第十三或第十四卷）中，不仅将预知未来事件归因于耶稣（尽管在涉及伯多禄的某些事上混淆不清，仿佛这些事关乎耶稣），而且证实了结果与祂的预言相符。因此，他本人也藉这些关于预知的承认，仿佛违背己意，表达了其见解：我们体系的父辈所传的教义并非缺乏神性之能。\par

% structure: p0032 source=h2 target=subsection reason=relative-heading-rank; base=h2

\subsection{第十五章}

\Person{策尔苏斯}继续说道：\Quote{耶稣的门徒，因无确凿事实可依，便虚构了祂预知一切将发生之事；}——他们未观察，或不愿观察，作者们所怀的真理之爱：他们承认耶稣曾预先告诉门徒：\BibleQuote{今夜，你们为我的缘故都要跌倒}——此语应验，因他们全都跌倒了；祂又向伯多禄预言：\BibleQuote{鸡叫以前，你要三次不认我}——随后伯多禄三次否认。若说他们并非热爱真理，而如策尔苏斯所揣度，是虚构捏造者，他们便不会记载伯多禄的否认，也不会记载门徒的跌倒。尽管这些事确实发生，然而谁能证明它们果然如此？然而，按一切可能性，这些事本应被那些欲教导福音读者为宣认基督信仰而轻看死亡的人所略过。但如今，既然圣言藉其大能将得胜于人，他们便如实记述了这些事；而这些事，不知何故，既不会伤害读者，也不为否认提供任何借口。\par

% structure: p0034 source=h2 target=subsection reason=relative-heading-rank; base=h2

\subsection{第十六章}

他的论断极其薄弱，他说：\Quote{耶稣的门徒写下这些关于他的记述，是为了减轻对他的指控：就好像}，他说道，\Quote{有人说某人是义人，却又显示他行了不义；或者说他是虔诚的，却又犯了谋杀；或者说他是不死的，却又死了；并在所有这些陈述之后加上一句话，说他预言了这一切。}现在，他的比喻立即显得不合适；因为那位决定成为世人在生活规范上的活榜样的人，同时也表明他们应当如何为宗教而死，这并无荒谬之处，更不必说他为人类而死是对整个世界的恩惠，正如我们在前书中所证明的。此外，他想像救主所受的一切苦难的宣认，不仅没有削弱，反而加强了他的反对意见。因为他既不了解\Person{保罗}的哲学论述，也不了解先知们关于此事的陈述。他也忽略了某些异端者宣称\Person{耶稣}所受的苦难只是外表，而非真实。因为他若知道，就不会说：\Quote{你们甚至没有声称，他似乎在恶人眼中受了这刑罚，而实际上并未遭受；相反，你们承认他公开受了苦。}但我们并不认为他的苦难只是外表，这样他的复活也不会是虚假的，而是真实的事件。因为真正死了的人，如果他确实复活了，就是真正复活了；而只是表面死了的人，并没有真正复活。但由于\Person{耶稣基督}的复活成为不信者的笑柄，我们将引用\Person{柏拉图}的话，说\Person{阿尔美尼乌斯之子厄鲁斯}在火葬堆上十二天后复活，并叙述了他在冥府所见；既然我们是在回答不信者，在此引用\Person{赫拉克利德斯}关于那位失去生命的妇人的记载，也并非全然无用。许多人被记载从坟墓中复活，不仅是在下葬之日，还有次日。那么，如果那位行了许多奇妙之事的人，既超越人的能力又有如此充分的证据，以至于无法否认其作为的人试图将它们比作巫术来诽谤，他在死亡中也显示出更大的天主性能力，使他的灵魂（如果愿意）可以离开身体，在身外执行某些职分，然后随意返回，这又有什么奇怪的呢？\Person{耶稣}在\BibleBook{若望}{福音}中这样说过：\BibleQuote{谁也不能夺去我的性命，而是我甘心情愿舍掉它；我有权舍掉它，我也有权再取回它来。}\BibleRef{若 10:18}也许正是为此，他加速了离开身体，以保全身体，使他的腿不被折断，如同与他同钉的强盗一样。\BibleQuote{兵士便来了，把第一个人的腿打断，又把另一个与他同钉的人的腿打断；可是，当他们来到耶稣跟前，看见他已经死了，就没有打断他的腿。}\BibleRef{若 19:32-33}因此，我们已经回答了\Quote{耶稣如何能够预言这些事？}这个问题。至于\Quote{死人如何能不灭？}，愿理解的人知道，不灭的不是死人，而是从死者中复活的那一位。事实上，死人远非不灭，甚至\Person{耶稣}在去世之前——那个复合的、将要受死的人——也不是不灭的。因为注定要死的人没有不灭的；只有当他不复受死亡管辖时，他才是不灭的。但\BibleQuote{基督既从死者中复活，就不再死；死亡不再统治他了。}\BibleRef{罗 6:9}虽然那些不能理解这样的事如何可能的人可能不愿承认这一点。\par

% structure: p0036 source=h2 target=subsection reason=relative-heading-rank; base=h2

\subsection{第十七章}

他的言论也极其愚蠢，他说：\Quote{哪一个神明、神灵或有智慧的人，在预知这些事将要临到他时，不会尽可能避免它们？但他却一头扎进那些他预先知道将要发生的事中。}然而，\Person{苏格拉底}知道喝了毒药后会死，如果他听从\Person{克里托}的劝说，从监狱逃走，他本可以避免这些灾难；但他仍然决定，在他看来符合正确理性，宁可以符合哲学家身份的方式死去，也不要以不体面的方式保住生命。\Person{列奥尼达}，那位拉刻代蒙将军，知道他将与部下在\Place{温泉关}战死，却没有试图以可耻的方式保全生命，而是对同伴们说：\Quote{我们去吃早餐吧，因为我们将在冥府用晚餐。}那些有兴趣收集这类故事的人，会发现很多这样的例子。那么，如果\Person{耶稣}知道一切将要发生的事，却没有避免，反而迎上他所预知的事，这有什么奇怪呢？当\Person{保罗}，他的门徒，听到他上\Place{耶路撒冷}将遭遇什么时，他仍前去面对危险，斥责那些在他周围哭泣、想要阻止他上\Place{耶路撒冷}的人？我们同时代的许多人，也深知如果宣认基督信仰就会被处死，如果否认就能获释、财产也被归还，但他们却蔑视生命，自愿为宗教选择死亡。\par

% structure: p0038 source=h2 target=subsection reason=relative-heading-rank; base=h2

\subsection{第十八章}

此后，那犹太人又提出一个愚蠢的说法，说：\Quote{如果耶稣预先指出叛徒和发假誓者，他们怎么会不敬畏祂如同天主，而停止，一个放弃意图的背叛，另一个放弃他的背誓呢？}在这里，博学的策尔苏没有看到他陈述中的矛盾：因为如果耶稣作为天主预知事件，那么祂的预知不可能不真实；因此，祂所知道将要背叛祂的人不可能不执行其意图，祂所斥责将要否认祂的人也不可能不犯那罪行。因为如果一个人能够因事先被警告这些行为的后果而避免背叛行为，另一个人也能避免否认行为，那么祂的话就不再真实，祂预言一个人将背叛祂，另一个人将否认祂。因为如果祂预知叛徒，祂就知道背叛所源于的那邪恶，而这邪恶绝不会因预知而被除去。再者，如果祂确知一个人将否认祂，祂是从看出那否认行为将要产生的软弱而做出那预言的，然而这软弱也不会因预知而立即被除去。但我不知道他从哪里得出这样的说法：\Quote{这些人背叛祂和否认祂，却没有表现出对祂的丝毫担忧。}因为关于叛徒，已经证明说他没有表现出对主的担忧而背叛主是假的。同样，关于否认祂的人也证明了这一点：因为在否认之后，他出去痛哭。\par

% structure: p0040 source=h2 target=subsection reason=relative-heading-rank; base=h2

\subsection{第十九章}

他提出的异议也很肤浅：\Quote{通常的情况是，当一个人被密谋对付，而他知道了这阴谋，并向共谋者表明他知道他们的计划时，后者就会改变主意，并保持警惕。}因为许多人即使在知道他们计划的人面前也继续密谋。然后，他好像要总结自己的论点，说：\Quote{不是因为这些事情被预言了它们才发生，那是不可能的；但因为它们已经发生了，它们被预言就表明是假的：因为那些事先听说他们的阴谋被发现的人，要执行他们背叛和否认的计划，是完全不可能的！}但如果他的前提被推翻，那么他的结论也站不住脚，即\Quote{我们不应该因为这些事情被预言了就相信它们已经发生。}现在我们主张，它们不仅作为可能的事情发生了，而且因为它们发生了，它们被预言的事实被表明是真实的；因为关于未来事件的真理是通过结果来判断的。因此，正如他所说，这些事件的预言被证明是不真实的，这是假的；他说\Quote{那些事先听说他们的阴谋被发现的人，要执行他们背叛和否认的计划，是完全不可能的}，这是徒劳的。\par

% structure: p0042 source=h2 target=subsection reason=relative-heading-rank; base=h2

\subsection{第二十章}

让我们看看他在这之后如何继续。他说：\Quote{这些事件，他作为神预言了，预言必定要实现。因此，那位比所有人都更应该施恩于人类、尤其是自己家里的人的天主，竟带领自己的门徒和先知——他惯常与他们一同吃喝——陷入如此严重的邪恶，使他们成为不虔敬和不圣洁的人。事实上，与人同桌共食的人不会犯谋害之罪；但与天主共宴之后，他却成为了同谋。更荒谬的是，天主自己竟图谋陷害与他同席的人，将他们变成叛徒和恶徒！}现在，既然你希望我回答那些甚至在我看来无稽的指控，以下就是我们对这些说法的回复。克耳苏斯以为，一件因预知而被预言的事，是因为被预言才发生的；但我们不承认这一点，主张预言者并非事件发生的原因，因为他预言了它会发生；而未来的事件本身，即使没有预言也会发生，却为具有预知能力的人提供了预言其发生的时机。当然，当事件可能发生也可能不发生，即二者之一会发生时，这结果就呈现在预言者面前。因为我们并不主张，预知事件的人通过暗中消除其可能发生或不发生的可能性，就作出这样的宣告：\Quote{这事必定发生，绝无其他可能。}这一论述适用于一切取决于我们自身的事件的预知，无论是在圣经中，还是在希腊人的历史中。逻辑学家所谓的\OriginalTerm{懒惰论证}{idle argument}，是一种\OriginalTerm{诡辩}{sophism}，在克耳苏斯看来它可能不是诡辩，但根据正确的推理，它确实是诡辩。为了说明这一点，我将从圣经中取出关于犹达斯的预言，或我们的救主预知他为叛徒的事；并从希腊历史中取出给予拉伊俄斯的神谕，暂且承认其真实性，因为它不影响论证。在圣咏一〇八篇中，救主藉着圣咏的话论及犹达斯，开头说：\BibleQuote{我赞美之天主，求莫缄默无言；因为邪恶人的口和欺诈人的口，已张开攻击我。}如果你仔细查看这圣咏的内容，就会发现：正如预知他将出卖救主，他也被视为出卖的原因，并因他的邪恶而应受预言中的咒诅。因为，圣咏作者说，让他受这些事，\BibleQuote{因为他不纪念施恩，却迫害穷困的人。}因此，他本可以施恩，不迫害他所迫害的人。但他虽可能这样做，却没有做，而是执行了叛卖的行为，以致应得预言中对他宣布的咒诅。\par

我们回答希腊人时，将引用悲剧诗人所记载的、给予 \Person{拉伊奥斯} 的神谕，或用神谕的原话，或用同义的话。神谕向他预告未来之事说：\Quote{不要试图违背众神之意生子。因为如果你生了一个儿子，你的儿子必将杀死你；你全家将血流成河。}由此可知，\Person{拉伊奥斯} 有能力不去试图生子，因为神谕不会命令不可能之事；他也有能力做相反的事，所以两者都不是必然的。他没有防止生子的结果，便是因此遭受了悲剧中关于 \Person{俄狄浦斯}、\Person{约卡斯塔} 及其儿子们所描述的灾祸。所谓的\Quote{\OriginalTerm{懒惰论证}{idle argument}}是一种诡辩，例如可用于病人的情况，以诡辩的方式阻止他请医生来促进康复；其大意如下：\Quote{如果命定你会痊愈，那么无论你请不请医生，你都会痊愈；但如果命定你不会痊愈，那么无论你请不请医生，你都不会痊愈。但命定要么痊愈要么不痊愈，所以你请医生是徒劳的。}与此巧妙类比的是：\Quote{如果命定你生子，那么无论你与女人交合与否，你都会生子；但如果命定你不生子，那么无论你与女人交合与否，你都不会生子。现在，命定要么生子要么不生子，所以与女人交合是徒劳的。}因为，在后一例中，与女人交合并非徒劳，因为不交合则完全不可能生子；同样，在前一例中，如果疾病康复需通过医术实现，那么必须请医生，因此说\Quote{你请医生是徒劳的}就是错误的。我们举这些例子，都是因为这位博学的 \Person{塞尔苏斯} 断言：\Quote{既然祂是天主，祂预言了这些事，那么预言\Emph{必然}应验。}如果\Quote{\Emph{必然}}是指\Quote{\Emph{必定}}，我们不能接受。因为完全可能不应验。但如果他用\Quote{\Emph{必然}}是指\Quote{\Emph{\OriginalTerm{单纯未来性}{simple futurity}}}，即没有什么妨碍它为真（尽管可能不发生），那么他根本不影响我的论点；耶稣预言了叛徒或背誓者的行为，并不等于祂是这种不敬不圣之事的起因。因为祂在我们中间，知道人心里有什么，看到那人的邪恶性情，因他的贪财之心和他对师尊缺乏坚定的职责观念，预见他将会做什么，于是与许多其他宣告一起，也说了这话：\BibleQuote{同我蘸手在盘子里的，就是他要出卖我。}\par

% structure: p0045 source=h2 target=subsection reason=relative-heading-rank; base=h2

\subsection{第二十一章}

也要看到 \Person{塞尔苏斯} 论断的肤浅和明显的虚假，他断言\Quote{同席之人不会谋害主人；如果他不谋害人，更不会在与人共餐后谋害天主。}谁不知道许多人在分享餐桌上的盐后，竟谋害了款待他们的人呢？全部希腊和蛮族历史中充满了这样的实例。帕罗斯岛的抑扬格诗人斥责 \Person{吕坎贝斯} 违背了以餐桌之盐确认的盟约，对他说：\par

\Quote{但你违背了一个重大的誓言——即以餐桌之盐所立的誓。}\par

对那些爱好历史学问、全身投入其中而忽略生活中更必要的其他知识的人，可以举出无数例子，证明那些受到他人款待的人却谋害了他们。\par

% structure: p0049 source=h2 target=subsection reason=relative-heading-rank; base=h2

\subsection{第二十二章}

他还加上，仿佛已经用无可辩驳的论证和结论集结了一个论点：\Quote{而且，更荒谬的是，天主自己谋害了坐席的人，把他们变成叛徒和不敬的人。}然而，要证明耶稣如何谋害或将门徒变成叛徒或不敬的人，他是不可能的，除非用那种任何人都极易驳斥的推论。\par

% structure: p0051 source=h2 target=subsection reason=relative-heading-rank; base=h2

\subsection{第二十三章}

\Quote{如果他决定了这些事，并且顺服圣父而受惩罚，那么很显然，身为天主且自愿服从，那些按他自己决定所做的事既不是痛苦的也不是苦恼的。} 但他没有注意到这里他立即自相矛盾。因为如果他承认祂受惩罚是因为祂决定了这些事并且将自己交托给祂的圣父，那么显然祂确实受了刑罚，而惩罚者加给祂的，不可能不是痛苦的，因为痛苦是不由自主的。但是，如果因为祂甘愿受苦，祂所受的既不是痛苦也不是苦恼，那么他如何承认\Quote{祂受了惩罚？} 他没有觉察到，当\Person{耶稣}一旦藉着诞生取了身体，祂所取的身体是能够承受痛苦以及人类所遭遇的那些苦恼的，如果我们把苦恼理解为无人自愿选择的东西。因此，既然祂自愿取了身体，并非完全与人类血肉之躯不同性质，那么与身体一同，祂也取了它的苦难和苦恼，这些是祂无法避免承受的，因为施加苦难的人有能力把苦恼和痛苦加在祂身上。并且在前面的篇幅中我们已经表明，如果祂不愿意，祂就不会落入人的手中。但祂确实来了，因为祂愿意来，并且因为原先已经显明祂为人类而死对全人类有益。\par

% structure: p0053 source=h2 target=subsection reason=relative-heading-rank; base=h2

\subsection{第二十四章}

此后，为了证明所发生在祂身上的事是痛苦和苦恼的，并且如果祂愿意，祂也无法使之不同，他继续说：\Quote{\Person{耶稣}为何哀恸、哭泣、祈祷，以求逃脱死亡的恐惧，用这样的话表达自己：\InnerQuote{父啊，若是可能，就让这杯离开我吧？}} 现在，在这些话中，请观察\Person{刻尔苏斯}的恶意：他不接受促使福音作者们爱真理的心（这些人本可以缄默那些\Person{刻尔苏斯}认为是可指责的要点，但他们出于多种原因并未省略，任何人在阐释福音时都可在适当之处给出这些原因），却指控福音的陈述，大肆夸大事实，并引用福音中没有记载的话，因为从未见过\Person{耶稣}哀恸。而且他改变了措辞，在表达\Quote{父啊，若是可能，就让这杯离开我吧}中，却没有给出紧接着的话，这些话立即显明\Person{耶稣}对祂圣父的顺服以及祂的心胸宽广，经文如下：\Quote{然而，不要照我的意愿，而照你的意愿成就吧。} 不仅如此，他还假装没有注意到\Person{耶稣}对祂圣父旨意的欣然顺服，即在祂被判受苦的事上所表现的，在宣告\Quote{若是这杯不能离开我，除非我喝了它，就照你的旨意成就吧}中，他的行为就像那些以恶毒心肠聆听圣经的邪恶之人，\Quote{他们昂首谈论恶事。} 因为他们似乎听到了宣告\BibleQuote{我杀死}，并且常常以此作为向我们责备的话；但\BibleQuote{我使人活}这话他们却不记得——整个句子表明那些生活在公开邪恶中并作恶的人被天主处死，并且将一种更好的生命注入他们，就是天主赐给那些死于罪恶之人的生命。同样，这些人听到了\BibleQuote{我击打}的话，却没有看到\BibleQuote{我必医治}，这些话像医生的言语，他剖开身体，施加重伤，以便从中取出有害健康的物质，并不以痛苦和撕裂结束工作，而是通过治疗使身体达到他所期望的健康状态。此外，他们也没有听见宣告的全文\BibleQuote{他撕裂，又包缠}，而只听见这部分\BibleQuote{他撕裂}。同样，\Person{刻尔苏斯}的这个犹太人也这样行，他引用\Quote{父啊，但愿这杯离开我}，却没有加上接下来的话，这些话显明\Person{耶稣}的坚定和祂对受苦的准备。但这些事为从天主的智慧中解释提供了极大空间，并且可以被\Person{保禄}称为\Quote{成全的人}的人合理深思，当他说\Quote{我们在成全的人中讲智慧}，我们暂且略过，只就当前目的有用的事略讲一些。\par

% structure: p0055 source=h2 target=subsection reason=relative-heading-rank; base=h2

\subsection{第二十五章}

我们在前面几页已经提到，耶稣的一些言论指向祂内那\Quote{一切受造物的首生者}，例如\Quote{我是道路、真理、生命}之类的话；另一些言论则属于祂内被认为属于人的那一面，例如\Quote{但现在你们想要杀我，一个把从父那里听到的真理告诉你们的人}。在此，祂相应地描述了人性肉身的软弱性，以及祂人性中精神的准备状态：软弱性体现在\Quote{父啊！若是可能，就让这杯离开我吧}这句话中；精神的准备状态则体现在\Quote{但是，不要照我的意愿，而要照你的意愿}这句话中。既然要遵循我们引文的顺序，请注意，首先只提到一个单一的实例，可以说，表明了肉身的软弱；随后其他更多的实例，表明了精神的意愿。因为\Quote{父啊！若是可能，就让这杯离开我吧}只有一句；而其他更多的话，如\Quote{不要照我的意愿，而要照你的意愿}以及\Quote{我父啊！如果这杯不能离开我，非要我喝不可，就照你的意愿成就吧}。还要注意，话不是\Quote{让这杯离开我}，而是整个表达带有虔诚和敬畏的语气：\Quote{父啊！若是可能，就让这杯离开我吧}。我确实知道对这一段落还有另一种解释：——救世主预见到犹太民族和耶路撒冷城将因犹太人对祂所犯的恶行而遭受的痛苦，纯粹出于对他们的博爱，并希望他们能逃脱即将来临的灾祸，才发出这样的祈祷：\Quote{父啊！若是可能，就让这杯离开我吧}。就好像祂说：\Quote{因为我喝这惩罚之杯，整个民族将被你抛弃，我祈祷，如果可能，就让这杯离开我，好使那对我犯下如此罪行的你的子民，不至于完全被你抛弃。}但是，如果像塞尔苏斯所主张的那样，\Quote{那时耶稣并未遭受任何痛苦或苦难}，那么人们怎能后来引用耶稣为宗教忍受苦难的榜样呢？如果祂并没有经历人类的苦难，而只是\InnerQuote{看起来}如此？\par

% structure: p0057 source=h2 target=subsection reason=relative-heading-rank; base=h2

\subsection{第二十六章}

塞尔苏斯笔下的这位犹太人仍然指控耶稣的门徒编造了这些说法，对他们说：\Quote{即使你们犯了谎言，也没能给你们编造的东西增添可信的色彩。}对此我们不得不说，有一种简单的方法可以掩盖这些事件——那就是根本不记载它们。因为如果福音书没有包含这些事情的记述，谁能责备我们说耶稣在世时说过这样的话呢？塞尔苏斯确实没有看到，同一些人既被耶稣所欺骗，相信祂是天主和预言的主题，又明知这些说法是假的而编造虚构，这是矛盾的。因此，他们的确没有犯编造谎言的罪，而是他们的真实印象如此，并且如实记载了；否则他们就是伪造历史，并且没有持有这些观点，也没有在承认祂为天主时受骗。\par

% structure: p0059 source=h2 target=subsection reason=relative-heading-rank; base=h2

\subsection{第二十七章}

之后他说，某些基督徒信徒，就像醉酒后对自己施以暴力的人一样，已将福音从原有的纯正内容篡改成三倍、四倍甚至多倍，并重新编排，以便能够回应反对意见。我不知道有谁修改过福音，除了马尔西翁的追随者、瓦伦提诺的追随者，以及我想还有路济安的追随者。但这种指控并非针对基督教体系本身，而是针对那些敢于如此戏弄福音的人。正如哲学并不因存在智者派、伊壁鸠鲁派、逍遥派或任何其他持错误观点的人而受指责，真正的基督教也不因有人篡改福音历史并引入反对耶稣教义意义的异端而受指责。\par

% structure: p0061 source=h2 target=subsection reason=relative-heading-rank; base=h2

\subsection{第二十八章}

既然塞尔苏斯笔下的这位犹太人指责基督徒使用预言基督生活事件的先知，我们除了已经谈论过的内容之外，还必须说，他本应像他所说的那样饶恕个人，并解释预言本身，在承认基督徒解释的合理性之后，指出他们如何使用这些预言是可以推翻的。因为这样他就不会显得轻率地基于微小的理由就采取如此重要的立场，尤其是当他断言\Quote{预言与成千上万的其他事情更可信地相符，而不是与耶稣相符}时。他本应认真对待基督徒这一有力的论点，因为这是他们提出的最强有力的论点，并应针对每个具体的预言证明，它比其他事件更有可能适用于耶稣。然而他没有看到，这只是一个反对预言书的人针对基督徒提出的似是而非的论点；但塞尔苏斯在此把一种犹太人不会提出的异议放在了一个犹太人的口中。因为犹太人不会承认预言可以以更大的可能性适用于无数其他事件而不是耶稣；相反，他会努力在给出每个预言的他认为的含义之后，反对基督徒的解释，当然不是以任何令人信服的方式，而只是试图这样做。\par

% structure: p0063 source=h2 target=subsection reason=relative-heading-rank; base=h2

\subsection{第二十九章}

在前文中，我们已经讲到了这一点，即预言基督将有两次来临于人类，因此我们无需回应由一个犹太人提出的反对意见，他说：\Quote{先知们预言将要来的那位是一位大能的君主，万国万军的主。}但我认为，出于犹太人的精神，以及他们苦毒的仇恨，对耶稣毫无根据甚至不可能的诽谤，他补充道：\Quote{先知们没有预言这样的瘟疫。}因为无论是犹太人、塞尔苏斯，还是任何其他人都无法提出任何论证，证明一场瘟疫能使人从作恶转变成合乎自然、以节制及其他美德为特征的生活。\par

% structure: p0065 source=h2 target=subsection reason=relative-heading-rank; base=h2

\subsection{第三十章}

塞尔苏斯也以此反对我们：\Quote{从这样的征兆和错误的解释，以及如此微不足道的证据，没有人能证明他是天主，是天主子。}但他本应列举所谓的错误解释，并证明它们确实是错误，同时通过推理指出证据的微不足道，这样，基督徒如果发现他的某些反对意见似乎可信，便能回答并驳斥他的论点。然而，他关于耶稣所说的话确实应验了，因为耶稣是一位大能的君主，尽管塞尔苏斯拒绝承认这是事实，尽管最清晰的证据证明这在耶稣身上是真实的。他说：\Quote{因为太阳既照亮其他一切物体，便首先使自己可见，所以天主子也应当如此。}我们要回答说，他正是这样做的；因为正义在他的日子里升起，和平丰盛，从祂诞生开始，天主就为祂的教导预备万民，使他们同在一王之下，即罗马人的皇帝，这样，由于许多王国并存而导致的万民不和，就不会使耶稣的宗徒们更难以完成他们师傅所命令的任务，当祂说：\Quote{你们要去使万民成为门徒}。此外，确实耶稣诞生于\Person{奥古斯都}统治时期，奥古斯都可谓将地上众多民族融合为一个君主国。如果存在许多王国，那将阻碍耶稣的教导传播到全世界；不仅由于上述原因，还因为各地的人必须从事战争，为自己的祖国而战，这在奥古斯都时代之前以及更遥远的时期都是必要的，例如当伯罗奔尼撒人和雅典人互相争战，以及其他民族也类似地争战时。那么，和平的福音教导，甚至不允许人对仇敌进行报复，如何在全世界得胜呢？除非在耶稣来临之际，一种更温和的精神已被引入万事的运行中。\par

% structure: p0067 source=h2 target=subsection reason=relative-heading-rank; base=h2

\subsection{第三十一章}

他接着指责基督徒\Quote{犯了诡辩推理的罪，说天主子就是\OriginalTerm{圣言}{\GreekText{Λόγος}}本身。}他认为这加重了指控，因为\Quote{当我们宣称\OriginalTerm{圣言}{\GreekText{Λόγος}}是天主子时，我们呈现的并不是一个纯净圣洁的\OriginalTerm{圣言}{\GreekText{Λόγος}}，而是一个最卑劣的人，受到鞭打和钉十字架的惩罚。}关于这一点，我们在前文已简短回应了塞尔苏斯的指控，那里表明基督是一切受造物中的首生者，祂取了身体和人的灵魂；天主在世上创造了如此伟大的事物，并给了命令，万物便受造；接受命令的那一位就是天主圣言。既然在塞尔苏斯的作品中说这些话的是一个犹太人，那么引用这句话并不不合宜：\BibleQuote{祂派遣了祂的圣言，医治了他们，救他们脱离了毁灭。}——我们刚才也提到过这一段。现在，虽然我与许多自称博学的犹太人交谈过，但从未听到任何人赞同圣言是天主子的说法，就像塞尔苏斯声称的那样，他借着犹太人的口这样说：\Quote{如果你们的\OriginalTerm{圣言}{\GreekText{Λόγος}}是天主子，我们也同样表示同意。}\par

% structure: p0069 source=h2 target=subsection reason=relative-heading-rank; base=h2

\subsection{第三十二章}

我们已经表明，耶稣既不能被看作傲慢的人，也不能被看作巫术师；因此没有必要重复我们之前的论点，以免在回应塞尔苏斯的重复时，我们自己也不必要地重复。现在，在挑剔我主耶稣的族谱时，有些点甚至令基督徒感到困难，并且由于族谱之间的差异，一些人对族谱的正确性提出质疑，但塞尔苏斯甚至没有提到这些点。因为塞尔苏斯，一个真正的吹嘘者，自称通晓关于基督信仰的一切，却不知道如何巧妙地质疑圣经的可信性。但他断言，\Quote{族谱的编纂者出于骄傲，使耶稣成为从第一个人和犹太诸王的后裔。}他认为自己提出了一个显著的指控，补充道：\Quote{木匠的妻子不可能不知道这个事实，如果她真有如此显赫的出身的话。}但这与问题有何相干？就算她不知道自己的出身，那又如何？假设她\Emph{是}不知道，她的无知如何能证明她不是从第一个人所出，或不能从犹太诸王追溯她的起源？塞尔苏斯是否认为穷人必定总是从穷人祖先所出，或国王总是生于国王？但花费时间在此类论据上似乎是愚蠢的，因为众所周知，即使在今天，有些比玛利亚更穷的人却出自富贵显赫的祖先，而国家的统治者和国王却来自没有名声的人。\par

% structure: p0071 source=h2 target=subsection reason=relative-heading-rank; base=h2

\subsection{第三十三章}

\Quote{然而，}\Person{克尔苏斯}继续说道，\Quote{耶稣作为一位神，行了什么伟大的事迹？他使他的敌人蒙羞，还是使他针对他的阴谋落得可笑的下场？} 对于这个问题，虽然我们能够指出发生在他身上的事那些显著而神奇的性质，但我们还能从其他来源提供答案，除非从福音叙述中，其中说：\Quote{地大震动，磐石崩裂，坟墓开了，圣殿的幔子从上到下裂为两半，白日黑暗，太阳不发光。} 但如果\Person{克尔苏斯}在认为能从福音叙述中找到指控基督徒的把柄时就相信它们，而当它们确立耶稣的神性时却拒绝相信，我们对他的回答是：\Quote{先生，要么全然不信福音叙述，那么你就不能再设想能基于它们提出指控；要么，如果你相信它们的陈述，就应惊奇地注视那降生成人的天主圣言，并愿意施恩于全人类。而这一点显明了耶稣工作的崇高：直至今日，凡天主所愿意的人，都因祂的名而获得治愈。至于\Person{提庇留凯撒}时代（耶稣似乎是在他的统治下被钉死的）的日蚀，以及当时发生的大地震，我想，\Person{弗勒贡}也在他的\WorkTitle{编年史}第十三或十四卷中写到了。}\par

% structure: p0073 source=h2 target=subsection reason=relative-heading-rank; base=h2

\subsection{第34章}

\Person{克尔苏斯}笔下的这位犹太人，照他的想象嘲笑耶稣，被描述为熟悉\Person{欧里庇得斯}的\WorkTitle{酒神的女信徒}，其中\Person{狄俄倪索斯}说：\par

\Quote{每当我想，神自己就会释放我。}\par

犹太人并不十分熟悉希腊文学；但假设有一个如此精通它（以致引用它对他合适）的犹太人，（那么）耶稣\Emph{不能}释放自己，因为祂没有这样做，这又如何（能成立）呢？因为他相信我们自己的圣经：\Person{伯多禄}在被囚禁在监狱中后获得自由，一位天使解开了他的锁链；\Person{保禄}在\Place{马其顿的斐理伯}与\Person{息拉}一同被枷锁束缚，被神力释放，监狱的门开了。但很可能\Person{克尔苏斯}嘲笑这些记述，或者他从未读过它们；因为他大概会回答说，有些术士能够通过咒语解开锁链和打开门，这样他就会把我们的历史中所述的事比作术士的作为。\Quote{但是，} 他继续说，\Quote{连定祂罪的人也没有遭遇任何灾难，像\Person{彭透斯}那样，即疯狂或被撕裂。} 然而他不知道，定祂罪的与其说是\Person{比拉多}（他知道\Quote{犹太人因嫉妒而把祂解来}），不如说是犹太民族，它被天主定罪、撕裂、分散到全地，程度远超\Person{彭透斯}所遭遇的。此外，他为什么故意忽略关于\Person{比拉多}妻子的记载，她看见了一个异象，深受感动，就派人去对她丈夫说：\Quote{你不要管那义人的事，因为今天我因祂的缘故在梦中受了许多苦？} 再者，\Person{克尔苏斯}默然略过耶稣神性的证据，却试图从福音叙述中侮辱祂，提到那些嘲笑耶稣的人，给祂穿上紫袍、戴上茨冠，把芦苇放在祂手里。现在，\Person{克尔苏斯}，你从哪里得到这些叙述，除非是从福音叙述？你因而认为它们是适合侮辱的材料；而记录它们的人并不认为你和你这样的人会嘲笑它们，而是认为其他人会从中得到一个榜样，如何鄙视那些因祂的信仰而嘲笑和讥讽祂的人，祂正当地为信仰舍生。更应钦佩他们的真诚，以及那甘愿为人类承担这些事的存有，祂以恒心和忍耐忍受了这一切。因为并没有记载祂发出任何悲叹，或者在被定罪后做了什么或说了什么不当的话。\par

% structure: p0077 source=h2 target=subsection reason=relative-heading-rank; base=h2

\subsection{第35章}

但针对这一反对意见：\Quote{若不是从前，那么现在至少为何祂不显示一些神性，从这种羞辱中解救自己，并向那些侮辱祂和祂的圣父的人报仇？} 我们必须回答说，这就好比我们向那些接受天主眷顾论并相信异象的希腊人说：为什么天主不惩罚那些侮辱神性并颠覆天主眷顾论的人？因为正如希腊人会回答这样的反对意见，我们也会以同样或更有效的方式回答。不仅有从天而来的异象——日蚀——还有其他的奇迹，表明被钉者拥有某种神圣的、比大多数人更伟大的东西。\par

% structure: p0079 source=h2 target=subsection reason=relative-heading-rank; base=h2

\subsection{第36章}

继此，刻尔苏斯说：\Quote{被钉十字架的耶稣身体中所流的\OriginalTerm{神血}{\GreekText{ἰχώρ}}是什么性质？是否\InnerQuote{如不朽诸神身体中所流者}？}他以嘲讽的态度提出这个问题；但我们将根据福音书的严肃叙述（尽管刻尔苏斯可能不喜）证明，从耶稣身体流出的并非神话式荷马式的神血，而是祂死后，\Quote{有一个兵士用枪刺开祂的肋旁，立时流出了血和水。那看见的，就作证，而他的见证是真实的，并且他知道自己所说的是真实的。}\BibleRef{若 19:34-35}如今，在其他尸体中，血液凝固，纯水不会流出；但耶稣尸体的奇妙之处在于，围绕尸体，从肋旁流出了血和水。然而，如果这个刻尔苏斯——他为了找控告耶稣和基督徒的借口，甚至从福音中摘录被错误解释的段落，却对耶稣神性的证据缄默不语——愿意听从神圣的异兆，就让他读福音，看看连百夫长和与他一同看守耶稣的人，见了地震和所发生的事，都非常害怕，说：\Quote{这人真是天主子。}\BibleRef{玛 27:54}\par

% structure: p0081 source=h2 target=subsection reason=relative-heading-rank; base=h2

\subsection{第三十七章}

其后，这位从福音叙述中摘录那些他认为可以据以指控的陈述的人，以醋和苦胆作为责备耶稣的口实，说\Quote{他张着嘴急忙去喝它们，不能像任何普通人经常忍受的那样忍受口渴。}此事本身有一种独特而比喻性的解释；但在此场合，先知们预言了这件事这一事实可被接受为对这一异议的更常见的回答。因为在\BibleBook{}{圣咏}第六十九篇中，指着基督写着：\Quote{他们拿苦胆给我当食物；我渴了，他们拿醋给我喝。}\BibleRef{咏 69:22}现在，让犹太人说说，这先知性的文字代表谁在说这些话；并让他们从历史中举出一个人，他得了苦胆作为食物，并有人给他醋喝。他们敢断言他们仍期待来临的基督可能会遭遇这样的境况吗？那么我们就说：有什么阻止预言已经应验呢？因为正是这个预言在许多世纪以前就发出了，并且连同其他先知性的言论，足以引导公正审查全部事实的人得出结论：耶稣就是曾被预言为基督和天主子的一位。\par

% structure: p0083 source=h2 target=subsection reason=relative-heading-rank; base=h2

\subsection{第三十八章}

接下来说的几句话是：\Quote{你们这些诚恳的信徒责备我们，因为我们不承认这人是天主，也不同你们一样相信祂忍受了这些（苦难）是为了人类的益处，为使我们也能藐视惩罚。}对此，我们说：我们责备犹太人，他们曾受法律和先知（预言基督来临）的训练，因为他们既不反驳我们摆在他们面前证明祂是默西亚的论证（提出这样的反驳作为他们不信的辩护），又不在不提出反驳的情况下相信祂——祂是预言的对象，并借着祂的门徒即使在祂显现于肉身之后也清楚表明，祂为了人类的益处而忍受了这些事；祂第一次来临的目的，不是要在祂教训他们并向他们指明他们的责任之前定他们和他们的行为的罪，也不是要惩罚恶人、拯救善人，而是要以非凡的方式，借神圣大能的证据，将祂的道理传布给全人类，正如先知们也描绘了这些事情。此外，我们还责备他们，因为他们不相信那位显示内在大能者，反而声称祂借着魔王贝耳则步从人的灵魂中驱逐魔鬼；并且因为他们毁谤祂的仁爱性格——祂不忽略犹太地任何一座城，甚至不忽略任何一个村庄，为要在各处传扬天主的国——指控祂过着漂泊无定的生活，在可耻的身体中过着忧虑的存在。但为了一切能理解祂的人的益处而忍受这样的劳苦，并无羞耻可言。\par

% structure: p0085 source=h2 target=subsection reason=relative-heading-rank; base=h2

\subsection{第三十九章}

而刻尔苏斯笔下的这个犹太人所说的以下断言，即耶稣\Quote{生前没有赢得任何人，甚至连祂自己的门徒也没有赢得，就遭受了这些刑罚和苦难}，怎能显得不是明显的谎言呢？因为若不是许多群众听从祂、跟随祂，并被领到荒野——不仅因祂的话总是适合听众而具有说服力，而且祂也借着奇迹打动了那些不被祂的话感动而相信的人——那么犹太人的司祭长、长老和经师们对祂所激起的嫉妒从何而来？而说\Quote{祂甚至连祂自己的门徒也没有赢得}，这岂不是明显的谎言吗？祂的门徒当时确实表现出一些因怯懦恐惧而来的人性软弱——因为他们尚未被训练到表现出完全勇敢的程度——但他们绝没有放弃他们对祂是基督的判断。因为伯多禄在否认后，意识到自己堕入了何等深的邪恶，\Quote{出去痛哭。}\BibleRef{玛 26:75}而其他门徒，虽然因耶稣所遭遇的事而惊慌（因为他们仍钦佩祂），却因祂荣耀的显现，比以前更加坚定地相信祂是天主子。\par

% structure: p0087 source=h2 target=subsection reason=relative-heading-rank; base=h2

\subsection{第四十章}

此外，\Person{瑟尔苏}以一种极不具哲学精神的方式想象，我们的主的超越性不在于宣讲救恩和纯洁的道德，而在于与祂所取的人性相悖的行为，即不死——尽管祂已取 \OriginalTerm{有死性}{mortality}；或即便死，至少也不应以那种死亡作为榜样，使那些要学习如何为宗教而死的，以及那些在宗教与不敬神的错误观点面前，凭借其帮助勇敢行事的人——这些人视宗教人为完全不敬神，却认为那些对天主有错误认识的人最敬神，并且将不可磨灭的（植入人心中的）天主的观念应用于一切而非天主本身，尤其是当他们急切地摧毁那些全心（甚至至死）顺服于超越万物的独一天主之清晰证据的人时——视为榜样。\par

% structure: p0089 source=h2 target=subsection reason=relative-heading-rank; base=h2

\subsection{第四十一章}

在犹太人的角色中，\Person{瑟尔苏}继续指责\Person{耶稣}，声称\Quote{祂并未显明自己纯洁无恶}。愿\Person{瑟尔苏}说明，我们的主在何种\Quote{恶}上未显明自己纯洁。若他意指祂在严格意义上的\Quote{恶}上不洁，就请清楚证明祂身上有任何恶行。但若他认为贫穷、十字架以及恶人的阴谋是恶，那么显然他会说\Person{苏格拉底}也遭遇了恶，因为他未能显明自己免于恶。而希腊人中还有许多穷人，他们献身哲学研究，自愿接受贫穷生活，这一点许多希腊人都从记载中知晓：\Person{德谟克利特}任由财产成为羊群的牧草；\Person{克拉特斯}通过将卖物所得献给底比斯人而获得自由；甚至\Person{第欧根尼}因极端贫穷而住在木桶中。然而，在稍有理智的人看来，\Person{第欧根尼}并不因此被认为处于罪恶状态。\par

% structure: p0091 source=h2 target=subsection reason=relative-heading-rank; base=h2

\subsection{第四十二章}

此外，既然\Person{瑟尔苏}认为\Quote{\Person{耶稣}并非无可指责}，就让他指出任何一位信奉祂教义的人，所记载的任何能真正提供指责\Person{耶稣}之理由的事情；或者，若他并非从这些人获得指控祂的材料，就让他说明从何处得知使他断言祂并非无可指摘之事。然而，\Person{耶稣}履行了祂所应许的一切，并使祂的信徒受益。我们不断看到祂在事情发生前所预言的一切都应验了：即祂的福音要传遍全世界，祂的门徒要往万民中去宣讲祂的教义；而且，他们要因祂的教导而被带到总督和君王面前。我们因此惊叹祂，并且每日坚定对祂的信德。我不知道\Person{瑟尔苏}会要求以何种更大或更令人信服的证据来使祂的预言得以确认；除非，正如似乎所发生的，他不理解圣言成了人\Person{耶稣}，就要求祂不受任何人性软弱的影响，也不成为人们应当如何忍受生活灾难的光辉榜样——尽管这些在\Person{瑟尔苏}看来是最可悲可耻的事情，因为他视劳苦为最大之恶，以快乐为完善之善——这是任何承认天主眷顾之教义、并承认勇敢、坚韧与宽宏为德行的哲学家所不接受的观点。因此，\Person{耶稣}并未因祂所受的苦难而玷污祂所是的信仰对象，反而在那些赞同男子气概的人中，以及在那些被祂教导真正而恰当的幸福生活并不在此世、而是在祂所说的\Quote{来世}的人中，巩固了信仰；而在所谓的\Quote{现世}中，生活是一种灾难，或至少是灵魂的第一场也是最大的争战。\par

% structure: p0093 source=h2 target=subsection reason=relative-heading-rank; base=h2

\subsection{第四十三章}

\Person{瑟尔苏}接着对我们说了如下话：\Quote{你们大概不会说，他在未能赢得世上之人后，便下到冥府去赢得那里的人。}但无论他是否喜欢，我们断言：不仅\Person{耶稣}在肉身时赢得了不少人——甚至人数如此之多，以致因跟随者众多而有人密谋反对祂；而且当祂成为没有肉身遮盖的灵魂时，祂也住在那些没有肉身遮盖的灵魂中间，使其中愿意者归向祂，或那些祂出于自己所知的理由认为更适合于这样做的灵魂归向祂。\par

% structure: p0095 source=h2 target=subsection reason=relative-heading-rank; base=h2

\subsection{第四十四章}

\Person{瑟尔苏}接着以无法形容的愚蠢说道：\Quote{如果你们捏造出荒谬的辩护，并因此被荒唐地迷惑，却自以为真正作了好的辩护，那么是什么阻止你们把那些同样被定罪、并悲惨死去的人视为（比\Person{耶稣}）更伟大、更神圣的来自天上的使者呢？}现在，显然并清楚地，\Person{耶稣}所受的如所述之事与那些因邪术或任何其他指控而悲惨死去的人之间毫无相似之处，这对每个人都是显而易见的。因为没有人能指出任何一个邪术者的行为，能使人灵远离那些在人间盛行的许多罪过，以及（世上）罪恶的洪流。但是，既然\Person{瑟尔苏}笔下的这个犹太人将他与盗贼相比，并说：\Quote{任何一个同样无耻的家伙，都可能就一个甚至被惩罚所追上的强盗和杀人犯说，这样的人不是强盗，而是一个神，因为他向他的同伙预言他会遭受他实际所受的惩罚。}那么，首先可以回答：我们之所以持有关于\Person{耶稣}的看法，并因此信赖祂，视祂为从天主那里降临到我们中间的一位，并非因为祂预言了祂会遭受这些事。其次，我们断言，这比较本身早已在\WorkTitle{福音}中被预言；因为天主被恶人归为罪犯，他们宁愿要一个\Quote{杀人犯}（一个因叛乱和杀人而被下在监里的人）释放给他们，而把\Person{耶稣}钉在十字架上，并将祂钉在两个强盗之间。的确，\Person{耶稣}在祂的真门徒和为真理作证者中间，常常与盗贼一同被钉十字架，并遭受他们在人中所受的同样谴责。并且我们说，如果那些因对天主的虔诚而遭受各种伤害和死亡，以便按照\Person{耶稣}的教导保持纯全无染的人，与盗贼有任何相似之处，那么显然，\Person{耶稣}——这种教导的创立者——被\Person{瑟尔苏}比作一伙盗贼的首领，也就很有道理了。但是，那位为全人类的公共福祉而死的，和那些因信仰而受苦、并独自在所有人中被逼迫，因为他们认为这是尊敬天主的正确方式的人，都不是按照正义被处死的；\Person{耶稣}受迫害，祂的迫害者也并非没有犯下不虔敬的罪责。\par

% structure: p0097 source=h2 target=subsection reason=relative-heading-rank; base=h2

\subsection{第四十五章}

但请注意他论证的肤浅之处，关于\Person{耶稣}先前的门徒，他说：\Quote{其次，那些在他活着时与他为伍、聆听他的声音、并享受他作为老师的教导的人，一看到他遭受惩罚和死亡，既没有与他同死，也没有为他而死，甚至没有被诱导去藐视惩罚，反而否认他们是他的门徒，而现在你们却与他同死。}他在这里相信门徒们还是初学和不完善时所犯的罪——这罪记录在\WorkTitle{福音}中，并这罪记录在\WorkTitle{福音}中，实际上是犯了，以便他有指控\WorkTitle{福音}的借口；但他们在犯罪后的正直行为，那时他们在犹太人面前勇敢行事，遭受了他们无数残酷的对待，并最终为\Person{耶稣}的教义而死，他却沉默不语。因为他既不愿听\Person{耶稣}对\Person{伯多禄}的预言：\BibleQuote{当你年老时，你要伸出手来，}等等，\WorkTitle{圣经}对此补充道：\BibleQuote{他说这话，是指明\Person{伯多禄}要以怎样的死来光荣天主；}他也不愿听\Person{若望}的兄弟\Person{雅各伯}——一位宗徒，一位宗徒的兄弟——是如何为基督的教义被\Person{黑落德}用刀所杀；他也不愿听\Person{伯多禄}和其他宗徒因\WorkTitle{福音}所表现的许多勇敢事例，以及\BibleQuote{他们怎样从公议会面前出去，因被鞭打而欢喜，因为他们配为\Person{耶稣}的名受辱，}并且超越了希腊人所记述的关于他们的哲学家的刚毅和勇敢的许多事例。因此，从一开始，\Person{耶稣}就在听者中灌输这一训诫，教导人轻视世人所热切追求的生命，而认真活出那类似天主的生命。\par

% structure: p0099 source=h2 target=subsection reason=relative-heading-rank; base=h2

\subsection{第四十六章}

但是，\Person{瑟尔苏}笔下的这个犹太人怎能逃脱说谎的指控呢？他说\Person{耶稣}\Quote{在世时只赢得了十来个最卑微的船夫和税吏，而且甚至连这些人都没有完全赢得。}现在，犹太人自己也会承认，祂吸引过来的人不止十个，也不止一百、一千，而是一次五千，另一次四千；祂吸引他们到这样的程度，以至于他们跟随祂进入旷野——只有旷野才能容纳那些通过\Person{耶稣}而信仰天主的聚集群众——在那里祂不仅对他们讲论，还向他们显示祂的作为。现在，由于他的重复，他也迫使我们重复，因为我们谨慎地防止被认为忽略了他提出的任何指控；因此，关于我们面前的事，按照他著作的顺序，他说：\Quote{坚持认为，如果他本人在世时没有赢得一个人支持他的观点，死后任何愿意的人却能赢得如此众多的人，这难道不是荒谬至极吗？}而他应该按照真理说：如果在他死后，不仅仅是那些愿意的人，而是那些既有意愿又有能力的人，能够赢得如此众多的皈依者，那么，当他活着的时候，通过他话语和行为的更大能力，他自己赢得更多倍的追随者，岂不是更合情理吗？\par

% structure: p0101 source=h2 target=subsection reason=relative-heading-rank; base=h2

\subsection{第四十七章}

此外，他提出一种他自己的说法，仿佛是对他某问题的答复，他问：\Quote{你有什么推理过程使你们视他为天主子？}因为他让我们回答说：\Quote{我们被赢到祂一边，是因为我们知道祂所受的受苦刑是为了消灭邪恶之父。}然而，我们是因无数其他理由而被赢到祂的教导中，我们在前文只提到其中最小部分；但若天主准许，我们将继续列举它们，不仅是在处理策尔苏斯所谓的\WorkTitle{真言}时，也在许多其他场合。而且，好像我们说过我们认为祂是天主子是因为祂受苦刑，他问：\Quote{那么怎样？难道不是有许多其他人也受过刑罚，而且并不少耻辱吗？}在此策尔苏斯表现得如同福音最可鄙的敌人，如同那些想象从我们关于被钉十字架的耶稣的历史中必然得出我们应当敬拜那些被钉十字架的人的人一样。\par

% structure: p0103 source=h2 target=subsection reason=relative-heading-rank; base=h2

\subsection{第四十八章}

此外，策尔苏斯无法否认耶稣所行的奇迹，已多次诽谤地称它们为巫术；而我们也多次尽力回应了他的说法。现在他让我们说：\Quote{我们认为耶稣是天主子，因为祂治好了瘸子和瞎子。}他又补充说：\Quote{而且，如你们所断言，祂复活了死人。}祂治好瘸子和瞎子，因此我们视祂为基督、天主子，这从先知书中的话对我们显而易见：\BibleQuote{那时瞎子的眼睛要打开，聋子的耳朵要开启；那时瘸子要像鹿一样跳跃。}祂复活死人也并非撰写福音之人的虚构，这可由以下事实表明：如果那是虚构，就会有\Emph{许多}人复活，而且是那些已埋葬多年的人。但既然不是虚构，有关此事所记载的人很容易数清：即会堂长的女儿（我不明白为何祂说：\BibleQuote{她不是死了，而是睡着了}，论到她说了些不适用于所有死者的话）；还有寡妇的独子，祂动了怜悯之心使他复活，使抬棺的人站住；第三个例子是拉匝禄，他已埋葬了四天。现在，关于这些事例，我们要对所有心地正直的人，特别是对犹太人这样说：正如在厄里叟先知时代有许多癞病人，除了叙利亚的纳阿曼之外没有一个被治好；在厄里亚先知时代有许多寡妇，厄里亚只被派往漆冬的匝尔法特（因为那里的寡妇因天主旨意而配得先知在饼的事上行的奇迹）；同样，在耶稣时代有许多死人，但只有那些圣言知道适合复活的人从坟墓中起来，为使主所行的事不只是某些事物的象征，而是借这些行动本身，祂赢得许多人归向福音的奇妙教导。我还要说，按照耶稣的应许，祂的门徒行了比这些仅能感知的耶稣奇迹更大的事。因为灵魂瞎子的眼睛常常被开启；对美言充耳不闻的耳朵乐意聆听天主的教导和与祂共度的幸福生活；还有许多在内心里（如圣经所称\OriginalTerm{内心的人}{inner man}）脚瘸的人，如今被圣言医治，不只是跳跃，而是像鹿一样跳跃，鹿是一种与蛇为敌的动物，比一切毒蛇的毒液更强。这些被医好的瘸子从耶稣获得能力，用他们从前瘸的脚践踏邪恶的蛇和蝎子，并普遍践踏仇敌的一切能力；虽然他们践踏它，却不受伤害，因为他们也变得比一切邪恶和魔鬼的毒更强。\par

% structure: p0105 source=h2 target=subsection reason=relative-heading-rank; base=h2

\subsection{第四十九章}

耶稣为要使祂的门徒不仅不留意一般的术士，以及以任何其他方式声称能行奇迹的人——因为祂的门徒不需要受这种警告——而且也不留意那些自称为天主的基督、并企图藉某些表面的奇迹将耶稣的门徒吸引到自己一边的人，在某处说了这样的话：\BibleQuote{那时，若有人对你们说：看，基督在这里！或说：基督在那里！你们不要相信。因为将有假基督、假先知出现，显大奇事和异迹，以致如果可能，连被选的人也要被欺骗。看，我预先告诉你们了。所以，若有人对你们说：看，祂在旷野里，你们不要出去；看，祂在内室里，你们不要相信。因为闪电怎样从东方发出，直照到西方，人子的来临也要这样。}在另一处又说：\BibleQuote{到那一天，许多人要对我说：主啊，主啊！我们不是因祢的名吃过喝过，因祢的名驱过魔鬼，行过许多奇迹吗？那时我就要向他们声明说：你们这些作恶的人，离开我！}但塞耳苏想要把耶稣的奇迹比作人类巫术的行为，便明确地说：\Quote{啊，光明与真理！祂用自己的声音明确宣告，正如你们自己所记载的，将会有别的人到你们那里去，行类似的奇迹，他们是恶人和术士；祂称运用这种手段的那一位为撒殚。所以耶稣自己并不否认这些作为至少并非神圣，而是恶人之行；祂被真理的力量所迫，同时不仅揭露了别人的行为，也定了自己同样的行为。那么，从同样的行为得出结论，说这一个是天主，而其他的是术士，岂不是可怜的推论吗？为什么别人因这些行为应被视为恶人，而这个人却不是？他们竟有祂作为反对祂自己的见证！因为祂自己已经承认，这些不是神圣本性的作为，而是某些欺骗者和彻底邪恶之人的发明。}现在请观察，塞耳苏是否显然藉这些言论诽谤了福音，因为耶稣关于那些将行奇事异迹的人所说的话，与塞耳苏所声称的完全不同。因为如果耶稣只是简单地告诫祂的门徒要提防那些声称能行奇迹的人，而没有说明他们会自称是什么，那么他的怀疑或许还有一点根据。但既然耶稣要我们提防的那些人是自称为基督——这不是术士提出的要求——而且祂说甚至一些过邪恶生活的人也会因耶稣的名字行奇迹，从人身上驱逐魔鬼，那么这些人身上的巫术，或者任何对这事的怀疑，反而可以说完全被排除了，而基督的神性得以确立，祂门徒的神圣使命也得以确立；因为有可能一个使用祂的名字、并被某种未知力量推动而假装自己是基督的人，似乎行出类似耶稣的奇迹，而另一些人则因祂的名字行出类似祂真门徒的作为。\par

此外，保禄在《得撒洛尼后书》中指示了那恶人将怎样在适当的时候被显现出来：\BibleQuote{那无法无天的人，即丧亡之子，他反对并高举自己于一切被称为神或受崇拜者之上，以致坐在天主的殿中，显示自己就是神。}他又对得撒洛尼人说：\BibleQuote{你们也知道那阻止他的是什么，使他到了自己的时候才被显现。因为那不法的奥秘已经运行，只待那阻止者从中间被除去；那时，那不法的人就要被显现，主耶稣要以口中的气杀死他，以来临的显现消灭他。他来临的，是依靠撒殚的能力，具有各种异能、征兆和欺诈的奇迹，以及各种不义的欺骗，在那些丧亡的人身上。}保禄在说明为何允许那不法的人继续存在的原因时说：\BibleQuote{因为他们没有接受对真理的爱，为得救恩。为此，天主给他们一种强烈的错觉，使他们相信谎言，为使所有不相信真理、反而喜爱不义的人都受审判。}现在让任何人说，福音或宗徒的著作中是否有任何陈述能引起怀疑，认为其中包含对巫术的预言。凡愿意的人，也可以在达尼尔书中找到关于假基督的预言。但塞耳苏歪曲了耶稣的话，因为耶稣并没有说将有别的人行类似祂的奇迹，而他们是恶人和术士，尽管塞耳苏断言祂说了这样的话。因为埃及术士的能力与梅瑟的天赐恩宠并不相似，结果清楚证明了前者的行为是巫术的效果，而梅瑟的行为是藉神能完成的；同样，假基督以及那些假称作为基督门徒能行奇迹之人的作为，被称为谎言的奇事和征兆，以各种不义的欺骗在那些丧亡的人身上得胜；而基督及其门徒的作为，其果实不是欺骗，而是人灵魂的救恩。谁又能理性地主张，一种改善的道德生活——它日渐减少人过犯的数目——能出自一个欺骗的体系呢？\par

% structure: p0108 source=h2 target=subsection reason=relative-heading-rank; base=h2

\subsection{第五十一章}

诚然，\Person{切尔苏斯}对圣经略有所知，因而让\Person{耶稣}说：\Quote{是某个\Person{撒殚}策划了这些诡计}；但他却回避了问题，断言\Quote{\Person{耶稣}并未否认这些工作毫无神圣性，而是来自恶人}，因为他将不同种类的事物混为一谈。正如狼与狗虽在体形和声音上有些相似，却并非同类，普通斑鸠与鸽子也非同类，同样，因天主大能所行的事与巫术的效果之间毫无相似之处。我们还可进一步回应\Person{切尔苏斯}的诽谤：难道藉巫术由邪恶魔鬼所行的才被视为奇迹，而由神圣、圣洁的本性所行的就不是奇迹吗？难道人类生活只能忍受更坏的事物，却从未接受更好的事物吗？在我看来，我们必须确立一个普遍原则：正如无论何处，任何邪恶之物想要与善物具有相同本性，就必定有某种善物与之对立；同样，与巫术所引起的事物相对，人类生活中也必然存在一些由神圣能力所产生的事物。由此推论，我们要么否定两者，断言两者都不存在；要么承认一方，尤其是邪恶的一方，同时也要承认善的现实。若有人断言有些工作是藉巫术而行，却不承认也有工作是神圣能力之果，在我看来，他就像承认存在诡辩和貌似有理的论证——这些论证看似确认真理，实则暗中破坏真理——却否认真理在人类中任何地方有家园，或否认有一种不同于诡辩的辩证法。但我们一旦承认，魔术和巫术（它们从邪恶魔鬼获得力量，这些魔鬼被复杂的咒语束缚，并服从于巫术师）的存在必然意味着人类中有些工作来自神圣能力，那么，我们为何不通过那些自称行奇迹者的生活、道德以及其奇迹的后果（即它们是有害于人还是有助于行为的改善）来检验他们呢？例如，哪个邪恶魔鬼的仆役能行这样的事？又靠什么咒语和巫术？另一方面，谁拥有纯洁圣洁的灵魂、精神（我想还有身体），接受神圣的精神，并施行这样的工作以造福人类、引导他们信仰真天主？如果我们必须探究（不被奇迹本身所迷惑）谁藉善的能力、谁藉恶的能力行这些事，以便我们不致不加区别地诽谤一切，也不致把一切都当作神圣的来赞美和接受，那么，从\Person{梅瑟}和\Person{耶稣}时代所发生的事（整个民族因他们的奇迹而建立）中，岂不显然可见这些人是藉神圣能力完成了他们所记载的事迹吗？因为邪恶和巫术不会引导整个民族不仅超越人手所造的偶像和像，而且超越一切受造物，上升到宇宙天主的非受造之源。\par

% structure: p0110 source=h2 target=subsection reason=relative-heading-rank; base=h2

\subsection{第五十二章}

但既然在\Person{切尔苏斯}的论著中是一个犹太人说了这些话，我们要对他说：请吧，朋友，你为什么相信你们著作中记载的、天主藉\Person{梅瑟}之手所行的那些工作是真正神圣的，并努力反驳那些诽谤说它们就像\Place{埃及}术士的工作一样是巫术的产物；而你却效仿你们的\Place{埃及}对手，指控\Person{耶稣}所行的、你也承认确实发生的那些工作不是神圣的呢？如果\Person{梅瑟}奇迹的最终结果——一个完整民族的建立——明显证明是天主在希伯来立法者\Person{梅瑟}时代使之成就，那么，为什么\Person{耶稣}所行的远超过\Person{梅瑟}的更大工作，不更应该显明是神圣的吗？前者带领他自己民族的人——\Person{亚巴郎}的后裔，他们遵守遗传的割损礼，并小心持守\Person{亚巴郎}的习俗——出离\Place{埃及}，为他们制定了你们认为神圣的法律；而后者则敢于承担更伟大的事业，在已有的体制、祖传习俗以及符合现有法律的生活方式之上，再加添一种合乎福音的体制。正如为了\Person{梅瑟}不仅在长老中而且在民众中获得信任，必须行那些记载的奇迹；同样，\Person{耶稣}为了被那些已学会要求征兆和奇迹的民众所信，不也需要行这样一些奇迹——因它们比\Person{梅瑟}的奇迹更伟大、更神圣——以致能使人们从犹太人的寓言和流传于他们中的人类传统中转变过来，承认这位教导并行这样事情的人比先知更伟大吗？因为那被先知宣告为\Person{基督}并为人类救主的人，怎么会不比先知更伟大呢？\par

% structure: p0112 source=h2 target=subsection reason=relative-heading-rank; base=h2

\subsection{第五十三章}

事实上，克耳苏斯笔下的这个犹太人所提出的反对信仰耶稣之人的一切论点，依同理也可作为反对梅瑟的控诉理由：因此，声称耶稣和梅瑟所行的巫术彼此相似，这并没有区别——就克耳苏斯笔下这个犹太人的言辞而言，两者都同样可被指控；例如，当这个犹太人论及基督时说：\Quote{但是，光明与真理啊！耶稣亲口明说，正如你们自己所记载的，将另有他人出现在你们当中，他们要行像我一样的奇迹，但他们是恶人和巫士。}那么，某个希腊人或埃及人，或其他任何不信犹太人的一方，可能会论及梅瑟说：\Quote{但是，光明与真理啊！梅瑟亲口明说，正如你们所记载的，将另有他人出现在你们当中，他们要行像我一样的奇迹，但他们是恶人和巫士。因为你们的律法上记载：\BibleQuote{如果你们中出现一位先知，或作梦的人，给你提供一种神迹或奇事，而那神迹和奇事发生了，他对你说：我们去追随别的神，你们素不相识的神，事奉他们吧！你不要听信那个先知或作梦的人的话。}\BibleRef{申 13:1-3}}等等。此外，他曲解耶稣的话说：\Quote{而他把想出这些事的人称为撒殚。}若将此应用于梅瑟，有人可能会说：\Quote{而他把想出这些事的人称为作梦的先知。}正如这个犹太人所断言关于耶稣的，即\Quote{连他自己也不否认这些行为毫无神圣性，而是恶人的行为；}同样，任何不信梅瑟著作的人，引用已经说过的话，也会说同样的事，即\Quote{连梅瑟也不否认这些行为毫无神圣性，而是恶人的行为。}他也会同样对待这一点：\Quote{为真理的力量所迫，梅瑟同时既暴露了别人的行为，也定了自己的罪。}当这个犹太人说：\Quote{从相同的行为中，得出一个是天主，其余是巫士的结论，这难道不是可悲的推论吗？}有人可能会根据已经引用的梅瑟的话反驳他：\Quote{那么，从相同的行为中，得出一个是先知和天主的仆人，其余是巫士的结论，这难道不是可悲的推论吗？}但是，除了我已经提到的这些比较之外，克耳苏斯又详细论说，提出这一点：\Quote{为什么因这些行为，别人应被视为恶人，而非这人，既然他们以他为反对自己的见证？}——我们也要在已经说过的话之外提出以下一点：\Quote{为什么从那几段梅瑟禁止我们相信那些显示神迹奇事之人的经文中，我们应当视这些人而非梅瑟为恶人，因为他诬蔑其中的一些人，关于他们的神迹奇事？}他还进一步论证同样的观点，以显示加强他的意图，他说：\Quote{他自己承认这些不是神圣本性的工作，而是某些欺骗者和极恶之人的发明。}那么，谁是\Quote{他自己？}你，犹太人，说那是耶稣；但指控你同样应受这些指控的人，会将这个\Quote{他自己}转移到梅瑟身上。\par

% structure: p0114 source=h2 target=subsection reason=relative-heading-rank; base=h2

\subsection{第五十四章}

此后，诚然，克耳苏斯笔下的犹太人为保持自始至终赋予犹太人的角色，在向那些成为信徒的同胞讲话时说：\Quote{那么，是什么促使你们（成为他的追随者）呢？是因为他预言死后会复活吗？}如今，这个问题像其他问题一样，可以反过来针对梅瑟。因为我们可以对犹太人说：\Quote{那么，是什么促使了\Emph{你}（成为梅瑟的追随者）呢？是因为他关于自己之死做了如下记载：\BibleQuote{上主的仆人梅瑟死在摩阿布地，正如上主所说的。人将他葬在摩阿布，靠近培敖尔庙；直到今日没有人知道他的坟墓。}\BibleRef{申 34:5-6}}因为正如这个犹太人不相信\Quote{耶稣预言死后复活}的说法，另一个人可能会对梅瑟提出类似的断言，并回答说，梅瑟也曾记载（因为\WorkTitle{申命纪}是他的著作）说\Quote{直到今日没有人知道他的坟墓}，以便夸大并提升他的埋葬地的重要性，因为它不为人知。\par

% structure: p0116 source=h2 target=subsection reason=relative-heading-rank; base=h2

\subsection{第五十五章}

犹太人继续向那些同族中归信的人讲话，如下：\Quote{来吧，我们姑且承认那预言确实被说出了。可是，还有多少人行这样的骗术，以欺骗他们天真的听众，并借欺骗获利呢？——据说，如\Place{斯基提亚}的\Person{匝摩克息斯}，\Person{毕达哥拉斯}的奴隶；\Place{意大利}的\Person{毕达哥拉斯}本人；\Place{埃及}的\Person{兰普西尼托斯}（后者据说在\Place{阴府}与\Person{德默特耳}掷骰子，并带着从她那里得到的一块金色手帕回到了阳间）；还有\Place{俄德律萨人}中的\Person{俄耳甫斯}，\Place{忒萨利亚}的\Person{普洛忒西拉俄斯}，\Place{泰那罗斯海角}的\Person{赫拉克勒斯}，以及\Person{忒修斯}。但问题是，是否有真正死了的人，真的带着一个真实的肉身复活了？或者，你们以为别人的叙述不仅是神话，而且看起来就像神话，而你们却从十字架上的声音、他断气时的地震和黑暗，为你们的戏剧找到了一个体面且可信的结局？他活着时不能帮助自己，死后却复活了，并显示出他受刑的痕迹和他的双手如何被钉子刺穿：谁看见了这些？一个半疯癫的妇人，如你们所说，或许还有另一个也许参与了同一欺骗体系的人，他要么因特殊的心态而做了这样的梦，要么在游荡的想象影响下，按自己的愿望形成了一种幻象，无数人都遇到过这种情况；或者，最可能的是，一个人想要用这个奇事来打动别人，并借这样的谎言为像他自己一样的骗子提供机会。}\par

现在，既然是一个犹太人在说这些话，我们将为我们的耶稣辩护，仿佛我们在回答一个犹太人，仍然继续引用关于\Person{梅瑟}的记载进行比较，并对他说：\Quote{还有多少人行类似梅瑟的骗术，以欺骗他们愚蠢的听众，并借欺骗获利呢？} 这个反对意见更适合出自一个不信梅瑟之口（我们可以引用\Person{匝摩克息斯}和\Person{毕达哥拉斯}的事例，他们曾行这样的骗术），而不是出自一个犹太人之口，因为犹太人并不十分通晓希腊人的历史。此外，一个不信梅瑟奇迹的埃及人，可以可信地引用\Person{兰普西尼托斯}的事例，说兰普西尼托斯下到\Place{阴府}并与\Person{德默特耳}掷骰子，然后偷了她的金色手帕，并以此作为他曾到过阴府并从那里返回的证据，这比梅瑟记载他进入有天主所在的黑暗，并且唯独他超越众人亲近天主，更可信。因为梅瑟这样说：\BibleQuote{惟独梅瑟可以走近天主；其余的人不可走近。} 然后，我们这些耶稣的门徒，对提出这些反对意见的犹太人说：\Quote{你们一面攻击我们对耶稣的信仰，一面为自己辩护，并回答埃及和希腊的反对者：对于你们针对我们的耶稣所提出的那些指控，你们将说什么？这些指控同样可以首先针对梅瑟提出。如果你们竭力为梅瑟辩护——他的历史确实容许清晰有力的辩护——那么你们在支持梅瑟时，将不自觉地成为帮助确立耶稣更大神性的不情愿的助手。}\par

% structure: p0119 source=h2 target=subsection reason=relative-heading-rank; base=h2

\subsection{第五十六章}

既然犹太人说这些所谓英雄下到\Place{阴府}并从那里返回的传说都是骗人的把戏，声称这些英雄消失了一段时间，秘密地从众人眼前隐退，后来又宣称自己从阴府回来了——因为他的话似乎对\Place{俄德律萨}的\Person{俄耳甫斯}、\Place{忒萨利亚}的\Person{普洛忒西拉俄斯}、\Place{泰那罗斯}的\Person{赫拉克勒斯}以及\Person{忒修斯}有这样的含义——让我们努力证明，耶稣从死者中复活的记载绝不能与这些相比。因为所提到的每一位英雄，如果他们愿意，都可以秘密地从人眼前隐退，如果他们决定的话，再回到他们离开的人那里；但耶稣是在所有犹太人面前被钉十字架，他的尸体在他的民族面前被杀，他们怎么能说他对那些传说中下到阴府并返回的英雄施行了类似的欺骗呢？但我们认为，也许可以提出以下考虑作为耶稣公开被钉十字架的辩护，特别是联系那些被认为被迫下到阴府的英雄的故事：如果我们假设耶稣是以一种不为人知的死亡而死，以至于他死亡的事实并没有对整个犹太民族公开，然后他确实从死者中复活了，那么在这种情况下，对于他本人也会产生与那些英雄相同的怀疑。因此，除了耶稣被钉十字架的其他原因之外，或许这也促成了他死在十字架上这个显眼的死亡，以便没有人可以说他自愿从人眼前隐退，只是看似死了，实际上并没有；并且当他再次出现时，把从死者中复活变成了一场骗人的把戏。但我认为，一个清晰无误的证据是他的门徒们的事业，他们献身于教导一种对人的生命有危险的教义——如果他们是编造了耶稣从死者中复活，他们就不会如此勇敢地教导这种教义；同时，他们不仅使别人轻视死亡，而且他们自己首先表现出对死亡恐惧的蔑视。\par

% structure: p0121 source=h2 target=subsection reason=relative-heading-rank; base=h2

\subsection{第五十七章}

但请观察，这位刻尔苏斯的犹太人是否非常盲目地说，他说任何人都不可能用真正的身体从死人中复活，他的原话是：\Quote{但问题是，是否真有死人曾用真正的身体复活过？}然而，一个相信\BibleBook{列王纪}{上}和\BibleBook{列王纪}{下}中关于两个小孩记载的犹太人，是决不会说这些话的；其中一个孩子被\Person{厄里亚}复活，另一个被\Person{厄里叟}复活。因此，我想\Person{耶稣}之所以只向犹太人显现，是因为他们已经习惯了奇迹；这样，通过将他们自己所相信的与他所做的以及有关他的记述相比较，他们可以承认那一位，行了更伟大的事，行了更奇妙的奇迹，比所有在他之前的人都更伟大。\par

% structure: p0123 source=h2 target=subsection reason=relative-heading-rank; base=h2

\subsection{第五十八章}

此外，在刻尔苏斯的犹太人引用了那些关于行骗者假装复活的希腊故事之后，他对那些皈依基督教的犹太人说道：\Quote{你们是否认为别人的说法不仅是神话，而且看起来也像神话，而你们却在十字架上祂断气时的声音中，为你们的戏剧找到了一个合宜而可信的结局？}我们回答那位犹太人：\Quote{你视为神话的，我们也视为神话；但那些我们共同的圣经的记述，不仅是你，也是我们所引以为傲的，我们绝不视为神话。因此，我们相信其中记载的一些人从死人中复活，并不认为它们是欺人之谈；我们尤其相信其中所记载的那位复活者，祂自己预言了这件事，也是别人预言的对象。而祂的复活在这方面比其他人更奇妙：他们是被先知\Person{厄里亚}和\Person{厄里叟}复活的，而祂却不由任何先知，而是由祂在天上的父所复活。因此，祂的复活也产生了比他们更大的结果。因为通过\Person{厄里亚}和\Person{厄里叟}复活孩童给世界带来了什么巨大的好处，能比得上宣扬\Person{耶稣}复活——作为信条接受并借由神力实现——所产生的好处呢？}\par

% structure: p0125 source=h2 target=subsection reason=relative-heading-rank; base=h2

\subsection{第五十九章}

\Person{刻尔苏斯}还臆想地震和黑暗都是捏造的；但关于这些，我们在前文已尽力进行了辩护，引用了\Person{弗勒贡}的见证，他记述这些事件发生在我们救主受难之时。他接着说：\Quote{耶稣活着的时候，对自己毫无帮助，但他死后复活，并展示了他受罚的痕迹，显示了他的双手如何被钉子刺穿。}我们问他，\Quote{对自己毫无帮助？}这句话是什么意思？因为如果他所指是缺乏美德，我们回答祂\Emph{确实}提供了极大的帮助。因为祂没有说过或做过任何不当之事，而是真的\BibleQuote{祂被牵去待宰，像只绵羊；在剪毛人面前不出声，像只羔羊}；福音证实祂没有开口。但如果刻尔苏斯把这表述应用于无关紧要的和肉体的事，（意思是耶稣在那些事上无法帮助自己），我们说我们已从福音书中证明，祂是自愿去面对受苦的。接下来，他谈到福音书中关于祂复活后展示受罚痕迹和双手被刺穿的记载，他问：\Quote{谁看到了这个？}而且他对有关\Person{玛利亚玛达肋纳}（据记载她看见了祂）的叙述表示怀疑，他回答：\Quote{一个半疯的女人，如你们所说。}而且因为她不是唯一被记载看见复活救主的人，还有其他人的记载，这位刻尔苏斯的犹太人也诽谤这些叙述，补充说：\Quote{还有另一个从事相同欺骗体系的人！}\par

% structure: p0127 source=h2 target=subsection reason=relative-heading-rank; base=h2

\subsection{第六十章}

其次，仿佛这是可能的，即一个死人的形象可以像活人一样出现在另一个人面前，他以伊壁鸠鲁派的观点采纳了这一说法，并说：\Quote{有些人因特殊的精神状态而做梦，或因扭曲的想象力而形成了他自己所希望看到的形象，便报告说看到了这样的人；而这，}他继续道，\Quote{在无数人身上都发生过。}但即使他的这个说法似乎有很大说服力，它也只不过适合确认一个必要的道理，即死者的灵魂以分离的状态存在（脱离身体）；而采纳这种意见的人，并非毫无理由地相信灵魂的不朽，或至少是持续存在，正如\Person{柏拉图}在\WorkTitle{论灵魂}一书中也说，已死之人的阴影般的幻影曾出现在一些人身边，围绕他们的坟墓。这些关于死者灵魂的幻影是由某种实体产生的，而这个实体存在于灵魂中，这灵魂以分离的形式存在于一个被称为光辉身体的形体中。但刻尔苏斯不愿承认任何这类观点，他坚持认为有些人是做了白日梦，并在扭曲的想象力影响下，为自己形成了他们所希望的形象。相信人在睡眠中做梦并非不合理；但假设那些并非完全丧失理智、也没有处于谵妄或疑病症影响之下的人有清醒的异象，则是不可信的。刻尔苏斯看到这一点，便称那女人\Quote{半疯}——这一说法并非来自记录此事的史书，而是他借此指控这些事件为虚假。\par

% structure: p0129 source=h2 target=subsection reason=relative-heading-rank; base=h2

\subsection{第六十一章}

因此，按照克尔苏斯的想象，耶稣在死后只显示了祂在十字架上所受伤痕的外观，而实际上并未如所描述的那样受伤；然而，根据福音的教导——克尔苏斯任意接受其中某些部分以寻找控告的理由，而拒绝其他部分——耶稣召唤了祂的一位陷入怀疑的门徒，此人认为这个奇迹是不可能的。那个人确实相信了那妇女所说她看见了祂的陈述，因为他并不认为看见死者的灵魂是不可能的；但他尚未相信祂以肉身复活，这肉身是先前身体的\OriginalTerm{对范}{antitype}。因此他不仅说：\Quote{除非我看见，否则我不信}；而且补充道：\Quote{除非我把我的手探入钉痕，并把我的手放在祂的肋旁，否则我不信。}这些话是托马斯说的，他认为灵魂的身体可能被感官之眼看见，并且完全相似于其先前的形貌，\par

\Quote{无论大小、眼睛的美丽，\newline{}和声音方面，\newline{}}\par

\Quote{并且常常，}\par

\Quote{身上也穿着（生前）那样的衣裳。}\par

因此，耶稣召来托马斯，说：\Quote{伸出你的指头，看我的手；伸出你的手，探入我的肋旁；不要无信，而要信。}\BibleRef{若 20:27}\par

% structure: p0135 source=h2 target=subsection reason=relative-heading-rank; base=h2

\subsection{第六十二章}

那么，根据所有关于祂的预言——其中也包括这复活预言——以及根据祂所行的一切和所遭遇的一切，这个事件应比其他一切更为奇妙。因为先知曾以耶稣的口吻预言说：\BibleQuote{我的肉身要安息于希望中，祢不会将我的灵魂遗弃在\OriginalTerm{阴府}{Hades}，也不让祢的圣者见到腐朽。}\BibleRef{咏 16:9-10} 实在，祂复活后，存在于一个介于祂受难前身体的粗重性与不被如此身体包裹的灵魂显现象之间。因此，当祂的门徒们聚在一处，托马斯也在其中时，\BibleQuote{耶稣来了，门都关着，站在他们中间说：愿你们平安。然后祂对托马斯说：伸出你的指头，等等。}\BibleRef{若 20:26-27} 同样在路加福音中，当西满和克罗帕彼此谈论所发生的一切时，\BibleQuote{走近他们，与他们同行。他们的眼睛却被蒙蔽，以致认不出祂来。祂对他们说：你们走路时彼此谈论的是什么事？}\BibleRef{路 24:15-18} 当他们的眼睛开了，认出祂来，那时圣经明确地说：\BibleQuote{祂就从他们眼前隐没了。}\BibleRef{路 24:31} 尽管克尔苏斯可能想把关于耶稣以及祂复活后看见祂的人的事，与某种想象的不同显现和那些编造这些事的人等量齐观，但对于那些进行公正而明智考察的人，这些事件只会显得更加奇妙。\par

% structure: p0137 source=h2 target=subsection reason=relative-heading-rank; base=h2

\subsection{第六十三章}

在这些论点之后，克尔苏斯接着提出一个对福音叙述的控诉，不可轻率放过。他说：\Quote{如果耶稣想要显示祂的能力确实是神圣的，祂就应该向那些虐待过祂的人、向定祂罪的人并向所有人显现。}因为按照福音的记载，我们认为这也是真实的：祂复活后并没有像从前那样公开向所有人显现。但《宗徒大事录》记载：\BibleQuote{祂四十天之久显现，向门徒们讲论天主之国的事。}\BibleRef{宗 1:3} 而在福音书中，并没有说祂一直与他们在一起；只说有一次，八天后，当门都关着时，祂出现在他们中间，另一次则以类似的方式。保禄在《格林多前书》的结尾部分，论到祂没有像受难前那样公开显现时，如此写道：\BibleQuote{我当日所领受又传给你们的，最重要的就是：基督照圣经所说，为我们的罪死了，而且埋葬了，又照圣经所说，第三天复活了，并且显给刻法看，然后显给那十二位看；此后一次显给五百多位弟兄看，其中多半到现在还在，但也有已经安眠的；以后显给雅各伯看，再后显给众宗徒看；最后也显给我这个像流产儿的人看。}\BibleRef{格前 15:3-8} 我认为，这段经文包含某些伟大而奇妙的奥秘，不仅超出大多数普通信友的领悟，甚至超出那些在基督徒知识上大有长进之人的领悟，其中解释了为何祂从死人中复活后，没有像复活前那样显现。在这部为回应反对基督徒及其信仰的著作而写的论著中，请注意我们能否从更多的理性论证中提出少数几个，并对这篇辩护的听众产生影响。\par

% structure: p0139 source=h2 target=subsection reason=relative-heading-rank; base=h2

\subsection{第六十四章}

虽然\Person{耶稣}只是单独一个人，但祂却依人们审视祂的不同角度，不止是一件事；也不是所有看见祂的人都以同样的方式看见祂。现在，依不同观点，祂不止是一件事，这从以下陈述清楚可见：\BibleQuote{我是道路、真理、生命}；\BibleQuote{我是食粮}；\BibleQuote{我是门}，以及其他无数的话。并且祂被看见时并不以同样的形貌出现在所有看见祂的人面前，而是依照每个人接受祂的能力，这一点对那些注意以下事的人将是清楚的：当祂即将在高山上改变容貌时，祂并没有让祂的所有宗徒都看见这景象，而只有\Person{伯多禄}、\Person{雅各伯}和\Person{若望}，因为只有他们在那时能够看见祂的荣耀，并观察\Person{梅瑟}和\Person{厄里亚}的荣耀形貌，聆听他们的谈话和来自天上云彩的声音。我也认为，在祂登上那座只有祂的门徒来到祂跟前、祂教导他们真福八端的山之前——当祂在山下某处、天色已晚时，祂治愈那些被带到祂面前的人，使他们摆脱一切疾病和痛苦——祂对那些患病和需要祂医治帮助的人，与对那些因自身力量而能随祂上山的人，所显现的不同。不但如此，即使当祂私下给自己的门徒解释那些对外面群众所说的比喻（他们没有得到解释）时，那些听到解释的人被赋予较之那些听到却未得解释者更高的听觉器官，同样也完全如此，关于他们灵魂的眼睛，并且我认为，也关于他们身体的眼睛。以下陈述表明祂并非总是一样的形貌：即当\Person{犹达斯}即将出卖祂时，对与他同行的群众说话（因为他们不认识祂）：\BibleQuote{我亲吻谁，谁就是祂。} 我认为救主自己也用以下话语指示了同一件事：\BibleQuote{我天天在圣殿里教导你们，你们却没有拿住我。} 因此，怀着如此崇高关于\Person{耶稣}的观点——不仅关于内在的、对群众隐藏的神性，而且关于祂身体的改变容貌（这发生在祂愿意的时间和对象）——我们说，在\Person{耶稣}\BibleQuote{脱去宰制者和掌权者}之前，且当祂尚未死于罪恶时，所有人都能看见祂；但当祂\BibleQuote{脱去宰制者和掌权者}之后，且不再有任何东西能被群众看见时，所有先前见过祂的人现在不能看见祂了。因此，为了顾惜他们，祂在自己从死者中复活后并没有向所有人显示自己。\par

% structure: p0141 source=h2 target=subsection reason=relative-heading-rank; base=h2

\subsection{第六十五章}

并且我为什么说\Quote{向所有人}？因为即使对祂自己的宗徒和门徒，祂也不是一直临在，也不是不断向他们显示自己，因为他们不能不间断地接受祂的神性。因为祂的神性在祂完成救恩计划之后更加辉煌：而这位\Person{伯多禄}，别号\Person{刻法}，似乎是宗徒的初果，能够看见祂，和他一起的还有那十二人（\Person{玛弟亚}代替了\Person{犹达斯}的位置）；之后祂同时显现给那五百弟兄，然后显现给\Person{雅各伯}，随后又显现给其余所有宗徒之外的其他人，或许还有那七十人，最后显现给\Person{保禄}，如同给一个流产儿，他很知道怎样说：\BibleQuote{我比众圣徒中最小的还小，但恩宠赐给了我；} 可能\BibleQuote{最小的}这个表达与\BibleQuote{流产儿}有相同的意思。因为正如没有人能合理地责备\Person{耶稣}没有让祂所有的宗徒上高山，而只有那三位已经提到的，在祂改变容貌时，当祂即将显示祂衣服上出现的荣耀、以及\Person{梅瑟}和\Person{厄里亚}同祂交谈的荣耀时；同样，也没有人能合理地反对宗徒们的陈述，他们介绍\Person{耶稣}复活后的显现不是向所有人，而只向那些祂知道已经领受了能看见祂复活的眼目的人。此外，我认为以下关于祂的陈述对于我们的话题具有护教价值，即：\BibleQuote{基督为此死了，又复活了，为要做死人和活人的主。} 因为请注意，这些话传达的是：\Person{耶稣}死了，为要做死人的主；祂复活了，为要做不仅是死人，也是活人的主。无疑，宗徒所理解的基督要为之做主的死人，是在\WorkTitle{格林多前书}中所说的那些人：\BibleQuote{因为号筒要响，死人要复活成为不朽坏的；} 至于活人，则是那些将要被改变的人，他们不同于将要复活的死人。关于活人的话是这些：\BibleQuote{我们将要被改变；} 这个表达紧随于\BibleQuote{死人先要复活}的陈述之后。此外，在\WorkTitle{得撒洛尼前书}中，他用不同的词语描述同一改变，说那些睡了的人与活着的人不同；他的话是：\BibleQuote{弟兄们，论到睡了的人，我不愿意你们不知道，免得你们忧伤，像那些没有望德的人一样。如果我们信\Person{耶稣}死了，又复活了，同样，那些在\Person{耶稣}内睡了的人，天主也要把他们与\Person{耶稣}一同带来。我们照主的话告诉你们这件事：我们这些活着还存留到主来的人，决不会在那些睡了的人之先。} 我们以为适合这段经文的解释，已经在\WorkTitle{得撒洛尼前书}的释经注释中给出了。\par

% structure: p0143 source=h2 target=subsection reason=relative-heading-rank; base=h2

\subsection{第六十六章}

若你们因所有相信耶稣的群众没有看见祂的复活而感到惊讶，那么请记住，保禄在写给格林多人的书信中，对他们说，因为他们不能领受更深的道理：\BibleQuote{我曾决定，在你们中不知道别的，只知道耶稣基督，这被钉在十字架上的耶稣基督}；这等于说：\BibleQuote{你们以前不能，如今也不能，因为你们仍是属血肉的。}因此，圣经按照天主的安排，记载了耶稣在受难前向所有人无差别地显现，但并非总是如此；而在受难后，祂不再以相同的方式向所有人显现，而是有分别地按各人所应得的量度。正如经上记载\BibleQuote{天主显现给亚巴郎}或某位圣人，这种\Quote{显现}并非经常发生，而是间或发生，且并非向所有人；同样，你们要理解，天主子的显现原则，与天主向后者的显现相同。\par

% structure: p0145 source=h2 target=subsection reason=relative-heading-rank; base=h2

\subsection{第六十七章}

因此，我们尽己所能，在这样的论著中回应了这样的异议：\Quote{如果耶稣真的愿意显示祂的神性能力，祂应当向那些虐待祂的人、向定祂罪的法官，并向所有人毫无保留地显现自己。}然而，祂并无义务向定祂罪的法官或虐待祂的人显现。因为耶稣宽恕了这双方，免得他们像索多玛人那样被击打失明，正如罗特所接待的天使的美貌遭人谋害时发生的事。关于此事的记载是：\BibleQuote{那两个人伸出手来，将罗特拉进屋里，关上门；又使门外的人，无论大小，都眼目昏眩，以致他们找不着门口。}因此，耶稣愿意向每一个有能力看见祂的人，按照其能力的大小，显示祂的神性能力。我不认为祂避免被看见有其他原因，而只是顾及那些无法看见祂的人是否适合。刻尔苏斯补充说：\Quote{祂死后不再需要惧怕任何人，正如你们所说，祂是天主；祂来到世间也根本不是为了隐藏自己。}这是徒劳的。然而，祂来到世间不仅是为了被认识，也是为了隐藏。因为祂的一切，甚至那些认识祂的人也未能全部知晓，他们还有一部分未被揭示；而有些人则完全不了解祂。祂为那些本是黑暗和黑夜之子、并立志成为光明和白昼之子的人敞开了光明之门。因为我们的救主，如同一位良医，更是为了我们这些充满罪恶的人而来，而非为了那些义人。\par

% structure: p0147 source=h2 target=subsection reason=relative-heading-rank; base=h2

\subsection{第六十八章}

然而，让我们看看刻尔苏斯笔下的这个犹太人是如何宣称的：\Quote{如果这至少有助于彰显祂的神性，那么祂就应当立即从十字架上消失。}这在我看来，就像那些反对天主上智安排理论的人所提出的论点一样，他们以不同于现实的方式安排事物，并声称如果按照他们的安排，世界会更好。在那些他们能够安排的情况下，人们证明他们所作的安排实际上使世界变得更糟；而在那些他们没有将事物描绘得比实际更糟的情况下，他们被证明是在追求不可能的事；因此，在任何一种情况下，他们都应受嘲笑。那么，祂作为一位更神圣的存在，能够随心所欲地消失，这是否可能，从事情的本质来看是清楚的；而且，根据那些不单单挑选某些叙事片段以攻击基督教、并将其余视作虚构之人的判断，这一点从关于祂的记载中也可以确定。因为\BibleBook{路加}{福音}记载，耶稣在复活后拿起饼来，祝福了，擘开，递给西满和克罗帕；当他们接过饼时，\BibleQuote{他们的眼睛开了，这才认出耶稣来；但祂却从他们眼前隐没了。}\par

% structure: p0149 source=h2 target=subsection reason=relative-heading-rank; base=h2

\subsection{第六十九章}

但我们愿指出，祂瞬间从十字架上消失（比祂留在其上）并不更有利于整个救恩计划的实现。因为关于耶稣所发生之事的单纯字面和叙述，并未呈现真理的全貌。因为每一件事，对于那对圣经有明智理解的人，都可表明为另一事物的象征。因此，正如祂的钉十字架包含着一个真理，体现在\BibleQuote{我已与基督同钉十字架}这句话中，也隐含在\BibleQuote{我绝不夸耀，唯夸我们主耶稣基督的十字架，藉着祂，世界已被钉于十字架，而我也被钉于世界}这句话里；正如祂的死是必要的，因为\BibleQuote{祂既死了，就一次而永远地为罪死了}，以及\BibleQuote{效法祂的死}，以及\BibleQuote{如果我们与祂同死，也必与祂同活}；同样，祂的埋葬也适用于那些已效法祂的死、已与祂同钉十字架并与祂同死的人，正如保禄所宣告的：\BibleQuote{我们藉着圣洗与祂一同埋葬，也与祂一同复活。}然而，这些关乎祂埋葬、坟墓和埋葬祂的人的事，我们将在更适当的时机，当正式以这类题目为目的时，更详细地阐述。但目前，注意到包裹耶稣纯洁身体的洁白细麻布，以及若瑟从磐石中凿出的新坟墓——那里\BibleQuote{还没有人躺过}，或如若望所说\BibleQuote{里面从来没有放过人}——就足够了。且看三位福音作者的和谐是否不令人印象深刻：他们以为应当将坟墓描述为\Quote{从磐石中凿出或开凿}的；这样，考察叙述话语的人可从中看到值得深思之处，既在于坟墓的\Emph{新}（玛窦和若望提到这一点），也在于路加和若望的陈述：里面从未葬过人。因为祂不像其他死人（即使在死亡中，也在水和血中显出生命的迹象），可以说是一个\Emph{新}的死人，理当被安放在一个崭新洁净的坟墓中，以便正如祂的诞生比任何其他诞生更纯洁（因为祂不是通过普通生育而生，而是由童贞女所生），祂的埋葬也可能具有象征性的纯洁，体现于祂的身体被安置在一个新坟墓中，这坟墓不是用各处收集的石头建造、没有自然合一，而是从\Emph{一块}磐石中凿出、各部分连在一起的。然而，对这些问题的解释，以及从叙述本身上升到它们所象征之事的方法，可在更适当的时机，在专门论述此类主题的著作中，以更深刻、更适合它们神圣特质的方式进行探讨。至于字面叙述，则可这样解释：即对于决意忍受悬在十字架上的祂而言，持守祂所取性格的一切伴随物是合适的，以便祂作为人被杀、作为人而死、也作为人被埋葬。但即使根据策尔苏斯的观点，福音书中记载耶稣立即从十字架上消失，他和别的非信徒也会挑剔这叙述，并会提出类似这样的反对：\InnerQuote{究竟为什么祂在被钉十字架后消失，而不是在受苦之前消失？}如果他们在从福音书得知祂并没有立即从十字架上消失之后，还认为可以挑剔这叙述，认为它没有像他们认为应当的那样编造任何这样立即的消失，而是如实记述，那么，他们不也应该相信祂的复活，并相信祂按自己的意愿，一次在门关着时站在门徒中间，另一次在分饼给两位熟识者后，对他们说了某些话之后立即消失不见吗？\par

% structure: p0151 source=h2 target=subsection reason=relative-heading-rank; base=h2

\subsection{第七十章}

但是，策尔苏斯笔下的这位犹太人怎么能说耶稣隐藏了自己？因为关于祂，他的话是这样的：\Quote{被派为使者的人，在应当传达他的信息时，有谁会隐藏自己？} 但祂并没有隐藏自己，祂对那些想要捉拿祂的人说：\BibleQuote{我天天在圣殿里公开教训人，你们并没有拿住我。} 然而，我们已经有一次答复过策尔苏斯的这项指控，如今他又重复，我们就以先前所说的为满足。在前面的篇幅中，我们也答复了这项异议：\Quote{当祂在肉身中时，没有人相信祂，祂却不停地向众人宣讲；但当祂从死者中复活后，本可以产生强烈的信仰，却只秘密地向一个女人和自己亲密的朋友显现。} 说祂只向一个女人显现，这并不真实；因为根据《玛窦福音》记载：\BibleQuote{在安息日过后，一周的第一天天快亮时，玛利亚玛达肋纳和另一个玛利亚来看坟墓。看，发生了大地震，因为主的天使从天降下，前来把石头滚开。} 紧接着，玛窦又补充说：\BibleQuote{看，耶稣迎接\Emph{她们}——显然指上述的两位玛利亚——说：\InnerQuote{愿你们平安！}她们就上前抱住祂的脚，朝拜了祂。} 我们也答复了这项指控：\Quote{当祂受惩罚时，众人看见了祂；但复活后，只被一个人看见。}那时我们为这个事实辩护：\Quote{祂并没有被所有人看见。}现在我们可以说，祂那纯粹人性的一面是所有人可见的，但那些属于神性本质的一面——我所说的不是相关属性，而是区别属性——并非所有人都能接受。但请看这里策尔苏斯陷入的明显矛盾。因为他稍早前说耶稣曾秘密地向一个女人和自己亲密的朋友显现，接着又立即补充说：\Quote{当祂受惩罚时，所有的人都看见了祂；但复活后，只有一个人看见了祂，而情况本应相反。}让我们听听他说的\Quote{本应相反}是什么意思。在受惩罚时被所有人看见，但在复活后只被一个人看见，这二者是相反的。那么，就他的话所传达的意思而言，他希望发生既不可能又荒谬的事，即：当祂受惩罚时，只被一个人看见；但复活后，却被所有人看见！不然，你将如何解释他的话：\Quote{情况本应相反？}\newline{}\par

% structure: p0153 source=h2 target=subsection reason=relative-heading-rank; base=h2

\subsection{第七十一章}

耶稣教导我们是谁派遣了祂，祂说：\BibleQuote{除了圣子，没有人知道圣父；}又说：\BibleQuote{从来没有人看见天主；只有在圣父怀里的独生子，祂将天主表明出来。}祂谈论神性，向祂的真正门徒传授关于天主的道理。我们在圣经叙述中发现了这类教导的痕迹，并借此帮助我们形成神学概念。一处经文说：\BibleQuote{天主是光，在祂内毫无黑暗；}另一处说：\BibleQuote{天主是神，朝拜祂的人应当以心神和真理朝拜祂。}但圣父派遣祂的目的不可胜数；凡愿意的人，可以从预言祂的先知们，以及从福音书的叙述中得知。从宗徒们，尤其是保禄那里，也会学到不少。此外，祂引导虔诚的人走向光明，而犯罪的人祂将惩罚——这个事实策尔苏斯没有注意到，因而把祂描述为\Quote{一位将要引导虔诚的人走向光明，并且怜悯其他的人，无论他们是犯罪还是悔改。}\newline{}\par

% structure: p0155 source=h2 target=subsection reason=relative-heading-rank; base=h2

\subsection{第七十二章}

在以上陈述之后，他继续说：\Quote{如果祂想隐藏自己，为什么有声音从天上宣告祂是天主子？如果祂不寻求隐藏，为什么祂受惩罚？或者为什么祂死？}他用这些问题企图证明记载之间存在矛盾，却没有注意到耶稣既不愿意有关祂的一切都被所有遇到祂的人知道，也不愿意一切都不被知道。因此，从天上宣告祂是天主子的声音，说：\BibleQuote{这是我的爱子，我所喜悦的，}经上并没有说这声音被群众听到，正如策尔苏斯笔下的这位犹太人所设想的那样。此外，在高山上从云中发出的声音，也只有那些与祂一同上山的人听到。因为天主的声音之性质，只有说话者愿意听见的人才能听见。我认为，所提到的天主的声音，既不是被击打的空气，也不是空气的任何震动，也不是在论声音的论著中提到的任何其他东西；因此，它是通过比感官更好的、更神圣的听觉器官被听见的。当说话者不愿祂的声音被所有人听见时，那拥有更精细耳朵的人能听见天主的声音，而灵魂的耳朵沉睡的人却察觉不到是天主在说话。我说这些是因为他问：\Quote{为什么有声音从天上宣告祂是天主子？}至于这个疑问：\Quote{如果祂想隐藏，为什么祂受惩罚？}在前面的篇幅中关于祂受苦所作的更详尽的论述应当足够了。\newline{}\par

% structure: p0157 source=h2 target=subsection reason=relative-heading-rank; base=h2

\subsection{第七十三章}

此后，犹太人接着从前提推导出一个并非必然的结论；因为从\Quote{祂愿意藉所受的刑罚，也教导我们轻视死亡}并不能推出：祂复活后应公开召唤众人到光明中，并教导他们祂来临的目的。因为祂早已用这些话召唤众人到光明中：\BibleQuote{凡劳苦和负重担的，你们都到我这里来，我要使你们安息。}而祂来临的目的，已在关于真福的讲论、随后的宣告、比喻以及与经师和法利塞人的对话中详尽说明了。\BibleBook{若望}{福音}给予我们的教导表明，耶稣的雄辩不在于言语，而在于行动；福音叙述也清楚显示，祂的宣讲是\Quote{具有权威的}，正因如此，众人都惊讶祂。\par

% structure: p0159 source=h2 target=subsection reason=relative-heading-rank; base=h2

\subsection{第七十四章}

此外，犹太人还进一步说：\Quote{这些话都是从你们自己的书中摘取的，此外我们不需要别的证人；因为你们是自相残杀。}\par

然而，我们已经证明，在犹太人的陈述中，无论是关于耶稣还是关于我们，都有许多愚蠢的主张与我们的福音叙述相悖。而且，我认为他并未证明\Quote{我们自相残杀}，只是他这样想象而已。当犹太人在他先前的话之后一般性地补充说：\Quote{至高天上的神啊！哪一位神在显现于人时竟被怀疑？}我们必须对他说：根据梅瑟法律记载，天主曾以极其公开的方式访问希伯来人，不仅在埃及所行的征兆和奇迹、红海的穿越、火柱和光云中，而且在十诫向全体民众颁布时，但那些看见这些事的人仍然不信；因为他们若相信所见所闻，就不会铸造牛犊，不会将自己的荣耀变为吃草的牛犊的像，也不会彼此指着牛犊说：\BibleQuote{以色列，这就是领你出埃及地的神。}请观察，同一民族过去在整个旷野漂流时期拒绝相信如此奇迹和神圣显现（正如他们在犹太法律中所记载的），如今在耶稣荣耀来临之际，他们是否也因祂带着权威所说的大能话语和在全民众前所行的奇事而不愿信服？\par

% structure: p0162 source=h2 target=subsection reason=relative-heading-rank; base=h2

\subsection{第七十五章}

我认为上述足以使任何人相信，犹太人不信耶稣与他们自古以来对天主的不信是一致的。因为我要答复此克耳苏笔下的犹太人，当他问：\Quote{哪一位神显现于人中竟被怀疑，而且是在那些期盼祂的人中显现？又请问，为何那些长久等待祂的人却不认识祂？}朋友们，你们希望我们如何回答你们的问题？照你们看，哪一类奇迹更大？是在埃及和旷野所行的，还是我们宣称耶稣在你们中间所行的？如果你们认为前者更大，那么由此不正表明，轻视大奇迹的人也会鄙视小奇迹吗？这是对我们的耶稣奇迹叙述的看法。而如果耶稣所行的奇迹被认为与梅瑟所记的同样伟大，那么在一个对两种盟约开端都表现出不信的民族中，发生了何等奇异的事？因为立法的开端是在梅瑟时代，其中记载了你们中间不信者和恶人的罪；而我们立法和新盟约的开端则公认是在耶稣时代。你们因不信耶稣而证明自己是那在旷野中不信神圣显现者的子孙；这样，我们的救主所说的话也适用于你们这些不信祂的人：\BibleQuote{因此你们证明并赞同你们祖先的行径。}还有那预言也在你们中间应验了：\BibleQuote{你的生命必在你们眼前悬而未定，你必昼夜恐惧，对自己的生命没有把握。}因为你们不信那来拜访人类生命的主。\par

% structure: p0164 source=h2 target=subsection reason=relative-heading-rank; base=h2

\subsection{第七十六章}

\Person{塞尔苏斯}扮演犹太人的角色时，未能发现任何可能反对福音的异议，而这些异议若不也同样适用于反对律法和先知，便不会被他用来攻击耶稣。因为他用这样的话指责耶稣：\Quote{祂使用威吓，并因轻率的原因辱骂人，说\InnerQuote{祸哉，你们}和\InnerQuote{我预先告诉你们}。因为用这样的表达，祂显然承认自己无法说服人；而一位天主，甚至一个明智的人，都不会这样做。} 现在请看，这些指责岂不显然反过来落在犹太人身上吗？因为在律法和先知著作中，天主使用威吓和辱骂，使用与福音中一样严厉的语言，例如\Person{依撒意亚}的这些话：\BibleQuote{祸哉，那些使房屋相连，使田地相接的人；}\BibleRef{依 5:8} 和 \BibleQuote{祸哉，那些清晨早起，追求烈酒的人；}\BibleRef{依 5:11} 和 \BibleQuote{祸哉，那些用长绳牵引罪过的人；}\BibleRef{依 5:18} 和 \BibleQuote{祸哉，那些称恶为善，称善为恶的人；}\BibleRef{依 5:20} 和 \BibleQuote{祸哉，你们这些饮酒有能的人；}\BibleRef{依 5:22} 以及无数其他类似的段落。以下这些话岂不也像他所说的威吓：\BibleQuote{祸哉，犯罪的民族，恶贯满盈的民族，作恶的种子，败坏的子女？}\BibleRef{依 1:4} 等等，随后他加上与他说耶稣所用的同样严厉的威吓。因为那宣布\BibleQuote{你们的国土荒凉，你们的城市被火焚毁，你们的地在你们眼前被外邦人吞噬，荒凉如同被外邦人倾覆？}\BibleRef{依 1:7}的话，岂不是一个威吓，而且是一个大威吓吗？\BibleBook{厄则克耳}{先知书}中不是也有对人民的辱骂，当时主对先知说：\BibleQuote{你住在蝎子中间？}\BibleRef{厄 2:6} 那么，\Person{塞尔苏斯}，你认真地让犹太人这样说耶稣：\Quote{祂因轻率的原因使用威吓和辱骂，当祂说\InnerQuote{祸哉，你们}和\InnerQuote{我预先告诉你们}时？} 你难道看不出，你这位犹太人对耶稣的指控，也可用来指控天主吗？因为那在先知著作中说话的天主，显然也面临同样的指控，如\Person{塞尔苏斯}所认为的，即不能说服人。此外，我还可以对这位犹太人（他认为通过这样的说法对耶稣提出了好的指控）说，如果他根据经文记载，为\BibleBook{肋未}{纪}和\BibleBook{申命}{纪}中记载的众多诅咒辩护，那么我们也能为被认为是耶稣所说的辱骂和威吓做出同样好、甚至更好的辩护。至于\Person{梅瑟}律法本身，我们比犹太人更能辩护，因为我们从耶稣那里学到了对律法著作更明智的理解。不仅如此，如果犹太人理解了先知经文的意义，他就能表明天主使用威吓和辱骂并非出于轻率的原因，当祂说\Quote{祸哉，你们}和\Quote{我预先告诉你们}时。\Person{塞尔苏斯}认为即使是明智的人也不会诉诸这样的表达，而天主又怎能使用这样的表达来使人们悔改呢？然而，基督徒只认识一位天主——即那位在先知和主（\Person{耶稣}）里说话的天主——能够证明这些威吓和辱骂（如\Person{塞尔苏斯}所认为并称呼的）的合理性。在此，要对这位自称既是哲学家又了解我们整个体系的\Person{塞尔苏斯}说几句话。朋友，当\Person{荷马}的史诗中的\Person{赫尔墨斯}对\Person{奥德修斯}说，\par

\Quote{为什么，现在，可怜的人，你独自一人游荡在山顶上？}\par

你满足于这个解释，即\Person{荷马}史诗中的\Person{赫尔墨斯}对\Person{奥德修斯}说这样的话是为了提醒他的责任，因为塞壬的特点是谄媚和说讨人喜欢的话，围绕她们，\par

\Quote{是一堆巨大的骨头，}\par

并且她们说，\par

\Quote{到这里来，备受赞誉的\Person{奥德修斯}，希腊人的伟大荣耀；}\par

然而，如果我们的先知和\Person{耶稣}自己，为了引导听者远离邪恶，使用了诸如\Quote{祸哉，你们}以及你所认为的辱骂这样的表达，那么这种语言难道不是迁就了听者的处境，并将这些话作为医治的良药应用在他们身上吗？除非，你确实认为天主，或分享神性的一位，在与人交谈时，只顾及自己的本性和与自己相称的事，而不顾及适合向那些处于祂圣言的安排和引导之下的人提出什么，也不顾及祂应当根据每个人的个别特质与之交谈。那么，说\Person{耶稣}无法说服人，岂不是一个荒谬的断言吗？当你比较情况：不仅在犹太人中，他们有许多这样的实例记录在预言中；而且在希腊人中，所有那些因其智慧而获得巨大声誉的人，却无法说服那些阴谋反对他们的人，也无法促使他们的法官或控告者停止作恶，并努力通过哲学的道路追求美德？\par

% structure: p0172 source=h2 target=subsection reason=relative-heading-rank; base=h2

\subsection{第七十七章}

之后，那犹太人依犹太信仰发言说：\Quote{我们确实希望将有身体的复活，并享有永生；这事的榜样与原型将是那位被派遣到我们这里来的那一位，祂将显示天主没有什么是不可能的。}我们不知道那犹太人是否会说那期待的基督在自己身上显示复活的榜样，但假定他如此想、如此说。那么，我们就要这样回答那位告诉我们他是从我们自己的著作中获取信息的人：\Quote{朋友，你读了那些你认为在其中发现了控告我们的材料的著作，难道没有在那里找到耶稣的复活，以及宣告祂是死者中的首生者吗？还是因为你不会承认这些事曾被记载，它们就真的没有被记载吗？}但因那犹太人仍承认身体的复活，我认为现在不是与一个相信并说自有身体复活的人讨论这个问题的合适时机，无论他对这个问题是否有清晰的理解并能为之辩护，还是他只是出于传说而赞同。以上便是我们对这位塞尔苏斯笔下的犹太人的答复。当他补充说：\Quote{那么，祂在哪里，好让我们看见并相信祂？}我们回答：那曾在预言中说话、行奇迹的祂如今在哪里，好让我们看见并相信祂是天主的一部分？难道允许\Emph{你}提出异议说天主不持续向希伯来民族显现自己，却不允许我们为耶稣作同样的辩护吗？耶稣曾一次复活了自己，并使祂的门徒相信祂的复活，如此彻底地说服他们相信这真理，以至于他们通过自己的苦难向所有人显示他们如何能够嘲笑生活中的一切烦恼，因为他们亲眼目睹了永生与复活在言语和行为上清楚地向他们彰显。\par

% structure: p0174 source=h2 target=subsection reason=relative-heading-rank; base=h2

\subsection{第七十八章}

那犹太人继续说：\Quote{耶稣来到世界，就是为了叫我们不信祂吗？}我们立刻回答：祂来不是要造成犹太人的不信；但祂预知会有这样的结果，就预言了这事，并利用他们的不信来呼召外邦人。因为因着他们的罪，救恩临到了外邦人，关于他们，在预言中说话的基督说：\BibleQuote{我不认识的民族成了我的臣民，他们一听见就服从了我；}以及\BibleQuote{我没有寻找的人，他们找到了我；我没有求问的人，我向他们显现了。}此外，犹太人在那样对待耶稣之后，甚至在此生也受到了惩罚。当我们指控他们有罪时，任凭犹太人说什么，比如：\Quote{天主的眷顾与美善岂不是最奇妙地显示在你们的惩罚中，以及你们被剥夺耶路撒冷、圣所和你们辉煌的敬拜中吗？}无论他们如何回答关于天主眷顾的事，我们都能更有效地回应，指出天主的眷顾奇妙地彰显在利用那民族的过犯上，以便通过耶稣基督，从外邦人中那些与盟约无关、与恩许无缘的人，召叫他们进入天主的国。这些事早已被先知们预言，他们曾说，由于希伯来民族的过犯，天主将拣选，不是民族，而是从各地选出的个人；并拣选了世上的愚拙事物，使一个无知的民族认识神圣的教导，天主的国从前者被夺去，赐给后者。在众多预言中，这里仅引述\WorkTitle{申命记之歌}中关于呼召外邦人的预言，如下，是以上主的口吻说的：\BibleQuote{他们以不是神的惹动我的妒火，以他们的偶像激怒了我的怒气；我也要以不是子民的惹动他们的妒火，以愚昧的国民激怒他们的怒气。}\par

% structure: p0176 source=h2 target=subsection reason=relative-heading-rank; base=h2

\subsection{第七十九章}

关于耶稣的这一切论证的结论，被那犹太人这样陈述：\Quote{他因此是一个人，且具有这样的本性，正如真理本身所证明，理性所昭示他的那样。}然而，我不知道，一个敢于将他的宗教崇拜和教导的教义传播到全世界的人，能否在没有天主援助的情况下成就他所愿，并能超越所有抵挡他教义传播的人——君王与统治者、罗马元老院、各地的总督以及平民百姓。一个本身毫无卓越本性的人的本性，又怎能转化如此庞大的群众？如果只是智者被感召，那并不稀奇；但最无理性的人、沉溺于情欲的人，以及由于他们的无理性而极难改变以采取更有节制生活的人，也被感化了。然而，正是因为基督是天主的德能和圣父的智慧，祂才成就了，并仍在成就这样的结果，尽管不相信祂话语的犹太人或希腊人不愿承认。因此，我们不会停止按照耶稣基督的训诫信仰天主，并力图转化那些在宗教问题上盲目的人，尽管正是那些真正盲目的人指责我们盲目；而且，他们，无论是犹太人还是希腊人，误导那些跟随他们的人，却指责我们迷惑人——一种美好的迷惑，真的！——使他们从放荡变为节制，或至少向节制迈进；从不义变为正义，或至少趋于如此；从愚昧变为明智，或正在成为明智；并且，代替怯懦、卑鄙和胆怯，展现出刚毅和勇气的德行，尤其是在为他们敬拜万物造物主天主的宗教而奋斗中所彰显的。因此，耶稣基督的来临并非由一位先知所预告，而是由所有先知预告；而策尔苏斯让一个犹太人说只有一位先知预言了基督的来临，正是他无知的表现。但是，由于策尔苏斯所引入的这个犹太人，在声称这些事确实符合他自己的律法之后，在某处结束了他的论述，并提到了其他不值得记忆的事，我也将在此结束我对策尔苏斯论著的第二卷答复。但如果天主允准，且基督的德能存于我的灵魂，我将在第三卷中尽力处理策尔苏斯随后的论述。\par

\SourceRecord{Translated by Frederick Crombie.；Ante-Nicene Fathers；Vol. 4.；Edited by Alexander Roberts, James Donaldson, and A. Cleveland Coxe.；Buffalo, NY: Christian Literature Publishing Co.,；1885.；<http://www.newadvent.org/fathers/04162.htm>.}

% structure: p0001 source=h1 target=section reason=page-title

\section{\WorkTitle{驳斥切尔苏斯}, 第三卷}

% structure: p0002 source=h2 target=subsection reason=relative-heading-rank; base=h2

\subsection{第一章}

在回应切尔苏斯著作的第一卷书中——他曾狂妄地将他所撰反对我们的论著题为\WorkTitle{真道}——我们已遵照您的吩咐，忠实的安布罗修，尽我们所能，逐一审查了他的序言及紧随其后的部分，一路检验他的每一个断言，直至结束他所杜撰的那位犹太人针对耶稣的滔滔谩骂。在第二卷书中，我们又尽可能回应了那位犹太人的诋毁中针对我们这些因基督而信从天主的人所提出的所有指控；现在，我们进入本次论述的第三部分，旨在驳斥他以自己名义提出的那些主张。\par

他发表意见说：\Quote{犹太人与基督徒之间的争论是最愚蠢的}，并断言\Quote{我们彼此之间关于基督的讨论与俗语所谓的\InnerQuote{驴影之争}毫无区别}，还认为\Quote{犹太人和基督徒的研究中没有什么重要的东西：因为两者都相信圣神曾预言会有一位救主来到人类中间，但在所预言的那一位是否已经来临这一点上尚未达成一致。}因为我们基督徒确实相信耶稣就是按照先知预言而来的那一位。但大多数犹太人远不信他，以至于那些在他来临当时活着的犹太人合谋反对他；现今的犹太人则赞同那些古代犹太人对他所行的事，说他的坏话，声称他凭借巫术冒充先知所预言要来者，并依照犹太传统被称为基督。\par

% structure: p0005 source=h2 target=subsection reason=relative-heading-rank; base=h2

\subsection{第二章}

但让切尔苏斯及其赞同他指控的人告诉我们，犹太先知预言那位要作过义德生活之人统治者的诞生地——这些义人被称为天主的\Quote{产业}——又预言厄玛奴耳应由贞女怀孕，并由那位预言对象行出如此征兆和奇事，且他的话语如此迅速传播，以致宗徒们的声音传遍全地，他还将受犹太人定罪后的某些苦难，并从死者中复活，这难道与'驴影'有任何相似之处吗？因为先知们是偶然作出这些宣告，心中没有对真理的确信，不仅促使他们说话，还使他们认为这些宣告值得记录下来吗？而一个如此伟大的民族，即犹太人，他们早已拥有自己的土地安居，在一无所知的情况下承认某些人为先知，而拒绝其他人为假预言者，却没有承认这种区分的正确性吗？难道没有任何动机促使他们将那些后来被视为先知的人的话与被视为神圣的梅瑟书卷归为一类吗？那些指控犹太人和基督徒愚昧的人，能否向我们表明，如果犹太民族中间没有对未来知识应许的期盼，他们如何能持续存在？当周围每个民族都按照其古老制度相信自己从他们视为神祇的那里得到神谕和预言时，唯独这个民族——他们受教蔑视外邦人视为神祇的一切，认为它们不是神，而是魔鬼，正如先知所宣告的：\BibleQuote{外邦人的一切神祇都是魔鬼}——难道他们中间没有一个人自称是先知，能够阻止那些由于渴望知道未来而准备投靠外邦民族魔鬼的人？那么请判断，难道不是必然地，既然整个民族受教轻视其他地方的鬼神，他们就必须有众多先知，宣告那些本身远为重要并且超越所有其他国家神谕的事件吗？\par

% structure: p0007 source=h2 target=subsection reason=relative-heading-rank; base=h2

\subsection{第三章}

其次，奇事在各处，或至少在众多地方发生过，正如切尔苏斯自己所承认的，他举了埃斯科拉庇俄斯为例，此人曾惠及许多人，并向众多城市——如特里卡、埃皮达鲁斯、科斯和佩尔加蒙——预言未来事件；除了埃斯科拉庇俄斯，他还提到普罗孔尼索斯的阿瑞斯特亚斯、某位克拉佐梅尼人以及阿斯提帕莱亚的克莱俄墨得斯。但唯独在犹太人中间——他们声称自己是奉献给万有之天主的——竟然没有发生任何奇事或征兆，用以帮助他们坚定对创造万物的天主的信仰，并增强他们对另一种更美好生命的希望！但他们怎能设想这样一种情形呢？因为他们会立即转而崇拜那些给予神谕并施行医治的魔鬼，而离弃那位只在言语上被认为帮助他们、却从未以可见方式向他们显示自己的天主。但如果这种结果并未发生，相反，他们遭受了无数灾难，却仍不愿放弃犹太教及其律法，并曾在亚述、波斯和安提约古统治下遭受残酷对待，那么对于那些不轻信他们的奇迹历史和预言的人而言，岂不应当按事情的可能性设想：这些事件不可能是虚构，而是在那些圣洁的先知灵魂中，如同那些为德行而付出任何劳苦的人一样，有一种神圣的灵确实感动他们预言一些关于当时的事，另一些关于后代的事，但主要关于某一位要来作人类救主的人物吗？\par

% structure: p0009 source=h2 target=subsection reason=relative-heading-rank; base=h2

\subsection{第四章}

若是上述情形，那么犹太人和基督徒如何是在追寻\Quote{驴的影子}，试图从他们共同相信的预言中确定那被预言的是已经来临，还是尚未到来，仍是期待的对象？但即使我们向策尔苏斯让步，承认预言所宣告的不是耶稣，那么即便在此假设下，探究预言书的意义也不是追寻\Quote{驴的影子}，只要能够清楚地指出所谈论的那一位，并表明他被人预言为何等人物、要做何事，且可能的话，预言他何时来临。然而在前文中我们已经就耶稣是被预言为基督的那一位进行了论述，引用了众多预言中的一部分。因此，犹太人和基督徒假设先知是受神圣感召而发言，这并不错；错的是那些对所期待的先知预言的来临者怀有错误观念的人：先知们的\Quote{真言论}已揭示了那来临者的位格与特征。\par

% structure: p0011 source=h2 target=subsection reason=relative-heading-rank; base=h2

\subsection{第五章}

紧接这些论点之后，策尔苏斯臆想犹太人是埃及人的后裔，他们离开了埃及，因反叛埃及国家而离开，并轻视了该民族在敬拜方面的习俗，他说\Quote{他们遭到耶稣的信徒（这些信徒相信他是基督）的对待，正如他们曾加在埃及人身上的相同待遇；并且导致两种新状况的原因都是背叛国家。}现在让我们看看策尔苏斯在此做了什么。古代埃及人曾对因犹太发生大饥荒而定居埃及的希伯来民族施以许多残暴，后来因他们对客居者和求助者的不义而遭受了惩罚，这是天主眷顾所定的、要降在整个民族身上的惩罚，因为他们合谋反对整个曾是他们的客人、并未伤害过他们的民族；在受到来自天主的灾祸打击之后，他们艰难地允许希伯来人在短时间后到他们愿意去的任何地方，因为这些希伯来人被不公正地扣为奴隶。因此，因为他们是一个自私的民族，尊崇任何与他们稍有关系的人远超过尊崇品行更好的陌生人，他们对梅瑟和希伯来人提出了一切可能的指控——并非完全否认他确实行过奇迹奇事，而是声称这些是巫术所成，而非出自神圣能力。然而梅瑟并非巫师，而是一个虔诚的人，他献身于万有之天主，是圣神的分享者；他按照天主的启示为希伯来人颁布法律，并完全忠实地记载了所发生的事件。\par

% structure: p0013 source=h2 target=subsection reason=relative-heading-rank; base=h2

\subsection{第六章}

策尔苏斯并非以公正精神调查这些事实（埃及人以一种方式叙述，希伯来人以另一种方式叙述），而是仿佛被迷住了似的偏向前者，他将那些压迫客居者的话当作真理接受，宣称曾受不公待遇的希伯来人在反叛埃及人之后离开了埃及——他没有看到如此一大群反叛的埃及人怎么可能成为一个民族：这个民族以那次反叛为起源，并在反叛时改变了自己的语言，以致那些直到那时都使用埃及语的人，竟能一下子完全采用希伯来人的语言！然而，就算按他的假设，他们在离开埃及后也对母语产生了憎恶，那么为什么他们这样做之后，没有转而采用叙利亚语或腓尼基语，却选择了与这两种都不同的希伯来语呢？但我认为理性证明这一说法是虚假的，即那种声称原本是埃及人的人反叛了埃及人、离开国土、前往巴勒斯坦并占据了今日称为犹太之地的说法。因为希伯来语是他们祖先下埃及之前的语言；并且梅瑟在撰写被犹太人视为神圣的五卷书时所使用的希伯来文字，与埃及文字不同。\par

% structure: p0015 source=h2 target=subsection reason=relative-heading-rank; base=h2

\subsection{第七章}

同样，\Quote{希伯来人（原本）是埃及人，他们（政治存在）的开端始于他们反叛之时}这一说法是虚假的，同样虚假的是\Quote{在耶稣的时代，另一些犹太人反叛了犹太国家，成为了他的追随者}；因为策尔苏斯以及与他想法一致的人都无法指出基督徒有任何反叛行为。然而，如果一次反叛导致了基督徒团体的形成，以至于它这样从犹太人那里获得了存在（这些犹太人被允许拿起武器保卫他们的家庭成员并杀死敌人），那么基督立法者本不会完全禁止杀人；然而他从未教导他的门徒可以暴力对待任何人，无论那人多么邪恶。因为他认为，像他这样的源于神圣根源的法律，不允许杀害任何个人。如果基督徒起源于反叛，他们也不会采纳如此极其温和的法律，以致当他们注定要像羊一样被宰杀时，他们从不抵抗迫害者。事实上，如果我们稍微深入思考，我们可以就出埃及一事说：如果一个整个民族\Emph{一下子}采用了被称为希伯来语的语言，好像它是来自天上的恩赐，那就是一个奇迹，因为他们的一个先知说：\Quote{当他们从埃及出来时，他们听到了他们不懂的语言。}\par

% structure: p0017 source=h2 target=subsection reason=relative-heading-rank; base=h2

\subsection{第八章}

我们也可以通过以下方式得出结论：那些与梅瑟一同离开埃及的人并非埃及人；因为如果他们曾是埃及人，他们的\Emph{名字}也会是埃及式的，因为在每一种语言中，（人及事物的）名称都与该语言同源。但若从名字是希伯来式的这一点可以确定这些人并非埃及人，——而圣经中充满了希伯来名字，且这些名字是在他们于埃及时就赐予其子女的——那么显然，埃及人的说法是虚假的，该说法声称他们是埃及人，并与梅瑟一同离开埃及。现在绝对确定的是，正如梅瑟历史所记载的，他们从希伯来祖先延续而来，使用一种语言，他们也从该语言中取用名字赋予其子女。但关于基督徒，因为他们被教导不可向敌人复仇（并因此遵守了温和而仁爱的律法）；并且因为他们即使有能力，也不会发动战争，即使他们获得了这样做的权柄——他们从天主那里得到了这一赏报：天主常为他们作战，并在某些时刻阻止那些起来反对他们、想要消灭他们的人。为了提醒他人，使他们通过看到\Emph{少数}人为其宗教而奋斗，也能更好地藐视死亡，有些人，在特殊时机，且这些人数目有限，为基督信仰的缘故忍受了死亡——天主不允许整个民族被消灭，而是希望它延续，并使整个世界充满这有益而神圣的教义。另一方面，为了使意志软弱的人重获勇气，超越对死亡的思虑，天主以其眷顾介入信徒之事，仅凭祂意志的行动就驱散了一切针对他们的阴谋；因此，无论是君王、统治者还是民众，都不能超过一定限度对他们施暴。这就是我们对策尔苏斯陈述的回答，他说\Quote{叛乱是古代犹太民族的最初开端，随后也是基督信仰的开端。}\par

% structure: p0019 source=h2 target=subsection reason=relative-heading-rank; base=h2

\subsection{第九章}

但既然他在接下来的陈述中明显犯有虚假之罪，让我们审视他所说的断言：\Quote{如果所有人都愿意成为基督徒，后者却不希望有这样的结果。} 上述陈述是虚假的，这一点由此而明：基督徒在他们能力范围内，并没有忽视采取各种措施将他们的教义传播到全世界。因此，他们中的一些人不仅致力于周游城市，甚至走访村庄和乡舍，以便为天主赢得皈依者。没有人会认为他们这样做是为了谋求利益，因为有时他们甚至不接受必要的生计；或者如果任何时候受此需要所迫，他们只满足于所需之物的供给，尽管许多人愿意与他们分享（自己的富足），并给予远超他们所需的帮助。诚然，在当今，由于基督徒信友众多，不仅富人，还有有地位的人、娇贵而高贵的女士们都接待基督信仰的教师，有些人或许敢于说，某些人担任基督信仰教师之职是为了些许荣耀。然而，对基督信仰最初之时理性地产生这种怀疑是不可能的，因为那时所涉及的危险，尤其对教师而言是巨大的；而在当今，它在其余人类中所蒙受的耻辱，远大于那些持有相同信仰者中所享有的任何所谓的荣誉，尤其是当这种荣誉并非为所有人所共享时。因此，从事情的本质来看，说\Quote{如果所有人都愿意成为基督徒，后者却不希望有这样的结果}是虚假的。\par

% structure: p0021 source=h2 target=subsection reason=relative-heading-rank; base=h2

\subsection{第十章}

但请注意他作为其论断的所谓证据：\Quote{基督徒起初人数稀少，且持相同见解；但当他们增长为庞大群体时，便分裂四散，各人都想拥有自己的派别：因为这是他们从一开始的目的。} 基督徒起初人数稀少，相比后来成为基督徒的众多群众而言，这是毋庸置疑的；然而，综观一切，他们并非如此稀少。因为激起犹太人嫉妒耶稣并煽动他们密谋反对祂的，正是跟随祂进入荒野的众多人数——一次有五千男人，另一次有四千男人在那里跟随祂，还不包括妇女和儿童。耶稣言语的魅力如此之大，以至于不仅\Emph{男人}愿意跟随祂到荒野，而且\Emph{妇女}也忘记了性别的软弱和外出礼貌，就这样跟随他们的导师进入旷野。孩子们也同样不受这类情绪影响，或跟随父母，或也可能被祂的神性所吸引，为使这神性植入他们内，便与父母一同成了祂的跟随者。但假定基督徒起初人数稀少，这如何有助于证明基督徒不愿使所有人都相信福音的教义呢？\par

% structure: p0023 source=h2 target=subsection reason=relative-heading-rank; base=h2

\subsection{第十一章}

他接着又说，\Quote{所有的基督徒都同心合意，}却未注意到，即使在这一方面，从一开始，信徒们对于被视为神圣的经卷之含义就已经存在意见分歧。当宗徒们还在宣讲，亲眼见证\Person{耶稣}（事迹）的人还在教导祂的道理时，从犹太教皈依的人之间就外邦信徒是否应当遵守犹太习俗，或者应当摒弃洁净与不洁食物的重担——因为那些已经放弃祖传外邦习俗并成为\Person{耶稣}信徒的人并无义务遵守这一点——发生了不小的争论。甚至在与曾见过\Person{耶稣}的人同时代的\Person{保禄}的书信中，也提到了一些成为争论主题的具体事项——即关于复活，以及复活是否已经过去；还有主的日子，是否临近。就连那劝勉要\BibleQuote{远避世俗的虚谈和那假知识的反对；有人自称有这知识，却偏离了信德}\BibleRef{弟前 6:20-21}，也足以表明从一开始，当（据\Person{刻尔苏斯}想象）信徒人数稀少时，某些教义就已有不同解释。\par

% structure: p0025 source=h2 target=subsection reason=relative-heading-rank; base=h2

\subsection{第十二章}

其次，既然\Person{刻尔苏斯}指责我们中间存在异端作为反对基督信仰的理由，说\Quote{当基督徒人数大增时，他们分裂成派别，各人都想有自己的党派；}又说\Quote{如此因人众多而分裂，他们互相驳斥，却仍有共同的一个\Emph{名}，如果他们确实还保留这个名的话。而这是他们尚羞于放弃的唯一事物，其他事则由各派别以不同方式决定。}作为回应，我们说：不同种类的异端从未源于任何在其原则无关紧要或对人类生活无益的问题上。因为医学科学对人类有用且必不可少，而在如何治疗身体方面存在许多争论点，因而在希腊人当中以及在我也认为的那些自称使用医学的蛮族当中，医学科学中显然流行着众多异端。同样，既然哲学以真理为职业，并许诺为调节生活而认识现存事物，且努力教导何为有益于人类种族，而探讨这些问题伴随着极大的意见分歧，因此哲学中便产生了无数异端，其中一些比另一些更著名。就连犹太教本身，由于对\Person{梅瑟}和先知著作的不同理解，也为异端的产生提供了借口。那么，鉴于基督信仰在人们眼中成为可敬之物，不仅对\Person{刻尔苏斯}所认为的更为卑微的阶层，而且对许多从事文学研究的希腊人来说，都必然产生异端——这根本不是出于党派之争或纷争，而是因为许多文人热切希望了解基督信仰的教义。其结果是，人们对那些被普遍视为神圣的言论作了不同理解，于是产生了异端，这些异端以那些确实钦佩基督信仰起源但因某种似是而非的理由而陷入不一致观点之人的名字命名。然而，没有人会因医学有异端而理性地回避医学；也没有追求善好的人会憎恨哲学，并以哲学有许多异端作为其反感的借口。同样，也不应因犹太教中的异端而谴责\Person{梅瑟}和先知的神圣经卷。\par

% structure: p0027 source=h2 target=subsection reason=relative-heading-rank; base=h2

\subsection{第十三章}

如果这些论证成立，那么我们为何不以同样方式为基督信仰中的异端辩护呢？关于这些，\Person{保禄}似乎以非常引人注目的方式说到：\BibleQuote{因为在你们中间不免有异端，好叫那些被证验的人，在你们中间显明出来。}\BibleRef{格前 11:19}因为正如在医学中，那个人因对多种医学异端的经验以及对其中大多数的诚实审视而选择了更优的体系，而被\Quote{证验}；又如哲学中的大成就者，是在经验性地了解各种观点之后，赞同了最好的；同样，我会说，最智慧的基督徒是仔细研究过犹太教和基督信仰异端的人。然而，因基督信仰有异端就指责它的人，也会指责\Person{苏格拉底}的教导，从他的学派中产生了许多持不同观点的其他学派。甚至\Person{柏拉图}的意见也可能因\Person{亚里士多德}脱离其学派并创立新学派而被指为错误——关于这一点，我们在前书已有论述。但在我看来，\Person{刻尔苏斯}已经熟悉了某些异端，它们甚至在\Quote{耶稣}之\Emph{名}的共同使用上都不与我们一致。也许\Person{刻尔苏斯}听说过称为\OriginalTerm{敖非特派}{Ophites}和\OriginalTerm{加音派}{Cainites}的派别，或其他类似性质的派别，这些派别在所有方面都已偏离了\Person{耶稣}的教导。然而，这当然不能成为指控基督信仰教义的理由。\par

% structure: p0029 source=h2 target=subsection reason=relative-heading-rank; base=h2

\subsection{第十四章}

此后他接着说：\Quote{他们的结合越是显得没有实质理由，就越令人惊奇。然而，反叛是一条实质理由，还有从反叛中得来的利益，以及对外敌的恐惧。这些便是使他们的信仰得以稳固的原因。}对此我们回答：我们的结合确实基于一个理由，或者更确切地说，并非基于一个理由，而是基于天主的作为，因而其开端是天主在先知著作中教导人类期待基督的来临，祂将是人类的救主。因着这一点并未被真正驳倒（尽管在不信者看来\Emph{似乎}如此），同理，此教义被称赞为天主的教义，耶稣也被显明为天主子，无论是在祂降生成人之前还是之后。此外，我主张，即使在祂降生成人之后，那些拥有最敏锐属灵眼光的人总是发现祂最具神性，确实从天主那里降临到我们中间，祂的起源或后来发展并非来自人的智慧，而是来自天主在祂内的显示，天主以祂多方面的智慧和奇蹟先建立了犹太教，随后又建立了基督教；而认为反叛及其带来的利益是这一教义（它归化了并改善了如此多人）的起源原因的主张，已被有效地驳斥了。\par

% structure: p0031 source=h2 target=subsection reason=relative-heading-rank; base=h2

\subsection{第十五章}

然而，再次强调，并非对外敌的恐惧巩固了我们的结合，这一点是显而易见的，因为这一原因，出于天主的旨意，已经有一段时间不存在了。而且，就世俗方面而言，信徒们目前享有的安全存在很可能会结束，因为那些从各方面诽谤基督教的人再次将当前频繁的反叛归咎于信徒众多，以及他们没有像往昔那样受到当局的迫害。因为我们已经从\WorkTitle{福音}中学到：既不要在和平的日子里放松努力，放纵休息，也不要在世界向我们开战时变成懦夫，背弃那在\Person{耶稣基督}内的、对万有之天主的爱。并且，当我们向首批归化者灌输对偶像及各类肖像的鄙视，此外，还将他们的思想从崇拜受造物而非天主提升到宇宙的创造者时，我们清楚地彰显了我们起源的辉煌本质，并没有（如\Person{塞耳苏}所想象的那样）隐藏它。我们清楚地表明祂是预言的题旨，这既来自关于祂的众多预言，也来自那些能够明智地理解\WorkTitle{福音书}和宗徒们的宣告的人所仔细探究的传统。\par

% structure: p0033 source=h2 target=subsection reason=relative-heading-rank; base=h2

\subsection{第十六章}

\Quote{然而，我们聚集的各类传说，或我们编造的恐怖事情是什么？}正如\Person{塞耳苏}毫无根据地断言，谁喜欢展示就可以展示。我确实不知道他说\Quote{编造恐怖事情}是什么意思，除非是指我们关于天主作为审判者、人因行为受惩罚的教义，并且有各种证据，部分来自\WorkTitle{圣经}，部分来自或然理由。然而——因为真理是宝贵的——\Person{塞耳苏}在结尾说：\Quote{但愿我、这些人，或任何其他个体，永远不要排斥有关恶人将来受罚、善人得赏的教义！}那么，如果除去惩罚的教义，我们编造并强加给人类什么恐怖事情呢？如果他回答说：\Quote{我们编造从古代来源得出的错误意见，大声鼓吹，在人们面前吹嘘，如同\Person{库柏勒}的司祭在那些被纳入其秘仪的人耳边敲击铙钹一样；}我们就要反问他：\Quote{错误的意见来自哪些古代来源？}\par

% structure: p0035 source=h2 target=subsection reason=relative-heading-rank; base=h2

\subsection{第十七章}

他确实希望将我们信仰的信条与埃及人的相比；\Quote{在埃及人中，当你走近他们的圣殿时，可以看到辉煌的围地、丛林、又大又美的门廊、奇妙的圣殿、周围宏伟的帐幕，以及充满迷信与秘仪的崇拜仪式；但当你进入并走进去之后，所崇拜的对象被看见是一只猫、一只猿、一条鳄鱼、一只山羊或一条狗！}那么，我们与那些走近埃及圣殿的人所见的辉煌之间有什么相似之处？在穿过辉煌的门廊之后，在里面被崇拜的那些无理性的动物之间又有什么相似？难道我们的预言、万有之天主、禁止圣像的命令，在\Person{塞耳苏}看来也是敬畏的对象，而钉在十字架上的\Person{耶稣基督}，与崇拜无理性的动物是类似物吗？但如果他断言这一点——我认为他不会坚持别的——我们将回答说，我们在前文已经为那些影响耶稣的指控作了更详细的辩护，表明那些似乎以祂的人性身份发生在祂身上的事，富于对全人类的益处，以及对普世的救恩。\par

% structure: p0037 source=h2 target=subsection reason=relative-heading-rank; base=h2

\subsection{第十八章}

在接下来的段落中，谈及埃及人的说法，他们夸夸其谈地谈论无理性的动物，并断言它们是天主的某种象征，或按他们所谓的先知惯常称呼它们的任何名称，\Person{刻尔苏斯}说：\Quote{那些学了这些事的人心中产生一种印象：他们并未徒然领受入门仪式；} 然而对于基督教研究中已进步的人（借着\Person{保禄}称之为藉圣神所赐的\Quote{智慧的言语}和按圣神\Quote{知识的言语}的恩赐）所教导的真理，\Person{刻尔苏斯}似乎连概念都没有形成，这不但从他已说的话，也从他后来攻击基督教体系时所加添的话可以判断，他断言基督徒\Quote{将一切有智慧的人从他们的信仰教义中驱赶出去，只邀请无知和粗俗的人；} 关于这些断言，我们将在适当之时，当轮到适当位置时加以评论。\par

% structure: p0039 source=h2 target=subsection reason=relative-heading-rank; base=h2

\subsection{第十九章}

他确实说：\Quote{我们嘲笑埃及人，尽管他们提出许多绝不卑微的奥秘供我们思考，他们教导我们说，这种仪式是对永恒理念的敬拜，而不是像大众所想的那样是对短暂动物的敬拜；而我们却是愚蠢的，因为我们在关于耶稣的叙述中没有引入比埃及敬拜中的山羊和狗更高尚的东西。} 对此我们回复：\Quote{仁兄，假设您正确称赞了埃及人提出许多绝不卑微的奥秘以及对在他们中间（被敬拜的）动物的隐晦解释，然而您指责我们时却言行不一，仿佛您相信我们没有任何值得考虑的东西可说，而是所有我们的教义都是可鄙的、无足轻重的，因为我们按照\InnerQuote{智慧之言}向在基督徒中\InnerQuote{成全的人}阐述关于耶稣的叙述。关于这些人，\Person{保禄}认为他们能够理解在基督教中的智慧，说：\InnerQuote{我们在成全的人中讲智慧，但不是今世的智慧，也不是今世有权位者将要消灭的智慧；我们讲的乃是天主奥秘的智慧，那隐藏的智慧，天主在万世以前预定为使我们得光荣的；今世有权位者没有一个认识这智慧。}}\par

% structure: p0041 source=h2 target=subsection reason=relative-heading-rank; base=h2

\subsection{第二十章}

我们对那些与\Person{刻尔苏斯}持相似看法的人说：\Quote{那么\Person{保禄}，我们假定，当他宣称在成全的人中讲智慧时，他心中没有卓越的智慧的观念吗？} 既然他以其惯常的坦率如此声明，说\Emph{他}没有智慧，我们将首先回答：检查说出这些话的人的书信，仔细查看其中每个表达的意义——例如在\WorkTitle{厄弗所书}、\WorkTitle{哥罗森书}、\WorkTitle{得撒洛尼书}、\WorkTitle{斐理伯书}、\WorkTitle{罗马书}中——并说明两点：既理解\Person{保禄}的话，又能证明其中任何一点是愚蠢的或愚拙的。因为若有人专心勤读，我确信他要么会惊叹那能用普通语言表达伟大思想的人的理解力；要么若不惊叹，他只会显出自己可笑，无论他是像已经理解那样简单地陈述作者的意思，还是尝试反驳和驳斥他自以为理解的东西！\par

% structure: p0043 source=h2 target=subsection reason=relative-heading-rank; base=h2

\subsection{第二十一章}

我尚未谈及对\WorkTitle{福音}所写一切内容的遵守，每部\WorkTitle{福音}都包含许多难以理解的教义，不仅对大众，甚至对某些更聪明的人也如此，包括耶稣对\Quote{外面的人}所讲的比喻的极其深奥的解释，而将其完全含义的展示保留给那些已经超越外表教导阶段、在屋内私下到祂面前的人。当人理解这一点时，他会惊叹为什么有些人被称为\Quote{外面的人}，而另一些\Quote{在屋内}。再者，谁能理解耶稣的行动而不充满惊奇？祂时而登山以传授某些讲论或行某些奇迹，或为祂自己的显圣容，又下山医治病人和那些不能跟随祂到祂门徒所去之处的人。但目前不适合描述\WorkTitle{福音}中真正可敬和神圣的内容，或基督的心意——即智慧和圣言——包含在\Person{保禄}的著作中。但我们所说的已足够回答\Person{刻尔苏斯}那非哲学的嘲讽，他将天主的教会的内在奥秘与埃及的猫、猿、鳄鱼、山羊和狗相比较。\par

% structure: p0045 source=h2 target=subsection reason=relative-heading-rank; base=h2

\subsection{第二十二章}

但是，这位卑劣的戏谑者\Person{塞耳苏}，对我们不放过任何形式的嘲弄和取笑，在他的论著中提到了\Person{狄奥斯库里}、\Person{赫拉克勒斯}、\Person{埃斯库拉庇乌斯}和\Person{狄奥尼索斯}，希腊人相信他们是从人变成的神，并说\Quote{我们无法称这类存在为神，因为他们最初是人，但他们却显露出许多为人类利益而展示的高贵品质，而我们却断言耶稣死后被祂自己的追随者看见}；他还额外指控我们，好像我们说\Quote{祂确实被看见了，但只是一个影子}！对此我们回应道：塞耳苏非常狡猾，他没有在这里明确表示他不把这些存在视为神，因为他害怕那些可能阅读他论著的人的意见，那些人可能会认为他是无神论者；然而，如果他尊重在他看来是真理的事物，他就不会\Emph{伪装}把他们视为神。现在，我们对这两项指控都准备好了回答。因此，让我们对那些\Emph{不}把他们视为神的人这样回答：那么，这些存在根本不是神；但是，按照那些认为人的灵魂（死后）立即消亡的人的观点，这些人的灵魂也消亡了；或者，按照那些说灵魂持续存在或不朽的人的观点，这些人继续存在或是不朽的，他们不是神而是英雄——或者甚至不是英雄，而只是灵魂。因此，一方面，如果你假设他们\Emph{不}存在，我们将必须证明灵魂不朽的教义，这对我们来说是一个极为重要的教义；另一方面，如果他们\Emph{确实}存在，我们\Emph{仍然}必须证明不朽的教义，不仅通过希腊人所对此做出的精彩论述，而且以符合圣经教导的方式。我们将证明，那些生前是多神论者的人在离开这个世界后不可能获得更好的地方和地位，我们将引用关于他们的故事，其中记载了\Person{赫拉克勒斯}极大的放荡以及与\Person{翁法勒}的软弱奴役，以及关于\Person{埃斯库拉庇乌斯}的说法，说他们的\Person{宙斯}用闪电将他击死。至于\Person{狄奥斯库里}，将会说他们常常死去——\par

有一天活在交替的日子里，另一天\newline{}死去，并与诸神同等地获得尊荣。\par

那么，他们怎能合理地认为其中一个是应当被视为神或英雄呢？\par

% structure: p0049 source=h2 target=subsection reason=relative-heading-rank; base=h2

\subsection{第二十三章}

然而，我们通过先知性的圣经证明与我们的\Person{耶稣}有关的事实，并将祂的历史与圣经比较，证明没有记载祂有任何放荡。因为连那些密谋反对祂、寻找假见证来帮助祂的人，也没有找到任何看似有理的借口来提出虚假指控，以控告祂淫荡；但祂的死确实是阴谋的结果，与\Person{埃斯库拉庇乌斯}被闪电击死毫无相似之处。而疯癫的\Person{狄奥尼索斯}和他的女性服装有什么可敬之处，以至于\Emph{他}应当被当作神崇拜？如果那些为这类存在辩护的人诉诸寓意解释，我们必须逐一检查每个实例，确定它是否有根据，并且在每个具体情况下，那些被\Person{泰坦}撕裂、从天上宝座摔落的存在是否能够真实存在，是否值得尊敬和崇拜。而我们的\Person{耶稣}，祂向祂自己那队成员显现——因为我要采用\Person{塞耳苏}使用的词——确实\Emph{真正}显现了，塞耳苏对福音作了虚假指控，说所显现的是一个影子。让他们的历史和耶稣的历史仔细比较。塞耳苏会认为前者是真实的，而后者虽是目击者所记载，他们用行动表明他们清楚知道所见之事的性质，并愉快地为祂的福音受苦来显示他们的心态，难道是捏造的吗？现在，有谁渴望始终按正确理性行事，会随意同意关于一方的叙述，却冲向耶稣的历史，不经审查就拒绝相信关于祂的记载呢？\par

% structure: p0051 source=h2 target=subsection reason=relative-heading-rank; base=h2

\subsection{第二十四章}

再者，当说到\Person{埃斯库拉庇乌斯}时，说许多希腊人和外邦人承认他们经常看到，并且仍然看到，不是幻影，而是\Person{埃斯库拉庇乌斯}本人，在医治、行善并预言未来；\Person{塞耳苏}要求我们相信这一点，并且对我们这些相信\Person{耶稣}的人没有挑剔，当我们表达对这些故事的相信时，但当我们赞同耶稣奇迹的目击者和见证人时，他们清楚显明他们信仰的诚实（因为我们看到他们的纯朴，如同在文字中可见的良心的揭示），他却称我们为一群\Quote{愚蠢}的人，尽管他无法证明如他宣称的那样，无数的希腊人和外邦人承认埃斯库拉庇乌斯的存在；而我们，如果我们认为这是重要的事情，可以清楚地显示无数的希腊人和外邦人承认耶稣的存在。有些人通过这种信德接受了神奇能力的证据，通过他们在需要帮助的人身上所行的治愈，不引用任何其他名字，只引用万有的天主和耶稣的名字，并提及祂的历史。因为藉着这些方法，我们也曾见过许多人摆脱了严重的灾难、精神错乱、疯狂以及无数其他疾病，这些疾病既不能被人也不能被魔鬼治愈。\par

% structure: p0053 source=h2 target=subsection reason=relative-heading-rank; base=h2

\subsection{第二十五章}

现在，为了承认确实存在一位名叫\Person{阿斯克勒庇俄斯}{Æsculapius}的治愈之灵，他曾治愈人的身体，我要对那些对这类事件，或对\Person{阿波罗}{Apollo}的先知知识感到惊讶的人说：既然治愈身体是无关紧要的事，不仅善人，而且恶人也能做到；既然预知未来也是无关紧要的事——因为拥有预知能力的人未必表现出拥有德行——那么你必须证明，那些从事医治或预言未来的人在任何方面都不是邪恶的，而是展示出完美的德行典范，且离被视为神明不远。但他们将\Emph{不}能证明那些从事医术或拥有预知天赋的人是良善的，因为有许多不配活着的人据记载曾被治愈；而这些人的生活不正当，明智的医生不会希望治愈他们。在\Place{皮提亚神谕}{Pythian oracle}的答复中，你也会发现一些不合情理的命令，我们在此举出其中两条：即，当它命令\Person{克莱奥梅德斯}{Cleomedes}——我想是拳击手——应被尊以神圣的荣耀，认为他的拳击技艺有某种重大价值，却没有将授予拳击的荣誉授予\Person{毕达哥拉斯}{Pythagoras}或\Person{苏格拉底}{Socrates}；还有，当它称\Person{阿尔基洛科斯}{Archilochus}为\Quote{缪斯的仆人}——这个人将其诗才用于最邪恶、最放荡的主题，且其公共品行堕落不洁——并称他为\Quote{虔诚的}，因他是\Person{缪斯}{Muses}的仆人，而\Person{缪斯}{Muses}被视为女神！现在我倾向于认为，没有人会断言一个不具备一切节制和美德的男子是\Quote{虔诚的}，或者一个庄重的人会说出\Person{阿尔基洛科斯}{Archilochus}不雅抑扬格中包含的那种言辞。如果\Person{阿斯克勒庇俄斯}{Æsculapius}的医治技能或\Person{阿波罗}{Apollo}的先知能力本身没有任何神圣之处得到证明，那么即使我承认事实如所述，一个人又怎能合理地视他们为纯粹的神明来崇拜呢？——尤其是\Person{阿波罗}{Apollo}的先知之灵，纯净而无任何尘世之体，竟秘密地通过私处进入被称为女祭司的人的身体，当她坐在\Place{皮提亚洞窟}{Pythian cave}的入口时！然而关于\Person{耶稣}{Jesus}和祂的能力，我们没有这样的概念；因为\Person{童贞玛利亚}{Virgin Mary}所生的身体是由人性材料构成，并能承受人的创伤和死亡。\par

% structure: p0055 source=h2 target=subsection reason=relative-heading-rank; base=h2

\subsection{第二十六章}

让我们看看\Person{色尔苏斯}{Celsus}接下来说了什么，当他从历史中举出一些本身似乎难以置信的奇异事件，但至少从他的话语来看，他并未否认这些事件。首先，关于\Place{普罗科尼索斯}{Proconnesus}的\Person{阿里斯特亚斯}{Aristeas}，色尔苏斯这样说道：\Quote{那么，关于普罗科尼索斯的阿里斯特亚斯，他以一种如此显示神圣干预的方式从人群中消失，又以如此明确的方式再次显现，并在后来的许多场合访问了世界的许多地方，宣布了奇妙的事件，阿波罗命令梅塔蓬图姆的居民将他视为神，却没有人认为他是神。}这个叙述他似乎取自\Person{品达}{Pindar}和\Person{希罗多德}{Herodotus}。然而，此刻只引用后者从其史集第四卷的陈述就足够了，其大意如下：\Quote{作出这些诗句的阿里斯特亚斯是哪个国家的人，已经提到过；我现在要讲述我在普罗科尼索斯和库齐库斯听到的关于他的叙述。他们说，阿里斯特亚斯，出生不低于任何市民，进入普罗科尼索斯一家洗衣店，突然死亡；洗衣工关上店门，去通知死者亲属。当阿里斯特亚斯已死的消息传遍全城时，一个从阿尔塔刻来的库齐库斯人，与那些报告者发生了争执，声称他在去库齐库斯的路上遇见并与阿里斯特亚斯交谈过；他激烈地争辩说报告不真实；但死者的亲属去了洗衣店，带去了必要的东西以便运走尸体；但当房子打开时，阿里斯特亚斯既不见尸体也不见活人。他们说，后来，在第七年，他在普罗科尼索斯出现，创作了那些被希腊人称为\WorkTitle{阿里马斯比亚}的诗句，创作完成后，第二次消失了。这就是在这些城市流传的故事。但据我所知，这些事情发生在意大利的梅塔蓬图姆人身上，是在阿里斯特亚斯第二次消失后340年，这是我通过在普罗科尼索斯和梅塔蓬图姆计算发现的。梅塔蓬图姆人说，阿里斯特亚斯本人出现在他们的国家，劝诫他们为阿波罗建造一座祭坛，并在它旁边放置一座刻有普罗科尼索斯人阿里斯特亚斯之名的雕像；因为他说，阿波罗访问了他们的国家，在全体意大利人中唯独访问了他们；而他本人，即现在的阿里斯特亚斯，陪伴着阿波罗；并说当他陪伴这位神时，他是一只乌鸦；说完这些后，他就消失了。梅塔蓬图姆人说，他们派人去德尔斐询问神，这个人的显现意味着什么；但\Place{皮提亚}{Pythian}命令他们服从这个显现，如果他们服从，将对他们有益。他们于是得到了这个答复，就执行了命令。现在，一座刻有阿里斯特亚斯之名的雕像被放置在阿波罗像旁边，周围种植了月桂树：雕像放置在市集广场上。关于阿里斯特亚斯就说到这里。}\par

% structure: p0057 source=h2 target=subsection reason=relative-heading-rank; base=h2

\subsection{第二十七章}

如今，针对关于阿里斯提亚斯的这段叙述，我们必须回答说：如果凯尔苏斯仅将此作为历史提出，而不同意其真实性，那我们本会以另一种方式回应他的论点。但他断言他\Quote{因神性之介入而消失}，\Quote{又以无可置疑的方式再次显现}，\Quote{访问了世界许多地方}，\Quote{并作出了惊人的宣告}；此外，还有\Quote{一条阿波罗的神谕，命令墨塔蓬提人将阿里斯提亚斯当作神明崇拜}，他以此作为自己的权威并完全同意这些叙述。既然如此，我们要问：怎么可能？你，凯尔苏斯，一方面认为耶稣门徒们关于其师所传的奇迹完全是虚构，并指责那些相信他们的人，另一方面却不把这些你所说的故事视为戏法或编造？你一方面指责他人盲目相信耶稣的奇迹，另一方面却毫无证据地相信上述说法，又怎能证明自己的可信？难道希罗多德和品达在你看来是说实话的，而那些如此关心为耶稣的教义而死，并为后人所相信的真理留下了如此卓越著作的人，为了你所认为的\Quote{虚构}、\Quote{神话}和\Quote{戏法}而投身于充满危险的生活和暴力死亡吗？那么，将有关阿里斯提亚斯和耶稣的记载置于中间立场，从结果以及从道德改善和崇拜至上天主所带来的益处来看，你是否能看到：我们必须相信耶稣所发生的事迹并非没有神性干预，而关于普洛孔涅索斯人阿里斯提亚斯的故事则不然？\par

% structure: p0059 source=h2 target=subsection reason=relative-heading-rank; base=h2

\subsection{第二十八章}

因为，如果阿里斯提亚斯的相关奇迹是出于上智安排，那是为了什么目的？这些在你看来如此非凡的事件，给人类带来了什么益处？你无法回答。但我们叙述耶稣生平的事件时，并非没有普通的辩护理由，即：天主愿意推荐耶稣的教义作为拯救人类的教义，这教义确实基于宗徒们，如同正在兴建的基督教大厦的基础，但在随后的世代中也不断增长，其间不少治愈以耶稣之名完成，还发生了某些其他不小的显现。那么，阿波罗是哪一类神明？他命令墨塔蓬提人将阿里斯提亚斯当作神明崇拜？他这样做有何目的？这种崇拜对墨塔蓬提人有什么好处？如果他们将不久前还是凡人的人当作神明，那有什么益处？然而，阿波罗（被我们视为一个获得了奠祭和祭香之荣耀的魔鬼）关于阿里斯提亚斯的建议在你看来值得考虑；而万有之天主及其圣天使通过先知——不是在耶稣出生之后，而是在祂显现于人类之前——预先宣告的启示，却没有激起你的赞叹，不仅赞叹那些领受天主圣神的先知，也赞叹那位预言的对象：祂进入生命的时刻早被许多先知清楚预言，以致整个犹太民族一直期待着祂的到来，在耶稣降临后彼此激烈争论；其中许多人承认了基督，相信祂就是预言之对象，而另一些人不相信祂，反而不屑于那些因耶稣教导而不愿引起丝毫骚动的人的温良，竟敢对耶稣施加宗徒们如此真实坦率地记载的暴行，并未从他们对祂的奇妙历史中秘密省略那些在众人看来会给基督教教义带来耻辱的事。但耶稣自己和祂的门徒都希望信徒不仅相信祂的天主性和奇迹，仿佛祂没有分享人性，没有取那\Quote{相反圣神}的人性肉身；他们也看到，那降入人性、进入人类苦难之中、并取了人灵魂和肉身的权能，通过信德连同其神圣要素，有助于信徒的救恩，当他们看到从祂开始了神性与人性的结合，为使人性通过与神性的共融而上升为神性，不仅在耶稣内，而且在所有不仅相信、而且活出耶稣教导的生活的人身上，这生活提升每一个按照耶稣诫命生活的人，使之与天主为友并共融。\par

% structure: p0061 source=h2 target=subsection reason=relative-heading-rank; base=h2

\subsection{第二十九章}

根据\Person{采尔苏斯}，\Person{阿波罗}曾愿意\Place{麦塔蓬托}人将\Person{阿里斯特阿斯}视为神明。然而，麦塔蓬托人认为证明阿里斯特阿斯是人的证据——并且很可能不是一个有德行的人——比神谕宣称他是神明或值得神圣尊荣的宣告更为有力，因此他们不听从阿波罗，结果无人将阿里斯特阿斯当作神明。但关于耶稣，我们愿意说：既然接受他为天主之子——取了人的灵魂和身体的天主——对人类有益，而这对喜爱身体之魔鬼的贪食欲以及那些因此将他们视为神明的人似乎无益，那么地上的魔鬼（那些未受魔鬼本性教导的人以为它们是神），连同它们的崇拜者，都渴望阻止耶稣教义的传播；因为他们看见他们贪恋的奠酒和香气正因耶稣训诲的盛行而被扫除。但派遣耶稣的天主挫败了魔鬼的一切阴谋，并使耶稣的福音在全世界得胜，以使人悔改和革新，并到处建立了教会，以对抗那些迷信、放荡和邪恶之人的集会；因为构成各城集会之市民的群众正是这种性格。而由基督训导的天主教会，若仔细与它们所在地区的集会相比，就像世上的明灯；因为谁会不承认，即使是教会中较差的成员，以及与较好者相比不太相称的人，仍然比许多属于各地区集会的人更为优秀呢？\par

% structure: p0063 source=h2 target=subsection reason=relative-heading-rank; base=h2

\subsection{第三十章}

因为天主的教会，例如雅典的教会，是一个温顺、坚定的团体，因为它渴望取悦那统治万有的天主；而雅典人的集会则喜爱纷争，完全不能与那城中的天主教会相比。你也可以对哥林多的天主的教会和哥林多人民的集会，以及亚历山大的天主的教会和亚历山大人民的集会，说同样的话。若听到此事的人是个坦诚的人，且愿意探究以弄清真相，他必会充满惊叹，因着那不仅构想出这计划，并且得以在各处，在各城的人民集会旁，建立天主的教会的那一位。同样，若将天主的教会中的议会与任何城市中的议会相比较，你会发现，教会中的某些议员堪当治理天主之城（如果全世间真有这样一座城的话）；而在其他地方的议员，其品格中毫无足以显示他们在表面上优于同城公民的那种优越性。同样，你也必须将每城中\Emph{教会}的\Emph{统治者}与\Emph{人民}的统治者相比较，以便看到，即使在天主教会中那些远未尽本分、比更活跃的人更懒散的议员和统治者中，仍然有可能在德行进步的各方面，发现他们普遍优于各城中议员和统治者的品格。\par

% structure: p0065 source=h2 target=subsection reason=relative-heading-rank; base=h2

\subsection{第三十一章}

如果事情如此，那么，对于能够成就如此巨大成果的耶稣，若说在\Emph{祂}内居有非凡的神性，为何不合情理呢？而此事既未发生在\Place{普罗科内斯}人\Person{阿里斯特阿斯}身上（尽管\Person{阿波罗}愿意他被视为神明），也未发生在\Person{采尔苏斯}枚举的其他个体身上，他说：\Quote{无人将那位拥有如箭离弦般被带走的奇能的\Place{海波尔波雷}人\Person{阿巴里斯}视为神明。}因为那位赐予这位\Place{海波尔波雷}人\Person{阿巴里斯}如箭般被带走奇能的神明，为何要赐予他这样的恩赐呢？是为了使人类因此得益，还是他自己因拥有这种能力而获得什么好处？——即使我们假设这些说法并非全然虚构，而是在某个魔鬼的配合下确实发生了。但如果记载说，我的耶稣被接到荣耀中，我在这行为中看到天主的安排，即：成就此事的天主，借此方式将这位导师推荐给那些目击者，以便他们作为那些争战的人——不是为人的教义，而是为神圣的教导——能尽己所能奉献给那统治万有的天主，并做一切事以讨祂喜悦，因为这些人要在天主的审判中，接受在此生所行之善或恶的报应。\par

% structure: p0067 source=h2 target=subsection reason=relative-heading-rank; base=h2

\subsection{第三十二章}

但是，既然塞耳苏斯接着提到了克拉佐美尼人的事例，并在关于此人的叙述后加上了这样的评论：\Quote{他们不是报告说，他的灵魂常常离开身体，以无形的方式四处游荡吗？然而人并不把他视为神。} 那么我们必须回答：大概是某些恶鬼策划了让这样的陈述被写下来（因为我不相信他们策划了这样的事竟然真的\Emph{发生}了），目的是为了使关于耶稣的预言和祂所说的话，要么被毁谤为与这些虚构类似，要么不引起惊讶，因为它们并不比其他事件更非凡。但是我的耶稣关于祂自己的灵魂（这灵魂离开身体不是由于任何人类的必然性，而是因为正是为了这个目的赋予祂的奇迹般的能力）说：\BibleQuote{没有人夺去我的性命，是我自己舍弃的；我有权舍弃，也有权再取回。} 因为祂有权舍弃，所以祂在说\BibleQuote{圣父啊，祢为什么舍弃我？} 并大声呼喊后，交付了灵魂，提前于那些为了不让刑罚延续而打断被钉者腿的刽子手。并且祂\Quote{取回祂的性命}，当祂向门徒们显现，曾在他们面前向不信的犹太人说\BibleQuote{拆毁这座圣殿，三天之内我要把它重建起来}，而\BibleQuote{祂这话是指祂身体的圣殿说的}；此外，先知们在他们的著作中多处预言了这样的结果，并且在这句\BibleQuote{我的肉身也要安居在希望中，因为祢必不将我的灵魂遗弃在地狱，也不让祢的圣者见到腐朽}。\par

% structure: p0069 source=h2 target=subsection reason=relative-heading-rank; base=h2

\subsection{第三十三章}

然而，塞耳苏斯显示他读过不少希腊历史，当他进一步引述关于\Place{阿斯提帕莱亚}的\Person{克莱奥梅德斯}的叙述时：\Quote{他}他讲述说\Quote{进入了一个方舟，虽然被关在里面，却没有被发现，而是由于某种神意的安排，当某些人切开方舟要捉拿他时，他飞了出去。} 现在，这个故事，如果是虚构的（正如它看起来那样），不能与关于耶稣的叙述相比，因为在这种人物的生平中，没有迹象表明他们拥有归于他们的神性；而耶稣的神性既由得救的教会存在所确立，也由关于祂的预言、因祂的名所行的医治、祂里面的智慧和知识，以及那些知道如何从简单的信德上升、并按照耶稣的诫命（祂说\BibleQuote{搜索圣经}）和保禄的意愿（他教导\BibleQuote{我们应该知道如何回答每个人}）探究神圣圣经含义的人所发现的更深真理所确立，还有那说\BibleQuote{要常常准备好，以温和的态度回答每个询问你们心中盼望理由的人}的人。然而，如果他想让人承认这不是虚构，就让他说明，这种超自然的力量通过某种神意的安排，使他从方舟中飞走的目的何在。因为如果他能提出任何值得考虑的理由，并指出赋予\Person{克莱奥梅德斯}这种力量符合天主任何值得敬拜的目的，我们将决定我们应该给出的回答；但如果他不能就此说出任何令人信服的话（显然是因为找不到任何理由），那么对于那些没有接受这个故事的人，我们要么轻蔑地谈论它，指责它是假的，要么我们会说，某种邪魔的力量迷惑了眼睛，在阿斯提帕莱亚人身上产生了像变戏法者所产生的结果，而塞耳苏斯认为，当他以神谕般的口吻说\Quote{他因某种神圣的安排从方舟飞走了}时，他对这个人是说得像神谕的。\par

% structure: p0071 source=h2 target=subsection reason=relative-heading-rank; base=h2

\subsection{第三十四章}

然而，我认为塞耳苏斯所知道的只有这些例子。而且，为了显得他自愿忽略其他类似事例，他说：\Quote{还可以举出许多其他同类例子。} 就算承认有过许多这样的人，他们对人类没有带来任何益处，那么他们每个人的行为，与耶稣的工作以及关于祂所行的奇迹（我们已经详细讨论过）相比，又能算得了什么呢？他接着想象，\Quote{在崇拜那位被俘虏并处死的人时，} 如\Emph{他}所说，\Quote{我们就像革特人崇拜\Person{扎摩尔克西斯}，基利家人崇拜\Person{摩普索斯}，阿卡纳尼亚人崇拜\Person{安菲洛科斯}，底比斯人崇拜\Person{安菲阿拉俄斯}，勒巴狄亚人崇拜\Person{特罗福尼俄斯}一样。} 在这些事例中，我们将证明他毫无根据地拿我们与上述相比。因为这些不同的部落为以上列举的个别人物建造了庙宇和雕像，而我们在敬拜神性时却避免了任何这样的方式（因为它们更适合于某类邪魔，这些邪魔固定在他们偏爱的特定地方，或者通过某些仪式和咒语，似乎被转移后定居下来），我们只对耶稣心怀敬畏，祂将我们的心思从一切可感的事物（不仅是可朽的，而且是注定要朽坏的）中召回，提升到用祈祷和正义的生活——这是我们献给祂的，作为介于非受造与一切受造物本性之间的中保——来敬拜掌管万有的天主，祂将来自圣父的恩惠赐予我们，并作为大司祭将我们的祈祷呈献给至高天主。\par

% structure: p0073 source=h2 target=subsection reason=relative-heading-rank; base=h2

\subsection{第三十五章}

但我愿意，针对那个出于某种未知理由而提出上述言论的人，以随便谈话的方式说一些如下似乎于他也不无恰当的言辞。那么，你所提到的那些人难道是无足轻重的吗？在\Place{勒巴狄亚}与\Person{特洛福尼奥斯}有关的地方，或是在\Place{底比斯}的\Person{安菲阿拉奥斯}神庙，或是在\Place{阿卡纳尼亚}的\Person{安菲洛科斯}，或是在\Place{基利家}的\Person{摩普索斯}，难道就没有能力吗？还是在这些人物中有某种存在，要么是一位\OriginalTerm{魔鬼}{demon}，要么是一位\OriginalTerm{英雄}{hero}，甚至是一位\OriginalTerm{神明}{god}，在运行超出人力的事？因为如果他回答说，这些个体身上并没有什么比其他人更魔鬼般的或更神圣的东西，那么让他立即表明他自己的看法，即作为一个\OriginalTerm{伊壁鸠鲁派}{Epicurean}，一个不与\Person{希腊人}持相同见解的人，既不承认\OriginalTerm{魔鬼}{demons}，也不像\Person{希腊人}那样敬拜\OriginalTerm{神明}{gods}；并且让他表明，他之前列举的例子（好像他相信它们是真的）以及他在后面几页中加上的例子，都是毫无目的地提出的。但是，如果他将断言所谈到的人物是\OriginalTerm{魔鬼}{demons}、\OriginalTerm{英雄}{heroes}，甚至是\OriginalTerm{神明}{gods}，那么让他注意，他将通过自己所承认的东西建立一个他不希望的结果，即\Person{耶稣}也是某种类似的存在；正因如此，他也能够向不少人证明，祂是从天主降临来探访人类的。如果他一旦承认了这一点，且看他是否将被迫承认，祂比那些他将其与祂并列的人物更强有力，因为后者没有一个是禁止向其他人献上尊崇的；而祂，因自信，因为祂比所有这些其他人物都更有能力，便禁止将后者当作神圣者来接纳，因为他们是邪恶的\OriginalTerm{魔鬼}{demons}，占据了地上的地方，由于无力上升到更纯洁、更神圣的区域，那些地上的粗俗之物和无数邪恶无法到达之处。\par

% structure: p0075 source=h2 target=subsection reason=relative-heading-rank; base=h2

\subsection{第三十六章}

但既然他接下来提到哈德良宠臣的案例（我指的是关于少年\Person{安提诺乌斯}的传说，以及\Place{埃及}的\Place{安提诺波利斯}居民给予他的尊崇），并想象给予他的尊崇几乎不亚于我们给予\Person{耶稣}的，那么让我们表明这一陈述出于怎样的敌意。因为生活于哈德良宠臣之中、甚至不回避非自然淫欲者的生活，与可敬的\Person{耶稣}的生活有什么共同之处呢？连那些提出无数其他指控、说了许多谎言的人，也不能指控祂甚至有一丝放荡的倾向。不仅如此，如果人本着真理和公正的精神研究关于\Person{安提诺乌斯}的传说，他会发现，由于埃及人的\OriginalTerm{巫术法术}{magical arts}和仪式，在他死后，在以其名字命名的城市里，甚至\Emph{外表}上他行了一些事——这种效果据说也被埃及人以及那些精通他们所行法术的人在其他神庙中产生。因为他们在某些地方设立宣称拥有预言或医治能力的\OriginalTerm{魔鬼}{demons}，这些\OriginalTerm{魔鬼}{demons}常常折磨那些在普通食物种类或触摸死人身体上似乎犯了错误的人，以便他们能够吓唬无知的群众。在\Place{埃及}的\Place{安提诺波利斯}中被认为是一位神的存在，便是这种性质，其（所谓的）德能是一些依靠从中获利而生活的人的谎言捏造的；而其他人，被安置在那里的\OriginalTerm{魔鬼}{demon}所欺骗，又因软弱的良心被定罪，实际上认为他们正在承受由\Person{安提诺乌斯}施加的神圣惩罚。他们所举行的奥秘以及他们发出的看似预言的东西，也是这种性质。\Person{耶稣}的那些却远非如此。因为并不是一群\OriginalTerm{巫士}{sorcerers}向国王或统治者献媚，听其命令，而似乎使他成神；而是宇宙的创造者自身，依照祂话语的惊人说服力，举荐祂为值得尊崇的，不仅向那些善意的人，也向\OriginalTerm{魔鬼}{demons}和其他看不见的权势，这些权势即使在现今也表明，他们要么惧怕\Person{耶稣}的名，将其视为一位能力更高者的名，要么恭敬地接纳祂作为他们合法的统治者。因为如果这举荐不是由天主给予祂，那么\OriginalTerm{魔鬼}{demons}就不会仅仅因提祂的名而顺从地离开那些他们所攻击的人。\par

% structure: p0077 source=h2 target=subsection reason=relative-heading-rank; base=h2

\subsection{第三十七章}

埃及人既被教导崇拜安提诺乌斯，若你将他与阿波罗或宙斯相比，他们会忍受这种比较，安提诺乌斯因被列于这些神祇之中而受到抬高；因为凯尔苏斯说\Quote{他们不能忍受将他与阿波罗或宙斯相比}，这显然被证明是谎言。然而基督徒（他们已学到，永生在于认识独一的真天主，即超越万有的那一位，以及祂所派遣的耶稣基督；他们也学到，所有外邦人的神都是贪婪的魔鬼，盘旋于祭献、血和其他的祭品周围，为欺骗那些没有投靠超越万有的天主的人；但天主的圣洁天使与地上所有魔鬼的本性和意志不同，只有极少数的、仔细而明智地研究过这些事的人才知道他们）不会容忍将他们与阿波罗或宙斯，或任何通过气味、血和祭品而受敬拜的存在相比较。他们中有些人出于极度的单纯，不能为自己的行为提出理由，但真诚地遵守所领受的诫命；另一些人则出于不可轻视的理由，甚至是深刻的理由，并且（如希腊人所说）出自事物的内在本质。在后一批人中，天主是经常谈论的主题，那些通过祂的独生圣言而获得分享祂神性（因此也分享祂的名号）而被天主尊荣的人也是常谈的话题。他们也谈论很多关于天主的天使以及那些反对真理但受了欺骗的人；由于受欺骗，他们称这些人为神、天主的使者、善灵或英雄，这些英雄因一个善良人灵魂的迁入而成为如此。这样的基督徒也将指明，正如在哲学中有许多人看似拥有真理，却或通过似是而非的论证欺骗了自己，或轻率地同意他人提出和发现的东西；同样，在那些脱离肉身的灵魂中，无论是天使还是魔鬼，有些被似是而非的理由所诱导而自称为神。因为这样的事的理由不可能被人完全精确地发现，所以被认为安全的是，任何人都不应将自己托付给任何存在如同天主，除了耶稣基督，祂就像是万有的统治者，既看见了这些重大奥秘，又将它们启示给少数人。\par

% structure: p0079 source=h2 target=subsection reason=relative-heading-rank; base=h2

\subsection{第三十八章}

那么，对安提诺乌斯或任何其他类似之人的信仰，无论在埃及人还是在希腊人中，都可以说是不幸的；而对耶稣的信仰则似乎要么是幸运的，要么是彻底研究的结果，对大众而言显得像前者，对极少数人则显得像后者。当我谈到某种信仰是（如大众所称的）不幸时，我将原因归于天主，因为祂知道进入人生命的每个人所分得的各种命运的原因。此外，希腊人会承认，即使是在那些被认为最具智慧的人中，好运也起了很大作用，例如在选择某类教师而非另一类教师时，在遇到更好的导师时（有些导师教授最相反的学说），以及在更好的环境中成长时；因为许多人在这样的环境中长大，以至于他们从未能接受任何更好事物的观念，而是从幼年起就一直过着要么作为放荡者或暴君宠儿的生活，要么处于某种其他悲惨状态中，使灵魂无法向上仰望。这些不同命运的原因，按一切可能性来看，都可在天主眷顾的旨意中找到，尽管人不易查明这些原因；但我之所以说这些，是出于对我主题主要部分的离题，因为那句谚语说：\Quote{这就是信仰的力量，因为它抓住了首先呈现的东西。}由于教育方法的不同，有必要谈论人信仰的差异，其中一些人在信仰上更幸运，另一些人则较不幸运；并由此进而表明，所谓的好运或厄运，即使在最有天赋的人身上，也会使他们显得更富有理性，并基于理性而对大多数人的意见表示赞同。然而关于这些就足够了。\par

% structure: p0081 source=h2 target=subsection reason=relative-heading-rank; base=h2

\subsection{第三十九章}

我们也必须注意凯尔苏斯接下来说的言论，他对我们说：\Quote{信德占据了我们的心灵，使我们对我们给予耶稣道理的赞同表示屈服；}因为的确，是信德产生了这样的赞同。然而，要注意，那信德本身是否不展现出值得称赞之处，因为我们把自己托付给那统御万有的天主，向那引导我们达到如此信德的祂表示我们的感恩，并宣称祂不可能在没有天主助佑的情况下尝试或完成这样的结果。我们也信赖福音作者的意图，观察到他们的虔诚和良心，这在他们著作中表现出来，其中没有任何虚假、欺骗或狡诈的内容；因为对我们来说，那些不熟悉希腊人狡诈诡辩所教导的技巧（这种技巧以极大的貌似合理和尖锐为特征）以及法庭上流行的修辞的灵魂，不可能如此编造出那些本身适合引导人走向信德和与信德相称的生活的事件。而且我认为，正因为如此，耶稣愿意使用这样的人作为祂道理的教师，即，这样就不会有任何怀疑似是而非的诡辩的根据，而是清楚地向所有能理解的人表明，作者们那毫无诡诈的意图，可以说是带着极大的单纯，被认为配得上神性的更大力量陪伴，这力量成就的远远超过似乎可能通过词语的迂回和句子的编织，并辅以希腊艺术的一切区分所能成就的。\par

% structure: p0083 source=h2 target=subsection reason=relative-heading-rank; base=h2

\subsection{第四十章}

但请注意，我们信德的原则与生来植入我们心灵的普遍观念是否和谐，是否对那些坦诚聆听其陈述的人产生改变；因为尽管一种对事物的歪曲看法，借助许多同样效果的教导，已经能够在群众的心中植入这样的信念：即像是神，以及由金银、象牙和石头所造的东西值得崇拜，但常识却禁止假设神是任何可朽的物质，或者当人按照某种祂外形的形象或象征赋予祂一个体现在死物质中的形式时，祂受到尊崇。因此我们立刻说像不是神，这些（艺术的）受造物不能与造物主相比，与那统御万有、创造、维持并管理宇宙的天主相比是渺小的。理性的灵魂，可以说，认识到它（与神性）的关系，立刻抛弃它一度以为是神的东西，重新恢复它对造物主的自然之爱；并且因为对祂的情感，它也接受了那通过祂所指定的门徒首先向万民提出这些真理的那一位，祂派遣这些门徒，赋予他们神性的能力和权柄，去宣讲关乎天主及其国度的道理。\par

% structure: p0085 source=h2 target=subsection reason=relative-heading-rank; base=h2

\subsection{第四十一章}

但既然他指控我们（不知已经多少次了）\Quote{把这耶稣——不过是一具可朽身体——当作天主，并以为我们这样做是虔诚的，}那么再对此回答就是多余的，因为在前文中已经说了很多。然而，让那些提出这种指控的人明白：我们所尊崇并相信从起初就是天主和天主子的是那位圣言本身、智慧本身和真理本身；至于祂的可朽身体及其所含有的灵魂，我们断言，不仅通过与祂的共融，更通过结合与混合，它们接受了最高的能力，并且在分受了祂的神性之后，被改变为天主。如果有人对我们关于祂身体的这种说法感到困难，就让他注意希腊人关于质料所说的，严格地说，质料没有性质，接受造物主所愿意赋予的，并且常常脱去它以前的性质，而接受另一种不同且更高的性质。如果这些意见正确，那么，耶稣身体的这种可朽性质，如果天主的眷顾如此愿意，被改变为一种以太般的、神性的性质，又有什么可奇异的呢？\par

% structure: p0087 source=h2 target=subsection reason=relative-heading-rank; base=h2

\subsection{第四十二章}

塞尔苏斯并非良好的辩论者，当他将耶稣终有一死的肉身与金、银、石头相比，声称前者比后者更易朽坏时。因为，正确地说，不朽之物并不比另一不朽之物更不易朽坏，可朽之物也不比另一可朽之物更易朽坏。但是，假设有朽坏程度的不同，我们可以回答说：如果潜在于一切性质之下的质料能够交换其中某些性质，那么耶稣的肉身为何就不能交换性质，成为适合居于以太和其上天界的身体，不再具有属于肉体的软弱以及塞尔苏斯称为\Quote{不洁}（他用这词时不像一位哲学家）的那些特性呢？因为真正不洁的是因邪恶而不洁。身体的本性并非不洁；就其为身体的本性而言，它并不拥有作为不洁之源的恶习。但是，既然他猜到了我们会作的回答，他就针对耶稣身体的改变说：\Quote{好了，在他脱去这些性质之后，他将成为一位神；}（若如此），那么为什么不是埃斯科拉庇俄斯、狄俄尼索斯和赫拉克勒斯呢？对此我们回答说：\Quote{埃斯科拉庇俄斯、狄俄尼索斯或赫拉克勒斯行了什么伟大的事迹？}他们又能指出哪些人是在言行上得到改善、变得更好，从而可以证明他们为神呢？因为让我们阅读许多关于他们的记载，看看他们是否免于放荡、不义、愚蠢或怯懦。如果在他们身上找不到这类事，那么塞尔苏斯的论证可能就有分量，它将上述诸人与耶稣等同。但如果可以肯定，尽管他们被记载了一些可敬的事，然而仍被记录做了无数违反理性的事，那么你怎能再有任何理由说这些人抛弃了有死的身体后成为神，更胜于耶稣呢？\par

% structure: p0089 source=h2 target=subsection reason=relative-heading-rank; base=h2

\subsection{第四十三章}

接着他论到我们说：\Quote{我们嘲笑那些敬拜朱庇特的人，因为他的坟墓在克里特岛被指出；然而我们却敬拜那从坟墓中复活的那一位，尽管不知道克里特人遵守这一习俗的理由。}现在请看，他如此为克里特人、为朱庇特和他的坟墓辩护，暗指那些寓言式的观念，据称有关朱庇特的神话就是依照这些观念编造的；而他攻击我们，我们承认我们的耶稣确实被埋葬了，但坚持他也从坟墓中\Emph{复活}了——克里特人尚未就朱庇特提出这一陈述。但既然他似乎承认朱庇特的坟墓在克里特，当他说\Quote{我们不知道克里特人遵守这一习俗的理由}时，我们回答说：昔兰尼的卡利马科斯读过无数诗歌作品和几乎全部希腊历史，并不知晓有关朱庇特及其坟墓的故事中包含任何寓意；因此他在献给朱庇特的颂诗中指责克里特人说：——\par

\Quote{克里特人常说谎：因为你的坟墓，哦君王，\newline{}克里特人建造了；然而你并未死去，\newline{}因为你永远活着。}\par

那说\Quote{你并未死去，因为你永远活着}的人，虽否认朱庇特的坟墓在克里特，却仍然记载朱庇特身上有死亡的开端。但在地上出生就是死亡的开端。他的话是：——\par

\Quote{而瑞亚在帕拉西亚人中生下了你；}\par

然而他本应看到，在因坟墓而否认朱庇特生于克里特之后，他生于阿卡迪亚就完全与他生而必死相合。卡利马科斯谈及这些事的方式如下：\Quote{哦朱庇特，有人说你生在伊达山上，\newline{}有人说在阿卡迪亚。\newline{}哪一个是说谎者，哦父亲？\newline{}克里特人常说谎，}等等。是塞尔苏斯迫使我们讨论这些题目，他以不公正的方式对待耶稣，同意关于他死亡和埋葬的记载，却将他的复活视为虚构，尽管这不仅被无数先知预言，而且有诸多证据表明他死后曾显现。\par

% structure: p0095 source=h2 target=subsection reason=relative-heading-rank; base=h2

\subsection{第四十四章}

在这些论点之后，塞尔苏斯引用了一些反对耶稣教训的异议，这些异议出自极少数被认为是基督徒的人，依他看来，并非较有智慧者，而是较无知的一类，并断言\Quote{以下是他们制定的规条：凡受过教育、或聪明、或谨慎的人都不可到我们这里来（因为这些品质被我们视为邪恶）；但若有任何无知、不聪明、不受教育或愚昧的人，让他们放心前来。通过这些话语，他们承认这样的人配得上他们的天主，从而明显表明他们只渴望并且能够争取愚笨、卑贱、愚蠢的人，以及妇女和儿童。}对此我们回答说，正如耶稣教导节制，说\BibleQuote{凡注视妇女，有意贪恋她的，他已在心里奸淫了她}，若有人看见少数被称为基督徒的人过着放荡的生活，他最公正地责备他们生活违背耶稣的教导，但若将他们的可指责行为归咎于福音，则最不合理；同样，若他仍发现基督徒的教训邀请人走向智慧，那么责备就应归于那些停留在自己无知中的人，他们所说的并非塞尔苏斯所述（因为尽管他们中有些简单无知，但他们不像他声称的那样无耻），而是其他严重性小得多的事情，但这些事情却使人们转离智慧之道。\par

% structure: p0097 source=h2 target=subsection reason=relative-heading-rank; base=h2

\subsection{第四十五章}

但基督信仰的目标是使我们成为智慧的人，这不仅可以从我们也在使用的古代犹太著作中证明，而且尤其可以从耶稣时代以后撰写的、在教会中被视为神圣的那些著作中证明。现在，在\BibleBook{}{圣咏}第五十篇中，\Person{达味}在他的祈祷中向天主这样说：\BibleQuote{祢将祢智慧的隐密和奥秘之事启示了我。}\Person{撒罗满}也因求智慧而获得了智慧；如果有人细读\BibleBook{}{圣咏}，他会发现该书充满了许多智慧格言；他的智慧的证据可见于他的著作，这些著作以简短的词句包含了大量的智慧，在其中你会发现许多对智慧的赞美和获得智慧的鼓励。\Person{撒罗满}如此智慧，以至于\BibleQuote{\Person{舍巴女王}听见了撒罗满的名声和主的名声，就来用难题试探他，对他倾诉她心里所有的一切；撒罗满回答了她所有的问题。君王没有省略任何问题而不回答她。\Person{舍巴女王}看见了撒罗满的一切智慧和所拥有的财产，便神不守舍。她对君王说：我在本国所听到关于你和你的智慧的传闻，的确是真的；我没有相信那些告诉我的人，直到我来了，亲眼看见了。看，他们连一半都没有告诉我。你添加的智慧和财产超过了我所听到的传闻。}关于他，也记载说：\BibleQuote{天主赐给\Person{撒罗满}极大的智慧和明悟，以及像海边的沙那样宽阔的心胸。\Person{撒罗满}的智慧超过所有古代人和所有埃及智者的智慧；他比所有的人更智慧，比则辣黑人\Person{厄堂}、玛曷耳的儿子\Person{赫曼}、\Person{加耳苛耳}、\Person{达尔达}更智慧。他在四周所有的民族中都很有名。\Person{撒罗满}说了三千句箴言，他的诗歌有五千首。他讲论树木，从黎巴嫩的香柏直到墙上长出的牛膝草；也讲论鱼和走兽。各地的人都来听\Person{撒罗满}的智慧，并来自地上听过他智慧名声的众君王。}\par

福音渴望在信徒中有智慧的人，以至于为了锻炼听众的理解力，它以谜语、所谓的\Quote{隐晦}言语、比喻和问题说出某些真理。一位先知——\Person{欧瑟亚}——在他的预言结尾说：\BibleQuote{谁是智慧的，他会明白这些事？谁是明智的，他会知道它们？}此外，\Person{达内尔}和他的同伴们在巴比伦君王周围的智者所培养的学问上取得了如此大的进步，以至于显示出他们比所有人优秀十倍。在\BibleBook{}{厄则克耳}中，对那因自己的智慧而极为自傲的\Place{提洛}君王说：\BibleQuote{你比\Person{达内尔}更智慧吗？没有秘密向你不启示的吗？}\par

% structure: p0100 source=h2 target=subsection reason=relative-heading-rank; base=h2

\subsection{第四十六章}

如果你来看耶稣时代之后写的书，你会发现那些听比喻的众多信徒，好像是在\Quote{外面}，只配得外来的道理，而门徒则私下学习比喻的解释。因为耶稣私下向祂自己的门徒启示一切，认为那些渴望知道祂智慧的人比群众更尊贵。祂应许那些信祂的人，要派遣智者和经师给他们，说：\BibleQuote{看，我派遣智者与经师到你们这里来，其中有些你们要杀死，钉在十字架上。}\Person{保禄}也在天主赐予的\OriginalTerm{神恩}{charismata}列表中，将\Quote{智慧的言语}置于首位，其次（作为较低者）是\Quote{知识的言语}，第三位（更低的是）\Quote{信德}。因为他认为\Quote{言语}高于奇迹能力，所以他将\Quote{奇迹的运作}和\Quote{治病的恩赐}置于言语恩赐较低的位置。在\BibleBook{宗徒}{大事录}中，\Person{斯德望}为\Person{梅瑟}的大学问作证，这完全是从古代著作中获得的，不为大众所知。因为他说：\BibleQuote{\Person{梅瑟}学习了埃及人的一切智慧。}因此，关于他的奇迹，有人怀疑他行奇迹也许不是因为他声称来自天主，而是凭借他精通的埃及知识。因为\Person{法郎}怀着这样的怀疑，召集了埃及的术士、智者和行邪术的人，但这些人在\Person{梅瑟}的智慧面前毫无用处，因为\Person{梅瑟}的智慧胜过了埃及人的一切智慧。\par

% structure: p0102 source=h2 target=subsection reason=relative-heading-rank; base=h2

\subsection{第四十七章}

但很可能，保禄在\BibleBook{格林多}{前书}中，针对那些以希腊智慧为傲的希腊人所说的话，已使有些人相信福音并非以赢得智者为目的。现在，凡持此见解的人，当明白福音在责备恶人时，论到他们说：他们并非在关于理智、且不可见与永恒的事上成为智者；而只是在忙于感觉事物、并将这些视为至关重要时，才成为属世的智者。因为存在众多意见，有些支持物质与形体的主张，断言一切具有实质性存在的都是形体的，此外别无他物，无论称之为不可见或无形体，福音也称这些构成属世的智慧，它消亡衰败，只属于今世；而那些将灵魂从今世之事提升至与天主同在的福乐及其王国，并教导人轻视一切感觉和可见之物（因其只是暂时的），而奔向不可见之事、注视那些未见之事的意见，福音称这些构成天主的智慧。但保禄，作为真理的热爱者，论及某些希腊智者（当他们的说法为真时）说：\BibleQuote{他们虽然认识天主，却不以祂为天主而光荣祂，也不感谢祂}。他为他们认识天主作证，也说这事并非没有天主的许可临到他们，他说：\BibleQuote{因为天主将这事显明了给他们}；我想这隐约暗示那些从感觉上升到理智的人，当他加上：\BibleQuote{自从创世以来，天主那看不见的事，即祂永远的大能和神性，藉着所造之物就可以明白看见，以致他们无可推诿；因为虽然他们认识天主，却不以祂为天主而光荣祂，也不感谢祂}。\par

% structure: p0104 source=h2 target=subsection reason=relative-heading-rank; base=h2

\subsection{第四十八章}

或许也有人因这些话——\BibleQuote{弟兄们，你们看看你们的蒙召，按肉体有智慧的不多，有能力的不多，有尊贵的也不多。但是天主偏偏选了世上愚拙的，来使有智慧的羞愧；又选了世上卑贱的、被人轻视的、以及那无有的，为要废掉那有的，使一切有血气的，在天主面前一个也不能自夸}——而以为没有受过教育、有智慧或审慎的人会接受福音。现在，为回答这样的人，我们要说：经上并未说\BibleQuote{\Emph{没有}一个按肉体的有智慧的人}，而是说\BibleQuote{按肉体的有智慧的人\Emph{不多}。}此外，显而易见，在被称为\Quote{主教}的人应有的资格中，保禄在描述主教应为何种人时，列出一个条件：他也应善于教导，说：他应当能驳倒反对的人，用他里面的智慧堵住愚妄和欺骗之人的口。正如他选择那些只结过一次婚的、而非再婚的人为主教，选择无可指责的而非有可责之处的，选择节制的而非不节制的，选择审慎的而非不审慎的，选择行为端正的而非即使有轻微不端之处的人，同样，他愿意那些优先被选为主教职的人应善于教导，并能驳倒反对的人。既然如此，刻尔苏斯如何能公正地指责我们说\Quote{不要有任何受过教育的、有智慧的或审慎的人来到我们这里}？相反，凡愿意的，无论是\Quote{受过教育的}、\Quote{有智慧的}、\Quote{审慎的}，都请来到我们这里；同样，如果有人无知、不敏、未受教育或愚笨，也请他来吧：因为福音应许要医治这些前来的人，使他们全堪当得起天主。\par

% structure: p0106 source=h2 target=subsection reason=relative-heading-rank; base=h2

\subsection{第四十九章}

以下言论也不真实，即所谓\Quote{只有愚笨和卑贱的人、缺乏知觉的人、奴隶、妇女和儿童，才是天主圣言的教师想要赢得的人。}福音确实邀请这样的人，为使他们变得更好；但它也邀请其他与此截然不同的人，因为基督是众人的救主，更是信徒的救主，无论他们是聪明还是淳朴；并且\BibleQuote{祂为我们作了赎罪祭，不仅为我们的罪，也为普世的罪。}在此之后，我们无需再费心回答刻尔苏斯下面这类话：\Quote{为什么受过教育、研究过最好的意见、既拥有智慧之实又拥有其貌，是件坏事？这对认识天主有何妨碍？难道这不是一种帮助，反而更能使人达到真理吗？}确实，受过教育并非坏事，因为教育是通向德行的道路；但将持错误见解的人列入受教育者之列，就连希腊的智者也不会这样做。另一方面，谁不承认研究最好的意见是福气？但什么是最好的意见呢？除了那些真实的、能激励人修德行之意见外，还有别的吗？此外，一个人\Emph{是}有智慧的，而非如刻尔苏斯所说的只是\Emph{显得}有智慧，这是一件极好的事。受过教育、研究过最好的意见、并有智慧，非但不妨碍认识天主，反而是一种帮助。这话由我们说来，比由刻尔苏斯说来更为恰当，尤其如果证明他是一个伊壁鸠鲁派。\par

% structure: p0108 source=h2 target=subsection reason=relative-heading-rank; base=h2

\subsection{第五十章}

但让我们看看他在接下来的话中说了什么：\Quote{不仅如此，我们确实看到，那些在市场上表演最可耻的把戏并聚集人群的人，决不会接近智者的集会，也不敢在他们面前炫耀自己的把戏；但无论他们在哪里看到年轻人、一群奴隶和一群无知的人，他们就会挤进去，炫耀自己。} 现在，请看他是如何用这些话诽谤我们，将我们比作那些在市场上表演最可耻的把戏并聚集人群的人！请问，我们表演了什么可耻的把戏？或者\Emph{我们的}行为中有什么与他们的相似，因为我们通过诵读和解释所读的内容，引导人敬拜宇宙的天主，并修养相应的德行，使他们远离轻慢神明和一切违背正确理性的事？的确，哲学家们愿意召集那些听他们宣讲德行的听众——尤其是某些犬儒学派的人，他们公开与所遇之人交谈。那么，他们会坚持认为，这些不聚集那些被认为受过教育的人，而是从街上邀请和聚集听众的人，与那些在市场上表演最可耻的把戏并聚集人群的人相似吗？然而，无论是刻尔苏斯，还是任何持相同意见的人，都不会责备那些按照他们所认为的仁爱精神向无知的民众发表言论的人。\par

% structure: p0110 source=h2 target=subsection reason=relative-heading-rank; base=h2

\subsection{第五十一章}

如果他们这样做不应受责备，那么让我们看看基督徒是否比他们更大、更好地劝勉众人行善。因为公开演讲的哲学家不挑选听众，而是愿意听的人站着听。而基督徒，在事先尽可能检验那些希望成为他们听众之人的灵魂，并事先私下教导他们之后，当他们（在进入团体之前）似乎已经充分表明了对美德生活的渴望时，才引入他们，并且只在这之后才引入，私下形成一类初学接受但尚未获得完全净化标记的人；另一类是那些尽力表明其意向，除了基督徒认可的事物外别无他求的人；在这些当中，有一些被指定的人负责调查加入者的生活和行为，以便阻止那些犯下丑行的人进入他们的公开聚会，而完全不同的人则全心全意地接纳他们，以便日复一日地使他们变得更好。他们对那些犯罪的人，尤其是过着放荡生活的人，也是如此行事的，他们将这些人逐出团体，尽管按刻尔苏斯的说法，他们就像那些在市场上表演最可耻把戏的人。然而，可敬的毕达哥拉斯学派曾为那些背弃他们哲学体系的人建造衣冠冢，视他们为死人；而基督徒则哀叹那些被放荡或任何其他罪恶所征服的人为死人，因为他们已迷失且于天主已死，并视他们如从死中复活（若他们表现出合宜的改变），然后在稍后的时间接纳他们，比起初接纳的人间隔更大，但不会在教会中授予任何职位或高位给那些在公开宣认福音后又堕落跌倒的人。\par

% structure: p0112 source=h2 target=subsection reason=relative-heading-rank; base=h2

\subsection{第五十二章}

现在，关于刻尔苏斯的以下陈述：\Quote{我们也看到那些在市场上表演最可耻的把戏并聚集人群的人，} 看看是否没有明显的虚假，以及比较了毫无相似之处的事物。他说，这些被他比作我们的人，即那些\Quote{在市场上表演最可耻的把戏并聚集人群的人，决不会接近智者的集会，也不敢在他们面前炫耀他们的把戏；但无论他们在哪里看到年轻人、一群奴隶和一群愚蠢的人，他们就会挤进去炫耀自己。} 他这样说，无非就是像街头女人一样用辱骂诬蔑我们，她们的目的是互相诽谤；因为我们尽一切努力使我们的聚会由智者组成，我们当中那些特别卓越和神圣的事物，只有在拥有大量有智慧的听众时，我们才敢在讨论中公开提出，而当我们看到听众是较单纯的人，需要那种比喻性地称为\InnerQuote{奶}的教导时，我们就隐藏和默默略过更深奥的真理。\par

% structure: p0114 source=h2 target=subsection reason=relative-heading-rank; base=h2

\subsection{第五十三章}

因为我们的\Person{保禄}在写给\Person{格林多人}的信中用到了这句话——格林多人是希腊人，他们的道德尚未净化——：\BibleQuote{我用奶喂养了你们，而不是肉；因为那时你们还不能承受，就是现在也不能，因为你们仍是属血肉的：因为在你们中间有嫉妒和纷争，你们岂不是属血肉的，并且照着人的样子行事吗？}同一位作者知道有一种更适合灵魂的滋养，并且那些被接纳的年轻人在饮食上被比作奶，他接着说：\BibleQuote{你们成了需要奶，而不是强壮食物的人。因为凡是用奶的，在正义的话语上都是不熟练的，因为他是婴孩。但强壮的食物是属于成年人的，就是那些因习惯而感官受到操练，能辨别善恶的人。}那么，那些认为这些话是善说的人，会相信我们信仰的高贵教义从不会在智者的集会中被提及，反而每当（我们的导师）看到年轻人、一群奴隶和一伙愚蠢的人时，就会公开提出神圣而可敬的真理，并在这些人面前炫耀自己谈论这些真理吗？但仔细检查我们全部著作精神的人都很清楚，\Person{塞爾蘇斯}被一种类似无知民众的仇恨所激动，不假思索地说出这些谎言。\par

% structure: p0116 source=h2 target=subsection reason=relative-heading-rank; base=h2

\subsection{第五十四章}

我们承认，尽管\Person{塞爾蘇斯}不愿如此，但我们确实渴望以天主之言教导所有的人，以便给予年轻人适宜的劝勉，并向奴隶们展示如何恢复思想的自由，并因圣言而变得高贵。在我们当中，那些担任基督宗教使者的人充分宣告，他们欠了希腊人和化外人、智者和愚人的债（因为他们不否认甚至有义务治愈愚人的灵魂），以便他们能尽可能摆脱无知，并努力获得更大的明智，同时也聆听\Person{所罗门}的言语：\BibleQuote{哦，愚昧人哪，要有一颗明悟的心}，以及\BibleQuote{你们中最愚笨的人，让他转向我}；智慧也用这些话劝勉那些缺乏理解的人：\BibleQuote{来，吃我的饼，喝我为你调和的酒。抛弃愚昧，好使你存活，并在知识中修正理解。}对于这一点，我也要说（因为它关系到要点），以回应\Person{塞爾蘇斯}的陈述：难道哲学家们不邀请年轻人听他们的讲座吗？难道他们不鼓励年轻人以更好的生活取代邪恶的生活吗？难道他们不希望奴隶学习哲学吗？那么，我们是否必须责备那些曾劝勉奴隶修德的哲学家呢？责备\Person{毕达哥拉斯}曾这样对待\Person{札莫尔克西斯}，责备\Person{芝诺}对待\Person{珀耳修斯}，责备那些最近鼓励\Person{爱比克泰德}研究哲学的人？哦，希腊人哪，你们确实可以召唤青年、奴隶和愚昧的人来研究哲学，但如果我们这样做，我们就不是出于博爱的动机，希望用理性的药物治愈每一个理性的本性，并将他们带入与造物主天主的共融之中吗？这些评论足以回应\Person{塞爾蘇斯}的诽谤——与其说是控告，不如说是诽谤。\par

% structure: p0118 source=h2 target=subsection reason=relative-heading-rank; base=h2

\subsection{第五十五章}

但由于\Person{塞爾蘇斯}喜欢堆砌对我们的诽谤，而且在他已经说过的之外又添加了其他，让我们也审视这些，看看是基督徒还是\Person{塞爾蘇斯}有理由为所说的话感到羞耻。他声称：\Quote{我们确实看到，在私人住宅里，羊毛工人、皮革工人、漂布工，以及最未受教育、最粗鲁的人，在他们年长和更聪明的主人面前不敢说一句话；但当他私下找到孩子们，以及某些像他们一样无知的妇女时，他们便倾泻出美妙的言论，大意是：他们不应该听从父亲和老师，而应该听从他们；说前者是愚蠢和迟钝的，不知道也不能做任何真正好的事情，被空洞的琐事所占据；说\Emph{他们}独自知道人应该怎样生活，如果孩子们服从他们，孩子们自己就会幸福，也会使家庭幸福。当他们这样说时，如果看到一位青年导师、或一个更聪明的人、甚至父亲本人走近，他们中比较胆怯的就会害怕，而比较大胆的则怂恿孩子们挣脱束缚，低声说：在父亲和老师面前，他们既不愿意也不能向他们解释任何好事情，因为这些人厌恶地回避那些人的愚蠢和迟钝，认为他们完全败坏，在邪恶上已经走得很远，而且会惩罚他们；但如果他们想要（利用他们的帮助），就必须离开父亲和老师，和妇女以及他们的玩伴一起去妇女的居室，或去皮革店，或去漂布店，这样才能达到完美；——通过这样的话，他们赢得了这些孩子。}\par

% structure: p0120 source=h2 target=subsection reason=relative-heading-rank; base=h2

\subsection{第五十六章}

现在请注意，他是如何借着这类陈述来贬低我们当中那些宣讲圣言的人，这些人竭尽全力将灵魂提升至万有的造物主，并教导我们应当轻视\Quote{可感觉的}、\Quote{暂时的}、\Quote{可见的}事物，尽己所能达到与天主的共融、领悟\Quote{可理解的}和\Quote{不可见的}事物，以及与天主及其友人们共享的永生福乐；他竟然将这些人比作\Quote{私家作坊里的羊毛工、皮匠、漂洗工，以及最粗鲁的乡下人，这些人处心积虑地把少年男孩引向邪恶，并鼓动妇女离弃父亲和老师，跟随他们}。现在，让刻尔苏斯指出，我们使孩童和妇女远离了哪一位智慧的父亲或哪一位老师；并让他透过比较那些信奉我们教义的孩童和妇女，看看他们从前所听过的见解是否比我们的更好，以及我们究竟以何种方式将孩童和妇女从高尚可敬的学问中拉开，诱导他们走向更坏的事物。然而，他绝不可能对我们提出任何此类指控，因为恰恰相反，我们使妇女远离放荡的生活、远离与她们所生活之人的不和、远离对剧场和舞蹈的一切狂热欲望、远离迷信；而我们则训练正值青春萌动、对性欲感到渴望的男童养成节制的习惯，不仅向他们指出那些罪恶所伴随的耻辱，也指出恶人的灵魂因这类行为所陷入的境地、它将遭受的审判以及将要承受的惩罚。\par

% structure: p0122 source=h2 target=subsection reason=relative-heading-rank; base=h2

\subsection{第五十七章}

但是，那些我们称之为无聊者和愚人、而刻尔苏斯却为之辩护的教师，仿佛他们是传授更好道理的教师，究竟是谁呢？（我不知道，）除非他认为那些引诱妇女走向迷信和猥亵表演的人是妇女的良师而非无聊者，以及那些领导和引诱年轻人做出我们所知的种种放荡行为的人才是并非缺乏理智的教师。我们确实也尽可能地召唤这些人脱离哲学的教条，归向我们对天主的崇拜，通过展示其卓越与纯洁。但是，既然刻尔苏斯通过他的言论声称我们并非如此行，而是只召唤愚人，我要对他说：\Quote{如果你指控我们使那些已经专注于哲学研究的人放弃哲学，你虽未说真话，但你的指控还有几分貌似合理；然而，你现在说我们将信徒从良师那里拉开，那么请指出，除了哲学教师或被任命讲授某些有用学问的人之外，还有其他什么良师？}\par

然而，他无法指出任何这样的人；而我们则公开地、毫不隐秘地应许：\Emph{他们}——那些按照天主圣言生活、凡事仰望天主、在一切所行之事上都如同在天主面前一般的人——将是有福的。难道这是羊毛工、皮匠、漂洗工和未受过教育的乡下人的教导吗？这样的断言他无法证明。\par

% structure: p0125 source=h2 target=subsection reason=relative-heading-rank; base=h2

\subsection{第五十八章}

但是，那些在刻尔苏斯看来类似私家作坊里的羊毛工、皮匠、漂洗工和未受过教育的乡下人的人，他断言，在父亲或老师面前，他们将不愿开口，或无法向男童解释任何有益的事。对此，我们要说：我亲爱的先生，你指的是什么样的父亲和什么样的老师呢？如果你指的是一个赞赏美德、厌恶邪恶、欢迎善事的人，那么要知道，我们将以最大的胆量向孩子们宣讲我们的见解，因为我们在这样的审判者面前将获得好评。但是，如果在一个憎恨美德与善良的父亲面前，以及在那些教导与纯正道理相悖的人的面前，我们保持沉默，请不要为此责备我们，因为你责备我们是没有理由的。无论如何，在那些父亲认为哲学奥秘对于他们的儿子以及一般年轻人是无聊无益之事的情况下，你在传授哲学时，就不会在无价值的父母面前透露其中的秘密；而你，想要把那些转向哲学研究的恶父母的儿子们区分开来，会遵守适当的时机，以便哲学的道理能触及年轻人的心灵。我们对我们的教师也持同样的说法。因为，如果我们把（听众）从那些教授猥亵喜剧、放荡的抑扬格诗歌以及许多其他既无益于说话者也无益于听者（因为后者不知道如何以哲学的心态聆听诗歌，前者也不知道如何对每个年轻人说出有益的话）的教导者那里拉开，那么遵循这样的做法，我们并不羞于承认我们所行的。但是，如果你能向我指出训练年轻人从事哲学并使他们操练于哲学的教师，我就不会从这样的人那里拉开年轻人，反而会努力提升他们，如同那些先前已修习了整个学艺范围和哲学科目的人，使他们达到那寓于多数基督徒之中的崇高、可敬的雄辩高峰，在那里讨论最重要的问题，并证明和显示出这些问题已由天主的先知和耶稣的宗徒以哲学的方式处理过。\par

% structure: p0127 source=h2 target=subsection reason=relative-heading-rank; base=h2

\subsection{第五十九章}

随即，\Person{策尔苏斯}察觉自己过于刻薄地毁谤了我们，仿佛是为辩护，便表述如下：\Quote{任何人都可从以下话语看出，我所提出的指控并不比真理所强迫的更重。那些邀请人参与其他秘仪的人，如此宣告：\InnerQuote{凡手洁而舌慎的人}；另一些则如此说：\InnerQuote{凡从一切污秽中纯净、灵魂自认无恶、且生活良善公义的人}。这就是那些应许净化罪过的人所发的通告。但让我们听听这些基督徒邀请的是何种人。他们说：凡罪人、凡愚昧无知者、凡幼童，并且一般说来，凡不幸者，天主的国都要接纳。那么，你岂不称不义者、窃贼、入室行窃者、下毒者、亵渎圣物者、盗墓者为人吗？若有人为盗贼集会发布通告，他还会邀请其他人吗？}现针对这些言论，我们回答说：邀请那些\Emph{灵魂患病}的人来\Emph{治愈}，与邀请那些\Emph{健康}的人来\Emph{认识与研究}神圣事物，并非同一回事。然而，我们兼顾这两点：首先邀请众人来获得医治，劝诫罪人归向教导人不犯罪的道理，劝诫愚昧无知者归向产生智慧的道理，劝诫幼童在思想上成长到成人，劝诫单纯的不幸者归向好运，或更恰当的说法——归向真福。而当那些转向德行的人有了进步，表明自己已因圣言而净化，并已尽可能度更善的生活时，我们才邀请他们参与我们的奥迹，而非在此之前。\Quote{因为我们是在成全的人中讲智慧。}\par

% structure: p0129 source=h2 target=subsection reason=relative-heading-rank; base=h2

\subsection{第六十章}

再者，正如我们所教导的：\Quote{智慧不会进入邪恶的灵魂，也不居住于罪身内}\BibleRef{智 1:4}，我们说：凡手洁的人，因此举起圣洁的手向天主，因专注于崇高与天上的事物，能说：\Quote{我举起我的手，如同晚祭}，愿他来我们这里；凡通过日夜默思主的法律而拥有智慧口舌的人，并因\Quote{习惯使他的官能受到锻炼，能分辨善恶}，愿他不迟疑地来接受那适应于在虔敬与一切德行上操练之人的有力而理性的滋养。既然天主的恩宠与所有怀着纯爱爱慕不朽教义导师的人同在，凡不仅远离一切玷污，而且远离被视为较小过犯的人，让他勇敢地领受耶稣的奥迹，这些奥迹本只启示给圣洁与纯净的人。策尔苏斯的领受奥迹者于是说：\Quote{让那灵魂自认无恶的人来。}但按照耶稣的训诫执行入门礼的人，会对那些心中已得净化的人说：\Quote{让那灵魂长久以来自认无恶的人，尤其是他自从把自己交托给圣言的医治以来，听耶稣私下对他的真门徒所讲的道。}因此，他在希腊秘仪的入门者与耶稣教义的教师之间所做的比较中，他不知道邀请恶人来得医治与引介已得净化者进入神圣奥迹之间的区别！\par

% structure: p0131 source=h2 target=subsection reason=relative-heading-rank; base=h2

\subsection{第六十一章}

我们并非邀请邪恶的人、窃贼、入室行窃者、下毒者、亵渎圣物者、盗墓者以及所有那些策尔苏斯以其夸张风格列举的其他人类参与\Emph{奥迹}、进入那隐藏在奥迹中的\Emph{智慧共融}——这智慧是天主在万世以前为祂圣徒的荣耀所预定的——而是邀请这类人来\Emph{得医治}。因为在圣言的神性中，有一些有助于治愈患者的帮助，关于这些帮助，圣言说：\Quote{健康的人不需要医生，有病的人才需要；}此外，还有一些帮助向灵魂与身体纯洁的人显示\Quote{奥迹的启示，这奥迹自世界起始以来一直隐藏，如今却借着先知们的圣经显明出来，并借着我们的主耶稣基督的显现}，这\Quote{显现}向每位成全的人显示，并光照理性获得事物的真知。但由于他夸大了对我们的指控，在他列举完那些卑劣之徒后，加了这一句：\Quote{盗贼还会通过通告召唤什么人到自己身边？}我们回答这个问题说：盗贼召集这类人到他周围，是为了利用他们的恶行去对付那些他们想杀戮和掠夺的人。相反，基督徒即使邀请盗贼所邀请的人，也是邀请他们到一种完全不同的召叫：即借祂的圣言包扎这些伤口，将取自圣言的药物——类似于酒、油、膏药和医学技术中其他治疗用具——应用到在罪恶中溃烂的灵魂上。\par

% structure: p0133 source=h2 target=subsection reason=relative-heading-rank; base=h2

\subsection{第六十二章}

其次，他诋毁那些劝告生活邪恶者悔改并改过迁善的讲论与著作，声称我们说\Quote{天主被派遣到罪人那里}。他的这一说法，无异于责怪某些人说：因为城中有病人，一位极仁慈的君王便派遣了一位医生给他们。的确，天主圣言作为医生被派遣到罪人那里，但对于那些已经洁净且不再犯罪的人，则是作为神圣奥秘的教师。但凯尔苏斯无法看出这一区别——因他无意以热爱真理的精神来行动——便说：\Quote{为什么祂不被派遣到那些无罪的人那里？不犯罪有什么不好？}我们回答：如果所谓\Quote{无罪的人}是指那些不再犯罪的人，那么我们的救主耶稣也被派遣到这样的人那里，但不是作为医生。而如果所谓\Quote{无罪的人}是指那些从未犯过罪的人——因为他在陈述中没有加以区分——我们回答：一个人不可能如此无罪。我们说这话时，当然除去了那在基督耶稣内的人，祂\BibleQuote{没有犯过罪}。凯尔苏斯确实怀着恶意说我们声称\Quote{天主将接纳那不义的人，若他因自己的邪恶而自卑，却不接纳那义人，即使他从起初就以德性（仰望）天主。}现在我们断言：一个人不可能从起初就以德性仰望天主。因为邪恶必然首先存在于人中。正如保禄也说：\BibleQuote{当诫命来到时，罪复活了，我死了。}此外，关于那不义的人，我们并不教导说：他因自己的邪恶而自卑就足以被天主接纳；而是说：如果他因过去的行为而自责，并因此谦卑地行走，且以合宜的方式生活于将来，天主就会接纳他。\par

% structure: p0135 source=h2 target=subsection reason=relative-heading-rank; base=h2

\subsection{第63章}

此后，他不理解\BibleQuote{凡高举自己的，必被贬抑}这句经文，也不理解（尽管连柏拉图也教导）\Quote{良善且德性的人谦卑且有序地行走}，而且他还不知道我们诫命说\BibleQuote{所以，你们要谦卑自下，服在天主大能的手下，这样到了时候，祂必高举你们}；他便说：\Quote{那些适当主持审判的人，会叫在他们面前哀哭自己恶行的人停止可怜的哭号，以免他们的判决更多出于怜悯而非真实；而天主却不按真实审判，而是按阿谀奉承。}那么，圣经中有什么阿谀和可怜哭号的话呢？当罪人在祈祷中向天主说：\BibleQuote{我承认了我的罪，没有隐瞒我的过犯。我说：我要向主告明我的罪过}等等？因为，他能证明这种做法不适合那些在祈祷中谦卑自下于天主手中的罪人悔改吗？而且，他在指控我们的努力中变得困惑，自相矛盾：一方面似乎认识一个\Quote{无罪}的人和\Quote{一个义人，他从起初就以德性（装饰）仰望天主}；另一方面又接受我们的说法：没有全然义的人，或没有无罪的人；因为，仿佛他承认其真实性，便说：\Quote{这确实显然是真实的：人类天性倾向于犯罪。}其次，仿佛不是所有的人都受到圣言的邀请，他说：\Quote{所有的人都应无差别地被邀请，因为所有人确实都是罪人。}然而，在前面几页中，我们已经指出耶稣的话：\BibleQuote{凡劳苦和负重担的，你们\Emph{全都}到我这里来，我要使你们安息。}\Emph{所有}的人，因此，因罪性而劳苦且负重担的，都被邀请到天主的言语中所说的安息，因为\BibleQuote{天主发出了祂的言语，医治了他们，拯救他们脱离了毁灭。}\par

% structure: p0137 source=h2 target=subsection reason=relative-heading-rank; base=h2

\subsection{第64章}

但既然他还说\Quote{罪人比他人更受偏爱是什么意思？}并说了其他类似的话，我们必须回答：绝对没有一个罪人比不是罪人的人更受偏爱；但是有时一个罪人，认识到自己的罪，并因此悔改，因自己的罪而谦卑，他比一个被认为罪过较小的人更受偏爱，而那个人却不认为自己有罪，反倒因自己认为的某些好品质而自高，并因此沾沾自喜。这一点，凡愿意以公平精神阅读福音书的人，都可以从那个税吏的比喻中清楚看到，税吏说\BibleQuote{天主，可怜我这个罪人吧！}而法利塞人却怀着某种邪恶的自负夸口说：\BibleQuote{我感谢祢，因为我不像其他的人，勒索、不义、奸淫，也不像这个税吏。}因为耶稣在叙述他们二人后加上了这句话：\BibleQuote{这个人下去回家，成了正义的，而不是那个；因为凡高举自己的，必被贬抑；凡谦卑自己的，必被高举。}因此，当我们教导说：每个人，无论他是谁，面对天主的伟大时，都意识到人性的软弱，我们必须从那唯一能弥补我们欠缺者那里，向我们（必死的）本性所缺乏的祈求时，我们并没有亵渎天主，也没有说谎。\par

% structure: p0139 source=h2 target=subsection reason=relative-heading-rank; base=h2

\subsection{第65章}

但他却想象，我们发出这些劝言是为了使罪人悔改，是因为我们无法赢得任何真正良善和正义的人，因此才向最不圣洁和堕落的人敞开大门。然而，如果有人公正地观察我们的集会，我们能够呈现更多从并非极为邪恶的生活中悔改的人，而非那些犯下最可憎罪行的人。因为那些自知有更好事物的人，自然渴望天主关于义人赏报的应许是真实的，从而比那些生活极其邪恶的人更轻易地赞同圣经的陈述，后者因自知有罪而拒绝承认他们将受到审判万民者以适合如此大罪人的惩罚，且审判万民者不会违反正理施加惩罚。有时，当极其堕落的人因基于悔改的希望而愿意接受未来惩罚的教义时，他们却被犯罪的习惯所阻止，仿佛不断浸染并染上了邪恶，不再拥有轻易转向正当生活并合乎正理的生活的能力。尽管\Person{刻尔苏斯}观察到这一点，但他不知为何，仍以下列方式表达自己：\Quote{然而，确实，每个人都清楚，没有人能通过惩罚，更不用说通过仁慈的处理，彻底改变那些因本性和习惯而犯罪的人，因为改变本性是极其困难的事。但那些无罪的人享有更好的生活。}\par

% structure: p0141 source=h2 target=subsection reason=relative-heading-rank; base=h2

\subsection{第六十六章}

现在，我认为\Person{刻尔苏斯}在此犯了一个大错误，他拒绝承认那些因本性和习惯而犯罪的人有彻底转变的可能性，声称他们甚至无法通过惩罚得到治愈。因为很明显，所有人都有犯罪的本性倾向，有些人不仅出于本性，而且出于习惯，但并非所有人都无法彻底转变。因为在每一个哲学派别和天主的圣言中，都发现有人经历了如此巨大的转变，以至于他们可以作为卓越生活的典范。在英雄时代的名字中，有人提到\Person{赫拉克勒斯}和\Person{尤利西斯}，在后来的时代，\Person{苏格拉底}，以及在最近生活的人中，\Person{穆索尼乌斯}。因此，\Person{刻尔苏斯}不仅对我们发出诽谤，说\Quote{每个人都清楚，那些因本性和习惯而犯罪的人，无论如何——即使通过惩罚——也无法彻底变好，}而且也对哲学中最崇高的名字发出诽谤，这些哲学家并未否认人恢复德性是有可能的。尽管他没有精确表达他的意思，但我们仍将更善意地理解他的话，并证明他的推理不健全。因为他的话是：\Quote{那些因本性和习惯而倾向于犯罪的人，没有人能通过惩罚彻底改造他们；}我们按理解反驳了他的话，并尽力而为。\par

% structure: p0143 source=h2 target=subsection reason=relative-heading-rank; base=h2

\subsection{第六十七章}

然而，他可能想表达这样的意思：那些既因本性又因习惯而犯下最为堕落之人所犯之罪的人，即使通过惩罚也无法彻底转变。然而，这一点从某些哲学家的历史中可以看出是错误的。因为谁不会把那个以某种方式屈服于他的主人、将他置于妓院、任由任何人玷污的人视为最为堕落的人之一呢？这样的故事正与\Person{斐多}有关！又有谁不会同意，那个带着吹笛手和一伙狂欢的放荡同伴闯入可敬的\Person{色诺克拉底}的学园、侮辱一位备受朋友钦佩的人，不是人类中最恶劣的恶徒之一呢？然而，尽管如此，理性仍有足够的力量使他们悔改，并使他们在哲学上取得如此进步，以至于其中一人被\Person{柏拉图}认为值得叙述\Person{苏格拉底}关于不朽的言论，并记录他在狱中的坚定，当时他藐视毒芹，以无所畏惧和平静的心态探讨如此众多而重要的话题，即使是那些全神贯注且不受干扰的人也很难跟上；而另一方面，\Person{波勒蒙}从一个放荡者变成了最节制生活的人，接替了以可敬品格著称的\Person{色诺克拉底}在学园中的位置。因此，\Person{刻尔苏斯}说\Quote{因本性和习惯而犯罪的人，即使通过惩罚也无法彻底改造}，这话不真实。\par

% structure: p0145 source=h2 target=subsection reason=relative-heading-rank; base=h2

\subsection{第六十八章}

凡以条理井然、文辞优雅著称的哲学论辩，在方才列举的那些人以及曾行恶事的其他人身上产生如此结果，这丝毫不足为奇。然而，当我们看到克尔苏斯称之为\Quote{粗俗}的这些言论充满了能力，仿佛咒语一般，它们立时使众多人从放荡不羁的生活转向极端有节，从邪恶生活转向更善，从怯懦或柔弱转向如此高尚的勇气，以至于使人为内心所彰显的虔诚而藐视死亡本身，我们岂不应当公正地钦佩它们所蕴含的能力吗？因为那些起初担任（基督徒）使徒职分、辛劳建立天主教会之人的言语——甚至他们的宣讲——也伴随着一种劝服的力量，虽然不像那些自诩为柏拉图哲学或任何其他纯属人的哲学之信徒的言论那样，后者除了人性的特质外别无所有。然而，伴随耶稣宗徒话语的明证来自天主，并由圣神和能力所证实。因此，\Emph{他们的}话语迅速而飞快地传播，或者更确切地说，\Emph{天主的}圣言藉着他们为工具，转变了许多生来及习惯上都是罪人的人，这些人无人能以惩罚改造，却被那按照己意塑造并转变他们的圣言所改变。\par

% structure: p0147 source=h2 target=subsection reason=relative-heading-rank; base=h2

\subsection{第六十九章}

克尔苏斯照其惯常方式继续断言：\Quote{彻底改变本性极其困难。}然而，我们深知每个理性灵魂只有一种本性，并主张没有一个是万有创造者所造为邪恶的，但许多人因教育、邪僻榜样及周围影响而\Emph{变成}了邪恶，以致邪恶在某些人身上成为了第二天性。我们确信，天主的圣言改变那以邪恶为第二天性的本性，不仅并非不可能，甚至不是十分困难的事，只要人相信他必须将自己交付于万有的天主，并行事全为取悦祂，因为在祂那里不可能\par

善恶同享一份尊荣，\newline{}或懒惰者与勤劳者\newline{}同等消亡。\par

但即便改变某些人极其困难，其原因也应当归于他们自身的意志，这意志不愿接受相信万有的天主是今生一切行为的公义审判者。因为审慎的选择和修习，对于成就那些看似极其困难——夸张地说，几乎不可能——的事情大有助益。当人渴望在剧场中央凌空绷紧的绳索上行走，同时携带众多重物时，他难道不是通过修习和专注而成就了如此技艺吗？而当人渴望按照德行生活时，尽管从前可能极其邪恶，他难道就发现这是不可能的吗？试问，持这种观点的人，难道不是在谴责理性动物创造者的本性，而非谴责被造物本身——如果祂使人性有能力成就如此困难且毫无用处之事，却使其无力获得自身真福？但这些话足以作为对\Quote{彻底改变本性极其困难}这一断言的答复。其次，他断言：\Quote{无罪之人分享更美好的生命。}却没有澄清他所谓\Quote{无罪的}是指从生命之初就是如此的人，还是指因转变而成的人。从生命之初就是如此的人，不可能有任何一位；而转变（心灵）后成为这样的人，则为数甚少，即那些在归依救恩之言后成为这样的人。他们归依时并非如此。因为若非圣言的帮助——而且是圆满的圣言——人不可能成为无罪之人。\par

% structure: p0151 source=h2 target=subsection reason=relative-heading-rank; base=h2

\subsection{第七十章}

其次，他反对那似乎为我们所持的断言：\Quote{天主将能行一切事。}他甚至在此也不明白这些话的用意，以及其中所包含的\Quote{\Emph{一切事}}指什么，以及为何说天主\Quote{将能}。现在不必谈论这些事；尽管他本可以有理由反对这一命题，但他并未这样做。或许他不理解可以貌似合理地用来反对它的论证，或者即使理解，他也看到了可能的反驳。按我们的判断，天主能行一切祂在不停止为天主、为善、为智的前提下所能行的事。但克尔苏斯——不明白\Quote{天主能行一切事}这句话的含义——断言：\Quote{祂不愿行任何邪恶的事，}承认祂有行恶的\Emph{能力}，但无\Emph{意愿}。反之，我们认为，正如那本性上因其内在的甜性而使其他事物变甜的东西不能产生与其特有本性相反的苦味，又如那本性上因其光明而产生光的东西不能造成黑暗，同样，天主也不能行邪恶，因为行恶的能力与祂的神性和全能相悖。而如果现存物中有人因本性倾向邪恶而能行恶，那是因为其本性没有不行恶的能力。\par

% structure: p0153 source=h2 target=subsection reason=relative-heading-rank; base=h2

\subsection{第七十一章}

接着，他假设了一件更理性的信徒们尚未认可的事，但也许一些缺乏智慧的人认为那是真的——即：\Quote{天主如同那些被怜悯所征服的人一样，自己也被征服，因怜悯恶人的哀号而减轻他们的痛苦，却抛弃了不做这类事的善人，这是最大的不义。}然而，依我们的判断，天主不减轻任何不转向德善生活的恶人的痛苦，也不抛弃任何已是善人的人，也不仅仅因为某人发出哀叹就减轻他的痛苦，或对他施以怜悯——这里使用的\Quote{怜悯}是较常见的含义。但那些因自己的罪恶而严厉谴责自己，并因此哀叹、痛悔自己为迷途者（就其先前行为而言），且表现出令人满意的转变的人，天主因他们的悔改而接纳他们，将他们视为已从大恶生活中转化的人。因为德善住进这些人的灵魂，驱逐先前占据他们的邪恶，使他们忘却过去。即使德善未能完全进入，但只要灵魂中有相当大的进步，也足以按进步的程度驱除和摧毁邪恶的洪流，使之几乎不再存留于灵魂中。\par

% structure: p0155 source=h2 target=subsection reason=relative-heading-rank; base=h2

\subsection{第七十二章}

其次，他仿佛以我们教义的导师身份发言，如此表达：\Quote{智慧之人拒绝我们的言论，被他们的智慧引入歧途、陷入网罗。}对此我们回复：既然智慧是关于神性与人性及其原因的知识，或者按天主的圣言所定义的，是\BibleQuote{天主的德能之气，全能者荣耀的纯粹流露；永远光明的辉耀，天主德能的无瑕明镜，以及祂美善的肖像}\BibleRef{智 7:25}，那么，真正有智慧的人绝不会拒绝一位通晓基督教原则的基督徒所说的话，也不会被引入歧途或陷入网罗。因为真正的智慧不会误导，只有无知才会误导；而在现存事物中，唯有知识是持久的，以及从智慧而来的真理。但如果与智慧的定义相反，你把任何以诡辩意见武断立论的人都称为智慧之人，那么，按照\Emph{你}所称呼的\Quote{智慧}，这样的人会拒绝天主的言语，被似是而非的诡辩误导并陷入网罗。而且，按照我们的教义，智慧不是对恶的知识，而对恶的知识（可以说）存在于那些持错误观点并被其欺骗的人身上，因此，我会称这样的人为无知而非智慧。\par

% structure: p0157 source=h2 target=subsection reason=relative-heading-rank; base=h2

\subsection{第七十三章}

此后，他再次诽谤基督信仰的使者，宣称他讲述\Quote{荒谬之事}，尽管他没有指出或清楚说明他所称的\Quote{荒谬}之事是什么。在他的诽谤中，他说：\Quote{没有智慧的人相信福音，因为被众多追随者所驱赶。}他这样做就像一个说：由于众多无知的人服从法律，没有智慧的人会服从例如\Person{Solon}、或\Person{Lycurgus}、或\Person{Zaleucus}、或任何其他立法者的人，尤其是如果他将\Quote{智慧之人}理解为（按德善生活而）有智慧的人。因为，正如这些立法者根据他们的效用观念，使这些无知之人被适当的引导和法律所包围一样，天主通过耶稣基督为世上众人立法，甚至将那些不按智慧方式生活的人也带到祂面前，以可能的方式引导他们走向更好的生活。天主深知此事，正如我们在前几页已指出的，在\Person{梅瑟}的经卷中说：\BibleQuote{他们以那不是天主的激起我的妒火，以他们的偶像惹我发怒；我也要以那不是子民的激起他们的妒火，以愚昧的民族惹他们发怒。}\Person{保禄}也知道此事，说：\BibleQuote{但天主拣选了世上愚拙的，叫有智慧的羞愧。}他一般称所有那些似乎在知识上有长进、却陷入无神论的多神崇拜的人为有智慧的，因为\BibleQuote{他们自称聪明，反成了愚拙，将不朽坏天主的荣耀变为偶像，仿佛必朽坏的人、飞禽、走兽和爬虫的样式。}\par

% structure: p0159 source=h2 target=subsection reason=relative-heading-rank; base=h2

\subsection{第七十四章}

他更进一步指责基督信仰的教师\Quote{追求无智之人}。对此我们要问：你所说的\Quote{无智之人}是指谁？因为确切地说，每一个恶人都是\Quote{无智之人}。如果你所说的\Quote{无智之人}就是指恶人，那么你在引导人追求哲学时，是设法赢得恶人还是贤人？但贤人是不可能获得的，因为他们已经献身于哲学了。那么你（设法赢得）的是恶人；但如果他们是恶人，他们就是\Quote{无智之人}吗？你这样许多人便被你争取到哲学那里，所以你也是在追求\Quote{无智之人}。如果我确实追求那些被如此称为\Quote{无智之人}的人，那么我的行为就像一位仁慈的医生，他寻求病人以便帮助和治愈他们。然而，如果你所说的\Quote{无智之人}是指那些不聪明的人，即智力低下的人，那么我的回答是：我也竭尽全力改善这样的人，尽管我不愿意用这样的材料来建造基督徒团体。因为我更偏爱那些较为聪明敏锐的人，因为他们能够理解难懂的语句，以及律法、先知和福音中那些表达隐晦的段落——这些段落你蔑视了，认为它们不值得注意，因为你没有弄清它们所包含的意思，也没有尝试进入作者的意图。\par

% structure: p0161 source=h2 target=subsection reason=relative-heading-rank; base=h2

\subsection{第七十五章}

然而，他后来又说：\Quote{基督信仰的教师就像一个许诺使病人恢复身体健康、却阻止他们咨询有技能的医生（以免他的无知被揭露）的人}，对此我们应回应：\Quote{你所指的是哪些医生？你指的是我们使无知之人远离的那些人？因为你并不认为我们劝勉那些投身哲学的人接受福音，所以你认为后者就是那些医生，而我们却使那些被邀请来接受天主圣言的人远离他们？}他当然不会回答，因为他提不出医生的名字；否则他必然求助于那些无知的医生，他们自动奴性地屈服于多神的崇拜以及无知之人所持有的其他任何意见。无论哪种情形，都能表明他在论证中徒劳地使用了\Quote{阻止他人接近有技能的医生}的比喻。但是，如果是为了使那些受欺骗的人远离伊壁鸠鲁的哲学以及被视为其体系的医生，那么我们这样做岂不是极其合理——使这些人远离由刻尔苏斯的医生所带来的危险疾病，即导致否定眷顾并以快乐为善的疾病？不过，姑且承认我们确实使那些我们鼓励成为门徒的人远离了其他哲学医生——例如逍遥学派，他们否认眷顾的存在以及神性对人的关系——那么，我们为何不能以虔诚的方式训练和医治那些被鼓励的人，说服他们献身于万有的天主，并使那些服从我们的人从那些被视为哲学家的话语所造成的重大创伤中解脱出来？不仅如此，再承认我们使人们远离斯多葛派学说的医生，他们引入一个可朽坏的天主，断言天主的本质是身体，能够在各部分中被改变和变化，并且主张万物终将消亡，唯独天主存留；那么，我们为何不能如此解放我们的受教者脱离邪恶，并用虔诚的论据引导他们献身于造物主，并赞美基督徒体系的圣父——祂如此安排了最仁慈的、适合灵魂皈依的训导，使之传遍全人类？不仅如此，如果我们治愈了那些因接受灵魂转世说而陷入愚蠢的人——这种说法来自某些医生的教导，他们认为理性本性有时下降进入各种无理性的动物，有时进入不能运用想象的状态——那么我们为何不能通过一种教义来改善我们受教者的灵魂？这种教义并不教导说，恶人受到的惩罚是一种无知或无理性的状态，而是表明天主加于恶人的劳苦和惩戒是一种导向皈依的药物。因为那些有智慧的基督徒，怀着这样的观点，对待头脑简单的人，就像父母对待非常幼小的孩子一样。因此，我们并不求助于年轻人或愚昧的乡下人，对他们说：\Quote{逃避医生吧。}我们也不说：\Quote{你们要小心，不要有人掌握知识。}我们也不断言\Quote{知识是邪恶的。}我们也不至于疯狂到说\Quote{知识使人失去健全的心智。}我们甚至不会说任何人因智慧而丧亡；尽管我们进行训导，但我们从不说\Quote{要听从我}，而是说\Quote{要听从万有的天主，以及耶稣，祂是有关天主之训导的授予者。}我们当中也没有人狂妄到说出刻尔苏斯借我们的教师之口对熟人所说的那句话：\Quote{只有我能拯救你们。}看看他所说的反对我们的谎言！此外，我们并\Emph{不}断言\Quote{真正的医生毁灭了他们许诺要治愈的人。}\par

% structure: p0163 source=h2 target=subsection reason=relative-heading-rank; base=h2

\subsection{第七十六章}

他接着提出了第二个对我们不利的比喻，说：\Quote{我们的教师就好像一个醉酒的人，走进一伙醉汉当中，却指责清醒的人喝醉了。}但是，请他从\Person{保禄}的书信中说明，耶稣的宗徒是否曾纵酒，以及他的言语是否不是清醒之言；或者从\Person{若望}的书信中说明，他的思想是否没有散发出节制和远离邪恶之醉意的精神。因此，没有哪个心智健全、教授基督信仰道理的人会醉酒；但\Person{凯尔苏斯}怀着一种与哲学家截然不同的心态对我们说出这些诽谤之词。此外，请\Person{凯尔苏斯}指出，那些基督信仰的传播者所指责的\Quote{清醒}之人究竟是谁。因为在我们看来，所有向无生命之物祈祷如同向天主祈祷的人都是沉醉的。我为何说\Quote{沉醉}？\Quote{疯狂}是更合适的词，用来形容那些急忙赶往神庙，崇拜雕像或将动物视为神祇的人。那些认为由卑劣且常常是极恶之人所制作的雕像能给真正的神祇带来任何尊荣的人，其疯狂程度也不亚于前者。\par

% structure: p0165 source=h2 target=subsection reason=relative-heading-rank; base=h2

\subsection{第七十七章}

接着，他把我们的教师比作一位患眼炎的人，把他的门徒也比作患同样疾病的人，并说：\Quote{在患眼炎的人群中，这样一个人会指责那些视力敏锐的人是瞎子。}那么，我们要问，哦，希腊人啊，在我们看来，除了那些不能从世界的无比伟大及其内容、从受造物的美丽中举目仰望，不能看出他们应当唯独崇拜、赞美、敬畏那位创造这些事物的天主，并且认为任何由人设计、用于尊崇天主的事物，不论是指向造物主还是不指向造物主，都不应被尊崇的人之外，还有谁是有不见的呢？因为，拿那些本不应与之相比较的事物，去与那超越一切受造存在的无限卓越相比，乃是那些理智昏暗之人的行为。因此，我们不是说那些视力敏锐的人患了眼炎或瞎了；而是断言那些因不认识天主而委身于神庙、雕像以及所谓的圣节的人，他们的心智被蒙蔽了，尤其当他们在不敬之外还过着放荡的生活，甚至不追求任何高尚的工作，反而做尽一切可耻之事时。\par

% structure: p0167 source=h2 target=subsection reason=relative-heading-rank; base=h2

\subsection{第七十八章}

在提出如此严重的指控之后，他想让人觉得，尽管他还有其他指控可以提出，他却默不作声地略过了。他的话如下：\Quote{我必须对他们提出这些指控以及其他类似性质的指控，不必一一列举；我断言他们是错误的，他们对天主行为傲慢，目的是用空洞的希望引导邪恶之人，并说服他们轻视更好的事物，告诉他们如果戒除那些事物，对他们会更好。}对此可以回答说，从那些皈依基督信仰者身上显出的力量来看，被福音赢得的绝非\Quote{邪恶}之人，而是较为单纯的一类人，也是许多人所说的\Quote{粗鲁}之人。因为这些人，由于害怕所威胁的惩罚，这些惩罚激发并劝诫他们避免那些会招致惩罚的行为，他们努力将自己交付给基督宗教，受到圣言的力量影响如此之深，以至于因害怕圣言中所说的\Quote{永罚}，他们蔑视世人发明用以折磨他们的一切酷刑——甚至死亡本身以及无数其他灾难——没有一个智者会说这是邪恶之人的行为。邪恶之人怎能实践节制、清醒，或仁爱、宽宏呢？不，就连对天主的敬畏，这样的人也无法感受到；由于敬畏对众人有益，福音鼓励那些尚不能为了事物本身而选择它的人，把它当作最大的祝福、超乎一切应许的恩赐来选择；因为这种原则无法植入那个宁愿生活在邪恶中的人心中。\par

% structure: p0169 source=h2 target=subsection reason=relative-heading-rank; base=h2

\subsection{第七十九章}

如果有人在这些事情上认为，在那些相信圣言的人群中表现出来的是迷信而非邪恶，并且指责我们的道理使人迷信，那么我们要回答他说：正如某位立法者曾回答一个人问他是否为自己的公民制定了最好的法律，说他并没有给他们绝对最好的，而是他们所能接受的最好的；同样，基督信仰道理之父可以说，我给了众人所能接受的最好的法律和教导来改善道德，不是用想象的辛劳和惩罚来威胁罪人，而是用真实的、必要的惩罚，用于纠正那些抵抗的人，尽管他们完全不理解施加惩罚者的目的，也不理解这些辛劳的效果。因为惩罚的道理既有益处，又符合真理，并且以隐晦的词语表述是有利的。此外，基督信仰的传播者所赢得的多半不是邪恶之人，我们也不侮辱天主。因为关于祂，我们既讲述真理，也讲述那些在众人看来显而易见、但对他们来说并不像少数以哲学方式探究福音真理的人那样清楚的事。\par

% structure: p0171 source=h2 target=subsection reason=relative-heading-rank; base=h2

\subsection{第八十章}

然而，既然塞耳苏声称\Quote{我们藉着虚幻的希望来争取基督徒}，我们就这样答复他，因为他批评我们关于永福生活以及与天主相通的教义：至于您，好先生，那些接受了\Person{毕达哥拉斯}和\Person{柏拉图}关于灵魂的学说的人，也同样是藉着虚幻的希望被争取的，他们认为灵魂的本性是升到穹苍之上，在超天空间观赏那些有福的观众所见到的景象。按照您的说法，啊，\Person{塞耳苏}，那些接受了灵魂（死后）持续存在的学说，并过着一种使他们成为英雄、与诸神同住的生活的人，也是藉着虚幻的希望被争取的。大概\Person{塞耳苏}也会说，那些相信灵魂是从外进入（身体）、并且将从死亡权势下被释放的人，是被空洞的希望所争取的。那么，让他出来辩论吧，不再隐瞒他所属的学派，而是承认自己是一名\OriginalTerm{伊壁鸠鲁派}{Epicurean}，让他面对那些在希腊人和外邦人中被非轻率提出的关于灵魂不朽、或死后持续、或思想原则不朽的论证；让他证明这些言论是用空洞的希望欺骗那些接受它们的人；而他的哲学体系的追随者则纯净地没有空洞的希望，并且他们确实导向善的希望，或者——更符合他意见的是——由于灵魂（死后）立即完全消灭，根本产生不出任何希望。除非，也许\Person{塞耳苏}和伊壁鸠鲁派会否认他们对自己终向——快乐——所抱持的希望是虚幻的，因为他们认为快乐是至善，在于身体的永久健康，以及\Person{伊壁鸠鲁}所抱持的关于它的希望。\par

% structure: p0173 source=h2 target=subsection reason=relative-heading-rank; base=h2

\subsection{第八十一章}

并且不要认为我接受那些哲学家关于灵魂不朽或死后持续的观点来反驳\Person{塞耳苏}，是不符合基督宗教的；因为，与他们持有某些共同见解，我们将更方便地确立我们的立场，即未来的永福生活仅属于那些接受了按照耶稣而来的宗教，并对万物的创造者怀有纯洁、真诚、不掺杂任何受造物的敬爱之心的人。并且，让任何愿意的人指出，我们说服人轻视的所谓\Quote{更好的事物}是什么，并让我们将在基督内与天主结合的永福终向——即圣言、智慧和一切德性——按照我们的观点，这将是藉天主的恩赐，赏赐给那些度过纯洁无瑕生活、并对万有的天主怀有单一而无分割的爱的人——与各哲学派别（无论是希腊还是外邦的）的教导，或宗教奥秘的宣称所随从的终向相比较；并让他证明任何其他派别所预言的终向优于我们所应许的，因此，那个才是真实的，而我们的终向不符合天主的恩赐，也不配那些善度生活的人；或者，让他证明这些话并非由充满圣先知灵魂的圣神所说。并且，让任何愿意的人指出，那些被所有人承认为出于人的话，优于那些被证明为神圣且受灵感而发的话。而所谓的\Quote{更好的}事物是什么，我们教导接受它们的人最好戒绝呢？因为，如果这样说并不傲慢，那么显而易见，没有什么比将自己交托给万有的天主，并顺服于那提升我们超越一切受造物、并通过有生命且生活的圣言——他也是生活的智慧和天主子——引领我们到超越万有的天主面前的教义更美好的事了。然而，既然我们答复\Person{塞耳苏}著作的第三卷已经延伸至足够长度，我们在此将结束当前的论述，而在后续中，我们将处理\Person{塞耳苏}随后所写的内容。\par

\SourceRecord{Translated by Frederick Crombie.；Ante-Nicene Fathers；Vol. 4.；Edited by Alexander Roberts, James Donaldson, and A. Cleveland Coxe.；Buffalo, NY: Christian Literature Publishing Co.,；1885.；<http://www.newadvent.org/fathers/04163.htm>.}

% structure: p0001 source=h1 target=section reason=page-title

\section{\WorkTitle{驳斥塞耳苏斯，卷四}}

% structure: p0002 source=h2 target=subsection reason=relative-heading-rank; base=h2

\subsection{第一章}

在前三卷中，我们已充分陈述了针对塞耳苏斯论著的答复，如今，可敬的\Person{盎博罗削}，我们通过基督向天主祈祷，将本卷作为对以下内容的答复呈上。我们祈祷能获得言语，正如\BibleBook{耶肋米亚}{书}所载，主对先知说：\BibleQuote{看，我将我的话放在你口中，如同火。看，我今天立你于万民和列国之上，为要拔除、拆毁、毁灭、推翻、建造和栽种。}因为我们现在需要言语，能从每个受伤的灵魂中拔除塞耳苏斯这本论著对真理发出的辱骂，或类似他意见的言论。我们也需要思想，能拆毁所有基于错误意见的建筑，特别是塞耳苏斯在他的作品中所建的那座，类似那些说\Quote{来吧，让我们建造一座城和一座塔，塔顶通天}的人所建的建筑。是的，我们甚至需要一种智慧，能推倒所有高抬自己反对天主知识的障碍，特别是塞耳苏斯反对我们的那种傲慢的高峰。接下来，因为我们在拔除和拆毁刚才提到的障碍之后，不能止步，而必须在被拔除的地方栽种\Quote{天主的田园}的植物，在被拆毁的地方建起天主的建筑和祂荣耀的圣殿——所以我们为此也当向主祈祷，祂赐予了\BibleBook{耶肋米亚}{书}上命名的恩赐，使我们也能获得适合建造基督（圣殿）和栽种属神法律及其相关预言言语的言语。最重要的是，必须证明，针对塞耳苏斯在已有断言之后的论述，关于基督的预言都是真实的预言。因为他同时装备自己反对两方——一方是犹太人，否认基督的来临已经发生，却期望将来；另一方是基督徒，承认耶稣就是预言所说的基督——他作了如下的陈述：——\par

% structure: p0004 source=h2 target=subsection reason=relative-heading-rank; base=h2

\subsection{第二章}

\Quote{但是，某些基督徒和（所有）犹太人竟主张，前者说某位神或某位神之子\Emph{已经}降临地上，后者说祂\Emph{将要}降临，并使大地居民成义，这是最无耻的断言，其驳斥不需多言。}此处他似乎正确地评论，不是\Quote{某些}犹太人，而是\Emph{全体}犹太人，都想象有某位（神）将降临地上；关于基督徒，则是他们中的\Emph{某些}人说祂已经降临。因为他指的是那些从犹太圣经证明基督的来临已经发生的人，他似乎知道有些异端派别否认基督耶稣被先知所预言。然而在前文中，我们已经尽己所能讨论了基督是预言主题的问题，因此为避免重复，我们不重复本可在此论述的许多内容。现在注意到，如果他希望以一种表面的力量颠覆对预言著作的信心，无论关于基督未来或过去的来临，他本应提出我们基督徒和犹太人在彼此讨论中引用的预言本身。因为这样他就会显得使那些被预言陈述的似是而非的性质（如他所认为的那样）所吸引的人，偏离同意它们的真理，也偏离因为这些预言而相信耶稣是基督。但如今，他无法回答关于基督的预言，或者根本不知道哪些关于祂的预言，便没有提出任何预言宣言，尽管有无数预言指涉基督；但他以为他是在控告预言圣经，而他甚至没有陈述他自己所称的它们的\Quote{似是而非的性质！}然而他并不知道，并非犹太人（如我们在前文所示）说基督将作为一位神或一位神之子降临。当他断言\Quote{他被我们说是已经来了，但被犹太人说是作为默西亚的来临尚在未来}，他似乎在用这个指控批评我们的陈述是最无耻的，并且无需冗长的驳斥。\par

% structure: p0006 source=h2 target=subsection reason=relative-heading-rank; base=h2

\subsection{第三章}

他又接着说道：\Quote{天主这样的降下有什么意义？}他没有注意到，按照我们的教导，降下的意义首先在于转化福音中所称的\InnerQuote{以色列家迷失的羊}，其次在于因他们的背命而夺去所谓的\Quote{天主的国}，并把它交给别的园户——不是古代犹太人，而是基督徒——他们将在适当的季节把天国的果实交给天主（每个行动都是\Quote{天国的果实}）。因此，我们将从众多的回答中选择一些话来答复策尔苏斯的问题，当他说：\Quote{天主这样的降下有什么意义？}策尔苏斯在这里自己给出一个答案，这个答案既不会被犹太人也不会被我们给出，当他问：\Quote{是为了了解世上发生的事吗？}因为在我们中没有人说基督进入此生是为了了解人间的事。然而，他立刻好像有人会回答说\InnerQuote{是为了了解人间的事}，便对自己说法提出反对：\Quote{他不是全知的吗？}然后，好像我们要回答\Quote{他全知}，他又提出一个新问题，说：\Quote{那么他全知，却不改善（人），也不能藉他的神能改善（人）。}他所有这些话都是愚妄之言；因为天主藉着他的圣言——圣言不断从一代进入圣善的灵魂，使他们成为天主的友人和先知——确实改善了那些听从其话的人；并且藉着基督的来临，祂通过基督信仰的教义改善人，不是那些不愿意的，而是那些选择了更好生活、中悦天主的人。此外，我不知道策尔苏斯在提出反对时，期望发生什么样的改善，他问：\Quote{难道祂就不能藉祂的神能使（人）改善，除非祂特意差遣某人吗？}他是否希望改善藉着天主以新思想充满人心、立即除去（内在的）邪恶、并植入美德来实现？另一个人现在可能会问，这在事物的本性中是否不一致或不可能；然而我们会说：\Quote{就算如此，且让它可能。}那么，我们的自由意志在哪里？赞同真理有什么可称赞的？拒绝虚假又怎么会值得赞美？但即使一旦承认这不仅是可能的，而且可以（由天主）适当地完成，为什么人们不更应该追问（提出类似策尔苏斯的问题）：为什么天主不能藉祂的神能造出不需要改善、而是本身有德和完美的人，而邪恶完全不存在？这些问题会困惑愚昧无知的人，却不会困惑那洞悉事物本性的人；因为如果你取消了德行的自发性，你就摧毁了它的本质。然而，这需要整个论文来讨论；希腊人在他们关于天主眷顾的著作中已经对此作了大量论述。他们确实不会说策尔苏斯用话语所表达的：\Quote{天主确实全知，但不改善（人），也不能藉祂的神能这样做。}我们自己在我们的著作的许多部分已经尽可能谈论了这些要点，而圣经也给那些能理解的人确立了同样的道理。\par

% structure: p0008 source=h2 target=subsection reason=relative-heading-rank; base=h2

\subsection{第四章}

策尔苏斯用来攻击我们和犹太人的论点，将这样反过来针对他自己：我的好先生，那统万有的天主知道人间发生的事吗？还是不知道？如果你承认天主和眷顾的存在，正如你的论文所表明的，祂必然知道。如果祂知道，为什么祂不使（人）改善？那么，我们有义务为天主不使（人）改善（尽管祂知道他们的状况）的做法辩护，而\Emph{你}——你并没有清楚地在你的论文中表明你是伊壁鸠鲁派，却假装承认眷顾——不也同样有义务解释为什么天主虽然知道人间发生的一切，却不改善（人），也不藉神能救众人脱离邪恶？我们并不羞于说，天主不断地差遣（导师）来使人改善；因为在人中间有天主赐下的道理，劝勉他们进入更好的生活。但在事奉天主的人中，有许多不同，而完美和纯洁的真理使者，以及带来彻底改造的人，如梅瑟和先知，是很少的。但超越这所有的，是耶稣所带来的伟大改造，祂愿意治愈不仅生活在世界一隅的人，而是尽可能治愈每个国家的人，因为祂来是为了\Emph{所有人}的救主。\par

% structure: p0010 source=h2 target=subsection reason=relative-heading-rank; base=h2

\subsection{第五章}

\Person{色尔苏斯}，不知何故，接着向我们提出另一项反对，仿佛我们声称\Quote{天主自己要降来人间}。他还想像由此推得\Quote{祂离开了自己的住所}；因为他不知道天主的德能，也不知道\BibleQuote{主的神充满了世界，那支持万有的，通晓言语}。他也不明白\BibleQuote{我岂不充满天地？上主说}。他也不明白，按照基督信仰的教义，我们都\BibleQuote{在祂内生活、行动、存在}，正如\Person{保禄}在向\Place{雅典人}的讲论中所教导的。因此，即使宇宙的天主通过自己的德能与\Person{耶稣}一同降来进入人的生活，即使那在起初与天主同在的圣言，祂自己就是天主，来到我们中间，祂并不让出位置或腾出座位，以致一个地方空无祂，而另一个原来没有祂的地方被充满。相反，天主的德能与神性通过祂所拣选的人而来，并在祂所找到的地方居住，并不改变位置，也不离开自己的地方而充满另一个地方；因为当我们说祂离开一个地方而占据另一个地方时，我们并不是按\OriginalTerm{地方}{topically}的意义来理解这些说法；而是说恶人的灵魂以及被罪恶淹没之人的灵魂被天主所遗弃，而愿意生活正直者、正在进步（在道德生活中）者、或已按正直生活之人的灵魂则被充满，或成为天主圣神的分享者。因此，基督的降来或天主来到人间，并不需要祂放弃一个更大的座位，也不需要地上的事物被改变，如\Person{色尔苏斯}所想像的，他说：\Quote{如果你改变地上的一个事物，哪怕是最小的，一切都会被颠覆并消失。}如果我们必须谈论因天主德能的显现和圣言进入人间而发生在任何人身上的改变，我们不会不情愿地说，那允许天主圣言进入其灵魂的人，从邪恶变为正直，从放荡变为节制，从迷信变为虔敬。\par

% structure: p0012 source=h2 target=subsection reason=relative-heading-rank; base=h2

\subsection{第六章}

但如果你们要我们面对\Person{色尔苏斯}最荒谬的指控，请听他这样说：\Quote{如今，天主在人群中不被认识，并因此认为自己受到亏待，祂就想要使自己被人知道，也要考验那些相信祂的和那些不相信的，如同人类中那些新近获得财富并炫耀其财富的人一样；这样，他们就证明了天主方面有一种过分的、但完全是凡俗的野心。}我们于是回答说：天主不被恶人所认识，祂愿意使自己被认识，并非因为祂认为自己受到亏待，而是因为对祂的认识将使那拥有这认识的人脱离不幸。不仅如此，祂并非出于想考验那些相信或不相信祂的人的意愿，藉祂那不可言说的神圣德能亲自住在某些人内，或派遣祂的基督；祂这样做，是为了使那些相信祂并接受祂神性的人摆脱一切不幸，并使那些\Emph{不}相信的人不再以此为借口，即他们的不信是由于没有听到训诲之言。那么，有什么论证能证明根据我们的观点，天主如我们所描述的，是\Quote{像人类中那些新近获得财富并炫耀其财富的人}？因为天主并不向我们炫耀，并不出于希望我们理解和思量祂的优越；而是希望由祂被我们所知而产生的福乐植入我们的灵魂，祂便通过基督和祂永存的圣言成就这事，使我们与祂亲密交往。因此，基督信仰教义并不证明天主有任何凡俗的野心。\par

% structure: p0014 source=h2 target=subsection reason=relative-heading-rank; base=h2

\subsection{第七章}

我不知道这是怎么回事：在就我们刚才讨论的主题发表了一番愚蠢言论之后，他竟然又加上以下的话：\Quote{天主不希望为自己而使人认识祂，而是为了我们的救恩，愿意将祂的知识赐给我们，使那些接受的人能成为善良并获得救恩，而那些不接受的人则显明为邪恶并受惩罚。}然而，在说了这样的话之后，他又提出一个新的反对：\Quote{那么，在如此漫长的时间之后，天主现在才想起要使人们过正义的生活，而以前却忽略了这样做吗？}我们回答说：从来没有一个时候天主不愿望人们过正义的生活；祂不断显明对理性动物进步的关心，为他们提供修德的机会。因为在每一代，天主的智慧进入那些它确认为圣洁的灵魂，使他们成为天主的朋友和先知。而且在\WorkTitle{圣经}中可以找到每一代的圣洁者、天主圣神的领受者，以及尽力在其能力范围内感化同时代的人的名字。\par

% structure: p0016 source=h2 target=subsection reason=relative-heading-rank; base=h2

\subsection{第八章}

在某些世代中存在过先知，他们因更坚定、更有力的（宗教）生活，在领受神圣感召方面超越了同时代的其他先知，也超越了之前和之后的先知，这并不足为奇。因此，曾经有一个时期，某种卓越超群之物居于人类当中，并超越了一切先前或后来的事物，这也丝毫不令人惊讶。然而，这些事情蕴含着深奥的奥秘，是普通理解力所无法接受的。为了使这些困难消失，并回应那些针对基督来临所提出的反对意见——即\Quote{难道天主在如此长久的时期之后，现在才想到要使人过正义的生活，而之前却忽略了吗？}——必须谈及关于（万民）分开的叙述，并阐明为何\BibleQuote{至高者划定万民的疆界时，分开亚当子孙时，祂按天主天使的数目立定万民的边界，上主的份子是祂的子民雅各，以色列是祂产业的量绳；}还必须说明为什么每个人出生在特定的边界内，属于那按阄分得边界的人，以及为什么\BibleQuote{上主的份子是祂的子民雅各，以色列是祂产业的量绳}是合理的，并且以前上主的份子是祂的子民雅各，以色列是祂产业的量绳。至于后来的人，圣父对救主说：\BibleQuote{你向我请求，我就将外邦人赐给你作产业，将地极赐给你作基业。}因为有一些关于人类灵魂不同待遇的相互关联的原因，难以陈述和研究。\par

% structure: p0018 source=h2 target=subsection reason=relative-heading-rank; base=h2

\subsection{第九章}

因此，虽然\Person{刻尔苏斯}可能不愿承认，但在那著名的以色列的众多先知之后，作为全世界的革新者的基督来临了，祂不需要像先前之经济时代那样对人使用鞭子、锁链和酷刑。因为当撒种者出去撒种时，道理就足以将圣言撒遍各处。但如果有一个必然要限制世界存续的时期（因为世界有开端），如果世界将有终结，终结后将有万物的公义审判，那么，以哲学方式处理福音宣示的人，就必须用各种论证来确立这些教义，这些论证不仅直接来自圣经，而且来自由圣经推导出的推论；而人数众多、较为单纯的信徒，以及那些不能理解上智之诸多不同层面的人，必须将自己交托于天主和我们人类的救主，并满足于祂的\Quote{ipse dixit}，而不需要任何其他论证。\par

% structure: p0020 source=h2 target=subsection reason=relative-heading-rank; base=h2

\subsection{第十章}

其次，\Person{刻尔苏斯}，一如他的习惯，既没有证明也没有确立任何事，就继续说，好像我们以一种既不神圣也不虔诚的方式谈论天主，以致\Quote{很明显，他们以一种既不神圣也不虔敬的方式谈论天主；}他想象我们这样做是为了引起无知者的惊讶，并且我们没有就犯罪者必须受惩罚这一必要性讲真话。因此，他把我们比作那些\Quote{在酒神秘仪中引入幻影和恐怖之物的人。}至于酒神秘仪，是否存在任何可靠的记载，或者没有，让希腊人去说吧，让\Person{刻尔苏斯}和他的宴饮同伴去听吧。但我们为自己的做法辩护，我们说我们的目标是改造人类，要么通过刑罚的威胁——我们确信刑罚对整个世界都是必要的，并且对遭受它们的人也许并非无用——要么通过向那些度德性生活的人所许下的诺言，这些诺言中包含着关于有福结局的陈述，这结局将在天国中找到，是为那些堪当成为祂子民的人所保留的。\par

% structure: p0022 source=h2 target=subsection reason=relative-heading-rank; base=h2

\subsection{第十一章}

此后，他想要表明，我们关于洪水或大火的论述既不新奇也不惊人，而是由于误解了希腊人或野蛮民族中流传的这些事件的记载，我们在谈及它们时相信了我们自己的圣经。他写道：\Quote{这个信念在他们中间传播开来，源于对这些事件的误解，即经过漫长的周期、行星的回归与会合，大火和洪水惯常发生；而且因为上次洪水发生在\Person{德乌卡利翁}时代，时间的流逝，随着万物的交替，要求一场大火，这使他们说出一个错误的意见，即天主将降下，像施刑者一样带着火来。}现在，作为对此的回应，我们说：我不明白，\Person{刻尔苏斯}虽然博览群书，并且表明他研读过许多历史，却怎么没有注意到\Person{摩西}的古老性。\Person{摩西}被某些希腊史学家叙述说生活在\Person{福洛纽斯}之子\Person{伊那科斯}的时代，被埃及人以及那些研究腓尼基历史的人承认为极其古老的人。任何愿意的人都可以查阅\Person{弗拉维奥·约瑟夫斯}关于犹太古代的两卷书，以便了解\Person{摩西}如何比那些主张在大段间隔时间后世界上发生洪水和大火的人更为古老；而\Person{刻尔苏斯}断言，犹太人和基督徒误解了这些言论，并且由于不理解关于大火的论述，就宣称\Quote{天主将降下，像施刑者一样带着火来。}\par

% structure: p0024 source=h2 target=subsection reason=relative-heading-rank; base=h2

\subsection{第十二章}

那么，不论是否有时间周期、洪水或定期发生的火灾，也不论圣经是否知道这种事，不仅在许多经文中，尤其在\Person{撒罗满}说：\Quote{已有的事，后必再有；已行的事，后必再行}等等时，不属于当前讨论的场合。因为只须指出，\Person{梅瑟}和某些先知是远古之人，他们并未从他人领受有关（未来）世界火灾的论述；相反（如果我们要注意时间问题），倒是其他人误解了这些论述，没有准确查考他们的话，从而捏造了同一事件按一定间隔重演的虚构，且在本质和偶然属性上毫无区别。但我们并不把洪水或火灾归因于周期和行星周期；而是宣告它们的原因是邪恶的广泛盛行，以及由此导致的洪水或火灾将其清除。如果先知的声音说天主\Quote{降临}，而那位曾说\Quote{我岂不是充满天地吗？——主说}，这个词语是用于比喻意义。因为天主从祂自己的崇高和伟大中\Quote{降临}，当祂安排人的事务，尤其是恶人的事务时。正如习俗使人说教师\Quote{俯就}孩童，智者俯就那些刚投身哲学的青年，并不是以\Emph{身体}的方式\Quote{下降}；同样，如果圣经中任何地方说天主\Quote{降临}，应理解为是符合使用该词语的惯用法，同样，\Quote{上升}一词也是如此。\par

% structure: p0026 source=h2 target=subsection reason=relative-heading-rank; base=h2

\subsection{第十三章}

但因为\Person{才尔苏}嘲笑地说我们谈论\Quote{天主像行刑者带火一样降临}，从而不恰当地迫使我们探究更深奥的词语，我们将作几点说明，足以使听众对我们的辩护有所认识，以消除才尔苏对我们的嘲弄，然后我们将转向下文。神圣的圣言说我们的天主是\Quote{吞噬的火}，且\Quote{祂在祂面前引来火河}；甚至祂进入如同\Quote{炼金者的火和漂布者的草}，以净化祂的子民。但当祂被称为\Quote{吞噬的火}时，我们探究哪些事物适于被天主吞噬。我们断言它们是邪恶及其产生的行为，这些行为借喻地被称为\Quote{木、草、禾秸}，天主如同火一样吞噬它们。因此，恶人被说成在先前奠定的理性根基上建造\Quote{木、草、禾秸}。那么，如果有人能证明这些话被作者以不同方式理解，并能证明恶人\Emph{字面上}建造\Quote{木、草、禾秸}，那么显然火必须被理解为物质的、感官的对象。但相反，如果恶人的行为以\Quote{木、草、禾秸}的名称借喻地说出来，为什么不同时想到（探究）\Quote{火}一词应如何理解，以使这类\Quote{木}被焚毁呢？因为（圣经）说：\Quote{火要试验各人的工程是哪种。如果人在上面所建造的工程存留，他就要获得赏报；如果人的工程被烧了，他就要受到损失。}但除了所有源于邪恶的行为外，还有什么工程可以被这些话称为\Quote{被烧}的呢？因此，我们的天主是\Quote{吞噬的火}，是按我们所理解的这个词的意义；因而祂作为\Quote{炼金者的火}进入，以精炼那充满邪恶之铅的理性本性，并使其摆脱其他不纯的物质，这些物质掺杂了灵魂天然的、可以说是金银的东西。同样，\Quote{火河}也被说成是在天主面前，祂将彻底清洗那混合于整个灵魂中的邪恶。但这些话足以回答这个断言：\Quote{他们如此表达是为了宣扬一种错误观点，即天主将像行刑者一样带着火降临。}\par

% structure: p0028 source=h2 target=subsection reason=relative-heading-rank; base=h2

\subsection{第十四章}

但让我们看看策尔苏斯随后如何以极大的炫耀以如下方式宣布：\Quote{再说，}他说道，\Quote{让我们从开始重新论述这个主题，并拿出更多的证据。我并没有提出新的说法，而是陈述早已确定的事。天主是良善的、美丽的、有福的，而且是在最好和最美丽的程度上。但如果祂降世到人间，祂必定经历变化，即从善变为恶，从美德变为恶习，从幸福变为痛苦，从最好变为最坏。那么，谁会选择这样的变化呢？凡人的本性确实会经历变化和重塑，但不朽者的本性是保持同一且不变。因此，天主不可能接受这样的变化。}如今在我看来，当我已经叙述了圣经中所说的天主对人类事务的\Quote{屈尊}时，对这些异议已给出了适当的回答；为此目的，祂不需要经历转变，如策尔苏斯所想我们宣称的那样，也不需要从善变为恶，从美德变为恶习，从幸福变为痛苦，从最好变为最坏。因为，在祂的本质中保持不变，祂藉祂眷顾的经纶屈尊于人类事务。因此，我们表明，圣经将天主描述为不变的，既有如\Quote{祢是相同的，}这样的话，以及我不改变；而伊壁鸠鲁的众神，由原子组成，就其结构而言可分解，却竭力摆脱含有毁灭元素的原子。不，甚至斯多葛派的神，作为有形体的，在（宇宙）大火发生时，有时其整个本质由主导原则构成；而在万物重组时，又再次部分地成为物质。因为连斯多葛派也无法清晰理解关于天主的自然概念，即祂是完全不朽坏、单纯、非复合且不可分的。\par

% structure: p0030 source=h2 target=subsection reason=relative-heading-rank; base=h2

\subsection{第十五章}

至于祂曾降世到人间，祂曾\Quote{先前具有天主的形象，}并出于慈爱，舍弃了（祂的荣耀），以便能被人们接受。但我想，祂并没有经历从\Quote{善变为恶，}因为\Quote{祂没有犯\Emph{罪}；}也没有从\Quote{美德变为恶习，}因为\Quote{祂不认识\Emph{罪}。}也没有从\Quote{幸福变为痛苦，}而是祂贬抑自己，却仍然有福，甚至当祂的贬抑是为了我们人类的益处时。祂也没有从\Quote{最好变为最坏，}因为良善和慈爱怎能是\Quote{最坏的呢？}若说那位医生，他为了治愈病人而观看可怕景象、处理不洁物体，是从\Quote{善变为恶，}或从\Quote{美德变为恶习，}或从\Quote{幸福变为痛苦，}这合适吗？然而医生在观看可怕景象和处理不洁物体时，并不能完全避免自己陷入同样命运的可能性。但那位藉着内在于祂的天主圣言医治我们灵魂创伤的祂，自身不会接受任何邪恶。但如果那不朽的天主——圣言——藉着取了可朽的身体和人的灵魂，在策尔苏斯看来是经历了变化和转变，那么让他知道，圣言在本质上仍然是圣言，并不遭受身体或灵魂所遭受的任何事物；相反，祂偶尔屈就于那不能观看天主的光辉和灿烂的人，祂成为好像是肉身，以字面的声音说话，直到那以如此形式接受祂的人，因圣言的教导而被稍稍提升，能够凝视，可以说是，祂真实而卓越的显现。\par

% structure: p0032 source=h2 target=subsection reason=relative-heading-rank; base=h2

\subsection{第十六章}

因为圣言有，可以说是，不同的显现，根据祂向每个前来接受祂教导的人显现的方式；而这是依照那刚成为门徒者、或稍有进步者、或更进一步者、或已经\Emph{几乎}达到美德者、或甚至\Emph{已经}达到美德者的状况以相应的方式。因此，并非如策尔苏斯及其同类所认为的那样，我们的天主被转变了，并登上高山，显示了祂的真实显现是与那些留在山下、不能跟随祂上山的人所见的不同且远为卓越的。因为山下的人没有能看见圣言转变为祂荣耀和更神圣状况的眼睛。他们只能艰难地接受祂本来的样子；所以那些不能看见祂更卓越本性的人可以这样说祂：\Quote{我们看见了祂，祂没有美貌，也无俊美；但祂的形象是卑微的，次于人子。}让这些言论作为对策尔苏斯臆测的回答，他并不理解历史中所记载的耶稣的变化或转变，也不理解祂有死和不死的本性。\par

% structure: p0034 source=h2 target=subsection reason=relative-heading-rank; base=h2

\subsection{第十七章}

但那些叙述，尤其是当它们被按正确意义理解时，岂不比\Person{狄俄尼索斯}被\Person{提坦}欺骗、被赶出\Person{朱庇特}宝座、被他们撕成碎片、后来遗体被重新拼合、仿佛再次复活并升天的故事更值得尊敬吗？希腊人是否可随意将这类故事归于灵魂学说并作比喻解释，而针对\Emph{我们}，一扇一致的、且处处与住在纯洁灵魂中的圣神著作协调的解释之门却关闭着？那么，\Person{克尔苏斯}完全不了解我们著作的目的，因此他是在按自己的理解贬低它们，而非按它们的真实含义；反之，如果他曾思考一个要享永生之灵魂的适当之处，以及我们对其实质和原则应有的看法，他就不会如此嘲笑那不朽进入可朽身体的降临——它并非按照\Person{柏拉图}的\OriginalTerm{灵魂转世}{metempsychosis}，而是符合另一种更高层次的看法。他也会注意到一次以其伟大仁爱著称的\Quote{降下}，为了挽回（如圣经神秘地称呼的）\BibleQuote{以色列家迷失的羊}，它们从山上迷途下来，而牧人据某些比喻所说，曾下山寻找它们，却将那些\BibleQuote{没有迷失的}留在山上。\par

% structure: p0036 source=h2 target=subsection reason=relative-heading-rank; base=h2

\subsection{第十八章}

\Person{克尔苏斯}耽搁于他不了解之事，导致我们犯下赘言之罪，因为我们甚至不愿在表面上放过他任何一项异议。他于是继续写道：\Quote{天主要么真的像他们所称那样，将自己改变成一个可朽的身体（而此事之不可能已被宣称），要么他并\Emph{未}经历改变，而只是让观看者如此想象，从而欺骗他们，犯了虚妄之罪。但欺骗与虚妄不过是邪恶，只可用作药物：要么用于患病和癫狂的朋友，以图治愈，要么用于敌人，当一个人设法躲避危险时。但没有任何病人或癫狂者是天主的朋友，天主也从不惧怕任何人到需要以误导他来躲避危险的地步。}现在对这些陈述的回答或许部分与圣言（他是天主）的本性有关，部分与耶稣的灵魂有关。就圣言的本性而言，正如护士中食物的质量会参照婴儿的本性转变为乳汁，或医生为了病人健康之故予以安排（或专门为较强健者预备，因为他有更大的活力），同样，天主也会适合地改变——针对每个个体——圣言的能力，圣言天然具有滋养人类灵魂的禀性。有人领受（如圣经所称）\BibleQuote{圣言的纯奶}；另有一个较弱者领受\BibleQuote{蔬菜}；还有一个已长大成人者领受\BibleQuote{干粮}。我想，圣言不会违背自己的本性，在按每个人接受祂的能力提供滋养时欺哄或误导。但若有人将这种改变理解为指向耶稣灵魂进入身体后的事，我们就要问\Quote{改变}一词作何解。如果是指本质的改变，这种假设不可接受，不仅对耶稣的灵魂如此，对任何其他有理智的灵魂亦然。如果认为灵魂与身体结合时或进入某处时会受什么影响，那么那位以极大仁爱将救主带给人类世界的圣言会遭受什么不便？——过去自诩治病的人中，没有一个能像耶稣的灵魂那样，以其所行表明\Emph{它}所能做的，甚至为了我们族类的益处而自愿降至人类命运的水平。圣言深知此事，在圣经许多段落中也有类似表达，尽管目前引用保禄的一段证言便已足够：\BibleQuote{你们该怀有基督耶稣所怀有的心情：祂虽具有天主的形体，并没有以自己与天主同等为应当把持不放的，反而使自己空虚，取了奴仆的形体，成了人的模样；既有了人的样式，就贬抑自己，服从至死，且死在十字架上。因此，天主极其举扬祂，赐给了祂一个名字，超越一切名字。}\par

% structure: p0038 source=h2 target=subsection reason=relative-heading-rank; base=h2

\subsection{第十九章}

那么，别人或许会向\Person{策尔苏}承认天主并没有改变，只是使旁观者以为祂改变了；然而我们这些深信耶稣在人间来临并非仅是显现，而是真实的显示，并不受\Person{策尔苏}此项指控的影响。我们仍将尝试答复，因为，\Person{策尔苏}，你不是断言有时可以像用药一样使用欺骗和谎言吗？那么，如果竟要成就如此救恩的结果，发生一些类似的事件又有何荒谬呢？因为某些带有虚假意味的言辞对这类性格的人产生纠正作用（如同医生对病人作相似的宣告），胜过出于真理的言词。然而，这必须成为我们对抗其他反对者的辩护。因为那位医治患病朋友的主，用祂不会优先选用、仅依情况而用的疗法，医治亲爱的人类的，并无荒谬之处。而且，人类处于精神错乱时，必须用圣言视为有助于使那些受苦者恢复健全心智的方法来医治。但\Person{策尔苏}也说：\Quote{人向敌人这样做是为了躲避危险。但天主谁也不怕，以致于通过误导阴谋反对祂的人来逃避危险。}现在，根本无需且荒谬地答复一项没有人针对我们救主提出的指控。我们在答复其他指控时，已经回应了\Quote{无论是处于疾病或精神错乱状态的人，都不是天主的友人}这一说法。因为答复是：这样的安排并非为那些已经是友人、后来生病或患精神病的人的缘故，而是为使那些因灵魂疾病和自然理性的疏离而仍是敌人的人，成为天主的友人。因为明确记载着，耶稣为罪人承受一切，为使他们脱离罪恶，并归向正义。\par

% structure: p0040 source=h2 target=subsection reason=relative-heading-rank; base=h2

\subsection{第二十章}

接下来，由于他描述犹太人以他们特有的方式解释他们相信基督在他们中的来临仍在将来，而基督徒则以\Emph{他们}的方式坚持圣子进入人类生命已经发生，让我们尽所能简要思考这些要点。据\Person{策尔苏}说，犹太人声称\Quote{（人类的）生命充满一切邪恶，需要一个由天主派遣者，以便恶人受罚，万物得以净化，类似于发生的第一次洪水。}既然基督徒据说作了附加的陈述，显然他声称他们承认这些。那么，一位因邪恶泛滥而来净化世界并按各人应得对待各人的来临，有何荒谬呢？因为邪恶的蔓延不停止、万物不更新，不符合天主的属性。而且，希腊人也知道大地有时会被洪水或火灾净化，正如\Person{柏拉图}在某处也曾说道：\Quote{当众神淹没大地、用水净化它时，其中一些在山上，}等等。那么，难道必须说，如果希腊人作这样的断言，他们便值得尊敬和考虑，但如果我们持守某些被希腊人引述认可的观点，这些观点就不再可敬了吗？然而，关心我们所有记录的关联和真理的人，不仅会努力确立作者的古老性，还会确立其著作的庄严性及各部分的一致性。\par

% structure: p0042 source=h2 target=subsection reason=relative-heading-rank; base=h2

\subsection{第二十一章}

但我不明白，他怎能设想（\Place{巴贝耳}）塔的倾覆，其目的与洪水相似，洪水根据犹太人和基督徒的记载，净化了大地。因为，为了使\BibleBook{创世}{记}中关于塔的叙述不隐含任何秘密含义，而是如\Person{刻尔苏斯}所假设的那样，按字面理解，那么从这种观点来看，那件事似乎并非为了净化大地而发生的；除非他认为所谓的语言混乱就是这样一种净化过程。但在这一点上，有机会的人将在更合宜的时候加以讨论，其目的不仅要说明这段叙述在其历史关联中的含义，还要说明可以从中推衍出什么隐喻意义。然而，既然他想象撰写关于塔和语言混乱的记载的\Person{梅瑟}，歪曲了\Person{阿罗厄斯}儿子的故事，并将其归于那座塔，我们必须指出，我认为在\Person{荷马}时代之前，没有人提到过\Person{阿罗厄斯}的儿子，而我深信，\Person{梅瑟}所记载的关于塔的事情，不仅比\Person{荷马}，甚至比希腊文字的发明都要古老得多。那么，是谁歪曲了对方的叙述呢？是那些讲述\Person{阿罗亚代}故事的人歪曲了那个时代的历史，还是那位撰写关于塔和语言混乱记载的人歪曲了\Person{阿罗亚代}的故事？对于公正的听众而言，\Person{梅瑟}似乎比\Person{荷马}更古老。此外，\Person{梅瑟}在\BibleBook{创世}{记}中记载的因罪恶而毁灭\Place{索多玛}和\Place{哈摩辣}的火刑，被\Person{刻尔苏斯}比作\Person{法厄同}的故事——他所有这些说法都源于一个错误，即他没有注意到\Person{梅瑟}的（更）古老性。因为讲述\Person{法厄同}故事的人似乎甚至比\Person{荷马}还要年轻，而\Person{荷马}比\Person{梅瑟}年轻得多。因此，我们不否认，净化之火和世界的毁灭是为了扫除邪恶，更新万物；因为我们声称我们从先知的神圣经书中学会了这些。但是，正如我们在前面几页所说，先知们在预言许多未来事件时，表明他们对许多过去的事情说了真话，从而证明圣神的居驻，那么很明显，对于未来之事，我们应当信赖他们，或者更确切地说，信赖他们内住的圣神。\par

% structure: p0044 source=h2 target=subsection reason=relative-heading-rank; base=h2

\subsection{第二十二章}

但是，根据\Person{刻尔苏斯}的说法，\Quote{基督徒在犹太人的说法之外又添加了一些陈述，断言天主子已经为了犹太人的罪而被派遣；犹太人鞭打了\Person{耶稣}，并给他喝苦胆，从而为自己招来了天主的愤怒。}任何人若愿意都可以证明这种说法是虚假的，除非整个犹太民族在\Person{耶稣}从他们手中遭受这些苦难之后的一代人之内就被推翻了。因为我认为，在\Person{耶稣}被钉十字架之后四十二年，\Place{耶路撒冷}的毁灭发生了。自从犹太民族开始存在以来，从未记载过他们被逐出他们可敬的圣殿崇拜和服务如此之久，并被更强大的民族奴役；因为即使他们有时因自己的罪而被看似遗弃，他们仍然受到（天主的）眷顾，返回自己的土地，恢复自己的财产，并无阻碍地遵守他们的法律。因此，证明\Person{耶稣}是某种神圣和圣洁之物的一个事实是，犹太人因祂的长久遭受如此严重的灾难。我们充满信心地说，他们将永远不会恢复昔日的状况。因为他们犯下了最不圣洁的罪行，密谋反对人类救主，就在那座他们向上主献上包含伟大奥迹象征的崇拜的城市里。因此，那座\Person{耶稣}遭受这些苦难的城市理应彻底毁灭，犹太民族也应被推翻，而天主赐予他们的幸福邀请应传递给他人——我是指基督徒，纯正圣洁的崇拜教义已传给他们，他们获得了新的法律，与所有国家的既定政体相和谐；因为先前作为对一个单一民族（由其同族和相似习俗的王子统治）施加的法律，如今已无法完整地遵守。\par

% structure: p0046 source=h2 target=subsection reason=relative-heading-rank; base=h2

\subsection{第二十三章}

接下来，他按照惯常的方式嘲笑犹太人和基督徒的种族，将他们全都比作\Quote{一群蝙蝠，或一窝从巢中出来的蚂蚁，或沼泽中商议的青蛙，或粪堆角落中爬在一起的蠕虫，彼此争吵谁罪更大，并断言天主向我们预示并宣告万事；祂舍弃整个世界、天界和这广阔的大地，唯独在我们中间成为公民，唯独向我们传达祂的旨意，并不停地派遣和询问，我们怎样才能永远与祂联合。}在他虚构的描述中，他将我们比作\Quote{蠕虫，断言有一位天主，并且我们这些由祂所造的人立即像天主一样，万物——地、水、空气、星辰——都服从我们，并且万物都是为我们而存在，并被注定服从我们。}根据他的描述，这些蠕虫——也就是我们自己——说：\Quote{现在，既然我们中间某些人犯罪，天主将降临，或派遣祂的儿子，用火消灭恶人，以便我们其余的人能与祂同享永生。}所有这些之后，他又补充说：\Quote{这样的争论在蠕虫和青蛙之间比在犹太人和基督徒之间更可忍受。}\par

% structure: p0048 source=h2 target=subsection reason=relative-heading-rank; base=h2

\subsection{第二十四章}

针对这些，我们反问那些接受此类针对我们的诽谤的人：你们是否因天主的卓越而视所有人为一群蝙蝠、青蛙或蠕虫？或者你们并不将其他人类纳入此拟议的比较，反而因其拥有理性与既定法律而视\Emph{他们}为人，却轻视\Emph{基督徒}和\Emph{犹太人}，因为他们的意见令你们不悦，并将他们比作上述动物？无论你们如何回答我们的问题，我们都会努力表明，这种论断极不相称，无论是泛指全人类还是特指我们。因为，假设你们恰当地认为，所有人都与天主相比（正确地）被比作这些毫无价值的动物，因为他们的微小完全无法与天主的优越相比，那么你们所谓的微小是什么意思？诸君，请回答我。若你们指身体之微小，那么要知道，若以真理为判断，优劣并非由身体标准决定。因为，依此观点，秃鹫和大象会比我们人类优越；因为它们比我们更大、更强壮、更长寿。但明智者不会认为这些无理性受造物仅仅因身体而优于理性存在：因为拥有理性使理性存在远超前于一切无理性受造物。即使是德性而蒙福的存在族类也会承认这点，无论你们称之为善神，或者如我们惯常所称的天主的天使，或其他任何高于人的本性，因为他们内的理性官能已臻完美，并被赋有全部德性品质。\par

% structure: p0050 source=h2 target=subsection reason=relative-heading-rank; base=h2

\subsection{第二十五章}

但如果你们贬抑人之微小，不是因身体，而是因灵魂，认为它劣于其他理性存在，尤其是那些有德者，并因邪恶居于其中而低劣——那么，为何基督徒中的恶者与犹太人中的罪人应被称为一群蝙蝠、蚂蚁、蠕虫或青蛙，而不是其他民族中行恶的个人？因为在此方面，任何个人，尤其是被邪恶潮流冲走者，便是蝙蝠、蠕虫、青蛙、蚂蚁。虽然一个人可以是如\Person{狄摩西尼}那样的演说家，却若沾染像他那样的邪恶，犯下像他那样出于邪恶本性的行为；或者像\Person{安提丰}，他也被视为真正的演说家，却在其题为\WorkTitle{论真理}的著作中（类似于\Person{塞尔苏斯}的那篇论说）摧毁了天主眷顾的教义，——这样的人仍然是蠕虫，在愚昧与无知的粪堆角落打滚。的确，无论理性官能的本性为何，它都不能被合理地比作蠕虫，因为它拥有德性的能力。因为这些朝向德性的朦胧倾向不允许那些拥有获取德性之能且不能完全丧失其种子者被比作蠕虫。因此，看来一般人也无法在与天主比较时被视为蠕虫。因为理性，其起源在天主的理性中，不允许理性动物被视为完全疏离于神性。基督徒与犹太人中的恶者（他们实际上既非基督徒亦非犹太人），也不能比其他恶人更多地被比作在粪堆角落打滚的蠕虫。而如果理性的本性不允许这样的比较，那么显然我们不可诋毁那为德性而形成的人性，即使它因无知而犯罪，也不可将其比作所述的那类动物。\par

% structure: p0052 source=h2 target=subsection reason=relative-heading-rank; base=h2

\subsection{第二十六章}

但是，如果因为基督徒和犹太人的那些意见（这些意见似乎完全不为\Person{塞理斯}所理解）令塞理斯不悦，而使他们被视为蠕虫和蚂蚁，其余人类则不同，那么让我们考察基督徒和犹太人公认的意见，并将其与其余人类比较，看看是否那些曾经承认某些人是蠕虫和蚂蚁的人，不会\Emph{他们}自己显出是蠕虫、蚂蚁和青蛙，因为他们偏离了对天主的正确见解，并在一种虚妄的虔诚外观下，敬拜无理性的动物、偶像或其他人类手造的物品；然而，从这些受造物的美丽，他们本该赞美造主并敬拜祂。而那些确实是人类、甚至比人更尊贵的人（若有任何如此者），他们顺从理智，能够从草木石块，甚至从被认为最珍贵的物质——金银——提升，也从世界上的美好事物上升至万有的创造者，并将自己托付给祂，唯有祂能满足一切存在之物，鉴察所有人的思念，俯听所有人的祈祷；他们向祂献上祈祷，并在祂面前——祂鉴察一切——行事，并谨慎地，犹如在万物聆听者面前，不说任何不适宜向天主报告的话。这样的虔敬——既不受劳苦、也不受死亡危险、更不受花言巧语所战胜——难道不能使获得它的人不再被比作蠕虫吗？即使在他们获得如此非凡的虔敬之前曾被如此描述。那些制服了对性欲的强烈渴求——这种渴求使许多人的灵魂变得软弱无力——并因确信除非通过节制升到天主面前，否则无法与天主共融，而制服它的人，在你看来岂不成了蠕虫的兄弟、蚂蚁的亲戚、与青蛙相似吗？什么！正义的光辉品质——它保持对邻人和亲属共同权利的不可侵犯，并保持公平、仁慈和良善——难道不能使实践它的人免于被称为夜间飞鸟吗？那些沉溺于放荡——正如人类大多数那样——并与普通娼妓随意交往，且教导这种行为并非完全有悖于体面的人，岂不正是打滚的蠕虫吗？——尤其当他们与那些受过教训、不可将\BibleQuote{基督的肢体}和\OriginalTerm{圣言}{Word}所居住的身体，变成\BibleQuote{娼妓的肢体}的人相比时；他们已经知道，理性受造物的身体，既然奉献给万有的天主，就是他们所敬拜的天主的殿，这殿因他们对造主纯洁的认知而成为这样，并且他们小心谨慎，不因非法的快乐亵渎天主的殿，从而将节制作为对天主的虔敬来实践。\par

% structure: p0054 source=h2 target=subsection reason=relative-heading-rank; base=h2

\subsection{第二十七章}

我还没有说到人类中盛行的其他罪恶，即使是那些貌似哲学家的人也不能迅速从中解脱，因为在哲学中有许多假充者。我也没有提到许多这样的罪恶也存在于既非犹太人又非基督徒之中。事实上，如果你正确地考察什么构成基督徒，那么这样的恶行在\Emph{基督徒}中根本不会盛行。或者，如果发现任何这类人，至少他们不会出现在那些参加聚会、来到公开祈祷场合的人中，除非（这种情况很少发生）某个这类人混在人群中未被察觉。因此，我们不是聚集在一起的蠕虫；我们根据犹太人相信为神圣的圣经与他们抗辩，指出预言中所说的那一位\Emph{已经}来临，而\Emph{他们}因罪大恶极已被抛弃，\Emph{我们}这些接受了\OriginalTerm{圣言}{Word}的人对天主怀有最高的希望，既因我们信靠祂，也因祂有能力接纳我们进入祂的共融，使我们纯洁，摆脱一切邪恶和生活的恶毒。那么，如果一个人自称是犹太人或基督徒，他不会无条件地说天主为了我们特别创造了整个世界和穹苍。但是，如果一个人如同耶稣所教导的那样，心地纯洁、温良、和平，并愉快地为信仰而甘冒危险，那么这样的人可以合理地信赖天主，并充分理解预言中的话语，还可以说：\Quote{这一切都是天主预先指示我们这些相信的人，并宣告给我们的。}\par

% structure: p0056 source=h2 target=subsection reason=relative-heading-rank; base=h2

\subsection{第二十八章}

然而，既然他将那些被他视为虫子的人，即基督徒，说成是在说：\Quote{天主抛弃了天上诸域，轻视这广阔大地，唯独在我们中间居住，唯独向我们发言，并不断传递祂的消息和询问，问我们如何能永远成为祂的同伴，}我们必须回答：他归于我们的是我们从未说过的话，因为我们既阅读也知道，天主爱一切存在之物，并不憎恨祂所创造的任何东西，因为祂不会怀着仇恨创造任何事物。此外，我们读到这样的宣告：\BibleQuote{祢爱护一切所有，因为都是祢的，灵魂之爱慕者。因为祢不朽的神在万物之内。因此，对于那些暂时堕落的人，祢也加以责备，告诫他们，提醒他们的罪恶。}我们怎能断言\Quote{天主离开天上诸域和整个世界，轻视这广阔大地，唯独在我们中间居住，}既然我们发现所有有思想的人都在祈祷中说：\BibleQuote{大地充满了主的仁慈，}以及\BibleQuote{主的仁慈普及所有血肉；}\BibleRef{德 18:13}并且，天主是良善的，祂\BibleQuote{使太阳升起，光照恶人也光照善人，降雨给义人也给不义的人，}并且，祂鼓励我们效法祂，好使我们成为祂的儿子，教导我们尽可能将我们所享的恩惠扩展到所有人。因为祂自己被说成是所有人的救主，尤其是信徒的；祂的基督是\BibleQuote{赎我们的罪，而且不只是赎我们的，也是赎全世界的罪。}这就是我们对凯尔苏斯指控的回答。某些其他说法，与犹太人的特征相符，可能出自该民族的一些人，但绝非出自基督徒，因为基督徒受教知道：\BibleQuote{基督在我们还是罪人的时候，为我们死了，这证明了天主对我们的爱，}而且\BibleQuote{为义人死，是少有的；为善人死，或许有人敢去死。}但是现在，耶稣被宣称是为全世界罪人的缘故而来（好使他们离弃罪恶，把自己托付给天主），并且按照这些圣经的古老习俗，被称为\Quote{天主的基督。}\par

% structure: p0058 source=h2 target=subsection reason=relative-heading-rank; base=h2

\subsection{第二十九章}

然而，凯尔苏斯或许误解了他所称的某些\Quote{虫子，}因为他们断言说：\Quote{天主存在，并且我们仅次于祂。}他这样做，就像那些因为某个鲁莽的青年听了三天哲学讲座后，就自视高于他人，认为他人低劣且毫无哲学素养，而对整个哲学派别吹毛求疵的人一样。因为我们知道，有许多受造物比人更尊贵；我们且读到：\BibleQuote{天主在诸神的会中站立，}但这些神并非列邦所敬拜的，因为\BibleQuote{列邦的一切神祇都是偶像。}我们也读到：\BibleQuote{天主在诸神的会中站立，在诸神之中施行审判。}此外，我们知道：\BibleQuote{虽有称为神的，或在天上，或在地上，就如那许多的神，许多的主；但在我们，却只有一位天主，就是父，万物出于祂，我们也归于祂；并有一位主耶稣基督，万物藉祂而有，我们也藉祂而有。}我们知道，天使就是这样优于人类；因此，当人达于完美时，就会变得像天使一样。\BibleQuote{因为在复活的时候，人也不娶也不嫁，但义人如同天上的天使，}也变成\BibleQuote{与天使同等。}我们也知道，在宇宙的安排中，有一些存在被称为\Quote{宝座、}\Quote{宰制、}\Quote{掌权}和\Quote{率领；}我们看到我们人类虽然远不及这些，却可以希望藉着有德行的生活和凡事合乎理性，得以提升而相似于这一切。最后，因为\BibleQuote{将来如何，还未显明；但我们知道，当祂显现时，我们将像天主，并要看见祂实在怎样。}如果有人主张某些人所说的（无论是那些拥有智慧还是不拥有智慧，但误识了健全理性的人），即\Quote{天主存在，而\Emph{我们}仅次于祂，}我就会用\Quote{我们这些按理性行事的人}来替代\Quote{我们}一词，或者更确切地说，\Quote{\Emph{有德行的、}按理性行事的我们。}因为在我们看来，\Emph{所有}蒙福者都拥有同一德行，所以人的德行与天主的德行是同一的。因此，我们受教要成为\BibleQuote{完美的，如同我们在天的父是完美的一样。}那么，没有一个善良有德行的人是\Quote{在粪土中打滚的虫子，}没有一个虔诚的人是\Quote{蚂蚁，}没有一个义人是\Quote{青蛙；}一个灵魂被真理之光光照的人，也不能合理地比作\Quote{夜间的鸟。}\par

% structure: p0060 source=h2 target=subsection reason=relative-heading-rank; base=h2

\subsection{第三十章}

在我看来，策尔苏斯也误解了这句话：\BibleQuote{让我们按我们的肖像和模样造人}；因此把\Quote{虫}们说成，既然由天主所造，我们完全相似于祂。然而，如果他懂得人被造\BibleQuote{是天主的肖像}与\BibleQuote{按祂的模样}之间的区别，并且知道天主记载说：\BibleQuote{让我们按我们的肖像和模样造人}，但祂造人\BibleQuote{是天主的肖像}，而当时并没有\BibleQuote{按祂的模样}，他就不会说我们宣称\Quote{我们完全相似于祂}。此外，我们并不说众星是我们隶属的；因为那被称为\Quote{义人的复活}、并被智慧者所理解的复活，被那说\BibleQuote{太阳有一种光荣，月亮有另一种光荣，星辰也有另一种光荣，而且星辰与星辰的光荣也有分别；死人的复活也是如此}的人比作太阳、月亮和星辰。达尼尔也早已预言了这些事。策尔苏斯又说，我们断言\Quote{万物都被安排来隶属我们}，也许他听过我们中间有些智者这样说，也许他也不懂那句\BibleQuote{我们中间最大的，是众人的仆役}的话。如果希腊人说\Quote{太阳和月亮是必死之人的奴仆}，他们赞同这话并解释其意义；但既然这样的话我们或者根本不说，或者以另一种方式表达，那么策尔苏斯在这里也是诬告我们。此外，我们这些按策尔苏斯说法是\Quote{虫}的人，被他描述说：\Quote{既然我们中间有人犯罪，天主就会来我们这里，或派遣祂的圣子，来烧死恶人，而我们这些其他青蛙将同祂一起享受永生。}请看这位可敬的哲学家，如同一个低劣的小丑，将神圣审判、恶人受罚和义人得赏的宣告变成嘲笑、讥讽和笑料；并且在此之后加上一句：\Quote{这样的话如果由虫和青蛙说出，比由彼此争吵的基督徒和犹太人说出更可容忍！}然而，我们不会模仿他的榜样，也不会对那些自称知道万物本性、并彼此讨论万物如何被造、天地及其中的一切如何起源、以及人的灵魂（或是非受生而非由天主所造却受祂治理而从一身体转移到另一身体，或是与身体同时形成而永存或消逝）的哲学家说类似的话。因为，若不尊重并接受那些献身于探求真理之人的意图，而是嘲笑并辱骂地说这些人是\Quote{虫}，他们不以人生中自己粪堆的一角来衡量自己，竟如此重要的问题上发表意见，仿佛他们理解这些事，并竭力声称他们已看到了那些没有更高启示和神圣能力就无法看见的事物。\BibleQuote{因为除了人内里的心神，有谁知道人的事呢？同样，除了天主的神，也没有人知道天主的事。}然而，我们并不疯狂，也不将这样的人类智慧（我用\Quote{智慧}一词的普通含义）——它不忙于大众事务，而是探求真理——比作虫的蠕动或其他任何此类生物；相反，在真理之神中，我们为某些希腊哲学家作证说，他们认识天主，因为\BibleQuote{天主将自己显示给他们，尽管他们不将天主当作天主来光荣，也不感谢祂，反而在思想中变成虚妄，并自称聪明，却成为愚昧，将不朽天主的光荣变成了可朽坏的人、鸟、四足兽和爬虫的像。}\par

% structure: p0062 source=h2 target=subsection reason=relative-heading-rank; base=h2

\subsection{第三十一章}

此後，他想證明猶太人與基督徒以及他先前列舉的那些動物之間沒有區別，便斷言猶太人是\Quote{從埃及逃出的人，從未做過任何值得注意的事，也從未享有任何聲譽或地位}。關於他們並非逃亡者，也非埃及人，而是定居在埃及的希伯來人，我們在前幾頁已經論述過。但如果他認為他的斷言——\Quote{他們從未享有任何聲譽或地位}——能夠成立，是因為希臘人沒有記載他們歷史中的任何重大事件，那麼我們要回答說：如果有人從最初開始審視他們的政體及其法律的安排，就會發現他們是在地上代表天上生活的影子，並且在他們中間，除了那位統管萬有的天主之外，沒有人被認為是神；在他們中間，沒有製造偶像的人被允許享有公民權利。因為在他們的國家中，既沒有畫家，也沒有造像者，法律將這一切人都驅逐出去，以免給製造偶像提供藉口——這種藝術吸引愚人的目光，並將靈魂的眼睛從天主拖向塵世。因此，他們中間有這樣一條法律：\BibleQuote{不可違犯法律，為自己製造雕像，任何男或女的形象；或地上任何生物的形像，或天上飛鳥的形像，或地上爬蟲的形像，或地下水中魚類的形像。}法律確實希望他們顧及每個事物的真理，不製造與現實相反的表現，虛構出真正是男或女的樣子，或動物、飛鳥、爬蟲、魚類的本性。他們的這一禁令也是可敬而偉大的：\BibleQuote{不要舉目向天，免得你看見太陽、月亮、星辰和天上萬象，就受到誘惑去朝拜它們，事奉它們。}此外，整個民族所處的\Emph{régime}是何等的嚴謹，以致任何柔弱的人都不能在公共場合出現；將妓女——這些青年情慾的刺激物——從國家中除去，也是令人欽佩的安排！他們的法庭由最正直的人組成，這些人在長時期內樹立了無可指摘的生活榜樣後，才被委以主持審判的職責，並且由於他們超人的純潔品格，按照猶太人古老的語言習慣，被稱為神。這裡可以看到一個全民從事哲學的景觀；為了有時間聆聽他們的神聖法律，那些被稱為\Quote{安息日}的日子以及在他們中間存在的其他節日被設立。我何必談論他們的司鐸階層和祭獻——這些對那些想要查明事物意義的人包含著無數（更深奧真理的）指示呢？\par

% structure: p0064 source=h2 target=subsection reason=relative-heading-rank; base=h2

\subsection{第三十二章}

但既然沒有什麼屬於人性的東西是永恆的，這個政體也必然逐漸敗壞和改變。天主眷顧，為了適應各國人民，重新塑造了他們可敬的體系中需要改變的部分，並將耶穌的可敬宗教賜給了萬國中所有相信的人，以代替猶太人。耶穌不僅以智慧，而且以分享天主性而裝飾自己，推翻了關於地上魔鬼——那些喜愛乳香、血和祭獻氣味，如同神話中的提坦或巨人，將人從對天主的思念中拖下來——的教義；祂親自蔑視了他們的陰謀（這些陰謀主要針對較好的人類階層），制定了法律，確保那些按照它們生活的人獲得幸福，並且不通過祭獻諂媚魔鬼，而是完全輕視它們，藉助天主聖言的助佑——這聖言幫助那些仰望祂的人。既然天主的旨意是耶穌的教義在人類中盛行，魔鬼雖然竭盡全力要消滅基督徒，卻一無所獲；因為牠們煽動各地的君主、元老院、統治者，甚至民族本身——他們沒有察覺魔鬼的荒唐邪惡行為——敵對聖言和相信聖言的人；然而，比一切事物更有力的天主聖言，即使在遭受反對時，也從反對中獲得了增長的動力，繼續前進，贏得了許多靈魂，這是天主的旨意。我們提出這些話，是作為必要的離題話。因為我們想回答凱爾蘇關於猶太人的斷言——他們是\Quote{從埃及逃出的人，這些受天主喜愛的人從未做過任何值得注意的事}。此外，針對\Quote{他們從未享有任何聲譽或地位}的論述，我們說：他們作為\BibleQuote{被揀選的民族和君尊的司祭團體}分別居住，避免與周圍眾多民族交往，以使他們的道德免受敗壞，享受天主之力的保護，既不貪圖如大多數人那樣奪取其他王國，也不因弱小而被遺棄從而易受他人攻擊而完全毀滅；只要他們值得天主保護，這種情況就一直持續。但當他們作為一個全然陷於罪的民族，需要通過苦難被帶回他們的天主時，他們就被（祂）遺棄，有時長些，有時短些，直到羅馬時期，他們因殺害耶穌而犯了最嚴重的罪，被完全遺棄。\par

% structure: p0066 source=h2 target=subsection reason=relative-heading-rank; base=h2

\subsection{第三十三章}

紧接此后，\Person{色尔苏斯}攻击\Person{梅瑟}首卷书（称为\WorkTitle{创世记}）的内容，断言说：\Quote{犹太人因而试图从第一代变戏法和欺骗者中追溯自己的起源，援引晦暗模糊的话语之见证，其含义隐藏于幽暗之中，而他们又向无知无识者错误地解释这些话语，况且在漫长的以往时期，这样一个论点从未被质疑过。}在我看来，\Person{色尔苏斯}用这些话极为含糊地表达了他想要传达的意思。确实，他在这主题上的含糊可能是故意的，因为他看到确立犹太人从其祖先一脉相承的论证之力量；而另一方面，他又不愿显得不知道关于犹太人及其家系的问题是一个不能轻率处理的问题。然而，确定的是，犹太人将其谱系追溯到三位祖先：\Person{亚巴郎}、\Person{依撒格}和\Person{雅各伯}。这些人的名字与\OriginalTerm{天主}{God}的名字结合时具有如此效力，以至于不仅属于这个民族的人在祈祷\OriginalTerm{天主}{God}和驱逐魔鬼时使用\Quote{\Person{亚巴郎}的\OriginalTerm{天主}{God}，\Person{依撒格}的\OriginalTerm{天主}{God}，\Person{雅各伯}的\OriginalTerm{天主}{God}}这些话，而且几乎所有从事咒语和魔法仪式的人也这样做。因为在许多国家的魔法论著中都发现有这样的对\OriginalTerm{天主}{God}的呼求和神圣名字的采用，表明这些人在与魔鬼交往时熟悉使用它。因此，由犹太人和基督徒提出以证明\Person{亚巴郎}、\Person{依撒格}和\Person{雅各伯}（犹太民族的祖先）之神圣性的这些事实，在我看来并非完全不为\Person{色尔苏斯}所知，而是他没有明确陈述，因为他无法回答可能基于这些事实的论证。\par

% structure: p0068 source=h2 target=subsection reason=relative-heading-rank; base=h2

\subsection{第三十四章}

因为我们询问所有使用这种对\OriginalTerm{天主}{God}的呼求的人，说：\Quote{朋友们，请告诉我们，\Person{亚巴郎}是谁，\Person{依撒格}是什么样的人，\Person{雅各伯}拥有何等能力，以致\Quote{\OriginalTerm{天主}{God}}这个称号，当与他们的名字连用时，能产生这样的奇事？}你们是从哪里学到或能学到关于这些人的事实？又有谁致力于撰写关于他们的历史，要么直接通过归给他们神秘力量来赞扬他们，要么隐约暗示他们拥有某些伟大而奇妙的品质，这些品质对于有资格看见它们的人来说是显而易见的？当没有人能够回答我们的询问，指出这些人的记载是从什么历史（无论是希腊还是蛮族的历史）——或者如果没有历史，至少从什么神秘叙事——中得来的，我们将拿出名为\WorkTitle{创世记}的书，其中包含这些人的事迹以及向他们发出的神圣神谕，并说：\Quote{你们使用这个种族的三位祖先的名字，以最清楚的方式表明，呼求他们产生了不可轻视的效果，难道这不证明这些人的神性吗？}然而，我们除了从犹太人的圣书中知道他们之外，再无其他来源！此外，\Quote{\Person{以色列}的\OriginalTerm{天主}{God}}、\Quote{\OriginalTerm{希伯来人}{Hebrews}的\OriginalTerm{天主}{God}}以及\Quote{在\Place{红海}淹没了埃及王和埃及人的\OriginalTerm{天主}{God}}这些短语是经常用于对付魔鬼和某些邪恶权势的\OriginalTerm{公式}{formulæ}。我们从那些\Person{希伯来人}那里得知这些名字的历史和解释，他们以自己的民族文学和民族语言自豪地讲述这些事情，并解释其含义。那么，犹太人如何能试图从\Person{色尔苏斯}所认为的变戏法和欺骗者的第一代中追溯自己的起源，并无耻地努力把自己和他们的开端追溯到这些人呢？——这些人的名字既然是希伯来语，就向那些拥有用希伯来语文字写成的圣书的希伯来人证明，他们的民族与这些人有亲缘关系。因为直到现在，属于希伯来语的犹太名字要么取自他们的著作，要么一般取自那些意义由希伯来语揭示的词语。\par

% structure: p0070 source=h2 target=subsection reason=relative-heading-rank; base=h2

\subsection{第三十五章}

任何阅读克耳苏斯论文的人，应当观察它是否不传达某种类似上述的暗示，当他这样说时：\Quote{他们试图从那最初的一群江湖骗子和欺骗者中追溯自己的起源，援引隐晦而含混的言辞为证，其意义蒙昧不明。}因为这些名字的确隐晦，不为多数人所理解和认知，虽然在我们看来其意义并非可疑，即使它们被那些与我们宗教无关之人所采用；但按照克耳苏斯的说法，它们不传达任何含混，我不明白他为何拒绝它们。然而，若他诚心想推翻他以为犹太人无耻自夸的谱系——即他们以亚伯拉罕及其后裔为祖先而自夸——他本应引用所有与此相关的段落；首先，用他认为可能具有说服力的论据来支持他的主张，其次，通过他看来是真实的意义和有利于它的论据，勇敢地驳斥此问题上存在的错误。但无论是克耳苏斯还是其他任何人，都无法通过他们对用于奇迹目的之名字性质的讨论，确立关于它们的正确教义，并证明那些仅凭名字就能在同胞甚至外邦人中成就（如此结果）的人应当被轻视。他还应当说明，我们是如何曲解包含这些名字的段落，照他想象欺骗听众，而他本人自夸并非无知或愚钝，却给出了对这些名字的正确解释。在谈到犹太人从中推导其谱系的名字时，他贸然断言：\Quote{在漫长的前期，从未有过关于这些名字的争论，但现在犹太人却与某些其他人争论它们，}而他并未提及这些人。现在，让愿意者指出这些与犹太人争论的人是谁，并提出即使可能成立的论据，证明犹太人和基督徒在这些名字相关的问题上并未正确决定，而是有其他人以极大的学识和准确性讨论过这些问题。但我们确信无人能确立任何此类事情，因为显然这些名字来自希伯来语，而这种语言仅在犹太人中发现。\par

% structure: p0072 source=h2 target=subsection reason=relative-heading-rank; base=h2

\subsection{第三十六章}

克耳苏斯接下来从神圣记录以外的历史中，引出许多民族（如雅典人、埃及人、阿卡迪亚人和弗里吉亚人）声称有极大古老性所依据的段落，他们断言他们中有某些人是从地中生出，并且各自提出这些断言的证据，他说：\Quote{因此，犹太人在巴勒斯坦某角落过着卑微的生活，完全是未受教育之民，未曾听说这些事情早已被赫西俄德和无数其他受默感的人写成诗歌，却编造了一些最难以置信且乏味的故事，即：某个人是由天主的\Emph{手}所形成，并向他吹入了生命的气息，一个女子取自他的肋骨，天主颁布了一些命令，一条蛇违背了这些命令，并战胜了天主的诫命；这样讲述了一些老妇的寓言，并且最不虔敬地描绘了天主在（事物）起初就软弱无力，甚至连祂自己所形成的一个人类都不能说服。}通过这些例子，这位博学多识的克耳苏斯（他指责犹太人和基督徒无知和缺乏教育）清楚显明了他对各历史家（无论是希腊的还是蛮族的）年代学的精确知识，因为他以为赫西俄德和那些被他称为\Quote{受默感}之人的\Quote{无数}他人比梅瑟及其著作更古老——而梅瑟却被证明比特洛伊战争的时代古老得多！因此，不是犹太人编造了关于人类从地而生难以置信且乏味的故事，而是克耳苏斯的这些\Quote{受默感}之人，赫西俄德和他的其他\Quote{无数}同伴，他们未曾学习或听闻在巴勒斯坦存在的远为古老且最可敬的记载，便写下了如他们的\WorkTitle{神谱}之类的历史，尽其所能将\Quote{生育}归于他们的神祇，以及无数其他荒谬之事。这些正是柏拉图从他的\WorkTitle{国家}中驱逐的作家，认为他们是青年的败坏者——即荷马和那些创作了类似诗歌的人！显然，柏拉图并不认为那些留下如此作品的人是\Quote{受默感}的。但也许是出于对我们加添侮辱的愿望，这位享乐主义者克耳苏斯（若这是同一位克耳苏斯还写了另外两本反对基督徒的书）比柏拉图更有判断力，竟然称那些他实际上不视为\Quote{受默感}之人为\Quote{受默感}的人。\par

% structure: p0074 source=h2 target=subsection reason=relative-heading-rank; base=h2

\subsection{第三十七章}

此外，他还指控我们引入\Quote{由\OriginalTerm{天主的手}{hands of God}\OriginalTerm{塑造}{fashioning}的人}，尽管\BibleBook{}{创世纪}在叙述创造或\Quote{\OriginalTerm{塑造}{fashioning}}人时，完全没有提到天主的\Quote{\OriginalTerm{手}{hands}}；反而是\Person{约伯}和\Person{达味}使用了\BibleQuote{祢的手造了我、塑造了我}这样的说法。关于这一点，需要长篇论述，以指出使用这些词语的人如何理解\Quote{\OriginalTerm{创造}{making}}与\Quote{\OriginalTerm{塑造}{fashioning}}之间的区别，以及天主的\Quote{\OriginalTerm{手}{hands}}。因为那些不理解圣经中这些及类似说法的人，认为我们将主宰万有的天主赋予像人一样的形态；按照他们的想法，我们必然认为天主的身体有翅膀，因为按字面理解，圣经将这样的附属物归于天主。然而，当前的主题并不要求我们解释这些说法；因为在我们对\BibleBook{}{创世纪}的注释中，这些事已尽我们所能作了专门的研究。接着请注意\Person{凯尔苏斯}在以下内容中的恶意。圣经论述人的\Quote{\OriginalTerm{塑造}{fashioning}}时说：\BibleQuote{将\OriginalTerm{生命的气息}{breath of life}吹在他的脸上，人就成了一个有灵的生物。}对此，\Person{凯尔苏斯}恶意地嘲笑\Quote{\OriginalTerm{将生命的气息吹在他的脸上}{inbreathing into his face of the breath of life}}的说法，并不理解这一表达所使用的意义，就说\Quote{他们编造了一个故事，说一个人由\OriginalTerm{天主的手}{hands of God}塑造，并被吹入气息所\OriginalTerm{膨胀}{inflated}}，以便利用\Quote{\OriginalTerm{膨胀}{inflated}}一词（与皮肤充气进行类比），来嘲笑\BibleQuote{祂将生命的气息吹在他的脸上}这句话——这些词语是比喻性的，需要解释以表明天主将祂的\OriginalTerm{不朽的圣神}{incorruptible Spirit}赐予人；正如经上所说：\BibleQuote{因为祢的不朽的圣神充满万有。}\BibleRef{智 12:1}\par

% structure: p0076 source=h2 target=subsection reason=relative-heading-rank; base=h2

\subsection{第三十八章}

其次，由于他旨在诽谤我们的圣经，他嘲笑以下叙述：\BibleQuote{天主使\OriginalTerm{沉眠}{deep sleep}落在\Person{亚当}身上，他就睡了；于是取下他的一条\OriginalTerm{肋骨}{rib}，再把肉合起来。天主用从那人身上所取的肋骨，\OriginalTerm{塑造}{fashioned}了一个女人}，等等；他没有引用那些词语，而若引用，听者会认为这些话具有比喻意义。他甚至不愿意承认这些词语是寓言性地使用的，尽管他后来又说，\Quote{犹太人和基督徒中较为谦虚的人对这些事感到羞耻，并试图以某种方式给它们一个\OriginalTerm{寓意解释}{allegorical signification}}。现在我们可以对他说：你所敬奉的\Quote{被灵感激励的}\Person{赫西俄德}所写的关于女人的神话，其中说女人是由\Person{朱庇特}赐给人类作为祸害、作为偷\Quote{火}的报应，这些陈述是否要用寓意来解释？而关于女人是从男人（在他沉眠之后）肋旁取出、并由天主塑造的叙述，在你看来却是毫无理性意义和隐秘含义的吗？但这不是不公平的吗？一方面不嘲笑前者为神话，反而赞赏它们为披着神话外衣的哲学思想；另一方面却蔑视后者，认为它们冒犯理智，并宣称它们毫无价值？因为如果仅仅因为措辞本身，我们就要找那些具有隐秘含义的内容的毛病，那么请你看看下面\Person{赫西俄德}（一个正如你所说的被灵感激励的人）的诗句，它们岂不更令人发笑？——\par

\Quote{\InnerQuote{伊阿珀托斯之子！集云者心怀愤怒地说：\newline{}哦，技艺无与伦比！\newline{}你为取回火焰而得意，\newline{}并因欺骗神明而洋洋得意？\newline{}但你，与人类的后裔，\newline{}必为那欺诈而后悔，更巨大的灾祸由此开始；\newline{}我要为你的盗火送去祸患，\newline{}而所有人都拥抱它，渴望他们的灾厄。}\newline{}那统辖大地、统治苍穹的圣父\newline{}说完，他内心深处充满了笑意。\newline{}他命工艺之神听从他的旨意，\newline{}用水调和可塑的泥土，加以塑造：\newline{}注入生命的气息，赋予形态，\newline{}灵巧的活力，以及人的声音：\newline{}容貌如天上女神般美丽，\newline{}少女的样貌，带着爱的眉宇。\newline{}他命\Person{密涅瓦}教授染色的技艺，\newline{}使布料随着织梭飞舞而色彩斑斓；\newline{}他召唤爱之女神的魔法，在她优雅的头上\newline{}洒下无名的优雅；\newline{}种下那永不满足的渴望，\newline{}以及侵蚀身躯的服饰牵挂：\newline{}最后命\Person{赫尔墨斯}植入精妙的造作，\newline{}娴熟的风度，以及无耻的心思。\newline{}他说；他们的君王下级诸神遵命：\newline{}一位羞怯少女的塑造形象\newline{}从调和泥土中升起，由\Person{朱庇特}的意旨，\newline{}在成型之神的手下；腰带和长袍\newline{}由\Person{密涅瓦}的手系结和折叠：\newline{}天所生之优雅，以及温柔的劝说\newline{}用金链装饰她的四肢：时辰女神们\newline{}用春天的花朵缠绕她的鬓角。\newline{}全套服饰都由\Person{密涅瓦}的精心\newline{}依照身材制作，合体贴切。\newline{}但在她的胸中，那来自上方的使者，\newline{}充满深沉雷鸣的\Person{朱庇特}的谋略，\newline{}造作了狡诈的风度，造作了欺诈的谎言，\newline{}以及使热血沸腾、让智者迷惑的言语。\newline{}这位众神的传令官宣布了她的名字，\newline{}称那女人为\Person{潘多拉}；\newline{}因为众神都献上礼物以迷惑\newline{}人类善于创造的后裔。}\par

此外，关于那箱子的话本身也足以引人发笑；例如：——\par

\Quote{古时大地上，人类子孙居住在\newline{}远离灾祸之处，远离苦辛的重担\newline{}和使人衰老的沉重疾病；\newline{}如今凡人的悲哀生命不过一瞬。\newline{}女人的双手托着一个巨大的箱子；\newline{}她抬起盖子；将忧愁散入空中：\newline{}唯独，在器皿的边沿下被扣留，\newline{}\OriginalTerm{希望}{Hope}仍留在未破损的牢房里，\newline{}并没有逃到外面；这就是集云者\Person{朱庇特}的意愿：\newline{}女人的手曾将盖子盖上。}\par

现在，若有人要给这些文字赋予一种深奥的寓意（无论是否真有这种寓意），我们就要说：难道只有希腊人才有自由，把哲学含义隐藏在秘密的外衣下？或者也许埃及人，以及那些以自己的奥秘和其中隐藏的真理而自豪的蛮族人也行？而唯独犹太人，连同他们的立法者和史学家，在你看来却是最愚昧的人？难道只有这个民族没有得到神圣力量的份，却还被如此宏伟地教导，要上升到天主那非受造的本性，只仰望祂，并只从祂那里期待（实现）他们的希望？\par

% structure: p0082 source=h2 target=subsection reason=relative-heading-rank; base=h2

\subsection{第三十九章}

但是，刻尔苏斯还嘲笑蛇破坏了天主给那人的诫命，把这段叙述看作老太婆的寓言，而且他故意不提及天主的乐园，也不说天主据说在伊甸东边栽种了乐园，随后地上长出各种悦人眼目、好作食物的树木，以及乐园中央的生命树和知识善恶树，以及其他随后陈述的内容——这些本身就能使公正的读者看出，这一切并非不恰当地具有寓意。那么，让我们把苏格拉底在柏拉图\WorkTitle{会饮篇}中关于厄洛斯的言论与之对比，这些话借苏格拉底之口说出，被认为比会饮中所有其他人关于厄洛斯的言论更为恰当。柏拉图的话如下：\Quote{当阿佛洛狄忒降生时，众神举行宴会，墨提斯之子波罗斯也和其他人一同在场。他们用餐后，珀尼亚前来乞讨（因为看到有宴席），她站在门口。与此同时，波罗斯因喝多了琼浆（因为那时还没有酒）而走进宙斯的花园，醉醺醺地躺下睡着了。于是珀尼亚生出一个秘密计谋，想摆脱自己的贫困处境，与波罗斯生个孩子，便躺在他身边，怀上了厄洛斯。因此，厄洛斯成了阿佛洛狄忒的随从和侍者，因为他在她的生日宴会上被怀上，同时本性爱美，因为阿佛洛狄忒也是美的。既然厄洛斯是波罗斯和珀尼亚的儿子，他的状况如下。首先，他总是贫穷，远不像大多数人想象的那样娇嫩美丽，而是干枯、晒黑、赤脚、无家，总是睡在地上，没有遮盖；露宿在门边和公共道路上；拥有母亲的本性，常与匮乏为伴。但另一方面，按照父亲的特征，他善于谋算美和善的事物，勇敢、急躁、激烈；是个敏锐的猎人，不断想出计谋；既多思虑，又智谋丰富；终生像哲学家般行事；可怕的巫师、药剂师和诡辩家；既非不死之神也非凡人，但在同一天内，有时富足时便兴旺生活，另有时死亡，又因拥有父亲的本性而再次复活。然而，他所得到的供应总是渐渐消失，以致他从不缺乏，也从不富有；另一方面，他处于智慧和无知之间的中间位置。}现在，若有人读这些话时模仿刻尔苏斯的恶毒——愿基督徒绝不做这事！——他们就会嘲笑这神话，并把这位伟大的柏拉图变成笑柄；但若以哲学精神研究以神话形式传达的内容，能够发现柏拉图的意思，他们就会钦佩他能够为了众人而将出现在他眼前的伟大思想以这种神话形式隐藏起来，并以适宜的方式对那些懂得如何从神话中推断出编织神话者真正含义的人说话。我之所以提出柏拉图著作中的这个神话，是因为其中提到宙斯的花园，这似乎与天主的乐园有某种相似之处，还提到珀尼亚与蛇的比较，以及珀尼亚对波罗斯的计谋，这与蛇对人的计谋可作比较。确实，不太清楚柏拉图是偶然遇到这些故事，还是像有些人认为的那样，他在访问埃及时遇到一些对犹太奥秘进行哲学思考的人，并从他们那里学到一些东西，因而保留了一些观念，又抛弃了其他一些，以免因完全采纳犹太哲学的要点而冒犯希腊人——犹太人的法律和独特的政体因其外邦性质而在众人中声名狼藉。然而，现在并非解释柏拉图神话或蛇与天主乐园的故事以及其中发生的一切的适当时机，因为在一本注释\BibleBook{}{创世纪}的书中最已尽力处理了这些问题。\par

% structure: p0084 source=h2 target=subsection reason=relative-heading-rank; base=h2

\subsection{第四十章}

但他断言\Quote{梅瑟的记叙最不虔敬地描述了天主从（万物的）起初就处于软弱状态，并且不能使祂自己所造的哪怕一个人顺服（于祂）}，我们回答说，这个异议与以下情况十分相似：即有人指责罪恶的存在，而天主甚至未能阻止它出现在一个人身上，以便从万物的起初就可以找到一个生来没有沾染罪恶的人。因为正如那些辩护天主眷顾论的人用不容轻视的论据来论证一样，\Person{亚当}及其儿子的问题也将会被那些知道在希伯来语中\Person{亚当}意指人、并且知道在叙述中那些似乎指个别\Person{亚当}的部分里，\Person{梅瑟}是在论述一般人性的人以哲学方式处理。因为\BibleQuote{在亚当内}（如圣经所说）\BibleQuote{众人都死了}，并因\Person{亚当}的相似悖逆而被定罪，天主的话如此断言，与其说是关于\InnerQuote{一个特定的个人}，不如说是关于\InnerQuote{全人类}。因为在一系列似乎适用于一个特定个人的陈述中，对\Person{亚当}的诅咒被视为对（人类）所有成员都是共同的，而针对女人所说的话也是毫无例外地对\InnerQuote{每一个}女人说的。男人和女人被逐出乐园，并穿上了皮制长衣（这是天主因人犯罪而为犯罪者做的），这些包含了一个隐秘而神秘的教导（远超过\Person{柏拉图}的教导），即灵魂失去翅膀，被带到地上，直到它能找到一个稳固的栖息之处。\par

% structure: p0086 source=h2 target=subsection reason=relative-heading-rank; base=h2

\subsection{第四十一章}

此后，他接着说道：\Quote{其次，他们谈论一场洪水，一只巨大的方舟，里面装着一切，以及一只鸽子和一只乌鸦作为信使，他们篡改并鲁莽地歪曲了\Person{丢卡利翁}的故事；我猜想，他们不指望这些事会暴露，而是以为他们只是在为年幼的孩子编造故事。}在这些评论中，请注意此人——如此不符合哲学家身份——对这段非常古老的犹太叙述所表现出的敌意。因为，他无法对洪水的历史提出任何反驳，也看不出他本可以就方舟及其尺寸提出什么异议——即，根据普遍意见，接受方舟长三百肘、宽五十肘、高三十肘的说法，不可能维持它容纳了地上所有动物，每样洁净的七对和不洁净的四对——他仅仅称它为\Quote{巨大的，里面装着一切}。那么，它哪里\Quote{巨大}了？因为据记载它建造了一百年，长三百肘、宽五十肘，直到高三十肘的顶部收缩为一肘长一肘宽。我们为什么不应该更欣赏一座像广阔城市一样的建筑呢？如果它的尺寸按其可能的意义来理解，那么底面长九万肘、宽二千五百肘？我们为什么不应该赞叹其设计之巧妙，它建得如此紧凑，能够承受导致洪水的暴风雨？因为它不是用沥青或类似材料涂抹，而是安全地涂上了柏油。难道不值得赞叹吗？由于天主的眷顾安排，所有族类的元素都被带入其中，以便大地可以再次接受所有生物的种子，而天主使用了一个最正义的人作为洪水后出生之人的祖先。\par

% structure: p0088 source=h2 target=subsection reason=relative-heading-rank; base=h2

\subsection{第四十二章}

为了表明他读过\BibleBook{}{创世纪}，\Person{策尔修}拒绝了关于鸽子的故事，尽管他未能提出任何可以证明它是虚构的理由。接着，按照他的习惯，为了使叙述显得更可笑，他把\Quote{乌鸦}改成了\Quote{渡鸦}，并想象\Person{梅瑟}是这样写的，他鲁莽地歪曲了关于希腊人\Person{丢卡利翁}的传说；除非他认为这叙述并非来自\Person{梅瑟}，而是来自\InnerQuote{几个人}，这从他使用复数形式\Quote{篡改并鲁莽地歪曲了\Person{丢卡利翁}的故事}以及\Quote{因为\InnerQuote{他们}不指望，我猜想，这些事会暴露}可以看出。但是，那些把自己的圣经交给\InnerQuote{整个}民族的人，怎么会不指望它们会暴露呢？而且他们还预言这宗教将传到\InnerQuote{所有}民族。\Person{耶稣}宣布：\BibleQuote{天主的国必从你们手中夺去，赐给那结果子的民族}；祂对犹太人说出这些话时，还有什么别的意思呢？不就是说，祂自己要借着祂的神能，将包含天主之国奥秘的整部犹太圣经揭示出来吗？那么，当他们阅读希腊人的\WorkTitle{神谱}和关于十二神的故事时，他们赋予它们一种庄严感，给它们披上寓意的外衣；但当他们想要蔑视我们的圣经叙述时，就断言它们是笨拙地为幼童编造的神话！\par

% structure: p0090 source=h2 target=subsection reason=relative-heading-rank; base=h2

\subsection{第四十三章}

\Quote{总而言之，荒谬而不合时宜的}，他继续说，\Quote{是（关于）生子的（叙述）}，这里虽然他未提任何名字，显然是指\Person{亚巴郎}和\Person{撒辣}的历史。他亦讥讽\Quote{兄弟们的阴谋}，或指\Person{加音}谋害\Person{亚伯尔}的故事，或再加上\Person{厄撒乌}谋害\Person{雅各伯}；谈到\Quote{父亲的悲痛}，大概是指\Person{依撒格}因\Person{雅各伯}离去而悲伤，也可能指\Person{雅各伯}因\Person{若瑟}被卖往\Place{埃及}而悲伤。在叙述\Quote{母亲的狡猾计谋}时，我想他是指\Person{黎贝加}的做法，她设法使\Person{依撒格}的祝福不落在\Person{厄撒乌}身上，而落在\Person{雅各伯}身上。如果我们断言在所有这些事例中，\Person{天主}以非常显著的方式介入，我们犯了什么荒谬呢？因为我们深信，祂从不撤回对那些以可敬而有力的生活献身于祂的人的眷顾。此外，他嘲笑\Person{雅各伯}在与\Person{拉班}同住时获得的财产，不明白这些话是指什么：\Quote{凡没有斑点的都是\Person{拉班}的，凡有斑点的是\Person{雅各伯}的}；他说\Quote{\Person{天主}将驴、羊、骆驼赐给他的儿子们}，却没有看到\BibleQuote{这一切事发生在他们身上，都是为鉴戒，并且写在经上，是为我们这末世的人}\BibleRef{格前 10:11}。各民族中不同的风俗变得闻名，由\Person{天主}的言语所规范，作为产业赐给那借喻预表称为\Person{雅各伯}的人。因为从外邦人中皈依\Person{基督}的人，在\Person{拉班}和\Person{雅各伯}的历史中已有预示。\par

% structure: p0092 source=h2 target=subsection reason=relative-heading-rank; base=h2

\subsection{第四十四章}

他完全误解了圣经的意义，说\Quote{\Person{天主}也赐给义人水井}。他没有察觉义人不挖蓄水池，而是掘井，试图发现可饮之福的内在根基和源泉，因为他们以比喻的方式领受那命令：\BibleQuote{从你自己的器皿中喝水，从你自己的活水井中取饮。不可让你的水溢出你的泉源，而要流到你的街市。只属于你一人，不可让外人与你同享。}\BibleRef{箴 5:15-17}圣经常常利用真实事件的历史，来呈现更重要的真理，这些真理只是隐约暗示的；这类叙述涉及\Quote{水井}、\Quote{婚姻}以及义人\Quote{性交}的各种行为，但更适宜在解释那些段落本身的注释性著作中加以说明。然而，义人在\Place{培肋舍特人}境内挖掘水井，正如\WorkTitle{创世纪}所载，这从\Place{阿市刻隆}所显示的奇妙水井可以明显看出，这些水井因其结构与其他水井迥异而值得一提。此外，少女和婢女都应从隐喻来理解，这不仅是\Emph{我们的}教义，也是我们从起初从智者领受的，其中一位智者鼓励听者探究比喻意义时说：\BibleQuote{请告诉我，你们这些读法律的人，难道不听取法律吗？因为经上记载：\Person{亚巴郎}有两个儿子，一个由婢女所生，一个由自由妇人所生。但那由婢女所生的，是按血肉生的；那由自由妇人所生的，却是因恩许生的。这些都是寓意：是指两个盟约，一个是出于\Place{西乃山}，生子为奴，那就是\Person{哈加尔}。}\BibleRef{迦 4:21-24}稍后又说：\BibleQuote{但那在上面的\Place{耶路撒冷}是自由的，她是我们的母亲。}\BibleRef{迦 4:26}任何人拿起\WorkTitle{迦拉达书}都可以学到，关于\Quote{婚姻}和与\Quote{婢女}交合的段落是如何被寓意化了的；圣经希望我们效法的，不是那些行此事之人的字面行为，而是（正如耶稣的宗徒惯于称之为）属灵的行为。\par

% structure: p0094 source=h2 target=subsection reason=relative-heading-rank; base=h2

\subsection{第四十五章}

\Person{塞耳苏}本应认识到\WorkTitle{圣经}作者所显示的爱真理之心，他们甚至没有隐瞒有损自己名声的事，从而被引导去接受其他更神奇的故事为真，但他却反其道而行之，把\Person{罗特}和他女儿们的故事（既不考察其字面意义，也不考察其寓意）说成是\Quote{比\Person{提厄斯忒斯}的罪行更恶劣}。目前无需解释这段历史的寓意，也无需解释\Place{索多玛}的含义，以及天使对从那里逃离的人所说的话：\BibleQuote{不要回头看，也不要停留在整个周围地区；逃到山上去，免得你被消灭}；也无需解释\Person{罗特}和他妻子的含义，她因回头而变成了盐柱；以及他的女儿们灌醉父亲，以便从他怀孕的含义。但让我们用几句话缓和一下这段历史令人反感的地方。希腊人研究过行为的性质——善、恶、无所谓——并且其中较成功的研究者提出原则：只有意图赋予行为善或恶的性质，而没有目的的行为严格来说是\OriginalTerm{无所谓}{indifferent}的；当意图指向一个恰当的目标时，行为值得称赞；反之则应受谴责。因此，他们在涉及\Quote{无所谓之事}的部分中说，严格来说，一个男人与自己女儿发生性关系是无所谓之事，尽管这种事情在已建立的社群中不应发生。为了假设，为了说明这种行为属于无所谓一类，他们设想一个智慧的人与唯一的女儿被留下，而全人类都已毁灭；他们提出问题：父亲是否适合与女儿发生关系，以便根据假设防止人类灭绝。这在希腊人中被视为合理推理，并被有影响力的斯多葛学派所赞扬；但两个年轻女子，听说世界将被焚毁（尽管不理解其全部含义），看见火焰摧毁了她们的城市和国家，并认为重燃人类生命火焰的唯一途径在于她们的父亲和自己，在这种假设下，产生出希望世界继续存在的欲望，她们的行为难道就比那位智慧的人根据斯多葛学派假设在所有人已被消灭的情况下与女儿发生关系的行为更恶劣吗？然而我并非不知道，有些人曾因\Person{罗特}女儿们的欲望而反感，认为她们的行为非常邪恶；并说两个被诅咒的民族——\Place{摩阿布}和\Place{阿孟}——就是从那种不圣洁的关系中产生的。但确实，\WorkTitle{圣经}从未明确认可她们的行为为善，也未明确判定她们的行为应受谴责。不过，无论真实情况如何，它不仅允许一种寓意解释，而且在自身价值上也可以得到辩护。\par

% structure: p0096 source=h2 target=subsection reason=relative-heading-rank; base=h2

\subsection{第四十六章}

\Person{塞耳苏}还对\Person{厄撒乌}对\Person{雅各伯}的\Quote{仇恨}（我猜他是引用这个）加以嘲讽，尽管根据\WorkTitle{圣经}，\Person{厄撒乌}被公认是个恶人；他又不清楚说明\Person{西默盎}和\Person{肋未}的故事（他们因妹妹被\Place{舍根}王子玷污而袭击舍根人），并谴责他们的行为。接着他说到\Quote{兄弟出卖（彼此）}，指的是\Person{雅各伯}的儿子们；\Quote{一个兄弟被卖}，即\Person{若瑟}；\Quote{一个父亲被欺骗}，即\Person{雅各伯}，因为他看见\Person{若瑟}的彩色长衣时没有怀疑儿子们，相信了他们的话，哀悼成为埃及奴隶的儿子，好像他已死。请注意塞耳苏是以何等仇恨和说谎的态度收集\WorkTitle{圣经}的记载；只要哪里在他看来有可指责的理由，他就引用那段经文；但哪里有值得一提的美德——比如\Person{若瑟}不满足女主人的情欲，拒绝她的诱惑和威胁——他甚至不提这件事！他应当看到，\Person{若瑟}的行为远比有关\Person{贝勒罗丰}的记载更为高尚，因为\Person{若瑟}宁愿被关进监狱也不违背他的美德。尽管他可以针对指控者提出公正的辩护，他却宽宏大量地保持沉默，把他的案件托付给天主。\par

% structure: p0098 source=h2 target=subsection reason=relative-heading-rank; base=h2

\subsection{第四十七章}

随后，凯尔苏斯大约出于形式，且极为不准确地说到\Quote{司酒长和司膳长的梦，以及法郎的梦，并对它们的解释，因此若瑟被提出监牢，以便法郎将埃及的次位托付给他。}那么，这段历史本身即使不看，又有什么荒谬之处，竟被凯尔苏斯引为指控的理由呢？他将自己的著作题为\WorkTitle{真确言论}，但其中并不包含教义，而是充满对犹太人和基督徒的指控。他又说：\Quote{曾被出卖的人，对那些出卖他的兄弟们，当他们因饥饿受苦，并带着驴子来购买粮食时，善待了他们。}尽管他并未在著作中叙述这些事。但他确实提到了若瑟向兄弟们表明身份的事，虽然我不知道他出于什么目的，或能从中指出什么荒谬之处；因为连默莫斯本人，我们可以说，也无法对这样的事件——撇开其比喻意义不谈——本身如此引人入胜，提出任何合理的指责。他还进一步叙述：\Quote{曾被卖为奴的若瑟恢复了自由，并庄严地列队前往父亲葬礼，}并认为这段叙述为我们提供了指控的理由，因为他如此评论：\Quote{藉着若瑟，那著名而神圣的犹太民族，在埃及成长为一群民众后，被命令住在王国边界之外的某个地方，并在不名誉的地区牧放他们的羊群。}其中\Quote{被命令在不名誉的地区牧放他们的羊群}这句话，是他因仇恨而添加的；因为他没有指出埃及的哥珊地是不名誉的地方。民众离开埃及，他称之为逃跑，完全忘了出谷纪所记载的希伯来人离开埃及地的事。我们列举这些例子，是为了表明，表面上看来可能提供指控理由的叙述，凯尔苏斯并未成功证明其可厌恶或愚蠢，他完全未能确立他眼中我们圣经的邪恶特性。\par

% structure: p0100 source=h2 target=subsection reason=relative-heading-rank; base=h2

\subsection{第四十八章}

接着，仿佛他全身心致力于显示对犹太和基督教义的仇恨与厌恶，他说：\Quote{犹太人和基督徒中较为温和的作者，都赋予所有这些事一种寓意意义。}又说：\Quote{因为他们对这些事感到羞耻，便逃向寓意。}对此可以回答他说：如果我们必须承认神话和虚构，无论是出于隐秘含义还是其他目的，按其字面理解都是可耻的叙述，那么还有什么比希腊人的历史更应被如此看待呢？在这些历史中，儿子之神阉割父亲之神，父母之神吞食自己的儿女，一位女神母亲将一块石头递给\Quote{众神与人类之父}代替自己的儿子吞下，一位父亲与女儿同寝，一位妻子捆绑自己的丈夫，以被缚之神的兄弟和自己的女儿为盟友！但何必列举希腊人关于他们神祇的这些荒谬故事呢？这些故事本身极为可耻，即使赋予寓意意义也是如此。（例如）索利的克利西波，他被视为斯多葛学派的瑰宝，因其众多博学的论著，他解释萨摩斯岛上的一幅画，画中朱诺与朱庇特犯下不可言说的丑行。这位可敬的哲学家在其论著中说，物质接受神的\OriginalTerm{精言}{spermatic words}并将其保留在自身内，以装饰宇宙。因为在萨摩斯岛的那幅画中，朱诺代表物质，朱庇特代表神。正是由于这些和无数其他类似的神话，我们连在言语上也不愿称万有的天主为朱庇特，或称太阳为阿波罗，或称月亮为狄安娜。我们向造物主献上纯洁的敬拜，并虔敬地谈论祂的杰作创世，甚至在言语上也不玷污天主的事；我们赞同柏拉图在\WorkTitle{斐莱布篇}中的话，他不愿承认快乐是一位女神，他说：\Quote{普罗塔尔科斯，我对神祇的名号如此敬畏。}我们确实对天主的圣名及其杰作创世怀有这样敬畏，以至于我们不愿在寓意意义的借口下，接纳任何可能伤害青年的神话。\par

% structure: p0102 source=h2 target=subsection reason=relative-heading-rank; base=h2

\subsection{第四十九章}

如果\Person{刻尔苏斯}以公正的精神阅读圣经，他就不会说\Quote{我们的著作不能容纳寓意。}因为从先知书（其中记载了历史事件）——而不是从历史书——可以确信，历史部分也是带着寓意目的写成的，并且被极为巧妙地改编，不仅适合众多的简单信徒，也适合少数能够或愿意以明智精神探究事情的人。如果当今那些被\Person{刻尔苏斯}视为\Quote{犹太人和基督徒中比较谦逊的}作者是我们圣经的（最初）寓意解经者，那么他的指控或许还有表面上的合理性。但是，既然教义本身的始祖和作者们都赋予它们寓意意义，那么还能得出什么其他结论，除了它们被撰写成以寓意方式理解其主要意义之外？我们将从极多的例子中举出几个，以显示\Person{刻尔苏斯}对圣经提出了空洞的指控，当他说\Quote{它们不能容纳寓意}时。耶稣的宗徒\Person{保禄}说：\Quote{经上记着说：\BibleQuote{牛在场上踹谷的时候，不可笼住牠的嘴。}难道天主所挂念的是牛吗？还是祂完全是为我们说的呢？的确，这是为我们写的，因为耕种的应怀着希望耕种，打场的应怀着希望分享打场的出产。}在另一处，同一\Person{保禄}说：\Quote{因为经上记着说：\BibleQuote{为此，人要离开自己的父母，与自己的妻子结合，二人成为一体。}这是一个伟大的奥秘；但我是指着基督和教会说的。}又在另一处说：\Quote{我们知道我们的祖先都曾在云柱下，都曾从海中经过，都曾在云中和海中受了洗归于梅瑟。}然后，他解释关于玛纳的历史以及关于从磐石中出水的奇迹，接着这样说：\Quote{他们都吃过同样的神粮，都饮过同样的神饮；因为他们所饮的，是出自那随着他们的神磐石；那磐石就是基督。}此外，\Person{亚撒夫}为了表明\WorkTitle{出谷纪}和\WorkTitle{户籍纪}中的历史充满了难解之处和比喻，开始以如下方式，如\WorkTitle{圣咏集}中所记载的，当他正要提及这些事时：\Quote{\BibleQuote{我的百姓，请听我的指教，请侧耳听我口中的言语。我要开口说比喻，我要说出从古以来的谜语，就是我们所听见、所知道的，我们的祖先告诉我们的。}}\par

% structure: p0104 source=h2 target=subsection reason=relative-heading-rank; base=h2

\subsection{第五十章}

而且，如果梅瑟法律没有任何需要理解为隐秘意义的内容，先知就不会在他的祈祷中向天主说：\Quote{\BibleQuote{求你开启我的眼睛，我要观察你法律中的奇事。}}然而他知道，有一种无知的蒙蔽压在那些阅读却不理解喻意意义的人的心上，这蒙蔽是由天主的恩赐除去，当他垂听那已经尽己所能、并且因习惯而感官操练得以辨别善恶、并不断发出祈祷的人说：\Quote{\BibleQuote{求你开启我的眼睛，我要观察你法律中的奇事。}}有谁在读到时，不立刻追问那用臭粪便充满埃及山岭的是谁，埃及的山岭是什么，埃及的河流是什么——上述法郎曾夸口说：\Quote{\BibleQuote{河流是我的，是我造了它们。}}——以及巨鳄是谁，它鳞甲中的鱼是谁？——（这一切）都要与对河流的解释相协调？但是，何必更详尽地论证已经不需要证明的事呢？因为以下这些话正适用于这些事：\Quote{\BibleQuote{谁明智，就会明白这些事？谁聪敏，就会知道它们？}}现在我花了很多篇幅讨论这个问题，因为我希望表明\Person{刻尔苏斯}这一断言的不正确性，他说：\Quote{犹太人和基督徒中比较谦逊的设法以某种方式给予这些故事一个寓意解释，尽管其中一些故事并不允许这样解释，反而是一些极其愚蠢的编造。}希腊人的故事不仅是极其愚蠢的编造，而且是极其不虔敬的编造。因为我们的叙事明确考虑到众多简单的信徒，而编造希腊神话的人却没有这样做。因此，\Person{柏拉图}不无道理地将所有这类神话和诗歌从他的城邦中驱逐出去，正如我们所讨论的那些一样。\par

% structure: p0106 source=h2 target=subsection reason=relative-heading-rank; base=h2

\subsection{第五十一章}

\Person{瑟尔苏}似乎听说过有些著作包含对\Person{梅瑟}法律的寓意解释。然而，他并未读过这些著作；因为如果他读过，就不会说：\Quote{然而，那些被编造出来的寓意解释，比寓言本身更加可耻和荒谬，因为它们试图以令人惊异且完全无知的愚蠢，将根本无法调和的事物结合起来。}他这些话似乎是指\Person{斐罗}的著作，或更早的作者如\Person{亚里斯托布鲁}的著作。但我猜想\Person{瑟尔苏}并未读过这些书，因为在我看来，他们在许多章节中如此成功地把握了（神圣作者）的意义，以至于连希腊哲学家也会被他们的解释所吸引；因为在他们的著作中，我们不仅发现优美的文体，还有精妙的思想和教义，以及对\Person{瑟尔苏}视为圣经中寓言的东西所做的合理运用。此外，我知道毕达哥拉斯派的\Person{努美纽斯}——一位卓越的\Person{柏拉图}解释者，\Person{毕达哥拉斯}学说的首席教师——在他的许多作品中引用\Person{梅瑟}和先知们的著作，并对所论及的段落给出并非不可能的寓意解释，例如在他的著作\WorkTitle{鹳鸟}以及论及\WorkTitle{数字}和\WorkTitle{地方}的作品中。在他的论文\WorkTitle{善}第三卷中，他还引用了一段关于\Person{耶稣}的叙述——但未提及祂的名字——并赋予它一种寓意意义，无论成功与否，我将另择时机说明。他还叙述了关于\Person{梅瑟}、\Person{雅乃斯}和\Person{雅姆勒斯}的记载。但我们并不因此沾沾自喜，尽管我们赞同\Person{努美纽斯}而非\Person{瑟尔苏}及其他希腊人，因为他出于求知欲愿意探究我们的历史，并（恰当地）将它们视为有待寓意理解的叙述，而非愚蠢的作品之列。\par

% structure: p0108 source=h2 target=subsection reason=relative-heading-rank; base=h2

\subsection{第五十二章}

此后，他从所有包含寓意解释和注释的著作中，选出那些以不可轻视的语言和风格表达的、最不重要的著作，这些著作确实有助于增强众多单纯信徒的信德，却不适合给那些更聪慧的人留下印象，他继续说：\Quote{我知道有这样一部著作，名为\WorkTitle{\Person{帕皮斯库斯}与\Person{雅松}之间的论战}，它更适合激起怜悯和仇恨，而非欢笑。然而，我并非要反驳这些作品中的陈述；因为它们的谬误对所有人都是显而易见的，尤其是如果有人愿意耐心阅读这些书本身。我更想表明，自然教导我们：天主没有创造任何可朽之物；祂的作品，无论是什么，都是不朽的，而它们的作品却是可朽的。灵魂是天主的作品，而身体的本性是不同的。在这方面，蝙蝠、蠕虫或青蛙的身体与人的身体没有区别；因为质料相同，它们可朽的部分也相似。}尽管如此，我希望每个听到\Person{瑟尔苏}宣称并断言题为\WorkTitle{\Person{雅松}与\Person{帕皮斯库斯}关于基督的论战}的著作更适合激起仇恨而非欢笑的人，都能拿起这部作品，耐心聆听其中的内容；这样，发现其中并无引发仇恨之物，他便能根据书本身谴责\Person{瑟尔苏}。因为如果公正地研读，就会发现一部描述基督徒与犹太人讨论犹太圣经，并证明关于基督的预言完全适用于耶稣的作品中，毫无引发欢笑之物；尽管另一位论辩者以不失尊严的方式，并以符合犹太人身份的态度维持着讨论。\par

% structure: p0110 source=h2 target=subsection reason=relative-heading-rank; base=h2

\subsection{第五十三章}

我实在不知道，他是如何把那些不可能结合、也无法同时存在于人性中的事物联结起来的，正如他所说：\Quote{上述著作理应既被怜悯又被仇恨。}因为每个人都会承认，被怜悯的对象不会同时被仇恨，被仇恨的对象也不会同时被怜悯。此外，\Person{瑟尔苏}说，他的目的并非反驳这类陈述，因为他认为其荒谬对所有人显而易见，甚至在提出任何逻辑反驳之前，它们就显得恶劣，并值得怜悯和仇恨。但我们邀请每一位阅读我们这篇回应\Person{瑟尔苏}指控的人，耐心聆听我们圣经本身的内容，并尽可能根据它们的\Emph{内容}对作者的目的、他们的良心和心态作出判断；因为他会发现，他们是竭力维护自己所持守的人，其中一些人表明，他们所叙述的历史是他们亲眼所见、亲身经历的，是奇妙的，值得为未来听众的利益而记录。难道有人敢说，相信万有的天主，并在一切事上以取悦祂为目的而行动，甚至不起任何不悦于祂的念头——因为不仅我们的言语和行为，连我们的思想都将是将来审判的对象——这不是（对人）一切祝福的源泉吗？还有什么论证比这样的信念或意见更能有效地引导人性度修德的生活，即至高天主洞悉一切，不仅是我们所说所做的，甚至是我们所思想的？让任何愿意的人比较一下其他任何体系，这体系同时改造和改善的不仅是一两个人，而是尽其所能地改造和改善无数人，以便通过两种方法的比较，他能对导向修德生活的论证形成正确的看法。\par

% structure: p0112 source=h2 target=subsection reason=relative-heading-rank; base=h2

\subsection{第五十四章}

但在从克耳苏斯的话（即我从其引用的一段，那是 \WorkTitle{蒂迈欧篇} 的意译）中，出现了某些表述，如：\Quote{天主没有创造任何可朽之物，而只创造了不朽之物；可朽之物是其他神祇的作品；灵魂是天主的作品，但身体的本性不同，人的身体与蝙蝠、蠕虫或青蛙的身体没有区别，因为质料相同，它们的可朽部分也一样。}——让我们对此稍作讨论；并让我们表明，克耳苏斯要么没有显露出他的伊壁鸠鲁式观点，要么（如某人可能说的）他已经将其换成了更好的观点，要么（如另一人可能说的）除了名字之外，与伊壁鸠鲁学派的克耳苏斯毫无共同之处。因为他本应在表达这些观点，并打算不仅反驳我们，也反驳并非默默无闻的、基提翁的芝诺的信徒那个哲学派别时，证明动物的身体不是天主的作品，并且构造中展现的巨大技巧并非来自至高的理智。他还应对植物的无数多样性——它们受一种内在的、不可理解的本性支配，并为一般人类以及为人类服务的动物提供了绝非可鄙的用途（无论可能提出什么其他理由来解释它们的存在）——不仅陈述他的看法，还应向我们表明，并非完善的理智将这些品质印在植物的质料上。而且，当他曾将（各种）神祇描绘为所有身体的创造者，而只有灵魂是天主的作品时，为什么他——将这些伟大的创造行为分开，并分配给多个创造者——接下来没有用某种令人信服的理由证明神祇中存在这些多样性，其中一些构造人的身体，另一些构造（比如说）驮兽的身体，还有一些构造野兽的身体呢？他既看到某些神祇是龙、毒蛇和蛇怪的创造者，而另一些是每种植物和草本按种类的创造者，便应当解释这些多样性的原因。因为很可能，倘若他仔细考察每个细节，他要么会注意到是一位天主创造了万物，并按一定目的和一定理由造了每一样东西；要么，若他未能注意到这一点，他本会发现他应当对那些主张可朽性在本性上是无差异的人作出怎样的回答；并且，由一位工匠——他为了整体的益处而构造了不同种类的东西——形成一个由不同材料构成的世界，这并无荒谬之处。最后，他根本不应就如此重要的教义发表任何意见，因为他并不打算证明他所声称要证明的东西；除非，这个责备他人因简单信仰而相信的人，竟要我们相信他的纯然断言，尽管他宣称自己不仅仅断言，还要证明其断言。\par

% structure: p0114 source=h2 target=subsection reason=relative-heading-rank; base=h2

\subsection{第五十五章}

但我主张，若他有耐心（用他自己的话）去听梅瑟和先知们的著作，他便会因以下情况而引起注意：\BibleQuote{天主造了}这一表述用于天地和所谓穹苍，也用于光体和星辰；此后，用于大鱼和水中滋生的一切有生命的爬物，各从其类，以及天空的飞鸟，各从其类；又用于地上的走兽，各从其类，牲畜，各从其类，和地上的一切爬物，各从其类；最后用于人。然而，\BibleQuote{造}这一表述并未用于其他事物；而论到光，只说：\BibleQuote{就有了光。}论到诸水（普天之下的诸水）聚在一处，说：\BibleQuote{就这样成了。}同样，论到地上的草木，经上说：\BibleQuote{地就生出了青草，结种子的蔬菜，各从其类，并结果子的树木，果子内含有核，各从其类，在地上。}他还应探究：经上所记的、关于世界各部分得以产生的天主命令，是对一个对象说的，还是对多个对象说的？并且他不会轻易指责这些书卷中——无论是梅瑟，还是（如我们所说的）在梅瑟内说话的圣神（梅瑟也从祂获得了预言的权能）——所讲述的记述为不可理解、无隐秘意义；因为圣神\Quote{知晓现在、未来和过去，}其程度远超诗人所说的那些祭司所号称拥有的知识。\par

% structure: p0116 source=h2 target=subsection reason=relative-heading-rank; base=h2

\subsection{第五十六章}

再者，既然塞尔苏斯断言\Quote{灵魂是天主的化工，但身体的本质是不同的；而且在这方面，蝙蝠、蠕虫或青蛙的身体与人的身体并无区别，因为物质是相同的，它们可朽坏的部分也相似}，——我们对他的这一论证必须作如下回应：如果，既然相同的物质作为蝙蝠、蠕虫、青蛙或人的身体的基料，这些身体彼此之间就毫无区别，那么显然这些身体与太阳、月亮、星辰、天空或任何被希腊人称为可感知的神明的事物也毫无区别。因为，作为\Emph{所有}身体基料的相同物质，严格来说，既无性质也无形式，并从某个（其他）来源获得其性质——我不知道来自何处——既然塞尔苏斯执意坚持，没有任何可朽坏的东西能是天主的化工。那么，无论何物的可朽坏部分，既然从相同的基料中产生，按塞尔苏斯自己的论证，必然也是相同的；除非，确实，当他发现自己在此处被逼无奈时，他抛弃柏拉图（柏拉图认为灵魂出自某种器皿），而投靠亚里士多德和逍遥派，他们声称以太是\Emph{非物质的}，由一种与另外四种元素分开的第五种本性构成，柏拉图学派和斯多葛学派对此观点都曾严正抗议。而我们这些被塞尔苏斯鄙视的人也要反驳它，因为我们需要解释并维护先知如下的话：\BibleQuote{诸天将要毁灭，你却伫立不变；它们都要如同衣服渐渐陈旧；你将它们卷起，它们就被更换；你却永不改变。}然而，这些评论足以回应塞尔苏斯，当他断言\Quote{灵魂是天主的化工，但身体的本质是不同的}之时；因为从他的论证可推出，蝙蝠、蠕虫或青蛙的身体与天上的存在的身体之间并无区别。\par

% structure: p0118 source=h2 target=subsection reason=relative-heading-rank; base=h2

\subsection{第五十七章}

那么，看看我们是否应当向一个持有如此观点且诽谤基督徒的人让步，从而放弃一个解释身体之间差异的理论，该理论认为差异是由于内在和外在的不同性质被植入它们之中。因为我们知道，有\BibleQuote{属天的身体，也有属地的身体}；也知道\BibleQuote{属天的光荣是一种，属地的光荣是另一种}；甚至属天身体的光荣也不相同：因为\BibleQuote{太阳的光荣是一种，星辰的光荣是另一种}；而在星辰自身之中，\BibleQuote{这星与那星的光荣不同}。因此，作为期待死者复活的人，我们主张身体中的性质会发生变化：因为有些身体，\BibleQuote{播种的是可朽坏的，复活起来的是不朽坏的}；另一些，\BibleQuote{播种的是可羞辱的，复活起来的是光荣的}；再一些，\BibleQuote{播种的是软弱的，复活起来的是强健的}；而那些\BibleQuote{播种的是属生灵的身体，复活起来的是属神的身体}。那为身体所依存的物质能够接受造物主乐意赋予的性质，这一点我们所有接受天主眷顾道理的人都坚定持守；因此，如果天主愿意，现在一种性质被植入物质的这一部分，而以后另一种不同且更好的性质会被植入。但是，既然从世界开始就有为调节身体变化而确立的法则，并且这些法则在世界存续期间会一直持续，我不知道，在世界毁灭之后，当一种新的、不同的秩序接续而来时——也就是我们的圣经所称的（诸时代的）终结——在现今，一条蛇从死人身上形成，如众人所说，从背髓中长出，一只蜜蜂从牛身上生出，一只黄蜂从马身上生出，一只甲虫从驴身上生出，一般来说，大多数身体都生出蛆虫，难道不令人惊奇吗？事实上，塞尔苏斯认为，这可以表明，这些身体中没有一个是天主的化工，而且这些性质（我不知道为何如此安排，使一种从另一种中生出）也不是神圣理智的化工，神圣理智不会产生物质性质中所发生的变化。\par

% structure: p0120 source=h2 target=subsection reason=relative-heading-rank; base=h2

\subsection{第五十八章}

但我们还有更多话要对塞尔苏斯说，当他宣称\Quote{灵魂是天主的化工，而且身体的本质是不同的}，并且提出这样的意见不仅没有证明，甚至没有清楚定义他的意思；因为他没有清楚说明，他所指的是每一个灵魂都是天主的化工，还是只有理性灵魂。那么，这就是我们要说的：如果每一个灵魂都是天主的化工，那么显然那些最卑微的无理性动物的灵魂也是天主的化工，所以所有身体的本质都与灵魂不同。然而，他在后续的言论中似乎主张，当他声称\Quote{无理性动物比我们更蒙天主喜爱，对神性有更纯粹的知识}时，不仅是人的灵魂，而且在更大程度上，无理性动物的灵魂才是天主的化工；因为这可以从它们被认为比我们更蒙天主喜爱这一说法中推出。现在，如果只有理性灵魂是天主的化工，那么，首先，他没有清楚表明这是他的意见；其次，从他关于灵魂的含糊措辞——即，是否不是每一个灵魂，而只有理性灵魂才是天主的化工——可推出如下结论：并非所有身体的本质都（与灵魂）不同。但如果并非所有身体的本质都不同，尽管每个动物的身体与其灵魂相对应，那么显然，那个灵魂是天主化工的动物的身体，与那个居住着非天主化工的灵魂的动物的身体，将是不同的。因此，断言蝙蝠、蠕虫或青蛙的身体与人的身体没有区别，将是假的。\par

% structure: p0122 source=h2 target=subsection reason=relative-heading-rank; base=h2

\subsection{第五十九章}

因为如果某些石头和建筑物，按照它们是为光荣天主而建造、或是为容纳可耻可诅之人而建造，就应当被视为比其他石头和建筑物更神圣或更不洁，那是荒谬的；而身体却不应按照所居住的是理性或非理性的存有、以及这些理性存有是人类中最有德或最无价值的人而彼此不同。实际上，这种区分原则曾使一些人神化杰出人物的身体，认为它们接受了有德的灵魂，而拒绝并羞辱极恶之人的身体。我并不主张这一原则始终被正确运用，但它起源于一个正确的观念。确实，有智慧的人在\Person{阿尼图斯}和\Person{苏格拉底}死后，会想以同等的荣誉埋葬两人的身体吗？他会为两人堆起同样的坟冢或墓碑吗？我们引用这些例子是因为\Person{凯尔苏斯}的话：\Quote{这些没有一样是天主的作品}（其中\Quote{这些}一词指人的身体，或从身体出来的蛇，以及牛的身体，或从牛的身体出来的蜜蜂；还有马或驴的身体，或从马出来的黄蜂和从驴出来的甲虫）；为此我们不得不回到对他的陈述的考察，即\Quote{灵魂是天主的作品，但身体的本性不同。}\par

% structure: p0124 source=h2 target=subsection reason=relative-heading-rank; base=h2

\subsection{第六十章}

他接着说道：\Quote{一个共同的本性贯穿所有前面提到的身体，并且这个本性在循环变化中来去相同。}对此，从已说的可以清楚看出，不仅前面列举的那些身体具有共同的本性，连天上的身体也是如此。如果是这样，那么显然，按照\Person{凯尔苏斯}（虽然我不知道是否符合真理），有一个本性在所有身体中循环变化中来去相同。这也显然是一些认为世界将要毁灭之人的意见；而那些持相反意见的人，也将试图在不假设第五元素的情况下表明，按照他们的判断，也有一个本性\Quote{在所有身体中循环变化中来去相同。}因此，连那可朽的也存留，以便经历变化；因为作为基础的物质，虽然其属性消亡，但仍持存，按照那些认为物质非受造之人的意见。然而，如果任何论证能表明物质不是非受造的，而是为某些目的受造的，那么显然它将不具有如在非受造假设下所具有的相同恒常本性。但我们现在在回答\Person{凯尔苏斯}的指责时，目的不是讨论这些自然哲学的问题。\par

% structure: p0126 source=h2 target=subsection reason=relative-heading-rank; base=h2

\subsection{第六十一章}

他进一步主张：\Quote{没有物质的产品是不朽的。}对此可以回答说：如果没有物质的产品是不朽的，那么要么整个世界是不朽的，因而它不是物质的产品，要么它并非不朽。因此，如果世界是不朽的（这与那些说独有灵魂是天主的作品、且是从某个调合碗产生出来的观点相符），就请\Person{凯尔苏斯}表明世界不是从无性质之物质产生的，同时记住他自己的断言：\Quote{没有物质的产品是不朽的。}然而，如果世界并非不朽（因为它既然是物质的产品），而是可朽的，那么它是毁灭呢，还是不毁灭？因为如果它毁灭，它作为天主的作品而毁灭；那么，如果\Emph{世界}毁灭，\Emph{灵魂}又会怎样呢？它也是天主的作品。让\Person{凯尔苏斯}回答这个！但如果他歪曲不朽的概念，断言虽然\Emph{可朽}，它却是不朽的，因为它并不\Emph{真正}毁灭；它\Emph{能够}死去，但不\Emph{实际}死去——那么显然，按照他，将存在某种同时是可朽和不朽的东西，因为它能同时处于两种状态；那不死的将是可朽的，而那本性上并非不朽的，将因它不死而被特称为不朽！那么，根据语词意义的何种区分，他将主张没有物质的产品是不朽的呢？由此你看到，他著作中的思想，被仔细审视和检验时，被证明\Emph{不}是健全和无可辩驳的。在做了这些断言之后，他补充说：\Quote{关于这一点，这些评论已足够；如果有人能进一步听取和考察，他将知道（真理）。}那么，让我们这些在他看来是无知的人，看看我们若能稍听他一会儿并继续研究，会得出什么结果。\par

% structure: p0128 source=h2 target=subsection reason=relative-heading-rank; base=h2

\subsection{第六十二章}

接著，他認為能以寥寥數語使我們了解惡的性體問題，這個問題曾在許多重要的論著中受到各種討論，且得到截然相反的解釋。他的話是：\Quote{世上從前沒有、現在沒有、將來也不會再有比從前更多或更少的惡。因為萬有的性體始終如一，惡的產生也始終相同。}他似乎轉述了《\WorkTitle{泰阿泰德篇}》中的討論，其中柏拉圖讓蘇格拉底說：\Quote{惡既不可能從人間消失，也不可能在神明中確立，}等等。但我認為他並沒有正確理解柏拉圖，儘管他自稱在這一部論著中包含了全部真理，並將他反對我們的書命名為《\WorkTitle{真言}》。因為在《\WorkTitle{蒂邁歐篇}》中，當它說：\Quote{當神明用水淨化大地時，}表明大地被水淨化後所含的惡比淨化前更少。而我們主張世上曾有一段時間有較少的惡，這與柏拉圖的意見一致，因為在《\WorkTitle{泰阿泰德篇}》中，他說：\Quote{惡不能從人間消失。}\par

% structure: p0130 source=h2 target=subsection reason=relative-heading-rank; base=h2

\subsection{第六十三章}

我不理解塞爾瑟斯，他既然承認天主眷顧的存在（至少從本書的言辭來看），怎能說從未有過更多或更少的惡，而只有一個固定的數目；從而否定了關於惡之無限本性的美妙教義，並斷言惡在其自身本性上便是無限的。現在，從他的立場似乎可以推出：既然世上從未有過、現在沒有、將來也不會有更多或更少的惡，那麼，按照那些主張世界不可毀滅者的觀點，元素間的平衡是由天主眷顧維持的（不允許任何一方佔優勢，以防止世界的毀滅），因此，一種天主眷顧似乎也掌管著惡（其數目是固定的），以防止它們增加或減少！在其他方面，那些研究善惡問題的哲學家也駁斥了塞爾瑟斯關於惡的論證，他們還從歷史中證明，從前有一些放蕩女子在城外，面戴面具，賣身給需要者；後來，她們變得更加厚顏無恥，摘下了面具，雖然法律不允許她們進城，她們仍留在城外，直到隨著風俗日益敗壞，她們甚至敢於進城。克呂西普在他的著作《\WorkTitle{善與惡}》的序言中提供了這樣的記述。由此也可以看出，惡有增有減，即那些被稱為\Quote{含糊者}的人從前公開出現，遭受和做出一切可恥之事，以滿足那些與他們交往者的情慾；但最近當局已將他們驅逐。至於因邪惡蔓延而在人類生活中出現的無數惡事，我們可以說，從前它們是\Emph{不}存在的。因為最古老的史書雖然對有罪的人提出了無數其他指控，卻對那些可憎罪行的實施者一無所知。\par

% structure: p0132 source=h2 target=subsection reason=relative-heading-rank; base=h2

\subsection{第六十四章}

現在，在這些論證以及其他類似論證之後，塞爾瑟斯如何能不顯得可笑？他竟想像過去從未有過、將來也不會有更多或更少的惡。因為雖然萬有的性體始終如一，但惡的產生卻絕非恆定不變。雖然某一個體的性體始終如一，但他的心智、理性與行為卻不總是相同：有一個時期他尚未獲得理智；另一個時期，他雖擁有理智，卻已染上邪惡，且罪惡或增或減；又有一個時期，他投身於德性，或快或慢地進步，有時經過或長或短的默觀，達至完美的巔峰。同樣，我們可以更大程度地斷言宇宙的性體：雖然它類別相同，但其中發生的事件既不完全相同，也不相似；因為我們並非總有豐收或歉收的季節，也沒有持續的降雨或乾旱。同樣，對於有德性的靈魂，也沒有註定的豐產或不豐產時期；惡的傳播或增或減也是如此。那些盡其所能探究萬物的人，必須留意對惡的這種估量：其數量並非始終相同，這是由於天主眷顧的作用，祂要麼保存地上的事物，要麼藉洪水與烈火淨化它們；而且這或許不僅關乎地上的事物，也關乎整個需要淨化的宇宙，當其中的邪惡變得巨大時。\par

% structure: p0134 source=h2 target=subsection reason=relative-heading-rank; base=h2

\subsection{第六十五章}

此后，\Person{刻尔苏斯}继续说：\Quote{对于非哲学家而言，确实不易查明恶的起源，但对大众来说，只需说恶并非来自天主，而是附着于物质，居住在有死之物中就够了；而有死之物的进程从始至终相同，因此按照所定的循环，过去、现在与未来必定总是重复同样的事物。}\Person{刻尔苏斯}在这里指出，对于非哲学家而言，查明恶的起源并不容易，仿佛哲学家获得这种知识是轻而易举的，而对于非哲学家则是困难的，尽管经过巨大努力仍有可能达到。现在，我们对此说：恶的起源连哲学家也不易掌握，也许即使对于他们，除非藉由神圣启示，否则不可能清楚理解恶是什么、它们如何起源以及如何消失。但是，尽管不认识天主是一种恶，而且其中最大的恶之一是不知道应如何事奉和敬拜天主，然而，正如\Person{刻尔苏斯}也会承认的，无疑有些哲学家对此一无所知，这从不同哲学学派的观点可以明显看出；然而，按照我们的判断，没有人能够查明恶的起源，除非他知道认为虔诚可以在不同国家按照普遍流行的政府观念所制定的法律中毫发无损地保存下来是邪恶的。此外，凡没有听说过被称为\Quote{魔鬼}的那一位及其\Quote{天使}的相关事迹、不知道他在成为魔鬼之前是什么、他\Emph{如何}变成这样的、以及那些被称为他的天使者同时背道的原因的人，都无法查明恶的起源。但愿意获得这种知识的人必须更准确地了解魔鬼的本性，知道就其魔性本质而言，他们并非天主的作品，而只是就其具有理性而言才是天主的作品；也要知道他们的起源，他们是如何成为这样的存在，即当变成魔鬼时，他们心智的能力仍然保留。如果说有任何人类研究的主题难以被我们的本性所掌握，那么恶的起源无疑可以被认为是这样的主题。\par

% structure: p0136 source=h2 target=subsection reason=relative-heading-rank; base=h2

\subsection{第六十六章}

\Person{刻尔苏斯}接着仿佛他能说出关于恶的起源的某些秘密，却宁愿保持沉默，只说适合大众的话，继续说道：\Quote{关于恶的起源，对大众只需说：它们并非来自天主，而是附着于物质，居住在有死之物中。}的确，恶确实并非来自天主；因为按照我们的一位先知\Person{耶肋米亚}所说，确实\BibleQuote{邪恶与福乐并不出自至高者的口。}但是，认为附着于物质、居住在有死之物中的物质是恶的原因，在我们看来并不正确。因为每个人的心灵才是他身上所生之恶的原因，而这就是恶（抽象意义上的）；由此产生的行为是邪恶的，严格来说，我们并不认为有其他东西是恶的。然而，我知道这个话题需要非常详尽的处理，而藉由天主的恩宠光照心灵，那些被天主视为堪当获得关于此主题必要知识的人，可以成功尝试。\par

% structure: p0138 source=h2 target=subsection reason=relative-heading-rank; base=h2

\subsection{第六十七章}

我不明白\Person{刻尔苏斯}在写作反对我们的论文时，怎么会认为采纳一种需要至少大量貌似合理的推理才能使其显得如他所愿的观点是有利的，这种观点就是\Quote{有死之物的进程从始至终相同，而且按照所定的循环，过去、现在与未来必定总是重复同样的事物。}现在，如果这是真的，我们的自由意志就被取消了。因为如果在有死之物的流转中，按照所定的循环，同一事件必须永远在过去、现在和未来发生，那么显然，\Person{苏格拉底}必然永远是哲学家，并因引进异神和败坏青年而被定罪。\Person{阿尼图斯}和\Person{梅利图斯}必然永远是控告他的人，\Place{亚略巴古}议会必定永远判他被毒堇汁处死。同样，按照所定的循环，\Person{法拉里斯}必然永远是暴君，\Person{斐赖的亚历山大}必然犯下同样的残暴行径，那些被判处投入\Person{法拉里斯}的铜牛的人必然不断从其中发出哀号。但如果这些被承认，我看不出我们的自由意志如何能被保存，或如何能恰当地给予称赞或责备。我们还可以对\Person{刻尔苏斯}进一步说，作为对这种观点的回应：\Quote{如果道德事物的进程从始至终总是相同，而且按照所定的循环，同一事件必须在过去、现在和未来总是发生，}那么按照所定的循环，\Person{梅瑟}必再次与犹太人民从埃及出来，\Person{耶稣}必再次降临人间生活，并施行同样的行动——按这种观点，他并非只做一次，而是随着周期的运转做了无数次。不，基督徒也将在所定的循环中相同；\Person{刻尔苏斯}也将再次写他这篇论文，他已经做了无数次了。\par

% structure: p0140 source=h2 target=subsection reason=relative-heading-rank; base=h2

\subsection{第六十八章}

可是，刻尔苏斯却说，只有\Quote{按既定周期，\Emph{可朽}事物的进程在过去、现在和将来必定始终相同}；而大多数斯多葛学派却坚持认为，不仅可朽事物的进程如此，不朽事物以及他们视为神的事物也是如此。因为在过去已经发生过无数次、将来还会发生无数次的世界焚毁之后，从开始到结束，万物的安排都相同且将继续相同。斯多葛学派在试图回避别人对他们观点的反驳时——不知何故——断言：随着一个又一个周期循环，所有人将与以往周期中的人完全相同；因此，苏格拉底不会重新活着，而是一个与\Person{苏格拉底}完全相似的人，他将娶一位与\Person{赞蒂佩}完全相似的妻子，并将被与\Person{安尼图斯}和\Person{梅利图斯}完全相似的人指控。然而，我不理解世界如何总是相同，个体彼此无差别，而其中的事物却不相同（尽管完全相似）。但是，针对刻尔苏斯和斯多葛派言论的主要论证更适宜在别处探讨，因为本次场合与我们的目的不符，不能详述这些要点。\par

% structure: p0142 source=h2 target=subsection reason=relative-heading-rank; base=h2

\subsection{第六十九章}

他接着说：\Quote{有形之物也并非（由天主）赐予人，而是每一具体事物出于整体安全的缘故，按照我先前提到的从一物到另一物的变化而产生和消灭。}然而，我们无需在此停留反驳这些论述，因为我已经尽力反驳过了。以下论点也已被回答，即\Quote{在可朽者当中，善与恶将不多不少。}这一论点也被提及，即\Quote{天主不需要重新修正祂的工程。}但是，天主并非像一个人设计作品不完善且执行不熟练那样，借着洪水或大火来洁净世界以施行纠正，而是为了防止邪恶浪潮涨得更高；而且我认为，祂是在事先准确决定的时期扫除邪恶，以利于整个世界的善。然而，如果他（刻尔苏斯）断言邪恶消失后又重新出现，这类问题将需要在一篇专门的论文中探讨。那么，天主希望重新修正祂的工程，总是为了修复已经变得有缺陷的事物。因为虽然在创造世界时，万物已被祂以最美丽和最稳定的方式安排，祂仍然需要对那些受邪恶疾病折磨的人们以及因而被污染的世界行使某种治愈的力量。但天主没有忽视任何事物，也不会忽视任何事物；因为祂在每一个特定时刻，在败坏而改变的世界中做祂适宜做的事。正如农夫根据一年中不同的季节对土壤和农作物进行不同的耕作，天主管理整个时间世代，就好像它们是许多个别的年份，在每一个世代中做合宜之事，合理关照世界；这一工作，正如只有天主真正理解，同样也由天主完成。\par

% structure: p0144 source=h2 target=subsection reason=relative-heading-rank; base=h2

\subsection{第七十章}

刻尔苏斯关于邪恶发表了如下性质的陈述，即\Quote{虽然一件事在你看来是恶的，但绝不肯定它就是如此；因为你不知道什么对自己、对他人或对整个世界有益。}这个断言带有某种谨慎；它暗示恶的本质并非全然邪恶，因为在个别事例中可能被视为恶的东西，可能包含对整体社群有益之处。然而，为避免任何人误解我的话，并以作恶为借口，仿佛他的邪恶对世界有益，或至少\Emph{或许}有益，我们必须说：尽管天主保全每个人的自由意志，可以利用恶人的邪恶来治理世界，如此安排它们以有利于整体；然而，这样的人仍应受谴责，并且因此被指派用于一种令每个个体憎恶但有益于整体社群的用途。这如同有人说着：在一个城市里，一个人犯了某些罪行，并因此被判处在对社群有益的公共工程中服役，他所做的事对整个城市有利，而他自己却从事一项可憎的任务，任何有理智的人都不愿从事。耶稣的宗徒保禄也教导我们，甚至最恶的人也会有助于整体的善，尽管他们本身是卑贱的，但最有德行的人也将对世界有最大的益处，并因此占据最高贵的地位。他说：\BibleQuote{在大户人家，不但有金器银器，也有木器瓦器；有作贵重的，有作卑贱的。所以人若自洁，脱离卑贱的事，就必作贵重的器皿，成为圣洁，合乎主用，预备行各样的善事。}我认为有必要说这些话以回应那个断言，即\Quote{虽然一件事在你看来是恶的，但绝不肯定它就是如此，因为你不知道什么对自己或他人有益。}以免有人从这个话题所述之言找到犯罪的借口，以他如此行将有益于世界为托辞。\par

% structure: p0146 source=h2 target=subsection reason=relative-heading-rank; base=h2

\subsection{第七十一章}

然而，接续下文，\Person{塞耳苏斯}不明白圣经中论及天主的语言是依照拟人情感的观点而说的，他嘲笑那些讲述对不敬虔者说话时带有忿怒以及对罪人发出威胁的段落。我们必须说明：正如我们自己与幼童交谈时，并非以施展自己的口才为目的，而是顺应我们受托之幼童的软弱，说那些在我们看来有益于纠正和改善孩童（作为孩童）的话，做那样的事；同样，天主的圣言在处理历史时也是如此，以听众的能力和他们将获得的益处作为其宣告（关于祂）是否合宜的标准。总括来说，关于论及天主的这类说话方式，我们在\BibleBook{申命}{纪}中读到以下的话：\Quote{上主你们的天主容忍你们的习气，如同人容忍自己的儿子一样。}仿佛是取了人的习气，为获致人的益处，圣经使用了这样的表达；因为倘若天主以更适合其自身尊严的方式向群众说话，那是不适合群众状况的。然而，凡切望真正理解圣经的人，必会发现圣经对那些被称为\Quote{属神的人}所说的属灵真理，通过比较对较软弱者所说的话与对较敏锐之人所宣告的话的意义——对于能领会同一段落的人来说，两者常常存在于同一段经文之中。\par

% structure: p0148 source=h2 target=subsection reason=relative-heading-rank; base=h2

\subsection{第七十二章}

我们确实说天主的\Quote{义怒}。但我们并不主张这表示祂有任何\Quote{激情}，而是说这是为以严厉手段管教那些犯了许多重大罪恶的罪人而采取的方式。因为所谓天主的\Quote{义怒}和\Quote{义怒}是一种管教方式；这种看法符合圣经，由第六篇 \BibleBook{圣咏}{}的话可以证明：\Quote{上主，求祢不要在祢的震怒中责备我，不要在祢的烈怒中惩罚我}；\BibleBook{耶肋米亚}{}中亦然：\Quote{上主，求祢惩戒我，但要以公正，不要在祢的震怒中惩治我，免得祢使我归于无有。}此外，谁若读到\BibleBook{列王}{下}中天主的\Quote{义怒}促使\Person{达味}统计人民，并从\BibleBook{编年纪}{上}发现是魔鬼提议此事，将两者比较，便容易看出\Quote{义怒}为何被提及——正如宗徒\Person{保禄}所说，众人都是这\Quote{义怒}之子：\Quote{我们生来就是义怒之子，和别人一样}。再者，\Quote{义怒}并非天主方面的激情，而是各人因自己的罪招致己身，这由\Person{保禄}进一步的话显明：\Quote{难道你轻视祂丰厚的慈爱、宽容与忍耐，不晓得天主的慈爱是引领你悔改？但你硬着心，不肯悔改，为自己积蓄忿怒，直到忿怒的日子，即正义审判的天主启示的日子。}那么，人怎能为自己积蓄\Quote{义怒}以迎接\Quote{忿怒的日子}，如果\Quote{义怒}理解为\Quote{激情}？或者\Quote{义怒的激情}怎能有助于管教？此外，圣经教导我们完全不要发怒，在第三十七篇 \BibleBook{圣咏}{}中说：\Quote{停止忿怒，放下怒火}，并且藉\Person{保禄}的口命令我们\Quote{戒绝这一切：忿怒、暴躁、恶毒、毁谤、污秽的言语}，它不会使天主陷入它要我们彻底摆脱的同一激情。此外，论及天主愤怒的话语应按\Emph{寓意}理解，这由有关祂\Quote{睡觉}的记载显明：先知如同唤醒祂说：\Quote{主，祢为什么睡觉？求祢醒来！}又说：\Quote{那时上主像从睡梦中醒起，像酒醉的勇士呼喊。}那么，如果\Quote{睡觉}必另有所指，而非字面第一义，为何\Quote{义怒}不可同理理解？至于\Quote{威胁}，乃是对恶人将要遭受（惩罚）的预告；正如一人称医生的话为\Quote{威胁}，当他告诉病人：\Quote{你若不听我的处方，不按我指示的饮食和生活方式节制，我就必须用刀和火。}因此，我们并未将人的激情归于天主，也未对祂抱有亵渎的看法；当我们根据圣经本身、经过认真比较后呈现关于祂的各种叙述时，我们并未犯错。因为那\Quote{圣言}的智慧使者们，其唯一目标就是尽己所能使听众脱离软弱的见解，并以理智装备他们。\par

% structure: p0150 source=h2 target=subsection reason=relative-heading-rank; base=h2

\subsection{第七十三章}

作为他不理解关于天主\Quote{义怒}之言论的后续，他接着说：\Quote{难道不是荒谬的吗？一个\Emph{人}向犹太人发怒，便从少年以上全部杀尽，并烧毁他们的城（他们无力抵抗他）；而他们所谓全能的\Emph{天主}，发怒、愤慨并发出威胁，却不惩罚他们，反倒派遣自己的圣子忍受他所受的苦难？}如果犹太人，在他们胆敢加于\Person{耶稣}的待遇之后，连同所有少年都灭亡了，他们的城市也被火烧毁，那么他们所遭受的惩罚，并非来自别的义怒，而是他们为自己积蓄的；因为天主对他们的审判，由天主旨意所决定，按照希伯来人的传统用语，被称为\Quote{义怒}。而全能天主的圣子所受的苦难，是祂自愿为人的救恩而受的，正如我在前文尽力阐述的。他接着说：\Quote{但我不愿只谈犹太人（因为那不是我的目的），而是按照我的承诺，谈谈整个自然界，我将更清楚地说明已经说过的话。}有哪个谦逊的人，读到这些话，并知道人性的软弱，不会对那承诺解释\Quote{整个自然界}的冒犯性质感到愤慨，并对那种狂妄感到愤慨——那种狂妄促使他在自己的书上题写他胆敢赋予的书名（\WorkTitle{真言}）？但让我们看看他关于\Quote{整个自然界}要说些什么，以及他要\Quote{更清楚地}说明什么。\par

% structure: p0152 source=h2 target=subsection reason=relative-heading-rank; base=h2

\subsection{第74章}

他接着用很多话责备我们，因为我们断言天主创造万物是为了人。因为从动物史和它们所表现出的智慧，他要证明万物的存在并不比为了人更多是为了无理性动物。在我看来，他说话的方式与那些人相似，那些人由于对敌人的厌恶，用他们自己朋友受称赞的同一件事来控告敌人。因为正如在上述例子中，仇恨使这些人看不清，他们正通过自以为诽谤敌人的手段，来控告自己最亲密的朋友；同样，瑟尔苏也在论证中陷入混乱，没有看出自己是在控告斯多葛派的哲学家，他们并不错误地将人置于首位，将理性本性一般地置于无理性动物之前，他们主张天主眷顾主要是为了理性本性而创造万物。因此，理性受造物作为主要的受造物，占据相当于子宫中婴儿的位置，而无理性、无灵魂的受造物则占据与婴儿同时创造的包膜的位置。我也认为，正如城中物资和市场的监督者履行职责，只是为了人类，而狗和其他无理性动物则从剩余中获益；同样，天主眷顾以\Emph{特别方式}提供理性受造物所需；由此也得出，无理性受造物同样享受为人类所做之事的益处。正如那个错误地声称市场监督者为人类和狗提供一样多的人，因为狗也得到货物的一份；同样，瑟尔苏和与他想法相同的人，在断言天主的安排为人类提供生计，并不比为植物、树木、蔬菜和荆棘提供得更多时，就在更大程度上对那为理性受造物提供一切的天主犯有不敬之罪。\par

% structure: p0154 source=h2 target=subsection reason=relative-heading-rank; base=h2

\subsection{第75章}

首先，他认为\Quote{雷电、闪电和雨水不是天主的作为}——这最终更清楚地显露出他的伊壁鸠鲁派倾向；其次，他断言\Quote{即使承认这些是天主的作为，它们被造出来不仅是为了供养我们人类，也是为了供养植物、树木、草本和荆棘}——他像一个真正的伊壁鸠鲁派那样，坚持认为这些是偶然的产物，而非天主眷顾的工作。如果这些东西对我们来说并不比植物、树木、草本和荆棘更有用，那么显而易见，它们要么根本不是出自天主眷顾，要么出自一种并不比眷顾树木、草本和荆棘更眷顾我们的天主眷顾。这两个假设本身都是不敬的，并且反驳任何以此指控我们不敬的人都是愚蠢的，因为从所说的话中，每个人都能清楚地看出谁才是不敬之徒。接着，他又补充道：\Quote{虽然你可能会说这些东西，即植物、树木、草本和荆棘，是为人的使用而生长，但为何你坚持认为它们是为人的使用而非为最野蛮的无理性动物而生长？}那么，让刻尔苏斯明确地说，大地上如此多样的产物不是天主眷顾的工作，而是某种原子的偶然结合产生了如此多样的性质，并且如此多种类的植物、树木和草本彼此相似纯属偶然，没有任何安排的理性赋予它们存在，它们\Emph{并非}来源于某种令人惊叹的理智。然而，我们基督徒，专一崇拜创造这一切的唯一天主，为这些事物而感谢祂，因为祂不仅为我们，也（因我们的缘故）为我们所管辖的动物预备了这样的家园，正如经上所载：\BibleQuote{祂使青草为牲畜生长，使蔬菜为人服务，为使大地出产食物，酒能悦人心，油能润人面，粮能壮人心。}至于祂甚至为最野蛮的动物提供了食物，这并不令人惊讶，因为这些动物被一些哲学（论及此事）的人认为是为理性受造物提供锻炼而造的。我们中的一位智者曾在某处说：\BibleQuote{不要说：这是什么？或那是为什么？因为万物都按它们的用途被造。也不要说：这是什么？或那是为什么？因为一切事物都将在其季节中被寻出。}\par

% structure: p0156 source=h2 target=subsection reason=relative-heading-rank; base=h2

\subsection{第七十六章}

此后，刻尔苏斯希望坚持天主眷顾创造了大地产物，并非更多地为了我们，而是为了最野蛮的动物，他继续说道：\Quote{我们确实通过劳动和受苦勉强获得微薄而艰辛的生存，而它们无需播种和耕作，一切为他们自然产生。}他没有注意到，天主为了在所有地方锻炼人的理智（使它不致闲散和不熟悉技艺），将人造为充满需要的存在，以便藉着其非常匮乏的状态，迫使他成为技艺的发明者，其中一些技艺服务于他的生存，另一些服务于他的保护。因为那些本不会追求神圣事物、也不会从事哲学研究的人，若因富足而全然忽略心智的培养，倒不如被置于匮乏之中，以便运用理智去发明技艺，这样更好。人类生活必需品的匮乏导致了农艺的发明，以及葡萄栽培的发明；又导致了园艺、木工和铁工技艺，通过这些技艺制造出服务于维持生活的工具。覆盖物的匮乏又引入了纺织技艺，紧接着是梳毛和纺纱；然后是房屋建造；这样，人类的理智甚至上升到建筑技艺。必需品的匮乏也使得其他地方的产物通过航海和领航技艺被运送到缺乏它们的人手中；因此，即使在这方面，人们也可以赞叹天主眷顾，它使理性受造物比无理性的动物更易于匮乏，然而一切都是为了其利益。因为无理性动物有食物供应，因为它们甚至没有发明技艺的动力。此外，它们还有天然的覆盖物，因为它们有毛发、翅膀、鳞片或外壳。那么，让以上是我们对刻尔苏斯论断的答复，他说：\Quote{我们确实通过劳动和受苦勉强获得微薄而艰辛的生存，而它们无需播种和耕作，一切为他们自然产生。}\par

% structure: p0158 source=h2 target=subsection reason=relative-heading-rank; base=h2

\subsection{第七十七章}

接着，他忘记了自己的目标是控告犹太人和基督徒，他引用了一句欧里庇得斯的抑扬格诗句，这句诗与他自己的观点相反，并与这些文字争辩，指责它们是一个错误的陈述。他的原话如下：\Quote{但如果你要引述欧里庇得斯的话，即}\par

\Quote{太阳与黑夜是凡人的奴仆}\par

\Quote{为什么它们为我们的服务就比给蚂蚁和苍蝇的更多？因为黑夜为他们造出来以便他们休息，白昼为了他们能看见并恢复工作。}现在无疑的是，不仅某些犹太人和基督徒宣称太阳和天体是我们的仆人；而且，根据一些人的说法，那个被称为舞台哲学家的人也说了这话，他曾听过阿那克萨哥拉的自然哲学讲座。但此人断言世上万物都服从于所有理性受造物——一个理性本性被取来代表全体，按照部分代表整体的原则；这再次清楚地从以下诗句显示出来：——\par

\Quote{日与夜为凡人之奴隶。}\par

也许悲剧诗人以太阳指代白昼，因太阳是白昼的原因——他教导说，那些最需要白昼和黑夜的事物是月下的事物，其他事物则不如地上之物那般需要。于是白昼与黑夜隶属凡人，是为理性受造物而受造的。若蚂蚁与苍蝇，它们白昼劳作、黑夜安息，此外还受益于那些为人类受造的事物，我们便不能说白昼与黑夜是为蚂蚁与苍蝇而存在，也不可认为它们毫无目的受造，而是按照天主眷顾的计划，为人类而形成的。\par

% structure: p0164 source=h2 target=subsection reason=relative-heading-rank; base=h2

\subsection{第七十八章}

他随后进一步反对自己为人类辩护之辞，即无理性动物是为人类而受造，说道：\Quote{若有人因我们捕猎其他动物并以它们的肉为食而称我们为动物界的主宰，那么我们要说，为何我们不是为它们而受造，既然它们也捕猎并吞噬我们？不仅如此，我们追逐猎物时需要网罗、武器以及多人协助，还有猎犬；而它们则天生立即拥有武器，轻易将我们置于其下。}在此我们可观察到，理解的恩赐已被赋予我们作为强大的助力，远胜野兽可能拥有的任何武器。我们这些在体力上远弱于野兽、比某些动物矮小的人，竟凭理解获得主宰，捕获巨大的大象。我们以温和的方式驯服那些天性可被驯化的动物，而对于那些天性不同或驯化后似乎无用的动物，我们采取防范措施，使我们在需要时将野兽关起来，且当我们需取用它们身上的肉为食时，便宰杀它们，如同对待那些非凶猛的野兽。造物主遂使万物成为理性受造物及其自然理解的仆役。为某些目的，我们需用狗，譬如守护羊圈、牛栏、羊群或住所；为其他目的，需用牛，如农耕；又需用负轭的或驮重的牲畜。因此，可以说狮子、熊、豹、野猪等兽类被赐予我们，为激发我们内在的勇毅品德。\par

% structure: p0166 source=h2 target=subsection reason=relative-heading-rank; base=h2

\subsection{第七十九章}

其次，为回应人类——他们意识到自身远超无理性动物的优越性——他说：\Quote{关于你们断言天主赐予你们捕捉野兽并随意使用的能力，我们要说，很可能在城邦建立之前、技艺发明之前、如今这类社会形成之前、以及使用武器和网罗之前，人通常被野兽捕获并吞噬，而野兽极少被人捕获。}对此，需注意：虽然人捕野兽，野兽也捕食人，但那些凭理解主宰野蛮凶残者的人，与那些不用理解保护自己免受野兽伤害的人，二者之间有天壤之别。而当塞耳苏斯说：\Quote{在城邦建立之前、技艺发明之前、如今这类社会形成之前，}他似乎忘记了他先前所说的：\Quote{世界是非受造且不朽坏的，只有地上之物经历洪水与火灾，且并非同时发生。}姑且承认他这些让步与我们的观点完全一致，与他本人及其上述言论却全然不合；然而这有什么用呢，证明最初人大多被野兽捕获吞噬，而野兽从未被人捕获？因为，世界是按天主眷顾的意愿受造的，天主统管万有，人类原祖在其生存之初必须得到更高权能的保护，以便从一开始神圣本性便与人类结合。阿斯克拉的诗人意识到此事，咏道：\newline{}\par

那时宴席共享，坐席也共享，\newline{}如同不死的神明与有死的凡人。\par

% structure: p0169 source=h2 target=subsection reason=relative-heading-rank; base=h2

\subsection{第八十章}

再者，那些以梅瑟命名的圣经记载最初的人听到神圣的声音与神谕，并有时看到天主的使者前来访问他们。因为在世界存在之初，人性很可能得到更大程度的帮助（较以后），直到他们在理解及其他德性上有所进步，且技艺被发明，从而能自行维持生计，不再需要监管者及那些以奇迹方式展示天主旨意工具来引导他们的人。由此可知，\Quote{最初人被野兽捕获吞噬，而野兽极少被人捕获}是虚假的。塞耳苏斯的以下说法也是不实的：\Quote{天主就这样使人受制于野兽。}因为天主并未使人受制于野兽，而是将野兽作为理解的猎物赐予人，并赐予针对野兽的技艺——这些技艺是理解的产物。因为人并未脱离天主的帮助，便为自己渴求保护自己免受野兽侵害、并主宰野兽的方法。\par

% structure: p0171 source=h2 target=subsection reason=relative-heading-rank; base=h2

\subsection{第八十一章}

我们高尚的对手却没有注意到，有多少位哲学家承认天主眷顾的存在，并认为天主眷顾为了理性受造物的缘故创造了万有；他竟尽其所能地推翻那些表明基督教与哲学在这事上和谐一致的教义。他也没有看出，如果接受在天主面前人与蚂蚁或蜜蜂没有区别这一说法，会对宗教造成多大损害。他进而补充说：\Quote{如果人因建立城邦、使用政体、政府形式和主权而似乎优越于无理性动物，这不足为凭，因为蚂蚁和蜜蜂也做同样的事。蜜蜂确实有一位君主，他有随从和侍者；它们中间发生战争、胜利、对失败者的屠杀、城邦和郊区、连续的劳作，以及对懒惰者和邪恶者的审判；因为雄蜂被赶走和惩罚。}现在，他没有注意到经由理性和思虑所成就的事，与出于无理性本性且纯粹机械性的事之间的区别。因为这些事的起源无法用制造者身上的任何理性原则来解释，因为他们并不具有任何理性原则；而是最古老的存有，祂也是天主子，万有之王，创造了一种无理性本性，这种本性作为无理性的，对被认为配得理性者起到帮助作用。因此，在人类中，凭借许多技艺和井然有序的法律建立起城邦；而人类中的政体、政府、主权，或是名副其实的，体现了某些德性的趋势和运作，或是不当得名的，是模仿前者尽可能设计出来的：因为通过观照这些，最成功的立法者确立了最好的政体、政府和主权。然而，在无理性动物中找不到任何这些东西，尽管克耳苏斯可能将理性的名称和属于理性受造物的安排——如城邦、政体、统治者、主权——甚至加于蚂蚁和蜜蜂；然而，就这些方面而言，蚂蚁和蜜蜂不值得称赞，因为它们的行动并非出于反省。我们应当钦佩天主性，它甚至赋予了无理性动物一种仿佛模仿理性受造物的能力，也许是为了使理性受造物羞愧；使得例如观察蚂蚁的人变得更加勤勉、更节俭地管理自己的财物；而通过观察蜜蜂，他们可以顺服于自己的统治者，并在那些有助于城市安全的合乎宪法的职责中各尽其分。\par

% structure: p0173 source=h2 target=subsection reason=relative-heading-rank; base=h2

\subsection{第八十二章}

也许蜜蜂中所谓的战争也教导人们，在人类中若有战争的必要，应当以公义有序的方式作战。但蜜蜂没有城邦或郊区；它们的蜂巢、六边形巢室和连续的劳作，都是为了人类，人类需要蜂蜜用于多种用途，既可治疗失调的身体，又可作为纯净的食物。我们也不应当将蜜蜂对付雄蜂的措施与城中惩罚懒惰者和邪恶者的审判和刑罚相提并论。但如前所述，我们应当一方面在这些事上钦佩天主性，另一方面表达对人的钦佩，人能够思考并钦佩一切事物（作为与天主眷顾合作），并且不仅执行天主眷顾所决定的工作，也执行由他自己远见所产生的工作。\par

% structure: p0175 source=h2 target=subsection reason=relative-heading-rank; base=h2

\subsection{第八十三章}

在\Person{克耳苏斯}论完蜜蜂之后，为了（尽其所能）贬抑不仅我们基督徒，而且全人类的城市、体制、政府、主权，以及人们为祖国而进行的战争，他转而离题去称赞蚂蚁，以便在赞美它们时，将人类为维持生存所采取的措施与这些昆虫所采取的措施相比较，从而表明他对为冬天储备的远见之蔑视，认为这不过是蚂蚁的非理性眷顾，正如他所认为的那样。那么，一些较为单纯、不知道如何探究万物本性的人，岂不会（至少就克耳苏斯所能达成的程度而言）被转而远离帮助那些被生活重担所压迫的人，并分享他们的劳苦，当他说到蚂蚁时，说\Quote{当它们看到同伴在重担下辛劳时，它们互相帮助搬运负载}？因为那需要被圣言管教、却完全不懂圣言之音的人会说：\Quote{既然如此，在我们和蚂蚁之间没有区别，即使我们帮助那些因背负沉重负担而疲惫的人，我们为何还要继续白白这样做呢？}而蚂蚁作为非理性的受造物，岂不会因为它们的工作被与人类的工作相比较而大大自傲、自视甚高吗？另一方面，人类凭借理性能够听到他们对他人的仁爱如何被藐视，就会被伤害，至少就克耳苏斯及其论点所能达到的程度而言：因为他没有察觉到，当他希望使他的论著的读者远离基督信仰时，他也使非基督徒对承受生活最重负担之人的同情远离了他们。然而，如果他是位哲学家，能够察觉人类之间可以互相施予的善，那么他除了不将人类中间的祝福与基督信仰一同移除之外，还应当（如果他有能力）协助配合那些基督信仰与其余人类共有的卓越原则。再者，即使蚂蚁将那些发芽的谷粒单独存放在一个地方，使它们不会膨胀成芽，而是全年作为食物，这也不应被视为蚂蚁拥有\Emph{理性}的证据，而是万有之母大自然的工作，她甚至装饰了非理性的动物，以致连最微不足道的也不被忽略，却带有自然界植入其中的理性的痕迹。除非，克耳苏斯确实想通过这些论述隐约暗示（因为在许多实例中他喜欢采纳\Person{柏拉图}的思想）所有灵魂都属于同一类属，人的灵魂与蚂蚁和蜜蜂的灵魂之间没有区别，这是那种将灵魂从苍穹降下、使其不仅进入人身而且进入畜体之人的行为。然而，基督徒不会赞同这些意见：因为他们之前已被教导，人的灵魂是按天主的肖像创造的；他们看到，由神圣肖像塑造的本性不可能使其（原始）特征完全泯灭，并接受其他那些我不知道是按非理性动物何等模样形成的特征。\par

% structure: p0177 source=h2 target=subsection reason=relative-heading-rank; base=h2

\subsection{第八十四章}

既然他断言\Quote{当蚂蚁死亡时，幸存者会预留一个特殊的地方（作它们的埋葬之所），而它们祖先的坟墓就是这样的地方}，我们须回答：他对非理性动物的赞颂越多，就越（尽管违反其本意）彰显了那安排万有的理性的工作，并指出了人类当中存在的技艺，这技艺能以它的理性修饰甚至大自然赐予非理性受造物的恩赐。但我为何说\Quote{非理性的}，既然克耳苏斯认为这些按所有人共同观念被称为非理性的动物，其实并非如此？那个自称谈论\Quote{宇宙本性}，并在其著作的题词中夸耀其真实性的\Emph{他}，并不认为蚂蚁缺乏理性。因为，谈到蚂蚁彼此交谈时，他使用了以下语言：\Quote{当它们相遇时，它们进入交谈，因此它们从不迷路；所以它们拥有充分的理性、关于某些一般主题的一些共同概念，以及一种表达偶性事物的声音。}如今，人与人之间的交谈是通过声音进行的，声音表达意图的意义，也对所谓的\Quote{偶性事物}进行表达；但说蚂蚁也是这种情况，将是一个最荒谬的断言。\par

% structure: p0179 source=h2 target=subsection reason=relative-heading-rank; base=h2

\subsection{第八十五章}

此外，他还不以为耻地说（以便他意见的粗鄙本质对于后来之人显明）：\Quote{来吧，若有人从天俯视大地，那么 \Emph{我们} 的行为与蚂蚁和蜜蜂的相比，究竟有什么不同呢？} 那么，按照他自己的设想，从天俯视人类和蚂蚁行为的那一位，难道只注视他们的身体，而不是更关心那通过反思而活动的控制性理性吗？另一方面，后者的主导原则是无理性的，且由冲动和幻想与某种自然机制一起无理性地驱动。但设想那位从天俯视尘世事物者，会愿意从那样的距离注视人类和蚂蚁的 \Emph{身体}，而不更考虑主导原则的本性和冲动的来源——无论那是理性的还是无理性的——这是荒谬的。如果他一旦注视一切冲动的来源，那么显然他也会看到存在的差异，以及人的优越性——不仅超过蚂蚁，甚至超过大象。因为那位从天俯视者，在无理性的受造物中，无论它们的身体多大，将看到除了所谓无理性之外别无其他原则；而在理性的存有中，他将发现理性——人类以及神圣和天上的存有、甚至或许至高天主本身的共同所有物——正因如此，人才被称为按天主的肖像受造，因为至高天主的肖像就是他的理性。\par

% structure: p0181 source=h2 target=subsection reason=relative-heading-rank; base=h2

\subsection{第八十六章}

紧接此后，仿佛他竭尽全力将人类降低到更低的地位，并使他们与无理性的动物平等，且不愿遗漏关于后者任何显示其伟大的情况，他宣称：\Quote{在无理性的受造物中的某些个体，存在着巫术的能力}；因此即使在这方面，人也不能特别自豪，也不愿自夸比无理性的受造物优越。以下是他所说的话：\Quote{然而，若人因拥有巫术的能力而怀有高傲的想法，那么即使在这方面，蛇和鹰在智慧上比他们更优越；因为它们知道许多对付人和疾病的预防药，也知道某些有助于保护幼崽的石头的功效。不过，若人偶然碰到这些，他们便以为获得了一件奇妙的宝藏。} 首先，我不知道他为何将动物所表现出的对天然预防药的知识称为巫术——无论那知识是来自经验，还是某种天生的理解能力；因为 \Quote{巫术} 一词通常被用于别的事物。也许，他确实想作为一个伊壁鸠鲁派，暗中批评这类技艺的全部使用，认为它们仅基于巫师的宣称。然而，就算他承认人们 \Emph{确实} 因这类技艺的知识而非常自豪，无论他们是不是巫师：蛇在这一点上怎能比人更有智慧呢？当它们使用众所周知的茴香来增强视力并产生快速运动时，它们获得这种天然能力并非来自反思的运用，而是来自身体的构造；而人不像蛇那样仅凭本性获得这类知识，而是部分通过经验，部分通过理性，有时通过反思和知识。同样，如果鹰为了在巢中保护幼雏而在找到鹰石后将其带入巢中，那么它们如何显得有智慧且比人更聪明，而人却通过运用其反思能力和理解力，发现了本性作为礼物赐予鹰的东西？\par

% structure: p0183 source=h2 target=subsection reason=relative-heading-rank; base=h2

\subsection{第八十七章}

姑且承认动物也知道其他解毒剂：但这有何助于证明发现这些事物并非出于本性，而是出于理智？因为如果理智是发现者，那么蛇类就只会发现这一种（或至多一两种）事物，而不会发现其他；鹰类则只会发现另一种，其余动物亦复如是；这样它们各自只能有一种发现，远不及人类之多。然而，从每种动物对某类助益的本性固定倾向可以明显看出，它们并不拥有智慧或理智，而是具有一种由\OriginalTerm{圣言}{Logos}所赋予的自然结构倾向，旨在确保动物的生存。事实上，如果我想在这些问题上与凯尔苏斯争论，我可以引用撒罗满在《箴言》中的话，经文如下：\BibleQuote{地上有四样东西虽小，却比智者更聪明：蚂蚁不是一个强大的民族，却在夏天预备粮食；石獾是软弱的一群，却在岩石中建造房屋；蝗虫没有君王，却整齐列队而出；壁虎虽用手支撑，且容易捕捉，却住在君王的宫殿里。}然而我引用这些话，并非按其字面意义，而是按这书的标题（因为书名题为\WorkTitle{箴言}）来探究它们所包含的奥秘含义。因为圣经作者的习惯是将那些字面表达一种意思、隐藏意义却不同的著作分为许多类别，而其中一类就是\WorkTitle{箴言}。为此，在我们的福音书中，救主也被记载说：\BibleQuote{这些事我用比喻对你们说了，但时候将到，我不再用比喻对你们说。}因此，并非\Emph{可见的}蚂蚁是\Quote{比智者更聪明}，而是那些以\Quote{箴言}文体所指示的。我们必须对其他动物做出这样的结论，尽管凯尔苏斯认为犹太人和基督徒的书籍极其简单平凡，并认为那些给予它们寓意解释的人是对作者原意的曲解。我们所说的话足以表明凯尔苏斯枉然诽谤我们，并反驳他关于蛇和鹰比人类更聪明的断言。\par

% structure: p0185 source=h2 target=subsection reason=relative-heading-rank; base=h2

\subsection{第88章}

他为了更详细地表明人类对天主的想法并不优于其他凡有死之物的想法，反而某些无理性的动物能够思考关于那一位——就是在最敏锐的人类（希腊人和外邦人）中对其有如此分歧意见的那一位——于是他接着说：\Quote{如果因为人类能够领会天主的观念就被认为优于其他动物，那么持此见解的人应当知道，这种能力也会为许多其他动物所主张，而且理由充分：因为有什么比预知和预言未来的能力更为神圣呢？人类正是从其他动物，尤其是从鸟类那里获得了这门技艺。听从它们所提供征兆的人便获得了预言的恩赐。因此，如果鸟类和其他有预言能力的动物，由于天主的恩赐能够预知未来，并通过征兆教导我们，那么它们似乎更接近天主的交往，拥有更大的智慧，并且更蒙祂的宠爱。此外，较聪明的人说，动物举行比我们集会更神圣的聚会，它们知道在这些聚会上所说的话，并显示它们确实拥有这种知识：当它们事先宣称鸟类已经决定前往某个特定地方，并要做这件事或那件事，而后鸟类的离去和所做的事证实了它们的断言，就表明了这一点。没有一种动物比大象更遵守誓言，或对神圣事物表现出更大的虔敬；我认为这完全是因为它们对天主有些认识。}请看，他如何立即抓住并视作公认的事实，那些在哲学家（不仅希腊人，也包括那些从某些邪魔那里发现或学到了关于占卜鸟和其他动物之事的外邦人）之间存在争议的问题，据称通过这些动物向人类发出某些预言性的暗示。因为首先，占卜术以及一般的动物预言之术是否存在，本就有争议。其次，那些承认鸟占预言之术的人，对预言的方式也不一致；因为有些人认为，动物——鸟儿进行各种飞翔和鸣叫，其他动物进行各种运动——是受到某些预言的邪魔或神祇的推动。另一些人则相信，动物的灵魂在本性上更为神圣，并且适合此类运作，这是一种最难以置信的假设。\par

% structure: p0187 source=h2 target=subsection reason=relative-heading-rank; base=h2

\subsection{第89章}

然而，\Person{策尔苏斯}既想借上述论断证明无理性动物比人类更神圣、更聪明，就本应更详细地论证这种占卜术的实际存在，然后积极为之辩护，有效地反驳那些想要消灭此类占卜术之人的论点，并以令人信服的方式推翻那些声称动物从魔鬼或神明那里获得引导它们占卜之动力之人的论点，接着证明无理性动物的灵魂比人的灵魂更神圣。因为，若他这样做了，并以哲学精神处理此类事情，我们就当尽力回应他的自信论断，首先驳斥无理性动物比人更聪明的说法，并指出他们声称动物对天主的概念比我们更神圣、且动物之间举行某种神圣集会的说法是虚假的。但现在，相反，\Emph{他}——因我们相信至高天主而控告我们的人——要求我们相信鸟类的灵魂对天主拥有比人类更神圣、更清晰的概念。然而，若这是真的，则鸟对天主的概念比\Person{策尔苏斯}本人更清晰；他如此贬低人类，以致如此也就不足为奇了。不，就\Person{策尔苏斯}所能证明的而言，鸟类拥有更伟大、更神圣的概念——我不说比我们基督徒，或比与我们使用同一\WorkTitle{圣经}的犹太人，甚至比希腊的神学家所拥有的也更多，因为他们仅仅是人类。按照\Person{策尔苏斯}的看法，从事占卜的鸟类，诚然，比\Person{斐瑞居德}、\Person{毕达哥拉斯}、\Person{苏格拉底}和\Person{柏拉图}更了解神性！那么，我们应当到鸟类那里作老师，以便按照\Person{策尔苏斯}的观点，它们借着占卜能力教导我们关于未来事件的知识，同样，它们也将自己所获得的关于神性的清晰概念传授给人类，从而使人摆脱对神性的怀疑！因此，\Person{策尔苏斯}既视鸟类优于人类，就应以它们为师，而不应以任何希腊哲学家为师。\par

% structure: p0189 source=h2 target=subsection reason=relative-heading-rank; base=h2

\subsection{第九十章}

但我们从众多回应中略述数端，以回答\Person{策尔苏斯}的这些言论，从而表明他持有这些错误意见所包含的对造物主的忘恩负义。因为\Person{策尔苏斯}虽是人，且\Quote{居于尊位}，却没有智慧，因此他不将自己与鸟类及其他他认为能占卜的无理性动物相比较；反而让它们居首，贬低自己，并尽可能将整个人类与他一同贬低（认为人类对天主的看法比无理性动物更低劣、更低下），低于那些将无理性动物当作神明崇拜的埃及人。然而，让调查的主要点在于：是否确实存在一种通过鸟类及其他被认为有这种能力的生物进行的占卜术？因为倾向于确立这两种观点的论证都不容轻视。一方面，我们被迫不接受这样的技艺，以免理性存在舍弃神圣神谕，转而求助于鸟类；另一方面，有许多人的有力见证，说无数人因信赖鸟占而从巨大危险中被拯救出来。然而，目前姑且承认存在占卜术，以便借此向对此有偏见的人表明：若承认这一点，则人相对于无理性动物——甚至相对于那些有占卜能力的——的优越性巨大，完全不可比拟。那么，我们必须说，若它们之中有任何神圣本性能够预言未来，且如此丰富（以那种知识）以致其富余足可向任何想知道未来的人传达，则显然它们对涉及自身之事知道得更快（比对他人之事）；若它们拥有这种知识，它们就会警惕不飞往人们设下圈套和网罗捕捉它们的地方，或弓箭手瞄准并射落它们的地方。尤其，若鹰事先知道针对其幼雏的计谋——无论是由爬蛇上巢毁灭它们，还是由人为娱乐或其他有用目的或服务而捕捉它们——它们就不会将幼雏置于会受攻击之处；一般说来，这些动物就不会被人捕获，因为它们比人类更神圣、更聪明。\par

% structure: p0191 source=h2 target=subsection reason=relative-heading-rank; base=h2

\subsection{第九十一章}

此外，若占卜鸟彼此交谈，如\Person{策尔苏斯}所坚持的那样，预言鸟具有神圣本性，而其他理性动物也有关于神性和未来事件的预知；且若它们曾将此知识传授给同类，那么\Person{荷马}所提到的麻雀就不会在蛇要吞吃它和它的幼雏之处筑巢，同一位诗人笔下的蛇也不会不防备被鹰捕获。因为这位奇妙诗人在关于前者的诗中说道：——\par

\Quote{一条巨蛇射出，是可怕的征兆；\newline{}宙斯亲自派出这令人惊惧的记号。\newline{}它直扑树身，盘绕血色螺纹，\newline{}蜷曲成层层盘旋的圈痕。\newline{}最高枝头有一只母鸟占据；\newline{}八只雏鸟填满苔藓的巢里；\newline{}母鸟是第九：那蛇悬垂，\newline{}张黑颚，咬碎垂死雏鸟的生命；\newline{}而母鸟在旁低徊，悲惨呻吟，\newline{}垂翅的哀鸣为失去的孩子哭泣。\newline{}最后，当母鸟绕巢飞扑，\newline{}被扑翅的翅膀抓住，怪物将其扼杀：\newline{}也不久存：化为石雕，他屹立\newline{}在\Place{阿乌利斯}的沙滩上，成为恒久的奇观。\newline{}这是宙斯的旨意；因此我们敢于\newline{}信赖他的预兆，支持战争。}\par

至于第二个——那鸟——诗人说：——\par

朱庇特的鸟鸣叫着拍打天空；\newline{}一条巨大的流血蛇，\newline{}它的爪子缠绕；活着的，蜷曲着的，\newline{}它刺伤了鸟，鸟的喉咙受了伤。\newline{}因刺痛而疯狂，它丢下致命的猎物，\newline{}在空中盘旋，痛苦地飞翔，\newline{}漂浮在风中，撕裂天空，发出哀鸣；\newline{}在军队中间，落下的蛇躺着。\newline{}他们因恐惧而苍白，看着它展开的螺旋，\newline{}心中怦怦跳地看着朱庇特的预兆。\par

那么，鹰拥有占卜的能力，而蛇（因为占卜者也使用这种动物）却没有吗？既然这种区分很容易被反驳，那么说两者都有占卜能力的断言难道也不能被反驳吗？因为如果蛇拥有这种知识，它难道不会防备自己遭受鹰带来的痛苦吗？还可以找到无数其他类似的例子，表明动物并不拥有预言的灵魂，而是像诗人和大多数人所认为的那样，是\Quote{奥林匹斯神自己把他送到光明里}。而且，阿波罗使用鹰作为他的使者具有象征意义，因为鹰被称为\Quote{阿波罗的迅捷使者}。\par

% structure: p0197 source=h2 target=subsection reason=relative-heading-rank; base=h2

\subsection{第九十二章}

然而在我看来，是某些邪恶的魔鬼（可以说是泰坦或巨人族）对真天主和天上的天使行了不敬，并从天上堕落，他们出没于身体较稠密的部分，常去地上不洁的地方；他们因为拥有脱离尘世物质的身体，而能分辨未来之事，从事这类职业，并渴望带领人类远离真天主。他们秘密进入最贪婪、最野蛮、最邪恶的动物的身体，随时煽动它们做任何事：要么扭转这些动物的幻想，使它们做出各种飞翔和动作，以便人类被非理性动物中的占卜能力所迷惑，忽略寻求包含万有的天主；要么使人寻求对天主的纯洁敬拜，却让他们的理性匍匐在地上，在鸟类、蛇类，甚至狐狸和狼中间。因为那些精通此类事情的人观察到，从这类动物身上得到最清晰的预兆；因为魔鬼不能在温和的动物中像在它们中那样有效地行动，由于它们在邪恶方面的相似性；然而这些动物中存在的并不是邪恶，而是像邪恶的东西。\par

% structure: p0199 source=h2 target=subsection reason=relative-heading-rank; base=h2

\subsection{第九十三章}

因此，在梅瑟的著作中，凡是令我惊叹的，我要说以下一点值得赞美：即梅瑟观察了动物的不同本性，并且要么从天主那里了解了它们各自的特点以及与每种动物有亲缘关系的魔鬼，要么凭自己的智慧亲自发现了这些。他在安排不同种类的动物时，宣布所有埃及人和其他人认为具有占卜能力的动物为不洁；一般来说，所有不属于那一类的动物为洁净。在梅瑟提到的不洁动物中，有狼、狐狸、蛇、鹰、隼，以及诸如此类。一般而言，你会在法律和先知书中发现，这些动物被用作一切最邪恶之事的例子；狼或狐狸从来没有被用于好的目的。因此，每一种魔鬼似乎与某一种动物有某种亲和力。正如在人类中，有些人比其他人更强壮，而这与他们的道德品格毫无关系；同样，一些魔鬼在无关紧要的事情上会比其他魔鬼更有能力；其中一类按照我们圣经中称为\Quote{今世的首领}的意愿，使用一种动物来欺骗人，而另一类则通过另一种动物预言未来。此外，注意魔鬼的邪恶达到了何种程度，他们甚至取用黄鼠狼的身体来揭示未来！现在，请你自己思考一下：是相信至高天主和祂的圣子（而非人类，虽然他们就在它们面前）激动鸟类和其他生物去占卜，还是相信那些激动这些生物的是邪恶的、按我们圣经所称为\Quote{不洁的魔鬼}呢？\par

% structure: p0201 source=h2 target=subsection reason=relative-heading-rank; base=h2

\subsection{第九十四章}

但是，如果鸟类的灵魂因为它们预言未来而被视为神圣，那么我们为什么不主张，当人类接受预兆时，通过其听到预兆的那些人的灵魂也是神圣的呢？相应地，在这些人中，荷马提到的那个磨谷物的女奴，当她说关于求婚者的话时，也应被列入其中：——\par

\Quote{因为这将是我们最后一次在这里用餐。}\par

这个女奴，因此是神圣的，而伟大的尤利西斯——荷马的帕拉斯·雅典娜的朋友——却\Emph{不是}神圣的，但他理解了这位\Quote{神圣的磨谷者}所说的话作为预兆，并因此欢喜，正如诗人所说：——\par

\Quote{神圣的尤利西斯因这预兆而欢喜。}\par

现在，既然鸟类拥有神圣的灵魂，并且能够认识天主，或者按照刻尔苏斯的说法，认识众神，那么显然，当我们人类打喷嚏时，也是由于我们内在的某种神性，它赋予我们的灵魂预言的能力。这个信念有许多见证证实，因此诗人也说：——\par

\Quote{当他祈祷时，他打了一个喷嚏。}\par

并且珀涅罗珀也说：——\par

\Quote{难道你没有注意到，我儿子每说一句话就打个喷嚏吗？}\par

% structure: p0210 source=h2 target=subsection reason=relative-heading-rank; base=h2

\subsection{第九十五章}

真正的天主并不使用无理性的动物，也不使用任何偶然遇到的人来传达对未来的知识；相反，祂使用最纯洁、最圣善的人灵，祂激励他们并赋予他们先知之恩。因此，在梅瑟著作中，无论其他什么可能引起我们的惊奇，以下内容必定被视为如此：\BibleQuote{不可占卜，也不可观察飞鸟}；在另一处：\BibleQuote{因为上主你的天主将从你面前消灭的异民，他们听信预兆和占卜；但上主你的天主却不容许你这样行。}他还补充说：\BibleQuote{上主你的天主要从你弟兄中兴起一位先知。}此外，有一次，天主意欲借一个占卜者使祂的子民远离占卜之事，便使占卜者心中的神灵这样说：\BibleQuote{因为在雅各伯中没有邪术，在以色列中没有占卜。到了适当的时候，必将向雅各伯和以色列宣告上主将行的事。}现在，我们这些知道这些和类似言论的人，愿意以奥秘的意义遵守这条诫命，即：\BibleQuote{你要全心守护你的心}，使任何属于魔鬼的本性不可进入我们的思想，也使我们仇敌的灵不可随意扭转我们的想象。但我们祈求，认识天主荣耀的光照进我们的心，祈求天主的圣神住在我们的想象中，引导它们默观天主的事；因为\BibleQuote{凡受天主圣神引导的，都是天主的子女。}\par

% structure: p0212 source=h2 target=subsection reason=relative-heading-rank; base=h2

\subsection{第九十六章}

然而，我们应当注意，预知未来的能力绝不是神性的证明；因为就其本身而言，它是一种中性的事物，在好人坏人身上都会出现。无论如何，医生凭借专业技能预知某些事情，尽管他们的品格可能不好。同样，舵手们，尽管也许是恶人，也能预兆（好或坏天气的）迹象，以及猛烈风暴的来临和大气变化，因为他们从经验和观察中积累了这些知识，尽管我不认为因此任何人会称他们为\Quote{神}，如果他们的品格恰好是坏的。那么，策尔苏斯的断言是假的，他说：\Quote{还有什么比预知和预告未来的能力更神圣的呢？}同样，他的这一断言也是假的：\Quote{许多动物声称对天主有认识}；因为没有任何无理性的动物拥有对天主的任何认识。他的另一个断言也是完全虚假的：\Quote{无理性的动物比（人）更接近天主的团体}；实际上，即使仍处于邪恶状态的人，无论他们在知识上有多大进步，都远离那个团体。因此，只有那些真正有智慧且真诚虔诚的人更接近天主的团体；例如我们的先知们，以及梅瑟，后者因其极其纯洁，圣经说：\BibleQuote{只有梅瑟可以走近主，其余的人不得走近。}\par

% structure: p0214 source=h2 target=subsection reason=relative-heading-rank; base=h2

\subsection{第九十七章}

这个指责我们不敬的人，他的断言是多么不敬啊！他说：\Quote{无理性的动物不仅比人类更聪明，而且比人类更蒙天主所爱！}有谁会不（因恐惧）而退避，不愿理会一个宣称蛇、狐狸、狼、鹰、隼比人类更蒙天主所爱的人呢？因为他坚持这种立场的结果是，如果这些动物比人类更蒙天主所爱，那么显然它们比苏格拉底、柏拉图、毕达哥拉斯、斐瑞居德以及那些他不久前还颂扬的神学家们更蒙天主所爱。有人可能向他祈祷说：\Quote{如果这些动物比人更蒙天主所爱，愿你与它们一同蒙天主所爱，并变得像你所认为比人类更蒙天主所爱的那样！}愿没有人认为这样的祈祷是一种诅咒；因为谁不愿意在各方面都像那些他认为比别人更蒙天主所爱的人，好使他像他们一样享受天主的爱呢？由于策尔苏斯渴望表明无理性动物的集会比我们的更神圣，他声称这个说法不是出自普通人，而是出自有智慧的人。然而，只有有德行的人才是真正有智慧的，因为恶人不是。因此，他用以下方式说话：\Quote{有智慧的人说，这些动物举行比我们更神圣的集会，并且它们知道集会上所说的话，实际上证明了它们并非缺乏这种知识，因为它们事先提到鸟已宣布它们要飞往某个地方，或者要做这件事或那件事，然后表明它们\Emph{确实}飞到了那个地方，并做了预言的那件事。}现在，确实，没有一个有智慧的人曾讲述过这样的事情；也没有任何一个智者说过无理性动物的集会比人的更神圣。但是，如果为了检验他的说法而看其后果，显然，在他看来，无理性动物的集会比可敬的斐瑞居德、毕达哥拉斯、苏格拉底、柏拉图以及一般哲学家的集会更神圣；这个断言不仅本身不合理，而且充满荒谬。然而，为了让我们相信某些人\Emph{确实}从鸟的模糊叫声中知道它们将要离开，并且要做这件事或那件事，并事先宣布这些事情，我们会说，这些信息是由魔鬼通过征兆传递给人的，目的是让人被魔鬼欺骗，使他们的理解力从天主和天堂拖到地上，甚至更低的地方。\par

% structure: p0216 source=h2 target=subsection reason=relative-heading-rank; base=h2

\subsection{第九十八章}

此外，我不知道塞尔苏斯是如何听闻到大象的（忠信于）誓言，以及它们对我们天主的巨大忠诚，和它们对祂所拥有的认识。因为我知道许多关于这种动物本性的奇妙之事，以及它温顺的性情。但我不知道是否有人谈到过它对誓言的遵守；除非确实，他称其对温顺性情和（可以说）一旦与人类达成契约的遵守为信守誓言，而这说法本身也是虚假的。因为尽管罕见，但有时仍有记载，它们看似驯服之后，却以最野蛮的方式攻击人类，并犯下杀人罪行，因而被判处死刑，因为不再有任何用处。看到这之后，为了建立（如他所认为的）鹳鸟比任何人类更虔诚的观点，他引用了关于那种生物以孝爱之情为父母觅食以供养它们的叙述。我们不得不回答说，鹳鸟这样做并非出于对正当的考虑，也不是出于反省，而是出于一种自然的本能；形成它们的本性希望在一个无理性动物中展示一个例子，这个例子可以在对父母感恩的表现上让人类感到羞愧。如果塞尔苏斯知道出于理性和出于无理性的自然冲动这样做之间的巨大差异，他就不会说鹳鸟比人类更虔诚。但进一步，塞尔苏斯，仍然为无理性受造物的虔诚辩护，引用了阿拉伯鸟凤凰的例子，它许多年后回到埃及，在那里带着它已死并被裹在没药球中的父母，将其身体放在太阳庙中。现在这个故事确实有记载，如果它是真的，那么它可能源于某种自然的安排；天主眷顾自由地向人类展示，通过生物之间的差异，世界上普遍存在的结构多样性，这种多样性甚至延伸到鸟类，并且与这种多样性相协调，祂创造了一个生物，它唯一的一种，以便通过它，人们可以被引导去钦佩，不是受造物，而是创造它的那一位。\par

% structure: p0218 source=h2 target=subsection reason=relative-heading-rank; base=h2

\subsection{第九十九章}

\Quote{因此，万物被造并非为了人，也不比为狮子、鹰或海豚造的，而是为了使这个世界，作为天主的工程，在各方面完美而完整。为此，一切事物都已调整，不是相互参照，而是关于它们对整体的影响。天主眷顾整体，祂的眷顾永远不会舍弃它；它不会变得更坏；天主也不会在某个时间后把它带回自己那里；祂也不会因为人而发怒，也不比为猿猴或苍蝇发怒更多；祂也不威胁这些存在，它们每个都在其适当位置接受了指定的命运。} 让我们然后简要地回答这些陈述。我认为，确实，我在前文中已经表明，万物是为人和每一个理性受造物而造的，并且创造主要是为了理性受造物的缘故。塞尔苏斯确实可能说，这并不比为狮子或他提到的其他受造物更多；但我们坚持，造物主形成这些东西不是为狮子、鹰或海豚，而是全部为了理性受造物的缘故，并且\Quote{为了使这个世界，作为天主的工程，在各方面完美而完整。}因为我们必须赞同这个观点，说得好。天主眷顾的，并非如塞尔苏斯所认为的仅仅是\Emph{整体}，而是超过整体，在特殊程度上眷顾每一个理性受造物。天主眷顾永远不会舍弃整体；因为即使整体因作为整体一部分的理性受造物的罪而变得更邪恶，祂也安排来净化它，并在一个时间后把整体带回祂自己那里。此外，祂不因猿猴或苍蝇发怒；但对于人类，作为那些违背自然律的人，祂派遣审判和惩罚，并通过先知的口，以及通过来到整个人类中的救主威胁他们，使那些听到威胁的人得以悔改，而那些忽视这些悔改呼唤的人应当受那些惩罚，这些惩罚符合天主的意愿，该意愿为了整体的益处，施加于那些需要这种痛苦管教和纠正的人。但既然我们第四卷现在已达到足够的篇幅，我们将在此结束我们的论述。愿天主通过祂的儿子，即\OriginalTerm{天主圣言}{God the Word}、\OriginalTerm{智慧}{Wisdom}、\OriginalTerm{真理}{Truth}、\OriginalTerm{正义}{Righteousness}，以及圣经论及天主时所称呼的一切，赐予我们为读者的益处良好地开始第五卷，并借由祂那常存于我们灵魂中的圣言，成功地完成它。\par

\SourceRecord{Translated by Frederick Crombie.；Ante-Nicene Fathers；Vol. 4.；Edited by Alexander Roberts, James Donaldson, and A. Cleveland Coxe.；Buffalo, NY: Christian Literature Publishing Co.,；1885.；<http://www.newadvent.org/fathers/04164.htm>.}

% structure: p0001 source=h1 target=section reason=page-title

\section{《驳策尔苏斯》卷五}

% structure: p0002 source=h2 target=subsection reason=relative-heading-rank; base=h2

\subsection{第一章}

我尊敬的\Person{盎博罗削}，我们开始第五卷回应\Person{策尔苏斯}论著的工作，并非因为我们追求多言——这是被禁止的，并且沉溺其中难免犯罪——而是尽力在我们能力范围内，不放过他任何未经审查的言论，尤其是那些可能让一些人觉得他巧妙攻击我们和犹太人的话。倘若可能，我确实希望能与我的话语一同进入每一位阅读此书的读者内心，拔出那支伤害尚未完全披戴\Quote{天主的全副武装}之人的箭，并施以理性的药物来治愈\Person{策尔苏斯}造成的创伤——这创伤使听从他话语的人无法保持\Quote{信德上的健全}。但既然唯独天主的工作，凭祂自己的圣神并藉基督的圣神，无形地居住在祂认为堪当受访的人心中；那么，另一方面，\Emph{我们的}目标是藉论证和著作努力坚定人的信德，并配得称为\Quote{无需羞愧的工人，正确分解真理之言}。有一件事我们认为尤其必须做，如果我们愿忠实地履行您托付给我们的任务，那就是尽我们所能推翻\Person{策尔苏斯}的自信断言。那么，让我们引用他那些在我们已驳斥之后的断言（读者须判断我们是否成功），并回应它们。愿天主恩赐我们不以缺乏神感、空虚的理智和理性来接近主题，以便我们所愿帮助之人的信德不倚靠人的智慧，而是从祂的父——唯独祂能赐予——领受基督的\Quote{心意}，并因参与天主圣言而坚强，从而推倒\Quote{一切高举起来反对天主知识的自高之事}，以及\Person{策尔苏斯}那高举起来反对我们、反对耶稣、也反对梅瑟和先知的想象；为的是那位\Quote{将话语赐给以大力宣扬它的人}，也供应我们，并赐给我们\Quote{大力}，使对天主圣言和德能的信德种植在所有阅读我们作品之人的心中。\par

% structure: p0004 source=h2 target=subsection reason=relative-heading-rank; base=h2

\subsection{第二章}

现在我们必须驳斥他如下的话：\Quote{噢，犹太人和基督徒，没有天主或天主子曾降临或将降临（地上）。但如果你们意指某些天使曾如此，那么你们称他们为什么？是神，还是另一种族？无疑是另一种族，而且很可能是魔鬼。} 既然\Person{策尔苏斯}在这里重复自己（在前文他已屡次提出此类断言），我们无需赘述，因我们已就此点所言已足。然而，我们将从更多论述中提及少数几点，我们认为与先前论证和谐，但又不完全同一角度；我们将由此表明，他普遍断言没有天主或天主子曾降临（人间），不仅推翻了多数人对天主显现的看法，也推翻了他自己先前承认的观点。因为如果\Quote{没有天主或天主子曾降临或将降临}这一普遍说法被\Person{策尔苏斯}确实坚持，那么显然我们在此推翻了相信有从天降世的神祇存在，他们或向人类预言未来，或藉神谕医治人类；那么\Person{皮提亚的阿波罗}、\Person{埃斯库拉庇乌斯}以及任何其他被认为如此行的神祇，都不会是天降之神。他们或许是为永远居于地上而抽得（义务）的神，因而像是从神界逃出来的；或是无力分享天上神群的神；否则\Person{阿波罗}、\Person{埃斯库拉庇乌斯}以及其他被认为在地上行事者，就不是神，而只是某些魔鬼，远比那些因德行升上天穹的人中之智者低劣。\par

% structure: p0006 source=h2 target=subsection reason=relative-heading-rank; base=h2

\subsection{第三章}

但请注意，当他渴望推翻我们观点时，这位在他整篇论著中从未承认自己是伊壁鸠鲁派的人，被证明是该派系的背弃者。现在，阅读\Person{策尔苏斯}著作并表示认同的（读者），你或者推翻相信一位眷顾人类并对每人施行天主眷顾的天主，或者承认这一点并证明\Person{策尔苏斯}断言为假。如果你完全消灭天主眷顾，那么你就要证明他那些承认\Quote{天主和眷顾}存在的断言为假，以维护你自己立场的真理；但另一方面，如果你仍承认天主眷顾存在，因为你不同意\Person{策尔苏斯}的断言——\Quote{没有天主或天主子曾降临或将降临人间}——那么为什么不在耶稣的言行和关于祂的预言中仔细查考，我们应当认为那一位作为天主和天主子降临人间的是谁？——是那说了并行了如此多事的耶稣，还是那些以神谕和占卜为借口、不改革崇拜者品德的，反而从对万有创造者的纯正圣洁崇拜和尊荣中背弃，并在假装敬奉众多神祇的幌子下，将听信他们者的灵魂从唯一可见的真天主那里撕裂开来的那些？\par

% structure: p0008 source=h2 target=subsection reason=relative-heading-rank; base=h2

\subsection{第四章}

但既然他在下一处说，仿佛犹太人或基督徒在回答关于那些降临访问人类者时，说他们是天使：\Quote{但若你说他们是天使，你称他们为什么？}他接着说，\Quote{他们是神，还是另外一类存有？}然后他又介绍我们，仿佛在回答：\Quote{是另外一类存有，大概是魔鬼，}——让我们继续注意这些评语。因为我们确实承认天使是\Quote{服役的神灵}，并且我们说\Quote{他们被派遣为那些将要承受救恩的人服务}；他们上升，带着人们的恳求，到达宇宙中最纯洁的天上之处，甚至更纯洁的超天领域；他们又从这些地方下降，按照各人的功过，传递天主命令他们赐予那些将接受祂恩惠之人的某些事物。我们由此从他们的职务学会称这些存有为\Quote{天使}，我们发现，由于他们是神圣的，在圣经中他们有时被称为\Quote{神}，但并非因此命令我们代替天主来尊敬和崇拜那些服侍我们、并将祂的祝福带给我们者。因为一切祈祷、恳求、转求和感恩，都应当通过那位超越所有天使的大司祭，就是生活的圣言和天主，上达于至高天主。而我们也应向圣言本身祈祷、转求、献上感恩和恳求，如果我们有分辨祈祷正用与滥用的能力。\par

% structure: p0010 source=h2 target=subsection reason=relative-heading-rank; base=h2

\subsection{第五章}

因为呼求天使，却没有获得超越人类认知的天使本性的知识，是不合理的。但按照我们的假设，让这关于他们的知识，即某种奇妙而奥秘的知识，被获得。那么，这知识使我们认识他们的本性和各自被委任的职务，就不会容许我们满怀信心地向任何其他对象祈祷，除非是向那全备一切的至高天主，而且是经由我们的救主天主子，即圣言、智慧、真理以及天主先知和耶稣宗徒的著作所赋予祂的一切称号。要使天主的圣天使对我们眷顾，并为我们做一切事，就足以让我们对天主的心态，在人性能力范围内，效法那些圣天使的榜样——他们又效法其天主的榜样；并且我们关于祂的圣子、圣言所具有的观念，就我们所能达到的程度而言，不应与圣天使所拥有的关于祂的更清晰的观念相抵触，而应在清晰度和明确性上日渐接近。但因为\Person{塞尔苏斯}没有读过我们的圣经，他自行给出一个仿佛出自我们之口的回答，说我们\Quote{断言那些从天上降临为人类施恩的天使是与神不同的种族}，并补充说\Quote{它们很可能被我们称为魔鬼}：没有注意到\Quote{魔鬼}这个名称并非像\Quote{人}那样是中性的词——人中有善有恶；也不是像\Quote{神}那样是卓越的词——它不被用于邪恶的魔鬼、雕像或动物，而是（被认识神圣事物的人）用于真正神圣和受祝福者；而\Quote{魔鬼}一词总是用于那些邪恶的势力，它们脱离较粗重身体的束缚，引诱人误入歧途，使人充满纷扰，并把人从天主和超天思想拖到下方的事物。\par

% structure: p0012 source=h2 target=subsection reason=relative-heading-rank; base=h2

\subsection{第六章}

接着，他对犹太人作了如下陈述：— \Quote{关于犹太人，第一个令人惊奇之处是，他们应当敬拜天和居住其中的天使，却忽略和漠视其最可敬和有力的部分，如太阳、月亮和其他天体，包括恒星与行星，仿佛可能\InnerQuote{整体}是神，而其各部分却不神圣；或者（仿佛合理）以最高敬意对待那些据说在某个黑暗处显现给某些人者——他们被某种歪曲的巫术弄瞎，或受阴影幻影的影响而做梦——而那些如此清晰而显著地向所有人预报者，藉着他们，雨、热、云、雷（他们向之献祭）、闪电、果实和一切丰产得以产生——藉着他们，天主向自己启示——那些在上方存在者中最杰出的使者——那些真正天上的天使——却被视为无足轻重！}在作这些陈述时，\Person{塞尔苏斯}似乎陷入了混乱，并且是基于对他所不了解之事的错误观念而写下它们；因为凡考察犹太人的实践并将其与基督徒相比校的人，都明显看出：那些遵循法律的犹太人——这法律以天主的口气说：\BibleQuote{你不可有其他神在我面前：你不可为自己雕刻偶像，也不可制造任何天上或地上或水中的形象；你不可向它们跪拜，也不可事奉它们}——除了创造天及万有的至高天主之外，不事奉任何其他对象。如今很明显，那些按法律生活、并敬拜天的造物主的人，不会同时敬拜天与天主。此外，任何服从梅瑟法律的人都不会向天上的天使跪拜；同样，他们既不向太阳、月亮和星辰——天上的万军——跪拜，他们也避免向天及其天使致敬，遵守那宣告的法律：\BibleQuote{恐怕你举目望天，看见太阳、月亮和星辰，就是天上的万军，就被勾引去敬拜事奉它们，那是上主你的天主分给万民的。}\par

% structure: p0014 source=h2 target=subsection reason=relative-heading-rank; base=h2

\subsection{第七章}

此外，他假定犹太人认为天是天主，便说这是荒谬的；他指责那些叩拜天、却不也叩拜太阳、月亮和星辰的人，说犹太人这样做，仿佛可能让\Quote{整体}应该是天主，而它的各部分却非神性。他似乎称天为\Quote{一个整体}，而太阳、月亮、星辰是其各部分。当然，犹太人和基督徒都不称\Quote{天}为天主。然而，姑且承认他所说的，天被犹太人称为天主，并假设太阳、月亮、星辰是\Quote{天}的部分——这绝非真实，因为地上的动物和植物也不是天的一部分——那么，即使按照希腊人的见解，如果天主是一个整体，祂的部分也必然是神性的吗？当然，他们说宇宙作为整体是天主，斯多葛主义者称之为第一天主，柏拉图学派称之为第二天主，其中一些称之为第三。那么，按照这些哲学家的说法，既然整个宇宙是天主，其各部分也是神性的；所以不仅人类是神性的，而且整个无理性的受造物，作为宇宙的\Quote{ \Emph{部分} }，也是神性的；此外，植物也是神性的。如果河流、山脉和海洋是宇宙的部分，那么既然整个宇宙是天主，河流和海洋也是神吗？但就连希腊人也不会断言这一点。然而，那些主管河流和海洋的（无论是他们称为的魔鬼或神），他们会称之为神。由此可知，塞尔苏斯的一般论述，即使按照持有天主眷顾教义的希腊人，也是错误的：如果某个\Quote{整体}是神，其各部分就必然是神性的。但根据塞尔苏斯的教义，如果宇宙是天主，那么宇宙中的一切，作为宇宙的部分，都是神性的。那么，按照这种观点，动物如苍蝇、蚊子、虫子以及各种蛇类，还有鸟类和鱼类，都将是神性的——这种断言连那些主张宇宙是天主的人也不会做出。但是，按照梅瑟法律生活的犹太人，虽然他们可能不知道如何接受以隐晦语言传达的法律的奥秘含义，却不会主张天或天使是天主。\par

% structure: p0016 source=h2 target=subsection reason=relative-heading-rank; base=h2

\subsection{第八章}

然而，我们断言，他由于吸收了错误观念而陷入混乱；来吧，让我们尽力指出这些，并表明虽然塞尔苏斯认为向天和其中的天使叩拜是犹太人的习俗，但这样的做法根本不是犹太人的，而是违反犹太教，正如向太阳、月亮、星辰以及偶像敬拜一样。至少在\WorkTitle{耶肋米亚书}中，你会找到天主藉先知的口责备犹太人民向这些对象敬拜，并向天后和天上万军献祭的话。此外，基督徒的著作在责备犹太人中犯的罪时，也表明当天主因某些罪而抛弃那个民族时，这些（偶像崇拜的）罪也是他们犯下的。因为在\WorkTitle{宗徒大事录}中关于犹太人记载说：\BibleQuote{天主就转身，任凭他们事奉天上的星辰，正如《先知书》上所写的：以色列家！你们在旷野四十年，何尝向我奉献过牺牲和祭品？你们抬着摩肋客的帐幕，和勒番星的偶像，就是你们为朝拜而制造的。}\BibleRef{宗 7:42-43}而在\Person{保禄}的著作中——他受过犹太习俗的精心训练，后来因\Person{耶稣}奇迹般的显现而皈依基督信仰——在\WorkTitle{哥罗森书}中可以读到以下话语：\BibleQuote{不要让那些专务谦卑和崇拜天使的人，凭着他们见到的幻象，妄自尊大，凭着血肉的意念，无端傲慢，剥夺你们的赏报；他们不持守头，可是整个身体，都由头借着关节和脉络，得到营养和联系，并赖天主的生长而生长。}\BibleRef{哥 2:18-19}但\Person{塞尔苏斯}既没有读过这些经文，也没有从任何其他来源了解它们的内容，我不知道他是如何认为犹太人并未违反他们的法律，向天和其中的天使叩拜。\par

% structure: p0018 source=h2 target=subsection reason=relative-heading-rank; base=h2

\subsection{第九章}

他仍然有点困惑，没有留意看什么与问题相关，便发表意见说，犹太人被幻术和巫术中所用的咒语（由于这些咒语，某些幻象出现，听从术士的咒语）诱导，向天上的天使叩拜，没有注意到这违反他们的法律，这法律对行这些事的人说：\BibleQuote{不可转向招魂的，也不可询问行巫术的，免得被他们玷污：我是上主你们的天主。}\BibleRef{肋 19:31}因此，他要么根本不应该将这种做法归于犹太人，因为他观察到他们遵守自己的法律，并称他们为\Quote{按照他们法律生活的人}；要么如果他归给他们，他就应该表明犹太人这样做违反了他们的法规。但另一方面，正如那些向据说在黑暗中和被巫术蒙蔽的人显现的幻象、以及做梦的人（由于模糊的幻象呈现）献上敬拜的人违反他们的法律一样，那些向太阳、月亮、星辰献祭的人也违反法律。因此，同一个人说犹太人格外遵守法律，不向太阳、月亮、星辰叩拜，而在天和天使的事上却不那么谨慎地遵守法律，这是很大的矛盾。\par

% structure: p0020 source=h2 target=subsection reason=relative-heading-rank; base=h2

\subsection{第十章}

如果我们必须为我们拒绝承认那些被希腊人称为\Quote{明显可见}的神祇为神（与天使、太阳、月亮和星辰同等）辩护，我们将回答：梅瑟的法律知道，这些后者已被天主分给天下万民，但未被分给那些被天主拣选为祂特选子民、超越地上万民的人。因为\WorkTitle{申命纪}上记载：\BibleQuote{当你举目望天，看见太阳、月亮、星辰，即天上万象时，切不要为它们所迷惑而朝拜事奉它们，这是上主你的天主分给了天下万民的。但上主将你们从铁炉中，从埃及领出来，使你们作祂特选的子民，如同今日一样。}因此，希伯来民族被天主称为\BibleQuote{特选的种族、王家的司祭、圣洁的邦国、属于主的子民}，关于这民族，主藉着对他说的声音向亚巴郎预言说：\BibleQuote{请你仰观天空，数算星辰，如果你能数算它们。}于是祂对他说：\BibleQuote{你的后裔也将这样。}他们既怀着这样的希望，即藉着理解和遵守天主的法律，要变得如同天上的星辰，便不会向那些他们本该类似的对象屈膝下拜。因为曾对他们说：\BibleQuote{上主我们的天主使我们人数增多；看，你们今日多如天上的星辰。}在\WorkTitle{达尼尔}书中也有以下关于那些将分沾复活之人的预言：\BibleQuote{那时，你的百姓中凡写在书上的都要得救。许多睡在尘土中的人要醒来，有的进入永生，有的蒙受羞辱和永久的厌恶。但凡有智慧的人，要发光如同穹苍的光辉；那教导多人归于正义的，要像星辰一样，直到永远。}等等。因此保禄在谈论复活时也说：\BibleQuote{有天上的形体，也有地上的形体；但天上形体的光荣是一种，地上形体的光荣是另一种。太阳有一种光荣，月亮有另一种光荣，星辰也有另一种光荣；而且星辰与星辰的光荣也有差异。死人的复活也是如此。}因此，那些曾受崇高教导要超越一切受造物、且因德行高尚而希望在天主那里享受最光荣的赏报的人，那些曾听见\BibleQuote{你们是世界的光}和\BibleQuote{你们的光也当这样照耀在人前，使他们看见你们的善行，光荣你们在天之父}的话的人，那些通过实践拥有这种灿烂不朽的智慧，或甚至获得了\BibleQuote{永恒之光的反映}的人，竟然因太阳、月亮和星辰的（仅仅）可见光而如此受影响，以致因着它们可感觉的光，自认为（尽管拥有如此伟大的理性的知识之光、真光、世界之光、人之光）在某种程度上低于它们，并向它们俯伏跪拜，这是不合情理的。因为如果它们要接受敬拜，应当被敬拜的原因并非那被众人赞赏的可感觉的光，而是由于理性和真光——如果天上的星辰确实是理性的、有德性的存在，并且已被那\BibleQuote{永恒之光的反映}的智慧以知识之光照亮的话。\par

% structure: p0022 source=h2 target=subsection reason=relative-heading-rank; base=h2

\subsection{第十一章}

但是，就连这理性的光本身，也不应由那看见并理解真光（藉分沾这光，这些也得以光照）的人来朝拜；也不应由那看见天主、即真光之父（论及他，经上说：\Quote{天主是光，在他内毫无黑暗。}）的人来朝拜。那些因太阳、月亮和星辰的光是可见的、属天的而敬拜它们的人，确实不会向地上的火花或灯盏下拜，因为他们看到那些被认为值得敬拜的对象，与火花和灯光相比，具有无可比拟的优越性。因此，那些明白天主是光，并领悟天主子就是\Quote{那照亮每一個人来到世上的真光}，且理解他如何说\Quote{我是世界的光}的人，理智上不会向那与天主（他是真光的光）相比如同一粒火花的太阳、月亮、星辰献上敬拜。我们这样说太阳、月亮和星辰，并非有意贬低天主创造之力的这些伟大工程，也不是像\Person{阿那克萨戈拉}那样称它们为\Quote{炽热的团块}，而是因为我们觉察到天主的性体及其独生子的性体那无可言喻的超越性，超越一切其他事物。并且，我们深信太阳本身、月亮和星辰都经由天主的独生子向至高的天主祈祷，因此我们认为向那些本身献上祈祷的存在者祈祷是不合适的；因为就连它们自己也宁愿我们将我们的恳求呈献给它们所祈祷的天主，而不是将恳求下达到它们自身，或者将我们的祈祷权分给天主和它们。在此，我可以就这一点使用一个比喻：我们的主和救主，有一次听到有人称他\Quote{善师}时，便将那说话的人指向他自己的父，说：\Quote{你为什么称我是善的？除了一位，即天主父，没有善的。} \BibleRef{玛 19:17} 既然圣父的爱子这样说，作为天主良善的肖像，是完全合理的，那么太阳为何不更有理由向那些向它下拜的人说：\Quote{你们为什么敬拜我？因为你们要敬拜主，你们的天主，惟独事奉他；}因为我和所有与我一起的，都事奉并敬拜他。虽然一个人可能没有太阳那样崇高，但尽管如此，让他向天主的圣言（他能医治他）祈祷，并更向他的父祈祷，因为父在古时也向义人\Quote{发出他的言语，医治了他们，拯救他们脱离了毁灭。}\par

% structure: p0024 source=h2 target=subsection reason=relative-heading-rank; base=h2

\subsection{第十二章}

因此，天主按其仁慈，并非以局部的方式，而是藉其眷顾，俯就人类；而天主子，不仅（在地上时），而且于\Emph{所有}时候，都与他的门徒同在，应验了这许诺：\Quote{看，我同你们天天在一起，直到今世的终结。}如果枝子若不留在葡萄树上，就不能结果实，那么显然，圣言的门徒，即圣言真葡萄树的理性枝子，若不留在真葡萄树——即天主的基督——内，就不能结出德行的果实；基督在地上以局部的方式与我们同在，并与世界各地依附他的人同在，也处处与那些不认识他的人同在。另一位由那写了福音的若望显明出来，当他以若翰洗者的身份说话时，他说：\Quote{有一位站在你们中间，是你们不认识的；他在我以后来。}并且，当那充满天地、曾说\Quote{我岂不充满天地吗？——上主说}的那位与我们同在，并临近我们时（因为我信他所说的：\Quote{我是临近的天主，不是遥远的天主——上主说}），却还寻求向太阳、月亮或星辰之一祈祷，其影响力并未达到全世界，这是荒谬的。但是，用刻尔苏斯自己的话说，姑且承认\Quote{太阳、月亮和星辰\Emph{确实}预言降雨、炎热、云彩和雷声，}那么，如果它们真的预言如此伟大的事物，我们岂不应该反而尊崇天主，因它们在发出这些预言时是他的仆人，并敬畏\Emph{他}而不是他的\Emph{先知}吗？那么，让它们预言闪电、果实和一切出产来临，并让所有这类事物受它们管理；然而我们不会因此敬拜那些自己献上敬拜的存在者，正如我们甚至不敬拜梅瑟和那些在他之后来自天主的先知，他们预言了比雨、炎热、云彩、雷声、闪电、果实和各种感官可见的出产更好的事物。不，即使太阳、月亮和星辰能预言比雨更好的事物，我们仍不会敬拜\Emph{它们}，而是要敬拜那些预言在它们内的\Emph{父}和天主的\Emph{圣言}，即它们的仆役。但就算它们是祂的传报者，并真是天上的使者，那么，我们岂不应该反而敬拜那惟独由它们宣告和传报的\Emph{天主}，而不是那些作为\Emph{传报者}和\Emph{使者}的存在者吗？\par

% structure: p0026 source=h2 target=subsection reason=relative-heading-rank; base=h2

\subsection{第十三章}

此外，\Person{采尔苏斯}认为我们视太阳、月亮和星辰为无物。关于这些，我们承认它们也\Quote{等待天主的子女显现}，目前因那使它们屈服于希望者的缘故，屈服于物质身体的\Quote{虚幻}。但如果\Person{采尔苏斯}读过我们论及太阳、月亮和星辰的无数其他经文，尤其是这些——\Quote{你们一切星辰，以及你，光明，请赞美祂}，以及\Quote{诸天诸天，请赞美祂}——他就不会说我们认为这些伟大的存在\Quote{大大赞美}上主天主为无物。\Person{采尔苏斯}也不知道这段经文：\Quote{受造之物等待天主的子女显现，怀有热切的期望。因为受造之物被屈服于虚幻之下，不是出于自愿，而是由于那使它们屈服于希望的那位；因为受造之物本身也要从腐败的束缚中得到释放，进入天主子女荣耀的自由。}让我们以这些话结束我们不崇拜太阳、月亮和星辰的辩护。现在，我们继续提出他随后的话，愿天主允许我们以真理之光照所暗示的论据来回应他。\par

% structure: p0028 source=h2 target=subsection reason=relative-heading-rank; base=h2

\subsection{第十四章}

以下是他所说的话：\Quote{他们愚蠢地认为，当天主如同厨师一般引入（要毁灭世界的）火时，所有其余的人类都会被烧尽，而只有他们能存留——不仅是当时还活着的人，还有早已去世的人，后者将从地中复活，穿着与生前相同的肉身；这样的希望不过是虫子所怀抱的。因为哪一种人的灵魂会渴望那曾屈服于腐败的身体？此外，你们有些基督徒也不认同这种观点，认为它极其卑劣、可憎且不可能；因为那种身体完全腐败后，怎能恢复到它原来的本性，回到它最初陷入瓦解的状态？既然无法回答，他们便逃到一个极其荒谬的托词，即天主无所不能。然而，天主\Emph{不能}做可耻的事，也不愿做违背祂本性的事；即使（按照你心中的邪恶）你渴望任何邪恶之事，天主也不会成就它；你不必立即相信它会实现。因为天主统治世界不是为了满足无度的欲望，或容许混乱无序，而是为了治理正直公义的本性。因为对于\Emph{灵魂}，祂确实能够赐予永生；而死的\Emph{身体}，正如\Person{赫拉克利特}所言，比粪土更无价值。然而，天主既不能也不会违背一切理性，宣称那充满不配提及之事的肉身将永远存在。因为祂是一切存在之物的理性，因此不能做任何违背理性或违背祂自身的事。}\par

% structure: p0030 source=h2 target=subsection reason=relative-heading-rank; base=h2

\subsection{第十五章}

现在请注意，在开头这里，他在嘲笑某些希腊人以不应轻视的哲学精神所探讨的世界焚毁学说时，如何试图让我们\Quote{仿佛将天主描绘为厨师一般，持有普遍焚毁的信仰}；他没有意识到，正如某些希腊人所认为的（或许他们是从古老的希伯来民族得到的信息），那降临到世界上的是一种炼净之火，并且也可能降临到每个需要火之惩戒和医治的人身上，这火确实\Emph{焚烧}，但不会\Emph{消灭}那些没有物质身体的人，而物质身体需要被那火消耗，这火焚烧并消耗那些在行为、言语和思想上用所谓的\Quote{建筑}中的木材、干草或秸秆建造了东西的人。而且，圣经说上主将像炼金之人的火和漂白者的碱皂一样，造访每个需要炼净的人，因为他们邪恶本性中涌出的罪恶物质之洪流混杂在他们里面；他们需要火，我是指，来炼净（渣滓）那些与铜、锡、铅混杂的人。愿意的人可以从\Person{厄则克耳}先知那里学到这一点。但我们说天主将火带给世界，不是像厨师，而是像天主，是那些需要火之纪律者的恩主，这将由\Person{依撒意亚}先知作证，在他的著作中记载了一个罪恶的民族这样被教训：\Quote{因为你有了炭火，坐在这火上；它们将是你的一种帮助。}圣经适切地适应了众多阅读它的人，因为它以隐晦的方式谈论悲伤和阴郁之事，以恐吓那些无法藉其他方法从罪恶洪流中被拯救的人，尽管如此，细心的读者仍会清楚地发现这些悲伤痛苦的惩罚对忍受它们的人所要达到的目的。然而，目前引用依撒意亚的话就足够了：\Quote{为了我名的缘故，我将显示我的愤怒，并将我的荣耀带给你，使我不毁灭你。}我们不得不以隐晦的话语提及不适合单纯信徒能力的问题，他们需要更简单的言语教导，以免我们显得忽视了\Person{采尔苏斯}的指控，即\Quote{天主引入（要毁灭世界的）火，如同一个厨师}。\par

% structure: p0032 source=h2 target=subsection reason=relative-heading-rank; base=h2

\subsection{第十六章}

根据上述所言，有明智的听众会明白，我们应当如何答复以下的话：\Quote{其余的所有人都将被彻底烧尽，只有他们独存。} 确实，如果这种事情被我们当中那些在圣经中被称为\Quote{世上愚拙的}、\Quote{软弱的}、\Quote{卑贱的}、\Quote{被轻视的}、\Quote{无有的}\BibleRef{格前1:27-28}的人所接受，那就不足为奇了，因为\Quote{天主乐意用愚拙的宣讲来拯救那些相信祂的人，因为世界凭自己的智慧没有认识天主}\BibleRef{格前1:21}——因为这种人不能清楚看出每段经文的意义，或者不愿花必要的功夫去查考圣经，尽管有\Person{耶稣}的诫命\Quote{你们查考经典}\BibleRef{若5:39}。此外，以下便是他关于天主将降于世界的火以及罪人将要遭受的惩罚的看法。也许，正如对儿童讲话时，有些事必须以适合他们幼稚状态的方式来讲述，以便把他们从年幼时就引向更好的生活；同样，对那些经文所称的\Quote{世上愚拙的}、\Quote{卑贱的}、\Quote{被轻视的}，惩罚经文那直接浅显的含义是合适的，因为他们无法接受任何其他方式的归化，只能通过恐惧和惩罚的呈现，从而从许多（将要降临的）灾祸中得救。因此，圣经宣告，只有那些未被火和惩罚伤及的人将存留——即那些观点、道德和心智已被净化到最高程度的人；而另一方面，那些本性不同的人——即那些按其所应得，需要受火惩罚的人——将卷入这些苦难中，其目的是天主适宜加给那些按祂肖像受造却违背那肖像本性之意志的人。这就是我们对\Quote{其余的所有人都将被彻底烧尽，只有他们独存}这一说法的答复。\par

% structure: p0034 source=h2 target=subsection reason=relative-heading-rank; base=h2

\subsection{第十七章}

接着，他要么自己误解了圣经，要么误解了那些（解释者）——他们也没有理解圣经——他进而断言，\Quote{我们声称，在将要通过净化之火降临世界的审判时期，不仅那时活着的人，而且早已死去的人，都将存留。}他没有注意到，\Person{耶稣}的宗徒以隐秘的智慧说过：\BibleQuote{我们众人都不死亡，却都要改变，在顷刻之间，眨眼之间，在末次号角吹响时：号角一响，死人要复活成为不朽的，我们都要改变。}\BibleRef{格前15:51-52} 他本应注意到说这些话者的意思，此人绝非死人，他区分了自己和与他相似的人以及死人，随后说：\Quote{死人要复活成为不朽的，}和\Quote{我们要改变。} 为证明宗徒写下我所引用的\WorkTitle{格林多人前书}那些话时就是这个意思，我还要引用\WorkTitle{得撒洛尼前书}，其中\Person{保禄}作为一个活着、醒着、不同于那些睡着的人，这样说道：\BibleQuote{我们照主的话告诉你们：我们活着还存留到主来临的人，决不会在已死的人以前。因为主必亲自由天降来，有呼喊声、总领天使的呼声和天主的号声。}\BibleRef{得前4:15-16} 然后，在此之后，他知道除了自己和像他这样的人之外，还有在基督内死了的人，就附加了这些话：\BibleQuote{那些在基督内死了的人要先复活；然后我们这些活着还存留的人，同时与他们一起被提到云中，在空中迎接主。}\BibleRef{得前4:16-17}\par

% structure: p0036 source=h2 target=subsection reason=relative-heading-rank; base=h2

\subsection{第十八章}

但既然他长篇大论地嘲笑那在教会中宣讲的、关于肉身复活的教义——这教义被更明智的信徒更清楚地理解——而且既然没有必要重复他的原话（已经引述过了），那么，让我们针对这个问题（就像在一部针对信仰陌生人的护教作品中，为着那些仍是\Quote{小孩子，被各种教义之风飘荡、卷去，中了人的阴谋和诡计}\BibleRef{弗4:14}的人），竭尽全力陈述并确立一些明确为读者所写的内容。因此，我们或圣经都不断言，那些早已死去的人将带着同样的身体、不经改变为更高状态\Quote{从地里起来重新活着}；因为这样说，是在错误地指控我们。我们可以聆听许多圣经段落以合宜于天主的方式论述复活，尽管目前引述\Person{保禄}在\WorkTitle{格林多人前书}中的话就足够了，他说：\BibleQuote{但有人会说：死人怎样复活呢？他们将带着什么样的身体来呢？糊涂人哪！你所播种的，若不先死去，决不会活过来；而且你所播种的，并不是那将来要生的形体，而是一颗赤裸的籽粒，比如麦粒或别的谷粒；但天主随自己的心意给它一个形体，给每个种子以它自己的形体。}\BibleRef{格前15:35-38} 现在，请注意，在这些话中，他说所播种的\Quote{并不是那将来要生的形体}；而是说，当身体被播种、赤裸地埋入土中（天主给每个种子以它自己的形体）时，就仿佛发生了一种复活：从播在地里的种子中长出茎秆，例如在芥菜等植物中，或更大的树，如橄榄树，或某一种果树。\par

% structure: p0038 source=h2 target=subsection reason=relative-heading-rank; base=h2

\subsection{第十九章}

\Person{天主}就按他自己的意愿给每个东西一个合适的身体：如同在播种的植物中那样，也如同那些在死亡中，可以说是被播种的存在者一样，他们在适当的时候，从所\Quote{播种}的，得到\Person{天主}按各人的功过所指定的身体。此外，我们还可以听到圣经用这些话语详细教导我们那可以说是\Quote{播种的}和从它\Quote{复活起来的}之间的区别：\BibleQuote{播种的是可朽坏的，复活起来的是不可朽坏的；播种的是可羞辱的，复活起来的是光荣的；播种的是软弱的，复活起来的是强健的；播种的是属生灵的身体，复活起来的是属神的身体。}愿有领悟能力的人理解这些话的含义：\BibleQuote{凡是属土的，就是属于那属土的；凡是属天的，就是属于那属天的。我们怎样带着那属土的形象，也要怎样带着那属天的形象。}虽然宗徒希望隐藏这段经文的隐秘含义——它不适合较单纯的信徒和平民的理解，他们因信德而进入更好的生活——但他后来仍不得不加上这些话（以免我们误解他的意思），在\Quote{让我们带着那属天的形象}之后，还说了这些：\BibleQuote{弟兄们，我告诉你们：血肉之身不能承受\Person{天主}的国，可朽坏的也不能承受不可朽坏的。}然后，知道经文中有一种隐秘和奥秘的含义，一位在他书信中为后人留下充满深意的话语的人，又附加了以下的话：\BibleQuote{看，我告诉你们一个奥秘。}这是他介绍更深刻、更奥秘的事物的惯用方式，这些事物正当地对群众隐藏，正如在\BibleBook{多俾亚}{传}中所写的：\BibleQuote{隐藏君王的秘密是好事，但揭示\Person{天主}的事工则是光荣的。}——以符合真理和\Person{天主}荣耀的方式，并有利于群众。因此，我们的希望不是\Quote{蠕虫的希望，我们的灵魂也不渴望一个已经腐朽的身体}；因为虽然它可能需要一个身体为了移动位置，但它明白——正如默想（从上而来的）智慧，遵从\Quote{义人的口要讲论智慧}——那\BibleQuote{地上的帐棚}，即那将要拆毁的建造的会幕，与那义人在其中叹息、被重压，不愿\BibleQuote{脱下}帐幕而是\BibleQuote{穿上}它，好使那有死的被生命吞没之间的区别。因为，由于整个身体的本性是腐朽的，那腐朽的帐棚必须穿上不朽；它的另一部分，即有死的并且变得容易遭受随罪而来的死亡，必须穿上不朽，以便当那腐朽的穿上了不朽，那有死的穿上了不朽时，先知们从前所预言的就会实现——死亡之\BibleQuote{胜利}（因为它曾征服我们并使我们受其统治）和它的\BibleQuote{毒刺}（用它刺那防备不周全的灵魂，并给它加上罪所带来的创伤）的消灭。\par

% structure: p0040 source=h2 target=subsection reason=relative-heading-rank; base=h2

\subsection{第二十章}

但是，关于复活，就时间所允许的程度，我们在本次场合已部分陈述了我们的观点（因为我们在我们著作的其他部分更系统地考察了这个主题）；如今，我们必须通过健全的推理驳斥\Person{塞尔斯}的谬误，他既不了解我们\WorkTitle{圣经}的意义，也没有能力判断我们智者的意义不应由那些除了（单纯的）对基督徒体系的信德外别无任何宣称的人来决定——让我们指出，那些因其推理能力和辩证技巧而不应被轻看的人，已经表达出了非常荒谬的意见。如果我们必须嘲笑它们是卑劣的老妇故事，那么我们嘲笑的对象应该是他们而不是我们的叙述。\OriginalTerm{廊下派}{Porch}的门徒断言，经过若干年后世界将发生一场大火，之后一切事物将重新安排，与先前的世界安排毫无差别。然而，他们中那些对此学说表示敬意的人说，在周期结束时，与先前周期相比，将会有一种变化，尽管极其微小。这些人主张，在接下来的周期中，同一事物将再次发生，\Person{苏格拉底}将再次成为\Person{索福隆尼斯科斯}的儿子，身为\Place{雅典}本地人；\Person{菲娜瑞特}与\Person{索福隆尼斯科斯}结婚后，将再次成为他的母亲。尽管他们没有提及\Quote{\OriginalTerm{复活}{resurrection}}一词，但实际上他们表明，源于种子的\Person{苏格拉底}将从\Person{索福隆尼斯科斯}的种子中生出来，并在\Person{菲娜瑞特}的子宫中成形；在\Place{雅典}长大，将从事哲学研究，仿佛他先前的哲学复活了，并且与先前毫无差别。\Person{安匿托士}和\Person{梅里托士}也将再次作为\Person{苏格拉底}的控告者出现，\Place{亚略巴古}{Areopagus}议事会将判处他死刑！但更可笑的是，\Person{苏格拉底}将穿上与他先前周期所穿完全相同的衣服，生活在同样不变的贫穷状态中，以及同样不变的\Place{雅典}城！\Person{法拉里斯}将再次施行暴政，他的铜牛将从受害者内部的惨叫声中再次发出吼声，与先前周期中被判刑者的叫声毫无二致！\Person{费雷的亚历山大}也将再次以与先前同样不变的残暴施行暴政，并判处与先前相同的\Quote{\OriginalTerm{不变}{unchanged}}个人死刑。然而，有什么必要详细讨论\OriginalTerm{斯多葛派}{Stoic}哲学家关于此类事物的学说呢？这些学说逃避了\Person{塞尔斯}的嘲笑，甚至也许被他尊崇，因为他认为\Person{芝诺}比\Person{耶稣}更有智慧？\par

% structure: p0042 source=h2 target=subsection reason=relative-heading-rank; base=h2

\subsection{第二十一章}

\Person{毕达哥拉斯}和\Person{柏拉图}的门徒，尽管他们似乎主张世界的不朽，却也陷入了类似的错误。因为如同行星在确定的周期之后回到相同的位置并彼此保持相同的关系，他们断言，地上一切事物将像在行星关系处于世界同一状态时那样。从这一观点必然得出，当经过悠长的周期后，行星彼此占据与\Person{苏格拉底}时代相同的关系时，\Person{苏格拉底}将再次由同一父母出生，并遭受同样的待遇，被\Person{安匿托士}和\Person{梅里托士}控告，并被\Place{亚略巴古}{Areopagus}议事会定罪！此外，\Person{埃及人}中的学者也持有类似观点，然而他们受到尊重，并未招致\Person{塞尔斯}及其同类之人的嘲笑；而我们，主张万有由\OriginalTerm{天主}{God}按每人自由意志的比例治理，并不断被带入更好的状态——只要它们容许如此——并且知道我们自由意志的本性容许偶然事件的发生（因为它不能接受\OriginalTerm{天主}{God}完全不变的本性），却似乎没有说出任何经得起检验的话。\par

% structure: p0044 source=h2 target=subsection reason=relative-heading-rank; base=h2

\subsection{第二十二章}

然而，不要有人怀疑，我们这样说时，属于那些确实被称为基督徒，却摒弃了\WorkTitle{圣经}所教导的复活教义的人。因为按照他们的原则，这些人完全无法说明从麦粒或其他（撒在地里的）东西长出的茎或树；而我们相信，那\Quote{播种的}若不先死去，决不会\Quote{活起来}，并且所播种的不是那将来的形体（因为\OriginalTerm{天主}{God}随自己的心意给它一个形体，使它复活起来——播种的是可朽坏的，复活起来的是不可朽坏的；播种的是可羞辱的，复活起来的是光荣的；播种的是软弱的，复活起来的是强健的；播种的是属生灵的身体，复活起来的是属神的身体），——我们既保存了\OriginalTerm{基督}{Christ}的\OriginalTerm{教会}{Church}的教义，也保存了\OriginalTerm{天主}{God}应许的伟大，并以论证而非仅仅断言来证明其成就的可能性；我们知道，即使天地及其中的万物都会过去，但祂关于每一事物的言语——作为整体的部分或属类的种——那起初与\OriginalTerm{天主}{God}同在的\OriginalTerm{天主}{God}\OriginalTerm{圣言}{Word}的言辞，决不会过去。因为我们愿意听从那一位的话，祂说：\BibleQuote{天地将要过去，但我的话决不会过去。}\par

% structure: p0046 source=h2 target=subsection reason=relative-heading-rank; base=h2

\subsection{第二十三章}

因此，我们不主张那已经腐败的身体恢复它原来的性质，正如那已经腐烂的麦粒不会恢复它先前的状态一样。但我们主张，正如在麦粒之上长出麦秆，在身体中有一种能力被植入，这能力不被毁灭，身体由此在不朽中复活。然而，廊下派的哲学家们，由于他们关于事物在某一周期后不变的观点，断言身体在完全腐败之后将恢复它原来的状态，并将再次采取那最初的、从中进入分解状态的本性，他们自以为用无可辩驳的论证确立了这些观点。但我们并不逃到最荒谬的托词中，说在天主一切都是可能的；因为我们知道如何理解这\Quote{一切}一词，它既不指那些\Quote{不存在}的事物，也不指那些不可理解的事物。但我们同时主张，天主不能做可耻的事，因为那样他就可能不再是天主；如果他做任何可耻的事，他就不是天主。然而，既然他提出一个原则，说\Quote{天主不愿违反本性的事}，我们就必须加以区分，并且说：如果任何人断言邪恶是违反本性的，而我们主张\Quote{天主不愿违反本性的事}——无论是出于邪恶还是出于非理性的原则——那么，如果这些事是按照天主的话语和旨意发生的，我们立刻就必须认为它们不是违反本性的。因此，天主所做的事，虽然对某些人来说可能是或\Emph{显得}难以置信，但并不违反本性。如果我们必须强调词语的力量，我们可以说，与一般所理解的\Quote{本性}相比，有些事是\Emph{超出}其能力的，天主随时可以做到；例如，将人提升到人类本性之上，使他进入更好、更神圣的状态，并使他保持在这种状态中，只要那被眷顾的对象通过他的行为表明他渴望（持续得到帮助）。\par

% structure: p0048 source=h2 target=subsection reason=relative-heading-rank; base=h2

\subsection{第二十四章}

此外，正如我们已说过，天主愿任何有损于他自身的事会破坏他作为天主的存在，我们将补充说：如果人，依照其本性的邪恶，愿任何可憎的事，天主不能允准。现在，我们回答塞耳苏斯的论断并非出于争辩的精神；而是本着真理的精神去研究它们，因为我们同意他的观点，即\Quote{他是天主，不是无节制的欲望、也不是错误和混乱的天主，而是正直和公义本性的天主}，因为他是众善之源。并且我们承认他能赐予灵魂永恒的生命；他不仅拥有\Quote{能力}，也拥有\Quote{旨意}。因此，鉴于这些考虑，我们一点也不为赫拉克利特被塞耳苏斯采纳的论断而困扰，即\Quote{尸体应当被抛弃，比粪土更无价值}；然而，即使关于这一点，也可以说，粪土的确应当被抛弃，而人的尸体，由于居住其中的灵魂，尤其是如果这灵魂曾是德性的，就不应当被抛弃。因为，与那些基于公平原则的法律相一致，尸体被认为配得埋葬，并给予通常的荣誉，以便尽可能不去侮辱那曾居住其中的灵魂，而将尸体（在灵魂离去后）像畜生的尸体一样抛弃。那么，不要不合理地认为，天主愿意宣告麦粒不是不朽的，而是从它长出的麦秆；那在朽坏中播种的身体不是不朽的，而是由他在不朽中复活的身体。但根据塞耳苏斯，天主自己就是万有的理性，而根据我们的观点，是他的圣子，我们用哲学的语言说：\BibleQuote{在起初已有圣言，圣言与天主同在，圣言就是天主}；而按照我们的判断，天主也不能做任何违反理性或违反他自己的事。\par

% structure: p0050 source=h2 target=subsection reason=relative-heading-rank; base=h2

\subsection{第二十五章}

现在我们接着注意塞耳苏前面的陈述之后的下述言论：\Quote{犹太人就成为一特殊的民族，并依照他们本国的风俗制定法律，至今仍加以遵守，并保持一种敬礼方式，不论其性质如何，总是从他们祖先传下来的。在这几方面，他们的行为与其他民族相同，因为每个民族都保留其祖传的习俗，不论这些习俗是什么，只要是在他们当中建立起来的。这样的安排似乎是有益的，不仅因为其他民族的思想也曾决定某些不同的做法，而且因为保护那些为公共利益而建立的传统是一种责任；此外，很可能世上的各个区域从起初就被分配给不同的监护神，并就这样分配给了某些治理的权力，而世界上的治理即以这种方式进行。在各民族中如此行的一切，凡是合乎（监护神的）旨意的，就是正确行事的；而摆脱各地从起初建立的制度，则是一种不虔敬的行为。}通过这些言词，塞耳苏表明：原为埃及人的犹太人，后来成为\Quote{特殊的民族}，制定了他们仔细保存的法律。我不重复已经在我们面前出现的他的言论，他说：对犹太人来说，遵守他们祖先的敬礼是有益的，正如其他民族细心遵守他们的敬礼一样。他还进一步提出一个更深入的理由，说明为何犹太人培养他们祖先的习俗是有益的，他隐约暗示：那些被指派监督正在立法的地区的（神灵），与各地的立法者合作，制定了各地的法律。因此，他似乎表示：犹太人的地区以及居住在那里的民族，由一个或多个神灵监护，无论是一个或多个，他们与梅瑟合作，并制定了犹太人的法律。\par

% structure: p0052 source=h2 target=subsection reason=relative-heading-rank; base=h2

\subsection{第二十六章}

\Quote{我们}，他说，\Quote{必须遵守法律，不仅因为其他人思想曾决定某些不同的做法，而且因为保护那些为公共利益而制定的法律是一种责任；此外，很可能世上的各个区域从起初就被分配给不同的监护神，并分配给了某些治理的权力，而世界上的治理即以这种方式进行。}如此，塞耳苏仿佛忘记了他曾对犹太人说过的话，现在将他们包括在他对一切遵守祖传习俗者的一般赞扬中，评论道：\Quote{在各民族中如此行的一切，凡是合乎（监护者）的旨意的，就是正确行事的。}且注意这里，他是否不公开地、尽他所能地表达一种愿望：犹太人应遵守自己的法律，而不背离它们，因为如果他背教，就会犯不虔敬的罪；因为他说：\Quote{摆脱各地从起初建立的制度，是一种不虔敬的行为。}现在，我愿问他以及那些持他观点的人：是谁从起初将世上的各个区域分配给不同的监护神？特别是，是谁将犹太人的地区和犹太民族本身，分配给一位或多位监护者，即他们所归属的？是如塞耳苏所说的，朱庇特将犹太民族及其地区指派给某个神灵或多个神灵吗？是\Emph{他的}（他们被分配归其监护的那一位）意愿要在他们当中实行那些盛行的法律，还是\Emph{违反}他的意愿而实行的？你会注意到，无论他回答什么，他都陷入困境。但如果世上的各个区域\Emph{不}是由某一位存在分配给各监护神，那么，每个区域就是随机地、没有更高权力的监护，而根据机遇被划分；然而，这种观点是荒谬的，并在很大程度上破坏了统御万有的天主的眷顾。\par

% structure: p0054 source=h2 target=subsection reason=relative-heading-rank; base=h2

\subsection{第二十七章}

的确，任何人若愿意，都可以讲述大地各方如何被分配给某些掌管权力，由监督者管理；但也请他告诉我们，在各民族中所行的事，当符合监督者的意愿时，如何才算是正当地行的。例如，请他告诉我们，西徐亚人允许弑亲的法律是否是正当的法律；或者波斯人的法律，不禁止儿子与母亲、女儿与父亲结婚的法律，是否是正当的。但我何需从那些在不同民族中制定法律的人中挑选例子，来探究法律如何按照监督的权力而正当地制定？然而，请\Person{瑟尔苏}告诉我们，废除那些允许母亲与女儿结婚，或宣告自缢而终的人有福，或宣称投入火中并以此结束生命的人得到完全净化，这样的祖传法律，如何会是一种不虔的行为？又如何废除那些在陶里刻松盛行的将陌生人祭献给狄安娜的法律，或某些利比亚部落中关于将儿童祭献给萨图尔努斯的法律，是一种不虔的行为？此外，从瑟尔苏的论断可以推出，犹太人废除那些禁止敬拜任何其他神祇、只敬拜万有的创造者的祖传法律，是一种不虔的行为。按照他的观点，虔敬并非因其本性而神圣，而是因某种外在的安排与任命。因为在某些部落中，敬拜鳄鱼是虔敬的行为，并且吃其他部落所敬拜的对象；而在另一些部落中，敬拜牛犊是虔敬的行为，在另一些部落中，则视山羊为神。这样，同一个人会依照一套法律被视为虔敬，依照另一套法律被视为不虔；这是所能想象的最荒谬的结果！\par

% structure: p0056 source=h2 target=subsection reason=relative-heading-rank; base=h2

\subsection{第二十八章}

然而，对上述言论可能作出的回答大概是：那遵守\Emph{自己的}国家法律的人是虔敬的，绝不因不遵守\Emph{其他的}土地的法律而负有不虔之罪；反之，在某些民族中被视为不虔的人，当他依照本国的法律敬拜自己的神时，实际上并非不虔，尽管他对他国法律所尊崇的神明进行战争，甚至以其为食。现在，请看这些说法是否没有表现出对正义、神圣和宗教本质的最大混乱？因为对这些事物没有下明确的定义，也没有描述它们具有自身独特的特征，并把遵从其诫命的人印上虔诚的印记。如果宗教、虔敬和正义属于那些仅因比较而成为如此的事物，以至于同一行为可能因不同的关系和法律而既虔敬又不虔，那么看看这是否也会推出节制、勇敢、明智以及其他德行也属于比较的范畴？那是最荒谬不过的了！然而，我们以上所说的，足以回答瑟尔苏的指控，这是一类更一般和简单的回答。但我们认为，有些习惯于更深探究的人可能会遇到这篇论文，所以我们斗胆提出一些更深刻的思考，传达一种关于大地各方最初分配给不同监督神灵的神秘而隐秘的观点，并尽我们所能证明，我们的教义没有上述的荒谬后果。\par

% structure: p0058 source=h2 target=subsection reason=relative-heading-rank; base=h2

\subsection{第二十九章}

在我看来，事实上，克耳苏斯误解了某些与世上事务安排有关的更深奥理由；这些事情甚至在希腊史中也略有提及，其中某些被视为神明的人物被描绘成曾为争夺阿提刻的占有而彼此争斗；而在希腊诗人的作品中，一些被称为神祇的也被描绘成承认此地某些地方比其它地方更受他们偏爱。此外，蛮族的历史，尤其是埃及的历史，也包含一些类似的对所谓埃及诸省划分的暗示，其中提到通过抽签获得赛伊斯的雅典娜，也是拥有阿提刻的那一位。埃及的学者可以列举无数这类例子，尽管我不知道他们是否将犹太人和他们的国家包括在这一划分中。现在，就这一问题上圣经之外的见证而言，目前引用的已足够。\newline{}此外，我们说，我们的天主先知和祂的真正仆人梅瑟，在《申命纪》的诗歌中，关于土地的分配作了如下陈述：\BibleQuote{当至高者分配万民时，当祂分散亚当的子孙时，祂按天主天使的数目定了万民的疆界；祂的份就是雅各伯，以色列是祂继承的产业。}\newline{}关于万民的分派，同一位梅瑟在其名为《创世纪》的著作中，以历史叙述的风格如此表达：\BibleQuote{当时全世界只有一种语言，说一样的话。当他们从东方迁移时，在史纳尔地方找到了一片平原，就在那里住下了。}稍后他接着说：\BibleQuote{上主下来要看世人所建造的城和塔。上主说：看，他们是一个民族，都说一样的语言。他们已开始这样做；现在他们想做的，没有办不到的。来，我们下去，混乱他们的语言，使他们彼此听不懂对方的话。于是上主将他们从那里分散到全地面上；他们便停止建造那城。因此那城名叫巴贝耳，因为上主在那里混乱了全地的语言，并从那里将他们分散到全地面上。}\newline{}此外，在撒罗满的论著《智慧篇》中，论及语言混乱和土地分配发生的事件时，我们发现关于智慧有如下的记载：\BibleQuote{当万民因邪恶的共谋而混乱时，智慧找到了义人，在天主面前保全了他无罪，并在他对儿子的慈爱中使他坚强。}\BibleRef{智 10:5}\newline{}然而关于这些问题，可以谈论许多，而且是神秘性质的；与此相应的是：\BibleQuote{隐藏君王的秘密是好事}——免得关于灵魂进入身体（但不是从一个身体转入另一个身体）的教义被抛给一般人的理解，也不要把圣物给狗，不要把珍珠丢在猪前。因为这样的做法是不虔敬的，相当于背叛天主智慧的神秘宣言，关于这智慧曾说得很好：\BibleQuote{智慧不进入恶意的灵魂，也不住在一个屈服于罪恶的身体内。}\BibleRef{智 1:4}\newline{}然而，以历史叙述的风格来表现那种意在借历史外衣传达隐秘意义的内容，就足够了，让有能力的人自己去探索与主题相关的一切。（那么，这段叙述可以这样理解。）\par

% structure: p0060 source=h2 target=subsection reason=relative-heading-rank; base=h2

\subsection{第三十章}

地上所有人民都应被视为使用同一种神圣的语言，当他们和谐共处时，就被保存使用这种神圣的语言，并且他们一直停留在东方，没有迁移，只要他们被那\Quote{光明}和永恒光明的\Quote{反映}的思想所充满。但当他们离开东方，开始怀有与东方相异的思想时，他们在史纳尔地找到了一处地方（这地方解释为\Quote{咬牙}，象征性地表示他们失去了赖以支持的食粮），并在那里居住下来。然后，他们想要聚集物质的东西，并将与天没有自然联系的东西与天连接起来，以便借物质的东西来反对非物质的东西，他们就说：\BibleQuote{来，我们做砖，用火烧透。}于是，当他们硬化并压缩这些粘土和物质的材料，表明他们想要使砖变成石头，使粘土变成沥青，并借这些材料建造一座城和一座塔，塔顶（至少在他们观念中）要通到天上，按那\BibleQuote{高抬自己反对天主知识的事物}的方式，各人（按他们离开东方的不同程度，以及砖变成石头、粘土变成沥青、并用这些材料进行建造的程度）被交付给性格或严厉或严厉程度不同的天使，直到他们为自己的胆大妄为付出代价；他们被那些天使——天使给每人印上自己的语言——按各人的功过带到地上不同的地方：例如，有的到炎热地区，有的到以寒冷惩罚居民的国度，有的到极其难以耕种的土地，有的到程度稍轻的，第五种被带到充满野兽的地方，第六种到野兽相对较少的国家。\par

% structure: p0062 source=h2 target=subsection reason=relative-heading-rank; base=h2

\subsection{第三十一章}

现在，接下来，如果有人有这个能力，就让他明白：在那些具有历史形式、包含一些字面真实但同时又传达更深含义的事物中，那些保存了自己原始语言的人，由于没有从东方迁移，继续拥有东方及其东方语言。让他注意，只有这些人成了主的产业，祂的子民被称为\Person{雅各伯}，\Person{以色列}成了祂继承的绳索；只有这些人受到一位统治者的治理，这位统治者接收那些置于他之下的人并不是为了惩罚，不像其他统治者那样。让他，也以凡人所能感知的限度，观察到：在那些被指派给主作为祂卓越产业的人的政治体中，所犯的罪最初是那些可被宽恕的，其性质不至于使罪人完全被抛弃；但随后罪越来越多，虽然仍属于可赦免的性质。并注意到这种情况持续了相当长的时间，且总有补救措施，经过一定间隔后，这些人又返回了他们的职责，让他注意到他们按所犯的过犯被交托给了被指派管理其他地域的人；起初受到轻微惩罚并做了补赎后，他们如同经历了管教一般，回到了自己适当的居所。让他也注意到，后来他们被交给了更严厉的统治者——按\WorkTitle{圣经}的说法，就是\Person{亚述}人和\Person{巴比伦}人。接下来，尽管施加了医治的手段，让他观察到他们仍然在增加过犯，因此被压迫他们的民族统治者驱散到其他地区。而他们自己的统治者故意忽视他们受其他民族统治者压迫的事，为的是他也可以有正当理由，作为报复自己，获得能力从其他民族中尽可能地夺走一些人，为他们制定法律，指出一种他们应当遵守的生活方式，从而引导他们达到那先前不曾犯罪的子民所达到的终点。\par

% structure: p0064 source=h2 target=subsection reason=relative-heading-rank; base=h2

\subsection{第三十二章}

借此，让那些有能力理解如此深奥真理的人明白：那位拥有从未犯罪之人份额的统治者远比其他的强大，因为他能够从整个份额中拣选个体，将他们与那些为了惩罚而接收他们的人分开，并将他们置于法律和一种有助于忘却从前过犯的生活方式之下。但是，正如我们先前所说，这些话应当被视为我们以隐晦的方式所说，用以指出那些人的错误，他们断言\Quote{大地各区域从起初就分配给不同的守护神灵，并分派给某些统治权力，以这种方式管理}；刻尔苏斯正是从这一说法中引出前面提到的言论。但是，由于那些从东方出走的人因他们的罪被交付给\BibleQuote{邪僻的心}\BibleRef{罗 1:28}，并\BibleQuote{可憎的情欲}\BibleRef{罗 1:26}，以及\BibleQuote{因心里情欲所导致的不洁}\BibleRef{罗 1:24}，为使他们饱尝罪恶而憎恨它，我们将拒绝刻尔苏斯的断言，即\Quote{因为分布在大地各区域的守护神灵，每个民族所做的事都是正当的}；因为我们渴望做那些不合这些神灵心意的事。因为我们看到，废除各地最初藉由更美好、更神圣的法律所建立的习俗，乃是宗教行为，这些法律是由耶稣制定的，祂拥有最大的权能，救我们脱离了\BibleQuote{这邪恶的当前世界}\BibleRef{迦 1:4}和\Quote{这即将灭亡的世界之首}；而拒绝俯伏在祂脚前，祂显明自己比所有其他统治者更神圣、更有权能，天父在多个世代之前藉先知预言时对他说：\BibleQuote{向我求，我必将外邦人赐给你作基业，将地极赐给你作产业。}\BibleRef{咏 2:8}因为祂也成了我们这些从外邦信靠祂和祂父——万有之上的天主——之人的\Quote{期望}。\par

% structure: p0066 source=h2 target=subsection reason=relative-heading-rank; base=h2

\subsection{第三十三章}

我们刚才所说的话，不仅回答了瑟尔苏关于监护神灵的论述，而且也在某种程度上预见了其后他提出的论点，当时他说：\Quote{让第二方出场；我将问他们从哪里来，他们把谁视为他们祖传习俗的创始人。他们会回答：没有人，因为他们与犹太人自己来自同一源头，并且他们的教导和监督并非来自其他地方，然而他们却背叛了犹太人。} 那么，我们每个人都在\Quote{末世}，当一位耶稣访问我们时，来到了\Quote{可见的主之山}，即超越一切言语的圣言，以及\Quote{天主的家}，即\Quote{生活的天主的教会，真理的柱石和基础}。我们注意到它如何建立在\Quote{众山之巅}上，即所有先知预言的根基之上。这房屋被高举于丘陵之上，即那些在智慧和真理方面自诩更高成就的人。万民都来到它那里，\Quote{万民}出来，彼此说，转向那在末世通过耶稣基督照耀出来的宗教：\Quote{来吧，让我们登上主的山，到雅各伯的天主家里；祂将教导我们祂的道路，我们要行走在其中。} 因为法律从住在熙雍的人中间出来，作为精神的法律定居在我们中间。此外，主的圣言从耶路撒冷出来，以便传播到各地，并在外邦人中施行审判，拣选祂看到顺服的人，拒绝不顺服的人，他们人数众多。对那些问我们从哪里来，或谁是我们的创始人的人，我们回答说，我们按照耶稣的劝告而来，为了\Quote{把我们敌意和傲慢的\InnerQuote{言语}之剑打成犁头，把以前用于战争的长矛改成镰刀。} 因为我们不再拿起\Quote{刀剑攻击国家}，也不\Quote{再学习战事}，成为了和平之子，为了耶稣，我们的领袖，而不是我们祖先所追随的那些，在这些人中我们曾是与盟约无关的外人，并领受了一条法律，为此我们感谢那将我们从（我们道路的）错误中拯救出来的主，说：\Quote{我们的祖先尊敬虚假的偶像，其中没有一位能降雨。} 那么，我们的监督者和教师，从犹太人中间出来，用祂教导的圣言治理全世界。我们通过预先这些评论，已尽我们所能驳斥了瑟尔苏的不实之词，附加了适当的回答。\par

% structure: p0068 source=h2 target=subsection reason=relative-heading-rank; base=h2

\subsection{第三十四章}

但是，为了不忽略瑟尔苏在这一段和前面段落之间所说的话，我们引用他的原话：\Quote{我们可以引用\Person{希罗多德}作为这一点上的见证，因为他这样说道：\InnerQuote{至于\Place{马雷阿}和\Place{阿皮斯}两城的居民，他们居住在\Place{埃及}靠近\Place{利比亚}的地区，并且自认为是\Place{利比亚}人而非埃及人，因厌恶祭献仪式，且不希望被禁止食用牛肉，就派人到\OriginalTerm{朱庇特·阿蒙}{Jupiter Ammon}的神谕处，说他们与埃及人没有关系，他们住在\Place{三角洲}之外，与埃及人没有共同感情，并希望允许他们吃各种食物。但神不允许他们如愿，说那地方是埃及的一部分，受\Place{尼罗河}泛滥灌溉，而居住在\Place{象岛}以南，并饮用\Place{尼罗河}河水的人是埃及人。}这就是\Person{希罗多德}的叙述。但是，} 瑟尔苏接着说，\Quote{\OriginalTerm{阿蒙}{Ammon}在神圣事物上不会比犹太人的天使更差的使者，所以每个民族遵守其既定的崇拜方式并无不对。的确，我们将会发现各国之间存在巨大差异，然而每个民族似乎都认为自己的方式最佳。居住在\Place{麦罗埃}的\Place{埃塞俄比亚}居民只崇拜\OriginalTerm{朱庇特}{Jupiter}和\OriginalTerm{巴克斯}{Bacchus}；\Place{阿拉伯}人只崇拜\OriginalTerm{乌拉尼亚}{Urania}和\OriginalTerm{巴克斯}{Bacchus}；所有埃及人崇拜\OriginalTerm{奥西里斯}{Osiris}和\OriginalTerm{伊西斯}{Isis}；\Place{赛斯}人崇拜\OriginalTerm{弥涅耳瓦}{Minerva}；而\Place{诺克拉提斯}人最近将\OriginalTerm{塞拉皮斯}{Serapis}列入他们的神祇，其他民族则根据各自的法律。有些人不吃羊肉，有些人不吃鳄鱼肉，还有人不吃牛肉，而他们厌恶猪肉。\Person{斯基泰人}确实认为吃人是高尚的行为。在\Person{印度人}中，也有一些人认为吃自己的父亲是神圣的职责，这件事在\Person{希罗多德}的某段话中提到。为了可信，我将再次引用他的原话，他这样写道：\InnerQuote{如果向所有人提出这个建议，即让他们从所有现有法律中选出最好的，每个人经过审查后都会选择自己国家的法律。人们都认为自己的法律最好，因此除了疯子，没有人会嘲笑这些事。但这是所有人在法律上的结论，可以通过许多其他证据来判断，尤其是以下例证：\Person{大流士}在位期间，传唤当时在场的希腊人，问他们愿意出多少钱来吃他们死去的父亲？他们回答，无论什么代价都不会这样做。之后，\Person{大流士}传唤那些称为\Person{卡拉提亚人}的\Person{印度人}，他们有吃父母的习惯，在希腊人面前（通过翻译得知情况）问他们愿意花多少钱来用火焚烧他们死去的父亲？他们大声呼喊，请求国王不要再说了。}这就是对这些事情的看法。\Person{品达}在我看来说得对，他说\InnerQuote{法律}是万物的王。}\par

% structure: p0070 source=h2 target=subsection reason=relative-heading-rank; base=h2

\subsection{第三十五章}

策尔苏斯的论证似乎通过这些例证得出以下结论：\Quote{所有的人都负有义务按照他们本国的习俗生活，这样他们就能避免责备；而基督徒却放弃了他们本土的习俗，他们不像犹太人那样是一个民族，因此他们因坚持耶稣的教导而应受责备。}那么，请告诉我们，对于哲学家以及那些被教导不要屈服于迷信的人来说，放弃他们本国的习俗，去吃那些在他们各自城市中被禁止的食物，这是否合乎体统？或者，他们的这种行为是否与体统相悖？因为，如果出于他们的哲学以及他们所受的反对迷信的教导，他们不顾本国的法律，去吃那些被他们祖先禁止的东西，那么基督徒（因为福音要求\Emph{他们}不要忙于雕像和偶像，甚至不要忙于天主所创造的任何受造物，而要上升到高处，将灵魂呈献给造物主。）当他们做出与哲学家类似的行为时，为什么要为此受责备呢？但是，如果为了维护他为自己提出的论点，策尔苏斯或与他想法一致的人会说，即使是研究过哲学的人也会遵守本国的法律，那么例如埃及的哲学家，为了避免吃洋葱以遵守本国的法律，或者避免身体某些部位如头和肩膀，以免违反祖先的传统，就会做出极其可笑的事情。我且不提那些因身体排气而战栗的埃及人，因为如果他们中的某个人成为哲学家，并且仍然遵守本国的法律，他会成为一个可笑的哲学家，行为非常不合哲学。同样地，一个人被福音引导去敬拜万有的天主，却由于顾及本国的法律，滞留在人间，流连于人的雕像和偶像，不愿上升到造物主那里，他就好比那些确实学习了哲学，却害怕那些本不应引起恐惧的事物，并且认为吃那些被列举出来的东西是不敬虔的人。\par

% structure: p0072 source=h2 target=subsection reason=relative-heading-rank; base=h2

\subsection{第三十六章}

但这个希罗多德所说的阿蒙是什么样的存在呢？策尔苏斯引用了他的话，似乎是要证明每个人都应遵守本国的法律。因为这位阿蒙不允许马雷阿和阿皮斯两座城市（它们位于利比亚附近区域）的居民把吃牛肉当作无所谓的事，而牛肉本身不仅是无关紧要的，而且不妨碍一个人品德高尚。如果阿蒙禁止吃牛肉，是因为这种牲畜在耕作土地方面有用，另外还因为雌性能够繁殖，那么这种说法就更有可信度了。但现在他只是简单要求那些饮用尼罗河水的人遵守埃及人关于牲畜的法律。于是，策尔苏斯借机嘲笑犹太人中的天使作为天主的使者，他说：\Quote{阿蒙作为神圣事物的使者，并不比犹太人的天使更差。}但他并没有研究这些天使的话语和显现的含义；否则他会看出，天主所关心的并不是牛，即使祂似乎为牛立法，或者为无理性的动物立法，而是在以无理性动物为表象所写的、为人类所写的内容中包含了某些自然的真理。此外，策尔苏斯说，任何希望遵守本国法律所认可的宗教崇拜的人，都没有犯下错误；按照他的观点，所以斯基泰人在遵守本国法律时吃人，也没有犯错误。而那些吃自己父亲的印度人，根据策尔苏斯的看法，是做了一件虔诚的，或者至少不是邪恶的事情。他确实引用了希罗多德的一段话，支持每个人应当出于体统而服从本国法律的原则；他似乎赞同那些被称为卡利亚提亚人的印度人的习俗，他们在达里乌斯时代吃掉自己的父母，因为当达里乌斯问他们愿意用多少钱来放弃这一习俗时，他们大声喊叫，请国王不要再说了。\par

% structure: p0074 source=h2 target=subsection reason=relative-heading-rank; base=h2

\subsection{第三十七章}

如此，通常有两个法律呈现给我们：一个是自然法律，\OriginalTerm{天主}{God}是其立法者；另一个是城邦的成文法。当成文法不与天主的法律相悖时，公民不当以异邦习俗为借口而离弃它；但当自然法律，即天主的法律，命令与成文法相反之事时，试看理性岂不告诉我们，要长久告别那成文法典及其立法者的意愿，将我们自己交付给立法者天主，并选择合乎其圣言的生活，尽管这样做时可能要遭遇危险、无数辛劳，甚至死亡与耻辱？因为当有些法律与天主的旨意和谐，却与城邦中生效的其他法律相抵触时，若无法取悦天主（以及那些管理此类法律的人），那么轻视那些能取悦万有创造主的行为，反而选择那些会令我们不悦天主的行为，尽管我们可能满足不洁的法律及其喜爱者，这岂不是荒谬？但既然在其它事情上，将自然法律——即天主的法律——置于成文法之上是合理的，而成文法是由人怀着与天主法律相悖的精神制定的，那么对于关乎天主的法律，我们岂不更应如此行？我们既不会像居住在\OriginalTerm{梅罗埃}{Meroe}附近的\OriginalTerm{埃塞俄比亚人}{Ethiopians}那样，随意只敬拜\OriginalTerm{朱庇特}{Jupiter}和\OriginalTerm{巴克斯}{Bacchus}；也不会以埃塞俄比亚人的方式尊崇埃塞俄比亚的神祇；也不会像\OriginalTerm{阿拉伯人}{Arabians}那样，只将\OriginalTerm{乌拉尼亚}{Urania}和\OriginalTerm{巴克斯}{Bacchus}视为神明；更不会崇奉那些以性别差异为区分的神祇（如阿拉伯人之敬拜乌拉尼亚为女神，巴克斯为男神）；也不会像所有\OriginalTerm{埃及人}{Egyptians}一样，视\OriginalTerm{奥西里斯}{Osiris}和\OriginalTerm{伊西斯}{Isis}为神；也不会像\OriginalTerm{赛斯人}{Saïtes}那样，将\OriginalTerm{雅典娜}{Athena}列于其中。即便对\OriginalTerm{瑙克拉提斯}{Naucratis}的古代居民而言，敬拜其他神祇是适宜的，而他们的现代后裔最近才开始尊崇\OriginalTerm{塞拉皮斯}{Serapis}（他从来就不是神），但我们不会因此断言，一位先前不是神、也完全不为人所知的新存有是神祇。因为天主子，\BibleQuote{一切受造物的首生者}，虽似乎最近才降生成人，但绝非因此而成为新近者。因为圣经知晓祂是万物创造中最古老的；因为关于人的创造，天主曾对祂说：\BibleQuote{让我们按我们的肖像，照我们的模样造人}。\par

% structure: p0076 source=h2 target=subsection reason=relative-heading-rank; base=h2

\subsection{第三十八章}

但我想指出，\OriginalTerm{凯尔苏斯}{Celsus}毫无根据地断言，每个人都敬拜其本族和本国的制度。因为他宣称：\Quote{那些居住在梅罗埃的埃塞俄比亚人只知道两位神祇，即朱庇特和巴克斯，并且只敬拜这两位；阿拉伯人也只知道两位，即巴克斯（他也是一位埃塞俄比亚神祇）和乌拉尼亚（对她的敬拜仅限于阿拉伯人）。}按照他的说法，埃塞俄比亚人不敬拜乌拉尼亚，阿拉伯人不敬拜朱庇特。那么，如果一个埃塞俄比亚人因故落入阿拉伯人手中，并且因不敬拜乌拉尼亚而被判为不敬之罪，因此面临死亡的危险，那么这位埃塞俄比亚人究竟是应该死去，还是违反其本国法律，向乌拉尼亚行礼？现在，如果他认为违反本国法律是恰当的，那么按照凯尔苏斯的话为标准，他就在做不当之事；而如果他选择走向死亡，就让他表明选择这种命运是合理的。我不知道，如果埃塞俄比亚的教义教导人思考灵魂不死以及宗教所受的尊荣，他们是否会敬奉那些被本国法律视为神祇者。类似的例证也可用于阿拉伯人，如果他们因故恰好拜访梅罗埃附近的埃塞俄比亚人。因为他们受教只敬拜乌拉尼亚和巴克斯，他们不会与埃塞俄比亚人一同敬拜朱庇特；如果因不敬之罪被判处死，就请凯尔苏斯告诉我们，他们做什么才是合理的。至于关于奥西里斯和伊西斯的传说，现在逐一列举是多余而不合时宜的。因为即使可以给这些传说赋予寓意，它们仍然教导我们向冷水、以及向服从于人的大地和一切受造动物献上神圣的敬拜。因为我认为，他们正是这样将奥西里斯关联于水，将伊西斯关联于地；至于塞拉皮斯，则有许多相互矛盾的记载，声称他是最近才公开出现，依据某些由\OriginalTerm{托勒密}{Ptolemy}授意表演的把戏，托勒密希望能向\OriginalTerm{亚历山大里亚}{Alexandria}人民展示一位可见的神。我们在毕达哥拉斯派\OriginalTerm{努美纽斯}{Numenius}的著作中读到关于他的造像：他分有受自然控制的一切动植物的本质，这样他看起来像是被塑造成了神，不仅借助制造偶像者的手艺、以及用于召唤魔鬼的凡俗秘仪和把戏，而且经由魔术师、巫师以及被他们咒语迷惑的魔鬼。\par

% structure: p0078 source=h2 target=subsection reason=relative-heading-rank; base=h2

\subsection{第三十九章}

因此，我们必须探究，对于理性的、温和的、始终按理性行动的动物来说，什么适合吃，什么不适合吃；而不是随意崇拜绵羊、山羊或牛；不食用这些动物是一种节制行为，因为人类从这些动物获得了许多益处。然而，难道还有比饶恕鳄鱼，并将其视为某些神话神明的圣物更愚蠢的事吗？因为饶恕那些不饶恕我们的动物，并照料那些捕食人类的动物，是极其愚蠢的标志。但是，\Person{凯尔苏斯}赞同那些按照本国法律崇拜和照料鳄鱼的人，他对此只字未提；而基督徒却似乎应受谴责，他们被教导憎恶邪恶，远离恶行，并敬畏和尊崇美德——认为美德由天主所生，是祂的圣子。因为，我们不应因其名称和本性的阴性而视智慧和正义为女性；因为在我们看来，这些是圣子，如祂的真门徒所示，论及祂时说：\BibleQuote{祂由天主成为我们的智慧、正义、圣化和救赎。}虽然我们可称祂为\OriginalTerm{第二位天主}{second God}，但让人们知道，我们用\OriginalTerm{第二位天主}{second God}一词无非是指一种能够包含所有其他德行的德能，以及一种能够包含一切事物中存在的所有理性的理性——这些事物是自然地、直接地、为普遍利益而生成的——并且，我们说，这\Quote{理性}曾居住在\Person{耶稣}的灵魂中，并与祂结合，其程度远超所有其他灵魂，因为只有祂得以完全领受绝对理性、绝对智慧和绝对正义的最高分。\par

% structure: p0080 source=h2 target=subsection reason=relative-heading-rank; base=h2

\subsection{第四十章}

但是，既然\Person{凯尔苏斯}在讲述了上述关于不同种类法律的话之后，又补充了以下评论：\Quote{\Person{品达}似乎对我说的是正确的，即法律是万物的君主}，让我们来讨论这个断言。您所说的，好先生，什么是\Quote{万物的君主}的法律？如果您指的是存在于各个城市中的法律，那么这样的断言是不正确的。因为并非所有人都受同一法律管辖。您本应说\Quote{法律是所有人的君主}，因为在每个民族中，某种法律是所有人的君主。但是，如果您指的是本来意义上的法律，那么它正是自然中\Quote{万物的君主}；尽管有一些人，像强盗一样抛弃了法律，否认其有效性，过着暴力和不义的生活。我们基督徒，既已认识到自然中\Quote{万物的君主}的法律，它等同于天主的法律，就努力按其诫命规范我们的生活，永远告别那类不敬的法律。\par

% structure: p0082 source=h2 target=subsection reason=relative-heading-rank; base=h2

\subsection{第四十一章}

让我们来注意\Person{凯尔苏斯}接下来提出的指控，其中与基督徒有关的极少，大部分与犹太人有关。他的话是：\Quote{如果，在这些方面，犹太人谨慎地遵守他们自己的法律，他们这样做不应受责备，而应受责备的是那些放弃了他们自己的习俗而采纳了犹太人的习俗的人。而且，如果他们以此自豪，认为自己拥有更高的智慧，并远离与别人的交往，好像自己不如他们纯洁，那么他们已经听说了他们关于天界的教义并非他们独有，而是——暂且不提其他——早已被波斯人接受，正如\Person{希罗多德}在某处提到的。\InnerQuote{因为他们有一个习俗，他说，登上山顶，向宙斯献祭，将整个天界称为宙斯。}而且，我认为，无论你把至高神称为\OriginalTerm{宙斯}{Zeus}，或\OriginalTerm{泽因}{Zen}，或\OriginalTerm{阿多纳伊}{Adonai}，或\OriginalTerm{萨巴奥特}{Sabaoth}，或像埃及人那样称为\OriginalTerm{阿蒙}{Ammoun}，或像斯基泰人那样称为\OriginalTerm{帕派乌斯}{Pappæus}，都没有区别。他们也不会因守割礼仪式而比别人更圣洁，因为埃及人和科尔基斯人在他们之前就实行了割礼；也不会因不吃猪肉，因为埃及人不仅禁食猪肉，还禁食山羊肉、绵羊肉、牛肉和鱼肉；而毕达哥拉斯及其门徒不吃豆子，也不吃任何有生命的东西。然而，他们不太可能蒙受天主恩宠，或被天主以不同于他人的方式所爱，或者天使从天上唯独被派到他们那里，好像他们被分得了\InnerQuote{某片有福之地}，因为我们看到了他们本人及他们认为配得的国家。那么，让这一帮人走吧，在为其夸耀付出代价之后，他们并不认识伟大的天主，而是被\Person{梅瑟}的诡计引入歧途而受骗，成了他的门徒却毫无益处。}\par

% structure: p0084 source=h2 target=subsection reason=relative-heading-rank; base=h2

\subsection{第四十二章}

显然，借由前面的论述，刻尔苏斯指责犹太人假称自己是至高的天主从万民中特选的一份。他又谴责他们夸口，因为他们自称认识伟大的天主，其实并不认识他，反被梅瑟的诡计所迷惑，上了他的当，成了他的门徒，却毫无益处。我们在前文已经部分论述了犹太人可敬而卓越的政制，当时它存在于他们中间，作为天主之城、他的圣殿以及在其中并在祭台上举行的祭献的象征。但如果有人留意立法者的意图和他建立的体制，研究与他相关的各项细节，并与今日异教民族的敬拜方式相比较，他将发现没有比之更令人钦佩的；因为，凡能在有死之人中实现的，一切对人类无益的事物都被禁止，只保留那些有益的东西。为此，他们没有竞技比赛，没有戏剧表演，没有赛马；也没有妇女向任何想要无生育的性交并蔑视人类生育本性的人出卖美色。从小被教导超越一切可见的本性，并相信天主并不固定于其中任何地方，而是仰望高处，超越一切有形实体的境界，这是何等的益处！几乎从出生起，刚刚能说话，就受到关于灵魂不朽、地下法庭的存在以及为义人所预备的赏报的教育，这又是何等巨大的益处！这些真理确实是披着寓言的外衣讲给儿童和持有幼稚观点之人的；而对于那些已经致力于探求真理并渴望进步的人，这些寓言可以说是变形为隐藏其中的真理。我认为，他们以配得上\Quote{天主的部分}这一称号的方式，蔑视各种占卜，认为那是徒然迷惑人的东西，且多来自恶神而非任何更好的本性；他们从那些因纯洁程度高而领受了至高天主之神的人的灵魂中寻求对未来的知识。\par

% structure: p0086 source=h2 target=subsection reason=relative-heading-rank; base=h2

\subsection{第四十三章}

然而，何须指出，那规定同一信仰者不得被奴役超过六年的法律，是多么合乎健全理性，且不伤害主人或奴隶？那么，不能说犹太人在这些方面与其他民族遵守同样的法律。因为如果他们相信自己的立法与异教民族相同，那就应受责备，并会被指控对自己法律的优越性麻木不仁。尽管刻尔苏斯不承认，但犹太人确实拥有一种智慧，不仅超过大众，也超过那些貌似哲学家的人；因为从事哲学研究的人在发表了最可敬的哲学言论之后，却堕入偶像和魔鬼的崇拜，而最卑微的犹太人却只仰望至高的天主；就这一点而言，他们确实应当以此自豪，并与他人保持距离，视之为可憎和不敬虔的。但愿他们从前没有犯罪，没有违犯法律，没有杀害先知，并且在近些日子没有密谋反对耶稣，这样我们就能拥有一个甚至柏拉图也试图描述的天上之城的典范；尽管我怀疑他能否像梅瑟和追随他的人那样成就那么多，他们以毫无迷信的言语养育了\Quote{特选的种族}和\Quote{圣洁的国民}，奉献给天主。\par

% structure: p0088 source=h2 target=subsection reason=relative-heading-rank; base=h2

\subsection{第四十四章}

因此，他认为关于\Quote{天}的教义与关于\Quote{天主}的教义没有区别；并且他说：\Quote{波斯人像犹太人一样，在山顶上向朱庇特献祭。}——却未注意到，正如犹太人认识一位天主，他们也只有一座神圣的祈祷所、一个全燔祭的祭台、一个香炉和一位天主的大司祭。因此，犹太人与波斯人毫无共同之处，后者攀登他们众多的山顶，献上与梅瑟律法所规定的毫无相似的祭品——犹太司祭正是按照这律法\BibleQuote{侍奉于天上之物的模型和影像}，以谜语般的方式解释关于祭献的法律的目的，以及这些祭献所象征的事物。因此，波斯人可以称\Quote{整个天穹}为朱庇特；但我们坚持\Quote{天}既不是朱庇特，也不是天主，因为我们确实知道，某些低于天主的存有上升到了诸天和一切可见本性之上：正是在这个意义上，我们理解这些话：\BibleQuote{赞美天主，诸天诸天，以及天上的水：愿它们赞美上主的名。}\par

% structure: p0090 source=h2 target=subsection reason=relative-heading-rank; base=h2

\subsection{第四十五章}

然而，既然\Person{塞耳苏斯}认为，至高的存在被称为\OriginalTerm{朱庇特}{Jupiter}、或\OriginalTerm{宙斯}{Zen}、或\OriginalTerm{阿多奈}{Adonai}、或\OriginalTerm{撒巴奥特}{Sabaoth}、或\OriginalTerm{阿蒙}{Ammoun}（如埃及人所称）、或\OriginalTerm{帕佩乌斯}{Pappæus}（如西徐亚人所称）都无关紧要，那就让我们稍加讨论，同时提醒读者，先前当\Person{塞耳苏斯}的言论引导我们思考这问题时已经说过的话。现在我们主张，名称的本性并不像\Person{亚里斯多德}所设想的那样，是命名者的规定。因为流行于人类中的语言并非源于人，这从那些能够查明咒语本性的人身上可以明显看出——这些咒语被语言的发明者根据不同语言以及名称的不同发音而不同地运用，关于这一点我们在前文已经简略谈及，并指出，当那些在某种语言中具有自然能力的名称被翻译成另一种语言时，它们不再能完成在母语中发出时所能成就的事。同样的特性也适用于人；因为如果我们把一个人从出生起在希腊语中被称为某个称呼的名字翻译成埃及语或拉丁语，或任何其他语言，我们无法使他做出或经历他在原先赋予他的称呼下本会做出或经历的那些事。甚至，即使我们把一个原先用拉丁语称呼的人的名字翻译成希腊语，我们也不能产生那咒语若保留最初赋予他的名字时自称能达成的效果。如果这些话对\Emph{人}之名是真实的，那么对于那些因任何原因而转用于\Emph{神明}之名，我们又当如何想？例如，从\Person{亚巴郎}之名被译成希腊文时有所转移，\Person{依撒格}之名和\Person{雅各伯}之名也有所意指；因此，如果有人无论是呼求还是起誓时，使用\Quote{\Person{亚巴郎}的天主}和\Quote{\Person{依撒格}的天主}和\Quote{\Person{雅各伯}的天主}这些表达，他就会产生某些效果，或由于这些名称的本性，或由于它们的能力，因为连恶魔也会被念这些名的人击败并顺服；而如果我们说\Quote{那回声之被选之父的神，那欢笑之神，那用脚跟踢人之神}，则提这名不会带来任何结果，如同其他没有能力的名称一样。同样，如果我们把\Quote{以色列}这个词译成希腊文或任何其他语言，我们不会产生任何结果；但如果我们保留它原样，并将它与那些精通此道者认为应当结合的表达连用，那么这个词的发音就会产生某种结果，与使用这类呼求者所自称的相符。关于\Quote{撒巴奥特}的发音，我们也同样可以说，这个词常被用于咒语；因为如果我们把这个术语翻译成\Quote{万军之主}或\Quote{万军的君主}或\Quote{全能者}（译者提出了不同的理解），我们将一无所成；而如果我们保留原发音，那么，正如精通此道者所维护的，我们就会产生某种效果。同样的观察也适用于\OriginalTerm{阿多奈}{Adonai}。那么，如果\Quote{撒巴奥特}和\Quote{阿多奈}即便被译成它们在希腊语中似乎对应的意思，也不能成就任何事，那么在那些认为至高的存在被称为\OriginalTerm{朱庇特}{Jupiter}、或\OriginalTerm{宙斯}{Zen}、或\OriginalTerm{阿多奈}{Adonai}、或\OriginalTerm{撒巴奥特}{Sabaoth}都无关紧要的人中，其结果更会是多么微乎其微！\par

% structure: p0092 source=h2 target=subsection reason=relative-heading-rank; base=h2

\subsection{第四十六章}

为了这些和类似的奥秘理由，\Person{梅瑟}和先知们是知道的，他们禁止那想要单独向唯一至高的天主祈祷的人说出其他神的名，也不可让那被教导要纯洁于一切愚昧思想和言语的心记住这些名。因此，我们宁愿忍受一切痛苦，也不肯承认\OriginalTerm{朱庇特}{Jupiter}是天主。因为我们不认为\OriginalTerm{朱庇特}{Jupiter}和\OriginalTerm{撒巴奥特}{Sabaoth}是相同的，也不认为朱庇特有任何神性，而是认为某个对人类和真天主不友善的恶魔，以这个称号为乐。而且即使埃及人以死亡威胁要我们承认\OriginalTerm{阿蒙}{Ammon}，我们宁可死也不称他为神，因为这名字极有可能用于某些埃及咒语，以召唤这恶魔。而且即使西徐亚人可能称\OriginalTerm{帕佩乌斯}{Pappæus}为至高的神，但我们不会对此表示同意；我们确实承认\Emph{有}一位至高的神性，但我们不以帕佩乌斯这名作为他的恰当称号，而是认为这称号取悦于那被分配了西徐亚旷野及其人民和语言的恶魔。然而，任何人若以西徐亚语、或埃及语、或他所成长的任何语言来称呼天主的称号，并不算犯罪。\par

% structure: p0094 source=h2 target=subsection reason=relative-heading-rank; base=h2

\subsection{第四十七章}

如今，犹太人施行割损礼的原因，与埃及人和\Place{科尔基斯}人中存在的割损礼原因不同，因此不应被视为相同的割损礼。正如献祭者并非向同一神明献祭，尽管他以相似的方式执行献祭仪式；祈祷者并非向同一神明祈祷，尽管他在祈求中求同样的东西；同样，如果有人施行割损礼，这绝不意味着这与另一人施行的割损礼是同一行为。因为执行此仪式的目的、法律和愿望，将行为置于不同的范畴。但为了让整件事更加清晰，我们需指出：在所有希腊人中，\Quote{正义}一词是相同的；但按照\Person{伊壁鸠鲁}的观点，正义是一回事；按照否认灵魂三分法的斯多葛派，又是另一回事；按照\Person{柏拉图}派，正义又是另一回事，他们认为正义是灵魂各部分的恰当职能。同样，\Person{伊壁鸠鲁}的\Quote{勇敢}是一回事，他们会忍受一些劳苦以逃避更多的劳苦；斯多葛派的\Quote{勇敢}是另一回事，他们为美德本身而选择一切美德；\Person{柏拉图}的\Quote{勇敢}又是另一回事，他认为美德本身就是灵魂中激情部分的行为，并将其置于胸部附近。因此，割损礼将根据施行者的不同观点而有所不同。但关于这个主题，在这样一篇论文中，此时无需多谈；因为任何想知道我们为何进入这主题的人，可以阅读我们在\BibleBook{罗马}{书}中所说的。\par

% structure: p0096 source=h2 target=subsection reason=relative-heading-rank; base=h2

\subsection{第48章}

然而，犹太人虽以割损礼为荣，但他们不仅将之与\Place{科尔基斯}人和\Place{埃及}人的割损礼分开，也与阿拉伯的\Person{依市玛耳}人的割损礼分开；而后者是从他们的祖先\Person{亚巴辣罕}（\Person{依市玛耳}的父亲）传下来的，他是与父亲一起接受了割损礼的。犹太人说，在第八天施行的割损礼是主要的割损礼，而根据情况施行的则是不同的；而且，这可能是由于某个天使对犹太民族的敌意，该天使有能力伤害未受割损的人，但对已经受过割损的人则无能为力。这似乎可以从\BibleBook{出}{谷纪}的记载看出来：在\Person{厄里厄则尔}受割损之前，那个天使能够攻击\Person{梅瑟}，但在他儿子受割损之后，便无能为力了。当\Person{漆颇辣}得知此事后，她拿了一块小石子，为她的孩子行了割损，并据通用版本的阅读，她说：\BibleQuote{我孩子的割损之血停止了，}\BibleRef{出 4:25-26}但根据希伯来文本，她说：\BibleQuote{你真是我的血郎。}\BibleRef{出 4:25-26}因为她知道那个天使在流血之前有力量，但在割损之血后变得无力。因此，对\Person{梅瑟}说的话是：\BibleQuote{你真是我的血郎。}\BibleRef{出 4:25-26}然而，这些看似有点好奇、不为大众理解的事情，我冒昧地如此详细地讨论；现在，作为一个基督徒，我再补充一点，然后继续下面的内容。因为我认为，这个天使对民众中未受割损的人，以及所有只敬拜造物主的人，可能具有力量；这种力量一直持续到\Person{耶稣}取了人性身体为止。但当\Person{耶稣}取了人性，并亲自接受了割损礼之后，这个天使对所有实行同样敬拜但未受割损之人的力量便消失了；因为\Person{耶稣}以其不可言喻的神性将它化为乌有。因此，他的门徒被禁止自行割损，并被（宗徒）提醒：\BibleQuote{如果你受割损，基督就与你无益了。}\BibleRef{迦 5:2}\par

% structure: p0098 source=h2 target=subsection reason=relative-heading-rank; base=h2

\subsection{第49章}

但犹太人也不以戒食猪肉为荣，好像是什么大事；而是因为他们已经认识到了洁净与不洁动物的本质，以及区别的原因，以及猪被列在不洁之中的原因。这些区别在\Person{耶稣}来临之前是某些事物的预表；在他来临之后，他对门徒说（门徒当时尚未理解有关这些事的道理，但说：\BibleQuote{凡俗和不洁的东西，从来没有进过我的口}\BibleRef{宗 10:14}）：\BibleQuote{天主所洁净的，你不可称为凡俗。}\BibleRef{宗 10:15}因此，无论是犹太人还是我们，都不受埃及司祭不仅戒食猪肉，还戒食山羊肉、绵羊肉、牛肉和鱼类的影响。但既然\BibleQuote{入口的不能污秽人}\BibleRef{玛 15:11}，且\BibleQuote{食物不能使我们中悦天主}\BibleRef{格前 8:8}，我们并不特别重视禁食，也不因贪食而吃。因此，就我们而言，\Person{毕达哥拉斯}的信徒（他们禁食一切有生命之物）可以按他们的意愿行事；只要注意\Person{毕达哥拉斯}派与我们苦修者禁食有生命之物的不同理由即可。因为前者禁食是基于灵魂转世的传说，正如诗人所说：——\par

\Quote{有人举起心爱的儿子，\newline{}祈祷后将他杀害；哦，多么愚蠢！}\par

然而，当我们禁食时，是因为我们\BibleQuote{克制自己的身体，使它服从}\BibleRef{格前 9:27}，并且渴望\BibleQuote{致死属于地上的肢体，即淫乱、污秽、邪情、恶欲}\BibleRef{哥 3:5}，我们尽力\BibleQuote{致死肉体的行为}\BibleRef{罗 8:13}。\par

% structure: p0102 source=h2 target=subsection reason=relative-heading-rank; base=h2

\subsection{第50章}

\Person{策尔苏斯}仍在表达他对犹太人的看法，说：\Quote{他们不大可能蒙天主大宠，或更受祂亲爱，仿佛有天使由祂单独派给他们，如同福地部分给了他们。因为我们能看到这民族本身，以及他们认为配得的土地。}我们将反驳此说，指出这民族显然蒙天主大宠，因为统管万有的天主甚至被我们信仰之外的异邦人称为希伯来人的天主。正因为他们蒙天主宠爱，所以没有被祂遗弃；虽然人数稀少，却继续享有天主威能的保护，以致在马其顿的\Person{亚历山大}统治期间，他们未受其伤害，尽管他们因某些盟约和誓言而拒绝起兵对抗\Person{达理阿}。据说在那次，犹太人的大司祭身披圣袍接受了\Person{亚历山大}的敬拜，\Person{亚历山大}宣称他曾看见一个身穿此衣的人在梦中向他宣布他将征服整个\Place{亚细亚}。因此，我们基督徒坚持认为\Quote{那民族极其幸运地享受天主的恩宠，并以一种与众不同的方式蒙祂所爱}；但这一切安排和天主恩宠却转移给了我们，乃因耶稣将曾在犹太人中间彰显的能力传授给那些从外邦归依祂的人。为此，尽管\Place{罗马}人企图对基督徒施加诸多暴行以消灭他们，却未能成功；因为有一只天主之手为他们争战，且旨在于使天主圣言从\Place{犹太}地的一角传遍全人类。\par

% structure: p0104 source=h2 target=subsection reason=relative-heading-rank; base=h2

\subsection{第五十一章}

但是，既然我们已经尽最大努力回答了\Person{策尔苏斯}对犹太人及其教义提出的指控，让我们继续考虑以下内容，并证明当我们宣称认识伟大的天主时，并非空言自夸，我们也没有被\Person{梅瑟}（如\Person{策尔苏斯}所想象）或甚至我们的救主\Person{耶稣}的任何欺骗伎俩引入歧途；而是为了一个良好的目的，我们听从在那\Person{梅瑟}内发言的天主，并接受了\Person{耶稣}——他证明祂是天主——作为天主之子，希望如果我们按照祂的话规范生活，就能获得最好的赏报。我们将甘愿略过那些在我们预料中已经陈述过的内容，即\Quote{我们从何而来，谁是我们领袖，从祂那里颁行了什么法律。}倘若\Person{策尔苏斯}坚持认为我们与埃及人没有区别——他们崇拜山羊、公羊、鳄鱼、牛、河马、狒狒或猫——他可以查证是否如此，其他有类似想法的人也可以。然而，我们已在前面几页中尽最大努力为自己辩护，论述了我们向我们的\Person{耶稣}所表达的敬意，指出我们已找到了更好的部分；并且，在证明\Person{耶稣基督}教导中所包含的真理是纯洁无瑕的，我们不是在赞美自己，而是在赞美我们的导师——至高天主曾通过许多见证人、犹太人的先知著作以及事情本身的清晰性为祂作证，因为祂不可能没有天主的帮助而完成如此伟大的事工。\par

% structure: p0106 source=h2 target=subsection reason=relative-heading-rank; base=h2

\subsection{第五十二章}

但我们现在要审查的\Person{策尔苏斯}的陈述如下：\Quote{那么，让我们略过那些可能针对他们导师的主张而提出的反驳吧，就让他真的被视为一位天使吧。但他是第一位也是唯一一位来到（人间）的吗，还是在他之前还有其他人？如果他们说他只有一位，他们就会被证明在自相矛盾。因为他们断言，在许多场合其他人来过，他们一共六十或七十人，并且这些人都变得邪恶，被扔到地下，被锁链惩罚，而温泉便源于此，那是他们的眼泪；此外，有一位天使曾来到这位所谓的墓前——按照某些人的说法是一位，但按照其他人的说法是两位——他回答妇女们说祂已经复活了。因为天主之子自己似乎不能打开坟墓，却需要另一位帮助才能滚开石头。而且，因\Person{玛利亚}怀孕，一位天使曾来到木匠那里，然后又来了另一位天使，以便他们抱起婴孩逃往\Place{埃及}。但还有什么必要一一列举，或者数清据说被派往\Person{梅瑟}以及他们中间其他人的天使数量呢？如果，那么，其他天使也被派来，显然他也来自同一天主。但他可能被认为具有宣告某种比前者更重大事项的样貌，好像犹太人一直在犯罪，或腐化他们的宗教，或行亵渎之事；因为这些事情都被隐约暗示了。}\par

% structure: p0108 source=h2 target=subsection reason=relative-heading-rank; base=h2

\subsection{第五十三章}

以上论述或许足以答复克尔苏斯的指控，就我们的救主耶稣基督被特别检视的那些问题而言。但为避免显得故意跳过其作品的任何部分，仿佛我们无力应对他，让我们即使冒着重复的风险（因克尔苏斯向我们发起挑战），尽可能以一切应有的简洁继续我们的论述，因为或许在某些主题上我们会有更精确或更新颖的想法。他确实说，\Quote{他省略了针对基督徒为其导师所提出的主张而引用的反驳}，尽管他\Emph{没有}省略任何他能够提出的东西，正如他先前的话所表明的，但这只是空泛的修辞手法。然而，我们并未在关于我们伟大救主的问题上被驳倒，尽管指控者可能\Emph{看似}驳倒我们，这一点对于以热爱真理的精神研读所有关于祂的预言和记载的人来说将是明显的。其次，既然他认为自己在说到救主时让步说\Quote{让他真的被视为一位天使}，我们回答说我们不接受克尔苏斯的这种让步；但我们着眼于祂的工作，祂来是为了以祂的言语和教导访问整个人类，正如祂的每一位追随者所能接受祂的程度。这是那位的工作，按照关于祂的预言所说，祂不仅仅是一位天使，而是\OriginalTerm{伟大的谋略天使}{Angel of the great counsel}：因为祂向人们宣告了天主和万有之父关于他们的伟大谋略，（说）那些献身于虔诚圣洁生活的人，通过他们伟大的事迹上升到天主那里；但那些不依附于祂的人，则使自己远离天主，并因不信祂而走向毁灭。然后他继续道：\Quote{如果天使来到人间，他是第一个也是唯一一个来的吗，还是以前也有别的来过？}他认为自己可以非常冗长地应对这两难困境中的任何一个，尽管没有一个真正的基督徒断言基督是唯一访问过人类的存有。因为，正如克尔苏斯所说，\Quote{如果他们说是唯一的一个}，那么就有别的出现在不同的人面前。\par

% structure: p0110 source=h2 target=subsection reason=relative-heading-rank; base=h2

\subsection{第五十四章}

接着，他以自认为合适的方式回答自己，说：\Quote{因此，他并不是唯一被记载访问过人类的，即使是那些以耶稣之名施教、却背弃了作为较低存有的\OriginalTerm{创造主}{the Creator}、转而依附于一位更高的天主和那位访问者的父的人，也断言在他之前有来自创造主的存有访问过人类。}现在，我们本着真理的精神探究所有与此主题相关的事，我们将指出，马尔西翁著名的门徒\Person{阿培肋斯}（他成为某个异端的创始人，并将犹太人的著作视为神话）断言耶稣是唯一访问过人类的。因此，即使针对他（他主张耶稣是唯一从天主来到人间的那一位），克尔苏斯引用关于其他天使降临的陈述也是徒劳的，因为正如我们已提到的，阿培肋斯不相信犹太圣经的奇迹叙述；他更不会接受克尔苏斯从他不理解的《哈诺客书》内容中引用的东西。因此，没有人能证明我们作假或做出矛盾的断言，仿佛我们既主张我们的救主是唯一来到人间的，却又说其他许多在不同时期来过。此外，在考察天使访问人类的问题时，他还极其混乱地引用了从《哈诺客书》内容中得来的东西，却不明白其含义；因为他似乎并未读过相关段落，也不知道名为《哈诺客》的书并未作为神圣著作在教会中流传，尽管他可能据此获得这个说法，即\Quote{同时有六十或七十位天使降下，陷入了邪恶状态}。\par

% structure: p0112 source=h2 target=subsection reason=relative-heading-rank; base=h2

\subsection{第五十五章}

但为了以坦诚的精神答应他（指\Person{塞尔苏斯}）在\WorkTitle{创世纪}内容中未曾发现的事，即\BibleQuote{天主之子看见人的女儿，见她们美丽，就随意挑选娶来为妻}，我们仍将在这一点上说服那些能理解先知含义的人：在我们之前就有人将这段叙述与关于灵魂的教义联系起来，这些灵魂被对人间肉体生命的欲望所占据，他说，用比喻的说法称之为\Quote{人的女儿}。但无论\Quote{天主之子渴望占有人的女儿}这句话的含义是什么，都丝毫不能证明耶稣不是唯一以天使身份访问人类、并明显成为所有脱离邪恶洪流之人的救主和恩主的那一位。接着，他将自己曾听到或随处看到的任何记载——无论基督徒是否认为其出自神启——混杂并混淆起来，加上说：\Quote{那共同降下的六十或七十位被投到地下，用锁链惩罚。}他又（作为引自\WorkTitle{厄诺克书}，但未提书名）引用以下的话：\Quote{因此，这些天使的眼泪就是温泉。}——这是\Place{天主的众教会}中从未提及或听说的事！因为从来没有如此愚蠢的人，将天使从天上降下所流的眼泪物质化为人类的眼泪。如果可以对\Person{塞尔苏斯}以严肃态度提出的反对意见开个玩笑，我们或许可以指出：没有人会说温泉（其中大部分是淡水）是天使的眼泪，因为眼泪本质是咸的，除非在\Person{塞尔苏斯}看来天使流的是淡水泪。\par

% structure: p0114 source=h2 target=subsection reason=relative-heading-rank; base=h2

\subsection{第五十六章}

紧接着，他将不同且无法统一的事物混杂并互相比较，在他关于那六十或七十位从天上降下、据他说以眼泪流出温泉的天使的陈述之后，附上以下的话：\Quote{又据说，来到\Person{耶稣}坟墓的，有些人说有两位天使，有些人说有一位。}我认为他没有注意到，\Person{玛窦}和\Person{马尔谷}说是一位，而\Person{路加}和\Person{若望}说是两位，这些叙述并不矛盾。因为那些说\Quote{一位}的，是指他（即那位天使）将石头从墓门滚开；而那些说\Quote{两位}的，是指向来到坟墓的妇女们显现身穿发光衣服的，或在坟墓内看见穿着白衣坐着的。这些事件中的每一个现在都可以证明是实际发生的，并且对于这些\Quote{现象}，对那些准备好看见圣言复活的人来说，隐含着一种比喻意义。然而，这样的任务不属于我们当前的目的，而是属于对福音的诠释。\par

% structure: p0116 source=h2 target=subsection reason=relative-heading-rank; base=h2

\subsection{第五十七章}

现在，希腊人曾记载说，有时人类目睹了奇迹般的显现；不仅那些可能被怀疑编造虚构故事的人，而且那些提供了一切证据证明自己是真正哲学家、并如实叙述了他们所遇到之事的人。我们曾在\Person{索利的克吕西波}的作品中读过这类记载，也读过一些关于\Person{毕达哥拉斯}的同类事情；还有近代一些作家（生活在很短时间以前）的作品，如\Person{凯罗涅亚的普鲁塔克}的论著\WorkTitle{论灵魂}，以及\Person{毕达哥拉斯学派努美尼乌斯}的著作\WorkTitle{灵魂的不朽}的第二卷。现在，当希腊人，特别是其中的哲学家叙述这类事情时，不应以嘲笑和讥讽来接受，也不应视为虚构和寓言；但是，当那些献身于万有的天主、忍受各种伤害甚至死亡、而不容许自己的嘴唇说出关于天主的谎言的人，宣告他们亲眼所见的天使显现时，却被认为不值得相信，他们的话不被视为真实！现在，以这种方式来判断一个人是说真话还是假话，是违背健全理性的。因为那些行事诚实的人，只有在长期仔细地考察了主题的细节之后，才会缓慢而谨慎地对这个人或那个人所叙述的奇迹之真实性表达意见；因为事实是，并非所有人都值得相信，也并非所有人都清楚地表明他们只是在向人类讲述虚构和寓言。此外，关于\Person{耶稣}从死者中复活，我们有这样的评论：在如此场合，若有一位或两位天使显现，宣告\Person{耶稣}已从死者中复活，并为那些相信此事而有益于其灵魂的人提供安全，这丝毫不令人惊奇。而且，在我看来，那些相信\Person{耶稣}复活的人，他们以不可轻视的信德果实——拥有美德的生活和脱离邪恶洪流——表现出他们并不缺乏天使的帮助，这些天使帮助他们完成归向天主的转变，这完全是合理的。\par

% structure: p0118 source=h2 target=subsection reason=relative-heading-rank; base=h2

\subsection{第五十八章}

但是，塞耳苏还质疑关于天使滚开那放置耶稣尸体的坟墓门口石头的叙述，他的做法就像学校里一个用一串陈词滥调指控别人的少年。而且，仿佛发现了叙述中某个巧妙的反对点，他评论道：\Quote{如此看来，天主子似乎不能打开自己的坟墓，却需要别人的帮助来滚开石头。}现在，不过度讨论此事，也不显得不合理地引入哲学评论去解释其寓意，我就只就叙述本身说一句：仆人、属下滚开石头，比那位其复活是为了人类益处的主做这事，在本身看来更加恭敬。我不提那些密谋反对圣言、希望置祂于死地并向所有人表明祂确实已死且不存在之人的愿望——他们希望坟墓不被打开，好让无人能在他们密谋之后见到活着的圣言；但是，那位为拯救人类而来到世界的\Quote{天主的使者}，在另一位天使的帮助下，比密谋者更有力，滚开了沉重的石头，为叫那些以为圣言已死的人确信，祂并非与\Quote{逝者}同在，而是活着，并走在前头，引导那些愿意跟随祂的人，以便向他们启示那些在初次进入（基督信仰学校）时，当他们尚不能接受更深教导时，祂先前向他们显示的真理之后的那些事。接下来，我不明白当他嘲笑关于\Quote{天使向若瑟显现告知玛利亚怀孕}的叙述，以及关于天使警告父母\Quote{抱起新生、生命有危险的婴孩，逃往埃及}的叙述时，他认为这对他的目的有何益处。然而，关于这些事，我们在前面的章节中已经回答了他的说法。但是，塞耳苏说\Quote{根据圣经，记载天使被派往梅瑟以及其他一些人}是什么意思？因为在我看来，这对他的目的毫无贡献，尤其因为其中没有任何一位曾尽力，在其能力范围内，使人类从罪中悔改。不过，就算有其他天使从天主那里被派遣，但祂来是宣布某种更重要的事（比任何祂之前的使者都重要）；当犹太人陷于罪恶、败坏了他们的宗教、行了不洁之事时，将天主的国转给了其他园户，这些园户在所有教会中特别关心自己，并竭尽全力借着圣洁的生活和与此相称的教义，将那些想从耶稣的教导中逃开的人赢到万有者天主面前。\par

% structure: p0120 source=h2 target=subsection reason=relative-heading-rank; base=h2

\subsection{第五十九章}

然后塞耳苏继续说道：\Quote{因此，犹太人和这些人（显然指基督徒）有同一位天主；}仿佛在提出一个不会被承认的命题，他接着断言：\Quote{确实，大公教会的成员承认这一点，并接受犹太人当中流传的关于世界创造的叙述，即关于六天和第七天；}在第七天，如圣经所说，天主从祂的工程中\Quote{安息}，退入对祂自己的默观；但在这一天，如塞耳苏所说（他不拘泥于历史的字面，也不理解其意义），天主\Quote{休息了}——这个词语在记载中并未出现。然而，关于世界的创造，以及\Quote{为天主之民所保留的安息}，这个主题是广阔的、奥秘的、深刻的、难以解释的。接下来，在我看来，出于填充他的书并使其看起来重要的欲望，他轻率地添加了一些陈述，例如以下关于第一个人的说法，他说：\Quote{我们给出与犹太人相同的叙述，并从他那推导出相同的族谱。}然而，关于\Quote{兄弟之间的彼此阴谋}，我们不知道有任何这样的阴谋，除了加音对亚伯尔的阴谋，以及厄撒乌对雅各伯的阴谋；但没有亚伯尔对加音的阴谋，也没有雅各伯对厄撒乌的阴谋：因为如果有，塞耳苏说我们给出与犹太人相同的关于\Quote{兄弟之间的彼此阴谋}的叙述才是正确的。不过，就算我们说的关于同一支下埃及和从那里返回（这不是像塞耳苏认为的\Quote{逃跑}）的叙述，这对指控我们或犹太人有什么用处呢？确实，在这里，当他谈论希伯来民族时，他以为可以嘲笑我们，称他们的出离为\Quote{逃跑}；但当他的职责是调查天主对埃及所施惩罚的叙述时，他却故意对此保持沉默。\par

% structure: p0122 source=h2 target=subsection reason=relative-heading-rank; base=h2

\subsection{第六十章}

如果为了精确回应刻尔苏斯（他认为我们对所讨论的问题与犹太人持相同意见），有必要精确表达我们的看法，我们会说：我们都同意（圣经）诸书是由天主圣神所撰写，但我们对其内容的含义\Emph{不}一致；因为我们不像犹太人那样规范我们的生活，因为我们相信法律的文字意义并非传达立法本意的那个意义。我们主张，\BibleQuote{当诵读梅瑟时，帕子还盖在他们的心上}，因为梅瑟律法的含义对那些未曾接纳经由耶稣基督的道路的人是隐藏的。但我们知道，如果一个人转向主（因为\BibleQuote{主就是那圣神}），帕子被除去后，他就\BibleQuote{以没有帕子遮蔽的脸，好像镜子，看见主的荣耀}，即那些隐藏在字面表达中的思想，并为了自己的荣耀，成为天主荣耀的分享者；这里\Quote{脸}一词是比喻用法，指\Quote{理智}，若不用比喻，可以这样说，在它里面是\Quote{内心的人}的面容，充满由对律法内容的真实理解而来的光明和荣耀。\par

% structure: p0124 source=h2 target=subsection reason=relative-heading-rank; base=h2

\subsection{第六十一章}

在上述评论之后，他接着说：\Quote{不要让人以为我不知道，他们中有些人会承认他们的天主与犹太人的天主相同，而另一些人会主张他是另一位，后者与之对立，并且圣子来自前者。}现在，如果他认为基督徒中存在众多异端是对基督宗教的指控，那么，同样，哲学中各种哲学流派在最重要的问题上而非细小无关紧要的问题上彼此分歧，难道不应该成为对哲学的指控吗？甚至医学也因其许多互相冲突的学派而应受到攻击。那么，就让我们承认，我们中间有一些人否认我们的天主与犹太人的天主相同；然而，因此不应责备那些从同一圣经中证明同一天主是犹太人和外邦人共同的天主的人，正如保禄（他从犹太教皈依基督宗教）清楚所说的：\BibleQuote{我感谢我的天主，我用纯洁的良心侍奉他，如同我祖先所侍奉的。}而且，也让我们承认有第三种人，他们称某些人为\Quote{属血肉的}，另一些人为\Quote{属神的}——我想他这里指的是\OriginalTerm{华伦提努}{Valentinus}的追随者——然而这对我们这些属于教会的人有何用处？我们正是以此指控那些认为有些本性得救、有些因其本性构成而丧亡的人。还让我们承认，有些人自称是\OriginalTerm{诺斯底派}{Gnostics}，就像那些自称为哲学家的\OriginalTerm{伊壁鸠鲁学派}{Epicureans}一样；然而，那些取消天主眷顾的教义的人不会被当作真正的哲学家，那些引入怪诞发明的人也不会被当作真正的基督徒，这些发明为耶稣的门徒所不齿。此外，还让我们承认，有些人接受耶稣，并以此为荣地自称基督徒，却像犹太群众一样按照犹太法律规范自己的生活——这是\OriginalTerm{厄彼安派}{Ebionites}的两种形式，他们要么与我们一样承认耶稣是童贞女所生，要么否认这一点，认为他像其他人类一样被生——这对我们这些属于教会的人（刻尔苏斯称我们为\Quote{群众中的那些人}）有何指控价值？他还补充说，某些基督徒相信\OriginalTerm{西比拉}{Sibyl}，大概是他误解了一些责备那些相信存在一位先知西比拉的人，并将这些相信者称为\OriginalTerm{西比拉派}{Sibyllists}。\par

% structure: p0126 source=h2 target=subsection reason=relative-heading-rank; base=h2

\subsection{第六十二章}

他接着向我们倾倒一大堆名字，说他认识一些\OriginalTerm{西满派}{Simonians}崇拜赫肋内或赫肋努斯作为他们的老师，并被称为\OriginalTerm{赫肋尼派}{Helenians}。但刻尔苏斯没有注意到，西满派根本不承认耶稣是天主之子，而是称西满为天主的\Quote{能力}，关于他他们讲述某些奇妙故事，说他以为如果他能拥有像他相信耶稣所具有的类似能力，他也会像耶稣在群众中那样在人们中间变得强大。但无论是刻尔苏斯还是西满，都不能理解耶稣如何像天主圣言的善牧一样，能够用他的教义播种希腊的大部分地区和野蛮人的土地，使这些地方充满那些使灵魂从一切邪恶中回转、并使之归向万有创造者的话语。刻尔苏斯还认识某些\OriginalTerm{玛尔色利纳派}{Marcellians}（得名于玛尔色利纳），和\OriginalTerm{哈波克拉提派}{Harpocratians}（得名于撒罗默），以及另一些得名于\OriginalTerm{玛利亚姆}{Mariamme}的人，还有得名于玛尔大的人。然而我们，由于热爱学习，在能力范围内不仅考察圣经的内容及其引起的分歧，而且为了热爱真理，我们尽可能研究了哲学家的观点，却从未遇到过这些派别。他还提到了\OriginalTerm{玛尔基翁派}{Marcionites}，其领袖是玛尔基翁。\par

% structure: p0128 source=h2 target=subsection reason=relative-heading-rank; base=h2

\subsection{第六十三章}

接下来，为了显得知道得比他已提到的还要多，他依照惯常的做法说：\Quote{还有其他一些人，邪恶地发明了某一位作为他们的老师和神明，并在极大的黑暗中翻滚，比埃及人安提诺乌斯的同伴更为不洁和可咒骂。}在我看来，他在谈到这些事情时，确实有几分真理地说，确实有一些其他的人，邪恶地发明了另一个神明，并发现他是他们的主，因为他们在那极大的无知黑暗中翻滚。然而，对于那个与我们的\Person{耶稣}相比较的\Person{安提诺乌斯}，我们将不再重复我们在前文中已经说过的话。\Quote{此外，}他接着说，\Quote{这些人彼此说出可怕的亵渎话，说各种羞于出口的事情；他们丝毫不为和谐而让步，以完全的仇恨彼此憎恨。}对此，我们已经说过，在哲学和医学中，派别与派别相互争斗。我们这些跟随\Person{耶稣}话语的人，操练自己思考、言说和行与祂话语相合的事，\BibleQuote{被人辱骂，就祝福；被人迫害，就忍受；被人毁谤，就劝慰；}我们不会对我们持不同意见的人说出\Quote{各种羞于出口的事情}，而是尽可能努力提升他们，使他们通过唯独归附造物主而达到更好的状态，并引导他们像那些将要受审判的人一样行事。如果那些持不同意见的人不肯信服，我们遵守针对这类人的训诫：\BibleQuote{异端者，第一次和第二次劝诫之后，就应当拒绝，知道这样的人已经偏离正道，犯罪，自己定自己的罪。}此外，我们懂得\BibleQuote{缔造和平的人是有福的}，以及\BibleQuote{温良的人是有福的}这些箴言，不会仇视基督信仰的败坏者，也不会称那些陷入错误的人为\Person{喀耳刻}和谄媚的骗子。\par

% structure: p0130 source=h2 target=subsection reason=relative-heading-rank; base=h2

\subsection{第六十四章}

我认为\Person{塞尔苏斯}误解了宗徒的陈述，其中宣告\BibleQuote{在末后的时期，有些人将背弃信德，听从诱惑人的神和魔鬼的教义；在虚伪中说谎，良心被烙铁烙过；禁止嫁娶，又禁戒吃那些天主所创造、为相信的人存感谢而领受的食物；}他也误解了那些用宗徒的这些声明来反对那些败坏基督信仰教义的人。正是由于这个原因，\Person{塞尔苏斯}说\Quote{基督徒中有一些人被称为\InnerQuote{耳朵被烙}；}还有一些被称为\Quote{谜语}——这个词我们从未见过。\Quote{绊脚石}这个表达确实在这些著作中频繁出现——我们习惯于用它来指那些使纯朴之人及易受欺骗者偏离健康教义的人。但是，无论是我们，还是我认为的任何其他基督徒或异端者，都不知道有被称为\Person{塞壬}的人，他们出卖人、欺骗人、堵住耳朵、把受骗的人变成猪。然而，这位自命无所不知的人却说了如下的话：\Quote{你可以听到，}他说，\Quote{所有那些分歧如此之大、在争论中用最无耻的语言相互攻击的人，口里说着：\InnerQuote{世界之于我已被钉在十字架上，我之于世界亦然。}}看来，\Person{塞尔苏斯}只能从\Person{保禄}的著作中记住这一句话；然而，我们为什么不可以也使用圣经中无数其他的引文呢？例如，\BibleQuote{因为我们虽然在血肉中行事，却不随从血肉作战；（因为我们作战的武器不是血肉的，而是靠着天主有能力攻陷堡垒，）摧毁计谋，以及一切反抗天主认识的高傲事物。}\par

% structure: p0132 source=h2 target=subsection reason=relative-heading-rank; base=h2

\subsection{第六十五章}

然而，既然他断言：\Quote{你们可以听到所有如此分歧的人说：\InnerQuote{世界于我已被钉在十字架上，我于世界亦然}}，我们要指出这种说法的不实。因为有一些异端派别不接受宗徒\Person{保禄}的书信，例如厄彼翁派和被称为恩克拉蒂派的派别。那么，那些不视宗徒为圣善智慧之人的人，就不会采用他的言语，说：\Quote{世界于我已被钉在十字架上，我于世界亦然。}因此，在这一点上，\Person{策尔苏斯}也同样犯了虚妄之罪。此外，他继续纠缠于他对现存各派分歧的指控，但在我看来，他所用的措辞似乎并不准确，也没有仔细观察或理解：那些在学习上有所进步的基督徒如何说自己比犹太人拥有更大的知识；以及他们是承认同一部圣经但作不同解释，还是不承认这些书卷为神圣。因为我们发现这两种观点在各派中都很盛行。他接着说道：\Quote{虽然他们的教义毫无根据，但让我们审视这体系本身；首先，让我们提及他们因无知和误解所导致的败坏，当他们在讨论基本原理时，竟以最荒谬的方式表达他们对所不了解之事的意见，例如以下这些。}然后，针对基督徒口中常说的某些表达，他提出哲学家著作中的某些其他表达与之对立，目的是要表明：\Person{策尔苏斯}认为基督徒所使用的崇高情感，早已被哲学家以更优美、更清晰的语言表达出来，以便他将那些被表现出崇高与虔敬品格的意见所吸引的人，拖向哲学的研究。然而，我们将在此结束第五卷，并从下文开始第六卷。\par

\SourceRecord{Translated by Frederick Crombie.；Ante-Nicene Fathers；Vol. 4.；Edited by Alexander Roberts, James Donaldson, and A. Cleveland Coxe.；Buffalo, NY: Christian Literature Publishing Co.,；1885.；<http://www.newadvent.org/fathers/04165.htm>.}

% structure: p0001 source=h1 target=section reason=page-title

\section{\WorkTitle{驳切尔苏斯，第六卷}}

% structure: p0002 source=h2 target=subsection reason=relative-heading-rank; base=h2

\subsection{第一章}

在开始这第六卷时，我敬爱的\Person{盎博罗削}，我们希望在其中回答\Person{切尔苏斯}对\Emph{基督徒}提出的那些指控，并非如人们可能认为的那样，是他从\Emph{哲学著作}中引用的那些反对意见。因为他引用了大量段落，主要来自\Person{柏拉图}，并将圣经的神圣宣言与这些并列，这些宣言甚至能打动理智的心灵；他还断言：\Quote{这些事情在希腊人那里（比在圣经中）被陈述得更好，并且以一种完全没有夸张和天主或天主之子的应许的方式。}现在我们主张，如果真理的使节们的目的是造福尽可能多的人，并尽他们所能，通过他们对人的爱，赢得每一个人——无论是有智慧的还是简单的，不仅仅是希腊人，还有蛮族人（确实，能够使粗野和无知的人皈依乃是一种伟大的仁爱），那么显然，他们必须采用一种对所有人都能受益的说话方式，并赢得各种类型的人。另一方面，那些远离无知的人，视他们为奴隶，无法理解流畅且有逻辑的演讲的华丽段落，因此只关注那些在文学修养中长大的人，他们限制了对公共利益的看法，使之变得非常狭窄。\par

% structure: p0004 source=h2 target=subsection reason=relative-heading-rank; base=h2

\subsection{第二章}

我之所以说这些，是为了回应\Person{切尔苏斯}和其他人对圣经语言简朴的指控，这种简朴似乎被华丽演讲的光彩所遮蔽。因为我们的先知、耶稣本人以及祂的宗徒们，都注意采用一种不仅传达真理，而且适合赢得大众的说话方式，直到每个人被吸引并引领，尽可能上升到对这些看似简单的话语中所包含奥秘的理解。因为，如果我敢这样说，很少有人（如果他们确实受益的话）从\Person{柏拉图}和像他那样写作的人优美而华丽的风格中受益；相反，许多人都从那些以简单实用、着眼于大众需要的方式写作和教导的人那里得到了益处。确实，很容易观察到，\Person{柏拉图}只出现在那些自称是文学人士的手中；而\Person{埃皮克提图}则受到普通能力的人们的钦佩，他们渴望得到益处，并察觉到可以从他的著作中获得的进步。我们现在说这些，并不是要贬低\Person{柏拉图}（因为广大的世界甚至发现他也有用），而是要指出那些人的目的，他们说：\Quote{\BibleQuote{我说的话、讲的道，不是用智慧委婉的言语，而是在于圣神和大能的明证，叫你们的信德不在乎人的智慧，只在乎天主的大能。}}因为天主的圣言宣告，宣讲（尽管本身是真实的，最值得相信）不足以触及人心，除非说话者从天主那里获得某种能力，并且他的话语上显现出恩宠；只有通过神圣的作为，这在有效说话的人身上才会发生。先知在第六十七篇圣咏中说：\Quote{\BibleQuote{上主将赐给宣讲者一句话，大有能力。}}因此，如果在某些方面，希腊人那里存在与我们的圣经相同的教义，但它们并不具有同样的吸引和配置人的灵魂去跟随它们的能力。因此，耶稣的门徒们，就希腊哲学而言是无知的人，却走遍了世界的许多国家，根据圣言的意愿，给每个听众留下深刻的印象，按照他们的功过，使他们获得道德上的改善，与他们的意志倾向接受善的程度成正比。\par

% structure: p0006 source=h2 target=subsection reason=relative-heading-rank; base=h2

\subsection{第三章}

那么，让古代的智者向那些能够理解他们的人宣讲他们的言论吧。假设例如\Person{柏拉图}，\Person{阿里斯顿}的儿子，在他的一封书信中谈论\Quote{至善}，并且说：\Quote{至善绝不能通过言语来描述，而是通过长期的习惯产生，并像从迸发的火焰中突然在灵魂中点燃一道光。}那么，我们听到这些话，承认它们说得很好，因为是天主向人启示了这些以及所有其他高贵的表达。正是因为这个原因，我们主张那些对天主持有正确观念，但没有按照真理向祂献上敬拜的人，应受针对罪人的惩罚。因为关于这样的人，\Person{保禄}明确地说：\Quote{\BibleQuote{天主的忿怒从天上显明在一切不虔不义的人身上，就是那些以不义压制真理的人。因为关于天主的知识，在他们里面是显明的，因为天主已经给他们显明了。自从创世以来，祂那看不见的属性和永恒的大能和神性，在受造之物上清楚可见，被人了解，所以他们没有借口。因为他们虽然认识天主，却不把祂当作天主荣耀祂，也不感谢祂，反而在思想中变为虚妄，他们无知的心就昏暗了。他们自称是聪明的，反而成了愚拙，将不朽坏的天主的荣耀变为偶像，仿佛必朽坏的人和飞禽、走兽、爬虫的形象。}}因此，真理确实被那些认为\Quote{至善不能通过言语描述}，但主张\Quote{经过长期的习惯和熟悉的使用，一道光突然在灵魂中点燃，仿佛由迸发的火焰，并且现在它自己支持自己}的人（以不义）压制了。\par

% structure: p0008 source=h2 target=subsection reason=relative-heading-rank; base=h2

\subsection{第四章}

虽然如此，那些以这种方式论及\Quote{至善}的人会前往\Place{比雷埃夫斯}向阿尔忒弥斯祈祷，好像她就是天主；他们会赞许愚昧者的庄严集会；在发表关于灵魂如此深奥的哲学言论并描述其于美德生活后进入更幸福世界的旅程之后，他们从天主所启示的那些伟大主题转向卑下琐碎的思想，并向埃斯科拉庇俄斯献上一只公鸡！尽管他们本已能从世界的创造和可感事物的默观中得知天主\BibleQuote{不可见的事}和事物的\Quote{原型}，并由此上升到仅凭理智理解的对象，且对祂的\BibleQuote{永恒的大能和神性}并非毫无认识，但他们却\BibleQuote{在思念中变为愚妄}，他们的\BibleQuote{愚昧的心}陷入了对天主（真正）敬拜的黑暗和无知。此外，我们看到那些以智慧和神学自傲的人，敬拜一个可朽坏之人的像——他们说这是为了尊荣祂——有时甚至与埃及人同流，敬拜飞鸟、四足走兽和爬虫！纵然有人似乎超越了这些行为，但仍会发现他们\BibleQuote{将天主的真理变为谎言}，敬拜事奉\BibleQuote{受造物胜过造物主}。既然希腊人中的智慧博学者在事奉天主时犯了错误，天主\BibleQuote{拣选了世上愚拙的，叫智慧的羞愧；又拣选了世上卑贱的、软弱的、被人轻视的、以及那无有的，为要废掉那有的}；这确实\BibleQuote{使一切血肉在天主前无所夸耀}。然而，我们的智者——其中最古老的梅瑟以及随后的先知们——知道至善绝不能用言语描述，于是首先写道：天主向应得之人、向有资格看见祂的人显现时，祂曾显现给亚巴郎、依撒格或雅各伯。至于显现的是谁、有什么形状、如何显现、像哪一位凡人，他们留给那些能证明自己与天主所显现的人相似的人去探究：因为天主不是被人肉眼所看见，而是被纯洁的心所看见。按我们的耶稣所言：\BibleQuote{心里洁净的人是有福的，因为他们要看见天主。}\par

% structure: p0010 source=h2 target=subsection reason=relative-heading-rank; base=h2

\subsection{第五章}

然而，灵魂中突然点燃光明，仿佛有火跃出，这在我们的圣经中早已是熟知的事实；如先知所说：\Quote{你们为自己点亮知识的光明。}此后生活的若望也说：\BibleQuote{在圣言内的生命，是人的光}，这\BibleQuote{真光普照每个人来到这世上}（即那被理智所感知的真世界），并使他成为世界之光：因为这光照在我们心中，为使我们得知天主光荣福音的光，这光显在基督耶稣的面容上。因此，那位在居鲁士统治前许多世纪就预言（他比居鲁士早十四代以上）的古老先知如此说：\BibleQuote{上主是我的光明，我的救援，我还畏惧谁呢？}以及\BibleQuote{祢的话是我脚前的明灯，是我路上的光}；又说\BibleQuote{上主，祢面容的光辉已向我们显现}；\BibleQuote{在祢的光明中，我们才能看见光明。}圣言劝勉我们来到这光明中，在依撒意亚的预言中说：\BibleQuote{耶路撒冷啊，兴起发光，因为你的光明已经来到，上主的荣耀已经升起来照耀你。}同一先知在预言耶稣来临——祂将使人远离对偶像、像和魔鬼的敬拜——时说：\BibleQuote{在黑暗中、死亡阴影之地居住的人，有光升起照耀他们}；又说\BibleQuote{那些坐在黑暗中的百姓，看见了大光。}现在请观察柏拉图关于\Quote{至善}的华丽辞藻与我们的先知关于蒙福之光的宣告之间的区别；并注意到柏拉图所含关于这主题的真理并未帮助他的读者达到对天主纯洁的敬拜，甚至他自己也未能如此——尽管他对\Quote{至善}发出了如此宏大的哲学论述——而圣经的简朴语言却使诚实的读者充满神圣的精神；这光靠油在他们内滋养着，正如某比喻中五个明智童女手中的灯得以保持明亮。\par

% structure: p0012 source=h2 target=subsection reason=relative-heading-rank; base=h2

\subsection{第六章}

然而，既然刻尔苏斯引用了柏拉图一封信中的另一段话，其内容如下：\Quote{如果我以为这些事能够以文字和口头演讲充分地向大众解释，那么我一生还有什么更高尚的追求，比把对人类如此有益、并引导所有人的本性朝向光明的东西写成文字更值得呢？}——那么让我们简要思考这一点，即柏拉图是否知道比他的著作中所包含的任何更深刻的教义，或比他留给我们的任何更神圣的事，让每个人按自己的能力去研究这个问题，同时我们要证明我们的先知们确实知道比圣经中任何内容更伟大的事，但他们并未将其写下来。例如，厄则克耳接受了一个里外都写满字的卷轴，其中包含\Quote{哀歌、歌曲、谴责}；但遵照圣言的命令，他吞下了那书卷，以免其内容被记录下来，从而让不配的人知道。若望也被记载看见了同样的事并做了类似的事。甚至保禄也听到了\BibleQuote{不可言传的言语，是人不能说出的}。关于比这一切都大的耶稣，记载说他私下与门徒交谈，尤其是在他们神圣的退隐处，谈论关于天主的福音；但他所说的话没有被保存下来，因为福音作者们认为这些话无法以文字或言语充分地向大众传达。如果重复这些杰出人物的真相不令人生厌的话，我要说他们比柏拉图更清楚地（借助他们从天主的恩宠所获得的智慧）看出哪些事应当\Emph{写}下来、如何写，哪些事绝不能写给大众，哪些事应当用\Emph{言语}表达，哪些事不应当如此传达。此外，若望在教导我们应当写和不应写之间的区别时，宣称他听到七道雷对他指示某些事，并禁止他将这些话写下来。\par

% structure: p0014 source=h2 target=subsection reason=relative-heading-rank; base=h2

\subsection{第七章}

在梅瑟和先知们的著作中——他们不仅早于柏拉图，甚至早于荷马和希腊文字的发明——也能找到一些段落，配得上天主赐予他们的恩宠，充满了伟大的思想，他们将这些思想表达出来，但并不是因为如刻尔苏斯所想象的，他们不完全理解柏拉图。因为他们怎么可能听到一个尚未出生的人的话呢？如果有人将刻尔苏斯的话应用于耶稣的宗徒们——他们比柏拉图年轻——那么请问，这表面上难道不正是难以置信的断言吗：即帐棚匠保禄、渔夫伯多禄以及离开父亲渔网的若望，竟然因误解柏拉图书信中的语言而表达出他们关于天主的那番话？但既然刻尔苏斯在多次要求我们立即同意（他的观点）之后，好像他现在又唠叨了一些超出他先前所讲的新东西，只不过是在重复自己，那么我们所答复的或许就够了。然而，既然他引用了柏拉图的另一段话，声称问答的方法启发了像他那样研究哲学的人的思想，那就让我们从圣经中证明，天主的圣言也鼓励我们实践辩证术：例如，撒罗满在一处宣称，\Quote{未经质询的教诲会走入歧途；}而留下称为\WorkTitle{智慧}的著作的耶稣、息辣之子，在另一处宣称，\Quote{愚人的知识如同经不起考验的话语。}然而我们的讨论方法较为温和；因为我们知道，主持宣讲圣言的人应当能够驳斥反对者。但如果有些人懒惰，不训练自己以便专注于阅读圣言、\BibleQuote{查考圣经}，并遵照耶稣的命令探究圣经的\Emph{含义}、为此向天主祈求、并持续\BibleQuote{敲击}其中可能关闭的部分，那么圣经并不因此被视为缺乏智慧。\par

% structure: p0016 source=h2 target=subsection reason=relative-heading-rank; base=h2

\subsection{第八章}

接下来，在其他一些柏拉图式的论述之后，这些论述证明\OriginalTerm{善}{the good}只能被少数人所认识，他补充道：\Quote{既然群众因轻蔑他人而妄自尊大，这远非正确，且充满虚荣和狂妄的期望，便断言，因为他们已得知某些可敬的教义，所以某些事情就是真理。}\Quote{然而，尽管柏拉图预言了这些事，他并不谈论奇迹，也不堵住那些希望就他承诺之事向他请教之人的口；他也没有命令他们立刻前来相信存在一个特定种类的天主，且祂有一位特定本性的儿子，这儿子曾降临（地上）并与我交谈。}\newline{}如今，针对这一点我们必须回应：关于\Person{柏拉图}，我想是\Person{阿里斯坦德}曾记载，他不是\Person{阿里斯顿}之子，而是一个幻影之子，这幻影以\Person{阿波罗}的形像接近\Person{安菲克提俄涅}。还有好几位其他追随柏拉图的人，在他们所写的其师传记中，也作了同样的陈述。\newline{}再者，关于\Person{毕达哥拉斯}，我们该说什么呢？他讲述了尽可能多的奇迹：在希腊人的大众集会上，他展示了他的象牙大腿，并声称认出了他从前作为\Person{欧福耳布斯}时曾佩戴的盾牌；他还被说成曾在同一天出现在两个不同的城市！\newline{}此外，凡宣称关于\Person{柏拉图}和\Person{苏格拉底}的传闻属于奇迹者，将会引用天鹅的故事：这天鹅被推荐给睡梦中的\Person{苏格拉底}，以及老师在遇到那年轻人时所说的话：\Quote{那么，这就是那只天鹅了！}\newline{}不仅如此，\Person{柏拉图}自认为拥有的第三只眼，他也会归入奇观之列。\newline{}但是，对于心怀恶意、并乐于诋毁那些超乎常人者之遭遇的人，诽谤的借口永远不会匮乏。\newline{}此类人甚至会嘲笑\OriginalTerm{苏格拉底的神明}{demon of Socrates}，视之为虚构。\newline{}因此，当我们叙述\Person{耶稣}的历史时，我们并不讲述奇迹，祂的真正门徒也没有记载任何此类关于祂的故事；然而这个自诩拥有普世知识、并引用许多柏拉图言论的\Person{凯尔苏斯}，我认为，他有意对\WorkTitle{柏拉图致赫尔墨亚斯和科里斯库斯的信函}中所论的关于\OriginalTerm{天主圣子}{Son of God}的言论保持沉默。\newline{}柏拉图的原文如下：\Quote{并且呼求万有之天主为见证——祂是现在和将来之事的统治者，既是\OriginalTerm{统治者和原因}{the ruler and cause}的父和主——如果我们是真哲学家，我们所有人都将清楚地认识祂，如同有福的人类所能达到的认识那样。}\par

% structure: p0018 source=h2 target=subsection reason=relative-heading-rank; base=h2

\subsection{第九章}

凯尔苏斯引用了柏拉图的另一句话，内容如下：\Quote{我曾想到就这些主题再稍作更为详细的阐述，或许我能够比我之前所做的更清晰地表达我自己；因为有一个\OriginalTerm{真实的原因}{real cause}，它阻碍了那些已冒险（哪怕仅在轻微程度上）就此类主题写作的人；而这一点我之前经常提到，现在似乎应该陈述它。在所有必然被运用于知识获取的现存事物中，有三个元素；知识本身是第四个；还应该确立第五个，它既是可知的也是真实的。在这些元素中，一个是\OriginalTerm{名称}{name}；第二个是\OriginalTerm{言语}{word}；第三个是\OriginalTerm{肖像}{image}；第四个是\OriginalTerm{知识}{knowledge}。}\newline{}如今，按照这一划分，\Person{若翰}是在\Person{耶稣}之前被引入，作为在旷野中呼喊者的声音，以致于对应于柏拉图所说的\OriginalTerm{名称}{name}；若翰之后所指出的第二个人是\Person{耶稣}，与耶稣一致的是\BibleQuote{圣言成了血肉}这一陈述；而这对应于柏拉图的\OriginalTerm{言语}{word}。柏拉图称第三个为\OriginalTerm{肖像}{image}；但我们，将\OriginalTerm{肖像}{image}这个词应用于不同的东西，会更精确地说：由圣言在灵魂中所造成的创伤印记，就是在我们每人内的基督；这印记是由\OriginalTerm{基督圣言}{Christ the Word}所印下的。至于基督，即在我们中那些完美者内的智慧，是否对应于第四个元素——\OriginalTerm{知识}{knowledge}，将被他有能力查明的人所知晓。\par

% structure: p0020 source=h2 target=subsection reason=relative-heading-rank; base=h2

\subsection{第十章}

他接着说：\Quote{你看\Person{柏拉图}，虽然主张至善无法用言语描述，但为了避免显得退缩到无可辩驳的立场，他附上了一个理由来解释这一困难，因为甚至连\InnerQuote{无}或许也可以用言语来解释。}但\Person{策尔苏斯}引述这一点，以证明我们不应只是简单地表示同意，而应为我们的信仰提供理由，因此我们也将引述\Person{保禄}的话，他在责备轻率的信徒时说：\Quote{除非你们轻率地相信了。}现在，因着反复自言的习癖，\Person{策尔苏斯}尽其所能迫使我们陷入同义反复，在引述过的夸耀之辞后，他再次重申：\Quote{\Person{柏拉图}并未犯有夸耀和虚谎的罪，宣称自己有了什么新发现，或从天上降下来宣布它，而是承认这些陈述从何而来。}现在，如果有人想答复\Person{策尔苏斯}，可以针对这类断言回答说，连\Person{柏拉图}也犯有夸耀之罪，当他在\WorkTitle{蒂迈欧篇}中借\Person{宙斯}之口说出以下话语：\Quote{诸神之神，我是他们的创造者和父亲，}等等。如果有人因着\Person{宙斯}之名所传达的意义而为这类话语辩护，既在\Person{柏拉图}的对话中如此说，那么，探究天主之子的话语或\WorkTitle{先知书}中造物主之言的人，为何不能表达出比\WorkTitle{蒂迈欧篇}中\Person{宙斯}之言更深刻的意义呢？因为神性的特征在于宣告未来的事件，这些预告并非出于人的能力，而是结果证明出于宣告者身上的神圣之灵。因此，我们并不对每一位听众说：\Quote{首先相信，我所引介给你们的这一位是天主之子；}而是将福音放在每个人面前，视其性情和倾向是否适合接受，因为我们已学会知道\Quote{当如何答复每一个人。}有些人只能接受劝勉而相信，我们就只对他们说这些；而对另一些人，我们则尽可能以问答方式通过论证来接近他们。我们也绝不说，如\Person{策尔苏斯}嘲弄地指控的那样：\Quote{相信我所引介给你们的这一位是天主之子，尽管他曾被可耻地捆绑、受辱，并且最近在众人眼前被极其轻蔑地对待；}我们也不附加说：\Quote{正因如此，更要相信。}因为我们努力在每一个具体论点上提出比前面篇章中更多的论据。\par

% structure: p0022 source=h2 target=subsection reason=relative-heading-rank; base=h2

\subsection{第十一章}

此后\Person{策尔苏斯}接着说：\Quote{如果这些人（指基督徒）引介这一位，而另一些人又引介另一位（作为基督），各派普遍即时的呼喊是：\InnerQuote{相信吧，如果你愿意得救，否则就离开，}那些切望得救的人该做什么呢？他们岂不要掷骰子来决定他们该投奔何处、加入谁吗？}现在我们要以如下方式来答复这一异议，因为事实的清晰性促使我们这样做。假如曾有记录说，有几个人以耶稣的方式作为天主之子出现在世上，并且每个人都吸引了一群追随者，以至于因着每个人自称是天主之子这种宣称的相似性，其追随者作证的那一位成了争论的对象，那么他所说的就有理由了：\Quote{如果这些人引介这一位，而另一些人引介另一位，各派普遍即时的呼喊是：\InnerQuote{相信吧，如果你愿意得救，否则就离开，}}等等；然而，耶稣基督是唯一降临人间的天主之子，这已向全世界宣告了：因为那些像\Person{策尔苏斯}一样、认为（耶稣的事迹）是一系列奇事、并为此希望行同样的事、以便自己也能获得类似对人的心灵控制的人，被证实是彻底的虚无。这样的人包括\Place{撒玛黎雅}的术士\Person{西满}，以及同地的\Person{多西修斯}；前者声称自己是那被称为大能的天主之能，后者则声称自己是天主之子。但西满派在全世界都找不到踪迹；而且，\Person{西满}为了为自己争取许多追随者，教导门徒视偶像崇拜为无关紧要之事，从而免除了基督徒所受教导宁愿接受的死亡危险。但即使在西满派开始存在之初，他们也未遭受迫害。因为那阴谋对抗耶稣教导的邪恶魔鬼很清楚，他自己的任何准则都不会因西满的教导而削弱。至于\Person{多西修斯}派，即使在前几个世纪也未曾显赫，如今已完全消失，据说他们总数不满三十人。\Place{加里肋亚}的\Person{犹达斯}，正如\Person{路加}在\BibleBook{宗徒}{大事录}中所记述的，也妄想自称某位大人物，正如他之前的\Person{特乌达斯}一样；但因为他们的道理不是出于天主，他们便被消灭，所有服从他们的人立即星散。所以我们并不\Quote{掷骰子来决定我们该投奔何处、加入谁，}仿佛有很多自称从天降临访问世人的声称者能够吸引我们跟随他们一样。然而，在这一点上，我们已经说得够多了。\par

% structure: p0024 source=h2 target=subsection reason=relative-heading-rank; base=h2

\subsection{第十二章}

因此，讓我們轉向\Person{塞爾蘇斯}提出的另一項指控。他甚至不熟悉（我們聖書的）言詞，卻因誤解而說：\Quote{我們宣稱在人間的智慧，在天主前是愚拙的。}\Person{保祿}曾說：\Quote{這\Emph{世界}的智慧，在天主前是愚拙的。}\Person{塞爾蘇斯}說：\Quote{這理由早已被陳述過。}他認為這理由是\Quote{我們希望藉此說法只贏得無知和愚笨的人}。但正如他自己所暗示的，他之前已經說過同樣的話；我們也盡力回應了。儘管如此，他仍想表明這句話是我們的發明，並借用自希臘哲人，他們宣稱人的智慧是一種，神的智慧是另一種。他引用\Person{赫拉克利特}的話，在一處說：\Quote{人的行為方式不受固定原則支配，但神的行為方式受固定原則支配}；在另一處說：\Quote{愚笨的人聽從魔鬼，如同孩童聽從成人}。此外，他還引用以下來自\Person{柏拉圖}所著的\WorkTitle{蘇格拉底的申辯}中的話：\Quote{雅典人啊，我獲得這個名聲，無非是因我的智慧。而這是哪一種智慧呢？大概是人的智慧；因為在那方面，我敢說我實際上是智慧的。}這些是\Person{塞爾蘇斯}援引的段落。但我還要補充以下\Person{柏拉圖}寫給\Person{赫爾墨亞斯}、\Person{埃拉斯托斯}和\Person{科里斯庫斯}的信中的話：\Quote{對\Person{埃拉斯托斯}和\Person{科里斯庫斯}，我雖然年邁，仍要說：除了他們擁有的關於\InnerQuote{形式}的高貴知識外，他們還需要一種智慧，針對邪惡和不義之人的階層，作為抵禦他們的保護和驅逐力量。因為他們不諳世事，由於他們大部分人生與我們這些溫和而非邪惡的人一起度過。所以我說他們需要這些，免得他們被迫忽視真正的智慧，而過度專注於必要和屬人的事務。}\par

% structure: p0026 source=h2 target=subsection reason=relative-heading-rank; base=h2

\subsection{第十三章}

根據上述，一種智慧是人的，另一種是神的。現在，所謂\Quote{人的}智慧，我們稱之為\Quote{世界的}智慧，這在\Quote{天主前是愚拙的}；而\Quote{神的}智慧——不同於\Quote{人的}，因為它是\Quote{神的}——是藉著賜予它的天主的恩寵，賜給那些證明自己有能力接受它的人，尤其是那些認識兩種智慧不同，並在祈禱中向天主說：\BibleQuote{縱使在人的兒子中有完美的，但若缺少從您而來的智慧，他仍將被視為無有。}\BibleRef{智 9:6}我們確實主張，\Quote{人的}智慧是靈魂的操練，而\Quote{神的}智慧是\Quote{終點}，也被人稱為靈魂的\Quote{乾糧}，如經上所說：\Quote{乾糧是屬於成全的人，他們因習慣而鍛煉了感官，以辨別善惡。}此外，這觀點的確是古老的，其古老性並非如\Person{塞爾蘇斯}所認為僅追溯到\Person{赫拉克利特}和\Person{柏拉圖}。因為在這些人物之前，先知們已經區分了兩種智慧。現在只需引用\Person{達味}關於那按神聖智慧而智慧的人所說的話：\Quote{他看見智者死去時，將不會見到腐朽。}因此，神聖的智慧不同於信德，是天主所謂\Quote{\OriginalTerm{神恩}{charismata}}中的\Quote{第一種}；在其之後的\Quote{第二種}——在那些知道如何準確區分這些事的人看來——是所謂的\Quote{知識}；而\Quote{第三種}——因為即便較單純的、盡其所能事奉天主的階層也必須得救——是信德。因此\Person{保祿}說：\Quote{聖神賜給此人智慧之言，同一聖神賜給另一人知識之言，又賜給另一人信德。}因此，你會發現並非普通人，而是那些卓越而傑出的基督徒，參與了\Quote{神聖的}智慧；因為關於神聖智慧的主題，不是對\Quote{人類中最無知、最卑賤或最未受教導的人}談論的。\par

% structure: p0028 source=h2 target=subsection reason=relative-heading-rank; base=h2

\subsection{第十四章}

\Person{凯尔苏斯}以\Quote{无知的、奴性的、愚昧的}等称号指称他人，我想，他是指那些不熟悉他的法律、未受希腊学识训练的人；而另一方面，我们认为那些不耻于向无生命之物祈求（恳求）、向没有力量者求健康、向死者求生命、向无助者求援助的人才是\Quote{无知的}。虽然有些人会说这些对象不是神，只是真神的仿像和象征，然而正是这些人在认为低贱工匠的手可以仿造神的形像时，是\Quote{无知的、奴性的、愚昧的}；因为我们断言，我们中最卑微的人已从这种无知和缺乏知识中解脱，而最聪明的人能理解并把握神圣的希望。然而我们\Emph{不}主张，一个未经世上智慧训练的人不可能接受\Quote{神圣的}；但\Emph{确实}承认，所有人类的智慧与\Quote{神圣的}相比都是\Quote{愚拙的}。接下来，他非但没有尽力提出理由，反而称我们为\Quote{巫术者}，并断言\Quote{我们仓皇逃避更文雅的人群，因为他们不适合成为我们欺骗的对象，而我们试图诱惑更粗野的人}。但他没有注意到，从一开始，我们的智者就受过外在学识的训练：例如\Person{梅瑟}，学尽了埃及人的一切智慧；\Person{达尼尔}、\Person{阿纳尼雅}、\Person{阿匝黎雅}、\Person{米沙耳}，学尽了亚述人的一切学识，以致他们在那国所有智者之上十倍。此外，如今教会中，相对于（普通信徒）的众多数量，也有一些\Quote{智慧}的人，他们从被我们称为\Quote{属血肉的}智慧归向教会；还有一些人由此进至那\Quote{神圣的}智慧。\par

% structure: p0030 source=h2 target=subsection reason=relative-heading-rank; base=h2

\subsection{第十五章}

\Person{凯尔苏斯}接着，像一个听闻过谦卑话题被大加谈论的人，却不曾费心理解它，想要对我们的谦卑实践加以恶评，并以为这是从\Person{柏拉图}某些未能领会的言辞中借来的，\Person{柏拉图}在\WorkTitle{法律篇}中如是说：\Quote{如今，按照古训，天主自身具有万物的开始、终结和中间，祂按照本性运行，径直前进。正义始终跟随祂，是违反神律之事的惩罚者；那将要得到幸福的人，谦卑地、得体地装束着，紧紧跟随正义。}他却没有注意到，在比\Person{柏拉图}古老得多的著作中，有一篇祷文包含这样的话：\BibleQuote{上主，我的心不骄横，我的眼不自高；我不去图谋伟大之事，也不去行超过我能力的事；我若不曾谦卑，等等。}现在这些话表明，那有谦卑心意的人，决不会用不雅或不吉的方式自卑，如屈膝跪地、一头栽倒在地、穿上可怜之人的衣服、或向自己撒灰。但按先知的意思，那有谦卑心意的人，虽然\BibleQuote{行在伟大和奇妙的事}中，即那些超越他能力的——也就是那些真正伟大的教义和那些奇妙的思想——他却\BibleQuote{在天主大能的手下自卑}。如果有些人因愚钝未能清楚理解谦卑的教义，而如此行事，那不应归咎于我们的教义；我们倒必须宽恕这些人的愚钝，他们追求更高的事，却因心智的昏聩未能达到。那\Quote{谦卑而得体地装束着}的人，比\Person{柏拉图}的\Quote{谦卑而得体地装束着}的人更甚：因为他得体地装束着，一方面是因为他\BibleQuote{行在伟大和奇妙的事}中，这些事超越他的能力；另一方面，他谦卑，是因为在身处这些事的同时，他自愿自卑，不是随便在某一位手下，而是在\BibleQuote{天主大能的手下}，藉着\Person{耶稣基督}——这教义的导师——祂\BibleQuote{不以自己与天主同等为应当把持不舍的，反而使自己空虚，取了奴仆的形体，被认作人的样式，就贬抑自己，听命至死，且死在十字架上。}这谦卑的教义是如此伟大，以至于它并非由普通人物教导；而是我们伟大的救主亲自说：\BibleQuote{你们向我学习，因为我是良善心谦卑的；这样你们的心灵必得安息。}\par

% structure: p0032 source=h2 target=subsection reason=relative-heading-rank; base=h2

\subsection{第十六章}

其次，关于耶稣对富人的宣示，当祂说：\BibleQuote{骆驼穿过针孔，比富人进入天主的国还容易。}克耳苏断言这句话显然出自柏拉图，而耶稣歪曲了这位哲学家的言论，即\Quote{不可能同时既因善行出众，又拥有财富。}现在，凡是能对事物稍加留意的人——不仅是在耶稣的信徒中，也在其余人类中——听到耶稣（祂生于犹太人中间，在犹太人中间长大，被认为是木匠若瑟的儿子，且从未学习过文学——不仅希腊文学，连希伯来文学也未学过，正如热爱真理的圣经为祂作证）竟然读过柏拉图，且喜欢他关于富人的观点（即\Quote{不可能同时既因善行出众，又拥有财富}），就歪曲并改成了\BibleQuote{骆驼穿过针孔，比富人进入天主的国还容易}，有谁会不嘲笑克耳苏呢？如果克耳苏不是怀着仇恨与厌恶的心态去阅读福音书，而是充满对真理的热爱，他就会思考：为什么选择骆驼——这种在体型结构上弯曲的动物——作为与富人比较的对象？以及那\BibleQuote{狭窄的针孔}对那看见\BibleQuote{通往生命的门是窄的，路是狭的}的人有什么含义？他还会思考：这动物按照法律被描述为\BibleQuote{不洁}，它有一个可接受的因素（即反刍），却有一个受谴责的因素（即不分蹄）。此外，他还会探究：骆驼在圣经中作为比较对象出现过多少次，针对什么对象？以便他了解\OriginalTerm{圣言}{Logos}论及富人的含义。他也不会不加考察这个事实：耶稣称\BibleQuote{穷人}为\BibleQuote{有福的}，而称\BibleQuote{富人}为\BibleQuote{可怜的人}；这些话是指感官可见的穷人和富人，还是\OriginalTerm{圣言}{Logos}所认识的某种贫穷应被视为\Quote{完全有福}，以及某种富人应被彻底谴责？因为即使是普通人也不会不加区别地赞扬穷人，其中许多人过着极其邪恶的生活。但关于这一点，我们已经说得够多了。\par

% structure: p0034 source=h2 target=subsection reason=relative-heading-rank; base=h2

\subsection{第十七章}

此外，由于克耳苏出于贬低我们的圣经所描述的天主之国的意图，没有引用任何一段（仿佛它们不值得他记载，或者也许因为他并不熟悉），而另一方面，他却引用柏拉图的言论，既从他的\WorkTitle{书信集}也从他\WorkTitle{费德鲁斯篇}中，仿佛这些是神圣启示的，而我们的圣经却不是；因此，让我们提出几点，以便与柏拉图这些貌似合理的声明作比较——这些声明并没有促使这位哲学家以与万有创造者相称的方式来敬拜祂。因为他本不应该用我们称之为\Quote{偶像崇拜}而大众称之为\Quote{迷信}的东西来玷污或污染这种敬拜。现在，按照一种希伯来的修辞方式，在\BibleBook{圣咏}{第十八篇}中论及天主说：\BibleQuote{祂使黑暗成为祂的掩蔽。}这是表明那些应恰当领受的关于天主的观念是不可见、不可知的，因为天主将祂自己隐藏在黑暗中，仿佛从那些不能承受祂知识的光辉、或者不能注视它们的人面前隐藏起来，部分是由于他们的理智被必死的卑微之躯所蒙蔽而受到玷污，部分是由于其理解天主的能力较为微弱。为了表明天主的知识极少赐予人，只在极少数人身上找到，记载说梅瑟进入了天主所在的黑暗。再者，关于梅瑟，经上说：\BibleQuote{惟梅瑟一人可走近上主，别人不可。}此外，为使先知能表明与天主有关的教义的深度，这深度是那些没有\BibleQuote{洞察一切，连天主的深奥事理也洞悉的圣神}的人所不能及的，他又加上：\BibleQuote{深渊如同衣服是他的遮盖。}不仅如此，我们的主和救主，天主的\OriginalTerm{圣言}{Logos}，表明认识圣父的伟大唯独由祂自己适当地理解和知道，其次是由那些心志被圣言和天主本人光照的人，祂宣布：\BibleQuote{除了圣父外，没有人认识圣子；除了圣子和圣子所愿意启示的人外，也没有人认识圣父。}因为没有人能像生养祂的圣父那样相称地认识那\OriginalTerm{无始}{uncreated}的、受造万物的首生者，也没有人能像生活的圣言、祂的智慧与真理那样认识圣父。凡分享祂（圣言）的人，祂从圣父那里除去所谓的\Quote{黑暗}，即祂\Quote{使成为祂的掩蔽}的，以及被称为祂的\Quote{遮盖}的\Quote{深渊}，如此揭开圣父的面纱，每个能够认识圣父的人就认识圣父。\par

% structure: p0036 source=h2 target=subsection reason=relative-heading-rank; base=h2

\subsection{第十八章}

我认为应当从大量段落中引用这几个例子，在这些段落中，我们的神圣作者表达了他们对天主的看法，以向那些有眼能看见圣经之可敬特性的人表明，先知的神圣作品包含着比\Person{刻尔苏斯}所欣赏的\Person{柏拉图}那些言论更值得敬畏的内容。现在，\Person{刻尔苏斯}所引用的\Person{柏拉图}宣言如下：\Quote{万物都围绕着万有之王，万物都为了他而存在，他是一切善事的原因。对于第二等级的事物，他是第二；对于第三等级的事物，他是第三。人的灵魂因此渴望知道这些事物是什么，它注视着与自身同类的事物，其中没有一样是完美的。但至于王和我所提到的那些事物，没有什么能与它们相比。} 此外，我本可以提及那些被希伯来人称为\OriginalTerm{色辣芬}{seraphim}、在\Person{依撒意亚}中所描述的存在者，它们遮盖天主的头和脚，以及那些被称为\OriginalTerm{革鲁宾}{cherubim}、由\Person{厄则克耳}所描述的存在者，它们的姿态，以及天主被描述为乘坐于革鲁宾之上的方式。但由于它们是以一种极其奥秘的方式提及的，为了那些不配和卑鄙的人，他们无法进入神学的伟大思想和可敬本质，我认为在这篇论文中谈论它们不合适。\par

% structure: p0038 source=h2 target=subsection reason=relative-heading-rank; base=h2

\subsection{第十九章}

\Person{刻尔苏斯}接着断言：\Quote{某些基督徒误解了\Person{柏拉图}的话，大声吹嘘一个\InnerQuote{超天}的天主，从而上升到犹太人的天之上。} 的确，通过这些话，他并没有清楚地说清楚他们是上升到犹太人的\Emph{天主}之上，还是只上升到他们所誓言的天空之上。然而，我们目前的目的不是谈论那些承认与犹太人所崇拜的神不同的神的人，而是为自己辩护，并表明犹太人的先知（他们的著作被列入我们的经卷）不可能从\Person{柏拉图}那里借用任何东西，因为他们比\Person{柏拉图}更古老。他们并没有从\Person{柏拉图}那里借用这个声明：\Quote{万物都围绕着万有之王，万物都因他而存在；} 因为我们知道，先知们、耶稣本人和他的门徒们已经说出了比这些更高尚的思想，他们清楚地表明了在他们内的圣神的意义，这圣神正是基督的圣神。这位哲学家也不是第一个展示\Quote{超天}地方的人；因为\Person{达味}很久以前就已经揭示了那些已经上升到可见事物之上的人对天主的深刻而丰富的思想，当他在\WorkTitle{圣咏集}中说：\BibleQuote{苍天之上的天，以及天上的水，请赞美天主，愿它们赞美上主的圣名。} 我确实不否认\Person{柏拉图}从某些希伯来人那里学到了引自\WorkTitle{费德若斯}的那些话，或者，正如一些人记载的那样，他通过阅读我们的先知著作引用了它们，当时他说：\Quote{下方的诗人从未歌颂过这超天的地方，也永远不会以恰当的方式歌颂，} 等等。在同一段落中还有以下内容：\Quote{因为那既无色又无形、无法触摸、真实存在的本质，是灵魂的舵手，仅被理智所看见；而真正的知识种类围绕它占据着这个位置。} 我们的\Person{保禄}，受这些话的教导，渴望\Quote{超现世}和\Quote{超天}的事物，并为了它们竭尽全力去获得，在\WorkTitle{格林多后书}中说：\BibleQuote{我们这轻微而短暂的苦难，为我们成就了极其重大、永恒的荣耀；我们并不注目所见的事物，而是注目那看不见的事物；因为所见的是暂时的，看不见的才是永恒的。}\par

% structure: p0040 source=h2 target=subsection reason=relative-heading-rank; base=h2

\subsection{第二十章}

现在，对于那些能够理解他的人，宗徒清楚地展示了\Quote{可感知的事物}，称它们为\Quote{所见之物}；而他称\Quote{看不见的事物}为理智的对象，只有理智才能认识。他也知道，\Quote{所见}之物是\Quote{暂时的}，而理智可认识的\Quote{看不见的}事物是\Quote{永恒的}；他渴望停留在对这些事物的默观中，并因热切渴望而得到帮助，他认为一切痛苦都是\Quote{轻的}和\Quote{无足轻重的}，在苦难和患难时期，他从未被它们压垮，而是通过默观（神圣的）事物，认为每一场灾难都是轻的，因为我们也有\BibleQuote{一位伟大的大司祭}，他以他的大能和智慧的伟大\BibleQuote{穿过了诸天，就是耶稣，天主子}，他曾应许所有真正学习了神圣事物并与它们和谐生活的人，为他们先行到超现世的事物那里；因为他的话是：\BibleQuote{我去哪里，你们也要在那里。} 因此，我们盼望，在经历此处的苦难和争战之后，到达最高的诸天，并按照耶稣的教导，接受那涌流到永生的水泉，充满知识的河流，并将与那些据说在诸天之上、赞美他圣名的水联合。我们所有赞美他的人，将不会被天的旋转所带走，而是永远沉浸在对天主不可见之事的默观中，这些事不再通过他从创世以来所造之物被我们认识，而是如耶稣的真门徒用这些话所表达的：\BibleQuote{然后面对面；} 以及：\BibleQuote{当那完全的来到时，那部分的就要消失。}\par

% structure: p0042 source=h2 target=subsection reason=relative-heading-rank; base=h2

\subsection{第二十一章}

在天主的众教会中通行的圣经，并未提及\Quote{七}层天，也没有任何确定数目，但它们确实教导存在着\Quote{诸天}，不论这指的是希腊人称为\Quote{行星}的那些天体的\Quote{球体}，或是某种更奥秘的事。策尔苏斯也附和\Person{柏拉图}的意见，断言灵魂能够穿过行星往返于大地；而我们最古老的先知\Person{梅瑟}说，我们的先知\Person{雅各伯}曾见到一个神视——一架梯子直通上天，天主的使者在梯子上上去下来，主则支撑在梯顶——通过这梯子的事，模糊地指向\Person{柏拉图}所见的同一真理，或是比这些更伟大的事。关于这个主题，\Person{斐洛}撰写了一篇论文，值得所有喜爱真理的人进行深思熟虑和明智的研究。\par

% structure: p0044 source=h2 target=subsection reason=relative-heading-rank; base=h2

\subsection{第二十二章}

在此之后，策尔苏斯渴望在他的攻击我们的论文中展现他的学识，也引用了某些波斯秘仪，他说：\Quote{这些事在波斯人的记述中，尤其是在他们中所举行的\OriginalTerm{密特拉}{Mithras}秘仪中，被模糊地暗示。因为在后者中有两个天体运行的图示——即恒星的运动和行星中的运动，以及灵魂穿过这些天体的通道。图示如下：有一架带有高大门的梯子，在梯顶有第八道门。第一道门由铅构成，第二道由锡，第三道由铜，第四道由铁，第五道由混合金属，第六道由银，第七道由金构成。第一道门他们归给土星，以\Quote{铅}指示这颗星的缓慢；第二道归给金星，将她比作锡的光辉和柔软；第三道归给木星，坚定而稳固；第四道归给水星，因为水星和铁都适于承受一切，且能赚钱而勤勉；第五道归给火星，因为由混合金属构成，多变而不均；第六道银门归给月亮；第七道金门归给太阳——从而模仿后二者的不同颜色。} 他接着探讨星辰按此顺序排列的原因，这由其余物质的名称所象征。此外，波斯神学还加入或引用了音乐方面的理由；并且他还试图加上第二个解释，也与音乐考量相关。但在我看来，引用策尔苏斯在这些事上的话是荒谬的，类似于他自己所做的：在他对基督徒和犹太人的指控中，他极不恰当地不仅引用了\Person{柏拉图}的话，而且对这些还不满足，还加上了波斯\OriginalTerm{密特拉}{Mithras}秘仪及其解释。现在，无论这些事如何——无论波斯人和那些主持\OriginalTerm{密特拉}{Mithras}秘仪的人对此给出的是虚假还是真实的解释——他为何选择引用这些，而不是其他秘仪及其解释？因为\OriginalTerm{密特拉}{Mithras}秘仪在希腊人中似乎并不比\Place{厄琉息斯}秘仪，或\Place{埃吉纳}岛上人们受启入\Person{赫卡忒}仪式的秘仪更著名。但如果他必须引入异邦秘仪及其解释，为何不引用埃及人的秘仪（这些被许多人高度尊重），或\Place{卡帕多细亚}人关于科马纳的\Person{狄安娜}的秘仪，或\Place{色雷斯}人的秘仪，甚至\Place{罗马}人自己的秘仪（他们使其元老院最高贵的成员受启）？但如果他认为与这些中的任何一个进行比较都是不恰当的，因为它们没有提供指控犹太人或基督徒的帮助，那么他为什么也不认为引用\OriginalTerm{密特拉}{Mithras}秘仪的例子是不恰当的呢？\par

% structure: p0046 source=h2 target=subsection reason=relative-heading-rank; base=h2

\subsection{第二十三章}

如果有人希望获得更深入默观灵魂进入神圣事物的途径，不是从他所引用的那个非常微不足道的派别的陈述，而是从书籍——一部分是犹太人的书（他们在会堂中诵读，并被基督徒采纳），另一部分是唯独基督徒的书——那么，让他仔细阅读\BibleBook{厄则克耳}{先知书}末尾，先知所见的神视，其中列举了不同种类的门，这些门模糊地指涉神圣灵魂进入更美好世界的不同方式；也让他阅读\BibleBook{若望}{默示录}中关于天主的城、天上的耶路撒冷及其基础和门的记载。如果他有能力找出那些将向神圣事物迈进的人由符号所指明的道路，就让他阅读\Person{梅瑟}所著的名为\BibleBook{}{户籍纪}的书，并寻求一位能够引导他理解关于以色列子民安营的叙述之含义的人的帮助；即：哪些是朝向东方布置的（正如第一个营那样）；哪些是朝向西南和南方；哪些是朝向大海；以及最后那些是朝向北方驻扎的。因为他将看到，各个地方都有一种不容轻率对待的含义，也并非如策尔苏斯所想，只要求愚蠢和奴性的听众。相反，他将在这些安营中分辨出与所列举的数目相关的某些事物，这些事物是特别适合每个支派的，但目前似乎不是讨论这些的适当时机。此外，让策尔苏斯以及所有读他的书的人知道，在真正且由天主认可的圣经中，没有任何一处提到\Quote{七}层天；我们的先知、耶稣的宗徒，以及天主子自己，也没有重复任何从波斯人或\OriginalTerm{迦比里}{Cabiri}那里借来的东西。\par

% structure: p0048 source=h2 target=subsection reason=relative-heading-rank; base=h2

\subsection{第二十四章}

在引用了密特拉教奥秘的例子之后，克耳苏宣称，谁若想探究基督信仰的奥秘，并与前述波斯人的奥秘相比，将两者对照，并揭开基督信仰的仪式，便会看出其间差异。然而，凡是他能列出各派名称的地方，他都毫不犹豫地引用他所认为熟悉的名称；但当他最应该这样做的时候——如果这些名称真的为他所知——并告知我们他所绘制的图表是哪个派别使用的，他却没有这样做。依我之见，他是从某个微不足道、被称为\OriginalTerm{奥菲特派}{Ophites}的派别的某些声明中——他误解了这些声明——部分借用了关于该图表的说法。我们素来好学，曾偶然见到这个图表，并在其中发现了那些人的图像；正如保禄所言，这些人\Quote{\InnerQuote{潜入人家，猎取那些满身罪恶、受各种私欲牵引的愚蠢妇女；她们常常学习，却永远无法认识真理。}}\BibleRef{弟后 3:6-7}然而，这图表如此缺乏可信度，以至于那些易受骗的妇女，乃至最粗鄙的乡民，或任何容易被花言巧语所迷惑的人，都从未接受过它。事实上，我们走遍许多地方，寻访各地自诩有知识的人，却从未遇到任何人相信这图表。\par

% structure: p0050 source=h2 target=subsection reason=relative-heading-rank; base=h2

\subsection{第二十五章}

在这图表中描绘了十个圆圈，彼此分开，但由一个圆圈连接，该圆圈被称为万有之魂，称为\Quote{\OriginalTerm{里维亚坦}{Leviathan}}。犹太圣经称这\OriginalTerm{里维亚坦}{Leviathan}——无论其意为何——是天主所造，作为玩物；因为我们在圣咏集中读到：\Quote{\BibleQuote{祢以智慧创造了万物，大地充满了祢的造物。那里有海洋，广大无边，其中有无数爬行的动物，大小活物都有；那里有船只航行，也有祢所造的\OriginalTerm{巨龙}{Leviathan}，好游戏其中。}}\BibleRef{咏 104:24-26}以\Quote{\OriginalTerm{里维亚坦}{Leviathan}}一词代替\Quote{巨龙}。于是，这邪恶的图表论及这被圣咏作者如此轻看的\OriginalTerm{里维亚坦}{Leviathan}，说它是遍历万物的灵魂！我们还看见图表中有一个称为\Quote{\OriginalTerm{贝赫莫特}{Behemoth}}的存在，似乎置于最底下的圆圈之下。这咒诅图表的绘制者将\OriginalTerm{里维亚坦}{Leviathan}刻在圆周和中心，使其名称出现在两个不同地方。此外，克耳苏说这图表\Quote{\Quote{被一条粗黑线分开}}，他断言这条线称为\Quote{\OriginalTerm{革赫纳}{Gehenna}}，即\Quote{\OriginalTerm{塔尔塔罗斯}{Tartarus}}。既然我们发现福音中提及\Quote{\OriginalTerm{革赫纳}{Gehenna}}为惩罚之地，便查考它是否在古代圣经中提到，尤其是因为犹太人也使用这个词。我们查明，在圣经希伯来文中，名为\Quote{欣嫩子之谷}的地方，本质上与\Quote{\OriginalTerm{革赫纳}{Gehenna}}意义相同，被称为\Quote{欣嫩谷}和\Quote{革赫纳}。继续研究，我们发现称为\Quote{\OriginalTerm{革赫纳}{Gehenna}}或\Quote{欣嫩谷}之地，位于本雅明支派所分得之地内，耶路撒冷也在其中。我们试图查明，属于本雅明支派的\Quote{天上的耶路撒冷}与\Quote{欣嫩谷}之间有何关联，便发现了关于惩罚之地的说法有某种确证，该惩罚之地旨在净化那些需要通过折磨来净化的灵魂，与经上之言相符：\Quote{\BibleQuote{主像炼金之人的火，像漂布之人的碱。祂必坐下如炼净银子和金子的人。}}\BibleRef{拉 3:2-3}\par

% structure: p0052 source=h2 target=subsection reason=relative-heading-rank; base=h2

\subsection{第二十六章}

因此，惩罚将在耶路撒冷附近对接受净化的人实施，这些人已将邪恶的要素——在某个地方被比喻性地称为\Quote{铅}——纳入灵魂的本质；正因如此，在匝加利亚书中，不义被描绘为坐在\Quote{一塔冷通铅}之上。\BibleRef{匝 5:7-8}然而，关于这个主题所能说的言论，既不可对所有人说，也不宜在此时宣述；因为将这类解释写下来并非毫无危险，因为群众只需知道关于罪人惩罚的教训便已足够；而超越这一点，对于那些即便惧怕永罚也难以克制而不陷入任何邪恶程度、不陷入因罪而生的罪恶洪流的人，并无益处。因此，关于\OriginalTerm{革赫纳}{Gehenna}的教义，对那图表和克耳苏而言都是未知的；否则，图表的制作者就不会夸耀他们关于动物和图画的图像，仿佛真理由此彰显；克耳苏在其反对基督徒的论文中，也不会针对基督徒提出指控时，引用他们从未说过的话，而有些话可能出自一些现已消失、或人数极少且屈指可数的人之口。正如那些自称信奉柏拉图学说的人，不适合为伊壁鸠鲁及其不虔诚意见辩护，同样，我们也不适合为这图表辩护，或反驳克耳苏对其的指控。因此，我们可以允许他在这方面的指控被视为多余和无用，因为任何被此类意见所迷惑的人，我们都会比克耳苏更严厉地谴责。\par

% structure: p0054 source=h2 target=subsection reason=relative-heading-rank; base=h2

\subsection{第二十七章}

在图式之事之后，他以问答形式提出了某些骇人听闻的陈述，涉及教会作者所称的\Quote{印记}；这些陈述并非源于信息不全；例如：\Quote{施印者被称为父，受印者被称为青年和儿子}；以及谁回答：\Quote{我已被生命树的白油膏所傅油}——这些事我们从未听说甚至发生在异端者中。其次，他甚至确定了那些交付印记者所提及的数目，即\Quote{七位天使，他们附着于将死之身的灵魂两侧；一方被称为光明天使，另一方被称为\InnerQuote{执政天使；}}；并且他断言，\Quote{那些被称为\InnerQuote{执政天使}之首领被称为\InnerQuote{可咒骂的神}}。然后，他抓住这个表述，并非无理地攻击那些敢于使用这类语言的人；因此，我们与那些谴责这类个体的人怀有相似的愤慨之情，如果确实存在有人称犹太人的天主——祂降下雨和雷霆，是这个世界的创造者，也是\Person{梅瑟}的天主及其所记载的宇宙起源论的天主——为\Quote{可咒骂的神}。然而，\Person{塞尔苏斯}使用这些表述时，似乎并非出于一个\Emph{理性的}目的，而是出于一种极不理性的目的，源于他对我们的仇恨，这如此不像一位哲学家。因为他的目的是，那些不熟悉我们习俗的人，在阅读他的论文后，会立即攻击我们，好像我们称这位高贵的创造者为\Quote{可咒骂的神}一般。在我看来，他确实像那些犹太人一样，当基督教开始被宣讲时，他们散布关于福音的虚假报告，例如：\Quote{基督徒献祭婴儿并吃其肉}；以及\Quote{自称基督徒的人，想要行\InnerQuote{黑暗的行为}，常熄灭（聚会中的）灯，然后每个男子与他遇到的任何女子发生性关系}。这些诽谤长期以来不合理地影响着许多人的思想，使那些与福音无关之人相信基督徒就是这样的人；甚至时至今日，它们仍误导一些人，阻止他们甚至与基督徒进行简单的谈话交流。\par

% structure: p0056 source=h2 target=subsection reason=relative-heading-rank; base=h2

\subsection{第二十八章}

\Person{塞尔苏斯}似乎就是怀着这样的目的，当他声称基督徒称创造者为\Quote{可咒骂的神}时；这样，那些相信他对我们的指控的人，如果可能，就会起来将基督徒作为最不敬神的人消灭。此外，他混淆了不同的事物，还陈述了为何\Person{梅瑟}宇宙起源论的天主被称为\Quote{可咒骂的神}的理由，断言\Quote{这就是他的品性，在那些这样看待他之人的眼中，他理当被咒骂，因为他宣布了对蛇的诅咒，而那蛇曾引导最初的人类认识善恶}。现在他本应知道，那些拥护蛇的事业的人，因为蛇给最初的人类提出了好的建议，并且远远超出神话中的提坦和巨人，因此被称为欧菲特派，他们离基督徒是如此之远，以至于他们对\Person{耶稣}的指控与\Person{塞尔苏斯}本人一样多；而且他们不允许任何人进入他们的集会，除非他首先咒骂\Person{耶稣}。那么，请看\Person{塞尔苏斯}的做法是多么不合理，他在反对基督徒的演说中，把那些甚至不愿听\Person{耶稣}的\Emph{名字}，或者否认\Person{耶稣}是智者和有德之人的人，描绘成这样的！那么，还有什么能显示出更大的愚蠢或疯狂呢，不仅是那些希望从蛇（作为善的创始者）得名的人，还有\Person{塞尔苏斯}，他认为对欧菲特派的指控也适用于基督徒！很久以前，确实有那位希腊哲学家，他偏爱贫穷的生活，并展现了幸福生活的典范，表明即使他一无所有，也不被排除在幸福之外，他称自己为犬儒主义者；而这些不敬神的可怜虫，作为不是人类（因为蛇是人类的敌人），而是作为蛇，他们以从蛇得名为欧菲特派而自豪，蛇是一种对人类最敌对、最令人恐惧的动物，并吹嘘有一个叫\Person{尤弗拉特斯}的人是这些亵渎观点的引入者。\par

% structure: p0058 source=h2 target=subsection reason=relative-heading-rank; base=h2

\subsection{第二十九章}

仿佛他是在诽谤基督徒，他继续控告那些称梅瑟及其律法的天主为\Quote{可咒骂的天主}的人；并且想象是基督徒那样说话，他这样说道：\Quote{还有什么比这种愚蠢的智慧更愚蠢或更疯狂的呢？犹太立法者究竟犯了什么错误？你为什么通过你所说的某种寓意和预表的解释方法接受他所给的宇宙起源论和犹太人的法律，而同时，哦，最不敬的人，你却不情愿地赞美那世界的创造主，他曾应许赐给他们一切，应许使他们的种族繁衍到地极，并以同样的血肉使他们从死者中复活，又曾启发他们的先知，而你却诽谤他！当这些想法迫使你时，你确实承认你敬拜同一位天主；但当你的导师耶稣和犹太的梅瑟作出矛盾的裁决时，你就寻求另一位天主来取代祂和圣父！}现在，通过这些陈述，这位显赫的哲学家刻尔苏斯明确地诽谤基督徒，断言当犹太人紧逼他们时，他们就承认与犹太人相同的天主；但当耶稣立法与梅瑟不同时，他们就寻求另一位天主来取代他。现在，无论我们与犹太人交谈，还是独自一人，我们知道只有一位同样的天主，就是犹太人古时所崇拜、至今仍宣认崇拜为天主的，我们对祂毫无不敬。然而，我们\Emph{不}断言天主会以同样的血肉使人从死者中复活，正如前几页所表明的；因为我们不主张那播种在腐朽、羞辱和软弱中的属血气的身体会像播种时那样复活。然而，关于这些主题，我们在前几页已经充分讨论过。\par

% structure: p0060 source=h2 target=subsection reason=relative-heading-rank; base=h2

\subsection{第三十章}

接着他回到七个执政邪魔的主题，他们的名字在基督徒中找不到，但我认为是被\OriginalTerm{敖非特派}{Ophites}所接受的。事实上，我们发现，在他们所依据的那个我们设法看到的图示中，相同的顺序正如刻尔苏斯所给出的那样。刻尔苏斯说\Quote{山羊被塑造得像狮子}，没有提及那些真正最不敬的人所赋予他的名字；然而我们发现，在圣经中被尊崇为造物主的天使的那位，被这个可咒骂的图示称为\OriginalTerm{似狮的弥额尔}{Michael the Lion-like}。再者，刻尔苏斯说\Quote{第二位是一头公牛}；而我们所拥有的图示说他是\OriginalTerm{似牛的苏里尔}{Suriel, the bull-like}。进一步，刻尔苏斯称第三位为\Quote{一种两栖动物，并能可怕地发出嘶嘶声}；而图示描述第三位为\OriginalTerm{似蛇的辣法耳}{Raphael, the serpent-like}。此外，刻尔苏斯断言第四位\Quote{有鹰的形状}；图示描述他为\OriginalTerm{似鹰的加俾额尔}{Gabriel, the eagle-like}。再者，据刻尔苏斯说，第五位\Quote{有熊的面容}；而根据图示，这是\OriginalTerm{似熊的陶塔巴奥特}{Thauthabaoth, the bear-like}。刻尔苏斯继续说道，第六位\Quote{被描述为有狗的脸}；图示称他为\OriginalTerm{厄辣塔奥特}{Erataoth}。他补充说，第七位\Quote{有驴的面容，并被命名为塔法巴奥特或奥诺厄尔}；然而我们发现，在图示中他被称为\OriginalTerm{奥诺厄尔}{Onoel}，或\OriginalTerm{塔塔拉奥特}{Thartharaoth}，外表有些像驴。我们认为精确陈述这些事是适当的，以免我们显得对刻尔苏斯自称知道的事情一无所知，而是我们基督徒比他知道得更清楚，可以证明这些不是基督徒的话，而是那些完全与救恩疏远、既不承认耶稣为救主、天主、导师，也不承认为天主之子的人的话。\par

% structure: p0062 source=h2 target=subsection reason=relative-heading-rank; base=h2

\subsection{第三十一章}

此外，若有人希望了解那些巫师的伎俩，他们试图借此以教导（仿佛拥有某些秘仪的知识）引诱人离开，却毫无成效，那么让他听听他们在经过所谓的\Quote{邪恶的栅栏}后的训示——这些栅门受制于统治灵体的世界。（以下是他们行事的方式）：\Quote{我向独形之王致敬，盲目之锁，完全遗忘，第一权能，由眷顾之灵和智慧所保存，我从他那里被纯然地派遣出来，已经分享了子与父的光明：愿恩宠与我同在；是的，父啊，愿恩宠与我同在。}他们还说，\OriginalTerm{奥格多}{Ogdoad}的起源由此而来。接下来，他们被告知在通过所谓的\OriginalTerm{阿尔达巴奥}{Ialdabaoth}时要说如下的话：\Quote{你，第一与第七位，生来充满自信地发令，你，阿尔达巴奥啊，你是纯净心灵的理性统治者，对子与父的完美工程，以类型的特征承载生命的象征，并为世界打开你曾为你的王国关闭的门，我再度自由地穿越你的领域。愿恩宠与我同在；是的，父啊，愿恩宠与我同在。}他们还说，名为\OriginalTerm{费农}{Phænon}的星与狮子状的统治者有感应。他们接着想象，那位已通过阿尔达巴奥并抵达\OriginalTerm{雅奥}{Iao}的人，应当如此说：\Quote{你，第二位雅奥，你在夜间闪耀，你是子与父隐秘奥秘的统治者，死亡的首位王子，无辜者的份，现在携带着我自己的胡须作为标志，我已准备好穿越你的领域，藉着活生生的话语坚固了那由你所生者。愿恩宠与我同在；父啊，愿恩宠与我同在。}他们接着来到\OriginalTerm{撒巴奥特}{Sabaoth}，他们认为应当向他说如下的话：\Quote{啊，第五领域的统治者，大能的撒巴奥特，你维护你受造物的法律，他们藉着你的恩宠通过更强有力的\OriginalTerm{五元组}{Pentad}的帮助而得释放，请接纳我，看见他们技艺的无暇象征，由一形象的印记所保存，一个由五元组释放的躯体。愿恩宠与我同在，父啊，愿恩宠与我同在。}在撒巴奥特之后，他们来到\OriginalTerm{阿斯塔费}{Astaphæus}，他们相信应当向他说如下的祈祷：\Quote{啊，阿斯塔费，第三道门的统治者，第一水元素的监督者，请视我为你的入门者之一，接纳我已由童贞之灵净化，你看见世界的本质。愿恩宠与我同在，父啊，愿恩宠与我同在。}之后是\OriginalTerm{阿洛埃}{Aloæus}，应当如此向他说话：\Quote{啊，阿洛埃，第二道门的治理者，请让我通过，因我带给你你母亲的象征，一份被各领域权能所隐藏的恩宠。愿恩宠与我同在，父啊，愿恩宠与我同在。}最后他们提到\OriginalTerm{霍雷}{Horæus}，并认为应当向他献上如下的祈祷：\Quote{你无所畏惧地跃过火墙，啊，霍雷，你获得第一道门的治理权，请让我通过，因你看到你自己权能的象征，刻在生命树的形象上，并按照这肖像形成，就像无辜的样子。愿恩宠与我同在，父啊，愿恩宠与我同在。}\par

% structure: p0064 source=h2 target=subsection reason=relative-heading-rank; base=h2

\subsection{第三十二章}

所谓的克尔苏斯的博学，其实不过是由好奇的琐事和愚蠢的闲谈构成，却使我们触及这些话题，目的是向每一位阅读他的论文和我们的答复的人表明，我们对这些题材并不缺乏信息，他正是以此为由诽谤基督徒，而基督徒既不了解也不关心这类事情。因为我们也渴望学习并阐述这些事，以免巫师们以比我们知道得更多为借口，欺骗那些容易被名称的光彩所迷惑的人。我本可以给出更多的例证，表明我们了解这些欺骗者的观点，并且我们否认它们，因为它们与我们的观点相异，是不虔敬的，与真基督徒的教义不相谐，我们甚至愿意为这教义而死。还必须注意，那些编造这一系列虚构的人，既不懂魔法，也不会区分圣经的意义，把一切都混淆了；因为他们从魔法中借用了\OriginalTerm{阿尔达巴奥}{Ialdabaoth}、\OriginalTerm{阿斯塔费}{Astaphæus}和\OriginalTerm{霍雷}{Horæus}这些名字，又从希伯来圣经中借用了希伯来语所称的\OriginalTerm{雅奥}{Iao}或\OriginalTerm{雅赫}{Jah}、\OriginalTerm{撒巴奥特}{Sabaoth}、\OriginalTerm{阿多奈}{Adonæus}和\OriginalTerm{厄娄}{Eloæus}。现在，取自圣经的名字是同一天主的名号；天主的敌人——连他们自己也承认——不理解这一点，导致他们以为雅奥是一个不同的天主，撒巴奥特是另一位天主，阿多奈（圣经中称为阿多奈）又是第三位，而厄娄（先知们在希伯来语中称为厄洛伊）也是不同的。\par

% structure: p0066 source=h2 target=subsection reason=relative-heading-rank; base=h2

\subsection{第三十三章}

克尔苏斯接下来讲述其他的寓言，大意是\Quote{某些人回到了\OriginalTerm{执政者}{archontics}的形状，以致有的被称为狮子，有的公牛，有的龙，或鹰，或熊，或狗。}我们在我们拥有的图表（克尔苏斯称之为\Quote{方形图案}）中也发现了这些不幸之人关于乐园之门的陈述。火剑被描绘为火焰之圆的直径，仿佛把守着知识和生命之树。然而，克尔苏斯要么不愿，要么不能复述那些根据这些不敬之人的寓言所描绘的、通过每道门的人所说的话；但我们这样做，是为了向克尔苏斯和那些阅读他论文的人表明，我们知道这些亵渎奥秘的深度，并且它们远远不同于基督徒向天主献上的敬拜。\par

% structure: p0068 source=h2 target=subsection reason=relative-heading-rank; base=h2

\subsection{第三十四章}

在结束了上述内容以及我们自己所添加的类似论述之后，塞尔索继续如下说：\Quote{他们继续将一件又一件事情堆积在一起——先知们的言论、圆圈套着圆圈、出自一个地上教会和割损的流溢、从一位\OriginalTerm{普律尼科斯}{Prunicos}——一位童贞女和活灵魂——流出的一种能力；一个被杀以活的天堂，一个被刀剑杀戮的大地，许多被杀以便能活的人，世界的罪死去时，世界上的死亡就停息了；还有一条窄路，自行敞开的门。在他们所有著作中都提到生命树和通过那\InnerQuote{树}的肉身复活，因为，我想，他们的老师被钉在十字架上，且是一名木匠；所以如果他碰巧是从悬崖上被扔下，或被推进坑里，或被吊死，或是一个皮匠、石匠、铁匠，那么就会（编造出）一个超越诸天的生命悬崖，或一个复活的坑，或一根不朽的绳索，或一块有福的石头，或一个爱的铁器，或一块神圣的皮革！如今有哪个老妇人不羞于在耳语中说出这样的事，即使是编故事哄婴儿入睡？}在使用了这样的言语之后，我认为塞尔索把他不完全听到的事情混淆在了一起。因为看起来，即使他可能听到了一些话，可追溯至某种现存的异端，他也并未清楚地理解所要传达的意思；而是把那些话堆积在一起，想要向那些既不认识我们的观点也不认识异端观点的人表明，他通晓基督徒的所有教义。这一点从前文也可看出。\par

% structure: p0070 source=h2 target=subsection reason=relative-heading-rank; base=h2

\subsection{第三十五章}

我们确实习惯于使用先知们的话，他们证明耶稣就是他们所预言的基督，并从先知书指出福音中关于耶稣的事件已经应验。但是当塞尔索说到\Quote{圆圈套着圆圈}时，（他或许是从）前述异端借用（了这个说法），该异端在一个圆圈（他们称之为万有之灵魂和\OriginalTerm{里外雅堂}{Leviathan}）中包含了执政鬼魔的七个圆圈；或者也许是由于误解了训道者，当他说：\BibleQuote{风在圆圈中旋转，又返回其圆圈}\BibleRef{训 1:6}时。至于\Quote{出自一个地上教会和割损的流溢}，大概是因为地上的教会被某些人称为来自天上教会和更美好世界的流溢；并且法律中所描述的割损是在一个专门用于净化之处所进行的割损的象征。此外，\Person{瓦伦廷}的追随者按照他们的错谬体系，将某种智慧称为\OriginalTerm{普律尼科斯}{Prunicos}，他们认为那位患血漏十二年的妇女是这智慧的象征；所以混淆了希腊、外邦和异端各种各样意见的塞尔索，在听说了她之后，便断言有一种能力从一位童贞女普律尼科斯流出。再说\Quote{活灵魂}，也许瓦伦廷的一些追随者神秘地将其指向他们称为世界灵魂创造者的存在；或者，与\InnerQuote{死的}灵魂相对，有些人不无优雅地将\InnerQuote{活的}灵魂称为\Quote{那个得救的人}的灵魂。但我不知道什么\Quote{被杀的天堂}，或\Quote{被刀剑杀戮的大地}，或许多被杀以便能活的人；因为这些很可能是塞尔索自己脑中编造出来的。\par

% structure: p0072 source=h2 target=subsection reason=relative-heading-rank; base=h2

\subsection{第三十六章}

此外，我们要说，当世界的罪死去时，世界上的死亡就停息了，这是指着宗徒的神秘话语说的，其文如下：\Quote{\BibleQuote{当祂将一切仇敌放在祂脚下之后，最后被毁灭的仇敌便是死亡。}\BibleRef{格前 15:25-26}} \Quote{\BibleQuote{当这必朽坏的穿上不朽坏的时候，那所记载的话就应验了：\InnerQuote{死亡被得胜吞灭。}}\BibleRef{格前 15:54}} 再者，\Quote{狭窄的下坡路}，或许那些持守灵魂转世学说的人会把它指向那种观点。而所谓自行敞开的门，有些人隐约地将其指向以下经文：\Quote{\BibleQuote{给我敞开仁义的门，我要进去称谢上主；这是上主的门，义人要进入这门。}\BibleRef{咏 118:19-20}}以及\Quote{\BibleQuote{你从死门把我提起，为叫我在熙雍女子门口宣扬你的一切赞美。}\BibleRef{咏 9:14}} 圣经还将那些导致灭亡的罪称为\InnerQuote{死门}，相反，将善行称为\InnerQuote{熙雍门}。同样地，\InnerQuote{仁义的门}，即\InnerQuote{德行的门}的同义词，这些门向追求德行的人开启。至于\Quote{生命树}的主题，当我们在解释创世纪中关于天主所种植的乐园的叙述时，将更恰当地说明。此外，塞尔索屡次嘲笑复活的主题——一个他不理解的道理；此刻，他不满足于先前所说，还加上：\Quote{据说有通过树的肉身复活；}我想，他是不理解那象征性的说法：\Quote{因树而来死亡，因树而来生命，}因为在亚当里有死亡，在基督里有生命。接着他嘲笑那\InnerQuote{树}，从两方面攻击它，说：\Quote{因此这树被引进来，要么是因为我们的老师被钉在十字架上，要么是因为他是一名木匠；}他没有注意到生命树在梅瑟著作中已经提到，他也瞎了眼看不到这点：在教会流传的任何福音书中，耶稣本人从未被描述为木匠。\par

% structure: p0074 source=h2 target=subsection reason=relative-heading-rank; base=h2

\subsection{第三十七章}

再者，凯尔苏斯以为我们虚构这\Quote{生命树}，是为给十字架一个寓意解释；由于在此点上的错误，他接着说：\Quote{如果他曾被丢下悬崖、或被推入坑中、或被绞死，那么就会有人发明一个远在天上的生命悬崖、或复活的坑、或不朽的绳索。}又说：\Quote{如果\InnerQuote{生命树}是虚构的，因为他——耶稣——（据说）是个木匠，那么如果他是个皮匠，就会有人提到圣皮；如果他是个石匠，就会有人提到有福的石；如果他是个铁匠，就会有人提到爱之铁。}然而，有谁看不出他指控的浅薄呢？他如此诽谤那些他自称要使其脱离迷惑的人。说完这些之后，他又以与那些虚构了狮形、驴首、蛇形统治天使及其他类似荒谬故事的人完全一致的语调说话，但这并不影响属于教会的人。的确，即使是醉醺醺的老妪，也会羞于对婴儿吟唱或低语这些故事，以哄其入睡，就像那些虚构驴头生灵及所谓在各门口发表的演说的人所做的那样。但凯尔苏斯并不了解教会成员的道理，即使那些遵照耶稣的命令，用毕生精力研究圣经、比希腊哲学家追求所谓智慧时更努力钻研圣书意义的人，也只有极少数能理解这些道理。\par

% structure: p0076 source=h2 target=subsection reason=relative-heading-rank; base=h2

\subsection{第三十八章}

此外，我们高贵的（朋友）并不满足于他从那图式中引出的反对意见，为了加强对我们（与它毫无关系的人）的指控，他还想从同一（异端者）那里引出另一些指控，却仿佛出自不同来源。他的话是：\Quote{他们的奇事中有一件不小，因为在上层圆圈之间——即那些在天之上的——有一些铭文，他们加以解释，其中特别有两个词：\InnerQuote{一个较大的和一个较小的}，他们将其归于圣父和圣子。}在所提到的图式中，我们发现了较大和较小的圆圈，其直径上刻有\Quote{圣父和圣子}；在较大圆圈（其中包含较小圆圈）与另一个由两个圆圈（外圈为黄色，内圈为蓝色）组成的圆圈之间，有一个形似斧头的屏障，上面刻有铭文。在它上面，有一个紧贴前两个圆圈中较大者的短圆圈，刻有\Quote{爱}；稍低处，有一个触及同一圆圈的圆圈，刻有\Quote{生命}。在第二个圆圈上，它与另两个圆圈交织并包含它们，有一个菱形图案（称为\Quote{智慧的先见}）。在它们的交叉点之内是\Quote{智慧的本性}。在交叉点之上有一个圆圈，刻有\Quote{知识}；再低处有一个圆圈，刻有\Quote{理解}。我们在答复凯尔苏斯时引入这些细节，是为了向读者表明，我们比他更了解这些事，并非仅凭传闻，尽管我们也反对它们。此外，如果那些以这些事为傲的人也自称懂得某种巫术和妖术（按他们的意见，这是智慧的顶峰），我们则不对此作任何断言，因为我们从未发现过此类事物。然而，凯尔苏斯，他已被屡次证实作假见证和无理指控，现在当自省，是否在这些事上也有虚假，或者他是否从那些与我们的基督信仰毫无关系的人的著作中摘录并引入了这些陈述。\par

% structure: p0078 source=h2 target=subsection reason=relative-heading-rank; base=h2

\subsection{第三十九章}

其次，谈到那些运用巫术和邪术的人，他们呼号野蛮的魔鬼名称，塞尔苏斯评论说，这样的人就像那些在相同事物上，在不知道希腊人所说的魔鬼名称与斯基泰人不同的众人面前表演奇迹的人。接着他引用希罗多德的一段话，说：\Quote{阿波罗在斯基泰人中被称为贡戈叙罗斯；波塞冬被称为塔吉马萨达；阿佛洛狄忒被称为阿尔金帕桑；赫斯提亚被称为塔比提。}现在，凡有能力者可以探究，在这些事情上，塞尔苏斯和希罗多德是否都错了；因为斯基泰人与希腊人对于那些被认为是神明者的事物，理解并不相同。因为怎能相信阿波罗在斯基泰人中被称为贡戈叙罗斯呢？我不认为贡戈叙罗斯转译为希腊语时，能得出与阿波罗相同的词源；或者阿波罗在斯基泰方言中具有贡戈叙罗斯的意义。关于其他名称，至今也未曾有过这样的断言，因为希腊人是根据不同的环境和词源，将他们认为神明者所拥有的名称赋予他们；而斯基泰人又是根据另一套环境；波斯人、印度人、埃塞俄比亚人、利比亚人，或者那些喜欢（凭想象）命名而不持守万有创造者那公义而纯洁之理念的人，也是如此。然而，我们在前面几页中已经说了足够多的话，当时我们想要证明撒巴奥特和宙斯不是同一位神祇，并且我们也根据圣经对不同的方言作了一些论述。因此，我们乐于略过这些点，因为塞尔苏斯会使我们重复。接着，他又将属于巫术和邪术的事情混在一起，也许不针对任何人——因为实际上并没有人假借这种崇拜行巫术——但也许他心中想到了某些人在单纯者面前\Emph{确实}运用这些做法（使他们显得像是因神圣能力而行），于是补充说：\Quote{何必一一列举那些教导净化方法、赎罪颂歌、驱邪咒语、造像或模拟魔鬼，以及衣服、数字、宝石、植物、根或一般各种事物中的各种解毒剂的人呢？}关于这些事情，理性并不要求我们作任何辩护，因为我们丝毫不受这类怀疑的牵连。\par

% structure: p0080 source=h2 target=subsection reason=relative-heading-rank; base=h2

\subsection{第四十章}

此后，塞尔苏斯在我看来，如同那些极端仇恨基督徒的人，在对基督教信仰全然无知的人面前，坚称他们确已查明基督徒吞噬婴儿的肉，并与他们的女人毫无节制地行淫。这些说法已被斥为针对基督徒捏造的谎言，并且被大众和完全不属于我们信仰的人承认；因此，塞尔苏斯下面的话也被发现是针对基督徒的诽谤，他说：\Quote{他在我们信仰的某些司铎手中见过野蛮的书籍，载有魔鬼的名字和奇迹般的作为；}他还断言：\Quote{我们信仰的这些司铎自称不做任何善事，只做一切旨在伤害人的事。}但愿塞尔苏斯反对基督徒的一切言论都能被大众驳斥，因为大众凭经验已确知这些事不属实，多数人与基督徒比邻而居，甚至连这些所谓做法的存在都未曾听闻！\par

% structure: p0082 source=h2 target=subsection reason=relative-heading-rank; base=h2

\subsection{第四十一章}

其次，仿佛他忘了自己的任务是反对基督徒，他说：\Quote{他结识了一位埃及乐师狄奥尼修斯，后者告诉他，关于巫术，它们只对未受教育的人和道德败坏的人有力量，而对哲学家则不能产生任何效果，因为哲学家谨慎地持守健康的生活方式。}如果我们的目的是讨论巫术，我们本可以在已说过的之外再补充几句；但既然我们只须回应塞尔苏斯最重要的论点，我们就关于巫术说：任何愿意探究哲学家是否曾被巫术俘虏的人，可以阅读莫伊拉格尼斯所写关于提亚纳的术士兼哲学家阿波罗尼乌斯《回忆录》的内容，其中这位非基督徒而是哲学家的人断言，一些颇具名望的哲学家被阿波罗尼乌斯的巫术力量所吸引，并作为术士投靠他；其中，我想他特别提到了尤弗拉特斯和一位伊壁鸠鲁派信徒。而\Emph{我们}则宣称并凭经验学到，凡按照来自耶稣的基督教敬拜万有者之天主，并按照祂的福音生活，昼夜持续合宜地使用规定的祈祷的人，不会被巫术或魔鬼掳走。因为确实\BibleQuote{主的使者在敬畏祂的人周围扎营，并拯救他们脱离一切邪恶}；并且教会中那些小子们的天使，被派来守护他们，据说总是看见他们在天之父的面容，无论\Quote{面容}或\Quote{看见}是什么意思。\par

% structure: p0084 source=h2 target=subsection reason=relative-heading-rank; base=h2

\subsection{第四十二章}

\Person{策尔苏斯}从另一个方面向我们提出如下指控：\Quote{某些极其不虔敬的错误，}他说，\Quote{是由他们因极端无知而犯下的，他们由此背离了神圣隐语的含义，制造了一个与天主对立的仇敌——魔鬼，并用希伯来语称他为\OriginalTerm{撒殚}{Satan}。诚然，这些说法纯属凡人的捏造，甚至不宜复述，即：全能的天主，渴望赐福于人，却有一个与他作对者，且无能为力。天主之子因而被魔鬼击败；受其惩罚，又教导我们也藐视他所施加的惩罚，预先告诉我们，撒殚在像他自己那样显现于人后，将行伟大惊异之事，为自己索取天主的荣耀，但那些希望远离他的人不应关注撒殚的这些作为，而应唯独信靠他。这些说法显然是一个欺骗者的话，他在策划并施计反对那些与他观点不同、并结成阵线对抗他们的人。}接着，为指出这些\Quote{隐语}（我们对其的误解导致引入关于撒殚的观点），他继续说：\Quote{古人隐约提到众神之间的一场战争，\Person{赫拉克利特}这样论及：\InnerQuote{若必须说有普遍的战争与不和，而万物皆在争斗中运行与治理。}\Person{斐瑞居德斯}，远较\Person{赫拉克利特}古老，讲述了一个神话，说一支军队与另一支敌对列阵，称\Person{克罗诺斯}为一方的首领，\Person{俄菲翁}为另一方，并叙述了他们的挑战与斗争，提到他们之间达成了协议，规定哪一方落入海洋即为战败，而逐出并征服他们的一方则拥有天堂。关于\Person{提坦}与\Person{巨人}的秘仪也有某种类似的（象征）意义，埃及关于\Person{提丰}、\Person{荷鲁斯}与\Person{奥西里斯}的秘仪也是如此。}在作了这些陈述之后，并未克服这些记述如何包含更高观点、而我们的记述是它们错误摹本这一困难，他继续辱骂我们，评论道：\Quote{这些不像那些关于魔鬼、恶神的故事，或者，如他更真实地评论，一个想建立对立教义的欺骗者。}同样，他理解\Person{荷马}，仿佛\Person{荷马}也隐约指涉类似\Person{赫拉克利特}、\Person{斐瑞居德斯}以及提坦与巨人秘仪的创始人所提及之事，即\Person{赫淮斯托斯}对\Person{赫拉}说的话如下：\par

曾因你的缘故，我领教了他无比的力量，\newline{}从以太高空头朝下被抛落。\par

以及宙斯对赫拉说的话：\par

你难道忘了，当被捆绑高悬，\newline{}从缀满星辰的穹苍广阔表面，\newline{}我将你颤抖地悬在黄金锁链，\newline{}所有狂怒的众神徒然反抗？\newline{}我头朝下将他们从奥林帕斯殿堂抛下，\newline{}在旋转中眩晕，坠落中窒息。\par

此外，解释\Person{荷马}的话，他补充道：\Quote{宙斯对赫拉说的话是天主对物质说的话；而对物质说的话隐约表明，起初与（天主）不和的物质，被天主取来，束缚并安排于规律之下，这些规律可类比为锁链；并且，作为对在其中制造混乱的恶神的惩罚，他将它们头朝下抛到这下界。}他声称，\Person{荷马}的这些话被\Person{斐瑞居德斯}如此理解，当时他说那区域之下是\Place{塔塔罗斯}区域，由\Person{玻瑞阿斯}之女\Person{哈耳庇厄}与\Person{风暴}守护，\Person{宙斯}将任何变得无序的神流放到那里。与此思想密切相关的还有\Person{雅典娜}的\OriginalTerm{佩普洛斯}{peplos}，它在\WorkTitle{泛雅典娜节}游行中被众人目睹。因为由此明显可见，他继续说，一个无母而纯洁的恶神掌控着\Person{巨人}的狂妄。此外，在接受希腊人虚构的同时，他继续堆砌对我们的如下指控，即：\Quote{天主之子被魔鬼惩罚，并教导我们，当我们受他惩罚时，也应当忍受。现在这些说法全然荒谬。因为我认为，倒是魔鬼应当受罚，而受他诬蔑的那些人不应当受到惩罚的威胁。}\par

% structure: p0090 source=h2 target=subsection reason=relative-heading-rank; base=h2

\subsection{第四十三章}

请看，那控告我们犯有最不虔敬的错误、并说我们偏离了神圣谜语（真义）的人，他自己岂不是显然错误，因为他没有注意到在梅瑟的著作中——那些著作不仅远早于赫拉克利图斯和斐瑞居德，甚至早于荷马——已经提到这个邪恶者及其从天上的堕落。因为那蛇——斐瑞居德所说的奥菲昂即源于它——成了人被逐出天主乐园的原因，以隐秘的方式影射了类似的事情，它借着应许神性和更大的祝福欺骗了女人；据说男人也追随了她的榜样。此外，在梅瑟的《出谷纪》中所提到的毁灭天使，除了那给服从他的人带来毁灭、且不抵挡他的恶行、也不与之抗争者，还能是谁呢？再者，在《肋未纪》中被放逐（到旷野）、在希伯来语中被称为阿匝则耳的那只羊，正是这个；必须把它放到旷野，当作赎罪祭品，因为阄落在了它身上。因为所有属于\Quote{较差}部分的人，因他们的邪恶而与作为天主产业的人对立，都被天主所遗弃。不仅如此，在《民长纪》中，贝里雅耳的子孙，他们被称为谁的子孙，若不是因他们的邪恶而成为他的子孙呢？除了所有这些例子之外，在比梅瑟本人还要古老的《约伯传》中，魔鬼被明确描述为出现在天主面前，要求对抗约伯的权柄，以便使他陷入最痛苦的试探；第一次试探包括失去所有财产和儿女，第二次则是用所谓的象皮病折磨约伯的全身。我略过从福音书中关于诱惑救主的魔鬼的引文，免得我为了答复责尔斯而引用较近期的著作。在《约伯传》的最后（一章）中，主在暴风和云彩中向约伯所说的话，记载在以其命名的书中，其中不乏关于蛇的内容。我还没有提及《厄则克耳书》中的段落，在那里他仿佛是在说法郎、或拿步高、或提洛王；也没有提及《依撒意亚书》中为巴比伦王所作的哀歌，从中可以学到不少关于邪恶的事，关于它的起源和产生的性质，以及它如何来自某些失去翅膀、跟随了那第一个失去自己翅膀者的存在。\par

% structure: p0092 source=h2 target=subsection reason=relative-heading-rank; base=h2

\subsection{第四十四章}

因为偶然或分享而来的善，不可能与那本性的善相同；然而，前者绝不会从那——可以说——为了保全自己而分享\Quote{生活的食粮}的人那里失去。但如果任何人失去了它，那必须是出于他自己的过错，因为他懈怠于分享这\Quote{生活的食粮}和\Quote{真实的饮料}，借着它们，翅膀得以滋养和浇灌，适合它们的目的，甚至正如最智慧的撒罗满关于真正富人所说的：\Quote{他为自己造了翅膀像鹰一样，回到他保护者的家中。}因为天主知道如何善用那些在邪恶中背弃祂的人，祂应当将这种邪恶放置在宇宙的某个部分，并设立一个德行的训练场，在那里那些渴望以\Quote{合法方式}恢复（所失去的）产业的人必须操练自己；为使他们因这些邪恶被考验，如同金子被火炼一样，并且竭尽全力防止任何卑劣之事损害他们的理性本性，他们才被证明配得上上升到神圣的事物，并被圣言提升到超越万物的福乐，可以说，达到善的顶峰。现在，那在希伯来语中被称为撒殚，被某些人称为撒殚纳斯——更符合希腊语的风格——译成希腊语时意为\Quote{对头}。但每一个偏爱恶习和邪恶生活的人，都是（因为行为与德行相悖）撒殚纳斯，即圣子的\Quote{对头}，而圣子就是公义、真理和智慧。然而，那首先在过着和平幸福生活的众生中失去翅膀、从福乐中堕落者，更适宜被称为\Quote{对头}；他按照厄则克耳所说，在他所有的道路上无瑕地行走，\Quote{直到在他身上发现了不义}，作为天主的乐园中\Quote{相似之印}和\Quote{华美冠冕}，充满了美善，却堕入毁灭，正如以神秘意义对他说的话：\Quote{你已堕落灭亡，将不再存留永远。}我有些冒昧地做了这些简短的评论，尽管这样做没有给这篇论文增加什么重要内容。然而，如果有谁有闲暇考查圣经，从各处收集关于邪恶的起源及其消解方式的论述，并将它们形成一个完整的道理体系，他就会看到梅瑟和先知们关于撒殚的看法，根本没有被责尔斯或任何一个被这个邪恶的魔鬼拉下去、从天主、从天主正确的看法、从天主圣言那里撕开的人梦见过。\par

% structure: p0094 source=h2 target=subsection reason=relative-heading-rank; base=h2

\subsection{第四十五章}

但既然\Person{策尔苏斯}拒绝接受关于所谓假基督的声明，既未读过他在\BibleBook{达尼尔}{书}中的论述，也未在\WorkTitle{保禄著作}中读过，亦未在\WorkTitle{福音书}中读过救主关于他降临的预言，我们就此主题也必须略作评论；因为\Quote{就如面孔不像面孔，}同样人的\Quote{心}也不彼此相似。因此，毫无疑问，人的心中将存在差异——那些倾向于德行的人并非以相同或类似的尺度被塑造和塑造向它；而另一些人因忽视德行，奔向相反极端。后者之中，有些人心灵根深蒂固地植根于邪恶，另一些人则较浅。那么，主张在人类中存在，可以说，两个极端——一端是德行，另一端是其反面——又有什么荒谬呢？这样，德行的完美就存在于那实现了在\Person{耶稣}身上给予之理想的人身上，从他那里向人类流出了如此巨大的归化、治愈与改善；而相反极端则在体现那被称为假基督之观念的人身上。因为天主以祂的预知理解一切，并预知这两者将产生的后果，愿意通过祂的先知向人类揭示这些，使那些理解他们言语的人熟悉善，并对善的反面保持警惕。此外，这两个极端中较好的一端，因祂的卓越而应被称为天主子；而相反的另一端则应称为邪恶魔鬼、\OriginalTerm{撒殚}{Satan}和魔鬼之子。其次，既然邪恶特别以扩散为特征，并当它模拟善的外表时达到其顶峰，因此罪恶伴随着征兆、奇事和虚假奇迹，通过其父魔鬼的合作而实现。因为魔鬼自己为欺骗人类所提供的帮助远远超过这些魔鬼给予变戏法者（为最卑劣目的欺骗人）的帮助。\Person{保禄}确实论及那被称为假基督者，虽有所保留，仍描述了其降临人类的方式、时间与原因。请注意他关于此事的言辞是否极为恰当，决不应受丝毫嘲笑。\par

% structure: p0096 source=h2 target=subsection reason=relative-heading-rank; base=h2

\subsection{第四十六章}

宗徒如此表述：\BibleQuote{我们恳求你们，弟兄们，因我们的主耶稣基督的来临，并因我们聚集到祂面前，使你们不致很快动摇思想，或受困扰，无论因言语、或因神谕、或因貌似来自我们的书信，以为主的日子已近。任何人都不应以任何方式欺骗你们：因为\Emph{那日子不会来临}，除非先有背叛，并显出那罪恶之人，丧亡之子；他反对并高举自己超过一切被称为神或受人敬拜者，甚至坐在天主的殿内，显示自己就是天主。你们不记得，我尚与你们同在时，已告诉你们这些事吗？现在你们知道那阻挡的，为使他到了时候才显露出来。因为罪恶的奥秘已在进行；只是那现今阻挡者将一直阻挡，直到他被除去。那时那恶者将要显露，主必以祂口中的气毁灭他，以祂来临的光耀除掉他：\Emph{他就是那按撒殚的运作而来的人}，带着一切能力、征兆和虚假奇迹，以及一切不义的欺诈，在那些丧亡的人身上；因为他们没有接受爱真理的心，以致得救。为此，天主给他们派遣了强大的迷惑，使他们相信谎言，使一切不信真理、反喜欢不义的人都被定罪。}解释此处所提及的每个细节不属于我们当前的目的。关于假基督的预言也在\BibleBook{达尼尔}{书}中陈述，它能使明智而坦诚的读者钦佩这些话为真正神圣和预言的；因为其中提到了与将要来临的国度相关的事，从达尼尔的时代开始，一直延续到世界的毁灭。凡愿意者皆可阅读。然而，请观察关于假基督的预言是否如下：\BibleQuote{在他们国度的后期，当他们的罪恶满盈时，必兴起一个面貌凶悍、通晓谜语的王。他的势力非常强大，他必奇妙地毁灭，顺利而行，并实行；他必毁灭勇士和圣民。他的锁链之轭必顺利，他手中有诡计，他心中自大，必用诡计毁灭多人；他必起来毁灭多人，并像用手压碎鸡蛋一样粉碎他们。}\Person{保禄}在他所引用的言辞中所说的，他说：\BibleQuote{因此他坐在天主的殿内，显示自己就是天主}，在\BibleBook{达尼尔}{书}中以如下方式提及：\BibleQuote{在圣殿上必有荒凉的可憎之物，在时期终末，荒凉将被终结。}如此众多（摘自更多经文）的段落，我认为应当引述，使听者能稍微理解圣经的意义，当它给予我们关于魔鬼和假基督的信息时；并对此目的所引用的感到满足后，让我们看看\Person{策尔苏斯}的另一项指控，并尽我们所能回答它。\par

% structure: p0098 source=h2 target=subsection reason=relative-heading-rank; base=h2

\subsection{第四十七章}

\Person{瑟尔苏}在所述之后接着说：\Quote{我能说明这事的原委，即他们何以称他为\InnerQuote{天主子}。古人将这个世界视为由天主所生，称之为祂的孩子和祂的儿子。于是，两者都称为\InnerQuote{天主子}，彼此极其相似。}因此他认为我们使用了\Quote{天主子}这一说法，曲解了关于世界由天主所生、是祂的\Quote{儿子}和\Quote{神}的论述。因为他未能如此考量\Person{梅瑟}和先知的时代，以致看到犹太先知普遍预言有一位\Quote{天主子}远在希腊人和\Person{瑟尔苏}所说的那些古人之前。他甚至不愿引用\Person{柏拉图}书信中的段落——我们前文曾提及——关于那位如此美妙安排这个世界者，承认祂是天主子；唯恐他自己也被\Person{柏拉图}（他常尊敬地提及）迫使承认这个世界的建造者是天主子，而祂的父是第一位神和万物的至高统治者。如果我们主张耶稣的灵魂通过与他最高度的合一而与如此伟大的\OriginalTerm{天主子}{Son of God}成为一体，不再与他分离，这一点也不奇怪。因为圣经的神圣语言也知晓其他事物，这些事物虽然按其本性是\Quote{二}，但彼此之间被认为是\Quote{一}，也确实是\Quote{一}。关于丈夫和妻子，经上说：\BibleQuote{他们不再是两个，而是一体了}；关于完人和那与真实的主、\OriginalTerm{圣言}{Very Word}、智慧与真理结合的人，经上说：\BibleQuote{那与主结合的，成为一灵}。如果那\BibleQuote{与主结合的，成为一灵}的人，谁曾与主——圣言、智慧、真理和正义——结合得比耶稣的灵魂更亲密，或至少接近那程度呢？如果这样，那么耶稣的灵魂和天主圣言——一切受造物的\OriginalTerm{首生者}{first-born of every creature}——就不再是二，而是一了。\par

% structure: p0100 source=h2 target=subsection reason=relative-heading-rank; base=h2

\subsection{第四十八章}

其次，当廊下派的哲学家断言天主和人的德行相同，并主张那掌管万有的天主并不比\Emph{他们的}智者更幸福，而是两者的幸福相等时，\Person{瑟尔苏}既不嘲笑也不嘲弄他们的见解。然而，如果圣经说完人通过德行与\OriginalTerm{圣言}{Very Word}结合，成为一体，以致我们推论耶稣的灵魂不与一切受造物的\OriginalTerm{首生者}{first-born of all creation}分离，他却嘲笑耶稣被称为\Quote{天主子}，没有留意圣经以隐秘和奥义的方式关于祂所说的。但为使那些愿意接受我们教义所推出的结论并由此获益的人能接受我们的观点，我们说圣经宣告由\OriginalTerm{天主子}{Son of God}所赋予生命的基督的身体，就是天主的整个教会，而这身体的各部分——被视为整体——由信徒组成；因为正如灵魂赋活并驱动身体——身体本身并没有像生物那样的自然运动能力——同样，\OriginalTerm{圣言}{Word}唤醒并推动整个身体，即教会，行适合之事，此外还唤醒属于教会的每个个别肢体，使他们离开圣言就什么也不做。既然这一切都经由不可轻视的推理链得出，那么主张正如耶稣的\Emph{灵魂}以完美和不可思议的方式与圣言结合，所以一般而言，耶稣这个人并不与那\OriginalTerm{独生者}{only-begotten}、一切受造物的首生者分离，也与祂不是一个不同的存有，这有何困难呢？但此话题已足够。\par

% structure: p0102 source=h2 target=subsection reason=relative-heading-rank; base=h2

\subsection{第四十九章}

让我们注意接下来的内容，凯尔苏斯用一句话表达了他对梅瑟宇宙起源论的看法，却未提供任何支持论据，就批评它，说道：\Quote{此外，他们的宇宙起源论极其愚蠢。} 现在，如果他提出了一些令人信服的证据来说明其愚蠢，我们会努力回答它们；但在我看来，仅仅因为他\Emph{断言}它不是\Quote{愚蠢的}，我就需要加以证明，这似乎不合理。然而，如果有人想知道我们接受梅瑟记述的理由以及可以为之辩护的论证，他可以阅读我们关于创世纪从书开头到\BibleQuote{这是人类的族谱书}这段的注释，在那里我们试图从圣经本身展示：起初创造的\Quote{天}是什么；\Quote{地}、\Quote{地的看不见的部分}、\Quote{无形}的是什么；\Quote{深渊}是什么，围绕它的\Quote{黑暗}；\Quote{水}是什么，运行在\Quote{水面上}的\Quote{天主的神}是什么；所造的\Quote{光}是什么，与起初创造的\Quote{天}有别的\Quote{穹苍}是什么；以及其他后续主题。凯尔苏斯还表达了他的观点，认为关于人类创造的叙述\Quote{极其愚蠢}，却没有陈述任何证明，也没有试图回答我们的论证；因为在我看来，他没有证据能够推翻\BibleQuote{人是按天主的肖像造的}这一声明。他甚至不明白天主所种植的\Quote{乐园}的意义，以及人最初在其中度过的生活；以及由于人的罪而被驱逐、安置在乐园对面之后所发生的事情。现在，既然他断言这些是愚蠢的说法，让他不仅关注每一个（一般）说法，尤其要注意这一句：\BibleQuote{祂安置了革鲁宾和四面转动发火焰的剑，把守生命树的道路}，并看看梅瑟写这些话是否没有严肃目的，而是像旧喜剧作家那样戏谑地叙述\Quote{\Person{Prœtus}杀了\Person{Bellerophon}}和\Quote{\Person{Pegasus}来自\Place{Arcadia}}。他们的目的是用这样的故事引人发笑；然而，不可能相信，一个为整个民族留下法律、并希望他的臣民相信这些法律是由天主赐予的人，会写出这样无关紧要的话，毫无意义地说了\BibleQuote{祂安置了革鲁宾和四面转动发火焰的剑，把守生命树的道路}，或者关于人类创造的其他任何陈述，而这些正是希伯来贤士们哲学研究的主题。\par

% structure: p0104 source=h2 target=subsection reason=relative-heading-rank; base=h2

\subsection{第五十章}

接下来，凯尔苏斯只是将一些古人关于世界和人类起源的各种相反意见堆砌在一起，然后断言：\Quote{梅瑟和先知们，他们给我们留下了书卷，却完全不知道世界和人的本性，编织了一张纯粹胡说八道的网。} 如果他现在展示\Emph{如何}在他看来圣经包含\Quote{纯粹胡说八道}，我们会努力摧毁那些在他看来证明其胡说八道性质的论证；但在此，仿效他的榜样，我们也戏谑地提出我们的意见：凯尔苏斯完全不知道先知们的含义和语言的性质，写了一部包含\Quote{纯粹胡说八道}的作品，并自夸地给它起了个标题\WorkTitle{真道}。由于他关于\Quote{创造的日子}的陈述——好像他清楚而正确地理解它们，其中一些发生在光、天、太阳、月亮和星辰的创造\Emph{之前}，一些发生在这些创造\Emph{之后}——成了他指控的根据，我们只作如下评论：梅瑟当时一定忘记了他刚才说过\BibleQuote{在六日内世界被创造完成}，并且由于这一疏忽，他在这句话后面加上了以下内容：\BibleQuote{这是人类的受造之书，在天主创造天地的那一天！} 然而，这完全不可信，在他说了关于六天的话之后，梅瑟会立即毫无特殊意义地加上\BibleQuote{在天主创造天地的那一天}；如果有人以为这些话可以指\BibleQuote{在起初天主创造了天地}，请留意，在\BibleQuote{有光，就有了光}和\BibleQuote{天主称光为昼}这些字句之前，已经说明了\BibleQuote{在起初天主创造了天地}。\par

% structure: p0106 source=h2 target=subsection reason=relative-heading-rank; base=h2

\subsection{第五十一章}

然而，在目前的情况下，我们并无意探讨理智与可感之物的问题，也不探究两类受造物如何被分配不同的日子，亦不深究与此主题相关的细节——因为若要阐释梅瑟的宇宙起源论，我们需要整部论著；而在开始回答塞耳苏斯之前，我们早已尽其所能地完成了这项工作，即在我们当时所拥有的能力范围内，讨论了梅瑟关于六日创造宇宙起源的问题。但我们必须牢记，圣言曾通过依撒意亚的口应许义人：日子将到来，那时太阳不再是他们的光，而是主自己要作他们永远的光，天主将是他们的荣耀。我认为，正是由于误解了某种有害的异端——该异端错误地解释\Quote{有光}这句话，仿佛那仅仅是造物主的一种\Emph{愿望}——塞耳苏斯才发表了这样的评论：\Quote{造物主并未从天上借光，不像那些从邻人的灯上点燃自己灯盏的人。}此外，由于误解了另一种不虔敬的异端，他又说道：\Quote{如果确实存在一个与伟大的天主敌对的咒诅之神，做了这违背天主意愿的事，那么天主为何要借给他光呢？}我们远非为这类幼稚言论辩护——相反，我们明确谴责这些异端者的论断是错误的，并着手反驳他们的观点——不是像塞耳苏斯那样反驳那些我们\Emph{不了解}的观点，而是反驳那些我们已准确认识的观点，这些认识部分来自他们的追随者的陈述，部分来自仔细阅读他们的著作。\par

% structure: p0108 source=h2 target=subsection reason=relative-heading-rank; base=h2

\subsection{第五十二章}

塞耳苏斯接着说道：\Quote{关于世界的起源与毁灭——无论它被视为无始无终、有始无终、还是相反——我目前暂不置评。}因此，我们也对这些点保持沉默，因为当前的工作并不需要。我们也不宣称，普遍天主的圣神（正如\Quote{天主的圣神运行在水面上}一语所暗示的）将自己混同于此世下界的事物，仿佛它们是异己之物；我们也不主张，某些针对祂圣神的邪恶阴谋——仿佛是由与伟大天主不同的另一位造物主所为，并且得到至上天主的容忍——需要被彻底挫败。因此，对那些说出此类荒谬言论的人，以及对未能巧妙驳斥这些言论的塞耳苏斯，我无话可说。因为他要么根本不应提及这些事，要么就应本着其所假装的仁爱性格，仔细阐述这些事，然后努力反驳这些不虔敬的主张。我们也从未听说，伟大的天主在将祂的神赐予造物主之后，又要求收回。接着，他愚蠢地攻击这些不虔敬的主张，问道：\Quote{有什么神会给予某物后又意图收回？因为要求收回（所给予之物）是匮乏之人的特征，而天主一无所缺。}对此，他仿佛说了什么巧妙的话攻击某些人，又补充道：\Quote{为何当他（借出祂的神）时，竟不知自己借给了一个邪恶的存在？}他又问道：\Quote{为何他无视一个与他目的相悖的邪恶造物主？}\par

% structure: p0110 source=h2 target=subsection reason=relative-heading-rank; base=h2

\subsection{第五十三章}

其次，他把各种异端混为一谈，没有注意到有些说法出自某一异端派别，有些则出自另一派别，就提出了我们对\Person{玛尔西昂}的异议。大概是从一些微不足道且无知的人那里听来的，他攻击了反驳这些说法的论点，但方式并不显得很明智。于是，他引用我们反驳\Person{玛尔西昂}的论点，却没有注意到他是在\Emph{反驳}\Person{玛尔西昂}，他问道：\Quote{他为什么暗中派遣，并毁灭他所创造的作品？他为什么暗中使用暴力、说服和欺骗？他为什么引诱那些据你们所说被他定罪或控告的人，像奴隶贩子一样把他们带走？他为什么教他们逃离他们的主？为什么逃离他们的父亲？为什么他违背父亲的旨意把他们据为己有？为什么他声称自己是陌生孩子的父亲？}对这些问题，他仿佛惊讶地补充了以下评论：\Quote{确实可敬啊，这位神渴望成为那些被另一位（神）定罪的罪人的父亲，也成为穷人的父亲，并且如他们自己所说，成为人类的渣滓的父亲，但他却无法抓住并惩罚从他那里逃脱的使者！}之后，仿佛是对那些承认这世界不是另一位陌生之神的作品的我们说话，他继续这样说道：\Quote{如果这些是他的作品，那么天主如何创造了恶？他为何不能说服和告诫（人）？他为何因人的忘恩负义和邪恶而懊悔？而且，他指责自己的手工作品，憎恨、威胁并毁灭自己的后代？他能把他们运出他自己创造的这个世界之外吗？}在我看来，这些话并没有清楚说明什么是\Quote{恶}；尽管希腊人中有许多派别对善恶本质看法不同，但他匆忙得出结论，仿佛从我们主张这世界也是宇宙的天主的作品这一观点，必然推出在我们看来\Emph{天主}是恶的创造者。然而，无论恶如何——是由天主创造的还是不是——它只是当你比较首要设计时得出的一个\Emph{结果}。我大为惊讶，如果他以为从我们主张这世界也是宇宙的天主的作品这一观点，可以推出天主是恶的创造者，那么这种推论是否也同样适用于他\Emph{自己}的陈述。因为有人可以对\Person{克尔苏斯}说：\Quote{如果这些是他的作品，那么天主如何创造了恶？他为何不能说服和告诫人？}的确，在推理中最大的错误就是指责持不同意见的人持有不健全的教义，而指责者本人在这方面却更应受到同样的指控。\par

% structure: p0112 source=h2 target=subsection reason=relative-heading-rank; base=h2

\subsection{第五十四章}

那么，让我们简要地看看圣经关于善与恶说了什么，以及我们应该如何回答这些问题：\Quote{天主如何创造了恶？}和\Quote{他怎么不能说服和告诫人？}根据圣经，严格来说，德行和善行是善的，严格来说，这些的反面是恶。我们在此满足于引用第三十四篇圣咏的一些经文，大意如下：\BibleQuote{寻求主的人不会缺少任何美物。孩子们，你们来听我；我要教导你们敬畏主。谁渴望生命，喜爱延年益寿，以得见美福？就要谨守口舌，不说恶言，唇舌不说欺诈的话。躲避罪恶，努力行善。}现在，\BibleQuote{躲避罪恶，努力行善}的训诫并不指有些人所说的\Emph{肉体}的恶或\Emph{肉体}的福乐，也不指外在事物，而是指\Emph{属灵}的福与恶；因为那躲避这类恶并实行这类善行的人，作为渴望真生命的人，将得享真生命；而作为喜爱看见\Quote{美善的日子}，在其中正义之言如太阳的人，他将看见这些日子，天主把他从这\Quote{邪恶的现世}和那些邪恶的日子中带走，关于这些日子，\Person{保禄}曾说：\BibleQuote{赎回时光，因为日子是邪恶的。}\par

% structure: p0114 source=h2 target=subsection reason=relative-heading-rank; base=h2

\subsection{第五十五章}

确实，可以发现一些段落，其中有形体和外在的（好处）被不恰当地称为\Quote{善，}——那些有助于自然生命的事物，而相反的事物则被称为\Quote{恶。}正是在这个意义上，\Person{约伯}对他的妻子说：\BibleQuote{我们从主手中领受了善，难道不也领受恶吗？}既然在圣经中，某一处经文里有这样的出自天主口中的陈述：\BibleQuote{我缔造和平，也创造恶；}又有另一处论及他说：\BibleQuote{恶从主那里降到耶路撒冷的城门，有战车和骑兵的喧嚣，}——这些经文曾使许多圣经读者困惑，他们不明白圣经所说的\Quote{善}和\Quote{恶}是什么意思——很可能\Person{塞尔苏斯}也因而困惑，便提出了\Quote{天主如何创造了恶？}这个问题；或者，也许他听人无知地讨论了与此相关的问题，便说出了我们注意到的这一说法。另一方面，我们主张，\Quote{恶}或\OriginalTerm{邪恶}{wickedness}以及由此而来的行为，\Emph{不是}由天主创造的。因为如果天主创造了\Emph{真正}恶的事物，那么关于（最后）审判的宣告又怎能自信地宣布呢？它告诉我们，恶人将按他们邪恶的多少受罚，而度圣善生活或行圣善行为的人，将享有福乐并得到天主的赏报。我很清楚，那些大胆断言这些恶是由天主创造的人，会引用圣经的某些表述（来支持他们），因为我们不能展示出一系列一致的经文；因为，尽管圣经（一般来说）责备恶人而赞许义人，但它仍包含一些说法，虽然数量相对较少，却似乎扰乱了蒙昧的圣经读者的心。然而，我认为在目前的论著中，不适于在此引用那些不协调的陈述，它们数量很多，而它们的解释需要大量的证明。因此，恶——如果指的是\Emph{真正}被称为恶的事物——\Emph{不是}由天主创造的；但有些恶，虽然相对于\Emph{整个}世界的秩序来说是\Emph{少数}，却\Emph{确实}出自他主要的工作，如同从木匠的主要工作中产生出螺旋刨花和锯屑，或者如同建筑师似乎成为他们建筑周围垃圾的原因，这些垃圾是石头和灰浆掉落的污物。\par

% structure: p0116 source=h2 target=subsection reason=relative-heading-rank; base=h2

\subsection{第五十六章}

然而，如果我们谈及所谓的\Quote{有形体的}和\Quote{外在的}恶——这些称呼不恰当——那么可以承认，\Emph{有些}时候其中一些恶是由天主创造出来的，为的是藉着它们使某些人得以悔改。这样的做法会有什么荒谬之处呢？因为，假如我们听到那些不恰当地被称为\Quote{恶}的苦难，是由父亲、导师和学监施加给他们照顾的人，或者由医生为了治愈而给病人动手术或烧灼，我们若说父亲虐待儿子，或学监和导师虐待学生，或医生虐待病人，那么操作者或责罚者都不会受到责备；同样地，如果天主为了悔改和治愈那些需要这种管教的人而将这样的恶降于人身上，那么这种观点就没有荒谬之处，也不会发生\BibleQuote{恶从主那里降到耶路撒冷的城门}——这些恶包括为了使他们悔改而由敌人加于以色列人的惩罚；也不会\BibleQuote{以棍棒惩罚那些离弃主律法的人，以鞭打责罚他们的过犯；}也不能说，\BibleQuote{你将有火炭放在他们身上；它们将成为你的帮助。}同样地，我们也解释了\BibleQuote{我，缔造和平，也创造恶}这句话；因为他使\Quote{有形体的}或\Quote{外在的}恶出现，用以净化并训练那些不受圣言和纯正道理管教的人。那么，这就是我们对\Quote{天主如何创造了恶？}这一问题的回答。\par

% structure: p0118 source=h2 target=subsection reason=relative-heading-rank; base=h2

\subsection{第五十七章}

关于\Quote{祂怎能不说服和劝诫人？}的问题，我们已经说过，如果这样的异议真的是一种指控的理由，那么，塞尔苏斯的异议就可以用来反对那些接受天主眷顾学说的人。任何人都可以回答这个指控，即天主是\Emph{无能}劝诫人的；因为祂借着整部圣经，并借着那些因天主恩宠的安排而成为听者导师的人，传达祂的劝诫。除非确实有某种特殊含义附着在\Quote{劝诫}一词上，好像它既意味着渗入被劝诫者的心智，又意味着使他听到指导者的话，这与该词的通常含义相反。对于\Quote{祂怎能无法说服人？}的异议——这也是可以针对所有相信天主眷顾的人提出的——我们必须提出以下看法。既然\Quote{被说服}属于那些被称为（可以这么说）\Quote{相互的}的词（比较短语\Quote{给人剃头}，当一个人努力使自己服从理发师时），因此，这不仅需要说服者的努力，还需要（可以这么说）对说服者的服从，即接受他所说的话。所以，不能说人未被说服是因为天主无法说服人，而是因为他们不愿接受天主的忠信之言。如果有人将这种说法应用于那些是\Quote{说服的工匠}的人，他也不会错；因为一个精通修辞学原理并恰当运用它们的人，有可能尽力说服，却似乎失败，因为他无法克服应当屈服于其说服技艺的人的意志。此外，说服并非来自天主，尽管说服之言可能由祂说出，这是保禄明确教导的，他说：\BibleQuote{这说服不是出于那召你们的人。}以下的话也表明了这种观点：\BibleQuote{你若愿意且顺从，就必吃地上的美物；你若拒绝且悖逆，刀剑必吞吃你。}因为若有人实际上渴望听劝诫者对他所说的话，并配得听闻天主的这些应许，就必须确保听者的意志和他对所说之事的倾向。因此，在我看来，在\WorkTitle{申命纪}中，以下的话是特别强调说出的：\BibleQuote{现今以色列啊，你的天主上主向你要求什么呢？岂不是敬畏你的天主上主，遵行祂的一切道，爱祂，遵守祂的诫命吗？}\par

% structure: p0120 source=h2 target=subsection reason=relative-heading-rank; base=h2

\subsection{第五十八章}

接下来要回答下面的质问：\Quote{祂怎么会在人变得忘恩负义和邪恶时后悔，并指责自己的手工作品，恨恶、威胁并毁灭自己的后代？}此时，塞尔苏斯诽谤并歪曲了\WorkTitle{创世纪}中所写的话，大意是：\BibleQuote{上主天主见人在地上的罪恶日益增多，人人心中终日图谋恶事，就后悔在地上造了人。天主心中思量说：我要将我造的人从地上除灭，从人到牲畜，到爬虫，到天空的飞鸟，都除灭，因为我后悔造了他们。}他引用了圣经中没有的话，好像它们传达了实际所写的意思。因为在那些话中并没有提到天主的后悔，也没有提到祂指责和恨恶自己的手工作品。如果天主似乎威胁要造成洪水的灾祸，从而在其中毁灭祂自己的儿女，我们必须回答说：既然人的灵魂是不朽的，那所谓的威胁旨在促使听者悔改，而洪水毁灭人则是大地的净化，正如某些颇有声望的希腊哲学家通过以下说法所表明的：\Quote{当诸神净化大地时。}至于将这些拟人化的短语归于天主，我们已经在前面几页中作过一些评论。\par

% structure: p0122 source=h2 target=subsection reason=relative-heading-rank; base=h2

\subsection{第五十九章}

其次，塞尔苏斯猜疑，或也许清楚看见，那些为洪水毁灭人类辩护者可能给予的答复，于是继续说：\Quote{但他若不毁灭自己的后代，又把他们迁移到他亲自创造的这个世界之外，到哪里去呢？}对此我们答复：天主决不会把那些因洪水而死亡的人从整个由天地构成的世界中移除，而是将他们从肉身的生活中移除，使他们脱离身体，同时摆脱地上的存在——在圣经许多地方，这通常被称为\Quote{世界}。尤其在\BibleBook{若望}{福音}中，我们经常发现地上的区域被称为\Quote{世界}，如经文：\BibleQuote{他是真光，照耀进入\InnerQuote{世界}的每个人；}又如：\BibleQuote{在世界上你们将有苦难，但放心，我已战胜了世界。}因此，如果我们将\Quote{移出世界}理解为从\Quote{地上区域}的转移，那么这一表述就毫无荒谬之处。反之，如果由天地构成的事物体系被称为\Quote{世界}，那么那些在洪水中灭亡的人决没有被移出所谓的\Quote{世界}。然而，如果我们注意经文\BibleQuote{我们并不注目那看得见的，而是注目那看不见的；}以及\BibleQuote{其实，自从天主创世以来，他那看不见的美善，即他永远的大能和神性，都可凭他所造的万物，被人辨认出来，}——我们可以说，那居住在\Quote{看不见的}事物，以及被称为\Quote{看不见者}的，已经离开了世界，圣言将他从此处移除，并迁移到天上区域，以便观赏一切美好之物。\par

% structure: p0124 source=h2 target=subsection reason=relative-heading-rank; base=h2

\subsection{第六十章}

但在考察他的这些断言之后，仿佛他的目的是要用许多话堆积其书，他用不同的语言重复了我们不久前已审查过的同样指控，说：\Quote{最愚蠢的事莫过于把世界的创造分配在几天之内，\Emph{在日子存在之前}：因为天尚未创造，地的基础尚未奠定，太阳尚未旋转，怎能有\Emph{日子}呢？}现在，这些话与下面的内容有何区别？\Quote{此外，从起初审视这些事，第一位和至大的天主发出命令：让这（第一件事）存在，然后这第二件事，然后这（第三件）；在第一天完成那么多之后，在第二天、第三天、第四天、第五天、第六天又做那么多，这不是荒谬吗？}我们尽可能地答复了这项反对天主\Quote{命令创造这第一、第二、第三件事}的异议，当时我们引用了经文：\BibleQuote{他说了，事就成就了；他命令了，万物就立定；}并指出那直接的创造者，可以说，世界的真正制造者是圣言，天主的圣子；而圣父，通过命令他自己的儿子——圣言——创造世界，是\Emph{首要}的创造者。至于第一天创造光，第二天创造穹苍，第三天将天下水聚在一处（大地由此生出那些仅受自然支配的果实），第四天创造（大）光体和星辰，第五天创造水生动物，第六天创造陆上动物和人，我们已在关于\BibleBook{}{创世纪}的注释中以及在前面几页中尽可能作了论述，当时我们责备那些人，他们按字面的\Emph{表面}意思，认为六天的时间用于创造世界，并引用经文：\BibleQuote{这是创造天地的来历，在上主天主创造大地和天空的那一日。}\par

% structure: p0126 source=h2 target=subsection reason=relative-heading-rank; base=h2

\subsection{第六十一章}

再者，由于他不明白这些话的意义：\Quote{到第六天，天主完成了祂所造的工作，就在第七天停止了祂所做的一切工作。天主降福了第七天，定为圣日，因为在这一天，祂停止了祂所开始的一切创造工作。}他以为\Quote{祂在第七天\Emph{停止了}}与\Quote{祂在第七天\Emph{安息了}}是同一回事，于是评论说：\Quote{此后，他确实疲倦了，像一个很糟糕的工匠，需要休息来恢复体力！}因为他完全不知道安息日和天主的安息，那是在世界受造完成后开始的，并持续整个世界的存续，在其中，凡在\Emph{自己的}六天内完成了\Emph{自己的}一切工作，并因没有忽略任何职责而得以升到（对天上事物的）默观和义人与有福者的集会中的人，都将与天主一同庆节。接着，好像圣经作了这样的陈述，或好像我们自己说天主因疲劳而安息，他继续说：\Quote{第一位天主感到疲劳、亲手工作或发出命令，是不合宜的。}赛尔苏说，第一位天主感到疲劳是不合宜的。而我们要说，天主圣言也不会感到疲劳，任何属于更美好、更神圣秩序的存在也不会，因为疲劳的感觉是那些在肉体内的人所特有的。你可以考察这一点是否适用于拥有任何类型身体的人，或拥有\Emph{属土的}身体、或比这稍好一点的身体的人。但\Quote{第一位天主亲手工作也不合宜。}如果你按\Emph{字面}理解\Quote{亲手工作}，那么这话既不适用于\Emph{第二位}天主，也不适用于任何分享神性的存在。但假如是在不当和比喻的意义上说的，以致我们可以把以下经文：\Quote{穹苍宣扬祂手的化工}和\Quote{诸天是祢手的工程}，以及其他类似的短语，以比喻的方式理解，就神性的\Quote{手}和\Quote{肢体}而言，那么\Quote{天主这样亲手工作}的话有什么荒谬呢？既然天主这样工作并无荒谬之处，那么祂发出\Quote{命令}也无荒谬之处；如此，凡奉祂命而行的事都应当是美丽而可赞美的，因为是命令它完成的天主。\par

% structure: p0128 source=h2 target=subsection reason=relative-heading-rank; base=h2

\subsection{第六十二章}

赛尔苏可能又误解了\Quote{因为主的口说了这话}这句话，或者因为有些无知的人曾轻率地解释这类事，而且他还不明白，按什么原则将根据身体肢体名称命名的部分归于天主的属性，于是他断言：\Quote{祂既没有口，也没有声音。}的确，如果声音是空气的震动、对空气的打击、或空气的一种形式，或由精通这类事物的人所给出的任何其他定义，那么天主不可能有声音；但所谓的\Quote{天主的声音}，在经文\Quote{众百姓看见了天主的声音}中，被认为是\Emph{看见}的，其中\Quote{看见}一词是按照圣经的习惯以属灵意义理解的。此外，他声称：\Quote{天主也不拥有任何\Emph{我们所知道的事物}}；但他没有指明\Emph{我们所知道的事物}是什么。如果他的意思是\Quote{肢体}，我们同意他，并理解\Quote{我们所知道的事物}是指那些被称为有形体的，且大致如此称呼的。但如果我们以普遍意义理解\Quote{\Emph{我们所知道的事物}}，那么有许多我们所知道的事物（可以归于天主）；因为祂拥有美德、福乐和神性。然而，如果我们赋予\Quote{\Emph{我们所知道的事物}}以更高的意义，既然我们所知道的一切都小于天主，那么我们也承认天主不拥有任何\Quote{\Emph{我们所知道的事物}}，这也没有荒谬之处。因为属于天主的属性远远超越所有事物，不仅是人性的本质所认识的，甚至那些已远远超越人性者的本质所认识的。如果他读过先知著作，一方面达味说：\Quote{但祢是同样的}，另一方面玛拉基亚说：\Quote{我是（主），且不改变}，他就会注意到我们中没有谁断言天主在行为或思想上有什么改变。因为祂保持不变，按万物的本性治理可变的事物，祂的圣言选择承担治理它们的职责。\par

% structure: p0130 source=h2 target=subsection reason=relative-heading-rank; base=h2

\subsection{第六十三章}

刻尔苏斯没有注意到\Quote{\Emph{按照}天主的肖像}与\Quote{天主的肖像}之间的区别，接着断言那\BibleQuote{一切受造物的首生者}是天主的肖像——圣言本身、真理本身，以及智慧本身，乃是祂美善的肖像，而人是\Emph{按照}天主的肖像受造的；此外，每一个以基督为头的男人，乃是天主的肖像和光荣——并且，他没有注意到\Quote{按照天主的肖像}这一说法属于人性的哪些特征，以及它在于一种从未有过、也不再拥有\BibleQuote{旧人及其行为}的本性，正是由于不具备这些品质，才被称为\BibleQuote{按照创造它的那一位的肖像}，他断言：\Quote{祂也没有把人造成祂的肖像；因为天主不是这样的，也不像任何其他（可见）存有的种类。}难道可以设想，那\Quote{按照天主的肖像}的元素应当存在于一个像人这样的复合存有的较低部分——我指身体——中吗？因为刻尔苏斯曾解释过，那按照天主的肖像受造就是这样吗？因为如果那\Quote{按照天主的肖像}仅在身体内，那么更好的部分，即灵魂，就被剥夺了那\Quote{按照祂的肖像}的东西，而这种（区别）存在于可朽坏的身体中——这是我们的任何人都没有提出的主张。但如果那\Quote{按照天主的肖像}是在两者一起之中，那么天主就必然是一个复合存有，仿佛由灵魂和身体构成，以便那\Quote{按照天主的肖像}的元素，即更好的部分，可以存在于灵魂中；而较低的部分，即\Quote{属于身体}的东西，可以存在于身体中——这又是我们的任何人都没有提出的主张。因此，剩下的就是：那\Quote{按照天主的肖像}必须被理解为存在于我们的\Quote{内在的人}中，这内在的人也被更新，其本性就是\BibleQuote{按照创造它的那一位的肖像}，当一个人成为\Quote{成全的}，如同\BibleQuote{你们在天之圣父是成全的}，并听到命令：\BibleQuote{你们应当是圣的，因为我，上主你们的天主，是圣的}，并学习训诫：\BibleQuote{你们该效法天主}，就将天主的肖像特征接受到他德性的灵魂中。此外，拥有这样一个灵魂的人的身体是天主的宫殿；并且天主居住在这灵魂中，因为它是按照祂的肖像受造的。\par

% structure: p0132 source=h2 target=subsection reason=relative-heading-rank; base=h2

\subsection{第六十四章}

刻尔苏斯又汇集了许多陈述，他以此作为我们所承认的话，但没有一个明智的基督徒会认可。因为我们中无人断言\Quote{天主分享形状或颜色}。祂甚至也不分享\Quote{运动}，因为祂坚定不移，祂的本性是恒常的，祂也邀请义人这样做，说：\BibleQuote{至于你，站在我这里。}如果某些表述暗示了祂身上有一种运动，比如\Quote{他们听见了上主天主在乐园的晚凉中\Emph{行走}的声音}，我们就必须这样理解：是罪人才将天主理解为运动，或者像我们理解天主的\Quote{睡眠}——这是喻意说法，或祂的\Quote{愤怒}，或其他类似属性。但是\Quote{天主连本体也不分享}。因为祂是被（他人）分享，而不是祂自己分享它们，并且祂被那些拥有天主的圣神的人所分享。我们的救主也不分享正义；但祂自己就是\Quote{正义}，祂是被义人所分享的。关于\Quote{本体}的讨论将会冗长而困难，尤其是当问题在于：那常存不变且非物质的东西是否可恰当地称为\Quote{本体}，以至于人们会发现天主是超越\Quote{本体}的，祂借着职分和能力，将其\Quote{本体}传达给那些祂通过圣言而通传自己的人，正如祂对圣言本身所做的那样；或者即使祂\Quote{是}本体，但祂的本性也被说成是\Quote{不可见}的，正如论及我们的救主时所说的话：祂被称为\Quote{\Emph{不可见}的天主的肖像}，而从\Quote{不可见}一词则表明祂是\Quote{非物质}的。还有一个值得探究的问题，即\Quote{独生子}和\Quote{一切受造物的首生者}是否应被称为\Quote{本体的本体}、\Quote{理念的理念}和\Quote{万有的根源}，而高于一切的是祂的圣父和天主。\par

% structure: p0134 source=h2 target=subsection reason=relative-heading-rank; base=h2

\subsection{第六十五章}

克耳苏接着说论天主的话：\Quote{万物都是出于祂。}（我不知道他这样说如何抛弃了他所有的原则。）而我们的保禄却宣称：\BibleQuote{万物都出于祂，借着祂，并归于祂。}显示祂是万物质体的开始（借着\Quote{出于祂}这话），又是其存续的纽带（借着\Quote{借着祂}这话），以及其终极的终向（借着\Quote{归于祂}这话）。的确，天主并非出于任何事物。但当克耳苏添加说：\Quote{祂不能借着言语达到}，我加以区分：若他指的是在我们\Emph{内}的言语——无论是心思中的言语，或是发出的言语——我也承认天主不能借着言语达到。然而若我们留意这段经文：\BibleQuote{在起初已有圣言，圣言与天主同在，圣言就是天主。}我们则认为天主要借着\Emph{这}圣言达到，并且不仅借着祂而被领悟，任何祂所启示父的人也都能领悟；这样我们将证明克耳苏断言\Quote{天主也不能借着言语达到}的虚假。此外，那说法\Quote{祂不能以名称表达}需要加以区分。若他意指确实没有言语或符号能代表天主的属性，这说法是真的，因为有许多品质不能用言语表明。例如，谁能用言语描述棕枣的甜味与无花果的甜味的区别？谁能用言语区分并说明每一事物的个别特性？因此若这样天主不能以名称描述，并不稀奇。但若你取这句意为可以用言语表达一些天主的属性，以便引导听者，犹如牵着手，使之能领悟关于天主的一些事——尽人性所能——那么说\Quote{祂\Emph{能}以名称描述}并无荒谬。关于\Quote{因为祂没有经历任何能用言语表达的痛苦}，我们也作类似区分。确实，天主性超出一切痛苦。关于这点就说到这里。\par

% structure: p0136 source=h2 target=subsection reason=relative-heading-rank; base=h2

\subsection{第六十六章}

让我们也看克耳苏的下一个陈述，在其中他引入了一个好像听到上述话的人，此人表达如下：\Quote{那么，我如何认识天主？我如何学习那通向祂的道路？你又如何将祂指示给我？因为现在你确实将黑暗抛在我眼前，我什么也看不清楚。}然后他好像回答那个困惑的人，并认为他给出了为何黑暗倾倒在说出上述话的人眼前的理由，他断言：\Quote{那些人们要带出黑暗进入光明辉煌的人，因无法忍受其光辉，视力受损，因而想象自己被打瞎了。}对此我们回答说：所有那些确实坐在黑暗中，并扎根于黑暗的人，是那些凝视画家、塑像者和雕刻家的邪恶手工，不愿向上看，并从一切可见和可感之物上升到万物的创造者——光明——的人；而另一方面，每一个人都在光明中，凡追随了\OriginalTerm{圣言}{Word}光辉的人，圣言指出了由于何种无知、不敬和对神圣事物的无知，这些对象被当作天主崇拜，并将那渴望得救的灵魂导向那超越万有的非受造的天主。因为\BibleQuote{那坐在黑暗中的百姓——外邦人——看见了一道大光，那坐在死荫之地的人，有光照亮了他们}——这光就是天主耶稣。因此，没有基督徒会对克耳苏或任何对神圣圣言的控告者回答：\Quote{我如何认识天主？}因为每一个基督徒都按自己的能力认识天主。没有人会问：\Quote{我如何学习那通向祂的道路？}因为他已听见那说：\BibleQuote{我就是道路、真理、生命}的那一位，并在旅程中尝到了由此而来的幸福。没有一个基督徒会对克耳苏说：\Quote{你如何将天主指示给我？}\par

% structure: p0138 source=h2 target=subsection reason=relative-heading-rank; base=h2

\subsection{第六十七章}

事实上，克耳苏斯所说的那句话是真确的：任何人听到他的话，都会回答说，既然他的话\Emph{是}黑暗的话，\Quote{你将黑暗抛在我眼前}。克耳苏斯和像他那样的人，确实渴望将黑暗抛在我们眼前；然而我们藉着圣言之光，驱散他们不敬意见的黑暗。基督徒确实可以反问克耳苏斯，他说的没有一样是清楚或真实的，\Quote{我在你的一切陈述中看不到任何\Emph{你的}东西}。因此，克耳苏斯并没有引领我们\Quote{从黑暗}进入\Quote{光明的辉煌}；相反，他想要将我们从光明带入黑暗，以黑暗为光，以光为黑暗，使自己承受依撒意亚先知以下述方式所描述的祸：\BibleQuote{祸哉，那些以黑暗为光，以光为黑暗的人}。但我们这些灵魂的眼睛已被圣言打开的人，看到了光与暗的区别，我们宁愿完全\Quote{站在光明中}，绝不与黑暗有任何关系。此外，真光具有生命，知道向谁显示他的全部光辉，向谁显示他的光明；因为他并不因接受者眼睛的薄弱而显示他的光芒。如果我们一定要谈论\Quote{视觉受影响和受损伤}，那除了陷入对天主无知、因情欲而不能看见真理的人的眼睛之外，我们还能说是谁的眼睛呢？然而，基督徒绝不认为他们被克耳苏斯或任何其他反对敬拜天主的人的言语所蒙蔽。但那些因追随众多迷途者和给魔鬼守节的人群而自觉被蒙蔽的人，应当亲近那能赐予视力的圣言，以便像那些贫穷瞎眼、躺在路边、因对耶稣说\BibleQuote{达味之子，可怜我吧}而被耶稣治愈的人一样，他们也能获得怜悯，恢复他们新鲜而美丽的视力，如同天主的圣言所能创造的那样。\par

% structure: p0140 source=h2 target=subsection reason=relative-heading-rank; base=h2

\subsection{第六十八章}

因此，如果克耳苏斯问我们，我们如何认为我们认识天主，以及我们如何靠祂得救，我们就要回答：天主的圣言，进到寻求祂的人当中，或进到接受显现的祂的人当中，能够使人认识并启示那位在圣言显现之前无人见过的圣父。除了天主圣言——祂\BibleQuote{在起初就与天主同在}，并为了那些依附于血肉、已成为血肉的人而成了血肉，以便让那些不能看见祂的人（因为祂是圣言，与天主同在，且是天主）能接受祂——还有谁能拯救并引导人的灵魂到达万有之神那里呢？祂以人形说话，宣布自己是血肉，祂呼唤那些血肉之人，以便首先使他们按照成了血肉的圣言得以转变，然后引导他们上升，看见祂在成为血肉之前是如何存在的；这样，他们领受恩惠，并从对祂按照血肉的大启蒙中上升，说：\BibleQuote{即使我们曾按血肉认识基督，但如今不再这样认识祂了。}因此祂成了血肉，且成了血肉后，\BibleQuote{在我们中间搭帐棚}，并非住在我们外面；在搭帐棚并住在我们\Emph{内}之后，祂没有停留在最初显现的形象中，而是使我们登上祂圣言的高山，向我们显示祂荣耀的形象和祂衣服的光辉；不仅祂自己的形象，还有属神的法律（即梅瑟）的形象，与耶稣一同在荣耀中被看见。此外，祂还向我们显示了全部预言，即使在祂降生成人后也没有消失，而是被接到天上，其象征是厄里亚。看见这些事的人可以说：\BibleQuote{我们看见了祂的荣耀，正如父独生者的荣耀，满溢恩宠和真理。}因此，克耳苏斯在他为我们设想的问题的回答中表现出相当大的无知：\Quote{我们如何认为我们能认识天主？以及我们如何知道我们将靠祂得救？}因为我们的答案正是我们刚才所说的。\par

% structure: p0142 source=h2 target=subsection reason=relative-heading-rank; base=h2

\subsection{第六十九章}

然而，策尔苏斯断定，我们所给予的答复是基于一种可能的推测，他承认他如下描述了我们的答复：\Quote{既然天主是伟大的且难以看见，祂便将自己的圣神放入一个与我们的身体相似的身体中，并把它派遣到我们这里，使得我们能够听见祂并认识祂。}但是在我们看来，万物的天主和父并非唯一伟大的存在者；因为祂将自己和祂的伟大（分享）分赐给了祂的独生子和一切受造物中的首生者，以便祂，作为不可见天主的肖像，甚至在其伟大中也保存了父的肖像。因为不可能存在一个可说匀称而美丽的不可见天主的肖像，却没有同时保存祂伟大的肖像。此外，在我们看来，天主的不可见，是因为祂不是形体，然而祂\Emph{能}被那些用心（即理智）看见的人看见；当然不是用任何的心，而是用纯洁的心。因为一颗受污染的心注视天主是不合宜的；因为那必须本身纯洁才能相称地观看那纯洁者。诚然，我们可以承认天主是\Quote{难以看见的}，然而祂并非唯一如此的存在者；因为祂的独生子也是\Quote{难以看见的}。因为天主圣言是\Quote{难以看见的}，祂的智慧也是如此，天主藉着这智慧创造了万有。因为谁能够看见那彰显在万有体系之每一部分中，并藉以创造每一事物的智慧呢？那么，并非因为天主是\Quote{难以看见的}，祂才派遣了天主祂的圣子，作为\Quote{容易看见的}对象。并且因为策尔苏斯不明白这点，他描述我们说了这样的话：\Quote{因为天主是\InnerQuote{难以看见的}，祂便将祂自己的圣神放在一个与我们的身体相似的身体中，并把它派遣到我们这里，使得我们能够听见祂并认识祂。}如今，正如我们已经阐明了的，圣子也是\Quote{难以看见的}，因为祂是天主圣言，万有藉着祂而受造，并且祂\BibleQuote{居住在我们中间}\BibleRef{若1:14}。\par

% structure: p0144 source=h2 target=subsection reason=relative-heading-rank; base=h2

\subsection{第七十章}

如果策尔苏斯确实明白了我们关于天主圣神的教导，并且知道了\BibleQuote{凡被天主圣神引导的，都是天主的子女。}\BibleRef{罗8:14}，他就不会把那个他描述为出自我们的答复归到自己身上，即\Quote{天主把自己的圣神放入一个身体中，并把它派遣到我们这里。}因为天主持续不断地将自己的圣神赐予那些能领受的人，尽管祂居住于配得之人的（心中）并非藉着分割和分离的方式。在我们看来，圣神也不是一个\Quote{形体}，正如火也不是一个\Quote{形体}——天主在经上被说成是\BibleQuote{我们的天主是吞噬之火。}\BibleRef{希12:29}因为所有这些都是比喻的表达，用来藉熟悉的和形体性的术语指示\Quote{有理智的存在者}的本性。同样，如果罪恶被称为\BibleQuote{木、草、禾秸}\BibleRef{格前3:12}，我们不会主张罪恶是有形体的；如果福分被称为\BibleQuote{金、银、宝石}\BibleRef{格前3:12}，我们不会主张福分是\Quote{有形体的}；同样，如果天主被说成吞噬木、草、禾秸和一切罪恶实质的火，我们不会理解祂为一个\Quote{形体}；同样，如果祂被称为\Quote{火}，我们也不会理解祂为一个形体。同样，如果天主被称为\Quote{神}，我们的意思并不是祂是一个\Quote{形体}。因为圣经的习惯是，将\Quote{有理智的存在者}称为\Quote{神体}和\Quote{属神的事物}，以区别于那些属于\Quote{感官}的对象；正如保禄所说的：\BibleQuote{但我们之所以够资格，是出于天主，祂也使我们成了新约的仆役，不是文字，而是神；因为文字叫人死，神却叫人活。}\BibleRef{格后3:5-6}在这里，藉着\Quote{文字}，他指的是\Quote{那明显的、感官性的圣经解释}，而藉着\Quote{神}，则是\Quote{理智}的对象。同样，对于\BibleQuote{天主是神}\BibleRef{若4:24}这个说法也是如此。并且因为法律的规条被撒玛利亚人和犹太人以形体和字面的方式遵守，我们的救主对撒玛利亚妇女说：\BibleQuote{时候将到，你们不在这山上，也不在耶路撒冷朝拜父。天主是神，朝拜祂的人应当以心神和真理朝拜祂。}\BibleRef{若4:21-24}藉着这些话，祂教导人们：天主必须在心神中朝拜，而不是在肉身中、带着肉身的祭献。祂被理解为神，与人们以心神和理智向祂献上的朝拜相称。然而，我们不是以形象朝拜父，而是\Quote{藉着真理}，这真理\BibleQuote{是经由耶稣基督而来的}\BibleRef{若1:17}，这是在梅瑟颁布法律之后。因为当我们归向主（主就是神）时，祂就除去那在诵读梅瑟时盖在心上的帕子。\BibleQuote{当我们归向主（主就是神）时，祂就除去那在诵读梅瑟时盖在心上的帕子。}\BibleRef{格后3:16}\par

% structure: p0146 source=h2 target=subsection reason=relative-heading-rank; base=h2

\subsection{第七十一章}

因此，\Person{塞耳苏}既不了解有关天主圣神的道理（\Quote{属血气的人不接受天主圣神的事，因为在他看来，这些都是愚妄；他也不能认识，因为这些事唯有藉圣神才可审断。}），就随意编织出一套（如他所愿的）谬论，以为我们称天主为神，在这点上与希腊的斯多葛派毫无区别，他们主张\Quote{天主是神，遍及万物，并将万有包含在自身之内}。确实，天主的眷顾与上智安排遍及万物，但并非像斯多葛派所说的神那样。眷顾确实包含它所眷顾的一切，并统摄它们，但不像有形的物体包含其内容那样——因为他们也认为那些是\Quote{形体}——而是以\Emph{神圣}的力量统摄其所包含的。诚然，根据廊下派哲学家的观点，他们宣称本原是\Quote{有形体的}，并因此认为万物都是可朽的，甚至胆敢使万有的天主成为可朽的；那么，天主圣言，那降临到最卑微之人中间的圣言，若非在他们看来显得过于荒谬，就不过是一种\Quote{有形体的}神而已。然而，在我们看来——我们力图证明理性灵魂超越一切\Quote{有形体的}本性，且是无形、非物质的实体——天主圣言，万物由祂而造，祂降临为使万物藉圣言而成，不仅为了人，也为了那些被认为仅受自然支配的最卑微之物，祂不是形体。斯多葛派可能将万物付之一炬；但我们知道没有任何非物质的实体能被火销毁，（也不相信）人的灵魂，或\Quote{天使}、或\Quote{宝座}、或宰制者、或\Quote{率领者}、或\Quote{掌权者}的本体能被火焚毁。\par

% structure: p0148 source=h2 target=subsection reason=relative-heading-rank; base=h2

\subsection{第七十二章}

因此，\Person{塞耳苏}作为一个不知道天主圣神本性的人，徒然断言：\Quote{寓居在人身体内的天主之子是一个神，故此这位天主之子不会是不死的。}接着他语无伦次地说，好像我们中有些人并不承认天主是一个神，而只承认关于祂的儿子，并且他认为他可以用这样的说法反驳我们：\Quote{没有一种神是永远持续的。}这差不多就像，当我们称天主为\Quote{吞噬之火}时，他却说\Quote{没有一种火是永远持续的}；他不明白我们称我们的天主是火的意义，以及祂所吞噬的是什么，即罪过和不义。因为一位良善的天主，在每个人藉其努力显示了他是何等样的斗士之后，以祂惩戒之火焚烧恶习，这是合宜的。接着，他又假设我们并不主张的论点，即\Quote{天主必然已经断气了}；由此也推出耶稣不能带着祂的身体复活。因为天主不会接收祂所交付的、已被身体接触玷污了的灵魂。然而，我们回答那些从未由我们提出的言论，如同是我们的，这是愚蠢的。\par

% structure: p0150 source=h2 target=subsection reason=relative-heading-rank; base=h2

\subsection{第七十三章}

他继续重复自己，在说了许多以前说过的话，并嘲笑了天主由童贞女诞生——对此我们已经尽我们所能答复了——之后，他加上了以下的话：\Quote{如果天主愿意从自己派遣祂的神降临，有什么必要将它吹入一个妇人的子宫呢？因为一个已经知道如何造人的人，祂也可以为这个人形成一个身体，而不将祂自己的神投入如此大的玷污之中；这样，如果祂直接从上面取得存在，祂就不会被怀疑不信了。}他说这些话，是因为他不知道那要服务于人类救恩的身体是纯洁的、没有玷污的童贞所生。因为，他引用斯多葛派的言论，假装不知道关于\OriginalTerm{中性事物}{things indifferent}的道理，他认为神性被投在玷污之中，并且要么因处于妇人的身体里直到身体形成，要么因取了身体，而被玷污。他这样做，就像那些人想象太阳的光线被粪便和臭气污染，而在这些东西中不再保持纯洁一样。然而，按照\Person{塞耳苏}的观点，如果耶稣的身体没有经过生育而形成，那么看见那身体的人就会立刻相信它不是由生育形成的；然而，一个被看见的对象并不会同时指明它起源的本性。例如，假设有一些蜂蜜（放在面前）不是由蜜蜂制造的，没有人能从味觉或视觉上断定它不是蜜蜂的作品，因为来自蜜蜂的蜂蜜并不通过感官表明其起源，只有经验才能告诉我们它不是来自蜜蜂。同样，经验告诉我们酒来自葡萄树，因为味觉不能区分（酒）是否来自葡萄树。因此，可见的身体并不表明它存在的方式。你会被诱导接受这个观点，藉着（观察）天上的天体，我们凝视它们时感知它们的存在和光辉；然而，我猜想，它们的外表并不提示我们它们是被造的还是非被造的；因此在这些问题上存在不同意见。然而，那些说它们是被造的人，对于它们被造的方式意见也不一致，因为外表并不提示它，尽管理性的力量可能发现它们是被造的，以及它们的受造是如何实现的。\par

% structure: p0152 source=h2 target=subsection reason=relative-heading-rank; base=h2

\subsection{第七十四章}

此后，他再次论及玛尔基翁的观点（此前已多次提及），其中一些他陈述正确，另一些则误解了；然而，我们无需回答或反驳这些。接着，他提出各种可能支持或反对玛尔基翁的论据，列举哪些观点为他开脱指控，哪些使他陷入指控；当他想要支持\Quote{耶稣是预言主题}这一陈述时——以便对玛尔基翁及其追随者提出指控——他明确问道：\Quote{一个如此受罚的人，若非这些事已对他预言，怎能被证明是天主之子？}然后他转而取笑，并依其惯常做法，对此加以嘲弄，引入\Quote{两个天主之子，一个是造物主之子，另一个是玛尔基翁之神之子；他描绘他们的单打独斗，说父辈间的神战如同鹌鹑间的争斗；或者说，父辈因年老无用、陷入老迈昏聩，完全不干预彼此，却让儿子们决一胜负。}我们之前对他的评论将反过来针对他：\InnerQuote{哪个老妇不会因用他题为\WorkTitle{真道论}的作品中插入的这类故事哄孩子入睡而感到羞耻？因为当他本该严肃从事论证时，他却置严肃论证于不顾，转而从事戏谑和滑稽表演，以为自己在写闹剧或嘲笑诗句；他没有注意到，这样的方法适得其反，其目的是使我们放弃基督信仰而赞同他的观点，而这些观点如果以某种庄严性陈述，或许会显得更有说服力；但既然他继续嘲笑、讥讽、扮演小丑，我们回答说，正是因为他没有有力的论证（因为他既没有，也无法理解），才转而从事如此胡言乱语。}\par

% structure: p0154 source=h2 target=subsection reason=relative-heading-rank; base=h2

\subsection{第七十五章}

在上述言论之后，他又补充道：\Quote{既然一位神圣的神灵居住在（耶稣的）体内，那么它必定在庄严、美丽、力量、声音、感人度或说服力方面与其他存有不同。因为那个拥有超越其他存有的神圣特质的人，不可能不拘泥于其他众人；然而此人却在任何方面与常人无异，据他们报告，他又矮、又丑、又卑贱。}很明显，从这些话可以看出，当责尔苏斯想控告耶稣时，他引用了圣经，仿佛他相信这些经文明显适合为控告祂提供把柄；但在同一经文中，当出现与所提出的控告相悖的陈述时，他却假装不知道它们！确实，经上承认记载了关于耶稣身体\Quote{丑陋}的一些陈述，但并非如他所说\Quote{卑贱}，也没有确凿证据表明祂\Quote{矮小}。以赛亚的话是这样的，他预言有关祂的话，说祂要来造访大众，并非以形貌俊美，也非以任何超群的美丽：\BibleQuote{主啊，谁相信了我们的报道？上主的手臂向谁显示？祂在祂面前发出宣告，如同孩童，如同干旱之地的根苗。祂没有形态，也没有荣耀，我们看见了祂，祂没有形态，也没有美丽；但祂的形态没有尊荣，且低于人子。}这些段落，责尔苏斯听到了，因为他认为这些对他控告耶稣有用；但他没有注意第四十五篇圣咏的话，以及为何其中说到：\BibleQuote{大能者啊，将你的剑佩在腿上，带着你的荣美和威严；继续前进，得胜为王。}\par

% structure: p0156 source=h2 target=subsection reason=relative-heading-rank; base=h2

\subsection{第七十六章}

尽管如此，让我们假设他没有读过那预言，或者他读过了，却被那些曲解它不指耶稣基督的人所误导。他对福音书有什么可说的？福音叙事中，耶稣上了高山，在门徒面前变了容貌，显现在荣耀中，那时梅瑟和厄里亚\BibleQuote{在光荣中显现，谈论祂在耶路撒冷将要完成的去世}；或者先知说\BibleQuote{我们看见了祂，祂没有形态，也没有美丽}等等？责尔苏斯接受这预言指耶稣，他在接受时被蒙蔽了，没有看到这一点是极大的证明：那显现为\BibleQuote{没有形态}的耶稣是天主之子，因为祂的外貌在祂出生前多年就成了预言的主题。但如果另一位先知谈到祂的荣美和威严，他将不再接受那预言指基督！如果从福音书中可以清楚得知\Quote{祂没有形态也没有美丽，但祂的容貌没有尊荣，且低于人子}，那么可以说，责尔苏斯的评论不是指着预言书，而是指着福音书说的。但现在，既然福音书和宗徒的著作都没有表明\Quote{祂没有形态也没有美丽}，显然我们必须接受先知关于基督的宣告为真实，这将阻止对耶稣的控告被提出。\par

% structure: p0158 source=h2 target=subsection reason=relative-heading-rank; base=h2

\subsection{第七十七章}

可是，那说\Quote{由于一个神圣的神灵居住在（耶稣的）身体内，它必定在伟大、或声音、或力量、或令人印象深刻、或说服力方面与其他存在的身体不同}的人，岂不也观察到，祂的身体根据观看者的能力（因而也根据其相应的效用）而改变关系，因为祂向每个人显现的样貌，正是那人应当看到的吗？再者，质料本性上能够被改变和变化，并能被转变为造物主所选择的任何形态，且能接受工匠所愿的一切性质，这并不足为奇：它时而具有一种性质，据此它被称为\Quote{祂没有形态，也没有华美}，时而具有另一种如此荣耀、庄严而奇妙的性质，以至于那些看到如此超凡可爱之美的观众——即与耶稣一同（登上）山峰的三位门徒——会脸伏于地。然而，他会说，这些都是捏造，与神话毫无区别，正如关于耶稣的其他奇迹一样；对于这一反驳，我们在前面已经更详细地回答过了。但在这一道理中也有某种神秘之处，它宣告耶稣变化的外貌应归因于\OriginalTerm{圣言}{Word}的本性，祂向群众显现的方式与向那些能跟随祂到我们所说的\OriginalTerm{高山}{lofty mount}的人不同；因为对于那些仍然留在山下、尚未预备好攀登的人，圣言\Quote{没有形态，也没有华美}，因为对这样的人来说，祂的形态是\Quote{没有尊荣的}，且低于比喻上称为\Quote{人子}的言语。因为我们可以说，哲学家们——他们是\Quote{人子}——的言语似乎比那向群众宣讲的天主之\OriginalTerm{圣言}{Word}美丽得多，后者也展示了（所谓）的\Quote{宣讲的愚妄}，而仅仅以此来看的人就说：\Quote{我们看见了祂；但祂没有形态，也没有华美。}确实，对于那些已获得能力跟随祂、以便甚至当祂登上\Quote{高山}时也能陪伴祂的人，祂\Emph{有}一个更神圣的容貌，他们看见它，如果（他们中间）有一位\Person{伯多禄}，他已在自己内接受了建立在圣言之上的教会建筑，并获得了如此（之善的）习惯，以致阴府的门都不能战胜他，且被圣言从死亡之门高举，以便他能在\Place{熙雍}女儿的门内\Quote{传扬天主的赞美}，以及其他那些从令人印象深刻的宣讲而生、并且丝毫不逊于\Quote{雷霆之子}的人。但\Person{策尔苏斯}和神圣圣言的敌人，以及那些没有本着真理的精神审查基督教道理的人，怎能知道耶稣不同外貌的意义呢？我也提及祂生命的不同阶段，以及祂在受难前和从死者中复活后所行的任何事。\par

% structure: p0160 source=h2 target=subsection reason=relative-heading-rank; base=h2

\subsection{第七十八章}

接着，\Person{策尔苏斯}提出了以下性质的观察：\Quote{再者，如果天主，像喜剧中的朱庇特一样，从长久的沉睡中醒来，想要拯救人类脱离邪恶，为什么祂要把你们所说的这个神灵只派遣到（地球的）一个角落？祂本应同样地将它吹入许多身体，并派遣它们到全世界。而喜剧诗人为了在剧场中引起笑声，写道朱庇特醒来后派遣\Person{墨丘利}到雅典人和拉栖代梦人那里；但你们不认为你们派遣\OriginalTerm{圣子}{the Son of God}到犹太人那里更加可笑吗？}请看这样的言论中\Person{策尔苏斯}的不敬特征，他，不像一位哲学家，竟把一个喜剧作家——其本分是引起笑声——拿来与我们的天主、万物的创造者相比较，把祂比作剧中那个醒来后派遣（信使）墨丘利的人物！诚然，我们在前面已经说明，天主并非如同从长久的沉睡中醒来一般派遣\Person{耶稣}到人类，祂如今为了正当的理由完成了降生成人的救恩计划，但祂一直赐福予人类。因为人类中从未有任何高尚行为，是神圣的\OriginalTerm{圣言}{Word}没有访问那些有能力、哪怕片刻接受神圣圣言如此运作的灵魂的。此外，\Person{耶稣}显然只到一个角落（的地球）的来临是有充分理由的，因为那作为预言主题的祂必须在那些已认识独一天主之道、阅读祂先知著作、并已知道基督宣告的人中出现，并且祂来到他们当中，正是在圣言将从这一个角落传遍全世界的时刻。\par

% structure: p0162 source=h2 target=subsection reason=relative-heading-rank; base=h2

\subsection{第七十九章}

因此，没有必要到处存在许多身体和许多像\Person{耶稣}一样的灵，为使整个人类世界能被天主的圣言光照。因为这一位圣言就足够了，祂如\BibleQuote{正义的太阳}升起，从\Place{犹太}发出祂来临的光芒，进入所有愿意接受祂者的灵魂。但若有人渴望看到许多充满神圣之灵的身体，相似于唯一的\Person{基督}，服务于各处人们的救恩，就请他注意那些在各地以纯正的教义和正直的生活宣讲\Person{耶稣}福音的人，他们自己也被圣经称为\Quote{受傅者}，经上记载：\BibleQuote{不要触摸我的受傅者，也不要伤害我的先知。}因为我们听说假基督要来，也得知世上已有许多假基督，同样，我们知道\Person{基督}已经来了，我们看到由于祂，世上有了许多受傅者，他们像祂一样爱好正义、憎恨不义，因此天主，\Person{基督}的天主，也用\BibleQuote{喜乐之油}傅抹了他们。但由于祂比祂的同伙更爱好正义、更憎恨不义，祂也获得了傅油的首批果实，并且，如果我们可以这样说的话，获得了喜乐之油的全部傅油；而祂的同伙则按各自的能力分享了祂的傅油。因此，既然\Person{基督}是教会的头，以致\Person{基督}和教会形成一个身体，这香膏从头下降到\Person{亚郎}的胡须——象征完美的人——并且这香膏下降直到他衣袍的边缘。这就是我对\Person{责尔苏}不敬言语的答复，他说道：\Quote{祂本应同样将（祂的圣神）吹入许多身体，并派遣祂进入整个世界。}确实，喜剧诗人为了引人笑，描绘\Person{朱庇特}睡着并从睡梦中醒来，派遣\Person{墨丘利}到希腊人那里；但圣言知道天主的本性不受睡眠影响，可以教导我们：天主在适当的时节，并按照正当理性的要求，管理世界的事务。然而，由于天主审判的伟大和不可理解，无知的人犯错并不奇怪，\Person{责尔苏}也在其中。因此，天主子被派遣到犹太人那里，——先知们曾出现在他们中间——为的是在犹太人中间以一个身体的形式开始，然后以威能和力量升起在灵魂的世界之上，这些灵魂不再愿意被天主遗弃，这一点并不荒谬。\par

% structure: p0164 source=h2 target=subsection reason=relative-heading-rank; base=h2

\subsection{第八十章}

在此之后，\Person{责尔苏}认为合适，将加色丁人称为从远古时代起就最受神灵感召的民族，正是从他们那里，迷惑人的占星术散布到人间。不仅如此，他还将术士归入同一类，从他们那里，巫术得名并传到其他民族，败坏和毁灭使用它的人。在本卷前面部分，（我们提到）即使按\Person{责尔苏}的看法，埃及人也犯了错误，因为他们虽然在他们所谓的庙宇周围有庄严的围墙，但里面除了猿猴、鳄鱼、山羊、毒蛇或其他动物外，什么都没有；但此刻，他乐意也将埃及民族称为最受神灵感召的，而且是从远古时代起——也许是因为他们从早期就与犹太人争战。此外，那些娶自己母亲、与自己女儿行淫的波斯人，在\Person{责尔苏}看来，也是一个受到感召的民族；甚至印度人也是如此，其中一些人在前面他被描述为食人者。然而，对于犹太人，特别是古代的犹太人，他们不采用这些习俗中的任何一种，他不仅拒绝给予他们受感召之名，而且还宣布他们将立即灭亡。他发出这个关于他们的预言，无疑是赋有先知的能力，却没有注意到犹太人的全部历史以及他们古老而可敬的政体都是由天主管理的；正是因他们的跌倒，救恩才来到外邦人那里，并且\BibleQuote{他们的跌倒就是世界的富足，他们的缺乏就是外邦人的富足}，直到外邦人数的圆满进来，然后全以色列，即\Person{责尔苏}所不知道的，都可能得救。\par

% structure: p0166 source=h2 target=subsection reason=relative-heading-rank; base=h2

\subsection{第八十一章}

然而，我不明白他如何论及天主，说祂虽然\Quote{知道万事，却不知道这一点：祂将自己的圣子派遣到邪恶的人中间，这些人既要犯罪，又要刑罚祂。}显然，他在此刻似乎忘记了\Person{耶稣}将要受的一切苦难早已被天主的神所预见，并由祂的先知所预言；由此并不推出\Quote{天主不知道祂将祂的圣子派遣到邪恶和犯罪的人中间，这些人也要刑罚祂。}然而，他立刻又补充说，\Quote{我们对此的辩护是，这一切都已被预言。}但既然我们第六卷已经达到足够的篇幅，我们就此止住，并将在第七卷中，若天主愿意，开始论证，我们将考虑他以为能反驳我们陈述的理由，即所有关于\Person{耶稣}的一切都已由先知预言；由于这些理由众多，需要详细回答，我们既不想因本书篇幅而缩短这个主题，也不想为了避免缩短而使第六卷超出适当篇幅。\par

\SourceRecord{Translated by Frederick Crombie.；Ante-Nicene Fathers；Vol. 4.；Edited by Alexander Roberts, James Donaldson, and A. Cleveland Coxe.；Buffalo, NY: Christian Literature Publishing Co.,；1885.；<http://www.newadvent.org/fathers/04166.htm>.}

% structure: p0001 source=h1 target=section reason=page-title

\section{\WorkTitle{驳塞尔苏斯 卷七}}

% structure: p0002 source=h2 target=subsection reason=relative-heading-rank; base=h2

\subsection{第一章}

在这前六卷中，我们按照我们的能力，可敬的兄弟\Person{安博罗削}，努力回应\Person{塞尔苏斯}对基督徒提出的指控，并尽可能没有忽略任何事情，而不先对其进行全面仔细的审查。现在，当我们进入第七卷时，我们因\Person{耶稣基督}——\Person{塞尔苏斯}控告的那位——呼求天主，愿祂，天主的真理，光照我们的心，驱散错误的黑暗，遵照我们先知所说的那句话，我们现在将其作为祈祷献上：\Quote{以祢的真理消灭他们。}因为显然是天主以祂的真理消灭与真理对立的话语和推理；如此，当这些被消灭时，所有脱离欺骗的人可以继续与先知同说：\Quote{我要向祢自由地祭献，}并向至高者奉献合理而无烟的祭献。\par

% structure: p0004 source=h2 target=subsection reason=relative-heading-rank; base=h2

\subsection{第二章}

\Person{塞尔苏斯}现在着手反驳那些声称犹太先知预言了基督耶稣生命中事件的人的观点。首先，让我们提及他的一个看法：那些假定存在另一位天主，不同于犹太人的天主的人，没有理由回应他的反对；而我们这些承认同一天主的人，则依靠关于耶稣基督所发出的预言来辩护。他的话是：\Quote{让我们看看他们如何能提出辩护。对于那些承认另一位天主的人，不可能有辩护；而那些承认同一天主的人，总会退回到同样的理由：\InnerQuote{这件事和那件事必定发生。}为什么？\InnerQuote{因为它很久以前就被预言了。}}对此，我们回答说：\Person{塞尔苏斯}最近对耶稣和基督徒提出的论据是如此软弱无力，以至于即使那些不虔敬到引入另一位天主的人，也能轻易推翻它们。的确，若不是给软弱者提供接受错误观念的任何借口是危险的，我们本可以自己提供答案，并向\Person{塞尔苏斯}表明他的观点是多么没有根据，即那些承认另一位天主的人无法面对他的论据。然而，目前让我们仅限于为先知辩护，继续我们之前对此主题所说的话。\par

% structure: p0006 source=h2 target=subsection reason=relative-heading-rank; base=h2

\subsection{第三章}

\Person{塞尔苏斯}继续论我们说：\Quote{他们不重视皮提亚女祭司的神谕，也不重视\Place{多多纳}、\Place{克拉鲁斯}、\Place{布兰奇达}、\Place{朱庇特·阿蒙}以及其他众多神谕所祭司的神谕；尽管在这些神谕的指引下，我们可以说殖民地被派遣出去，整个世界有人居住。但那些在\Place{犹太}，按照该地的风俗所说出或未说出的言论——正如它们在\Place{腓尼基}和\Place{巴勒斯坦}人中至今仍被传述——他们却视为奇妙的言论，且不可改变地真实。}关于这里列举的神谕，我们回答说，我们有可能从\Person{亚里士多德}和逍遥学派的著作中收集到不少论据来推翻皮提亚和其他神谕的权威。我们也可以引用\Person{伊壁鸠鲁}及其追随者的段落，表明即使在希腊人中间，也有人完全不信那些在整个希腊被承认和钦佩的神谕。但让我们承认，皮提亚和其他神谕所发出的回答不是假冒神灵启示之虚假人的话语；让我们看看我们是否能说服任何真诚的探究者，没有必要将这些神谕回答归属于任何神性，相反，它们可以追溯到邪恶的魔鬼——与人类为敌的邪灵，它们以这种方式想要阻碍灵魂上升，追求美德之路，并以真诚的虔诚回归天主。关于皮提亚女祭司，她的神谕似乎是最著名的，据说当她坐在\Place{卡斯塔利亚洞穴}的入口时，阿波罗的神谕之灵进入她的私处；当她被充满时，她便发出被敬畏地视为神圣真理的回答。请据此判断，那灵是否没有显示出它的亵渎和不洁本性，因为它选择进入女先知的灵魂，不是通过更合适的渠道——即既开放又不可见的身体毛孔，而是通过任何有节制的人都不会看见或提及的器官。这也不是发生一次或两次（那或许更可容许），而是每当她被认为接受阿波罗的灵感时都如此。此外，使女先知陷入如此癫狂和疯狂的状态以致失去自制，并不是神灵的本分。因为受神灵影响的人应当首先获得有益的效果；而这些效果不应首先由那些就自然或公民生活，或为了暂时的利益或好处而咨询神谕的人所享有；而且，当一个人与神性密切交往时，那应是最清晰的感知时刻。\par

% structure: p0008 source=h2 target=subsection reason=relative-heading-rank; base=h2

\subsection{第四章}

因此，我们可以从对圣经的考察中表明，犹太人的先知——他们因天主之神的光照而在先知工作上获得了必要的启示——是首先享受恩宠的人；借着圣神的接触——如果我可以这样说——他们的心思变得更加清明，灵魂充满了更明亮的光辉。身体不再成为德性生活的障碍；因为对于我们所称的\Quote{肉身的贪欲}，它已经麻木了。因为我们深信，圣神\Quote{歼灭身体的行为}，并摧毁那肉性情欲所激起的对天主的敌意。那么，如果皮提亚女祭司在预言时魂不守舍，那充满她心灵、以黑暗遮蔽其判断力的究竟是什么神灵呢？除非它属于那类邪魔，这类邪魔正为许多基督徒从附魔者身上驱逐出去。而且我们可以注意到，他们这样做并不借助任何奇异的魔术或咒语，而仅凭祈祷和简单的驱魔词，连最普通的人也能使用。因为执行这工作的大多是未受教育的人；这样便彰显了基督话语中的恩宠，以及邪魔的可鄙软弱；要克服并驱逐它们脱离人的身体和灵魂，并不需要那些在辩论上有力、在信仰事理上极为博学之人的能力和智慧。\par

% structure: p0010 source=h2 target=subsection reason=relative-heading-rank; base=h2

\subsection{第五章}

此外，如果不仅基督徒和犹太人，而且许多希腊人和异邦人也相信人的灵魂在脱离身体后仍然存活并存在；如果理性支持这样的观念：纯净的灵魂不被罪恶如铅般重压，便升到更高更纯净、更缥缈的身体领域，将它们较粗糙的身体连同不洁遗留在此下界；而那些被罪恶污染并拖曳到地上、甚至不能向上呼吸的灵魂，四处游荡，有时在坟墓周围显现为幽灵的影子，有时在其它地上的物体之间——如果事情如此，那么对于那些如我所说的，历经整个世代依附于特定的居所和地方（无论出于某种魔力还是它们自身的邪恶本性）的灵，我们该如何看待？我们岂非被理性驱使，将那些运用预言能力——这能力本身既不善也不恶——以欺骗人类、从而使人离开天主和祂纯洁事奉的灵视为邪恶的吗？此外，当我们补充说它们喜爱祭品的血和牺牲的烟雾气味，并以这些养育它们的身体，在这样出没之处感到快乐，仿佛在其中寻求生命的滋养，这点就更明显了：它们在这方面像那些堕落的男人，鄙视脱离感官的生命之纯洁，只倾向身体的快乐，以及那些快乐所在之属地且属身体的生命。如果德尔斐的阿波罗是一位神，如希腊人所认为的，难道他不会选择一个智者作他的先知吗？或者，如果找不到这样的智者，则会选择一个努力成为智者的人？为什么他不选择男人而选择女人来传达他的预言？如果他偏好后者，仿佛只有在女人的胸怀才能找到快乐，那么为什么他在女人中不选择一个童贞女来解释他的旨意呢？\par

% structure: p0012 source=h2 target=subsection reason=relative-heading-rank; base=h2

\subsection{第六章}

然而不然；那位在希腊人中备受赞赏的皮提亚，不认为任何智者——甚至不认为任何男人——配得所谓的\Quote{神圣附身}。在女人中，他没有选择一位童贞女，或一位因智慧或因其哲学成就而受推崇的；而是选择了一个普通的女人。也许更好一类的人太过良善，以致不适合作灵感的对象。此外，如果他是一位神，他本应运用他的预言能力如同诱饵，以便吸引人改变生活并实践德行。但历史上从未提及这类事。因为如果神谕曾称苏格拉底为所有人中最有智慧的，那么它借由关于欧里庇得斯和索福克勒斯所说的话，便削弱了那赞美之词的价值。原话是：——\par

\Quote{索福克勒斯有智慧，欧里庇得斯更有智慧，\newline{} 但比所有人更有智慧的是苏格拉底。}\par

因此，既然他将\Quote{有智慧的}这一称号赋予悲剧诗人，他尊崇\Person{苏格拉底}并非因为他的哲学，或因为他对真理和美德的热爱。说\Person{苏格拉底}优于那些为微薄报酬在舞台上竞争、并通过表演时而让观众泪流满面、伤心欲绝，时而让他们发出不雅笑声的人（因为这就是萨提尔剧的意图），这是对\Person{苏格拉底}的拙劣赞扬。而且，也许与其说是就哲学而言他将\Person{苏格拉底}称为所有人中最有智慧的，不如说是因他献给\Person{苏格拉底}和其他魔鬼的祭品。因为看来魔鬼在分配恩惠时，更注重献给他们的祭品而非美德的行为。因此，最优秀的诗人\Person{荷马}在描述通常发生之事时，为向我们展示什么最能影响魔鬼答应其崇拜者的愿望，引入了\Person{克律塞斯}，他为了一些花环和公牛、山羊的腿肉而获得了他为女儿\Person{克律塞伊斯}的祈祷的回应，以致希腊人被一场瘟疫逼迫归还她给他。我记得在某位毕达哥拉斯派学者的书中读到，他在论述那位诗人的隐义时写道，\Person{克律塞斯}向\Person{阿波罗}的祈祷以及\Person{阿波罗}随后降在希腊人身上的瘟疫，证明\Person{荷马}知道某些邪魔鬼怪喜爱祭品的烟雾，并为了奖赏献祭者而应允他们祷告毁灭他人。\Quote{他，就是\Person{朱庇特}，统治着寒冬的\Place{多多纳}，那里的先知们从不洗脚，睡在地上，}他拒绝了男性，并且，如\Person{塞尔苏斯}所说，任用\Place{多多纳}的妇女从事预言职务。就算有类似的神谕，如\Place{克拉洛斯}的神谕、\Place{布兰奇代}的神谕、\Place{朱庇特·阿蒙神庙}的神谕，或任何其他地方的神谕，但如何能证明这些是神，而不是魔鬼呢？\par

% structure: p0016 source=h2 target=subsection reason=relative-heading-rank; base=h2

\subsection{第七章}

关于犹太人中先知，有些人在成为天主启示的先知之前就是有智慧的人，而另一些人则在受天主启示时因心灵受到光照而变得有智慧。他们因过着几乎无可比拟的卓越生活——勇敢、高贵、不为危险或死亡所动——而被天主眷顾拣选，接受圣神，成为祂神圣圣言的保管者。因为理性教导我们，至高的天主的先知应当具有这样的品格，与之相比，\Person{安提斯泰尼}、\Person{克拉特斯}和\Person{第欧根尼}的坚定不过是儿戏。因此，正是由于他们坚定地持守真理，并忠实地责备恶人，\BibleQuote{他们被石头砸死，被锯锯死，受试探，被刀杀；他们披着绵羊山羊的皮各处奔跑；受穷乏、患难、苦害；在旷野、山岭、山洞、地穴漂流无定，世界配不上他们。}因为他们总是仰望天主和祂的祝福，这些祝福虽不可见、不能以感官感知，却是永恒的。我们有每位先知生平的历史，但目前只需关注\Person{梅瑟}的生平——他的预言记载在法律中；\Person{耶肋米亚}的生平——如那本以他命名的书中所记载；以及\Person{依撒意亚}的生平——他以无与伦比的苦行，赤身赤脚行走三年之久。请阅读并思考那些青年——\Person{达尼尔}及其同伴——的刻苦生活：他们如何戒绝肉食，只靠水和蔬菜生活。或者，如果你回溯更遥远的时代，想想\Person{诺厄}的生命，他曾预言；以及\Person{依撒格}，他给了儿子先知性的祝福；或者\Person{雅各伯}，他对他的十二个儿子逐一说话，开头是：\Quote{来，让我告诉你们在末日将发生何事。}这些先知和许多其他先知，代表天主说预言，预告了与\Person{耶稣基督}相关的事。因此，我们为此原因藐视\Person{皮提亚女祭司}的神谕，或那些在\Place{多多纳}、\Place{克拉洛斯}、\Place{布兰奇代}、\Place{朱庇特·阿蒙神庙}或其他无数所谓的先知所发出的神谕；而我们却以敬畏之心尊重犹太人的先知：因为我们看到这些人高贵、热忱而虔诚的生命，配得上圣神的默感，圣神奇妙的效果与魔鬼的占卜截然不同。\par

% structure: p0018 source=h2 target=subsection reason=relative-heading-rank; base=h2

\subsection{第八章}

我不知道是什么促使\Person{塞尔苏斯}在说\Quote{但按照那地的风俗，在\Place{犹太地}所说或未言说之事}时，使用了\Quote{或未言说}一词，仿佛暗示他心存怀疑，怀疑所记载的那些事从未被言说过。事实上，他对这些时代并不熟悉；他不知道那些预言基督来临的先知们，许多年前就预言了大量其他事件。他接着说，意在贬低古代先知：\Quote{他们预言的方式，就如我们今天在\Place{腓尼基}和\Place{巴勒斯坦}居民中仍然见到的那样。}但他没有告诉我们，他是指那些原则与犹太人和基督徒不同的人，还是指那些预言性质与犹太先知相同的人。无论如何，他的说法无论从哪个角度看都是错误的。因为那些没有接受我们信仰的人，从未做过任何接近古代先知所行之事；而且在较近的时代，自基督来临以后，犹太人中也没有兴起先知，他们公然地因对天主和对他们的先知所预言的那位的不敬而被圣神抛弃。此外，圣神在基督传教之初就显示了祂临在的迹象，并在祂升天后显示了更多；但自那时起，这些迹象已减少，不过仍有少数人身上留有祂临在的痕迹，这些人的灵魂已被福音净化，行为受其影响规范。\BibleQuote{因为受训诲的圣神必逃避诡诈，远离无知的思念。}\BibleRef{智 1:5}\par

% structure: p0020 source=h2 target=subsection reason=relative-heading-rank; base=h2

\subsection{第九章}

但是，既然\Person{塞耳苏斯}许诺要讲述预言在\Place{腓尼基}和\Place{巴勒斯坦}是如何传递的，说得好像这是他完全熟悉的事情，让我们看看他对此有什么要说的。首先，他假定预言有若干种类，但没有具体说明；事实上，他也做不到，这声明纯粹是虚张声势。然而，让我们看看他认为在这些民族中最完美的预言种类是什么。他说：\Quote{有许多人，虽然籍籍无名，却极其轻易地、在最微不足道的场合，无论在神庙内还是外，都表现出受灵感者的动作和姿态；另一些人在城市或军队中这样做，目的是吸引注意和引起惊奇。这些人习惯各自说：\InnerQuote{我是天主；我是天主子；或是，我是圣神；我来了，因为世界正在灭亡，你们，人啊，因你们的过犯而灭亡。但我愿意拯救你们，你们将看见我带着天上的能力再次回来。如今朝拜我的人是有福的。至于其余的人，我要降下永火，无论是城市还是乡村。那些不知道等待他们的惩罚的人，将徒然悔恨悲伤；而忠于我的人，我将永远保守。}}\newline{}\newline{}接着他又说：\Quote{在这些应许之上，又加上了奇怪、狂热、完全不可理解的言语，没有理智的人能够找到其意义：因为它们是如此隐晦，以至于毫无意义；但它们给每一个愚人或骗子以机会，按照自己的目的加以应用。}\par

% structure: p0022 source=h2 target=subsection reason=relative-heading-rank; base=h2

\subsection{第十章}

但若他在指控中诚实，他本应给出预言的准确措辞，无论是那些说话者自称是天主全能者的预言，还是天主子说话的预言，最后是那些以圣神名义的预言。因为这样他或许可以努力推翻这些断言，并表明那些敦促人离弃罪恶、谴责过去并预告未来的话语中并无神圣灵感。因为预言是由当时的人记录和保存的，使后来的人可以阅读并惊叹为天主的圣言，并且他们不仅从警告和劝诫中获益，也从预言中获益，这些预言通过事件证实出自天主之神，使人遵循法律和先知所载的虔诚实践。因此，先知们遵照天主的命令，将那些希望听者立即理解以指导其行为的事情完全坦白地宣告出来；而对于更深奥、更神秘的、超出常理理解范围的主题，他们则以谜语和寓言的形式，或所谓隐语、比喻或象征来表述。他们采用这种计划，是为了让那些为了真理和美德而不避辛劳、不辞辛苦的人可以探究其意义，并且找到后，可以按理性要求加以应用。但\Person{塞耳苏斯}，总是激烈地谴责，好像因无法理解先知的语言而愤怒，就这样嘲笑他们：\Quote{在这些伟大的应许之上，又加上了奇怪、狂热、完全不可理解的言语，没有理智的人能够找到其意义：因为它们是如此隐晦，以至于毫无意义；但它们给每一个愚人或骗子以机会，按照自己的目的加以应用。}\newline{}\newline{}\Person{塞耳苏斯}的这个陈述似乎巧妙设计，以劝阻读者进行任何探究或仔细研究其意义。在这方面，他并不像某些人，他们对一个被先知拜访以宣布未来事件的人说：\Quote{为什么这个疯子到你这里来了？}\par

% structure: p0024 source=h2 target=subsection reason=relative-heading-rank; base=h2

\subsection{第十一章}

我确实相信，比我能提出的更好的论据可以证明\Person{塞耳苏斯}的这些指控是虚假的，并展示预言的圣神灵感；但我们根据我们的能力，在注释\BibleBook{依撒意亚}{先知书}、\BibleBook{厄则克耳}{先知书}和\BibleBook{}{十二小先知书}的一些书卷中，已经逐字逐句详细解释了他所谓的\Quote{那些狂热而完全不可理解的段落}。如果天主在我们的时间内赐予我们恩宠，使我们在祂的圣言的知识上长进，我们将继续研究剩余的部分，或者至少那些我们能够阐明清楚的部分。其他有才智并希望研究\WorkTitle{圣经}的人也可以自己发现其意义；因为虽然有很多地方意义不明显，但绝没有像\Person{塞耳苏斯}所断言的那样\Quote{毫无意义}。也不像他说的\Quote{任何愚人或骗子都可以解释这些段落，使之符合自己的目的}。因为只有那些在基督的真道上有智慧的人（而且确实属于他们全体）才能展开甚至预言中隐晦部分的关联和意义，\BibleQuote{以属神的事解释属神的事}，并按照圣经作者的用法解释每一段话。\Person{塞耳苏斯}说他曾听过这样的人预言，这是不可信的；因为在\Person{塞耳苏斯}的时代，没有任何与古代先知相似的先知出现。如果有，那些听见并钦佩他们的人就会仿效古人的榜样，将预言记录下来。而且很显然，\Person{塞耳苏斯}在说谎，他说\Quote{他所听到的那些先知，在被他追问时，承认了他们真正的动机，并承认他们使用的模棱两可的话语实际上毫无意义}。他本应给出他所说的他所听过的人的名字，如果有的话，以便那些有判断力的人可以决定他的指控是真是假。\par

% structure: p0026 source=h2 target=subsection reason=relative-heading-rank; base=h2

\subsection{第十二章}

他又认为，那些援引先知著作来支持基督道理的人，对于其中归给天主的那恶、可耻或不洁的言论，无法给出恰当的答复；他假设无人能答复，便由此推出了一连串的推论，而这些推论没有一个是可以接受的。但他应当知道，凡愿意按照圣经教导生活的人，都明白这经文：\BibleQuote{愚昧人的知识，如同无意义的言语}\BibleRef{德 21:18}，并学会了\Quote{时常准备着，对任何问你们心中盼望原因的人，要予以答复。}他们不满足于仅仅断言某事某物曾被预言，而是努力消除任何表面的矛盾，并表明在这些预言中绝没有什么恶的、可耻的或不洁的事物，一切对于理解圣经的人都值得接受。但\Person{刻尔苏斯}本应从先知书中举出他认为恶的、可耻的或不洁的例子——若他看见这类章节的话；因为那样他的论证会更有力，也更有利于他的目的。然而他并未举出任何例子，只是大声坚持对圣经的无端指控，说其中有这些事。因此，我们没有理由为毫无根据的指控辩护，那些不过是空话而已，也不必费心去证明先知著作中毫无恶的、可耻的、不洁的或可憎的事物。\par

% structure: p0028 source=h2 target=subsection reason=relative-heading-rank; base=h2

\subsection{第十三章}

\Person{刻尔苏斯}说\Quote{天主行最无耻之事，或受最无耻之苦}，或说\Quote{祂纵容恶行}，这些话毫无真理；因为无论他说什么，这样的事从未被预言过。他本当从先知书中举出那些将天主描述为纵容恶行，或行且受最无耻之事的章节，而不该毫无根据地去误导读者的心思。事实上，先知们预言了基督所要受的苦，并说明了祂受苦的原因。天主因此也知道基督会受什么苦；但他从哪里得知天主的基督所要受的事是极其卑贱可耻的呢？他接着解释基督所受的最可耻最下贱的事，说了这样的话：\Quote{因为天主吃羊肉，或饮醋与苦胆，与吃污秽有何区别？}但按我们的信仰，天主并没有吃羊肉；\Person{耶稣}似乎吃了，但那只是因为祂有身体。至于预言中提到的醋和苦胆：\Quote{他们在我的食物中掺了苦胆，我渴时，他们给我喝酸醋}；我们已经讨论过这一点；由于\Person{刻尔苏斯}迫使我们再提，我们只说：那些抗拒真理之言的人，常常向基督天主子献上自己邪恶的苦胆和私欲的酸醋；但祂虽尝了，却不喝。\par

% structure: p0030 source=h2 target=subsection reason=relative-heading-rank; base=h2

\subsection{第十四章}

接下来，为了动摇那些因关于\Person{耶稣}的预言而相信\Person{耶稣}的人的信心，\Person{刻尔苏斯}说：\Quote{请问，如果先知们预言了伟大的天主——说得不苛刻——会变成奴隶，或生病，或死亡；难道因为这样的事被预言了，天主就必须死、生病或变成奴隶吗？祂必须死才能证明祂的天主性吗？但先知们绝不会发出如此邪恶不敬的预言。因此，我们不必查问某事是否被预言，而应查问那事本身是否可敬，是否配得上天主。在邪恶卑贱的事上，即使世上所有的人都疯狂地预言它，我们也不该相信。那么，虔诚的心灵怎能承认那些据说发生在祂身上的事，竟能发生在一位天主身上呢？}由此明显可见，\Person{刻尔苏斯}感到预言论证对于说服那些接受基督宣讲的人非常有效；但他似乎企图用一个相反的可能性来推翻它，即\Quote{问题不在于先知们是否发出了这些预言。}但若他愿意公正地推理而不回避，他倒应该说：\Quote{我们必须证明这些事从未被预言过，或者那些关于基督的预言从未应验在祂身上。}这样他就会确立自己的立场。这样一来，我们应用于\Person{耶稣}的那些预言究竟是什么，以及\Person{刻尔苏斯}如何能证明那种应用是错误的，便会清楚显明。我们也会看到，他是公正地反驳了所有我们为\Person{耶稣}从先知书中引用的证据，还是他自己被判定为无耻地以强词夺理抗拒最明显的真理。\par

% structure: p0032 source=h2 target=subsection reason=relative-heading-rank; base=h2

\subsection{第十五章}

在他假设有些预言本身不可能、且与天主的属性不符之后，他说：\Quote{如果这些事是关于至高天主的预言，我们是否仅仅因为它们被预言了，就必须相信它们关于天主呢？}因此他自以为证明了：即使先知们曾真实地预言了关于天主圣子的这类事，我们也不可能相信那些宣称祂会做或忍受这类事的预言。对此我们的回答是：这个假设是荒谬的，因为它结合了两条彼此对立的推理，从而相互否定。可以如下说明：第一条论证是：\Quote{如果至高天主的真先知说天主会成为奴隶、生病或死亡，这些事将会临到天主；因为伟大的天主的先知不可能说谎。}第二条是：\Quote{即使至高天主的真先知说这些事将要发生，但由于这些被预言的事本质上是不可能的，因此预言并不真实，所以那些被预言的事也不会发生在天主身上。}那么，当我们发现两个推理过程，其中大前提相同，却导致两个矛盾的结论时，我们使用称为\OriginalTerm{两命题定理}{theorem of two propositions}的论证形式来证明大前提是假的，在此案例中大前提是：\Quote{先知们曾预言伟大的天主会成为奴隶、生病或死亡。}于是我们得出结论：先知们从未预言过这样的事；该论证形式化表达如下：第一，如果第一件事是真的，那么第二件事是真的；第二，如果第一件事是真的，那么第二件事不是真的，因此第一件事不是真的。斯多葛派用来说明这种论证形式的具体例子如下：第一，如果你知道自己死了，那么你死了；第二，如果你知道自己死了，那么你没有死。结论是——\Quote{你不知道自己死了。}这些命题是这样展开的：如果你知道自己死了，你所知道的事是确定的，因此你死了。其次，如果你知道自己死了，你的死亡是一个认知对象；但由于死者一无所知，你知道了这一点证明你没有死。因此，将两个论证结合起来，你得出这样的结论——\Quote{你不知道自己死了。}现在，我们上面给出的塞尔苏斯的假设也大致相同。\par

% structure: p0034 source=h2 target=subsection reason=relative-heading-rank; base=h2

\subsection{第十六章}

然而，此外，他引入论证的预言与先知们实际关于耶稣基督所预言的截然不同。因为预言并没有预告天主会被钉十字架，当它们论及那将要受苦的一位时说：\BibleQuote{我们看见祂，祂无威仪，也无美貌；但祂的容貌被毁损，不像人子；祂是苦人，熟悉病苦。}那么，请注意他们多么明确地说是一个人要忍受这些人间的苦难。而耶稣自己，完全知道那将死的人必是一个常人，对控告祂的人说：\BibleQuote{如今你们却想要杀我，一个把从天主那里听到的真理告诉你们的人。}并且，在那个显现于人间的人当中，若有什么神圣的，即天主的独生子，一切受造的首生者，那自称\BibleQuote{我是真理}、\BibleQuote{我是生命}、\BibleQuote{我是门}、\BibleQuote{我是道路}、\BibleQuote{我是从天上降下的活粮}的一位，对于这一位及其本性，我们必须用与判断那在耶稣基督内所见的人的完全不同的方式来思考和推理。因此，你会发现没有一个基督徒，无论他多么单纯，多么不精于考据，会说出那死了的是\BibleQuote{真理}、\BibleQuote{生命}、\BibleQuote{道路}、\BibleQuote{从天上降下的活粮}、\BibleQuote{复活}；因为正是那以耶稣的人形向我们显现的那一位教导我们说：\BibleQuote{我是复活。}我们中间没有人，我说，如此狂妄以至于肯定\Quote{生命死了}、\Quote{复活死了}。塞尔苏斯的假设若成立，除非我们说先知们曾预言死亡会临到天主圣言、真理、生命、复活，或天主圣子所用的任何其他名号。\par

% structure: p0036 source=h2 target=subsection reason=relative-heading-rank; base=h2

\subsection{第十七章}

\Quote{先知们不会预言这事，因为这事涉及邪恶和不敬虔的事}——即伟大的天主竟成为奴隶或受死。但先知所预言的是与天主相称的：那神圣本性的光辉和印记，带着圣洁的人性灵魂——这灵魂要赋予\Person{耶稣}的身体以生命——来到世界，播撒祂圣言的种子，使所有接受并珍视它的人与至高的天主结合，并使那些在内里感受到天主圣言——祂要居于一个人的身体和灵魂中——之能力的人达到圆满的幸福。祂的确要居于其中，但并非让祂荣耀的一切光芒都局限在其中；我们也不应认为天主圣言的光只以这种方式照耀。因此，如果我们从\Person{耶稣}所怀的天主性来思考\Person{耶稣}，那么祂在此身份中所行的一切，都不冒犯我们对天主的观念，反而尽是圣洁；如果我们视祂为人，因与永生的圣言、绝对的智慧亲密无间而超乎众人，那么祂作为智慧和完美的一位，所忍受的正是那为了人类乃至所有理性受造物的益处而必须忍受的一切。一个人死了，祂的死不仅是为虔诚而死的榜样，也是推翻那恶灵魔鬼——它曾统治整个世界——权柄之战的第一击，这毫无荒谬之处。因为我们有迹象和凭证表明其帝国将覆灭，就是那些因\Person{基督}的来临而在各处摆脱魔鬼权势的人，他们在得赎后，将自己奉献于天主，日复一日努力在虔敬生活中进步。\par

% structure: p0038 source=h2 target=subsection reason=relative-heading-rank; base=h2

\subsection{第十八章}

\Quote{他们还会作此反思吗？如果犹太人的天主的先知预言那要来到世界的那位是同一位天主的圣子，那么祂怎能通过\Person{梅瑟}命令他们积聚财富、拓展疆域、充满大地，将各年龄的敌人置于刀剑之下，并彻底消灭他们——正如\Person{梅瑟}所说，祂自己确实这样做了——而且还威胁他们，若不服从祂的命令，祂就将他们视为公开的敌人？而另一方面，祂的圣子、\Person{纳匝肋那人}却颁布了与此完全相反的法律，宣称喜爱权力、财富或荣耀的人不能到圣父那里去；人不应比乌鸦更关心食物；应比百合花更不关心衣服；那打他们一记耳光的人，他们应递上另一边？到底是\Person{梅瑟}还是\Person{耶稣}教导错误？圣父派遣\Person{耶稣}时，是否忘了祂曾给\Person{梅瑟}的诫命？还是祂改变了心意，谴责了自己的法律，派来一位传达相反指示的使者？}\Person{赛尔苏斯}尽管自夸无所不知，但在这里陷入了最粗俗的错误，以为法律和先知中除了字面意思之外没有更深的意义。他不明白，奖赏正义生活以世俗财富是多么明显的不可信，因为最常见的是最优秀的人生活在极度贫困中。的确，那些因其纯洁生活而领受天主圣神的先知们，\BibleQuote{披着绵羊皮、山羊皮，到处漂流，受穷乏、患难、苦害，在旷野、山岭、山洞、地穴漂流。}因为，正如圣咏作者所说：\BibleQuote{义人的苦难何其多。}如果\Person{赛尔苏斯}读过\Person{梅瑟}的著作，我猜想，他会认为当对遵守法律的人说\BibleQuote{你要借给许多民族，你自己却不要借入}时，这应许是给予义人，使他世俗财富如此丰裕，不仅能借给犹太人，不仅借给两三个民族，而是\Quote{借给许多民族}。那么，按法律因义德而获得的义人财富该有多大，才能借给许多民族？我们是否也必须按照此解释，认为义人永远不会借入任何东西？因为经上记着：\BibleQuote{你自己什么也不要借入。}那么，那个民族竟如此长期坚持\Person{梅瑟}所教导的宗教，而按照\Person{赛尔苏斯}的假设，他们看到自己被那位立法者如此严重地欺骗了吗？因为从未提到任何人富有到能借给许多民族。如果人民按照\Person{赛尔苏斯}的理解来理解法律，他们不会如此热心地保卫一个其应许明显虚假的法律。如果有人要说，记载的人民所犯的罪证明他们藐视法律，无疑是出于他们觉得被法律欺骗了的感受，我们可以回答，只需读一读那段时期的历史，就能发现全体人民在主眼中行了恶之后，又回归了他们的本分，以及法律所规定的宗教。\par

% structure: p0040 source=h2 target=subsection reason=relative-heading-rank; base=h2

\subsection{第十九章}

如今，如果法律中的这些話，\BibleQuote{你要统治许多民族，没有一人管辖你}，仅仅是给他们统治的许诺，并且没有更深的意义，那么民众当然更有理由轻视法律的应许。刻尔苏斯又提出另一段经文，虽然他改变了其中的措辞，说全地要被希伯来民族充满；其实按历史的见证，这确实在基督来临后发生了，但与其说是天主的祝福，不如说是天主的愤怒。关于赐给犹太人要击杀他们仇敌的应许，可以回答说：任何人如果仔细考察这段经文的意义，会发现他无法按字面解释。目前只需提及圣咏集中义人所说的话，其中说：\BibleQuote{每早晨我要消灭地上的恶人，从上主的城剪除一切作孽的人。}那么，请根据说话者的话语和精神来判断：是否可能，他在圣咏的前一部分（任何人都可以读到）说出了最高尚的思想和意图之后，随后，按他话语的字面解释，他说要在早晨、一天中的其他时间，他要消灭地上所有的罪人，不留下一个活人，并且他要杀死耶路撒冷中所有作孽的人。法律中还有许多类似的表达，例如这句：\BibleQuote{我们没有留下任何活物。}\par

% structure: p0042 source=h2 target=subsection reason=relative-heading-rank; base=h2

\subsection{第二十章}

刻尔苏斯又补充说，曾预言给犹太人，如果他们不遵守法律，他们将遭受与对待仇敌同样的对待；然后他从基督的教导中引用了一些诫命，他认为这些诫命与法律的诫命相悖，并以此作为反对我们的论据。但在讨论这一点之前，我们必须先谈前面的内容。因此，我们主张法律具有双重意义——一是字面意义，一是属神意义——正如我们之前的一些人所表明的。关于第一种或字面意义，不是我们，而是天主，通过一位先知的口说：\BibleQuote{那些律例不好，那些典章不好；}然而，从属神意义理解，同一位先知使天主说：\BibleQuote{祂的律例是好的，祂的典章是好的。}显然，先知说的并非自相矛盾。保禄也同样说：\BibleQuote{文字叫人死，精神却叫人活，}用\Quote{文字}指代字面意义，用\Quote{精神}指代圣经的属神意义。因此，我们无论在保禄还是在先知那里，都能找到表面的矛盾。事实上，如果厄则克耳在一处说：\BibleQuote{我给了他们不好的诫命，不能使人得生的典章，}在另一处又说：\BibleQuote{我给了他们好的诫命和典章，人若遵行，就必因此得生，}保禄同样，当他想要贬低按字面理解的法律时，说：\BibleQuote{若那属死的职事，刻在石头上，尚且带有荣光，甚至以色列子民因梅瑟面容的荣光不能定睛看他，这荣光原是渐渐褪去的；那么，属神的职事岂不是更有荣光吗？}但在另一处，当他想要赞扬和推荐法律时，他称法律是\Quote{属神的}，并说：\BibleQuote{我们知道法律是属神的；}以及\BibleQuote{所以法律是圣洁的，诫命也是圣洁、公义、良善的。}\par

% structure: p0044 source=h2 target=subsection reason=relative-heading-rank; base=h2

\subsection{第二十一章}

那么，当法律的文字向义人应许财富时，刻尔苏斯可能追随那叫人死的文字，将其理解为世俗的、蒙蔽人眼的财富；但我们却认为，它指的是那些照亮眼睛的财富，使人在\BibleQuote{一切言论和一切知识上}富足。在这个意义上，我们\BibleQuote{嘱咐今世富足的人，不要自高，也不要倚靠无定的财富；只需倚靠永生的天主，祂厚赐百物给我们享受；又要他们行善，在善工上富足，甘心施舍，乐意与人分享。}因为，正如撒罗满所说，\BibleQuote{财富是买赎人生命的赎价；}但与此种财富相对的贫穷却是毁灭性的，因为\BibleQuote{穷人不能忍受责备。}关于财富所说的，也适用于统治权，对此经上说：\BibleQuote{义人追赶一千人，二人追赶一万人。}如果财富是按我们刚刚解释的意义来理解，那么请想一想，难道不是按照天主的应许，那在一切言论、一切知识、一切智慧、一切善工上富足的人，可以从他言论、智慧和知识的宝库中借给许多民族吗？保禄就是这样借给他所拜访的所有民族，\BibleQuote{从耶路撒冷带着基督的福音，直到依里黎古。}由于神圣的知识是藉启示赐给他的，他的心智被圣言光照，所以他本人不需要向任何人借取，也不需要任何人的职事来教导他真理之言。因此，正如经上所写：\BibleQuote{你要统治许多民族，他们没有管辖你，}他统治了他带到耶稣基督教导下的外邦人；他从来没有\BibleQuote{向人屈服，一时一刻也没有，}因为他自己比他们更有力量。就这样，他也\Quote{充满了全地。}\par

% structure: p0046 source=h2 target=subsection reason=relative-heading-rank; base=h2

\subsection{第二十二章}

如果我现在必须解释义人如何\Quote{杀死他的仇敌}并处处得胜，就应注意：当他说\BibleQuote{每早晨我要消灭地上的一切恶人，为要从上主的城里剪除所有作孽的人}时，所谓\Quote{地上}是指肉身，其私欲与天主为敌；所谓\Quote{上主的城}是指他自己的灵魂，其中有天主的殿，包含着对天主的真实认识和概念，使之受到所有观看者的赞叹。因此，一旦正义的太阳的光芒照进他的灵魂，他感到被其影响所加强和振作，便着手消灭所有肉身的私欲，这些私欲被称为\Quote{地上的恶人}，并从那位于他灵魂中的主的城里驱逐所有行不义的思想和所有与真理相悖的暗示。同样地，义人也把他们所有的仇敌——即他们的恶习——交付毁灭，甚至不放过婴孩，即邪恶的早期萌芽和冲动。我们也以此理解第137篇圣咏的语言：\BibleQuote{将要被毁灭的巴比伦城啊，那照你所待我们的报复你的，是有福的；那抓住你的婴孩摔在石头上的，是有福的。}因为巴比伦（意即混乱）的\Quote{婴孩}是灵魂中出现的那些烦扰的罪恶思想，而人若击打它们，仿佛将它们的头撞在理性和真理的坚固磐石上，他就是那\Quote{将婴孩摔在石头上}的人，因而他真是有福的。因此，天主可能命令人彻底消灭他们所有的恶习，甚至在它们初生之时，而绝没有命令任何与基督的教导相悖之事；并且祂自己也可能在那些是\Quote{内心的犹太人}的人眼前，将邪恶的所有后裔作为祂的仇敌消灭。同样地，那些不服从天主的法律和言语的人，可以恰当地比作被罪引入歧途的祂的仇敌；他们可以说要遭受与他们所应得的相同的命运，因为他们已成为天主真理的叛徒。\par

% structure: p0048 source=h2 target=subsection reason=relative-heading-rank; base=h2

\subsection{第二十三章}

从以上所述，显然可见，\Person{耶稣}——\Quote{纳匝肋人}——并没有颁布与刚才所考虑的关于财富的法律相对立的诫命，当祂说\BibleQuote{富人进入天主的国是难的}时；无论我们以最直接的意义理解\Quote{富人}一词，指其心思被财富所扰乱、如同被荆棘缠住以致不能结出属灵果实的人；还是指那充满虚假观念而富有的人，关于此人在\BibleBook{}{箴言}中记载：\BibleQuote{义人贫穷，胜于富而不义的人。}也许正是以下经文让\Person{塞耳苏斯}以为耶稣禁止祂的门徒追求野心：\BibleQuote{你们中间谁愿为首，就必作众人的仆人}；\BibleQuote{外邦人的君王主宰他们}，且\BibleQuote{那操权管治他们的被称为恩主}。但这些与\BibleQuote{你要管辖许多民族，他们却不能管辖你}的应许并无矛盾，尤其在我们对这话作出解释之后。\Person{塞耳苏斯}接着提到一句关于智慧的话，似乎他认为，按基督的教导，没有智慧人能到圣父那里去。但我们要问，他所说的\Quote{智慧人}指什么。因为如果他指的是那被称为\Quote{此世的智慧}、即\Quote{在天主前是愚妄}的智慧的人，那么我们就同意他的说法，即这种智慧人被拒绝接近圣父。但如果有人以智慧指基督，祂是\Quote{天主的德能和智慧}，那么这样的智慧人非但不会被拒绝接近圣父，我们反而认为，那被圣神以\Quote{智慧的言语}这一恩赐所装饰的人，远胜于所有未领受同一恩宠的人。\par

% structure: p0050 source=h2 target=subsection reason=relative-heading-rank; base=h2

\subsection{第二十四章}

我们坚持认为，追求人的荣耀不仅被\Person{耶稣}的教导所禁止，也被旧约所禁止。因此我们发现有位先知，当他在为自己因犯某些罪而施加某种惩罚时，便将世上的荣耀包括在这些惩罚之中。他说：\BibleQuote{上主我的天主，我若行了这事，若我手中有不义；我若以恶报那与我和好的人（是的，我反而交付了那无故与我为敌的人），就任凭仇敌追迫我的灵魂，捉拿它；任凭他将我的生命踏在地上，并使我的荣耀置于高处。}而我们的主这些训诫：\BibleQuote{不要为吃什么、喝什么而忧虑。你们看空中的飞鸟，或看乌鸦：它们不种不收，你们的天父尚且养活它们。你们比它们贵重得多！你们为何为衣裳忧虑呢？你们观察田野的百合花；}——这些训诫及其后的内容，与法律所应许的福分——即义人\BibleQuote{要吃他们的粮食而得饱足}——并无矛盾；也与\Person{撒罗满}所说的：\BibleQuote{义人吃得灵魂满足，恶人的肚腹却必缺欠。}并无矛盾。因为我们必须将法律中所应许的食物理解为灵魂的食物，它要满足的不是人性的两部份，而只是灵魂。福音的话，虽然可能包含更深的意义，却也可以按其更简单明显的含义理解，即教导我们不要因对饮食和衣着的焦虑而烦恼，而应生活俭朴，只追求必需的，信靠天主的眷顾。\par

% structure: p0052 source=h2 target=subsection reason=relative-heading-rank; base=h2

\subsection{第二十五章}

\Person{才尔苏}随后从福音中摘出诫命：\BibleQuote{谁若打你一次，你当让他再打你一次}，尽管他没有给出任何他视为与之相反的旧约段落。一方面，我们知道\BibleQuote{古时曾对他们说：以眼还眼，以牙还牙}；另一方面，我们读到：\BibleQuote{我对你们说：凡打你这一边的脸的，连另一边也转给他}。但既然有理由相信\Person{才尔苏}是提出了他从那些想要区别福音的天主与法律的天主之人那里听到的反对意见，我们必须回答说，这一诫命：\BibleQuote{凡打你这一边的脸的，连另一边也转给他}在更早的圣经中并非未知。因为如此，在\WorkTitle{耶肋米亚哀歌}中，他说：\BibleQuote{人最好是自幼背负上主的轭：他独坐、静默，因为他背负了它。他任人打他的脸颊，他受尽凌辱}。因此，福音的天主与法律的天主之间并无不一致，即使按字面理解关于打脸的诫命。所以，我们推断，\Quote{耶稣和梅瑟都没有虚假教导}。圣父派遣耶稣时，并没有\Quote{忘记祂曾赐予梅瑟的命令}；祂并没有\Quote{改变心意，谴责自己的法律，并借着祂的使者送来相反的指示}。\par

% structure: p0054 source=h2 target=subsection reason=relative-heading-rank; base=h2

\subsection{第二十六章}

然而，如果我们必须简要提及梅瑟古时给予犹太人的宪制与基督徒在基督教导下如今愿建立的宪制之间的区别，我们要指出，梅瑟的立法，若按字面理解，不可能与外邦人的召叫以及他们服从罗马政府相协调；另一方面，犹太人也不可能在假定他们接受福音的情况下，保持其民事体制不变。因为基督徒不能杀戮他们的敌人，或如梅瑟所命令的那样，将违反法律并因此被判应受这些惩罚的人处以火刑或石刑；因为犹太人自己，尽管多么渴望执行他们的法律，却不能施行这些刑罚。但在古代犹太人的情况下，他们有自己的土地和政府形式，如果剥夺他们向敌人作战、为国家而战、处死或以其他方式惩罚通奸者、杀人犯或其他犯有类似罪行者的权利，那么当敌人来攻击他们时，就会使他们遭受突然的彻底毁灭；因为他们的法律在那时会约束他们，阻止他们抵抗敌人。而那曾赐予法律、如今又赐予耶稣基督福音的同一上智安排，不愿犹太国继续存在，已摧毁了他们的城邑和圣殿：它废除了在圣殿以牺牲祭献和其他祂所规定的礼仪向天主奉行的敬礼。既然它摧毁了这一切，不愿它们继续存在，同样它也日复一日地扩展基督信仰，以致如今它被勇敢地四处宣讲，尽管有许多障碍反对基督教导在世界的传播。但由于天主的旨意是万民接受基督教导的益处，所有反对基督徒的人的计划都化为乌有；因为各地君王、统治者和人民越是迫害他们，他们的人数就越增加，力量就越强大。\par

% structure: p0056 source=h2 target=subsection reason=relative-heading-rank; base=h2

\subsection{第二十七章}

此后，\Person{才尔苏}详细陈述了他归给我们但我们并不持有的关于神体的观点，大意是\Quote{祂的本性是有形体的，并拥有如人一般的身体}。既然他着手反驳那些并非我们的观点，就没有必要给出这些观点本身或对其的反驳。的确，如果我们确实持有他归给我们并反对的那些关于天主的见解，我们就得引用他的话，提出我们自己的论据，并反驳他的观点。但是，如果他提出的是他从未听人说过，或假定他听说了，那一定是来自那些非常单纯、不明白圣经意义的人，那么我们就不必从事如此多余的任务去反驳它们。因为圣经清楚地讲到天主是没有身体的存有。因此经上说：\BibleQuote{从来没有人见过天主}；而且万物的首生者被称为\BibleQuote{不可见的天主的肖像}，这等于说祂是无形的。然而，我们在考察\BibleQuote{天主是神，朝拜祂的人应当以心神以真理朝拜祂}这句话的意义时，已经对天主的本性有所论述。\par

% structure: p0058 source=h2 target=subsection reason=relative-heading-rank; base=h2

\subsection{第二十八章}

\Person{凯尔苏斯}如此歪曲我们对天主本性的看法后，接着问我们\Quote{死后希望去哪里}；他让我们回答说\Quote{去一个比这更好的另一片土地}。对此他评论如下：\Quote{古代的神圣人物曾讲过为有福灵魂保存的福乐生命。有人称之为\InnerQuote{福岛}，有人称之为\InnerQuote{厄吕息翁原野}，如此称呼是因为他们将在那里脱离现今的灾祸。因此\Person{荷马}说：\InnerQuote{神明将送你们到厄吕息翁原野，在大地的边界，那里他们过着极安静的生活。}\Person{柏拉图}也相信灵魂不死，明确称灵魂所去的地方为\InnerQuote{土地}。他说：\InnerQuote{其广大无边，我们只占据一小部分，从\Place{法西斯河}到\Place{赫丘利柱}，我们沿着海岸居住，像沼泽旁的蚱蜢和青蛙。但还有许多其他地方同样被其他人居住。因为在大地各处有洞穴，形状和大小各异，水、云和空气流入其中。但那片纯净的土地位于纯净的天域。}}因此\Person{凯尔苏斯}以为，我们所说的比这片土地更美好更卓越的土地，是从某些他称为\Quote{神圣}的古代作者那里借用来的，尤其从\Person{柏拉图}而来，他在\WorkTitle{斐多篇}中谈论了位于纯净天域中的纯净土地。但他没有看到，远较希腊文学为古老的\Person{梅瑟}，曾应许那些按祂法律生活的人那圣地，就是\BibleQuote{一片美好而辽阔的土地，流奶流蜜之地}；这应许并不像某些人以为的那样，是指我们称为\Place{犹太地}的那部分大地；因为尽管它美好，它仍是有原罪之\Person{亚当}因过犯而受诅咒的大地的一部分。因为这句话：\BibleQuote{为了你所做的，大地受诅咒；你一生日日要凭忧愁，即劳苦，吃它的果实}，是对整个大地说的，每个在\Person{亚当}内死去的人都一生忧愁或劳苦地吃它的果实。由于整个大地都受了诅咒，它就在那些因\Person{亚当}被赶出乐园的人的一生中长出荆棘和蒺藜；人人必汗流满面才得糊口，直到归回所出的泥土。要详尽解释这一段所包含的一切，可以说很多；但我们目前只限于这几句话，旨在消除一种观念，即认为天主应许给义人的美好土地是指\Place{犹太地}。\par

% structure: p0060 source=h2 target=subsection reason=relative-heading-rank; base=h2

\subsection{第二十九章}

因此，如果整个大地因\Person{亚当}及在他内死去之人的行为而受了诅咒，那么显然大地各部分都分受诅咒，\Place{犹太地}也在其中；所以\BibleQuote{一片美好而辽阔的土地，流奶流蜜之地}这句话不能指它，尽管我们可以说，\Place{犹太地}和\Place{耶路撒冷}是那片纯净天堂中美好辽阔纯净土地的影子和预像，其中就是天上的\Place{耶路撒冷}。正是论及这\Place{耶路撒冷}，宗徒说话，如同一个\BibleQuote{与基督一同复活，寻求上面的事}的人，发现了不属于犹太神话的真理。他说：\BibleQuote{你们来到了\Place{熙雍山}，来到了永生天主的城，天上的\Place{耶路撒冷}，以及千万天使的集合。}为了确信我们对\Person{梅瑟}所说的\Quote{美好辽阔的土地}的解释不违背天主圣神的旨意，我们只需读所有先知关于那些离开\Place{耶路撒冷}、从它迷失、后来又回来定居在那称为天主居所和城市的地方的人所说的，如：\BibleQuote{祂的居所在圣所}；以及\BibleQuote{在我们天主的城中，在祂圣山上，上主伟大，应受极度的赞美；这城巍峨壮丽，是全地的喜悦。}现在只需引用\BibleBook{圣咏}{37}的话，它论及义人的土地这样说：\BibleQuote{仰望上主的人必继承大地}；稍后：\BibleQuote{但温良的人必继承大地，并因丰盛的平安而欢乐}；又说：\BibleQuote{祝福祂的人必继承大地}；以及\BibleQuote{义人必继承土地，永远居住其中。}请思考，以下同篇的话是否对聪明的读者而言显然指纯净天堂中的纯净土地：\BibleQuote{仰望上主，遵行祂的道路，祂必高举你，使你继承土地。}\par

% structure: p0062 source=h2 target=subsection reason=relative-heading-rank; base=h2

\subsection{第三十章}

在我看来，柏拉图的那个想象——即那些我们称为宝石的石头，其光泽来源于，可以说是，对那更好之地中石头的反射——也是取自依撒意亚描述天主之城的话语：\BibleQuote{我必使你的城垛为碧玉，你的石头为水晶，你的边界为宝石}；以及\BibleQuote{我必以蓝宝石奠定你的根基}。那些最尊敬柏拉图学说的人，将他的这个神话解释为寓言。而那些预言——我们推想柏拉图从中借用了——将由那些度先知般的虔诚生活、全心致力于研读圣经的人，向那些因纯洁生活和渴望神圣知识而配得学习的人解释。就我们而言，我们的目的只是要说：我们关于那圣地的肯定并非取自柏拉图或任何希腊人，而是他们——因为不仅活在最古老的梅瑟之后，甚至活在大多数先知之后——反而从他们那里借用，并且在此过程中，要么误解了他们对这类主题的隐晦暗示，要么在影射那更好之地时，试图模仿已落入他们手中的那部分圣经。哈盖明确区分了地与旱地，他将后者理解为我们所居住的地。他说：\BibleQuote{我必再一次震动天、地、旱地和海洋}。\par

% structure: p0064 source=h2 target=subsection reason=relative-heading-rank; base=h2

\subsection{第三十一章}

\Person{策尔苏斯}提到\Person{柏拉图}的\WorkTitle{斐多篇}中的段落，说：\Quote{并非人人都容易理解柏拉图的话，他说由于我们的软弱和迟钝，我们无法达到空气的最高区域；但若我们的本性能够进行如此崇高的观想，那么我们就能理解那才是真正的天，那才是真正的光。}由于策尔苏斯已推迟另寻机会解释柏拉图的思想，我们也认为目前不适合详细描述那圣善美好的地和其中的天主之城；而将对此的思考保留给我们的\WorkTitle{先知书注释}，尽管我们已以自己力所能及的方式，在对\BibleBook{圣咏}{第四十六篇}和\BibleBook{圣咏}{第四十八篇}的评注中部分地讨论了天主之城。梅瑟和先知们的著作——所有书籍中最古老的——教导我们，此处世上人类日常使用的一切事物，都有在名称上与之对应、而唯独为真实的其他事物。例如，有真正的光，有穹苍之外的另一个天，有一个不同于我们所见的太阳的正义的太阳。总之，为了区分这些事物与感官对象——它们没有真正的真实性——他们论及天主说：\BibleQuote{祂的作为是真理}；由此区别天主的作为与天主手的作为，后者是较低类别的。因此，天主在依撒意亚中抱怨人们说：\BibleQuote{他们不看重主的作为，也不关心祂手的化工}。关于这一点已经足够。\par

% structure: p0066 source=h2 target=subsection reason=relative-heading-rank; base=h2

\subsection{第三十二章}

\Person{策尔苏斯}接下来攻击复活的教义，这是一项高深而艰难的教义，比其他教义更需要高超而深邃的智慧，才能阐明它是如何相称于天主的；并阐明这是一项多么崇高的真理，它教导我们，有一种生殖原理存在于圣经称之为灵魂的\Quote{帐棚}中，义人就在其中\BibleQuote{叹息，受压，并非愿意脱去，而是愿意穿上}。策尔苏斯嘲笑这教义，因为他不理解它，并且他是从无知的人那里学来的，那些人无法以任何合理的理由支持它。因此，除了我们以上所说的，再作一点说明将是有益的。我们关于复活的教导，并非如策尔苏斯所想的那样，来自我们听过的关于\OriginalTerm{灵魂转生}{metempsychosis}的教义；而是我们知道，灵魂就其本性而言是非物质、不可见的，它不能没有一具适合其所在处所本性的身体而存在于任何物质之地。因此，它时而脱去一个先前必要、但已变化的状态中不再足够的身体，换上一个第二身体；时而它再取另一个身体作为前者之外的附加，需要它作为更好的遮蔽，适合天上更纯净的以太区域。当它出生来到世界时，它脱去在母腹中需要的包覆；而在做此之前，它穿上另一适合尘世生活的身体。然后，正如有\Quote{一个帐棚}和\Quote{一个地上的房子}在某种程度上为这帐棚所需，圣经教导我们说\BibleQuote{这帐棚的地上的房子将要被拆毁}，但这帐棚将要\BibleQuote{穿上天上非人手所造、永存的房屋}。天主的人还说：\BibleQuote{这可朽坏的将穿上不朽坏性}，这与\Quote{不朽坏的}不同；\BibleQuote{这必死的将穿上不死性}，这与\Quote{不死的}不同。的确，正如\Quote{智慧}之于\Quote{智慧者}、\Quote{正义}之于\Quote{正义者}、\Quote{和平}之于\Quote{和平者}那样，\Quote{不朽坏性}之于\Quote{不朽坏的}、\Quote{不死性}之于\Quote{不死的}也是如此。看哪，圣经鼓励我们期待怎样的前景，当它对我们谈论穿上不朽坏性和不死性时，它们就像是外衣，不会让那些被它们覆盖的人遭遇朽坏或死亡。到此为止，我冒昧地提及这个主题，以回答一位不了解而攻击复活教义的人，他仅仅因为对此一无所知，就使之成为轻蔑和嘲笑的对象。\par

% structure: p0068 source=h2 target=subsection reason=relative-heading-rank; base=h2

\subsection{第三十三章}

既然\Person{凯尔苏斯}以为我们坚持复活教义是为了能看见并认识天主，他便对此议题作出如下推论：\Quote{他们被彻底驳倒和征服之后，仍然仿佛不顾一切反对，又回到同一个问题：\InnerQuote{那么我们如何看见并认识天主呢？我们如何到祂那里去？}} 但任何愿意听我们的人要留意：如果我们为了其他目的需要身体，例如占据一个物质的空间，身体必须适应那个空间，因此\Quote{帐篷}以我们所示的方式被披上，那么为了认识天主，我们并不需要身体。因为看见天主的不是肉身的眼睛，而是按造物主肖像受造的心灵，天主以其眷顾使之有能力认识天主。看见天主属于纯洁的心，从这样的心不再发出\Quote{恶念、凶杀、奸淫、淫乱、偷盗、妄证、亵渎、恶眼}或任何其他恶事。因此经上说：\BibleQuote{心里洁净的人是有福的，因为他们要看见天主。}但由于我们意志的力量不足以获得完全洁净的心，我们需要天主创造它，因此那正确祈祷的人向天主献上这个祈求：\BibleQuote{天主，求你给我创造一个纯洁的心。}\par

% structure: p0070 source=h2 target=subsection reason=relative-heading-rank; base=h2

\subsection{第三十四章}

我们并不这样问：\Quote{我们如何到天主那里去？}好像我们认为天主存在于某个地方。天主的本性极为崇高，不为任何空间所限：祂掌管万有，自己却不受任何事物的束缚。因此，诫命\Quote{你要跟随上主你的天主}并不是命令肉身的接近；先知在他的祈祷中\BibleQuote{我的灵魂紧紧跟随你}也不是指身体上的接近。所以\Person{凯尔苏斯}歪曲了我们，他说我们期望用肉身的眼睛看见天主、用耳朵听见祂、用手感性地触摸祂。我们知道圣经提到眼睛、耳朵和手，这些名称与肉身器官只有名称相同，而更奇妙的是，它们提到一种更神圣的官能，与通常说的感官截然不同。因为当先知说\BibleQuote{求你开我的眼睛，使我看出你法律中的奇妙}，或\BibleQuote{上主的诫命是纯洁的，能光照眼睛}，或\BibleQuote{求你光照我的眼睛，免得我沉睡至死}，没有一个人会愚蠢到以为肉身之眼能看见神圣法律的奇妙，或上主的法律能光照肉身的眼睛，或死亡的睡眠落在肉身之眼上。当我们的救主说\BibleQuote{有耳听的，听吧！}，任何人都明白所说的耳朵是更神圣的。当说到上主的话在耶肋米亚或某位先知\Quote{手中}，或使用\Quote{藉梅瑟之手颁布的法律}，或\Quote{我用手寻求上主，没有受骗}，没有一个人愚蠢到不明白\Quote{手}是比喻的说法，正如若望所说\BibleQuote{我们的手曾摸过的生命之言}。如果你还愿意从圣书中学习到有一种比肉身感官更神圣的官能，只需听撒罗满所说：\Quote{你将找到一个神圣的官能。}\par

% structure: p0072 source=h2 target=subsection reason=relative-heading-rank; base=h2

\subsection{第三十五章}

因此，用这种方式寻求天主，我们不需要去访问\Person{特洛福尼乌斯}、\Person{安菲阿拉俄斯}和\Person{摩普索斯}的神谕所，\Person{凯尔苏斯}要我们去那里，向我们保证我们会在那里\Quote{看见以人形出现的神明，清清楚楚，毫无幻象。}因为我们知道这些是魔鬼，以牺牲的血、烟和气味为生，被卑劣的欲望囚禁在希腊人称为神庙的监狱里，但我们知道这些不过是欺骗人的魔鬼的住所。\Person{凯尔苏斯}恶意地补充说，关于这些在他看来人形的神明，\Quote{他们并非只显现一次或偶尔显现，像那欺骗人的一样，而是永远与渴慕交往的人敞开交往。}从这句话看来，\Person{凯尔苏斯}似乎以为基督在复活后向门徒显现如同一个掠过眼前的幽灵；而他所谓人形的神明永远向渴慕的人显现。但是，一个像他描述的那样飞来飞去欺骗观看者的幻影，怎能在它消失之后产生如此大的效果，使人心如此转变，以致按照天主的旨意规范自己的行为，仿佛自觉将来要受祂审判？一个幻影怎能驱赶魔鬼，并显示其他无可争辩的能力证据，而且不仅仅在一个地方（像这些人形的所谓神明那样），而是让自己的神圣能力传遍全世界，吸引并聚集所有愿意度善良高贵生活的人？\par

% structure: p0074 source=h2 target=subsection reason=relative-heading-rank; base=h2

\subsection{第三十六章}

在\Person{策尔苏斯}说了以上这些我们尽力回答的话之后，他继续谈道，论到我们说：\Quote{他们又要问：\InnerQuote{除非通过感官的感知，我们怎能认识天主？因为除了通过感官，我们如何能获得任何知识？}}他对此回答道：\Quote{这不是人的语言；它不是来自灵魂，而是来自血肉。如果这些没有精神、属于血肉的种族能够做到的话，就让他们听我们说：如果你们不运用感官，而是以灵魂向上仰望；如果你们转离身体的眼睛，开启心灵的眼睛，唯有如此，你们才能看见天主。如果你们寻找一位向导引导你们走上这条路，就必须躲避所有那些欺骗者和变戏法的人，他们会把你们引向幻影。否则，你们将扮演最可笑的角色：你们一边咒骂那些被公认为神的神祇，把祂们当作偶像，一边却向一个比这些更可怜的偶像致敬——这个偶像甚至不是偶像或幻影，而是一个死人，并且你们寻找一个像祂一样的父亲。}关于这段文字，我们首先要指出的是，他用了拟人法，让我们以这种方式为复活教义辩护。当被引入的人物性格和思想被忠实地保留时，这种修辞手法是得当的；但若不符合该说话者的性格和观点，便是滥用。因此，我们会公正地谴责一个人把哲学术语塞进野蛮人、奴隶或未受教育的人口中，因为我们知道这些哲学属于作者，而不属于那些根本不懂哲学的人。同样，我们会谴责一个人引入被描绘为智慧且精通神圣知识的人物，却让他们说出只能出自无知者或受世俗激情影响者之口的话。因此，\Person{荷马}之所以受到称赞，原因之一就在于他笔下的人物性格保持一致，如\Person{内斯特}、\Person{尤利西斯}、\Person{狄俄墨得斯}、\Person{阿伽门农}、\Person{特勒马科斯}、\Person{佩涅洛佩}等。相反，\Person{欧里庇得斯}在\Person{阿里斯托芬}的喜剧中被抨击为轻浮的说教者，常常把从\Person{阿那克萨哥拉}或其他哲学家那里学来的至理箴言放到一个野蛮女人或可怜的奴隶口中。\par

% structure: p0076 source=h2 target=subsection reason=relative-heading-rank; base=h2

\subsection{第三十七章}

如果以上是对拟人法正确与错误使用的准确描述，那么我们是否有理由嘲笑\Person{策尔苏斯}把基督徒从未说过的话加在他们头上？因为，如果他让说话的人是无学问的人，那么这样的人怎么可能区分\Quote{感官}和\Quote{理性}、\Quote{感官对象}和\Quote{理性对象}？要以这种方式辩论，他们必须曾在斯多亚学派门下学习过，因为斯多亚学派否认一切理智实体的存在，主张我们所领悟的一切都是通过感官获得的，而且一切知识都来自感官。但是，另一方面，如果他把这些话放进那些仔细探究基督教教义含义的哲学家口中，那么这些陈述就不符合他们的性格和原则。因为，凡学过天主是看不见的，并且祂的某些作品是看不见的（即由理性领会）的人，是不可能说这样的话的，仿佛是为了证明他们对复活的信仰：\Quote{除非通过感官的感知，他们怎能认识天主？}或\Quote{除了通过感官，他们如何能获得任何知识？}因为这些话并不是写在只有少数智者才读的秘密著作中，而是写在那最为广泛传播、最为民众所熟知的著作中：\BibleQuote{天主那看不见的事，从创世以来，就藉着所造之物被清楚看见，为人所领悟。}由此可以推断，尽管生活在世上的人必须首先运用感官感知可感对象，以便由此进展到对理智事物本性的认识，然而他们的知识不能停留在感官对象上。因此，虽然基督徒不会说没有感官就不可能认识理智对象，反而会说感官提供了获得知识的第一种途径，他们完全可以问这个问题：\Quote{没有感官，谁能获得任何知识？}而不该承受\Person{策尔苏斯}的辱骂，当他接着说：\Quote{这不是人的语言；它不是来自灵魂，而是来自血肉。}\par

% structure: p0078 source=h2 target=subsection reason=relative-heading-rank; base=h2

\subsection{第三十八章}

由于我们坚持伟大的天主本质上是单纯的、不可见的和非物质的，祂是纯粹理智，或超越理智与存在的某物，我们绝不能断言天主能通过任何其他方式被领悟，唯有通过按祂肖像而形成的理智，尽管现在，按保禄的话说，\BibleQuote{我们如今彷佛对着镜子观看，模糊不清，但那时就要面对面了。}如果我们使用\Quote{面对面}这个说法，愿无人曲解其意；而要借这段经文来解释：\BibleQuote{我们众人以揭开的脸面反映主的光荣，逐渐被改变成同一肖像，光荣上加光荣。}这表明我们在此语境中不使用这个词来指可见的脸庞，而是象征性地使用，正如我们已经表明眼睛、耳朵和其他身体部位的用法一样。而且，一个男人——我是指一个使用身体的灵魂，也称作\Quote{内在的人}，或简称为\Quote{灵魂}——他回答的方式，不是塞尔苏斯让我们回答的那样，而是天主自己的人所教导的。同样，一个基督徒不会使用\Quote{肉身的语言}，因为他已学会\BibleQuote{靠着圣神治死身体的行为}，并\BibleQuote{在身上常常带着耶稣的死}；也称\BibleQuote{治死你们在地上的肢体}，并且真正明白这些话：\BibleQuote{我的神不永远与人相争，因为人也是血肉}，以及\BibleQuote{属血肉的人不能讨天主的喜欢}，他竭力从各方面不再按血肉生活，而只按圣神生活。\par

% structure: p0080 source=h2 target=subsection reason=relative-heading-rank; base=h2

\subsection{第三十九章}

现在让我们听听他邀请我们学习什么，以便我们可以从他那里确定我们应如何认识天主，尽管他认为他的话超出所有基督徒的能力。他说：\Quote{让他们听，}\Quote{如果他们有这个能力的话。}但他并没有像他应该的那样教导我们，反而辱骂我们；他本应在演讲开始时对他所称呼的人表示善意，却称那些宁愿死也不愿用一个词背弃基督信仰、并准备忍受各种酷刑或任何死亡的人为\Quote{懦弱的种族}。他还给我们加上\Quote{属肉体的}或\Quote{纵欲的}称号，尽管我们惯于说，\BibleQuote{我们曾按血肉认识过基督，但从今以后不再这样认识祂了。}而且我们如此乐意为了宗教事业舍弃生命，以至于没有哲学家能更轻易地脱去他的长袍。然后他对我们说这话：\Quote{如果不运用你的感官，而是用灵魂向上看；如果转离身体的眼睛，打开心灵的眼睛，这样，只有这样才能看见天主。}他不知道，这种对两种眼睛（身体的眼睛和心灵的眼睛）的引用是他从希腊人那里借来的，在我们的作者中早已使用；因为梅瑟在创造世界的记述中，将犯罪前的人描述为既看见又看不见：看见时，论到女人说，\BibleQuote{女人见那棵树实在好吃，悦人眼目，且是可喜爱的，能使人获得智慧}；又看不见时，当他介绍蛇对女人说，好像她和她的丈夫是瞎子：\BibleQuote{天主知道，你们吃了那果子那一天，你们的眼就会开了}；还有当它说，\BibleQuote{他们吃了，两人的眼就开了。}感官的眼睛那时被打开了，它们本来该闭上，以免分心，妨碍用心灵的眼睛看；而正是心灵的眼睛，由于罪（照我想）那时被关闭了，他们直到那时一直用它们享受看见天主及其乐园的喜乐。我们救主熟悉我们身上的这种双重视觉，祂说：\BibleQuote{我为了审判来到这世界，使看不见的看见，看得见的反而变成瞎子。}——意思是，看不见的眼睛是指心灵的眼睛，因祂的教导而明亮；而看见的眼睛是感官的眼睛，祂的话使它们变瞎，以便灵魂能不受干扰地观看适当对象。因此，所有真正的基督徒都使心灵的眼睛敏锐，并闭上感官的眼睛；这样，每个人根据他更好的眼睛被激活、感官眼睛被暗化的程度，看见并认识至尊天主，以及祂的圣子，即圣言、智慧等。\par

% structure: p0082 source=h2 target=subsection reason=relative-heading-rank; base=h2

\subsection{第四十章}

继我们已评论过的\Person{塞耳苏斯}言论之后，接着是他对所有基督徒的另一些评论，但这些评论倘若适用于任何人，则应针对其教义与\Person{耶稣}所传授的完全相异之人。因为我们先前已指出，正是那些奥菲特派完全背弃了\Person{耶稣}，或许还有一些持有类似见解的人，他们是\Quote{骗子与魔术师，将人引向偶像与幻影}；正是他们以可悲的辛劳背诵天上守门者的名字。因此，这些话用来针对基督徒是完全不恰当的：\Quote{你若寻求一位向导引领这道路，就应避开一切欺骗者和魔术师，他们必会带你走向幻影。}而且，他似乎完全不知道这些骗子与他完全一致，且在诋毁\Person{耶稣}及其宗教方面不亚于他，于是他混淆我们与他们，继续说道：\Quote{否则，你们的行径将极为可笑：你们一面咒骂那些公认的其他神明，视之为偶像，一面却向比它们更卑劣的偶像致敬——那甚至不是偶像或幻影，而是一个死人——并且寻求一位与他自己相似的父。}他完全不明白我们的见解与那些编造这些故事之人的见解有天壤之别，并且他以为他对他们的指责也适用于我们，从以下段落可见：\Quote{为了如此怪异的妄想，为了支持那些奇妙的顾问，以及那些你们向狮子、两栖动物、驴形生物及其他对象所说的奇妙话语，为了那些你们辛劳铭记其名字的神圣守门者，为了如此缘由，你们忍受残酷的折磨，死于火刑。}诚然，他并不知道，那些视驴形、狮形或两栖动物为通往天堂的守门者或向导的人，没有一个会为捍卫他们所认为的真理而赴死。那种使我们为宗教而甘受各种死亡、死于火刑的过度热忱（如果可以这样称呼的话），却被\Person{塞耳苏斯}归给那些从未忍受此类苦难的人；他把我们这些因信仰而被钉十字架的人与相信虚幻生物——狮子、两栖动物及其他怪物——混为一谈。我们摒弃所有这些神话，并非出于对\Person{塞耳苏斯}的顺从，因为我们从未在任何时候持有此类幻想；而是按照\Person{耶稣}的教导，我们反对所有这类观念，不允许\Person{弥额尔}或任何其他被提及者具有那种形状和形态。\par

% structure: p0084 source=h2 target=subsection reason=relative-heading-rank; base=h2

\subsection{第四十一章}

但让我们审视\Person{塞耳苏斯}希望我们跟随的那些人是谁，以便我们不会缺少既因其古老性又因其圣德而受推荐的向导。他将我们引向所谓的受神启发的诗人、贤士和哲学家，却未指名道姓；因此，在承诺指出应引导我们的人之后，他只是笼统地将我们交给受神启发的诗人、贤士和哲学家。倘若他具体列出他们的名字，我们必会向他表明，他希望将那些盲目于真理的人作为向导给我们，从而必然引我们误入歧途；或者，他们即使并非完全盲目，也在许多信仰事务上错误。但无论他所谓的\Quote{受启发的诗人}是指\Person{奥菲斯}、\Person{帕默尼德斯}、\Person{恩培多克勒}，甚或\Person{荷马}本人与\Person{赫西奥德}，请任何人指出，那些跟随他们指引的人如何走在更好的道路上，或过着更卓越的生活，比起那些在\Person{耶稣基督}学校里受教导、摒弃所有像与雕像、甚至摒弃所有犹太迷信，以便通过天主的圣言仰视那唯一的天主——他是圣言之父——的人。那么，那些贤士和哲学家是谁，\Person{塞耳苏斯}希望我们向他们学习如此多的神圣真理，并为了他们而放弃天主的仆人\Person{梅瑟}、创造万物的天主的先知（他们真正受神感发了如此多的事），甚至那光照并教导全人类虔敬之道的那一位——以致若有人仍对他奥秘的知识一无所知，他也不能责备那人？他确实对人怀有如此丰盛的爱，以致赐予更有学识的人一种能将灵魂高举超脱一切世间之物的神学；同时，他同样体贴地俯就无知之人、简单妇女、奴隶，简言之，所有那些唯独从\Person{耶稣}那里获得帮助以更好规范生活的人——这帮助由他对天主的教导所提供，按照他们的需要和能力加以调整。\par

% structure: p0086 source=h2 target=subsection reason=relative-heading-rank; base=h2

\subsection{第四十二章}

\Person{刻尔苏斯}接下来又引述\Person{柏拉图}，认为他是神学真理更有效的导师，并从\WorkTitle{蒂迈欧篇}中引用了以下段落：\Quote{要发现这个宇宙的创造者与圣父，是一件困难的事；即便发现了，也不可能向所有人讲论祂。}接着他自己又加了这样的评述：\Quote{那么，你看到了神性之人如何追求真理之道，也看到\Person{柏拉图}深知并非所有人都能行走在这条道路上。但既然智慧之人已经找到了这条道路，其明确目的正是为了能将关于那第一位、那\OriginalTerm{不可言说的存有}{unspeakable Being}的一些概念传达给我们——一种概念，即通过其他对象向我们呈现祂——他们便努力通过综合（即不同性质的结合）、或通过分析（即某些性质的分离与排除）、或最终通过类比，以这些方式，我再说，他们努力将那不可言词表达的东西呈现于我们面前。如果你也能追随这条道路，我会感到惊讶，因为你如此完全地与肉体结合，以致除了不洁之物以外什么也看不见。}\Person{柏拉图}这些话是崇高且令人赞叹的；但是，请看圣经是否没有给我们一个例子，表明在\OriginalTerm{天主圣言}{God the Word}——祂\BibleQuote{在起初就与天主同在}，并\BibleQuote{成了血肉}——身上有一种对人类更大的关怀，目的是要向所有人启示那些按照\Person{柏拉图}所说即使他本人发现了也不可能向所有人讲论的真理。\Person{柏拉图}可能会说\Quote{要发现这个宇宙的创造者与圣父是一件困难的事}；他用这句话暗示，人性并非完全不可能获得一种或堪当于天主、或至少远超通常所达程度的知识（尽管如果\Person{柏拉图}或任何其他希腊人真的找到了天主，他们就不会向其他神祇致敬和崇拜，也不会将天主的名归于其他对象；他们必会放弃一切别的神明，不会将那些与这位伟大天主毫无共同之处的对象与祂相提并论）。至于我们，我们主张，人性若不寻求，并且没有所寻求的那位的帮助，就根本不可能寻求天主或获得对祂的清晰认识。祂向那些在尽己所能之后承认需要祂帮助的人自我启示；祂向祂所悦纳的人启示自己，其程度是人与仍居身内的灵魂所能认识天主的限度。\par

% structure: p0088 source=h2 target=subsection reason=relative-heading-rank; base=h2

\subsection{第四十三章}

请注意，当\Person{柏拉图}说\Quote{在发现了宇宙的创造者与圣父之后，不可能向所有人讲论祂}时，他并没有说祂是不可言说的、不能用言语表达的。相反，他暗示祂是可以被述说的，并且有少数人可以得知祂。但\Person{刻尔苏斯}似乎忘记了他刚刚引用的\Person{柏拉图}的措辞，立刻给天主起了\Quote{不可言说的}这个名称。他说：\Quote{既然智慧之人找到了这条路，其目的正是为了能给我们关于万有的第一位（祂是不可言说的）的一些概念。}至于我们，我们主张不仅天主是不可言说的，连那些低于祂的事物也是如此。那些正是\Person{保禄}努力要表达的事物，他说：\BibleQuote{我听到了不可言传的言语，是人不可说出来的。}此处\Quote{听到}一词用于\Quote{理解}的意义，正如经文所说：\BibleQuote{有耳朵听的，听吧！}我们还认为要看见宇宙的创造者与圣父是一件困难的事，但藉着以下所提及的方式，是可能看见祂的：\BibleQuote{心地纯洁的人是有福的，因为他们要看见天主。}不仅如此，还要按照那位\BibleQuote{是不可见的天主的肖像}的话的意义：\BibleQuote{看见我的，就是看见那派遣我的父。}没有一个明智之人会认为这些话是指着祂的身体临在说的，因为那是所有人都能看见的；否则那些喊\Quote{钉死祂，钉死祂}的人和那对\Person{耶稣}的人性拥有权柄的\Person{比拉多}，都成了看见天主父的人，这是荒谬的。此外，这些话：\BibleQuote{看见我的，就是看见那派遣我的父}，不应按其粗浅意义来理解，这从祂回答\Person{斐理伯}的话可以清楚看出：\BibleQuote{斐理伯，我同你们在一起这么久，你还不认识我吗？}此前\Person{斐理伯}曾问：\BibleQuote{请把父显示给我们，我们就心满意足了。}因此，谁若明白\BibleQuote{圣言成了血肉}这句话应如何理解天主的独生子、所有受造物的首生者，也将明白我们如何在看见不可见的天主的肖像时，看见了\Quote{宇宙的创造者与圣父}。\par

% structure: p0090 source=h2 target=subsection reason=relative-heading-rank; base=h2

\subsection{第四十四章}

\Person{塞尔苏斯}认为，我们若要获得对天主的认识，要么通过数学家称之为综合与分析的方法，对某些事物进行组合或分离，要么通过他们也使用的类推法，并且以这种方式，我们或许能进入至善之门。但当圣言说：\BibleQuote{除了圣子和圣子所愿意启示的人，没有人认识圣父。}祂宣告，除非借着由上而来的天主的恩宠，并带有某种神圣的默感，没有人能认识天主。的确，我们有理由认为，对天主的认识是人类本性所无法企及的，因此人类在对天主的看法上陷入了许多错误。于是，正是出于天主对人类的美善与慈爱，并借着天主恩宠的奇妙作为，对祂在预知中所看见、并知道他们必会相称于那使他们认识祂的天主、且永不背离对祂服务的忠实依附的人——即使他们被那些因不知真宗教为何物、而将此名加在理应被称为任何事物而非宗教之物的人判处死刑或加以嘲笑——天主无疑看到了那些傲慢自大的人，他们轻蔑一切他人，夸耀自己认识天主，并因哲学而对神圣事物有深刻了解，却仍然不亚于最无知的人，追随他们的偶像、庙宇和著名的神秘仪式；看到这一切，祂\BibleQuote{拣选了世上愚拙的}——那些最单纯的基督徒，然而他们过着比许多哲学家更有节制和更洁净的生活——\BibleQuote{叫有智慧的羞愧}，这些人竟不羞于向无生命的物体祈祷，把它们当作神或神像。因为当看到有人从哲学中学到了关于天主或众神如此深刻而高尚的思想，却转身向偶像祈祷，或以为凝视这些物质之物就能从可见的象征上升到精神的、非物质之物时，哪个理性的人能忍住不笑呢？但一个基督徒，即使是普通百姓，也确信每个地方都是宇宙的一部分，整个宇宙是天主的圣殿。无论他在世界的何处，他都祈祷；但他超越宇宙，\Quote{闭上感官的眼睛，而向上举目，观看灵魂。}他不停留在天穹；而是在圣神的引导下，思想越过诸天，从而如同超出可见宇宙，他向天主献上祈祷。但他不祈求琐碎的祝福，因为他从耶稣那里学会了不求微小或卑贱之物，即可感知的事物，而只求伟大且真正神圣的事物；天主将这些赐予我们，引导我们获得那唯独与祂同在、并通过祂的圣子——圣言——就是天主——才能找到的福乐。\par

% structure: p0092 source=h2 target=subsection reason=relative-heading-rank; base=h2

\subsection{第四十五章}

但让我们进一步看看他提议教导我们的是什么，如果我们真能理解的话，因为他称我们是\Quote{完全与肉身联姻的；}然而，如果我们生活得好，并按照耶稣的教导生活，我们听到这样对我们说：\BibleQuote{你们不属于肉身，而属于圣神，如果天主的圣神住在你们内。}他还说我们看不到任何纯洁之物，尽管我们努力使我们的思想也远离一切罪的玷污，并且在祈祷中我们说：\BibleQuote{天主，求祢在我内造一颗清洁的心，使我心内重有正直的精神，}以便我们能以那颗\BibleQuote{纯洁的心}看见祂，唯有这颗心才赐予看见祂的特权。那么，这是他提供给我们教导的内容：\Quote{事物要么是\Emph{可理解的}，我们称之为实体——存在；要么是\Emph{可见的}，我们称之为\Emph{生成}：前者中存有真理；从后者产生错误。真理是知识的对象；真理与错误构成意见。可理解的对象由理性认识，可见的对象由眼睛认识；理性的活动称为智性知觉，眼睛的活动称为视觉。因此，正如在可见事物中，太阳既不是眼睛也不是视觉，而是使眼睛看见、使视觉成为可能的东西，由于它，可见事物被看见，一切可感事物存在，它自身也成为可见的；同样，在可理解事物中，那既不是理性，也不是智性知觉，也不是知识的东西，却是使理性认识、使智性知觉成为可能的原因；由于它，知识产生，一切可理解的事物、真理本身和实体存在；而它自身，超越这一切，以某种不可言说的方式成为可理解的。这些是提供给有智力者思考的；如果你们也能理解其中任何一点，那就好了。如果你们认为有一位神圣的圣神从天主降临，向人宣布神圣的事物，无疑正是这位圣神启示这些真理，正是在同样的影响下，古人们宣告了许多重要的真理。但如果你们不能理解这些，那就保持沉默；不要暴露你们的无知，也不要以瞎子指控那些看见的人，或以跛子指控那些奔跑的人，而你们自己在心智上完全跛脚和残缺，过着纯粹动物般的生活——即身体的生活，那是我们本性中死亡的部分。}\par

% structure: p0094 source=h2 target=subsection reason=relative-heading-rank; base=h2

\subsection{第四十六章}

我们谨慎不反对公正的论证，即使它们出自非我信仰者；我们力求不吹毛求疵，也不试图推翻任何合理的推论。但在此，我们必须回应那些诽谤渴望尽己所能服事天主之人的品格者；天主接纳最卑微之人对祂所怀的信德，也同样接纳有学识者更精微明智的虔诚——因为两者都是通过那将纯正宗教的本质启示给人类的大司祭，向宇宙的造物主献上祈祷与感恩。那么，我们说，那些被污蔑为\Quote{精神上瘸腿和残缺}、\Quote{只为那已死的身体而活着}的人，正是努力真诚地说：\BibleQuote{因为我们虽然在血肉中行事，却不凭着血肉争战；我们争战的武器不是属血肉的，而是凭天主有力摧毁堡垒的}的人。那些对渴望成为天主仆役的人抛出如此恶毒指控的人，应当小心，免得他们因诽谤努力行善的人，而使自己的灵魂\Quote{瘸腿}，使内在的人\Quote{残缺}，从其中割断了造物主已栽植在一切理性受造物本性中的正义与心灵的节制。然而，那些借着圣言所给予的其他教训，学习并实践了\BibleQuote{被人辱骂就祝福，被人迫害就忍受，被人毁谤就劝导}的人，可谓在圣神内行走在正直的路上，净化并整顿整个灵魂。他们区分——而且这种区分对他们而言不只是言辞上的——\OriginalTerm{实体}{substance}与\OriginalTerm{生成}{becoming}，即那存在者与那生成者之间的区别；区分理性所领悟的事物与感官所领悟的事物；他们将真理联系于前者，而避免由后者产生的错误；他们受教不看那些\Quote{生成}的或现象的事物，即那些可见的、因此是暂时的，而看那些更好的事物，无论我们称之为\OriginalTerm{实体}{substance}，或\OriginalTerm{属神的}{spiritual}事物（因其为理性所领悟），或\OriginalTerm{不可见的}{invisible}（因其超出感官范围）。\Person{耶稣}的门徒看待这些现象事物，只是为了将其作为阶梯，升向理性事物的知识。因为\BibleQuote{天主的不可见之物}，即理性的对象，\BibleQuote{从创世以来就明明可见}，藉着理性，\BibleQuote{藉着所造之物就可以明白}。当他们从这世界的受造物上升到天主的不可见之物时，并不停留；而在充分运用心智于这些事物，并理解其本性之后，他们上升到\Quote{天主永恒的大能}，即祂的天主性。因为他们知道，天主出于对人的爱，已\Quote{显明}了祂的真理和\Quote{那可知的祂}，不仅赐给那些献身服事祂的人，也赐给那些远离祂所要求的纯洁敬拜与服事的人；并且，那些因天主眷顾而获得这些真理知识的人中，有些却做出与那知识不相称的事，\BibleQuote{以不义压制真理}，在天主赐予他们如此伟大真理的知识之后，他们无法在天主面前找到任何借口。\par

% structure: p0096 source=h2 target=subsection reason=relative-heading-rank; base=h2

\subsection{第四十七章}

因为\WorkTitle{圣经}对那些知道\Person{采尔苏斯}所说的那些事物、并信奉基于这些原则的哲学的人作证说：\BibleQuote{他们虽然知道天主，却不当作天主光荣祂，也不感谢祂，反而在思念中成为虚妄}；尽管天主以知识的光明照亮了他们，\BibleQuote{他们愚昧的心}却被掳去，变成\Quote{昏暗}。因此，我们可以看到，那些自以为有智慧的人在讲论天主和理性所领悟之事物上提出如此宏伟的论证之后，显示了极大的愚昧，他们\BibleQuote{将不可朽坏的天主的荣耀，变成了必朽坏的人、飞禽、走兽和爬虫形状的偶像}。既然他们生活的方式与他们从天主所接受的知识不相称，天主的眷顾任其自流，他们就被\BibleQuote{任凭他们随着心里的情欲，陷入不洁，以致彼此玷辱自己的身体}，在无耻和放纵中，因为他们\BibleQuote{将天主的真实变为虚谎，去敬拜事奉受造之物，而不敬拜事奉造物主}。\par

% structure: p0098 source=h2 target=subsection reason=relative-heading-rank; base=h2

\subsection{第四十八章}

但是那些因无知而被轻视、被看作愚人和卑贱奴隶的人，一旦将自己交托给天主，藉着接受耶稣的教导来引导，他们就远非因放纵或无耻情欲而玷污自己，反而在许多情况下，如同完美的司铎，对这些快乐毫无兴趣，在行为和思想上都保持童贞的纯洁。雅典人有一位祭司，他对自己控制情欲的能力没有信心，便决定用毒芹从根源上抑制它们；因此他在雅典人中被视为纯洁，适合主持宗教崇拜。但在基督徒中，可以找到不需要毒芹就适合纯净事奉天主的人，对于他们，圣言代替毒芹能够驱除所有邪恶意念，使他们能向神圣存有呈上祈祷。至于其他所谓的神祇，也有一群精选的童贞者，由男人看守，或者可能没有看守（这不是目前的问题），她们被认为为了所事奉的神的荣耀而生活在纯洁中。但在基督徒中，那些保持终身童贞的人，并非为了人的荣誉、酬劳或奖赏，也非出于虚荣心；而是\Quote{因他们选择将天主留在自己的知识中}，他们被天主保存在一个令祂喜悦的精神中，履行每一项职责，充满一切公义和良善。\par

% structure: p0100 source=h2 target=subsection reason=relative-heading-rank; base=h2

\subsection{第四十九章}

我现在所说的，不是为了挑剔希腊人持有的任何正确意见或健全教义，而是为了表明同样的事，甚至比这些更好、更神圣的事，已由那些神圣的人——天主的先知和耶稣的宗徒——说过了。这些真理由所有希望获得基督信仰完美知识的人充分探究，他们知道\BibleQuote{义人的口谈论智慧，他的舌头讲论公正；他天主的律法在他心里}。但即使那些因知识欠缺或缺乏意愿，或没有耶稣引领他们达到理性的宗教观，而未深入这些深奥问题的人，我们也发现他们相信至高天主和祂的独生子、圣言和天主，并且他们常在品格中表现出高度的庄重、纯洁和正直；而那些自称为智慧的人却鄙视这些美德，沉溺于鸡奸的污秽、无法无天的情欲，\BibleQuote{男与男行可耻之事}。\par

% structure: p0102 source=h2 target=subsection reason=relative-heading-rank; base=h2

\subsection{第五十章}

塞耳苏没有解释错误如何伴随\Quote{生成}或受造的产物；他也没有表达得足够清晰，使我们能够将他的观念与我们的进行比较并做出判断。但是先知们对于受造之物有一些明智的提示，他们告诉我们，甚至对新生婴儿也献了赎罪祭，因为他们并非无罪。他们说：\BibleQuote{我是在罪孽中成形的，我母亲在罪中怀了我}；还有：\BibleQuote{他们一出母胎就与主疏远}；随后还有独特的说法：\BibleQuote{他们一出世就走迷了路，说谎话}。此外，我们的智慧人对所有可感知的对象如此轻视，以至于他们有时将一切物质之事称为虚妄：\BibleQuote{因为受造之物被屈服于虚妄，并非出于自愿，而是由于那使它屈服的一位，带着希望}；有时称为虚妄的虚妄，\BibleQuote{虚妄的虚妄，传道者说，一切都是虚妄}。谁对人类灵魂地上此生的评价如此严厉，如同那说：\BibleQuote{人最好的境遇也完全是虚妄}的人？他毫不迟疑地指出灵魂今生的生命与来生的生命之间的区别。他没有说：\Quote{谁知道死不是生，生不是死？}而是大胆宣告真理，说：\BibleQuote{我们的灵魂屈身于尘土}；以及：\BibleQuote{你使我归于死亡的尘土}；同样地：\BibleQuote{谁能救我脱离这死亡的身体？}还有：\BibleQuote{谁将改变我们卑贱的身体}。也有一位先知说：\BibleQuote{你使我们降在受苦之地}；他以\Quote{受苦之地}指这个地上区域，就是亚当——即人——因罪被赶出乐园后所到之处。也要注意，那个说：\BibleQuote{现在我们对着镜子观看，模糊不清，但那时要面对面}和\BibleQuote{当我们住在身体内，就是离开了主的家}的人，对灵魂在此生和来生的不同生活有如此深刻的认识；因此\BibleQuote{我们情愿离开身体，与主同住}。\par

% structure: p0104 source=h2 target=subsection reason=relative-heading-rank; base=h2

\subsection{第五十一章}

但是，还需要引用更多反对\Person{塞尔苏斯}的语句来证明他所说的不过是我们早已说过的内容，因为前面已经充分证明了这一点。似乎接下来的话与此有些关联：\Quote{如果你们认为有一位圣神从天主那里降临，向人宣告神圣之事，那么毫无疑问，正是这位圣神启示了这些真理；古代的人也是在同样的影响下宣告了许多重要的真理。}但他不知道，那些对我们说\Quote{你那不朽的神在万物之中，因此天主逐渐责罚犯过的人}的人，以及在他们其他教导中教导我们说\Quote{领受圣神}这句话所指的属灵影响力高于\Quote{不多几日以后，你们要因圣神受洗}这句话中的人，这两者之间的差距有多大。然而，即使经过深思熟虑，也很难辨别那些在不同时间、短时间内获得真理知识和天主观念的人，与那些更完全地被天主感动、常与祂共融、始终被祂的圣神引导的人之间的区别。如果\Person{塞尔苏斯}致力于理解这一点，他就不会指责我们无知，也不会禁止我们把那些相信宗教表现在人手工制造的像中的人称为\Quote{瞎子}。因为凡是用灵魂之眼观看的人，侍奉天主的方式就是始终仰望万有的创造者，只向祂祈祷，并在那位完全看见我们、甚至看见我们思想的天主面前行事。因此，我们热切渴望既为自己看见，也作瞎子的向导，把他们带到天主的圣言那里，使祂从他们的思想中除去无知的黑暗。如果我们的行为与那位教导祂的门徒\BibleQuote{你们是世界的光}的祂相称，也与说\BibleQuote{光在黑暗中照耀}的圣言相称，那么我们将成为黑暗中人的光；我们将赐予智慧给缺乏智慧的人，我们将教导无知的人。\par

% structure: p0106 source=h2 target=subsection reason=relative-heading-rank; base=h2

\subsection{第五十二章}

如果\Person{塞尔苏斯}把我们描述为灵魂上的瘸腿和残缺的人，因为我们跑向庙宇如同跑向真正神圣的地方，却看不到人手工制造的东西不能真正神圣，但愿他不要生气。那些根基在\Person{耶稣}教导上的虔诚之人也奔跑，直到他们跑完赛程，那时他们可以满怀真实和信心说：\BibleQuote{那美好的仗我已经打过了，当跑的路我已经跑尽了，信德我已经守住了；从此以后，有公义的冠冕为我存留。}我们各人都\BibleQuote{不是无定向地}奔跑，与邪恶搏斗时\BibleQuote{不是打空气的}，而是与那些顺从\BibleQuote{空中掌权者的首领，就是现今在悖逆之子心中运行的灵}作战。\Person{塞尔苏斯}可能说我们\Quote{与那已死的身体一同生活}，但我们学会了：\BibleQuote{如果你们随从肉体生活，必要死亡；但如果你们依赖圣神，治死身体的恶行，必要生活；}以及\BibleQuote{如果我们因圣神生活，就让我们也随从圣神而行。}但愿我们能以行动让他相信，他说我们\Quote{与那已死的身体一同生活}是冤枉了我们！\par

% structure: p0108 source=h2 target=subsection reason=relative-heading-rank; base=h2

\subsection{第五十三章}

在\Person{塞尔苏斯}说了这些话（我们已尽力反驳）之后，他继续这样对我们说：\Quote{既然你们如此渴望新奇之事，那么如果你们选择一位光荣死去、其神性有某些神话支持以流传后世的人作为你们狂热崇拜的对象，岂不是好得多！为什么，如果你们不满足于\Person{赫拉克勒斯}、\Person{阿斯克勒庇俄斯}和古代其他英雄，你们还有\Person{俄耳甫斯}，一个公认受天主感动的人，他死于暴力。但也许别人在你们之前已经采纳了他。那么你们可以选取\Person{阿那克萨库斯}，当他被投入臼中遭受最野蛮的捶打时，他对痛苦表现出高贵的藐视，说：\InnerQuote{打吧，打吧，打阿那克萨库斯的壳，因为他本人你们是打不到的。}这确实是一种真正神圣精神的话语。但别人在你们之前已遵从他的自然律解释。那么，你们难道不能选取\Person{爱比克泰德}吗？当他的主人扭他的腿时，他微笑着镇定地说：\InnerQuote{你会扭断我的腿；}当腿被扭断时，他加了一句：\InnerQuote{我难道没有告诉过你你会扭断它吗？}你们的神在受苦时说过比这些更出色的话吗？如果你们甚至说\Person{西比拉}（你们中有些人承认她的权威）是天主的女儿，你们说的还更有道理。但你们竟敢在她的著作中加入许多不敬之事，并把一个以最可耻的生活结束于最悲惨死亡的人立为神。比起他，\Person{约纳}在鲸鱼腹中，或\Person{达尼尔}从野兽中被救出，或任何更离奇的事，岂不是更适合得多！}\par

% structure: p0110 source=h2 target=subsection reason=relative-heading-rank; base=h2

\subsection{第五十四章}

但他既向我们引荐赫拉克勒斯，就请他复述几句他的言论，并为他那对翁法勒的耻辱顺从进行辩护。请他证明，一个像拦路强盗那样强夺农夫的公牛然后吃掉，同时以主人的咒骂自娱的人，应当得到神圣的尊荣；为纪念此事，甚至到今天，向赫拉克勒斯邪神献祭时仍伴随着咒骂。他又向我们提出埃斯科拉庇俄斯，仿佛要迫使我们重复已说过的话；但我们克制不谈。至于奥尔甫斯，他欣赏他什么，以至于断言他经众人一致同意被视为神感之人，并度过了高贵的一生？若非塞耳苏斯有意反对我们、贬低耶稣，从而大加赞扬奥尔甫斯，我真是大错特错了；而且，当他了解那些关于诸神的不敬神话时，他难道没有将它们抛弃，认为它们甚至比荷马的诗歌更应被排除在秩序井然的城邦之外吗？的确，奥尔甫斯所说的关于他们所谓神祇的事情，比荷马所说的还要糟糕得多。阿那克萨尔库斯对塞浦路斯僭主阿里斯托克勒翁说\Quote{敲吧，敲阿那克萨尔库斯的壳}，这固然高贵，但那是希腊人所知的阿那克萨尔库斯一生中唯一令人钦佩的事迹；虽然有人像塞耳苏斯那样，据此可正当地因其勇敢而尊敬这人，但将阿那克萨尔库斯奉为神祇则不合情理。他还指引我们去看埃皮克提图，他的坚毅正当地受人钦佩，虽然他的主人打断他的腿时所说的话，不能与塞耳苏斯拒绝相信的耶稣那奇妙言行相比；那些言语伴随着如此神圣的能力，以至于直到今日，它们不仅使一些较愚昧简单的人改宗，也使许多最开明的人归信。\par

% structure: p0112 source=h2 target=subsection reason=relative-heading-rank; base=h2

\subsection{第五十五章}

当他在列举那些他要送我们到他们那里去的人时，又加上说：\Quote{你们的神在受苦时说过什么像这样的话？}我们要回答说，耶稣在鞭打下以及在所有苦难中的沉默，比希腊人在灾难中所说的一切话，更能表明祂的坚毅与顺服。也许塞耳苏斯会相信那些可信之人真诚记录的一切，他们在如实叙述耶稣所行的一切奇事时，特别指出祂在受鞭打时所保持的沉默；当人们给祂穿上紫红袍、戴上荆棘冠冕，并把一根芦苇放在祂手里当作权杖时，祂对所受的种种侮辱显出同样的独特温良；没有一句不当或愤怒的话从祂口中说出来，针对那些这样凌辱祂的人。既然祂以沉默的坚毅接受了鞭打，并以温良忍受了那些凌辱祂的人的一切侮辱，就不能像有些人所说的那样，祂是因怯懦软弱才说出这话：\Quote{父啊！若是可能，就让这杯离开我；但不要照我的意愿，而要照祢的意愿。}这祈祷看来是为求免除祂所谓的\Quote{杯}，其含义我们已在别处详加探讨和阐述。但就较明显的意义而言，试问这难道不是向天主献上的充满虔敬的祈祷么？因为没有人会天生认为临到自己的事是必然的、不可避免的；虽然他可能在必要时顺服并非不可避免的事。此外，这话：\Quote{但不要照我的意愿，而要照祢的意愿}，不是屈服于必然者的语言，而是对临到自己的一切感到满足并以敬畏之心顺从天主眷顾的安排的话。\par

% structure: p0114 source=h2 target=subsection reason=relative-heading-rank; base=h2

\subsection{第五十六章}

塞耳苏斯然后又加上——我不知道出于什么理由——与其称耶稣为天主之子，我们倒不如把这尊荣归给西比拉；他声称我们在她的书中插入许多不敬的言论，却不指出这些插句是什么。他本可以拿出一些较古老的抄本，证明这些抄本没有他归咎于我们的插句，从而证实他的断言，但他没有这样做，甚至没有证明这些段落具有不敬的性质。此外，他又多次谈到耶稣的\Quote{生活是最可耻的}，如同他以前说过的那样，不止一两次，而是许多次，尽管他没有停下来具体说明他认为最可耻的生活行为有哪些。他似乎认为他可以这样不作证明就断言，并辱骂一个他毫无所知的人。如果他着手说明他在耶稣的行为中发现哪种可耻，我们就会反驳对他提出的各项控诉。耶稣确实遭遇了最可悲的死亡；但同样的话也可用于苏格拉底、他刚提到的阿那克萨尔库斯以及许多其他人。如果耶稣的死亡是可悲的，其他人的不也是同样吗？如果他们的死并不可悲，能说耶稣的死是可悲的吗？你由此看出，塞耳苏斯的目的就是诋毁耶稣的品格；我只能认为，他是被某种与那已被耶稣击破并征服的权势类似的灵所驱使，这灵如今发现自己失去了赖以生存的烟火与血，它曾欺骗那些在地上透过偶像寻求天主、而不仰望治理万有的真天主的人。\par

% structure: p0116 source=h2 target=subsection reason=relative-heading-rank; base=h2

\subsection{第五十七章}

此后，仿佛他的目的是要扩充其书的篇幅，他劝告我们\Quote{要选择\Person{约纳}而非\Person{耶稣}作为我们的天主}；如此，他将宣讲悔改给单个城市\Place{尼尼微}的\Person{约纳}，置于向全世界宣讲悔改且效果远为巨大的\Person{耶稣}之前。他愿意我们视那位以奇特奇迹在大鱼腹中三天三夜的人为天主；他不愿意那位为人而死、天主通过先知为他作证、并在天地间行了大事的那一位，因此获得仅次于至高天主的尊荣。此外，\Person{约纳}被大鱼吞噬是因为拒绝照天主的命令宣讲；而\Person{耶稣}是遵行天主愿意他传授的训诲之后为人而死。更进一步，他又说，从狮子坑中获救的\Person{达尼尔}比\Person{耶稣}更值得我们崇拜，\Person{耶稣}却制伏了所有敌对势力的狂暴，并赐给我们\BibleQuote{权柄践踏蛇和蝎子，并胜过仇敌的一切能力}。最后，他再无别的名号可提供给我们，便加添说\Quote{还有其他一些更怪异的名号}，如此便同时轻蔑了\Person{约纳}和\Person{达尼尔}，因为\Person{色尔苏斯}里面的灵不能善待义人。\par

% structure: p0118 source=h2 target=subsection reason=relative-heading-rank; base=h2

\subsection{第五十八章}

\Quote{他们也有这样的诫命：我们不应报复那伤害我们的人，或者如他所说：\InnerQuote{谁打你的右脸，连左脸也转过来由他打。}这是一个古老的谚语，很久以前就被优美地表达过，他们只是以较粗俗的方式报道了它。因为\Person{柏拉图}记载\Person{苏格拉底}与\Person{克里同}交谈如下：\InnerQuote{\InnerQuote{我们决不可对任何人行不义吗？}\InnerQuote{当然不可。}\InnerQuote{既然我们决不可行不义，那么我们就不可回报别人对我们行的不义，就像多数人认为的那样吗？}\InnerQuote{在我看来我们不该。}\InnerQuote{但告诉我，\Person{克里同}，我们可以向任何人作恶吗？}\InnerQuote{当然不可，\Person{苏格拉底}。}\InnerQuote{那么，正如通常所说，遭受不义的人回报以不义是正义的，还是不义的？}\InnerQuote{是不义的。是的，因为加害于人等于行不义于人。}\InnerQuote{你说得对。因此我们不可用不义来回报不义，也不可向任何人作恶，无论我们从那人遭受了什么恶。}\Person{柏拉图}如此说；他又加添说：\InnerQuote{那么，请考虑你是否与我一致，并是否从这一原则出发，可以得出结论：永远不可行不义，即使是对所遭受的不义进行回报；或者，你是否与我不同，并不赞同我们出发的原则。这一直是我的意见，现在仍然是。}}这就是\Person{柏拉图}的见解，而且在他之前的神圣之人也持有这些见解。但这就足够了，以此作为其他真理被借用和歪曲的一个例子。任何愿意的人都可以轻易地查找更多。}\par

% structure: p0120 source=h2 target=subsection reason=relative-heading-rank; base=h2

\subsection{第五十九章}

当\Person{色尔苏斯}在此处或别处发现自己无法反驳我们所说的真理，却断言同样的事情已被希腊人说过，我们的回答是：如果那教义是纯正的，其效果是好的，那么无论它是经由\Person{柏拉图}或希腊的任何智者向希腊人启示，还是经由\Person{梅瑟}或任何先知交付给犹太人，还是经由\Person{耶稣基督}的记载中的教导或祂宗徒的训诲给予基督徒，这都不影响所传达的真理的价值。犹太人或基督徒的原则并不因同样的事情也被希腊人说过而受到反对，尤其是如果能证明犹太人的著作比希腊人的更古老。此外，我们不可想象一个用希腊语言的优美所装饰的真理必然优于同一个真理用犹太人和基督徒所用的更谦卑、不夸张的语言表达出来，尽管事实上犹太人用来写下流传给我们的书的语言，确实有独属于希伯来文天才的表达优美。即使有人要求我们证明同样的教义在犹太先知或基督徒著作中被表达得更好，无论这看起来多么自相矛盾，我们准备通过一个取自不同食物及其不同烹饪方式的比喻来证明这一点。假设有一种有益健康、有营养的食物，它被以某种方式烹饪和调味，只适合那些口味精致富足的人，而不适合农夫和穷苦劳动者的简单口味。再假设这同一种食物不是为迎合更精细的口味，而是为农夫、穷苦劳动者和一般平民预备的，简言之，让无数的人都能吃到。现在，根据假设，按前一种方式准备的食物只促进所谓的上等阶级的健康，而其他人都不能品尝它，而按后一种方式准备的食物却促进了广大民众的健康，那么，哪一种更有利于公共福祉呢？是为有身份的人准备食物的人，还是为大众准备食物的人？假设这两种情况下的食物都同样有益健康、有营养；因为很明显，致力于多数人健康的医生比只关注少数人健康的医生更能促进人类的福祉和公共善。\par

% structure: p0122 source=h2 target=subsection reason=relative-heading-rank; base=h2

\subsection{第六十章}

现在，理解了这一比喻后，我们必须将它应用于滋养人理性部分的神粮的特质。那么，请看：\Person{柏拉图}与希腊的智者们在其优美的言论中，是否不像那些只关注所谓上流社会而轻视大众的医生？而犹太人的先知们以及\Person{耶稣}的门徒，他们轻视单纯的文笔优雅以及圣经中所称的\Quote{人的智慧}、\Quote{属血肉的智慧}，这种智慧喜好奇诡，他们却像那些研究如何为最多的人提供最有益健康的食物的人。为此，他们使自己的语言和风格适应平民的能力，避免任何看似陌生的东西，免得因引入奇怪的表达方式而使他们对教导产生厌恶。事实上，若神粮的真正用途（继续这个比喻）是在取用者身上产生忍耐与温和的德行，那么，那个能在众人中产生忍耐与温和、或使他们在这些德行上成长的言论，难道不比那仅限于少数选民的言论（假设它真能使他们温和忍耐）更好吗？如果一个希腊人希望通过有益的教导去帮助只懂埃及语或叙利亚语的人，他首先要做的就是学习他们的语言；他宁愿在希腊人中被称为蛮夷，而说埃及人或叙利亚人的话，以便对他们有益，也不愿始终做希腊人而无法帮助他们。同样，神圣本性，其目的不仅是要教导那些被认为是精通希腊文学的人，也要教导其余的人类，因此它使自己适应它所面对的单纯大众的能力。它试图通过使用他们熟悉的语言来赢得较无知者的注意，使他们能在最初的引入之后，容易地被引导去追求隐藏在圣经中更深奥的真理。因为即使是圣经的普通读者也能看到，它包含了许多起初难以领会的事物；但那些致力于认真研读天主圣言的人却能够理解它们，并且随着他们在研究中所付出的努力和热忱，这些真理就变得明朗。\par

% structure: p0124 source=h2 target=subsection reason=relative-heading-rank; base=h2

\subsection{第六十一章}

从这些评论可以明显看出，当\Person{耶稣}说了\Person{塞耳苏斯}所称的\Quote{粗俗地}的话：\BibleQuote{有人打你的这一面颊，连另一面也转给他；有人想要告你，拿你的内衣，连外衣也让他拿去}，祂表达的方式使这条诫命比\Person{柏拉图}在\WorkTitle{克利托}中的话更有实际效果；因为后者如此不易为普通人理解，就连那些在学府中受过教育、被引入希腊著名哲学的人也很难理解他。也可以注意到，那条关于忍受伤害的诫命并没有因我主所使用的简朴语言而被败坏或贬低，但在这件事上，正如在其他事上一样，当他（\Person{塞耳苏斯}）说：\Quote{但让这一个例子足以说明这些真理和其他真理被借用和败坏的方式。任何人只要愿意，很容易通过搜索找到更多。}这纯属对我们宗教的诽谤。\par

% structure: p0126 source=h2 target=subsection reason=relative-heading-rank; base=h2

\subsection{第六十二章}

让我们现在来看看下面的话。\Quote{让我们继续，}他说，\Quote{另一个问题。他们不能容忍圣殿、祭台或雕像。在这方面，他们像斯基泰人、利比亚的游牧部落、不敬神的塞雷斯人，以及世上其他一些最野蛮和邪恶的民族。波斯人也有同样的看法，\Person{希罗多德}用以下的话证明：\InnerQuote{我知道在波斯人中，竖立雕像、祭台或圣殿被认为是不合法的；他们指责这样做的愚蠢，因为，我猜想，他们不像希腊人那样认为诸神具有人的本性。}\Person{赫拉克利特}也在一处说：\InnerQuote{向这些雕像祈祷的人，就像那些对墙壁说话的人，不知道诸神和英雄是谁。}那么，除了\Person{赫拉克利特}之外，他们还能教给我们什么更明智的教训呢？他确实足够清楚地暗示，一个人向雕像祈祷而不知道诸神和英雄是谁，是一件愚蠢的事。这是\Person{赫拉克利特}的看法；但至于他们（基督徒），他们更进一步，无一例外地鄙视所有雕像。如果他们只是说，由某个工匠加工的石、木、铜或金不能是一位神，那么他们的智慧是可笑的。因为，除非一个人极其幼稚天真，谁会把这些当作神，而不是当作献给诸神服务的祭品或代表它们的雕像呢？但如果我们不把这些视为代表神圣者，认为天主有不同的形状，正如波斯人同意他们所说的，那么他们要注意不要自相矛盾；因为他们说，天主按自己的形像造了人，并给了他一个像自己的形状。然而，他们会承认，这些雕像无论是否相似，都是为了尊荣某些存在而制作并奉献的。但他们将坚持认为，这些雕像所奉献的对象不是神，而是魔鬼，并且敬拜天主的人不应该敬拜魔鬼。}\par

% structure: p0128 source=h2 target=subsection reason=relative-heading-rank; base=h2

\subsection{第六十三章}

对此，我们的答复是：如果司库提亚人、利比亚的游牧部落、按塞尔苏斯之说没有神的赛里斯人，以及世上其他极野蛮而不敬神的民族，甚至波斯人，都不能忍受圣殿、祭台和神像，那并不因为我们与他们一样不能忍受它们，就意味着我们反对它们的理由与他们相同。我们必须探究反对圣殿和神像的原则，以便赞同那些基于正确原则反对的人，并谴责那些基于错误原则的人。因为同一件事可能出于不同的理由。例如，追随季蒂昂的芝诺的哲学家戒绝奸淫，伊壁鸠鲁的追随者也这样行，还有一些没有任何哲学原则的人也是如此；但请注意，促使这些不同类别的人如此行的理由是多么不同。第一类人考虑社会的利益，认为一个理性的人腐化一个法律已许配给他人的女子，从而破坏另一个人的家庭，这是自然所禁止的。伊壁鸠鲁派不这样推理；他们若戒绝奸淫，是因为他们把快乐视为人的至善，认为一个沉溺于奸淫的人为了这一种快乐会遇到许多快乐的障碍，如监禁、流放乃至死亡本身。实际上，他们往往在开始时冒着相当大的危险，为的是等待主人及其同伙离开屋子外出。因此，假如一个人可以犯奸淫而不被丈夫、他的仆人以及那些他会失去尊重的人发现，那么伊壁鸠鲁派就会为了快乐而犯罪。至于没有哲学体系的人，当有机会时也戒绝奸淫，他们通常是因为惧怕法律及其惩罚，而不是为了享受更多其他快乐。你看，一件被认为是同一行为的事——即戒绝奸淫——在不同的人身上并不相同，而是根据推动他们的动机而不同：一个人出于正确理由戒除，另一个人出于像伊壁鸠鲁派和我所说普通人的那样恶劣而不敬神的理由。\par

% structure: p0130 source=h2 target=subsection reason=relative-heading-rank; base=h2

\subsection{第六十四章}

因此，就像这种自我克制的行为，表面上看是同一的，实际上根据不同的人依其原则和动机而有所不同；同样，那些不能容忍在敬拜天主之事上有祭台、圣殿或神像的人也是如此。司库提亚人、游牧的利比亚人、无神的赛里斯人和波斯人在这一点上与基督徒和犹太人一致，但他们是被非常不同的原则所驱动的。这些前者中没有一个是基于害怕贬低对天主的敬拜、将其降低到敬拜人手所造的物质之物而厌恶祭台和神像的。他们反对这些也不是因为相信魔鬼选择了某些形式和地点，无论是因为它们被某种咒语束缚在那里，还是出于其他可能的原因，这些魔鬼选择了这些巢穴，在那里它们可以享用祭物的烟熏，从而追求它们邪恶的快乐。但基督徒和犹太人却遵循这条诫命：\BibleQuote{你要敬畏上主你的天主，惟獨事奉祂}\BibleRef{申 6:13}；以及另一条：\BibleQuote{除我以外，你不可有別的神。不可为你制造任何彷彿天上、或地上、或水中之物的雕像。不可叩拜这些像，也不可敬奉}\BibleRef{出 20:3-5}；以及：\BibleQuote{你要朝拜上主你的天主，惟独事奉祂}\BibleRef{玛 4:10}。正是考虑了这些和许多其他类似的诫命，他们不仅避免圣殿、祭台和神像，而且在必要时准备忍受死亡，而不愿以任何这样的不敬神之行为贬低他们对至高天主的观念。\par

% structure: p0132 source=h2 target=subsection reason=relative-heading-rank; base=h2

\subsection{第六十五章}

关于波斯人，我们已经说过，尽管他们不建造圣殿，但他们敬拜太阳和天主的其他受造物。这是我们所禁止的，因为我们被教导不要敬拜受造物而非造物主，而要认识到\BibleQuote{受造之物将脱离败坏的控制，进入天主儿女光荣的自由}\BibleRef{罗 8:21}，以及\BibleQuote{受造之物热切等待天主儿女的显扬}\BibleRef{罗 8:19}，并且\BibleQuote{受造之物被屈服于虚幻，并非出于自愿，而是出于那使它屈服的，在希望中}\BibleRef{罗 8:20}。因此，我们相信，那些\BibleQuote{处于败坏控制之下}、\BibleQuote{屈服于虚幻}、并\BibleQuote{在希望中}等待更好状态的事物，不应在我们的敬拜中占据天主——那全备者——和祂的儿子、一切受造物的首生者的位置。关于那些厌恶祭台和神像却敬拜受造物而非造物主的波斯人，我们在已有的话之外，再说这些就够了。至于塞尔苏斯所引用的\OriginalTerm{赫拉克利特}{Heraclitus}的话，他声称其要旨是\Quote{一个人向神像祈祷，却不知道诸神和英雄是谁，这是幼稚的愚蠢}，我们可以回答：很容易知道天主、天主的独生子，以及那些天主以天主的称号尊荣并分享祂神圣本性者，与那些外邦人的神——即魔鬼——截然不同；但不可能同时认识天主又向神像祈祷。\par

% structure: p0134 source=h2 target=subsection reason=relative-heading-rank; base=h2

\subsection{第六十六章}

愚昧的指责不仅适用于那些向偶像祈祷的人，也适用于那些假装随从大众榜样而行的人。属于这一类的有逍遥派哲学家以及\Person{伊壁鸠鲁}和\Person{德谟克利特}的追随者。因为在真正敬虔于天主的灵魂中，绝无虚假或伪装。我们避免尊崇偶像的另一个原因是，我们不想支持认为偶像就是神明的看法。正是基于这一点，我们谴责\Person{策尔苏斯}以及所有其他的人，他们虽然承认这些不是神明，却以智者的声誉向他们致以所谓的敬意。他们这样做便使跟随他们榜样的众人陷于罪恶，这些人崇拜这些偶像不仅是出于对习俗的尊重，而是因为他们陷入了一种信念，认为这些是真神明，并且不应听从那些认为他们所崇拜的对象不是真神明的人。\Person{策尔苏斯}确实说：\Quote{他们并不把它们当作神明，而只是作为献给神明的祭品。} 但他并没有证明它们与其说是献给人的，不如说像他所说的那样是献给神明本身的。因为很明显，它们是那些在对天主的看法上有错误的人所献的祭品。此外，我们并不认为这些偶像是天主的肖像，因为它们不能代表一位无形无像的存在。但是，由于\Person{策尔苏斯}认为我们陷入了矛盾，一方面我们说天主没有人形，另一方面我们却声称相信天主使人成为祂的肖像，并创造了人作为天主的肖像；我们的回答与已给出的相同，即我们认为与天主的相似性保存在理性的灵魂中，这灵魂被塑造向善，尽管\Person{策尔苏斯}看不出\Quote{是天主的肖像}与\Quote{是按照天主的肖像受造}之间的区别，竟断言我们说过：\Quote{天主使人成为祂自己的肖像，并给他一个与祂相似的形象。} 但这一点也已经在前面讨论过了。\par

% structure: p0136 source=h2 target=subsection reason=relative-heading-rank; base=h2

\subsection{第六十七章}

他对基督徒的下一评论是：\Quote{他们会承认这些偶像，无论是否相似，都是为尊敬某些存在而制作和奉献的；但他们也会坚持认为，这些存在所奉献的对象不是神明，而是魔鬼，并且敬拜天主的人不应敬拜魔鬼。} 如果他熟知魔鬼的本性以及它们的各种作为，无论是由精于此道者的咒语所引导，还是由于它们自己趋向依照其能力和倾向行动；如果，我再说一遍，他彻底理解了这门既广泛又难以被人类理解的学问，他就不会谴责我们说，那些敬拜至高者的人不应事奉魔鬼。至于我们自己，我们远非愿意事奉魔鬼，以至于我们通过使用从圣经学来的祈祷和其他方式，将魔鬼从人的灵魂中、从它们盘踞的地方、甚至有时从动物的身体中驱逐出去；因为甚至这些受造物也常常遭受魔鬼加给它们的伤害。\par

% structure: p0138 source=h2 target=subsection reason=relative-heading-rank; base=h2

\subsection{第六十八章}

在我们已经说了关于耶稣的一切之后，再回答\Person{策尔苏斯}的这句话将是无用的重复：\Quote{很容易证明他们敬拜的不是神，甚至不是魔鬼，而是一个死人。} 因此，我们暂且搁置这个异议，转向接下来的话：\Quote{首先，我想问为什么我们不事奉魔鬼？难道不是万有都按天主的旨意安排，祂的眷顾管理万有吗？难道宇宙中所发生的一切，无论是天主、天使、其他魔鬼，或英雄的工作，不都是由至高天主的法律所统治的吗？难道它们没有被分派各自相称的职份吗？因此，敬拜天主的人不也应该事奉那些天主赋予权力者吗？然而，他说，一个人不可能事奉许多主人。} 请再次注意，他如何一下子就解决了大量需要相当研究和对宇宙治理中最奥秘之事的深刻认识的问题。因为我们必须探究\Quote{万有按天主的旨意安排}这一说法的含义，并确定罪是否属于天主所安排的事物。因为如果天主的治理延伸到罪，不仅是在人身上，也在魔鬼和任何其他可能犯罪的灵体身上，那么说这种话的人就应明白\Quote{万有都由天主的旨意安排}这一表述有多么不妥。因为由此得出，所有罪及其后果都是由天主的旨意安排的，这不同于说它们是在天主的容许下发生的。因为如果我们按照\Quote{安排}一词的恰当含义，说\Quote{所有罪的结果都是被安排的}，那么显然万有都是按天主的旨意安排的，因此，所有作恶的人并没有违反祂的治理。同样的区别也适用于\Quote{眷顾}。当我们说\Quote{天主的眷顾管理万有时}，如果我们只把公正和正义的事归于那眷顾，我们就说出了一个大真理。但如果我们把任何事，无论多么不公正，都归于天主的眷顾，那么说天主的眷顾管理万有就不再真实，除非我们直接把那些作为祂安排的结果而流出的事归于天主的眷顾。\Person{策尔苏斯}还声称：\Quote{宇宙中所发生的一切，无论是天主、天使、其他魔鬼，或英雄的工作，都是由至高天主的法律统治的。} 但这也是不正确的；因为我们不能说，当犯法者犯法时，他们是在遵循天主的法律；而且圣经宣告，不仅是恶人是犯法者，邪恶的魔鬼和邪恶的天使也是犯法者。\par

% structure: p0140 source=h2 target=subsection reason=relative-heading-rank; base=h2

\subsection{第六十九章}

不只我们谈论邪恶的魔鬼，几乎所有承认魔鬼存在的人也都如此。因此，并非所有人都遵守至高者的法律；因为凡背离神圣法律的人，无论是由于疏忽、堕落与恶习，还是由于对正义的无知，都不遵守天主的法律，反而——使用我们在圣经中看到的一个新说法——遵守\Quote{罪恶的法律}。我说，在大多数相信魔鬼存在的人看来，其中有些是邪恶的；这些魔鬼非但不遵守天主的法律，反而违反它。但按照我们的信仰，所有魔鬼起初都不是魔鬼，而是因背离正道而成为魔鬼；因此\Quote{魔鬼}这个名称被赋予那些远离天主的存在。因此，敬拜天主的人不可事奉魔鬼。我们若观察那些用咒语召唤魔鬼以阻止某些事或达成许多其他目的的人的做法，也能了解魔鬼的真正本性。因为他们就是采用这种方法，通过咒语和巫术召唤魔鬼，诱使它们实现自己的愿望。因此，我们这些敬拜至高天主的人，不可能崇拜所有魔鬼；事奉魔鬼就是事奉所谓的众神，因为\Quote{外邦人的众神都是魔鬼}。同样的事实也体现在：最著名的所谓圣地（无论是神庙还是雕像）的奉献仪式，都伴随着古怪的魔法咒语，这些咒语由那些热心以巫术事奉魔鬼的人施行。因此，我们下定决心避免崇拜魔鬼，如同避免死亡一样；我们认为希腊人以为在祭坛、神像和神庙中向众神献上的崇拜，实际上是献给魔鬼的。\par

% structure: p0142 source=h2 target=subsection reason=relative-heading-rank; base=h2

\subsection{第七十章}

他的下一个评论是：\Quote{这些低级权力不是被天主分配了不同的部门，各自按其所配得的吗？}但这是一个需要非常深奥知识的问题。我们必须确定：统治万物的天主圣言是否指派邪恶的魔鬼担任某些职务，如同在国家里派定刽子手和其他执行残酷但必要职责的官员一样；或者，如同在旷野出没的强盗中，他们习惯从自己中间选出一位首领——因此，分散在各地如同成群结队的魔鬼，为自己选了一位首领，在其指挥下掠夺和蹂躏人的灵魂。要充分解释这一点，并说明基督徒拒绝崇拜除至高天主和一切受造的首生者（即祂的圣言和天主）之外的任何对象的行为是正当的，我们必须引用圣经中的话：\BibleQuote{凡在我以前来的，都是贼和强盗；羊却不听他们}；以及\BibleQuote{贼来，无非要偷窃、杀害、毁坏}；还有其他类似经文，如\BibleQuote{看，我已经给你们权柄去践踏毒蛇和蝎子，并胜过仇敌的一切能力；绝没有什么能伤害你们}；以及\BibleQuote{你要践踏狮子和毒蛇；少壮狮子和龙，你要踩在脚下}。但凯尔苏对这些事一无所知，否则他不会说出这样的话：\Quote{宇宙中发生的一切，无论是天主、天使、其他魔鬼还是英雄的工作，不都是受至高天主的法律所支配的吗？这些权力不是各自被分配了不同的部门，各按其配得吗？因此，事奉天主的人也应当事奉那些天主赋予如此权力的人，这不正义吗？}对此他补充道：\Quote{他们说，一个人不可能事奉多个主人。}最后这一点我们必须推迟到下一卷讨论；因为这本是我们为回应凯尔苏的论著而写的第七卷，已经足够长了。\par

\SourceRecord{Translated by Frederick Crombie.；Ante-Nicene Fathers；Vol. 4.；Edited by Alexander Roberts, James Donaldson, and A. Cleveland Coxe.；Buffalo, NY: Christian Literature Publishing Co.,；1885.；<http://www.newadvent.org/fathers/04167.htm>.}

% structure: p0001 source=h1 target=section reason=page-title

\section{\WorkTitle{驳斥克而索 第八卷}}

% structure: p0002 source=h2 target=subsection reason=relative-heading-rank; base=h2

\subsection{第一章}

我已完成了七卷，现在开始写第八卷。愿天主及其独生子——圣言，与我们同在，使我们能有效驳斥\Person{克而索}以欺骗性的标题\WorkTitle{真道}所发表的谎言，同时按我们的需要充分阐述基督信仰的真理。正如\Person{保禄}所言：\BibleQuote{我们是代基督作大使，好像天主借着我们劝你们一般；}我们也愿怀着同样的精神和言语，诚切渴望作基督的大使，去劝勉世人，因为天主圣言劝勉他们爱慕他自己，争取那些在认识\Person{耶稣基督}的道理之前，生活在对于天主的无知和黑暗之中的人，归向正义、真理和其他德行。再次，我愿说，愿天主赐予我们他纯净而真实的圣言，即\BibleQuote{那在争战中有能力的主}，以对抗罪恶。我们现在必须陈述\Person{克而索}的下一个反对意见，然后加以回答。\par

% structure: p0004 source=h2 target=subsection reason=relative-heading-rank; base=h2

\subsection{第二章}

在前文引述的一段中，\Person{克而索}问我们为什么不崇拜邪魔，我们就邪魔的问题，按照圣言的意思给出了回答。他提出这个问题，目的是引导我们崇拜邪魔，然后他描绘我们回答说，不可能服事许多主人。\Quote{这，}他接着说，\Quote{是叛乱的言论，只有那些脱离并远离所有人类社会的人才会使用。}\Quote{那些如此说话的人，}正如他所认为，\Quote{将自己的情感和激情归于天主。}\Quote{在人中间确实如此，服事一个主人的人不能好好服事另一个，因为他向一个服事会妨碍他向另一个应尽的义务；因此，已经受雇于一个主人的人，不能再接受另一个主人。}\Quote{同样，不可能同时服事不同本性的英雄或邪魔。}\Quote{但关于天主，他是不受任何痛苦或损失的，}他认为，\Quote{提防服事更多的神是荒谬的，好像我们是在与半神或其他此类神灵打交道。}他还说：\Quote{谁服事许多神，就是做取悦至高者的事，因为他尊崇属于至高者的事物。}他又补充说：\Quote{确实，尊敬天主没有给予光荣的人是错误的。}\Quote{因此，}他说，\Quote{在尊崇和敬拜所有属于天主的事物时，我们不会得罪那位万物所属者。}\par

% structure: p0006 source=h2 target=subsection reason=relative-heading-rank; base=h2

\subsection{第三章}

在进入下一点之前，最好我们先看看是否欣然接受这句话：\BibleQuote{没有人能服事两个主人，}加上\BibleQuote{因为他要么恨这一个而爱那一个，要么依附这一个而轻视那一个，}以及\BibleQuote{你们不能服事天主而又服事金钱。}对这段话的辩护将引导我们更深入、更探究性地查考\OriginalTerm{神}{gods}和\OriginalTerm{主}{lords}这两个词的意义和应用。神圣的圣经教导我们，有一位\BibleQuote{超乎诸神的大主}。以\OriginalTerm{神}{gods}这个名称，我们不是指外邦人所崇拜的对象（因为我们知道\BibleQuote{外邦人的神都是邪魔}），而是指先知们所提到的那些组成一个集会的\OriginalTerm{神}{gods}，天主\BibleQuote{审判}他们，并给每个神分派其适当的工作。因为\BibleQuote{天主在诸神的集会上站立：在诸神中施行审判。}因为\BibleQuote{天主是万神之主，}他藉其儿子\BibleQuote{从日出之地到日落之地召叫了大地。}我们也奉命\BibleQuote{要称谢万神之神。}此外，我们被告知\BibleQuote{天主不是死人的天主，而是活人的天主。}不仅如此，还有很多其他类似的经文。\par

% structure: p0008 source=h2 target=subsection reason=relative-heading-rank; base=h2

\subsection{第四章}

神圣的圣经教导我们同样思考万主之主。因为经上有一处说：\BibleQuote{要称谢万神之神，因为他的慈爱永远长存；要称谢万主之主，因为他的慈爱永远长存；}另一处说：\BibleQuote{天主是万王之王，万主之主。}因为圣经区分那些只是名称上的神和那些真正是神的神，无论是否被称作神；对于\OriginalTerm{主}{lords}这个词的使用也是如此。\Person{保禄}对此教训说：\BibleQuote{虽有称为神的，或在天上，或在地上，就如那许多的神，许多的主，}但是万神之神借着\Person{耶稣}随意召选人从东方和西方来承受他的产业，天主的基督借此显示他超越一切统治者，进入他们各自的领域，召唤他们出来归服于他，因此\Person{保禄}着眼于此，接着说道：\BibleQuote{然而我们只有一位天主，就是父，万物都出于他，也只有一位主耶稣基督，万物借着祂而有，我们也借着他而有；}他仿佛深刻体会到这道理的神奇奥秘，又加上说：\BibleQuote{然而这种知识，不是人人都有的。}当他说：\BibleQuote{我们只有一位天主，就是父，万物都出于他；只有一位主耶稣基督，万物借着他而有，}这里的\Quote{我们}指的是他自己和所有上升到万神之神、万主之主的人。现在，上升到至高的天主，就是通过他的儿子——在\Person{耶稣}身上显示出来的圣言和天主的智慧——把自己完全、专一的敬拜归给天主。因为唯有圣子引导那些努力在思想、言语和行为上保持纯洁以接近宇宙创造者天主的人。所以，我认为此世的首领，\BibleQuote{把自己扮成光明的天使，}曾经用这些话指向这一类的说法：\Quote{有一群神和邪魔跟随着他，分成十一个队伍。}他谈到自己和哲学家时说：\Quote{我们是朱庇特一派的人；其他人属于别的邪魔。}\par

% structure: p0010 source=h2 target=subsection reason=relative-heading-rank; base=h2

\subsection{第五章}

既然有如此众多的神和主，其中一些是真实存在的，另一些只是名义上的，我们不仅努力超越地上万民所崇拜的神，也超越圣经中所提及的神——那些对天主藉梅瑟和我们的救主耶稣所立的盟约一无所知、在祂藉他们赐给我们的应许中无份的人，完全不认识这些神。超越一切魔鬼崇拜的人，是不做任何取悦魔鬼之事的人；如果他像那些保禄称为\Quote{神}的人一样，或尽己所能，得以\BibleQuote{不看那看得见的，而看那看不见的}，那么他就升到比他们更有福的境界。那思想\BibleQuote{受造之物切望等候天主众子的显扬，不是出于自愿，而是由于那使它在希望中屈服的那位}，并赞美受造之物，看到它\BibleQuote{将完全脱离败坏的奴役，得享天主儿女光荣的自由}的人，不会受引诱在敬拜天主之外再敬拜任何别的，也不会侍奉两个主人。因此，持这些观点且拒绝侍奉多于一位主人的人，他们的话语毫无煽动或结党之意。对他们而言，耶稣基督是全然足够的主，祂亲自教导他们，为使他们完全受教后，能成为一个堪当天主的国度，并将他们呈献给天主圣父。的确，他们在某种意义上与那些远离天主的国度和祂的盟约的外人分离并保持距离，为能作为天国的公民生活，\BibleQuote{来到永生天主那里，来到天上的耶路撒冷，那里有千万的天使，有登记在天上的首生者的集会与教会}。\par

% structure: p0012 source=h2 target=subsection reason=relative-heading-rank; base=h2

\subsection{第六章}

但是，当我们拒绝侍奉除天主以外的任何一位——唯独侍奉祂，藉着祂的圣言和智慧——我们这样做，并非如同我们会因此对天主造成伤害或损失，像一个仆人去侍奉另一个主人时对人的伤害那样；而是我们害怕自己会受到伤害，因为我们会剥夺自己在天主内的份子，藉此我们在分享天主的福祉中生活，并充满那卓越的义子的精神，在天上圣父的众子心中，不是用言语，而是以内心的深切效果呼喊：\BibleQuote{阿爸，父啊！}拉栖第梦的使节被带到波斯王面前时，拒绝向他下拜，尽管侍从们试图强迫他们这样做，因为他们尊重那唯一对他们有权威和统治权的东西，即吕库尔戈斯的法律。然而，那些\BibleQuote{为基督作大使}的人，拥有更大得多和更神圣的使命，不应崇拜任何波斯人、希腊人、埃及人或任何民族的统治者，即使他们的官员和仆役、魔鬼和魔鬼的天使试图强迫他们这样做，并怂恿他们藐视那比地上所有法律都更强大的法律。因为那些\BibleQuote{为基督作大使}的人的主，就是基督自己，他们是祂的大使，而祂是\BibleQuote{圣言，在起初就与天主同在，且是天主}。\par

% structure: p0014 source=h2 target=subsection reason=relative-heading-rank; base=h2

\subsection{第七章}

但当克耳苏斯谈论英雄和魔鬼时，他提出了一个比他意识到的更深奥的问题。因为他在论及人间侍奉时声称，\Quote{第一个主人因他的任何一个仆人同时想侍奉另一位而受损害}，接着他又补充说，\Quote{英雄和其他类似魔鬼的情况也是如此}。现在我们必须问他，他认为那些英雄和魔鬼具有什么本性，因为他断言侍奉一位英雄的人不可再侍奉另一位英雄，侍奉一位魔鬼的人不可再侍奉另一位魔鬼，好像前一位英雄或魔鬼会像人一样，因起初侍奉他们的人后来献身于侍奉他人而受损害。也让他说明，他认为那些英雄或魔鬼会遭受什么损失。因为他将被迫要么陷入无尽的荒谬，先重复然后撤回他之前的陈述；要么放弃他轻率的猜测，承认他对英雄和魔鬼的本性一无所知。关于他的说法——当一个人的仆人转而侍奉第二个主人时，那个人遭受损害——问题来了：\Quote{对一个主人来说，一个侍奉他的仆人同时想要侍奉另一位主人，会造成什么性质的损害？}\par

% structure: p0016 source=h2 target=subsection reason=relative-heading-rank; base=h2

\subsection{第八章}

因为如果他像一个无学识、不知哲学的人那样回答，说所受的损害关乎外在事物，那么他将明显不知道苏格拉底那句名言：\Quote{阿尼图斯和梅勒图斯可以杀死我，却不能伤害我；因为较好者不可能被较差者所伤害。}但如果他说的损害是指邪恶的冲动或恶习，那么显然，一个人在不同地方侍奉两位智者，不会给智者带来这种损害。如果这个意思不符合他的意图，那么他竭力削弱\BibleQuote{没有人能侍奉两个主人}这句话的权威就是徒劳的；因为这句话只有在指我们藉着引领我们走向天主的圣子侍奉至高者时，才完全真实。我们侍奉天主，不是因为祂需要我们的侍奉，或者我们停止侍奉祂祂就会不快乐；而是因为我们自己因侍奉天主而获益，并且藉着侍奉至高天主——通过祂的独生子，圣言和智慧——我们得免忧愁和烦恼。\par

% structure: p0018 source=h2 target=subsection reason=relative-heading-rank; base=h2

\subsection{第九章}

请注意那种轻率的说法：\Quote{因为如果你敬拜宇宙中的任何其他事物}，好像要使我们相信，我们因事奉天主而敬拜属于天主的其他事物，对我们并无损害。但仿佛感到自己的错误，他纠正了那些话：\Quote{如果你敬拜宇宙中的任何其他事物}，并加以补充：\Quote{但我们不可尊敬任何人，除非是那些天主已授予此权的人。}我们要向凯尔苏斯提出这个问题，关于那些被作为神、魔鬼或英雄来尊敬的人：\Quote{先生，现在你能证明天主已将受尊敬的权利赐给了这些人，并且这并非起源于人类的愚昧和愚蠢，他们迷失了，远离了唯独应受敬拜和事奉的那一位？你不久前说，凯尔苏斯，哈德良的宠臣安提诺乌斯受到尊敬；但当然你不会说宇宙的天主曾将作为神受敬拜的权利赐给他。对于其他也是如此，我们请求证明，至高天主曾将受敬拜的权利赐给了他们。}但若有人就我们对耶稣的敬拜向我们提出同样的问题，我们将表明，受尊敬的权利是天主赐给祂的，\BibleQuote{为使众人尊敬子，如同尊敬父一样。}因为祂出生前的一切预言，都是为祂的敬拜作准备。祂所行的奇事——并非如凯尔苏斯所想的出于魔术，而是出于先知所预言的、神性的能力——是天主为基督的敬拜所作的见证。尊敬圣子，即圣言和理性，绝不违背理性，且为自己赢得极大的益处；尊敬那真理者，因尊敬真理而变得更好：我们也可以如此说尊敬智慧、公义以及圣经惯常用来称呼天主之子的所有其他称号。\par

% structure: p0020 source=h2 target=subsection reason=relative-heading-rank; base=h2

\subsection{第十章}

然而，我们献给天主之子的尊敬，以及我们献给天主圣父的尊敬，在于正直的生活，这由经文清楚教导我们：\BibleQuote{你仗恃法律，却因违犯法律而侮辱天主吗？}以及\BibleQuote{何况那践踏了天主之子，并将那使他成圣的盟约之血视为俗物，又侮辱了恩宠之神的人，你们想，他该受怎样更重的惩罚呢？}因为若违犯法律的人因违犯而侮辱天主，而践踏圣言的人践踏了天主之子，那么显然守法律的人尊敬天主，而敬拜天主的人，其生活由神圣圣言的原则和诫命所规整。若凯尔苏斯知道谁是属天主的子民，且唯有他们是有智慧的，以及谁是远离天主的陌生人，且这些都是所有不愿致力于德行的恶人，他必会在说出\Quote{一个人怎能因尊敬天主承认为祂自己所有的任何一位，而得罪天主，既然万有都属于祂？}这些话之前先深思熟虑。\par

% structure: p0022 source=h2 target=subsection reason=relative-heading-rank; base=h2

\subsection{第十一章}

他接着说：\Quote{此外，当谈到天主时，断言只有一位可被称为主的人，是在亵渎，因为他分裂了天主的国，在其中挑起叛乱，暗示神圣的国度中有不同派别，且存在一位与祂为敌的。}他如此说，若能以确凿的论据证明外邦人所敬拜的那些的确是神，而不仅仅是据说盘踞在雕像、庙宇和祭坛上的邪灵。但我们不仅渴望理解我们不断谈论和写作的那个神圣国度的本质，也渴望自己成为那唯独受天主统治的人，使天主的国成为我们的。然而，凯尔苏斯教导我们敬拜众多神明，他理当一致地不说\Quote{天主的国}，而说\Quote{众神的国}。因此，天主的国中没有派别，也没有任何神与祂为敌，尽管有一些人，像巨人和提坦族那样，存心邪恶，与凯尔苏斯一同要与天主争战，并与那些向耶稣的无数证据宣战的人，以及那作为圣言、为拯救我们人类而按照各人能接受的程度在全世界面前显现自己的那位宣战。\par

% structure: p0024 source=h2 target=subsection reason=relative-heading-rank; base=h2

\subsection{第十二章}

在接下来的内容中，有些人可能以为他说了一些对我们貌似有理的话。他说：\Quote{如果我们敬拜唯一天主，而没有别的神，也许他们会有充分的论据来反对敬拜其他神。但他们过分尊敬一位最近才在人中出现的那位，而且他们认为，如果他们也敬拜祂的仆人，这并不冒犯天主。}对此我们答复说，如果\Person{Celsus}知道那句话：\BibleQuote{我与父原是一体}\BibleRef{若 10:30}，以及天主圣子在祈祷中所用的词：\BibleQuote{愿他们合而为一，如同我们合而为一}\BibleRef{若 17:11}，他就不会以为我们敬拜任何别的神，除了那位至高天主之外。因为祂说：\BibleQuote{父在我内，我在父内}\BibleRef{若 14:10}。如果有人因这些话而担心我们会倒向那些否认圣父和圣子是两位格的人，那么让他仔细考虑那段经文：\BibleQuote{那许多信了的人，都一心一意}\BibleRef{宗 4:32}，以便他理解那句话的意思：\BibleQuote{我与父原是一体}。因此，正如我们解释过的，我们敬拜一位天主，即圣父和圣子；我们反对敬拜其他神的论据仍然有效。我们并非如祂不存在以前那样\Quote{过分尊敬一位最近才出现的那位}；因为我们相信祂自己所说的：\BibleQuote{在我还没有亚巴郎之先，我就有}\BibleRef{若 8:58}。祂又说：\BibleQuote{我是真理}\BibleRef{若 14:6}。我们当中决没有幼稚到认为真理在基督出现之前并不存在。因此，我们敬拜真理的圣父，以及作为真理的圣子；这两者就位格或自立体而言是两位，但在思想一致、和谐与意志同一上，却是合一的。他们如此完全合一，以致凡看见圣子——\BibleQuote{他是天主光荣的辉耀，是他本体的真像}\BibleRef{希 1:3}——的，就看见了祂内作为天主肖像的天主本身。\par

% structure: p0026 source=h2 target=subsection reason=relative-heading-rank; base=h2

\subsection{第十三章}

他还假设，\Quote{因为我们把敬拜天主与敬拜祂的儿子结合起来，所以在我们的观点中，不但天主，而且天主的仆人也应当被敬拜。}如果他是指那些真正是天主的仆人，在祂的独生子之后——即\Person{加俾额尔}、\Person{弥额尔}以及其他天使和总领天使——并且他说这些应该被敬拜——如果他也清楚地定义了\Quote{敬拜}一词的意义以及敬拜者的义务——我们或许就会提出我们就如此重要主题所想到的一些想法。但是，因为他把外邦人所敬拜的魔鬼也算在天主的仆人之中，他不能以一致性的借口诱使我们敬拜那些被圣言宣布为恶者的仆人、这个世界的首领、尽可能多地将人从天主那里引诱开来的那一位。因此，我们完全拒绝敬拜和侍奉他人所敬拜的那些，理由是他们不是天主的仆人。因为如果我们曾被教导将他们视为至高者的仆人，我们就不会称他们为魔鬼了。因此，我们以我们全部的力量，通过祈祷和恳求，敬拜独一天主和祂的独生子、圣言和天主的肖像；我们通过祂的独生子向宇宙的天主献上我们的祈求。我们首先将我们的祈求呈献给圣子，恳求祂，作为\BibleQuote{赎罪祭}\BibleRef{若一 2:2}和我们的大司祭，将我们的愿望、祭献和祈祷呈献给至高者。因此，我们的信德是通过祂的圣子——祂在我们内坚固信德——而指向天主的；\Person{Celsus}绝不能证明天主子是天主之国任何叛乱或不忠的原因。当我们钦佩祂的圣子——圣言、智慧、真理、正义，以及圣经所说关于祂那伟大圣父之子的一切——时，我们就是尊荣圣父。关于这点就说到这里。\par

% structure: p0028 source=h2 target=subsection reason=relative-heading-rank; base=h2

\subsection{第十四章}

Again \Person{塞尔苏斯} proceeds: \Quote{如果你告诉他们说耶稣不是天主子，而是天主是万有的父，唯独祂应当受真实的敬拜，他们不会同意停止崇拜那作为他们叛乱领袖的人。他们称他为天主子，并非出于对天主极大的敬畏，而是出于极度渴望抬高耶稣基督。} 然而，我们知道了天主子是谁，并知道祂是 \BibleQuote{祂荣耀的光辉，是祂本体的真像} 和 \BibleQuote{天主威能的气息，是从全能者荣耀流出的纯净影响}；此外，\BibleQuote{永恒之光的光辉，天主威能无瑕的明镜，以及祂良善的肖像}。所以，我们知道祂是天主子，天主是祂的父。在祂应有这样一位独生子的道理中，没有任何与天主的本性不相称或过分之处；没有人能说服我们，认为这样一位不是无生天主的子和父。如果\Person{塞尔苏斯}听说某些人主张天主子不是宇宙创造者的子，那是他与支持这种意见的人之间的事。因此，耶稣不是任何叛乱运动的领袖，而是和平的促进者。因为祂对门徒说：\BibleQuote{我把平安留给你们，我将我的平安赐给你们}；祂知道反对我们的将是世俗之人，而非天主之人，所以祂又说：\BibleQuote{不像世界所给的平安，我所赐给你们的平安}。即使我们在世上受压迫，我们仍信赖那说过的话的人：\BibleQuote{在世上你们将有苦难；但放心，我已战胜了世界}。祂就是我们称为天主子——即那位天主的子，按\Person{塞尔苏斯}的话，\Quote{我们最尊崇的}；祂是那被圣父高高举起的子。纵然在众多的信徒中可能有一些人与我们不完全一致，冒失地断言救主是至高天主；然而，我们不与他们同道，反而相信祂所说的话：\BibleQuote{派遣我的父比我大}。因此，我们不会如\Person{塞尔苏斯}所控告的那样，使我们称为父的那一位低于天主子。\par

% structure: p0030 source=h2 target=subsection reason=relative-heading-rank; base=h2

\subsection{第十五章}

\Person{塞尔苏斯}继续说：\Quote{为了真实呈现他们的信仰，我将使用他们自己的话，正如在所谓\WorkTitle{天上对话}中所说的：\InnerQuote{如果子比天主更强，而人子是祂的主，那么除了子之外，谁能做那统管万有的天主的主呢？怎么会有这么多人绕着井走，却没有一个人下井？你已经走了这么远的路，为什么还害怕？答：你错了，因为我既不缺乏勇气，也不缺乏武器。}那么，难道不清楚他们的观点正是我所描述的那样吗？他们设想另一位天上的天主是那被他们一致尊崇者的父，为要唯独尊崇这人子，他们以伟大天主的形式和名称高举他，并断言他比那统治世界的天主更强，且他统治那位天主。因此，他们的那句格言\BibleQuote{不可能侍奉两个主人}得以维持，目的是为了支持这位主的一方。} 在这里，\Person{塞尔苏斯}再次引用了某个最隐晦的异端派别的意见，并将其归给所有基督徒。我称之为\Quote{一个最隐晦的派别}；因为尽管我们常与异端辩论，但我们无法发现他是从哪一套意见中引用了这段文字，如果他确实引用自某位作者，而非自己编造或加上自己的推论。因为我们这些声明可见世界由那创造万物的天主管理者，由此宣告子并不比父更强，而是小于祂。我们将此信仰建立在耶稣自己所说的话上：\BibleQuote{派遣我的父比我大}。我们中没有一个人会疯狂到声称人子是天主的主。但当我们视救主为天主圣言、智慧、正义和真理时，我们确实说祂在这个身份上统治所有服从于祂的事物，但并不是说祂的统治延伸到处处为王的圣父和父。此外，正如圣言不违背任何人的意愿统治，仍然有邪恶的存在——不仅人类，还有天使和所有魔鬼——对于这些，我们说在某种意义上祂并不统治，因为他们不愿服从；但在另一种意义上，祂甚至统治他们，如同我们说人统治无理性的动物——不是通过说服，而是像驯服和降服狮子与驮兽者。尽管如此，祂仍然用尽一切方法说服那些不服从者顺服祂的权威。因此，就我们而言，我们否认\Person{塞尔苏斯}引用的我们的一句谚语的真实性：\Quote{除了祂，谁能做那万有之上的天主的主？}\par

% structure: p0032 source=h2 target=subsection reason=relative-heading-rank; base=h2

\subsection{第十六章}

塞尔苏斯摘录的剩余部分似乎取自某种其他形式的异端，整个内容奇怪地混杂在一起：\Quote{怎么回事，既然这么多人围着井口转，却没有一个人下到井里去？既然你已经走了这么远的路，为什么还要退缩害怕？回答：你错了，因为我既不缺乏勇气，也不缺乏武器。}我们这些属于以基督命名的教会的人断言，这些说法没有一个是真实的。因为他似乎只是为了使它们与他之前所说的话协调一致而编造了这些说法；但它们与我们无关。因为我们的原则是，不敬拜任何我们仅仅\Quote{假定}存在的神，而只敬拜这位宇宙的创造者，以及所有其他感官之眼看不见的事物。塞尔苏斯的这些评论可能适用于那些走另一条路、与我们步伐不同的人——即那些否认创造者，并以新形式为自己造另一个神的人，这神只有神的名称，他们却认为比创造者更高；与他们并列的，可能还有那些说圣子比统管万有的天主更伟大的人。关于我们不应事奉两个主人的诫命，我们已经揭示了我们认为其中包含的原则，那时我们证明，不能指控耶稣——他们的主——的追随者有任何叛乱或不忠，他们承认拒绝所有其他的主，唯独事奉那位作为圣子和圣言的天主。\par

% structure: p0034 source=h2 target=subsection reason=relative-heading-rank; base=h2

\subsection{第十七章}

塞尔苏斯接着说，\Quote{我们不敢建立祭坛、雕像和庙宇；而这，}他认为，\Quote{已被我们一致同意作为秘密和违禁团体的标志或识别特征。}他没有认识到，我们认为每个好人的精神就是一座祭坛，从那里升起真正属灵的馨香，即从纯洁良心中发出的祈祷。因此，若望在默示录中说：\BibleQuote{馨香是圣人们的祈祷}；圣咏作者说：\BibleQuote{愿我的祈祷如馨香上升到祢面前。}而适于献于天主的雕像和礼物，并非普通工匠的作品，而是由天主的圣言在我们内做工并塑造的，即我们效法\BibleQuote{一切受造物的首生者}，祂为我们树立了正义、节制、勇敢、智慧、虔诚及其他德行的榜样。因此，所有按照天主圣言，在自己灵魂内培植节制、正义、智慧、虔诚及其他德行的人，这些美德就是他们竖立的雕像，我们深信，我们应该以此来尊崇一切雕像的模型和原型：\BibleQuote{不可见的天主的肖像}，天主独生子。再者，那些\BibleQuote{脱去旧人和他的行为，穿上新人，这新人藉着知识按创造他者的肖像而更新}的人，在承担创造他们者的肖像时，就在自身内竖立了一座与至高天主自己所愿相像的雕像。正如在雕塑家中，有些人在技艺上达到惊人的完美，例如\Person{菲狄亚斯}和\Person{波利克里托斯}，在画家中，\Person{宙克西斯}和\Person{阿佩莱斯}，而另一些人则制作较差的雕像，还有一些人比二流艺术家更差——因此，总的来说，雕像和绘画的制作水平差异很大——同样，有些人以更好的方式和更完美的技巧塑造至高者的肖像；以致即使是\Person{菲狄亚斯}的奥林匹亚宙斯，也无法与按创造主天主的肖像所形成的人相比。但在整个受造物中，这一切肖像中最卓越的是我们的救主中的那个肖像，祂说：\BibleQuote{父在我内。}\par

% structure: p0036 source=h2 target=subsection reason=relative-heading-rank; base=h2

\subsection{第十八章}

每一个按自己能力效法祂的人，正是因这努力而竖立了一座按创造者肖像的雕像，因为藉着纯洁的心默观天主，他们成为祂的效法者。总的来说，我们看到所有基督徒都努力竖立祭坛和雕像，如我们所描述的那样，这些雕像并非无生命的、无感觉的，也不是为了容纳贪恋无生命事物的贪婪神灵，而是为了充满居住在我们所说德行肖像中的天主圣神，祂居住在按创造者肖像所形成的灵魂中。这样，基督的圣神就住在那些与祂自身在形式和相貌上相似的人身上。天主的圣言希望清楚地向我们表明这一点，便描绘天主向义人应许说：\BibleQuote{我要住在他们中间，在他们中行走；我要作他们的天主，他们要作我的子民。}并且救主说：\BibleQuote{如果有人听我的话并遵行，我和父要到祂那里去，并要在祂那里作我们的住所。}因此，让任何人，如果他愿意，将我所述的祭坛与塞尔苏斯所说的祭坛比较，将敬拜至高天主者灵魂中的肖像与\Person{菲狄亚斯}、\Person{波利克里托斯}等人的雕像比较，他就会清楚明白：后者是无生命之物，受时间摧残，而前者则居住在永不死去的灵魂中，只要理性灵魂愿意保存它们。\par

% structure: p0038 source=h2 target=subsection reason=relative-heading-rank; base=h2

\subsection{第十九章}

此外，如果还要将神庙与神庙比较，为要向那些接受\Person{克耳苏斯}意见的人证明，我们不反对建造与我们所说图像和祭坛相称的神庙，但我们拒绝为一切生命的赋予者建造无生命的神庙，那么，凡愿意学习的，就可以知道，我们受教导说：我们的身体是天主的殿堂，若有人以情欲或罪玷污天主的殿堂，他将被毁灭，因为他对真正的殿堂行了不敬之事。在所有这种意义上的殿堂中，最好、最卓越的是我们的救主\Person{耶稣基督}的纯洁神圣身体。祂知道恶人可能企图毁坏祂内天主的殿堂，但他们的毁灭意图不能胜过去建造那殿堂的神圣能力，于是对他们说：\Quote{\InnerQuote{你们拆毁这座圣殿，三天内我要把它重建起来。}}……祂这话是指祂身体的圣殿说的。在圣经的其他部分，当谈到复活的奥秘给那些耳朵被天主开启的人时，它将说那被毁坏的圣殿将用活而最珍贵的石头重新建造，由此让我们明白，凡被天主圣言引导、共同努力履行虔诚责任的人，都将成为天主唯一伟大圣殿中的一颗珍贵石头。因此，\Person{伯多禄}说：\BibleQuote{你们也就像活石，被建成一座属灵的殿宇，成为圣洁的司铎，藉着\Person{耶稣基督}奉献属灵的祭物。}\Person{保禄}也说：\BibleQuote{被建造在宗徒和先知的基础上，\Person{耶稣基督}自己为房角石。}在\BibleBook{}{依撒意亚}对\Place{耶路撒冷}说的这段话中，也有一个类似的隐藏暗示：\BibleQuote{看，我要用红宝石作你的基石，用蓝宝石奠定你的根基。我要用碧玉砌你的城垛，用水晶作你的城门，用各种宝石作你的边界。你所有的子女都要受上主的教导，你子女的平安将是伟大的。你将在正义中建立起来。}\par

% structure: p0040 source=h2 target=subsection reason=relative-heading-rank; base=h2

\subsection{第二十章}

因此，在义人当中，有些是红宝石，有些是蓝宝石，有些是碧玉，有些是水晶，因此在义人当中，有各种精选的珍贵宝石。至于这些不同宝石的属灵含义——它们的性质如何，以及每个珍贵宝石的名称特别适用于哪种灵魂——我们目前不能停下来仔细研究。我们只觉得有必要如此简要地表明我们对殿堂的理解，以及用珍贵石头建造的天主唯一圣殿真正意味着什么。因为正如在某些城市中，如果发生关于哪座城市有最辉煌的神庙的争论，那些认为自己的神庙最好的人会竭尽全力展示自己神庙的优越性和其他神庙的劣等——同样，当他们指责我们，认为不必通过建造无生命的神庙来崇拜神圣存在时，我们把我们的殿堂摆在他们面前，至少向那些不像他们无感觉的神祇那样盲目无感觉的人表明，我们的雕像和外邦人的雕像之间无可比较，我们的祭坛（我们可以称之为从其上发出的馨香）与外邦人的祭坛（其中充满牺牲的脂油和血）之间也无可比较；最后，由无感觉的人所赞赏的无感觉神祇的神庙，与配得上天主的殿堂、雕像和祭坛之间也无可比较。因此，我们反对建造祭坛、雕像和神庙，并非因为我们同意将此作为秘密和被禁止社团的标志，而是因为我们从\Person{耶稣基督}那里学到了侍奉天主的真正方式，我们避免任何在虔诚外表下导致那些放弃\Person{耶稣基督}为我们标示的道路的人陷于彻底不敬虔的事。因为只有祂是虔诚的道路，正如祂真实所说：\Quote{\InnerQuote{我就是道路、真理、生命。}}\par

% structure: p0042 source=h2 target=subsection reason=relative-heading-rank; base=h2

\subsection{第二十一章}

让我们看看\Person{克耳苏斯}对天主还说了些什么，以及他如何催促我们使用那些恰当称为偶像祭品，或者更好说，献给魔鬼的祭品，尽管他不了解真正的圣洁是什么，以及什么样的祭献使天主喜悦，却称它们为\Quote{圣洁的祭献}。他的话是：天主是所有人的天主；祂是良善的，一无所缺，且没有嫉妒。那么，什么能阻止那些最虔诚侍奉祂的人参加公共宴会呢？我看不出他所谓的\Quote{天主的良善、自足、无嫉妒}与祂的虔诚仆人参加公共宴会之间有什么关联。我承认，确实，从天主是良善的、一无所缺、且没有嫉妒这一事实，会推得我们可以参加公共宴会，只要证明这些公共宴会本身没有错误，并且建立在关于天主品格的正确观点之上，从而自然地源自对天主的虔诚侍奉。然而，如果所谓公共节庆绝不能证明与侍奉天主相合，反而可以被证明是人在时机出现时为了纪念某些人类事件，或为了表现水、土或地上果实的某些特性而发明的——那么，显然那些希望向神圣存在献上开明敬拜的人，将按健全理性行事，不参加公共节庆。因为\Quote{守节}，正如一位希腊智者所说，\Quote{\InnerQuote{无非就是履行自己的本分}}；那真正庆祝节日的人，就是履行本分并不断祈祷，在祈祷中向天主持续奉献无血祭献的人。因此，在我看来，\Person{保禄}有一句最崇高的话：\BibleQuote{你们谨守日子、月份、节期、年份。我为你们害怕，恐怕我为你们白费了工夫。}\par

% structure: p0044 source=h2 target=subsection reason=relative-heading-rank; base=h2

\subsection{第二十二章}

如果在这件事上有人反对我们，说我们自己习惯于遵守某些日子，例如主的日子、预备日、逾越节或五旬节，那么我必须回答：对于完美的基督徒，他总是在思想、言语和行为上事奉他的本性之主——天主圣言，他所有的日子都是主的日子，他常常遵守主的日子。同样，那不断为真实生命作准备，并戒绝这引人误入歧途的世俗享乐，不纵容肉体的情欲，但\BibleQuote{克制自己的身体，使它服从}的人，就常常遵守预备日。再者，那认为\BibleQuote{基督，我们的逾越节，已经被祭杀了}，并有责任通过吃圣言的肉来守庆节的人，就永不停止守逾越庆节；因为\Emph{\OriginalTerm{逾越}{\GreekText{πασχα}}}一词意为\Quote{逾越}，他不断在一切思想、言语和行为上努力从今生的事物逾越到天主那里，并奔向天主的城。最后，那能真正说\BibleQuote{我们与基督一同复活了}，并且\BibleQuote{他使我们与他一同复活，并在基督耶稣内使我们坐在天上}的人，就永远生活在五旬节期间；尤其是当他像\Person{耶稣}的宗徒那样，登上楼厅，专务祈祷恳求时，好能配得领受\BibleQuote{从天而来的狂风}，这风有力量摧毁众人中的罪恶及其果实，并且配得分享天主所赐的火舌之一部分。\par

% structure: p0046 source=h2 target=subsection reason=relative-heading-rank; base=h2

\subsection{第二十三章}

但大多数被视为信徒的人并不属于这个更高层次；由于不能或不愿以这种方式每天守节，他们需要一些有形的纪念物，以防止属灵的事物完全从他们的思想中消失。这种将某些日子与其余日子分别出来的做法，我认为正是\Person{保禄}在表达\Quote{部分庆节}时所指的；他用这些话指出，按照天主圣言的生活不在于\Quote{部分庆节}，而在于整个不停息的庆节。再者，请比较我们当中如上所述遵守的庆节与\Person{则利苏}及外教人的公开庆节，然后说，前者是否远为更神圣的礼仪，胜于那些肉欲放纵、导致酗酒和淫乱的庆节。我们目前要详细说明为什么天主法律要求我们吃\BibleQuote{困苦的饼}或\BibleQuote{无酵饼与苦菜}来守其庆节，或者说\BibleQuote{刻苦你们的灵魂}以及诸如此类的话，实在太长了。因为人是一个复合体，其中\BibleQuote{肉体的欲望相反圣神，圣神的欲望相反肉体}，不可能用全部本性来守节；要么他以圣神守节而折磨身体，由于肉体的情欲，身体不适合与圣神一同守节；要么他以身体守节，而圣神无法参与。但关于庆节，我们目前已经说得足够了。\par

% structure: p0048 source=h2 target=subsection reason=relative-heading-rank; base=h2

\subsection{第二十四章}

现在让我们看看，\Person{则利苏}基于什么理由敦促我们在公共庆节中使用祭偶像之物和公共祭献。他的话是：\Quote{如果这些偶像是虚无的，那么参加庆节有什么损害呢？另一方面，如果它们是魔鬼，那么它们也确定是天主的受造物，我们必须相信它们，按照法律向它们献祭，并祈求它们施恩。}关于这一陈述，我们最好拿起并清楚解释\Person{保禄}在\BibleBook{格林多}{前书}中论及祭偶像之物的整段内容，这将是有益的。宗徒从\BibleQuote{世上偶像算不了什么}这一事实推导出，使用祭偶像之物是有害的；并向那些有耳可听此类主题的人指出，参与祭偶像之物的人比杀人犯更糟糕，因为他毁灭了自己的弟兄，基督曾为他们而死。此外，他坚持认为祭献是献给魔鬼的；并由此指出，凡参与魔鬼筵席的人就与魔鬼联合；他得出结论，人不能既参与主的筵席，又参与魔鬼的筵席。但是由于需要一整篇论文才能充分阐述\BibleBook{格林多}{前书}中关于此主题的一切内容，我们将满足于这一简短的论证陈述；因为任何仔细思考过所说内容的人都会明白，即使偶像是虚无的，参与偶像庆节仍然是可怕的事。即使假定有这样的存在即魔鬼，祭献是献给它们的，也已经清楚地表明，当我们知道主的筵席与魔鬼的筵席之间的区别时，我们被禁止参加这些庆节。知道了这一点，我们尽可能努力常常参与主的筵席，并尽量避免在任何时候参与魔鬼的筵席。\par

% structure: p0050 source=h2 target=subsection reason=relative-heading-rank; base=h2

\subsection{第二十五章}

策尔苏斯说：\Quote{魔鬼属于天主，因此应当相信他们，应当按法律向他们献祭，应当祈求他们施恩。}凡愿意学习的人，必须知道天主的圣言从未说邪恶之物属于天主，因为天主判定它们不配这样一位主。因此，并非所有被称为\Quote{天主的人}的人都是如此，只有那些堪当天主的人——如\Person{梅瑟}、\Person{厄里亚}，以及任何其他这样称呼的人，或类似圣经中这样称呼的人——才这样称呼。同样，并非所有天使都称为天主的天使，只有那些有福的天使才这样称呼；那些堕落于罪恶中的被称为\Quote{魔鬼的天使}，正如坏人被称为\Quote{罪人}、\Quote{丧亡之子}或\Quote{不义之子}。既然在人间有善有恶，前者被称为属于天主，后者属于魔鬼，那么天使中有些是天主的天使，有些是魔鬼的天使。但在\OriginalTerm{魔鬼}{demons}中却没有这种区分，因为所有的都被认为是邪恶的。因此我们毫不犹豫地说，策尔苏斯说\Quote{如果他们是魔鬼，显然他们也必须属于天主}是虚假的。他必须要么证明天使和人类中善恶的区分毫无根据，要么证明在魔鬼中也能找到类似的区分。如果这是不可能的，那么显然魔鬼不属于天主；因为他们的王子不是天主，而是如圣经所说，是\Quote{\OriginalTerm{贝耳则步}{Beelzebub}}。\par

% structure: p0052 source=h2 target=subsection reason=relative-heading-rank; base=h2

\subsection{第二十六章}

我们不应相信魔鬼，尽管策尔苏斯敦促我们这样做；但如果我们必须顺从天主，我们宁可死或忍受任何事，也不顺从魔鬼。同样，我们不应向魔鬼求恩；因为无法向邪恶且谋求伤害人的存在求恩。此外，策尔苏斯要我们向魔鬼求恩所依据的法律是什么？如果是指国家制定的法律，他必须证明这些法律与神圣法律一致。但如果做不到这一点——因为许多国家的法律彼此完全不一致——那么这些法律必然要么根本不是严格意义上的法律，要么是恶人的法令；我们不可顺从这些法令，因为\Quote{我们必须顺从天主，而非世人。}所以，策尔苏斯劝我们向魔鬼祈祷的忠告应当摒弃，一刻也不能听从；因为我们的责任是只向至高天主和独生子、一切受造物的首生者祈祷，并请求祂作为我们的大司祭将我们从祂上升起的祈祷呈献给祂的天主和我们的天主，祂的父和那些按祂圣言生活之人的父。正如我们不愿得到那些希望我们追随他们邪恶生活之人的恩宠，且他们只在条件——我们不选择任何违背他们意愿之事——下才赐予我们恩宠，因为他们的恩宠会使我们成为天主的仇敌，而天主不会喜悦那些将这些人当作朋友的人——同样，那些了解魔鬼的本性、意图和邪恶的人，决不愿获得他们的恩宠。\par

% structure: p0054 source=h2 target=subsection reason=relative-heading-rank; base=h2

\subsection{第二十七章}

基督徒无所畏惧，即使魔鬼对他们不怀好意；因为他们受到至高天主的保护，天主喜悦他们的虔诚，并派遣祂的神圣天使守护那些堪当如此护卫的人，使他们不受魔鬼的任何侵害。谁因虔诚拥有至高者的恩宠，接受了耶稣——\Quote{伟大的谋略天使}——的引领，因通过基督耶稣获得天主的恩宠而满足，就能自信地说，他绝不会遭受全体魔鬼大军的任何侵害。\BibleQuote{主是我的光明，我的救援，我还畏惧谁？主是我生命的保障，我尚恐惧何人？虽然大军向我进攻，我的心仍然安泰。}以上是对策尔苏斯这些声明的答复：\Quote{如果他们是魔鬼，他们也显然属于天主，应当相信他们，应当按法律向他们献祭，应当向他们祈祷，求他们施恩。}\par

% structure: p0056 source=h2 target=subsection reason=relative-heading-rank; base=h2

\subsection{第二十八章}

现在我们继续讨论策尔苏斯的下一个声明，并仔细审视它：\Quote{如果由于他们祖先的传统，他们戒绝这类祭品，那么他们也必须依照毕达哥拉斯的见解戒绝一切动物食物，毕达哥拉斯以此表明他对灵魂及其身体器官的尊重。但如果如他们所说，他们戒绝是为了不与魔鬼同吃，我佩服他们的智慧，因为他们终于发现，每当我们吃时，他们总是与魔鬼同吃，尽管他们只在看见被宰杀的祭品时才拒绝这样做；因为当他们吃面包、喝酒或尝水果时，难道他们不是从某些魔鬼那里领受这些东西，以及他们喝的水和呼吸的空气吗？这些不同的自然领域被分配给了这些魔鬼？}在此我注意到，我无法理解他所说的某些人按照他们祖先的传统戒绝某些祭品，因而必须戒绝所有动物的肉。我们确实不否认神圣之言似乎命令某种类似的事，当它为了提升我们到更高更纯净的生活时，说：\Quote{不吃肉、不喝酒，不做任何使你的弟兄绊倒、或冒犯、或软弱的事，是好的；}又说：\Quote{不可因你的食物使基督为他而死的人丧亡；}又说：\Quote{如果食物使我的弟兄跌倒，我就永远不吃肉，免得使我的弟兄跌倒。}\par

% structure: p0058 source=h2 target=subsection reason=relative-heading-rank; base=h2

\subsection{第二十九章}

但应注意，那些自认为正确理解梅瑟法律的犹太人，谨慎地只吃被视为洁净的食物，并戒绝不洁净的食物。他们也不吃动物的血和被野兽撕裂的肉，以及其他一些目前我们不便详述的东西。然而，耶稣愿藉祂的教导引导众人以纯洁的方式敬拜和事奉天主，不愿因制定烦琐的饮食规则而妨碍许多可能因基督徒信仰而获益的人，因此祂规定：\BibleQuote{不是入口的污秽人，而是出口的才污秽人；因为凡入口的，是进到肚腹里，然后排泄到厕所里去。但出口的，是由心里发出来的恶念、凶杀、奸淫、邪淫、偷盗、妄证、亵渎。}保禄也说：\BibleQuote{食物不能使我们取悦于天主，因为我们不吃也无损，吃也无益。}因此，为免此事因未加说明而含混不清，耶稣的宗徒和长老们聚集在安提约基雅，并与圣神——正如他们自己所宣称的——一同决定写信给外邦信徒，禁止他们吃那些他们声称必须戒绝的东西，即：\BibleQuote{祭过偶像的食物、勒死的牲畜和血。}\par

% structure: p0060 source=h2 target=subsection reason=relative-heading-rank; base=h2

\subsection{第三十章}

因为祭过偶像的食物是祭献给邪魔的，天主的人不可共赴邪魔的筵席。至于勒死的牲畜，圣经禁止我们吃，因为血还在其中；而血，尤其是血的气味，据说是邪魔的食物。因此，我们若吃勒死的牲畜，或许会有这样的邪灵与我们一同进食。而禁止吃勒死牲畜的理由同样适用于吃血。提到这一点，不妨回忆一下塞克斯都著作中一句为大多数基督徒所熟知的美妙格言：\Quote{关于吃动物，}他说，\Quote{这是无所谓的事；但戒除它们更合乎理性。}因此，我们戒绝吃献给那些所谓的神祇、英雄或邪魔的祭品，并非仅仅出于我们父辈的一些传统，而是出于其他理由，其中一些我在此已提及。然而，不可认为我们应当像戒除一切丑行与邪恶那样戒除肉食：事实上，我们不仅戒除肉食，而且戒除所有其他种类的食物，如果我们不能享用它们而不招致邪恶及其后果的话。因为我们应当避免为贪吃或仅为满足食欲而吃，而忽视身体的健康与营养。我们不相信灵魂从一个身体转移到另一个身体，并且会下降到如此低的地步以致进入牲畜的身体。因此，我们若有时戒除吃动物的肉，显然并非出于与毕达哥拉斯相同的理由；因为我们尊崇的只是理性的灵魂，我们将它的身体器官以应有的尊重交付坟墓。因为理性灵魂的居所不应被随意丢弃而无尊荣，如同牲畜的尸体一般；当我们相信对身体的尊重也归于那位从天主领受灵魂、并以高尚方式使用它所寄居的身体器官的人时，就更是如此。关于\BibleQuote{死人怎样复活？带着什么身体来呢？}这个问题，我们已按我们的目的简短回答过。\par

% structure: p0062 source=h2 target=subsection reason=relative-heading-rank; base=h2

\subsection{第三十一章}

之后，策尔苏斯陈述犹太人和基督徒同样引以为据，以辩护他们不参与偶像祭献的理由，即：那些献身于至高天主的人与魔鬼一同进食是不对的。我们已经看到他对这一观点的反驳。依我们的意见，只有当一个人吃所谓祭献的肉，并喝那为尊崇魔鬼而倾奠的酒时，才能说他与魔鬼一同饮食。但策尔苏斯认为，我们无论如何吃饼饮酒，或品尝果实，甚至喝一口水，都不能不与魔鬼一同饮食。他还补充说，我们所呼吸的空气也是从魔鬼那里得来的；没有一只动物能不通过掌管空气的魔鬼而呼吸空气。如果有人要为策尔苏斯的这一说法辩护，就让他指出：并非天主的神圣天使，而是魔鬼（其整个种类都是邪恶的）被指派来传达上述一切恩惠。我们确实也主张，不仅地上的果实，而且每一条溪流、每一口空气，大地所出产的那些所谓自然生长的东西，以及泉水从泉源涌出、以流动的溪流滋润大地，空气保持纯净并维持呼吸者的生命，这些都只是由于某些存在者的运作和掌管（我们可以称它们为无形的农夫和守护者）才得以实现；但我们否认这些无形的存在者是魔鬼。如果我们可以大胆地说，我们会说：如果魔鬼在这些事中有什么份的话，那么属于它们的是饥荒、葡萄和果树的枯萎、人畜的瘟疫——所有这些正是魔鬼的适当工作，它们作为公开的行刑者，在某些时期获得权力来执行天主的审判，为的是使那些陷入邪恶深渊的人得到恢复，或为考验和锻炼智慧者的灵魂。因为那些在一切苦难中保持纯全无瑕的敬虔的人，向所有可见或不可见、观看他们的旁观者显明他们的真实品格；而那些心有别念、却隐藏其邪恶的人，当他们的真实品格因不幸而暴露时，既显明给他们自己，也显明给我们也可以称为旁观者的人。\par

% structure: p0064 source=h2 target=subsection reason=relative-heading-rank; base=h2

\subsection{第三十二章}

圣咏作者见证，天主的公义使用某些邪恶的天使向人降灾：\BibleQuote{祂将祂的烈怒、忿恨、恼怒和患难，藉着邪恶的天使降在他们身上。}是否魔鬼在被容许做它们随时准备做、但因受到约束而并非总能做的事情时，会越出这个范围？这是一个只能由那样的人来解决的问题：他能按人性所能允许的程度设想，如此众多的人类灵魂在走向通往必然死亡的道路时与身体分离，这如何与天主的公义相符。\BibleQuote{因为天主的判断是如此伟大，}以致一个仍披着可朽身躯的灵魂无法理解它们；\BibleQuote{并且它们不能言说；因此未受培育的灵魂}完全无法领悟它们。因此，鲁莽的灵因其在这些事情上的无知，并因鲁莽地对抗天主本身，而繁衍出不敬神的对眷顾的反对意见。所以，人并非从魔鬼那里获得任何满足生活必需的事物，尤其我们这些被教导要正确使用这些东西的人更不是如此。那些享用谷物、葡萄酒、树上的果实、水和空气的人，并非与魔鬼一同进食，反而更是与那些为此目的而被指派的神圣天使一同欢宴；这些天使仿佛被邀请到虔诚人的餐桌，这人听从圣言的诫命，说：\BibleQuote{所以你们或吃或喝，或无论做什么，一切都要为光荣天主而做。}又在另一处经上记着：\BibleQuote{你们无论做什么，都要因天主的名而行。}因此，当我们为光荣天主而吃、喝、呼吸，并在一切事上按正理而行时，我们并非与魔鬼一同欢宴，而是与神圣的天使一同欢宴：\BibleQuote{因为天主所造的样样都好，如以感恩的心领受，没有一样是可摈弃的；因为它藉着天主的话和祈祷得以圣化。}但如果这些事物，如策尔苏斯所认为的那样，是交托给魔鬼掌管，它们就不能是好的，也不能得以圣化。\par

% structure: p0066 source=h2 target=subsection reason=relative-heading-rank; base=h2

\subsection{第三十三章}

由此明显可见，我们已经回应了\Person{色尔苏斯}接下来的陈述，他这样说：\Quote{我们要么不活，甚至根本不该来到此生，要么就必须在这样一个条件下活着：即向那些被派驻管理世间万物的鬼神献上感恩、初果和祈祷；并且我们一生都要这样行，好使他们成为良善仁慈的。}我们当然必须活着，而且我们必须尽可能按照天主的圣言活着。当我们\BibleQuote{或吃或喝，都为光荣天主而行}时，我们就能这样活着；我们不应拒绝享用那些为我们受造的事物，却要怀着对造物主的感恩接受它们。正是在这些条件下，而非\Person{色尔苏斯}所想象的那样，我们被天主带入生命；我们不在鬼神的管辖之下，而是受至高天主的治理，借着那引领我们归向天主的——\Person{耶稣基督}。任何鬼神掌管世间事物，都不是按照天主的法律，而是因他们自己的不法，或许他们为自己寻得了缺乏天主知识和神圣生命的区域，或是充满天主仇敌的地方。也许，作为适合统治和惩罚他们者，他们被治理万物的圣言派去统治那些臣服于邪恶而非天主的人。为此，让\Person{色尔苏斯}这个不认识天主的人向鬼神献上感恩祭吧。但我们向万有的造物主献上感恩，并在感谢和为所受恩惠祈祷的同时，也吃那献给我们的饼；这饼藉着祈祷成为圣体，圣化那些真诚领受它的人。\par

% structure: p0068 source=h2 target=subsection reason=relative-heading-rank; base=h2

\subsection{第三十四章}

\Person{色尔苏斯}还希望我们向鬼神献上初果。但我们愿向那一位献上初果，他曾说：\BibleQuote{地要生出青草，结种子的菜蔬，并结果子的树木，各从其类，果子都包着核。}我们向他献上初果，也向他献上祈祷，\BibleQuote{因为我们有一位伟大的大司祭，已经进入诸天，就是天主子耶稣}，并且我们一生\BibleQuote{持定这个承认}；因为我们发现天主和他的独生子，在耶稣内彰显给我们，对我们真是慈爱和良善。如果我们还希望有众多永远友善的存有，我们被告知：\BibleQuote{有千万的服役者侍立在他面前，有亿万的侍从侍奉他。}这些存有把所有效法他们对天主的虔诚、并诚心祈祷呼求天主的人，都视为自己的亲人和朋友，与他们一同为他们的得救而工作，向他们显现，看顾他们是他们的职责和义务，并且如同有一项协议一般，他们以各种仁慈和拯救眷顾那些向天主祈祷的人，而他们自己也向天主祈祷：\BibleQuote{他们不都是奉职的神，被派遣给那些要承受救恩的人服务吗？}让博学的希腊人说，人的灵魂出生时就被置于鬼神的监管之下；但耶稣教导我们不要轻视教会中的小孩子，他说：\BibleQuote{他们的天使在天上，常见我在天之父的面。}先知说：\BibleQuote{上主的使者在敬畏他的人四围安营，拯救他们。}因此，我们不否认地上有许多鬼神，但我们坚持认为，他们存在于恶人中间并行使权柄，作为对他们邪恶的惩罚。但他们无权辖制那些\BibleQuote{穿戴了天主的全副武装}、获得力量\BibleQuote{抵挡魔鬼的诡计}、并且常常与他们争战的人，因为他们知道\BibleQuote{我们不是对抗血和肉，而是对抗率领者，对抗掌权者，对抗这黑暗世界的统治者，对抗天上的邪恶之神。}\par

% structure: p0070 source=h2 target=subsection reason=relative-heading-rank; base=h2

\subsection{第三十五章}

现在让我们来看\Person{色尔苏斯}的另一句话，他这样说：\Quote{波斯或罗马君王的总督，或统治者、将军、长官，甚至那些在国家中担任较低级职务或差役的人，都能严重伤害那些藐视他们的人；那么，地上和空中的总督和臣仆受到侮辱，难道就能不受惩罚吗？}现在请注意，他如何引入至高者的仆人——统治者、将军、长官，以及那些担任较低级职务和差役的人——像人一样，对那些侮辱他们的人施加伤害。因为他没有考虑到，一个智者不愿加害任何人，而是会竭尽全力去改变和纠正他们；除非\Person{色尔苏斯}所认为的由至高者任命的仆人和统治者，竟不如斯巴达的立法者\Person{来库古}或基提翁的\Person{芝诺}。因为当\Person{来库古}被人打瞎一只眼睛后，他将肇事者置于自己的权力之下；但他没有报复，而是不停地用一切劝说技巧，直到使他成为一个哲学家。而\Person{芝诺}，当有人说：\Quote{我宁可死，也要报复你}，他回答说：\Quote{宁可让我死，如果我不使你成为朋友。}我还没有提到那些受耶稣教诲塑造品格的人，他们听过这样的话：\BibleQuote{爱你们的仇敌，为迫害你们的人祈祷，好使你们成为你们天父的儿女，因为祂使太阳升起，光照恶人也光照善人，降雨给义人也给不义的人。}在先知书中，义人说：\BibleQuote{上主我的天主，如果我做了这事，如果我手中有不义，如果我以恶报那以恶待我的人，就让我无力地落在仇敌手中，让仇敌追杀我的灵魂，夺取它，将我的生命践踏在地上。}\par

% structure: p0072 source=h2 target=subsection reason=relative-heading-rank; base=h2

\subsection{第三十六章}

但是，天使——他们是天主真正的统治者和将领及仆役——并不像克耳苏所设想的那样，\Quote{伤害那些冒犯他们的人}；如果克耳苏所指的某个魔鬼确实施加恶事，那就表明他们是邪恶的，并且从未从天主那里接受过这类职分。他们甚至伤害那些在他们权下、并承认他们为主的人；因此，似乎那些违反任何地方关于食物规则的人，若处于该地的魔鬼权下，就要受苦；而那些不在他们权下、没有屈服于他们权力的人，则免受一切伤害，并对这些邪灵无所畏惧；尽管若因无知而落入其他魔鬼的权下，他们可能会受其惩罚。但是基督徒——我指的是真正的基督徒——他只屈服于天主和祂的圣言，不会受魔鬼的任何伤害，因为祂比魔鬼更强有力。基督徒不会受任何伤害，因为\BibleQuote{上主的天使要在敬畏祂的人四周扎营，并拯救他们}，而他的\Quote{天使}，\BibleQuote{常常在天上看见我天父的面}，通过唯一的大祭司向万有的天主献上祈祷，并将其祈祷与托付给他照管的人的祈祷相结合。因此，不要让克耳苏试图以魔鬼带来的祸害来吓唬我们，因为我们鄙视它们。魔鬼被鄙视时，对那些受祂——那位唯一能帮助所有配得祂帮助的人——保护的人无能为力；祂所做的，就是把自己的天使安置在虔敬的仆人身边，以至没有任何敌对的天使，甚至那称为\Quote{今世的首领}的，能对那将自己献给天主的人造成任何影响。\par

% structure: p0074 source=h2 target=subsection reason=relative-heading-rank; base=h2

\subsection{第三十七章}

接下来，克耳苏忘记了他是在对基督徒说话，他们只通过耶稣向独一天主祈祷；他将其他观念与他们的相混，荒谬地把这一切都归于基督徒。他说：\Quote{如果被呼求者以蛮族之名被召唤，他们就会有能力；但如果以希腊文或拉丁文被召唤，他们就再没有能力了。}那么，请他明确说出我们以蛮族之名呼求的是谁。任何人都会相信，这是克耳苏对我们的诬告，只要想到基督徒在祈祷时甚至不使用圣经指称天主的精确名称；而是希腊人用希腊名称，罗马人用拉丁名称，每个人都用母语尽己所能地祈祷和赞美天主。因为掌管地上万种语言的主垂听以不同语言向祂祈祷的人，如果我可以这样说，祂听到的只是一种声音，却用不同的方言表达出来。因为至高者不像那些只选择一种语言（蛮族语或希腊语）的人，对其他语言一无所知，也不关心说其他语言的人。\par

% structure: p0076 source=h2 target=subsection reason=relative-heading-rank; base=h2

\subsection{第三十八章}

接着，他让基督徒说出他从未从任何基督徒口中听到的话；如果他确实听到过，那一定是来自最无知和最无法无天的人。他们被说成：\Quote{看，}他到一座朱庇特或阿波罗或其他神祇的雕像前，说：\Quote{我辱骂它，击打它，它却不对我复仇。}他不知道，神圣律法的禁令中包括：\Quote{不可辱骂神祇}，这是为了防止养成辱骂任何人的习惯；因为我们已受教导：\BibleQuote{要祝福，不要诅咒}，且经上说：\BibleQuote{辱骂者不得继承天主的国}。我们中有谁如此愚蠢，以致说出克耳苏所描述的话，并且看不到这种轻蔑的言语无助于消除关于神祇的普遍看法？因为可以观察到，有些人彻底否认天主或掌管眷顾的存在，并以其不敬神和破坏性的教导在所谓哲学家中建立了派别，但他们自己以及接受其观点的人，都没有遭受人类普遍认为的祸害：他们身体健壮，财产丰富。然而，如果我们问他们遭受了什么损失，我们会发现他们确实受到了最确定的伤害。因为一个人若无法在世界秩序中看见创造它的那一位，还有什么更大的伤害能临到他呢？还有什么比心智的盲目——使他看不见每一个灵魂的创造者和父亲——更痛苦的磨难呢？\par

% structure: p0078 source=h2 target=subsection reason=relative-heading-rank; base=h2

\subsection{第三十九章}

在把这样的话放在我们口中，并恶意指控基督徒怀有他们从未有过的情感之后，他继而对这个所谓基督徒情感的表达作出答复，这与其说是答复，不如说是嘲弄，他说：\Quote{好先生，难道你没有看见，连你自己的魔鬼也不仅被辱骂，而且被驱逐出每一寸土地和海洋，而你自己，作为献给它的像，被捆绑并带去受罚，被钉在木桩上，而你的魔鬼——或如你所说的\InnerQuote{圣子}——却不对那作恶者复仇吗？}如果我们确实使用了他归于我们的那种语言，这个答复是可以接受的；尽管即使如此，他也没有权利称天主子为魔鬼。因为我们坚持所有魔鬼都是邪恶的，那使如此多人归向天主的一位，在我们看来不是魔鬼，而是天主圣言和天主子。我不知道克耳苏如何如此忘乎所以，竟称耶稣基督为魔鬼，尽管他从未提及任何邪恶魔鬼的存在。最后，至于对不敬神者所威胁的惩罚，这些惩罚将在他们拒绝一切补救、并且可以说患上了不可救药的罪恶之疾之后临到他们身上。\par

% structure: p0080 source=h2 target=subsection reason=relative-heading-rank; base=h2

\subsection{第四十章}

这就是我们关于惩罚的教义；而这一教义的灌输使许多人远离他们的罪。但让我们看看另一方面，赛尔苏斯提到的朱庇特或阿波罗的司铎对这个问题给出了什么答复。是这样的：\Quote{众神的磨磨得慢。}另一个人描述惩罚达于\Quote{子孙的子孙，以及那些在他们之后的人。}圣经的话是多么更好：\BibleQuote{父亲不可因儿女被处死，儿女也不可因父亲被处死；每个人都要因自己的罪被处死。}又说：\BibleQuote{凡吃酸葡萄的，自己的牙齿要发软。}并且：\BibleQuote{儿子不担当父亲的罪，父亲也不担当儿子的罪；义人的义必归自己，恶人的恶也必归自己。}若有人说那答复\Quote{直到子孙的子孙，以及那些在他们之后的人}与那段经文\BibleQuote{恨我的，我必追讨他的罪，自父及子，直到三四代}相符，让他从厄则克耳那里知道，这种说法不可按字面理解；因为他责备那些说\BibleQuote{我们的父亲吃了酸葡萄，儿女的牙齿就发软}的人，然后他说：\BibleQuote{我指着我的永生起誓，主说：每个人都要因自己的罪而死。}至于关于罪追讨到三四代的比喻性说法的正确含义，我们目前不能停留解释。\par

% structure: p0082 source=h2 target=subsection reason=relative-heading-rank; base=h2

\subsection{第四十一章}

然后他接着像老妇人一样地责骂我们。他说：\Quote{你们嘲笑和谩骂我们神的雕像；但如果你们当面辱骂巴克斯或赫丘利，也许就不会逃之夭夭了。然而，那些当你们的天主在人中时将他钉死的人，当时或他们一生中都未受任何惩罚。自从那时以来，发生了什么新的事情，使我们相信他不是骗子，而是天主之子？而且，那位派遣他的儿子带着某些指示给人类的天主，竟允许他遭受如此残忍的对待，让他的指示随他一起消亡，在这漫长的岁月中从未表露丝毫关心。哪有父亲如此不人道的？也许你会说，他之所以受这么多苦，是因为他愿意承受降临到他的事。但你们恶意诋毁的那些人也可以采用同样的说法，说他们愿意被诋毁，因此他们耐心忍受——因为最好公平对待双方——尽管这些神严厉惩罚嘲笑者，以至于他必须要么逃跑躲藏，要么被抓住而灭亡。}现在，我要回答这些说法：我们不诋毁任何人，因为我们相信\BibleQuote{诋毁者不能承受天主的国}。我们也读到\BibleQuote{祝福那咒骂你们的人；要祝福，不可诅咒}；还有\BibleQuote{被辱骂，我们就祝福}。即使我们向他人倾倒的辱骂似乎因我们所受的委屈而有某些借口，这样的辱骂仍不被天主的圣言所允许。何况我们若考虑辱骂他人是多么愚蠢，就更当避免！将辱骂用于石、金、银所制成的、被那些不认识天主的人当作神像的东西，同样是愚蠢的。因此，我们嘲笑的是那些崇拜它们的人，而不是无生命的像。此外，如果某些魔鬼居住在特定的像中，其中一个被称为巴克斯，另一个被称为赫丘利，我们也不诋毁它们：因为一方面，这毫无用处；另一方面，这不符合一个温和、和平、心灵温柔的人，他已学会不可诋毁任何人类或魔鬼，无论他多么邪恶。\par

% structure: p0084 source=h2 target=subsection reason=relative-heading-rank; base=h2

\subsection{第四十二章}

赛尔苏斯无意中陷入了一种矛盾，这很奇怪。他刚才还赞扬的那些魔鬼或神，现在却显明是最卑劣的受造物，他们的惩罚更多是为了报复而非为了改进那些辱骂他们的人。他说：\Quote{如果你们当面辱骂巴克斯或赫丘利，就不会安全逃脱。}任何一个人都会解释，一个人如何能在没有亲身在场的情况下听到，以及诸如此类的问题：\Quote{为什么他有时在场，有时不在场？}以及\Quote{是什么事务使魔鬼从一个地方到另一个地方？}又，当他说：\Quote{那些钉死你们的天主自己的人，并没有因此受害，}他假设我们称之为天主的，是耶稣被钉在十字架上的身体并被杀害，而不是他的神性；并且耶稣作为天主被钉死和杀害。我们已经详细论述过耶稣作为人所受的苦难，因此这里故意不再赘述，以免重复已说过的话。但他接着说：\Quote{那些给耶稣施加死亡的人，在如此长的时间里并没有遭受什么，}我们必须告诉他以及所有愿意学习真理的人：那个犹太民族曾高喊\Quote{钉死他，钉死他}要求钉死耶稣，宁愿释放因叛乱和谋杀被囚的强盗，而将因嫉妒被交付的耶稣钉死——这城市不久之后被攻击，经过长期围困后被彻底推翻和荒废；因为天主判定那地方的居民不配共同过公民的生活。然而，尽管说来难以置信，天主在将他们交给敌人时还是宽恕了他们；因为他看到他们无可救药地不愿悔改，且日益深陷邪恶。这一切降临到他们，是因为耶稣的血在他们的煽动下流在他们土地上；那土地不再能承受对耶稣犯下如此可怕罪行的人。\par

% structure: p0086 source=h2 target=subsection reason=relative-heading-rank; base=h2

\subsection{第四十三章}

自\Person{耶稣}受苦以来，发生了一些新事——我是指那临到这城、这整个民族的事，以及基督徒团体的突然而普遍兴起。还有一件新事，就是那些对天主的盟约是外人、与祂的应许无分、远离真理的人，竟因着一种神圣的能力得以拥抱真理。这些事不是骗子的事业，而是天主的事业，祂派遣了祂的圣言\Person{耶稣基督}，以显明祂的旨意。\Person{耶稣}以如此刚毅和温良所忍受的苦难与死亡，显示了施加者的残忍与不义，但它们并没有摧毁天主旨意的宣告；反而，我们可以说，它们倒有助于使这些旨意得以显明。因为\Person{耶稣}亲自教导了我们这一点，当祂说：\BibleQuote{除非一粒麦子落在地里死了，它仍然独自存在；但如果死了，就结出许多果实。}因此，作为这粒麦子的\Person{耶稣}死了，并结出了许多果实。而圣父一直在期待这粒麦子之死的结果，无论是现在正在产生的，还是将来要产生的。所以\Person{耶稣}的圣父是一位温柔慈爱的父，尽管\BibleQuote{祂没有怜惜自己的儿子，反而把祂交出}作为祂的羔羊\BibleQuote{为我们众人}，以至\BibleQuote{天主的羔羊}藉着为众人死去，能\BibleQuote{除免世罪}。因此，祂不是出于强迫，而是出于自愿，忍受了那些辱骂祂之人的谴责。然后，\Person{凯尔苏斯}回到那些对神像使用辱骂言语的人，说道：\Quote{对于那些你们加以侮辱的人，你们同样也可以说他们自愿接受这种对待，因此他们耐心忍受侮辱；因为最好对双方一视同仁。然而这些人严厉惩罚嘲笑者，以致他必须要么逃跑躲藏，要么被抓获而丧命。}因此，基督徒招惹魔鬼的报复，并不是因为他们对魔鬼加以侮辱，而是因为他们把魔鬼从神像中、从人的身体和灵魂中驱赶出去。这里，尽管\Person{凯尔苏斯}没有察觉到，他在这件事上却说了几分真话；因为确实如此：那些谴责基督徒、出卖他们、并以迫害他们为乐的人的灵魂，充满了邪恶的魔鬼。\par

% structure: p0088 source=h2 target=subsection reason=relative-heading-rank; base=h2

\subsection{第四十四章}

当那些为基督信仰而死的人的灵魂带着极大的荣耀离开身体时，他们摧毁了魔鬼的权势，挫败了他们针对人的阴谋。因此我认为，既然魔鬼从经验中得知他们被为真理而殉道者击败并压制，便害怕再次诉诸暴力。因此，在他们忘记自己所遭受的失败之前，世界很可能会与基督徒保持和平。但是当他们恢复力量，并以被罪恶蒙蔽的眼睛，想要再次向基督徒复仇并迫害他们时，那么他们又将失败，而那些虔敬之人的灵魂——他们为虔敬的事业舍弃生命——将彻底摧毁邪恶者的军队。既然魔鬼觉察到那些因宗教而胜利地面对死亡的人摧毁了他们的权威，而那些在苦难中屈服并否认信仰的人则落入他们的权势之下，我认为他们有时在基督徒受审时对他们深感兴趣，并竭力争取他们归向自己一边，因为他们感到基督徒的宣认对他们来说是折磨，而基督徒的否认则是他们的安慰和鼓励。同样的情感痕迹也可见于审判官的态度：他们看见那些忍耐暴行和折磨的人时非常苦恼，而看见基督徒在其中屈服时则极为得意。然而，这并非出于人性。他们清楚地看到，当那些被折磨压倒的人的\Quote{舌头}可以起誓时，\Quote{内心却未宣誓}。这可以作为对\Person{凯尔苏斯}评论的回答：\Quote{但他们严厉惩罚侮辱他们的人，以致他必须要么逃跑躲藏，要么被抓获而丧命。}如果一个基督徒有时逃跑，那不是出于恐惧，而是服从他师傅的命令，以便保全自己，并运用他的力量为他人谋益。\par

% structure: p0090 source=h2 target=subsection reason=relative-heading-rank; base=h2

\subsection{第四十五章}

让我们看看瑟尔苏接下来所说的话。内容如下：\Quote{有什么必要收集所有由祭司、女祭司以及其他人——无论是男人还是女人——在神圣感召下发出的神谕回答？所有从内殿传出的奇妙声音？所有启示给那些咨询祭品的人的知识？以及所有通过其他征兆和奇事传达给人的知识？有些神以可见的形式向人显现。世界充满了这样的实例。有多少城市是按照神谕的命令建造的？又有多少次，同样因神谕而从疾病和饥荒中得救？再者，有多少城市因忽视或遗忘这些神谕而悲惨地灭亡？有多少殖民地因遵循神谕的命令而建立并繁荣？有多少王公和普通人因此而得福或遭祸？有多少人为无子而哀伤，却得到了所求的祝福？有多少人使自己避开了魔鬼的愤怒？有多少肢体残缺的人得到了复原？还有，有多少人因对庙宇不敬而立即受到惩罚——有的立刻发狂，有的公开认罪，有的结束了自己的生命，有的染上不治之症！是的，有些人被内殿传来的可怕声音杀死。} 我不知道瑟尔苏为何把这些当作确凿的事实，而同时却把那些记载并流传给我们的、在犹太人中间发生或由\Person{耶稣}及其门徒所行的奇事仅仅视为神话。因为为什么我们的记载不能是真实的，而瑟尔苏的却是神话和虚构？至少，\Person{德谟克利特}、\Person{伊壁鸠鲁}和\Person{亚里士多德}的追随者并不相信后者，尽管这些希腊学派如果遇到行这些奇事的\Person{梅瑟}或任何先知，或者\Person{耶稣基督}本人，或许会被支持我们奇迹的证据所说服。\par

% structure: p0092 source=h2 target=subsection reason=relative-heading-rank; base=h2

\subsection{第四十六章}

据说阿波罗的女祭司有时会因贿赂而在回答中受影响；但我们的先知因其朴素的真理而受到同时代者和后来者的钦佩。因为通过先知所宣告的命令，城市得以建立，人得到治愈，瘟疫得以止息。事实上，整个犹太民族按照神圣的神谕从埃及迁往巴勒斯坦。当他们遵从天主的命令时，他们便兴旺；当他们偏离时，便遭受失败。有什么必要列举圣经历史中所有那些因遵从或蔑视先知的话而或好或坏的王公和普通人呢？如果我们要提到那些因无子而痛苦，但在向万有的创造主祈祷后成为父母的人，请阅读\Person{亚巴郎}和\Person{撒辣}的故事，他们在年老时生下了\Person{依撒格}，他是整个犹太民族之父：还有其他类似的例子。也请阅读\Person{希则克雅}的故事，他不仅按照\Person{依撒意亚}的预言从疾病中康复，而且还有勇气说：\Quote{以后我要生养儿女，他们要宣扬祢的公义。} 在\BibleBook{列王纪}{下卷}中，我们读到\Person{厄里叟}先知向一位好客地接待他的妇人宣告，因天主的恩宠她将生一个儿子；并且通过\Person{厄里叟}的祈祷，她成为了母亲。许多肢体残缺的人被\Person{耶稣}治愈。\BibleBook{玛加伯}{}也记载了那些胆敢亵渎耶路撒冷圣殿中犹太礼仪的人所受的惩罚。\par

% structure: p0094 source=h2 target=subsection reason=relative-heading-rank; base=h2

\subsection{第四十七章}

但希腊人会称这些记载为神话，尽管有两个完整的民族为它们的真实性作证。然而，我们为什么不能认为希腊人的记载是神话而不是我们的呢？也许有人不愿盲目接受自己的陈述而拒绝他人的，在经过仔细考察后，可能会得出这样的结论：希腊人所提及的奇事是由某些魔鬼执行的；犹太人中间的奇事是由先知或天使，或天主藉天使完成的；而基督徒所记载的奇事则是由\Person{耶稣}本人，或藉祂的能力在祂的宗徒身上完成的。那么，让我们比较所有这些记载；让我们考察行奇事者的目的和意图；并让我们探究接受这些善行的人身上产生了什么效果——是有益还是有害，或者两者都不是。古代犹太民族，在他们得罪天主、因大恶而被祂弃绝之前，显然是一个极其智慧的民族。但基督徒，他们以如此奇妙的方式形成了一个团体，起初似乎更多是被奇迹而不是被劝诫所诱导，从而离弃了他们祖先的制度，接受了对他们来说完全陌生的制度。事实上，如果我们从可能性来推理基督教社会的首次形成，我们不得不说，\Person{耶稣基督}的宗徒们，那些没有学识、出身卑微的人，若不是因所赐予他们的能力和伴随他们言语并使言语有效的恩宠，他们是不可能鼓起勇气向人们宣讲基督教真理的；而那些听到他们的人，若不是被某种非凡的能力和奇迹事件的力量所感动，也不会放弃他们祖先根深蒂固的习俗，并被诱导接受与他们所受教育如此不同的观念。\par

% structure: p0096 source=h2 target=subsection reason=relative-heading-rank; base=h2

\subsection{第四十八章}

接着，\Person{塞耳苏斯}在提及人们宁愿战死也不愿背弃基督教的热情之后，奇怪地加上了一些评论，想表明我们的教义与异教秘仪中司祭所传授的相似。他说：\Quote{仁兄，正如你相信永恒的惩罚一样，那些解释并引入神圣奥秘的司祭也是如此。你用来威胁他人的同样惩罚，他们也用来威胁你。现在值得考查的是，两者中哪一个更确凿地确立为真理；因为双方都同样确信真理在他们一边。但若要求证据，异教神祇的司祭提出了许多清晰而令人信服的证明，一部分来自恶魔所行的奇迹，一部分来自神谕的答复，以及其他各种占卜方式。}于是他想要我们相信，我们和奥秘的解释者同样教导永恒惩罚的教义，并且值得探究的是双方哪一边拥有真理。现在我要说，真理属于那些能使听众按照所听真理生活的人。但犹太人和基督徒正是被他们所持关于我们所说的来世、义人的赏报和恶人的惩罚的教义所感动。那么，让\Person{塞耳苏斯}或任何愿意的人指出，异教司祭和秘仪引导者的教导中，有谁曾这样被永恒惩罚所感动？因为揭示这一教义的目的，不仅是要论述惩罚并使人恐惧，还要使听到真理的人竭尽全力对抗那些导致惩罚的罪过。那些认真研究预言、不满足于粗略阅读其中预测的人，会发现这些预言足以使明智而真诚的读者相信，天主之神在那些人内，并且与他们的著作相比，恶魔的一切作为、神谕的答复或占卜者的言论，都无法相提并论。\par

% structure: p0098 source=h2 target=subsection reason=relative-heading-rank; base=h2

\subsection{第四十九章}

让我们看看\Person{塞耳苏斯}接着用些什么话对我们说：\Quote{此外，你们一方面如此重视身体——期望同样的身体复活，仿佛它是我们中最好最宝贵的部分；另一方面又将其暴露于酷刑，仿佛它毫无价值，这岂不是极其荒谬和矛盾的吗？但持这种见解、如此执着于身体的人，不值得与之争辩；因为在此及其他方面，他们表现出自己是粗俗、不洁，并毫无理由地背离共同信念。但我要向那些希望藉灵魂或心智（无论他们称其为精神实体、智能精神、圣洁蒙福的，还是活的灵魂，或是神圣非受造本性的属天不朽产物，亦或任何他们用以指称人精神本性的名称）与天主共享永生的人说话。他们正确地相信，生活善的人将受祝福，不义的人将遭受永远的惩罚。从这一教义中，他们和任何人都永不应偏离。}现在，正如他经常已经指责我们关于复活的意见，并且我们在这些场合已经以我们认为合理的方式辩护了我们的意见，我们不打算在每次重复一个反对意见时重复我们的辩护。\Person{塞耳苏斯}对我们提出无根据的指控，当他将这样的意见归于我们时：\Quote{我们复杂的本性中没有任何比身体更好更宝贵的东西；}因为我们坚持认为，远远超过所有身体的是灵魂，特别是理性的灵魂；因为正是灵魂，而非身体，带有造物主的肖像。因为，按照我们，天主不是有形体的，除非我们陷入\Person{芝诺}和\Person{克吕西普}追随者的荒谬错误。\par

% structure: p0100 source=h2 target=subsection reason=relative-heading-rank; base=h2

\subsection{第五十章}

但由于他指责我们过于关心身体，让他知道，当那种感觉是错误时，我们不分享它；当它是无关紧要时，我们只渴望天主应许给义人的东西。但\Person{塞耳苏斯}认为我们自相矛盾，既认为身体值得从天主那里得尊荣因而希望其复活，同时又将其暴露于酷刑仿佛不值得尊重。但当然，为虔敬而受苦、为美德而选择磨难，对身体来说并非不光彩；不光彩的是身体在恶习放纵中消耗力量。因为天主圣言说：\Quote{什么是可敬的种子？人的种子。什么是可鄙的种子？人的种子。}此外，\Person{塞耳苏斯}认为他不应该与那些盼望身体好处的人争辩，因为他们不合理地执着于一个永远无法满足期望的目标。他还称他们为粗俗不洁的人，一心制造不必要的纷争。但诚然，作为更高人性的一员，他应当甚至帮助粗鲁堕落的人。因为社会并不像排除非理性动物那样排除粗野未开化的人，而是我们的造物主使我们与全人类处于同一共同水平。因此，甚至与粗俗未开化的人争辩，并尽可能将他们提升到更高层次的教养——将不洁者带到可行范围内的最高纯洁度——将无理性的群众带到理性，将心灵患病者带到精神健康，这并不是有失身份的事。\par

% structure: p0102 source=h2 target=subsection reason=relative-heading-rank; base=h2

\subsection{第五十一章}

其次，他赞许那些\Quote{希望灵魂或心智（或按各种称谓：属神的本性、理性的灵魂、理智的、圣洁的、蒙福的）能与天主共享永生}的人；他也承认以下教理的正确性：\Quote{善人将获福乐，不义之人将受永罚。}然而，比\Person{色尔苏斯}所说过的任何话更令我惊异的是他随后所加的：\Quote{无论是他们或任何人，都不可背离这个教理。}因为在他撰写反对基督徒的著作时——基督徒信仰的本质就是天主，以及基督对义人的应许和对恶人等候惩罚的警告——他必定明白：如果一个基督徒因他的论点而被迫放弃基督教，那么毫无疑问，他在放弃基督教信仰的同时，也会抛弃他所说的任何基督徒和任何人都不可背离的那个教理。但我认为，就对待同胞的关怀而言，\Person{色尔苏斯}远逊于\Person{克吕西波}在其著作\WorkTitle{论克服情欲}中的表现。因为当\Person{克吕西波}试图对压迫和扰乱人类精神的情感与情欲施以疗治时，在使用了自认为有力的论证之后，他并不回避在第二和第三位使用一些他自己并不赞同的论证。他说：\Quote{因为假如有人主张有三种善，我们就必须按照那种假设来调节情欲；并且我们不可过于好奇地探究一个人在情欲影响下所持的意见，免得如果我们为了推翻占据心灵的意见而耽搁太久，疗治情欲的时机就会错过。}他又补充说：\Quote{因此，假若快乐是至善，或者一个被情欲支配的人持有这种意见，我们仍应同样帮助他，并向他表明：即使以快乐为人的最高和最终之善的原则来看，一切情欲也是不被允许的。}同样，\Person{色尔苏斯}在接受了\Quote{义人将蒙福，恶人将受永罚}这一教理之后，本应继续他的论题；在提出了自以为是主要论证之后，本应进一步用更多理由来证明这个真理：不义之人必定受永罚，而行善之人将蒙福。\par

% structure: p0104 source=h2 target=subsection reason=relative-heading-rank; base=h2

\subsection{第五十二章}

因为我们这些曾被许多（甚至可以说无数）论证说服去过基督徒生活的人，尤其渴望尽最大可能引导所有人接受基督徒真理的完整体系；但当我们遇到那些因针对基督徒的诽谤而怀有偏见，并因认为基督徒是一群不敬神之人而不肯听从任何愿意教导他们天主圣言原理的人时，我们就本着共同的人性原理，尽我们最大的努力说服他们接受恶人受罚的教理，并诱导甚至那些不愿成为基督徒的人也接受那个真理。当我们看到我们关于健康道德生活所教导的许多东西也被我们信仰的敌人所教导时，我们便如此急切地劝服他们接受正直生活的赏报。因为你会发现他们并未完全丧失关于是非、善恶的普遍观念。因此，愿所有的人，当他们仰望宇宙时，观察那恒定运转的星宿、行星的逆行、大气的构成及其对动物（尤其是人类）需要的适应，以及所有无数为人类福祉而设的巧妙安排；然后，在如此思考宇宙的秩序之后，愿他们谨防做出任何令这宇宙、灵魂及其理智本性的创造者不悦的事；并愿他们确信：惩罚将加于恶人，赏报将赐予义人，由那位按各人所应得对待一切，并按各人所行的善事及其所能呈上的自述而施予赏报的天主。愿所有的人都知道：善人将晋升到更高的境界，恶人将因他们的纵欲堕落、怯懦胆怯以及一切愚行而交付给痛苦与折磨。\par

% structure: p0106 source=h2 target=subsection reason=relative-heading-rank; base=h2

\subsection{第五十三章}

关于这个问题已经说了很多，现在让我们来看\Person{策尔苏斯}的另一个陈述：\Quote{既然人出生时与身体结合，无论是为了适合宇宙的秩序，还是为了使他们以那种方式遭受罪的惩罚；或者因为灵魂被某些情欲所压迫，直到在指定的时期中从这些情欲中得到净化——因为按照\Person{恩培多克勒}的说法，所有人类必须从福乐之所被放逐三万期——因此我们必须相信，他们被委托给某些存有，作为这个监狱的看守。}你会注意到，\Person{策尔苏斯}在这些话中，用可疑的人类猜测的语言谈论如此重大之事。他还添加了关于人类起源的各种意见，并且非常不愿意将这些意见中的任何一个判定为虚假。当他一旦得出结论，既不加区分地接受，也不轻率地拒绝古人所持的意见时，那么按照同样的判断规则，当他发现自己不愿相信犹太先知和\Person{耶稣}所教导的教义时，不是至少应该将它们视为尚待探究的事吗？难道他不应该考虑一下，一个忠心事奉至高天主、时常遭遇无数危险甚至死亡，也不愿牺牲天主的荣耀和他们所相信的天主旨意启示的民族，是否很可能完全被天主所忽视？难道不是更应认为，那些蔑视以人为艺术表现神圣存有、却努力在思想中提升自己以达到对至高者认识的民族，理应得到来自天主本身的某种启示吗？此外，他还应该考虑到，万有的共同圣父和创造者，他看见并听见一切，并适当评估每个寻求他、渴望事奉他的人的意图，也将赐予他们一些他宰制的益处，并将扩展他们对他本身的认识，这认识他曾一次赐予他们。如果\Person{策尔苏斯}和其他憎恨\Person{梅瑟}、犹太先知、\Person{耶稣}以及他的忠实门徒（他们为他的圣言忍受了如此之多）的人记起了这一点，他们就不会这样辱骂\Person{梅瑟}、先知、\Person{耶稣}和他的宗徒；他们也不会独独蔑视犹太人超过世上所有民族，说他们甚至比埃及人更坏——埃及人，无论出于迷信还是其他形式的错觉，尽可能地将神圣存有贬低到牲畜的水平。我们邀请人探究，并非因为我们希望引导任何人对基督宗教的真理产生怀疑，而是为了表明，那些从各方面辱骂基督宗教教义的人，至少暂停判断，不要如此轻率地陈述关于\Person{耶稣}和他的宗徒他们所不知道、也无法证明的事——无论用斯多葛派所谓的\Quote{\OriginalTerm{领悟知觉}{apprehensive perception}}，还是用不同哲学派别用作真理标准的方法。\par

% structure: p0108 source=h2 target=subsection reason=relative-heading-rank; base=h2

\subsection{第五十四章}

\Person{塞耳苏}又说：\Quote{因此我们必须相信，人被托付给某些看守这囚室的存在，}我们的回答是：那些被\Person{耶肋米亚}称为\Quote{地上的囚徒}的灵魂，当热切追求德行时，甚至在今生就脱离了邪恶的束缚；因为\Person{耶稣}宣告了这一点，正如先知\Person{依撒意亚}在祂来临之前许久所预言的，他说：\BibleQuote{囚徒要出去，黑暗中的人要显现。}\Person{耶稣}亲自，正如\Person{依撒意亚}也曾预言祂，升起为\BibleQuote{照亮坐在黑暗和死影中的人的光}，因此我们可以说：\BibleQuote{让我们挣断他们的捆绑，从我们身上甩掉他们的绳索。}如果\Person{塞耳苏}和那些像他一样反对我们的人，能够深入钻研福音的叙述，他们就不会劝我们信赖他们所谓的\Quote{囚室的看守者}。福音记载：一个女人佝偻着，完全不能直起腰来。\Person{耶稣}看见她，知道她是什么原因佝偻的，就说：\BibleQuote{这亚巴郎的女儿，被\Person{撒殚}捆绑了十八年，难道不该在安息日解开她的束缚吗？}还有多少人仍被\Person{撒殚}压制和捆绑，他阻止他们仰望高处，也使我们低头向下！除了那藉\Person{耶稣基督}而来的圣言，祂从前就曾感动先知，没有人能扶起他们。\Person{耶稣}来释放那些受魔鬼管辖的人；关于祂，祂意味深长地说：\BibleQuote{现在，今世的元首受了审判。}因此，我们并非无故诽谤魔鬼，而是谴责他们在世上的活动对人类具有毁灭性，并表明他们借着神谕、身体治愈和其他类似手段，试图使那下降到这\Quote{卑贱的身体}中的灵魂远离天主；那些感受到这种卑贱的人呼喊：\BibleQuote{我真不幸啊！谁能救我脱离这死状的身体？}因此，我们暴露我们的身体受鞭打和折磨并非徒然；如果通过这种方式，一个人可以避免将神的名号归于那些地上的神灵，他们与崇拜者结合，引人走向毁灭，那么忍受这样的苦难当然不是徒然的。实际上，我们认为，为德行而受苦，为敬虔而受酷刑，乃至为圣洁而受死，是合理且中悦天主的；因为\BibleQuote{在天主眼中，祂的圣人的死是宝贵的}；我们主张，克服贪生就是享受极大的善。但当\Person{塞耳苏}将我们比作那些因其罪行而公正受罚的臭名昭著的罪犯，并且毫不顾忌地将我们所持的这一可称赞的目的与罪犯的顽固相提并论时，他就成了那些将\Person{耶稣}视为罪犯之人的兄弟和同伴，应验了圣经所说：\BibleQuote{祂被列于叛逆者之中。}\par

% structure: p0110 source=h2 target=subsection reason=relative-heading-rank; base=h2

\subsection{第五十五章}

\Person{塞耳苏}接着说：\Quote{他们必须在两个选择中做出抉择。如果他们拒绝向众神履行应有的服务，并尊重那些被派管理这服务的人，那么就不要让他们长大成人，或娶妻生子，或参与任何生活事务；而是让他们尽快离开，不留后代，好让这类种族从地面上灭绝。或者，另一方面，如果他们娶妻、养育子女、品尝地上的果实、分享生活的一切祝福、并承受指定的苦难（因为大自然本身已为所有人分配了苦难；因为苦难必然存在，而世界是它们唯一的地方），那么他们就必须履行生活的职责，直到从束缚中释放，并向那些控制今生事务的存在表示应有的敬意，免得显得忘恩负义。因为他们在接受这些存在分施的好东西之后，却不回报他们，这是不公平的。}对此我们回答说，我们认为没有好的理由离开这个世界，除非当虔敬和德行要求如此时；例如，当那些被任命为审判官、自以为有权控制我们生命的人，将选择摆在我们面前：要么违反\Person{耶稣}的诫命而活着，要么继续服从祂而死亡。但天主允许我们结婚，因为并非所有人都适合那更高的、即完全纯洁的生活；天主要我们养育我们所有的孩子，不毁灭祂眷顾所赐的任何后代。这与我们不服从地上魔鬼的目的并不冲突；因为\BibleQuote{我们穿上了天主的全副武装，屹立不动}，作为虔敬的运动员，对抗那阴谋反对我们的魔鬼种族。\par

% structure: p0112 source=h2 target=subsection reason=relative-heading-rank; base=h2

\subsection{第五十六章}

因此，虽然刻尔苏斯用他自己的话说，要\Quote{尽快将我们逐出生命}，好使\Quote{这样一个种族从地上灭绝}；然而我们，连同那些崇拜造物主的人，将按天主的法律生活，绝不顺从罪恶的法律。我们若愿意就结婚，养育在婚姻中赐给我们的子女；若有需要，我们不仅会享受生命的福乐，也会承受命定的忧苦作为对灵魂的考验。因为天主圣经惯常这样谈论人的苦难，借此，如同金子要在火中受炼，人的精神也受试探，并被判定为堪受谴责或赞美。所以我们准备好面对刻尔苏斯称为邪恶的那些事，也准备说：\Quote{主啊，求祢试探我，考验我，炼净我的肺腑和心怀。}因为\Quote{若不按规矩争战}，就没有人会在这地上，在这卑微的身体里\Quote{得冠冕}。再者，我们并不向刻尔苏斯所说的那些管理世事者献上以为当得的尊敬。因为我们敬拜主我们的天主，单单事奉祂，并愿效法基督；当魔鬼对祂说：\Quote{你若俯伏朝拜我，我必将这一切都给你}时，祂用这话回答说：\Quote{你当朝拜主你的天主，单单事奉祂。}因此，我们不向那些按刻尔苏斯所说管理世事者献上以为当得的尊敬；因为\Quote{没有人能事奉两个主人}，我们\Quote{不能事奉天主而又事奉\OriginalTerm{玛门}{mammon}}，不论这名称是指一个还是多个。此外，若有人\Quote{因违犯法律而羞辱立法者}，在我们看来，如果天主的法律和玛门的法律完全对立，那么，我们宁可因违犯玛门的法律而羞辱玛门，以便藉遵守天主的法律而尊荣天主，胜过因违犯天主的法律而羞辱天主，以便藉顺从玛门的法律而尊荣玛门。\par

% structure: p0114 source=h2 target=subsection reason=relative-heading-rank; base=h2

\subsection{第五十七章}

刻尔苏斯以为，人在\Quote{被解除生命的束缚之前}，按照普遍接受的习俗，向国家所承认的每一位神明献祭，就是\Quote{履行了生活的责任}；他看不出真诚的虔诚所履行的真正责任。因为我们说，那真正履行生活责任的人，是常常记念谁是自己的创造者、什么事合乎祂的旨意，并在一切事上行事以讨天主喜悦的人。再者，刻尔苏斯要我们感谢这些邪魔，以为我们欠他们感谢之祭。但我们虽然承认感谢的责任，却坚持说，我们拒绝向那些不帮助我们、反而在我们既不献祭也不朝拜他们时与我们作对的东西表示感谢，并不是忘恩负义。我们更担心的是，我们不要对天主忘恩负义，祂以恩惠厚待我们，我们是祂的化工，无论我们处于何种境况，祂都眷顾我们，并赐给我们超越今生的希望。我们有一个向天主感恩的标记，就是在饼中，我们称之为圣体圣事。此外，正如我们先前已指明的，邪魔并不掌管那些为我们受造之物；因此，我们享用受造之物，却拒绝向那些与这些事物无关者献祭，并没有犯什么过错。况且，我们知道管理地上果实和动物生育的，不是邪魔，而是天使；我们赞美并祝福天使，因为他们被天主指派管理我们人类所需之物；然而，即使对他们，我们也不献上当属于天主的尊敬。因为这不中悦天主，也不会使受托管理这些事的天使们喜悦。事实上，我们不向他们献祭，比向他们献祭更令他们喜悦；因为他们并不贪求从地上升起的祭肉香气。\par

% structure: p0116 source=h2 target=subsection reason=relative-heading-rank; base=h2

\subsection{第五十八章}

刻尔苏斯接着说：\Quote{让人去询问埃及人，他会发现一切事物，甚至最微小的事物，都托付给某个邪魔照管。人的身体被分为三十六个部分，有同样多的空中邪魔被指派照管它，每一个负责不同的部分，虽然其他人使数目更大。所有这些邪魔在埃及语言中有不同的名称：如\OriginalTerm{克努门}{Chnoumen}、\OriginalTerm{克纳库门}{Chnachoumen}、\OriginalTerm{克纳特}{Cnat}、\OriginalTerm{西卡特}{Sicat}、\OriginalTerm{比乌}{Biou}、\OriginalTerm{埃鲁}{Erou}、\OriginalTerm{埃雷比乌}{Erebiou}、\OriginalTerm{拉马诺尔}{Ramanor}、\OriginalTerm{雷亚诺尔}{Reianoor}，以及其他类似的埃及名字。此外，他们呼求这些邪魔，就治好了身体特定部位的疾病。那么，有什么能阻止一个人尊敬这些或其他邪魔呢？如果他宁愿健康而不愿生病，宁愿顺利而不愿倒霉，并尽可能摆脱一切灾祸和麻烦的话。}刻尔苏斯就这样试图使我们的灵魂卑下，去崇拜邪魔，其假定是这些邪魔占有我们的身体，每一个对某个肢体拥有权柄。他以此为由要我们信靠他所说的这些邪魔，并事奉他们，以便我们健康而非生病、顺利而非倒霉，并尽可能避免一切灾祸和麻烦。至高天主的尊荣是不能分割或与他人分享的，他却如此轻看，以至于不能相信天主的能力——若呼求祂并崇高尊敬祂，祂就能赐给事奉祂的人一种权能，用以抵御邪魔对义人的攻击。因为他从未见过\Quote{耶稣之名}的能力——当真正信靠的人呼求时，能拯救不少人不只脱离邪魔和附着邪魔的情况，也脱离其他灾祸。\par

% structure: p0118 source=h2 target=subsection reason=relative-heading-rank; base=h2

\subsection{第五十九章}

大概那些赞同克耳苏斯观点的人，当听到我们说：\Quote{\BibleQuote{因耶稣之名，一切在天上的、地上的和地下的，无不屈膝，一切唇舌}都宣认\BibleQuote{耶稣基督是主，以光荣天主圣父。}}时，会嘲笑我们。但尽管他们可能嘲笑这样的说法，他们却会得到比克耳苏斯所举的克努门、克纳库门、克纳特、西卡特以及其他埃及名目（他提到他们被呼求，并医治人体各部分的疾病）更有说服力的论据。请注意，他如何试图使我们通过耶稣基督转离对万有之天主的信仰，却劝我们为了身体的福祉去相信三十六个蛮族魔鬼，这些魔鬼只有埃及术士以某种未知的方式呼求，并应许我们巨大的回报。按照克耳苏斯的说法，我们现在最好献身于魔法和巫术，而不是接受基督信仰，相信无数魔鬼，而不是相信全能的、生活的、自我启示的天主，祂藉着祂（那以伟大能力将真正圣洁原则传播给全世界所有人的人）显示了自己；是的，我可以毫不夸张地说，祂已将这种知识赐给所有有理智、需要从罪孽的瘟疫和败坏中获得解脱的受造物。\par

% structure: p0120 source=h2 target=subsection reason=relative-heading-rank; base=h2

\subsection{第六十章}

然而，克耳苏斯怀疑他此处所给教导的倾向是导向魔法，并担心这些话可能带来危害，于是补充说：\Quote{必须小心，免得任何人因使自己的心智习惯于这些事而过分沉迷其中，并因过度关注身体而使心智转离更高的事物，让它们被遗忘。因为或许我们不应轻视那些智者的意见，他们说大多数\OriginalTerm{地魔}{earth-demons}沉溺于肉欲、鲜血、气味、甜美的声音和其他诸如此类感官之物；因此他们只能医治身体，或预言人及城邦的命运，并做其他与此有死生命相关的事情。}如果在此方向上真有如此危险，甚至真理之敌也承认，那么，将我们自己完全交托给藉耶稣基督（祂给了我们这样的教导）而至高的天主，并向祂祈求所有帮助，以及圣善天使的保护以抵御沉迷于情欲、鲜血、祭品气味、奇异声音和其他感官事物的地灵，从而避免过分沉迷于这些魔鬼的能力、转离更高事物、因过度关注身体而让它们被遗忘的种种危险，岂不是更好吗？因为即使按照克耳苏斯的承认，它们除了医治身体外什么也不能做。但事实上，我要说，这些魔鬼即便受到极大的尊敬，是否真能医治身体也并不清楚。在寻求病愈时，一个人要么遵循更普通和简单的方法，求助于医术；要么，如果他想要超越通常的人类方法，就必须上升到一个更高更好的方式——通过虔敬和祈祷向万有之天主祈求祝福。\par

% structure: p0122 source=h2 target=subsection reason=relative-heading-rank; base=h2

\subsection{第六十一章}

因为请你自己思量，哪种心灵态度更能取悦于至高者（祂的权能至高无上且普遍，并为人类在身体、心智和外在事物上的福祉而引导一切）——是那在所有事上都委身于天主的人，还是那好奇探究魔鬼的名字、力量和作用、咒语、适合它们的草药、以及刻有符合它们传统形态（象征性或其他方式）铭文的石头的人？即使是最不聪明的人也明白，那心地单纯、不好奇探究、却在所有事上委身于神圣旨意的心，最蒙天主和所有肖似天主者的喜悦；但那为了身体健康、身体享乐和外在福乐而忙于魔鬼的名字，并探究用何种咒语来平息它们的心，将被天主谴责为邪恶和不虔敬，且更合于魔鬼而非人的本性，并会被交给魔鬼撕咬折磨。因为很可能它们——作为邪恶的受造物，并且如克耳苏斯所承认，沉溺于鲜血、祭品气味、甜美的声音等——不会信守它们对那些供给它们这些事物的人最庄严的承诺。因为如果别人为了对抗那些已经呼求它们的人而请求它们的帮助，并用更多的鲜血、气味以及它们所要求的供物来收买它们的青睐，它们就会反过来攻击那些昨天还向它们献祭并呈上悦人供物的人。\par

% structure: p0124 source=h2 target=subsection reason=relative-heading-rank; base=h2

\subsection{第六十二章}

在前一段中，塞尔苏斯曾详细谈论神谕之事，并指其答复为诸神之声；但现在他修正此说，承认\Quote{那些预言人类和城邦命运、关心凡俗事务者，乃是属地之灵，他们沉溺于肉欲、鲜血、气味、悦耳之声及其他此类事物，无法超越这些感官对象。}或许，当我们反对塞尔苏斯关于神谕及对所谓诸神崇敬的神学教导时，有人可能怀疑我们有不敬之罪，因为我们声称这些是魔鬼势力的诡计，为引诱人走向肉欲放纵。但任何对我们持此怀疑者，如今当看到上述来自一位公开反对基督教之人的作品段落——他终究被真理之灵所征服而写下此言——即可相信基督徒所提出的论述是有根据的。因此，尽管塞尔苏斯说\Quote{我们必须向他们献祭，只要有利于我们，因为不加辨别地献祭是不合理的}，但我们不可向沉迷于鲜血和气味的魔鬼献祭；也不可在思想中亵渎天主，将祂降低到邪恶魔鬼的水平。若塞尔苏斯曾仔细斟酌\Quote{有利}一词的含义，并考虑到真正的利益在于德行和善行，他就不会将\Quote{只要有利}这一短语应用于他所承认的这些魔鬼的侍奉。那么，如果身体的健康和人生的成功要以侍奉这类魔鬼为条件，我们宁可选疾病和不幸，同时保持真正忠于天主圣意的意识。因为，这比虽身体健康、世事亨通，但心灵患致命恶疾，因远离并弃绝天主而悲惨更好。我们宁愿向那毫不索取、只为人类及所有理性受造物的福祉而求助的天主求助，也不愿向那些喜悦鲜血和祭香气味的魔鬼求助。\par

% structure: p0126 source=h2 target=subsection reason=relative-heading-rank; base=h2

\subsection{第六十三章}

在如此谈论魔鬼及其对鲜血和祭香气味的喜爱之后，塞尔苏斯补充道，仿佛想收回他所提出的指控：\Quote{更正确的见解是，魔鬼无所欲求，一无所缺，但他们喜悦那些向他们尽虔敬之职的人。}若塞尔苏斯相信这是真的，他本应如此说，而非作出之前的陈述。但人性从未被天主及其独生子——真理——彻底遗弃。因此，就连塞尔苏斯也说了真话，指出魔鬼喜悦牺牲的血和烟；然而，由于他自身邪恶本性的力量，他又堕入错误，将魔鬼与那些严格履行一切义务——即使对不感恩者亦然——的人相比，而对感恩者则充满善行。此处，塞尔苏斯似乎陷入混乱。有时，他的判断被魔鬼的影响所蒙蔽，有时又从它们的迷惑势力中恢复，并获得一些真理的瞥见。因为他再次补充道：\Quote{我们绝不可在任何时候失去对天主的把握，无论白天黑夜，无论公开私下，无论言语行为，无论做什么或避免做什么。}依我理解，这是说：无论我们公开做什么，在所有行为中，在所有言语中，\Quote{让灵魂恒常专注于天主。}然而，似乎在与魔鬼的疯狂灵感进行辩论斗争后，他又被它们完全压倒，补充道：\Quote{既然如此，赢得地上统治者——无论是与我们的本性不同的存在，还是人间的君王和王子——的恩宠，有何害处呢？因为他们是通过魔鬼的媒介获得其尊位的。}在前一部分，塞尔苏斯竭尽全力贬低我们的灵魂，使之崇拜魔鬼；现在他又希望我们寻求君王和王子的恩宠，关于这些，世界和历史充满例证，我认为无需引用。\par

% structure: p0128 source=h2 target=subsection reason=relative-heading-rank; base=h2

\subsection{第六十四章}

因此，有一位我们应当寻求其恩宠、并向祂祈祷开恩于我们的——至高天主，祂的恩宠是藉着虔敬和实践一切德行获得的。若他想要我们在至高天主之后寻求他人的恩宠，让他想想：正如影子的运动跟随投射影子的物体，同样地，当我们获得天主的恩宠时，我们也获得所有与天主为友的天使和神灵的善意。因为他们知道谁配得上天主的赞许，他们不仅善待这些人，而且与他们合作，努力讨天主喜悦：他们代表这些人寻求天主的恩宠；他们将自己的祈祷和代祷与这些人的祈祷结合在一起。我们甚至可以大胆地说，向往更美好事物的人们，当他们向天主祈祷时，有成千上万的圣善力量站在他们一边。这些力量，即使未被请求，也与他们一同祈祷，为我们必死的种族带来援助，我敢说，甚至并肩作战：因为他们看到魔鬼激烈地争战，最猛烈地攻击那些献身于天主并蔑视魔鬼敌意的人的得救；他们看到魔鬼对那拒绝以牺牲的鲜血和烟气侍奉它们，反而极力在言语和行为上通过耶稣——祂曾四处\Quote{医治}并解救\Quote{所有被魔鬼压制的人}——与至高者和平共处的人，充满野蛮的仇恨。\par

% structure: p0130 source=h2 target=subsection reason=relative-heading-rank; base=h2

\subsection{第六十五章}

此外，我们应当鄙视讨好君王或任何其他人，不仅如果他们的恩宠是通过谋杀、放荡或残忍行为取得的，甚至如果这涉及对天主的不敬，或任何卑鄙的谄媚和阿谀之辞——这些事都不配勇敢和有原则的人，他们旨在将忍耐和刚毅这最高的德行与其他德行结合起来。然而，只要我们不做任何违反天主的法律和圣言的事，我们就不至于疯狂到激起君王和君侯的愤怒，从而招致苦难和折磨，甚至死亡。因为我们读到：\BibleQuote{每人要服从上级有权柄的人，因为没有权柄不是从天主来的，现有的权柄都是由天主规定的。所以谁反抗权柄，就是反抗天主的规定。}我们在对\BibleBook{}{罗马书}的阐释中，已尽己所能用了各种应用详细解释了这些话；但目前我们取其较明显且普遍接受的含义，以回应克耳苏的话，即\Quote{君王登上王权，并非没有魔鬼的力量。}关于君王和统治者的体制，这里可以谈很多，因为这是一个宽泛的主题，包括那些残酷暴虐地统治的统治者，以及那些将王权作为放纵奢侈和罪恶享乐之手段的统治者。因此，我们暂时跳过对此问题的全面讨论。然而，我们决不指着\Quote{君王的运气}或任何被认为等同于天主的事物起誓。因为如果\Quote{运气}一词不过是表达事件不确定过程的用语（如有些人所说，尽管他们似乎意见不一），我们就不指着那并不存在的事物作为天主起誓，好像它确实存在并能做什么，以免我们因誓言束缚于不存在的事物。另一方面，如果（如另一些人认为，他们说指着罗马君王的运气起誓就是指着他的魔鬼起誓）所谓的君王运气是在魔鬼的权下，那么在这种情况下，我们宁可死也不愿指着邪恶而背信弃义的魔鬼起誓，它常与它所控制的人一同犯罪，甚至犯罪更甚。\par

% structure: p0132 source=h2 target=subsection reason=relative-heading-rank; base=h2

\subsection{第六十六章}

然后克耳苏追随那些受魔鬼影响之人的榜样——时而痊愈，时而复发，仿佛他又变得清醒——说：\Quote{然而，如果任何敬拜天主的人被命令做任何不敬的事，或说任何卑鄙的话，这样的命令绝不应理会；我们必须忍受各种折磨，或接受任何形式的死亡，而不愿说甚或想任何不配天主的事。}然而，由于对我们的原则无知，且思想一片混乱，他又说：\Quote{但如果有人命令你颂扬太阳，或唱一首欢快的凯歌赞美弥涅耳瓦，你通过赞美它们，似乎就将更高的赞美归于天主；因为虔诚延伸到万物，就变得更加完美。}对此我们的回答是：我们并不等待任何命令才去颂扬太阳；因为我们受教导不仅要说得益于那些服从天主旨意的受造物，甚至也要说得益于我们的仇敌。因此我们赞美太阳，作为天主的荣耀化工，它服从祂的律法，听从\BibleQuote{太阳和月亮，请赞美上主}的召唤，并尽一切力量赞颂万有的圣父和造物主。然而，弥涅耳瓦——克耳苏将她与太阳并列——是各种希腊神话的主题，无论这些神话是否含有隐藏意义。他们说弥涅耳瓦全副武装地从朱庇特的脑中跳出；当伏尔甘追求她时，她为保持贞洁而逃走；伏尔甘情欲炽烈，种子落在地上，长出一个孩子，弥涅耳瓦抚养他，并称他为厄里克托尼俄斯，\par

他受那位蓝眼少女的抚养，\newline{} 但从丰饶的犁沟中诞生，\par

这位养育大地的大能后代。\par

因此显然，如果我们承认弥涅耳瓦是朱庇特的女儿，我们也必须承认许多寓言的虚构，这是任何摒弃寓言、寻求真理的人所不能接受的。\par

% structure: p0137 source=h2 target=subsection reason=relative-heading-rank; base=h2

\subsection{第六十七章}

至于以寓意方式看待这些神话，并视弥涅耳瓦为代表审慎，请任何人指出她历史中的实际事实，作为此寓意的基础。因为假设弥涅耳瓦作为古代的一个妇女，受到那些为追随者设立奥秘和礼仪、并希望她作为女神被称颂之人的尊崇，那么我们更被禁止对弥涅耳瓦献上神性的敬意，如果我们连如此荣耀的太阳都不准敬拜，尽管我们可以赞美它的荣耀。确实，克耳苏说\Quote{当我们唱赞美太阳和弥涅耳瓦的颂歌时，我们似乎对伟大的天主表示了更大的敬意}；但我们知道事实正好相反。因为我们只对至高者及其独生子——祂是圣言和天主——唱颂歌；我们赞美天主及其独生子，太阳、月亮、星辰和天上万军也同样如此。因为这些都组成神圣的合唱团，与世人中的义人一同赞美至高天主及其独生子。我们已经说过，我们不可指着人的君王或所谓的\Quote{君王的运气}起誓。因此无需再驳斥这些陈述：\Quote{如果你被命令指着人的君王起誓，这并没有错。因为地上一切都被赐给了他；你在此生所领受的一切，都是从他领受的。}然而我们否认地上万物都赐给了君王，或我们在此生所领受的一切都从他而来。因为我们所正当和光荣地领受的一切，都是从天主而来，并藉祂的眷顾，如成熟的果实，\BibleQuote{五谷——滋养人心，愉悦的葡萄树和葡萄酒——使人心里喜乐}，而且还有橄榄树的果实，使人容光焕发，这一切我们都从天主的眷顾领受。\par

% structure: p0139 source=h2 target=subsection reason=relative-heading-rank; base=h2

\subsection{第六十八章}

\Person{凯尔苏}接着说道：\Quote{我们必须服从那位古代作者，他很久以前说过：\InnerQuote{让一人为王，他是狡猾的\Person{克洛诺斯}之子所指定的；}}\newline{}并补充说：\Quote{如果你们放弃这一准则，你们将理应受国王的惩罚。因为如果所有人都像你们一样行事，那么他将完全陷入孤独和遗弃之中，地上的事务将落到最野蛮、最无法无天的蛮族手中；那时人间将不再有你们宗教的荣耀或真智慧的任何痕迹。}既然如此，\Quote{将只有一位主、一位王，}他不应是那个\InnerQuote{狡猾的\Person{克洛诺斯}之子所指定的}人，而是那赐予权柄者，祂\Quote{废王立王}，祂\Quote{在急难时兴起有用的人于地上}。\BibleRef{德 10:4}因为君王不是由那位\Person{克洛诺斯}之子任命的——根据希腊神话，他将父亲从王位上推翻并打入\Place{塔尔塔罗斯}（无论这寓意如何解释）——而是由统管万有的天主所任命，祂明智地安排与立王有关的一切。因此，我们确实放弃那行诗中的准则：\par

\Quote{他是狡猾的\Person{克洛诺斯}之子所指定的；}\par

因为我们知道，没有什么神或神之父会筹划任何弯曲或狡猾之事。但我们绝不否定上智的安排，以及万事直接或间接通过上智安排而发生的事实。国王不会\Quote{理应惩罚}我们，如果我们说，不是狡猾的\Person{克洛诺斯}之子赐给他王国，而是那\Quote{废王立王}者。但愿所有人都效法我，拒绝\Person{荷马}的准则，坚持王权的神圣起源，并遵守尊敬国王的诫命！在这种情况下，国王不会\Quote{完全陷入孤独和遗弃之中}，世界的\Quote{事务也不会落到最不虔敬和野蛮的蛮族手中}。因为，照\Person{凯尔苏}的话说，\Quote{他们若像我一样行事}，那么显然，即使是蛮族，当他们顺从天主的圣言时，也会成为最守法、最仁慈的人；每一种宗教崇拜都将被消灭，唯有基督的宗教得以存留，并最终得胜，因为它的原则一天天更多地占据人的心灵。\par

% structure: p0143 source=h2 target=subsection reason=relative-heading-rank; base=h2

\subsection{第六十九章}

于是，\Person{凯尔苏}似乎没有注意到他刚说过的话与现在的话自相矛盾，他之前说\Quote{如果所有人都像你们一样行事}，现在又补充说：\Quote{你们当然不是说，如果罗马人顺从你们的愿望，忽略他们对神和人的惯常义务，而去崇拜至高者，或随便你们怎么称呼祂，那么祂就会降临为他们争战，以致他们不需要其他帮助，只要祂的帮助。因为，正如你们自己所说，这位神曾向侍奉祂的人应许过这些和更多的事，且看祂是如何帮助了祂们和你们的！他们，非但没有成为全世界的统治者，反而连一块土地或一个家都没有留下；至于你们，如果你们中有人甚至在暗地里犯罪，就会被搜捕出来处以死刑。}所引发的问题是：\Quote{如果罗马人被说服采纳基督徒的原则，蔑视对公认之神和人的义务，而去崇拜至高者，那会发生什么？}以下是我对这个问题的回答。我们说：\BibleQuote{如果你们中间有两个人在地上同心合意地求什么事，义人的圣父在天上必为他们成就}，因为天主喜悦理性生物的合一，而厌恶纷争。那么，如果不仅是极少数人像现在这样同心合意，而是整个罗马帝国都如此，我们还能期待什么呢？因为他们将向圣言祈祷，圣言曾对受埃及人追击的希伯来人说：\BibleQuote{主必为你们争战，你们只管静默。}如果他们万众一心地祈祷，他们就能击退比那因\Person{梅瑟}祈祷主、以及与他一同祈祷的人所击败的多得多的敌人。那么，如果天主向遵守祂律法的人所应许的没有实现，这不归咎于天主的不信实。但祂使祂应许的成就取决于某些条件——即他们应遵守并按祂的律法生活；如果犹太人连一块土地或住所都没有剩下，尽管他们曾接受这些有条件的应许，那么全部责任都应归于他们的罪行，尤其是他们对待\Person{耶稣}的罪过。\par

% structure: p0145 source=h2 target=subsection reason=relative-heading-rank; base=h2

\subsection{第七十章}

但如果所有罗马人，如策尔苏所假设的那样，接受了基督信仰，那么当他们祈祷时，必将战胜敌人；更确切地说，他们将根本不会作战，因为有那神圣的力量保护他们，那力量曾应许因五十个义人的缘故而拯救五座整城。因为天主的人确实是地上的盐：他们维护世界的秩序；只要盐尚未腐败，社会就能维系下去：因为\BibleQuote{盐若失了味，便不配下地，也不配作肥料，只好丢在外面，被人践踏。有耳听的，听吧！}当天主允许试探者迫害我们时，我们就遭受迫害；当天主愿意我们免于受苦时，即使身处憎恨我们的世界，我们也享受奇妙的平安，信赖那说\BibleQuote{你们放心，我已战胜了世界}者的保护。祂确实战胜了世界。因此，世界只有在那位从父那里获得战胜世界之能力的意愿下才能得胜；因着祂的胜利，我们获得了勇气。即使祂愿意我们再次为信仰而争战奋斗，就让敌人来攻击我们吧，我们要对他们说：\BibleQuote{我靠着那加给我力量的基督耶稣，凡事都能做。}因为\BibleQuote{两只麻雀不是卖一个铜钱吗？但若没有你们在天之父的许可，它们中连一只也不会掉在地上}，正如圣经所说。天主的眷顾如此周全地涵盖一切，以至于我们连头发也都被祂数过了。\par

% structure: p0147 source=h2 target=subsection reason=relative-heading-rank; base=h2

\subsection{第七十一章}

策尔苏再次像往常一样糊涂，把一些我们中间从未有人写过的东西归咎于我们。他的话是：\Quote{你们说，如果你们现在的统治者接受了你们的意见，被敌人俘获，你们还能说服后来的统治者；等这些也遭俘获后，你们再说服他们的继任者，如此类推，直到最后，所有被你们说服的人都被俘获了，那时就会有一位精明的统治者出现，预见即将发生的事，在自我毁灭之前先把你们全部消灭。}对这些指控，无需任何答复：因为我们中间从未有人说，如果现在的统治者接受了我们的意见而被敌人俘获，我们还能说服他们的继任者；等这些也遭俘获，再说服后来的人，如此类推。但他根据什么断言，当一连串被我们说服的人因未能击退敌人而被俘获后，会有一位精明的统治者出现，预见即将发生的事，并将我们彻底消灭？在这里，他似乎乐于捏造并说出最荒谬的胡言乱语。\par

% structure: p0149 source=h2 target=subsection reason=relative-heading-rank; base=h2

\subsection{第七十二章}

其后他说：\Quote{如果可能，}同时暗示他认为最可取的，\Quote{亚细亚、欧罗巴、利比亚的所有居民，希腊人和外邦人，直到地极的尽头，都归于一法律；}但判断这根本不可能，他又说：\Quote{任何人若认为这可行，便是无知。}这需要仔细考虑和长篇论证，以证明这不仅可能，而且必然实现：所有有理性的受造物都将归于一法律。然而，若我们必须提及此事，将极为简要。斯多葛派确实认为，当最强的元素占据主导时，一切事物都将变成火。但我们的信念是：圣言将统御整个理性受造界，并将每一个灵魂改变成祂自己的完美；在这种状态下，每个人只需运用自己的能力，就能选择所欲之物，并得其所择。因为虽然身体上的疾病和创伤有些是医术无法治愈的，但我们认为，在心灵中没有一种邪恶强大到不能被至高的圣言和天主所克服。因为圣言比灵魂中的一切邪恶都更强，而治愈大能居于祂内；祂按照天主的旨意将这医治应用于每个人。万物的终局是邪恶的毁灭，至于它是否会被摧毁得再也无法在任何地方兴起，这超出我们当前的目的。预言中对邪恶的彻底毁灭和每个灵魂的正义恢复有很多隐晦的说法，但对我们当前的目的而言，引用\BibleBook{索福尼亚}{}的以下段落就足够了：\BibleQuote{你们要预备并早起；他们葡萄园的一切剩余都被毁坏了。因此，你们当等候我，直到我兴起作证的日子；因为我决意聚集列国，召集列邦，将我的愤恨，就是我的一切烈怒，倾倒在他们身上；因为全地必被我的妒火吞灭。那时我要转变万民，使他们有纯洁的语言，好叫他们都呼求上主的名，同心合意地事奉祂。从古实河外，我的祈祷者，就是我分散的子民，必向我献上供物。当那日，你必不因你一切所作的事，就是背叛我所行的，而羞愧；因为那时我必从你中间除掉那些因你的骄傲而欢喜的人；你也不再因我的圣山而自高。我必在你中间留下困苦贫寒的百姓，他们必投靠上主的名。以色列的余民必不犯罪，也不说谎言；口中也找不到诡诈的舌头；因为他们必吃喝躺卧，无人惊吓。}我把这事留给那些能够仔细研究整个主题的人，去阐明这预言的含意，尤其要探究这话的意义：\Quote{当全地毁灭时，必按各族的语言转向万民，}正如在语言混乱之前一样。让他们也仔细思考这应许：所有人都将呼求主的名，同心合意地事奉祂；并且一切藐视的辱骂都将被除去，不再有任何不义、虚谎的言语或诡诈的舌头。因此，在我看来，有必要简要地说明这些，而不详述细节，以回应\Person{凯尔苏斯}的言论，即他认为亚细亚、欧罗巴、利比亚的居民，无论是希腊人还是外邦人，达成任何一致都是不可能的。这样的结果对于尚在肉身中的人或许确实不可能，但对于那些已脱离肉身的人则不然。\par

% structure: p0151 source=h2 target=subsection reason=relative-heading-rank; base=h2

\subsection{第七十三章}

接下来，\Person{凯尔苏斯}敦促我们\Quote{全力以赴帮助君王，与他共同维护正义，为他而战；若他要求，则在他麾下作战，或与他一同率军。}对此我们的回答是：我们在必要时确实帮助君王，而且是神圣的帮助，\BibleQuote{穿上天主的全副武装。}我们如此行是顺从宗徒的训示：\BibleQuote{首先，我劝你，要为一切人恳求、祈祷、转求和谢恩，并为君王和一切有权位的人，}一个人越虔敬，就越有效地帮助君王，甚至超过那些出征杀敌的士兵。对于那些要求我们为共和国拿起武器、杀人的信仰仇敌，我们可以回答：\Quote{那些在某些神庙中任司铎的人，以及伺候你们所视为诸神的人，岂不保持双手不沾血，以便用无染人血的手向你们的神献上规定的祭品？即使战时，你们也绝不征召司铎入伍。如果这值得称赞，那么，当别人作战时，这些人作为天主的司铎和仆役，保持双手纯洁，为正义之战和秉公执政的君王祈祷，求天主消灭一切反对正义者的势力，岂不更应如此？}我们藉祈祷战胜所有挑动战争、诱人背誓、扰乱和平的魔鬼，因此我们帮助君王远胜于那些在战场上为他们作战的人。我们确实参与公众事务，在正义的祈祷之外，加上克己和默想，这些教导我们轻视享乐，不被它们牵引。没有人比我们更有效地为君王而战。我们虽不在他麾下作战，即使他要求，我们也不这样做；但我们为他而战，组成一支特殊军队——虔敬之军——藉向天主献上祈祷。\par

% structure: p0153 source=h2 target=subsection reason=relative-heading-rank; base=h2

\subsection{第七十四章}

如果\Person{克耳苏斯}要我们率领军队保卫祖国，那么请他知道，我们也做这事，而且不是为了被人看见或为了虚荣。因为\Quote{在暗中，}在我们心中，有如同司铎为同胞献上的祈祷。基督徒比他人更是祖国的恩人。因为他们培养公民，并向至高天主灌输虔敬；他们将那些在小城中生活良善而堪当的人提拔到一个神圣而属天的城，对这人可以说：\Quote{你在最小的城中忠信，请进入一座大城，}那里\Quote{天主站在众神的集会中，在众神中间审判；}祂将你列在他们中间，只要你不再\Quote{像世人一样死亡，或像首领中的一个跌倒。}\par

% structure: p0155 source=h2 target=subsection reason=relative-heading-rank; base=h2

\subsection{第七十五章}

\Person{克耳苏斯}也敦促我们\Quote{担任国家的政府职务，如果这是为维持法律和支持宗教所必需的。}但我们在每个国家中都承认存在着另一个由天主的圣言所建立的民族组织，我们劝勉那些在言语上有能力且生活无瑕的人去治理教会。我们拒绝那些有统治野心的人；但我们对那些因过度谦逊而不易接受在教会中担任公职的人加以劝勉。那些善治我们的人，是在大王的强力影响之下，我们相信大王是天主圣子、天主圣言。如果那些在教会中治理、被称为神圣民族——即教会——的统治者治理得好，他们是按照天主的命令治理，决不会让自己被世俗政策引入歧途。基督徒不担任公职，并不是为了逃避公众责任，而是为了将自己保留给在教会中更神圣和更必要的服务——为人的救恩。这种服务既是必要的，也是正当的。他们照管所有人——那些在教会内的，使他们日复一日活得更善，以及那些在教会外的，使他们得以充满神圣的言语和虔敬的行为；这样，当他们真正敬拜天主，并以同样的方式培养尽可能多的人时，他们就被天主的圣言和天主的法律所充满，从而通过祂的圣子——圣言、智慧、真理与正义——与至高天主联合；圣子将所有决心在一切事上使自己的生活符合天主的法律的人与天主联合。\par

% structure: p0157 source=h2 target=subsection reason=relative-heading-rank; base=h2

\subsection{第七十六章}

尊敬的\Person{安博罗削}，这就是我们蒙受赐予的能力、顺从您的命令所完成之事的结论。我们在八卷书中囊括了我们认为适宜回应\Person{克耳苏斯}那本他题为\WorkTitle{真言}的书的一切。现在，他的著作和我的回应，究竟哪一部更富真天主的圣神、对祂的虔敬，以及那以正确道理引导人达到最高尚生活的真理，这就有待读者来判断了。然而，您必须知道，\Person{克耳苏斯}曾许诺过一部续作，在其中他承诺为那些愿意接受他意见的人提供实际的生活规则。那么，如果他并未履行写第二卷的诺言，我们以这八卷书来回应他的言论，也就可以满足了。但如果他已开始并完成了那第二卷，请设法取来并寄给我们，这样我们可以按照真理的圣父所赐的能力予以答复，要么推翻其中可能的错误教导，要么排除一切嫉妒，对其中的任何真理予以认可。\par

光荣归于祢，我们的天主；光荣归于祢。\par

\SourceRecord{Translated by Frederick Crombie.；Ante-Nicene Fathers；Vol. 4.；Edited by Alexander Roberts, James Donaldson, and A. Cleveland Coxe.；Buffalo, NY: Christian Literature Publishing Co.,；1885.；<http://www.newadvent.org/fathers/04168.htm>.}
